% ############################################### PREAMBLE AND CONTROL PANEL #

% ########################################################### DOCUMENT CLASS #
\documentclass[  %   draft
             % , demo  % black boxes instead of figures
              , 8pt, twoside, openright]{extbook}

% ########################################################## AND CONTROL PANEL #

% true - hmarginration 1:1, urlcolor - darkgreen, version - for searching
\newif\ifscreen 	   \screentrue
% true - no apendix, no author of grammer, no backmatter ISBN, 
\newif\ifPDF               \PDFtrue
% true - rotate pages of morph verb tables
\newif\ifPDFrotate         \PDFrotatetrue
% true - adds white searchable value !
\newif\ifPDFsearchvalue    \PDFsearchvaluetrue
\newif\ifmakecovers        \makecoverstrue
\newif\ifmakefrontmatter   \makefrontmattertrue
\newif\ifinputintroduction \inputintroductiontrue
\newif\ifinputletters      \inputletterstrue
\newif\ifinputphon         \inputphontrue
\newif\ifinputmorpho       \inputmorphotrue
\newif\ifmakebackmatter    \makebackmattertrue
\newif\ifshowcrop           \showcroptrue
% corrections
\newif\ifC                  \Ctrue

\newif\ifPRINTHOUSE 	   \PRINTHOUSEfalse
\newif\ifPDFSEARCH  	   \PDFSEARCHfalse
\newif\ifPDFPRINT 	   \PDFPRINTfalse
\newif\ifPDFtest	   \PDFtesttrue

%versions
%PRINTHOUSE
\ifPRINTHOUSE
\makecoversfalse
\screenfalse
\PDFfalse
\PDFrotatefalse
\PDFsearchvaluefalse
\showcroptrue
\Ctrue
\fi

%PDFSEARCH
\ifPDFSEARCH
\screentrue
\PDFtrue
\PDFrotatetrue
\PDFsearchvaluetrue
\showcropfalse
\Ctrue
\fi

%PDFPRINT
\ifPDFPRINT
\screenfalse
\PDFtrue
\PDFrotatefalse
\PDFsearchvaluefalse
\showcropfalse
\Ctrue
\fi

%PDFtest
\ifPDFtest
\makecoversfalse
\makefrontmatterfalse
\inputintroductionfalse
\inputletterstrue
\inputphonfalse
\inputmorphofalse
\makebackmatterfalse
\screentrue
\PDFtrue
\PDFrotatefalse
\PDFsearchvaluefalse
\showcropfalse
\Cfalse
\fi

% 1. PDF (for printing)  
% a) searchable value, b) no apendix, c) version= PDF d) no ISBN and no ISBN bar
% e) automatic rotation of landscape pages, f) no author of Icelandic Grammer
% 2. PDF (for viewing offline and searching) + screen (horizontal ratio 1:1) 
% 3. PRINTHOUSE
% a) no searchable value, b) apendix, c) version ISBN d) ISBN bar
% e) no rotation of landscape pages, f)urllinkcolor black
% f) A4 paper dimentions, C5 layout, shows cropping marks

% ########################################################## PAGE GEOMETRIES #

\usepackage%[showframe]
           {geometry}
           

\ifshowcrop
% Basic geometry of document
\geometry
  { headsep    =   \baselineskip
  , textwidth  = 42\baselineskip
  , textheight = 60\baselineskip
  , hmarginratio = \ifscreen 1:1\else 2:3\fi
  , vmarginratio = 2:3
  , bindingoffset = 0cm
  , onecolumn
  , a4paper 
  , layout=c5paper,
  , layouthoffset=\dimexpr(\paperwidth-\csname Gm@layoutwidth\endcsname)/2\relax
   , layoutvoffset=\dimexpr(\paperheight-\csname Gm@layoutheight\endcsname)/2\relax
  , showcrop
  }
\else
\geometry
  { headsep    =   \baselineskip
  , textwidth  = 42\baselineskip
  , textheight = 60\baselineskip
  , hmarginratio = \ifscreen 1:1\else 2:3\fi
  , vmarginratio = 2:3
  , bindingoffset = 0cm
  , onecolumn,
  , c5paper
  }
\fi

\ifshowcrop
% Geometry of the dictionary section
\newcommand\dictionarygeometry{\newgeometry
  { textwidth  = 42\baselineskip
  , textheight = 60\baselineskip
  , hmarginratio = \ifscreen 1:1\else 2:3\fi
  , vmarginratio = 3:2
  , bindingoffset = 0cm
  , twocolumn
  , a4paper 
  , layout=c5paper
  , layouthoffset=\dimexpr(\paperwidth-\csname Gm@layoutwidth\endcsname)/2\relax
   , layoutvoffset=\dimexpr(\paperheight-\csname Gm@layoutheight\endcsname)/2\relax
  , showcrop
  }}
   % Geometry of the cover pages
\newcommand\covergeometry{\newgeometry
  { textwidth  = 42\baselineskip
  , textheight = 60\baselineskip
  , hmarginratio = 1:1
  , vmarginratio = 2:3
  , bindingoffset = 0cm
  , onecolumn
  , layout=c5paper
 , layouthoffset=\dimexpr(\paperwidth-\csname Gm@layoutwidth\endcsname)/2\relax
   , layoutvoffset=\dimexpr(\paperheight-\csname Gm@layoutheight\endcsname)/2\relax
  , showcrop
  }} 
\else
% Geometry of the dictionary section
\newcommand\dictionarygeometry{\newgeometry
  { textwidth  = 42\baselineskip
  , textheight = 60\baselineskip
  , hmarginratio = \ifscreen 1:1\else 2:3\fi
  , vmarginratio = 3:2
  , bindingoffset = 0cm
  , twocolumn,
  , c5paper
  }}
   % Geometry of the cover pages
\newcommand\covergeometry{\newgeometry
  { textwidth  = 42\baselineskip
  , textheight = 60\baselineskip
  , hmarginratio = 1:1
  , vmarginratio = 2:3
  , bindingoffset = 0cm
  , onecolumn
  , c5paper
  }} 
\fi

% ################################################################# LANGUAGES #

\usepackage[icelandic, latin, czech]{babel}
\usepackage{csquotes}


% ################################################################## ENCODING #

\usepackage[utf8]{inputenc}

% Smashing uppercase letters with diacritics
% ------------------ THIS IS EVIL AND SHOULD NEVER BE DONE -------------------
% Don't smash anything if we are on the first run to produce the index
\IfFileExists{figures.idx}{
  % Czech uppercase letters with diacritics:
  \DeclareUnicodeCharacter{00C1}{\strut\smash{\' A}} % Á
  \DeclareUnicodeCharacter{010C}{\strut\smash{\v C}} % Č
  \DeclareUnicodeCharacter{010E}{\strut\smash{\v D}} % Ď
  \DeclareUnicodeCharacter{00C9}{\strut\smash{\' E}} % É
  \DeclareUnicodeCharacter{011A}{\strut\smash{\v E}} % Ě
  \DeclareUnicodeCharacter{00CD}{\strut\smash{\' I}} % Í
  \DeclareUnicodeCharacter{0147}{\strut\smash{\v N}} % Ň
  \DeclareUnicodeCharacter{00D3}{\strut\smash{\' O}} % Ó
  \DeclareUnicodeCharacter{0158}{\strut\smash{\v R}} % Ř
  \DeclareUnicodeCharacter{0160}{\strut\smash{\v S}} % Š
  \DeclareUnicodeCharacter{0164}{\strut\smash{\v T}} % Ť
  \DeclareUnicodeCharacter{00DA}{\strut\smash{\' U}} % Ú
  \DeclareUnicodeCharacter{016E}{\strut\smash{\r U}} % Ů
  \DeclareUnicodeCharacter{00DD}{\strut\smash{\' Y}} % Ý
  \DeclareUnicodeCharacter{017D}{\strut\smash{\v Z}} % Ž
  % Icelandic uppercase letters with diacritics: ÁÉÍÓÚÝ (ÐÞÆ) and
  \DeclareUnicodeCharacter{00D6}{\strut\smash{\" O}} % Ö
}{}
% ------------------ THAT WAS EVIL AND SHOULD NEVER BE DONE ------------------

% Non-breakable space
\DeclareUnicodeCharacter{00A0}{~}

% Non-breakable hyphen
\newcommand*{\nobreakhyphen}{\mbox{-}}
\DeclareUnicodeCharacter{2011}{\nobreakhyphen}


% ############################################################ FONTS, SYMBOLS #

\usepackage[T1]{fontenc}

\usepackage{tgpagella}
\usepackage[scaled=0.88]{helvet}      % relative scale of the two fonts

\def\phvfamily{\fontfamily{phv}\selectfont}

%\usepackage{stmaryrd}
% Let's keep this minimal instead of using the package:
\DeclareSymbolFont{stmry}{U}{stmry}{m}{n}
%\SetSymbolFont{stmry}{bold}{U}{stmry}{b}{n}
\DeclareMathSymbol\shortrightarrow\mathrel{stmry}{"01}
\DeclareMathSymbol\shortuparrow\mathrel{stmry}{"02}

% Methinks these are not necessary
%\usepackage{dingbat}
%\usepackage{manfnt}
%\usepackage{latexsym}

\usepackage{pifont}




% #################################################################### LAYOUT #

% \normalsize should be {8pt}{9.6pt}
\def\HUGE {\fontsize{23.887872pt}{3\baselineskip}\selectfont}
\def\Huge {\fontsize{19.906560pt}{3\baselineskip}\selectfont}
\def\huge {\fontsize{16.588800pt}{3\baselineskip}\selectfont}
\def\LARGE{\fontsize{13.824000pt}{2\baselineskip}\selectfont}
\def\Large{\fontsize{11.520000pt}{2\baselineskip}\selectfont}
\def\large{\fontsize{ 9.600000pt}{2\baselineskip}\selectfont}

% Two columns layout ruler
\setlength\columnsep    {2\baselineskip}
\setlength\columnseprule{0.4pt}

% Necessary for baseline alignment
\topskip=\baselineskip
\raggedbottom
\setlength\parskip{0pt} % it's better to avoid glue

% Temporarily suppress warnings
\hbadness=5000
\vbadness=5000

\setlength\emergencystretch{17pt}

% Allow smaller emergencystretch in several cases
\newenvironment{xtolerant}[2]{%
  \par
  \ifx\empty#1\empty\else\tolerance=#1\relax\fi
  \ifx\empty#2\empty\else\emergencystretch=#2\relax\fi
}{%
  \par
}

% ######################################################### ASSORTED PACKAGES #

\usepackage{tikz}
\usetikzlibrary{calc}

\usepackage{tipa}                      % IPA phonetics

\usepackage{fourier-orns}              % ornaments
\usepackage{amsmath}                   % non-breakable dash
\usepackage{hyphenat}                  % no hyphen in abbreviations

\usepackage{hanging}
\usepackage{fix2col}
\usepackage{fixltx2e}                  % subscript
\usepackage{pagecolor}
\usepackage{afterpage}                 % Used in coloring trick for covers
\usepackage{pdfpages}

\usepackage{multirow}
\usepackage{multicol}
\usepackage{rotating}

\usepackage{etoolbox} % http://ctan.org/pkg/etoolbox
%\usepackage{newunicodechar} 		% new definition of middle dot
\usepackage[ISBN=978-80-260-2385-2,SC0]{ean13isbn}

\usepackage{pgffor}

\usepackage{newunicodechar}

\newunicodechar{·}{\chejnikcdot}
\makeatletter
\newcommand{\chejnikcdot}{%
  \ifnum\lastpenalty=10042
    \kern-0.06667em\relax % dots should always be kerned
  \fi
  \textperiodcentered
  \@ifnextchar|{\kern\chejnikkern\relax}{\penalty10042 }%
}

\newcommand{\closeupdotbar}{%
  \renewcommand\chejnikkern{-0.06667em}%
}
\newcommand\chejnikkern{0pt} % default

% Middle dot + | little bit closer, nicer
%\newunicodechar{·}{\advancedcdot}
%\makeatletter
%\newcommand*\advancedcdot{\textperiodcentered
%  \@ifnextchar|{\hspace{-0.03667em}}{\checkifcdotnext}}
%\newcommand*\checkifcdotnext[2]
%  {\def\tmpa{·}\def\tmpb{#1#2}\ifx\tmpa\tmpb\hspace{-0.06667em}\fi#1#2}
%\makeatother

% ############################################################ BASELINE GRID #

\usepackage{eso-pic}
\usepackage{tikzpagenodes}

% Command to draw a baseline grid
\newcommand\drawbaselinegrid{%
  \begin{tikzpicture}[overlay,remember picture]
    \draw [red!30!white, ultra thin, dashed]
      (0,0) grid [ ystep  = \baselineskip, xstep = \textwidth
                 , shift  = (current page text area.north west)
                 , yshift = -\dp\strutbox
                 ] ++(\textwidth,-\textheight);
    \draw [red!30!white, thin]
      (0,0) grid [ step = \baselineskip, xstep = \textwidth
                 , shift=(current page text area.north)
                 ] ++(0.5\textwidth,-\textheight)
      (0,0) grid [ step = \baselineskip, xstep = \textwidth
                 , shift=(current page text area.north)
                 ] ++(-0.5\textwidth,-\textheight);
    \draw[black!10!white]
      (current page text area.north west)
        rectangle (current page text area.south east);
  \end{tikzpicture}}

% Draw a baseline grid on every page
% \AddToShipoutPicture{\drawbaselinegrid}


% ################################################################## METADATA #

\title{%
  \textbf{Islandsko-český studijní slovník}%
  \thanks{Tato kniha byla vytvořena v \LaTeX{}u pod Ubuntu 14.04.%
    Poděkování patří všem autorům, kteří publikují pod svobodnými licencemi.}}
\author{Aleš Chejn, Ján Zaťko, Jón Gíslason}
\date{říjen 2011}


% ############################################################## BIBLIOGRAPHY #

\usepackage[style=authoryear]{biblatex}
\addbibresource{slovnik.bib}


% ##################################################################### INDEX #

\usepackage{imakeidx}

% Index of authors of photographs
\makeindex
  [ name = figures
  , title = {Seznam autorů fotografií a ilustrací}
  , intoc
  , columnseprule
  , columnsep = \columnsep
  , noautomatic                       % We handle this in make.sh
  ]

% icelandic rules start
\usepackage{filecontents}
\def\mygroup#1{{\bfseries#1}}
\begin{filecontents*}{style.xdy}
;; style.xdy
(markup-letter-group :open-head "~n\mygroup{" :close-head "}")
\end{filecontents*}
% icelandic rules end


% ############################################################### PAGE STYLES #

\usepackage{fancyhdr}

% -------------------------------------- BASIC PAGE STYLE
\fancypagestyle{plain}{
  \fancyhf\relax                       % Clear header/footer
  \renewcommand \headrulewidth {0.0pt} % No header rule
  \renewcommand \footrulewidth {0.0pt} % No footer rule
  \fancyfoot[C]{\thepage}              % Page number in footer, centred
}

% -------------------------------------- INDEX PAGE STYLE  
%\fancypagestyle{indexstyle}{
%  \fancyhf\relax                       % Clear header/footer
%  \renewcommand \headrulewidth {0.4pt} % Thin header rule
%  \renewcommand \footrulewidth {0.0pt} % No footer rule
%  \fancyhead[C]{\thepage}              % Page number in footer, centred
%}

% -------------------------------------- DICTIONARY PAGE STYLE
\fancypagestyle{myheadings}{
  \fancyhf\relax                       % Clear header/footer
  \renewcommand \headrulewidth {0.4pt} % Thin header rule
  \fancyhead[CO,CE]{\thepage}          % Page number in header, centred
  % NOTE: the page numbers will be printed when the dictionary is ready
  \fancyhead[LE,LO]{\phvfamily\bfseries\rightmark}
  \fancyhead[RE,RO]{\phvfamily\bfseries\leftmark}
}

\pagestyle{plain}  

% ##################################################################### LISTS #

\usepackage{expdlist}
\usepackage{enumitem}

\setlist{nosep, topsep=\baselineskip}
\setlist[itemize]{leftmargin=*,label=\scalebox{.6}{\textbullet}}


% #################################################################### FLOATS #

\usepackage{caption}
\captionsetup{labelformat=empty}

% Alter some LaTeX defaults for better treatment of figures:
% See p.105 of "TeX Unbound" for suggested values. 
% See pp. 199-200 of Lamport's "LaTeX" book for details.
% ------------------------------------- Parameters for ALL pages:
\renewcommand \topfraction    {0.9}   %   max fraction of floats at top
\renewcommand \bottomfraction {0.8}   %   max fraction of floats at bottom
% ------------------------------------- Parameters for TEXT (non float) pages:
\setcounter{topnumber}    {1}
\setcounter{bottomnumber} {1}
\setcounter{totalnumber}  {1}         %   2 may work better
\setcounter{dbltopnumber} {2}         %   for 2-column pages
\renewcommand \dbltopfraction   {0.9} %   fit big float above 2-col. text
\renewcommand \textfraction     {0.07}%   allow minimal text w. figs
% ------------------------------------- Parameters for FLOAT (non text) pages:
\renewcommand \floatpagefraction{0.7} %   require fuller float pages
% N.B.: floatpagefraction MUST be less than topfraction !!
\renewcommand{\dblfloatpagefraction}{0.7} % require fuller float pages
% remember to use [htp] or [htpb] for placement

\makeatletter
\setlength{\@fptop}{0pt}
\setlength{\@fpbot}{0pt plus 1fil}
\makeatother

% This kills spacing around floats
\setlength \intextsep        {0pt}
\setlength \floatsep         {0pt}
\setlength \textfloatsep     {0pt}
\setlength \dbltextfloatsep  {0pt}
\setlength \dblfloatsep      {0pt}

% This kills spacing around captions
\setlength \abovecaptionskip {0pt}
\setlength \belowcaptionskip {0pt}


% #################################################################### TABLES #

\usepackage{booktabs}
\usepackage{tabularx}
\usepackage{longtable}
\usepackage{ltxtable}
\usepackage{lscape}

\ifPDFrotate
\else
\makeatletter			       
\renewcommand\PLS@Rotate[1]{}          % this command prevents pdflscape 
\makeatother                           % to rotate pages
\fi

%\LetLtxMacro\oldtabularx\tabularx
%\LetLtxMacro\oldendtabularx\endtabularx
%\renewenvironment{tabularx}[1]
%  {\oldtabularx{\columnwidth}{#2}}
%  {\oldendtabularx}

\newcommand\LTXfw[1]{\LTXtable{\columnwidth}{#1}}

%%% \LetLtxMacro\oldtabularx\tabularx
%%% \LetLtxMacro\oldendtabularx\endtabularx
%%% \newenvironment{longtabularx}[2]
%%%   {\oldtabularx{#1}{#2}}
%%%   {\oldendtabularx}

% ######################################################### czech babel vs cmidrule
% ###################################### http://tex.stackexchange.com/a/112001/8380
\begingroup
  \makeatletter
  \catcode`\-=\active
  \AtBeginDocument{
  \def\@@@cmidrule[#1-#2]#3#4{\global\@cmidla#1\relax
    \global\advance\@cmidla\m@ne
    \ifnum\@cmidla>0\global\let\@gtempa\@cmidrulea\else
    \global\let\@gtempa\@cmidruleb\fi
    \global\@cmidlb#2\relax
    \global\advance\@cmidlb-\@cmidla
    \global\@thisrulewidth=#3
    \@setrulekerning{#4}
    \ifnum\@lastruleclass=\z@\vskip \aboverulesep\fi
    \ifnum0=`{\fi}\@gtempa
    \noalign{\ifnum0=`}\fi\futurenonspacelet\@tempa\@xcmidrule}
  }
\endgroup

% macro that supress cmidrule at page brake
% http://tex.stackexchange.com/questions/110548/using-longtable-with-cmidrule-page-breaking-issues
\makeatletter
\def\LT@output{%
  \ifnum\outputpenalty <-\@Mi
    \ifnum\outputpenalty > -\LT@end@pen
      \LT@err{floats and marginpars not allowed in a longtable}\@ehc
    \else
      \setbox\z@\vbox{\unvbox\@cclv}%
      \ifdim \ht\LT@lastfoot>\ht\LT@foot
        \dimen@\pagegoal
        \advance\dimen@-\ht\LT@lastfoot
        \ifdim\dimen@<\ht\z@
          \setbox\@cclv\vbox{\unvbox\z@
\copy\LT@foot\vss}%
          \@makecol
          \@outputpage
          \setbox\z@\vbox{\box\LT@head}%
        \fi
      \fi
      \global\@colroom\@colht
      \global\vsize\@colht
      \vbox
        {\unvbox\z@
\box\ifvoid\LT@lastfoot\LT@foot\else\LT@lastfoot\fi}%
    \fi
  \else
    \setbox\@cclv\vbox{\unvbox\@cclv

\setbox0\lastbox
\ifdim\ht0=0.29999pt \unskip % hope that was a cline we threw away
\else
\nointerlineskip
\box0     % put it back, whatever it was
\fi

\copy\LT@foot\vss}%
    \@makecol
    \@outputpage
      \global\vsize\@colroom
    \copy\LT@head\nobreak
  \fi}

\makeatother

% These values ensure that for top, mid and bottom rules
% space above + rule width + space below = baselineskip
\setlength \abovetopsep    {0.43\baselineskip}
\setlength \heavyrulewidth {0.10\baselineskip}
\setlength \belowbottomsep {0.43\baselineskip}
\setlength \aboverulesep   {0.47\baselineskip}
\setlength \lightrulewidth {0.06\baselineskip}
\setlength \belowrulesep   {0.47\baselineskip}

%\usepackage{letltxmacro}

%%% \LetLtxMacro\oldtabular\tabular
%%% \LetLtxMacro\oldendtabular\endtabular
%%% \renewenvironment{tabular}[2][t]
%%%   {\oldtabular[#1]{#2}}
%%%   {\oldendtabular}

%%% \LetLtxMacro\oldtable\table
%%% \LetLtxMacro\oldendtable\endtable
%%% \renewenvironment{table}[1][h]
%%%   {\leavevmode\oldtable[#1]\vspace{\dp\strutbox}}
%%%   {\vspace{-\dp\strutbox}\oldendtable}
%%%   % This aligns a [t] aligned tab-envo to baseline grid


%%% % declention and conjugation tables
%%% \usepackage{floatrow}
%%% \DeclareFloatFont{footnotesize}{\footnotesize}
%%% % "scriptsize" is defined by floatrow, "tiny" not
%%% \floatsetup[table]{font=footnotesize}

% ################################################################### FIGURES #

\usepackage{adjustbox} % loads also graphicx

\graphicspath{%
  {/home/chejnik/Dokumenty/hvalur.org/images/biolib/full/}%
  {/home/chejnik/Dokumenty/hvalur.org/images/uploaded_files/}}

\newlength\dicfigheight
\newlength\dicfigskip

\newcommand\dicFigure[3]{ % {filename}{caption}{index entry}
  \setlength\fboxsep{0pt}\setlength\fboxrule{0.5pt}
  \def\dicfigbox{\fbox{%
    \adjincludegraphics [ max height = 0.8\columnwidth
                        , max width  = 0.8\columnwidth ] {#1}}}
  \setlength\abovecaptionskip {\dp\strutbox}
  \setlength\belowcaptionskip {\ht\strutbox} 
  \setlength\dicfigheight {\heightof\dicfigbox+\fboxrule}
  \setlength\dicfigskip
    {\baselineskip*\numexpr1+\dicfigheight/\baselineskip\relax-\dicfigheight}
  \begin{figure}[ht]
    \centering\vspace*{\dicfigskip}\dicfigbox\caption{#2}\index[figures]{#3}
  \end{figure}
}


% #################################################################### COLORS #

\usepackage{color}

% http://www.colorschemer.com/online.html
% main color - grammar color
% #66023C
\definecolor {darkgreen}          {rgb} {0.40, 0.01, 0.24}

\definecolor {royalazure}         {rgb} {0.00, 0.22, 0.66}
\definecolor {brown}              {rgb} {0.40, 0.01, 0.24}

% COLORS FOR THUMB INDEXES
\definecolor {babyblueeyes}       {rgb} {0.63, 0.79, 0.95}
\definecolor {unitednationsblue}  {rgb} {0.36, 0.57, 0.90}
\definecolor {blue(ryb)}          {rgb} {0.01, 0.28, 1.00}
\definecolor {darkblue}           {rgb} {0.00, 0.00, 0.55}
\definecolor {screamingreen}      {rgb} {0.46, 1.00, 0.44}
\definecolor {limegreen}          {rgb} {0.20, 0.80, 0.20}
\definecolor {islamicgreen}       {rgb} {0.00, 0.56, 0.00}
\definecolor {upforestgreen}      {rgb} {0.00, 0.27, 0.13}
\definecolor {icterine}           {rgb} {0.99, 0.97, 0.37}
\definecolor {orange(colorwheel)} {rgb} {1.00, 0.50, 0.00}
\definecolor {orange-red}         {rgb} {1.00, 0.27, 0.00}
\definecolor {oucrimsonred}       {rgb} {0.60, 0.00, 0.00}
\definecolor {cottoncandy}        {rgb} {1.00, 0.74, 0.85}
\definecolor {orchid}             {rgb} {0.85, 0.44, 0.84}
\definecolor {vividcerise}        {rgb} {0.85, 0.11, 0.51}
\definecolor {patriarch}          {rgb} {0.50, 0.00, 0.50}

\definecolor {title}              {RGB} {  16,   13,   32}
\definecolor {golden}             {RGB} { 241,  184,   45}
\definecolor {freqstar}           {RGB} {  91,    3,   99}
\definecolor {newdev}             {RGB} {   3,   11,   99}

\definecolor {color1}             {RGB} { 182,   86,    0}
\definecolor {color2}             {RGB} { 143,    9,    6}
\definecolor {color3}             {RGB} {   3,   23,  118}
\definecolor {color4}             {RGB} {   0,   82,  168}
\definecolor {color5}             {RGB} {   0,   85,  142}
\definecolor {color6}             {RGB} {   0,  115,  162}
\definecolor {color7}             {RGB} {  34,  146,  186}
\definecolor {color8}             {RGB} {  40,  159,  153}
\definecolor {color9}             {RGB} {   0,  125,  111}
\definecolor {color10}            {RGB} {   4,  107,   60}
\definecolor {color11}            {RGB} {  71,  134,   81}
\definecolor {color12}            {RGB} { 109,  134,   42}
\definecolor {color13}            {RGB} { 205,  194,   18}
\definecolor {color14}            {RGB} { 204,  162,   24}

%\definecolor {color1}             {RGB} { 239,  222,  205}
%\definecolor {color2}             {RGB} { 251,  206,  177}
%\definecolor {color3}             {RGB} { 251,  185,  142}
%\definecolor {color4}             {RGB} { 248,  166,  112}
%\definecolor {color5}             {RGB} { 216,  231,  246}
%\definecolor {color6}             {RGB} { 186,  215,  244}
%\definecolor {color7}             {RGB} { 152,  197,  243}
%\definecolor {color8}             {RGB} { 118,  181,  245}
%\definecolor {color9}             {RGB} { 181,  249,  185}
%\definecolor {color10}            {RGB} { 160,  248,  167}
%\definecolor {color11}            {RGB} { 128,  247,  136}
%\definecolor {color12}            {RGB} {  87,  248,   98}
%\definecolor {color13}            {RGB} { 250,  241,  138}
%\definecolor {color14}            {RGB} { 250,  240,  111}
%\definecolor {color15}            {RGB} { 244,  239,   85}
%\definecolor {color16}            {RGB} { 245,  239,   57}

\definecolor {lightgrey}          {RGB} { 105,  105,  105}
\definecolor {darkgreen_real}     {rgb} {0.00, 0.50, 0.00}


% ############################################################# THUMB INDEXES #

% Thumb indexes' colors
%\newcommand\BoxColor{%
%  \ifcase\theletternum darkgreen!30\or babyblueeyes\or unitednationsblue\or blue(ryb)\or screamingreen\or limegreen\or islamicgreen\or upforestgreen\or icterine\or orange(colorwheel)\or orange-red%
%  \or oucrimsonred\or cottoncandy\or orchid\or vividcerise\or patriarch\or babyblueeyes\or unitednationsblue\or blue(ryb)\or screamingreen\or limegreen\or islamicgreen\or upforestgreen\or icterine\or orange(colorwheel)\or orange-red%
%  \or oucrimsonred\or cottoncandy\or orchid\or vividcerise\or patriarch\else darkgreen!30\fi}
  
% Thumb indexes' colors
%\newcommand\BoxColor{%
%  \ifcase\theletternum darkgreen!30\or color1\or color2\or color3\or color4\or color5\or color6\or color7\or color8\or color9\or color10%
%  \or color11\or color12\or color13\or color14\or color1\or color2\or color3\or color4\or color5\or color6\or color7\or color8\or color9%
%  \or color10\or color11\or color12\or color13\or color14\or color1\or color2\or color3\or color4\or color5\or color6\else color7\or color8\or color9\or color10\or color11\or color12\fi}

%\def\BoxColor#1{darkgreen!\the\numexpr103-#1-#1-#1\relax!blue}

\newcommand\BoxColor[1]{%
\ifcase#1 darkgreen!30\or color13\or color14\or color1\or color2\or color3\or color4\or color5\or color6\or color7\or color8\or color9\or color10%
\or color11\or color12\or color13\or color14\or color1\or color2\or color3\or color4\or color5\or color6\or color7\or color8\or color9%
\or color10\or color11\or color12\or color13\or color14\or color1\or color2\or color3\or color4\or color5\or color6\else color7\or color8\or color9\or color10\or color11\or color12\fi}

\ifshowcrop
%%%%%%%%%%%%%%%%%%%%%%%%%%%%%%%%%%%%%%%%%%%%%%%%%%%%%%%%%%%%%%%%%%%%%%%%%%%%%%%%

\newcounter{letternum}
\newcounter{lettersum}                 \setcounter{lettersum}{34}
% these margins separate top thumb from top of page and last thumb from
% bottom of page. BEWARE, it's from top of the (cut out) page, not the folio!
\newlength{\thumbtopmargin}            \setlength{\thumbtopmargin}{1cm}
\newlength{\thumbbottommargin}         \setlength{\thumbbottommargin}{1cm}
\newlength{\thumbwidth}                \setlength{\thumbwidth}{2.20cm}
\newlength{\thumbheight}
\pgfmathsetlength{\thumbheight}%
  {(\csname Gm@layoutheight\endcsname-\thumbtopmargin-\thumbbottommargin)/\value{lettersum}}
% inner xsep of letter from box margin
\newlength{\thumbsep}                  \setlength{\thumbsep}{.30cm}
% overlap over cutting line
\newlength{\thumboverlap}              \setlength{\thumboverlap}{3cm}

\tikzset{
  thumb/.style={
    text=white,
    outer sep=0pt,
    inner xsep=\thumbsep,
    minimum height=\thumbheight,
    text width=\thumbwidth-2\thumbsep,
    font=\sffamily\bfseries,},
  thumb east/.style={thumb,
    xshift=\dimexpr(+\thumboverlap-\paperwidth+\csname Gm@layoutwidth\endcsname)/2\relax,
    align=left,anchor=north east,},
  thumb west/.style={thumb,
    xshift=\dimexpr(-\thumboverlap+\paperwidth-\csname Gm@layoutwidth\endcsname)/2\relax,
    align=right,anchor=north west,},
}

\def\thumbnew{}
\def\thumbold{}
\usepackage{everypage}
\AddEverypageHook{\if\relax\thumbnew\relax\xdef\thumbnew{\thumbold}\fi}

\def\ethumbs#1,#2\relax%
  {\if\relax#1\relax\else\eventhumb{#1}\fi%
   \if\relax#2\relax\else\ethumbs#2\relax\fi%
   \gdef\thumbnew{}%
   \gdef\thumbold{#1,}}

\def\othumbs#1,#2\relax%
  {\if\relax#1\relax\else\oddthumb{#1}\fi%
   \if\relax#2\relax\else\othumbs#2\relax\fi%
   \gdef\thumbold{#1,}%
   \gdef\thumbnew{}}

\newcommand{\drawthumb}[2]{%
  \begin{tikzpicture}[remember picture, overlay]
    \node [thumb #2, fill=\BoxColor{#1}]
       at ($(current page.north #2)-%
            (0,\csname Gm@layoutvoffset\endcsname+\thumbtopmargin-\thumbheight+%
               #1*\thumbheight)$) {\csname Let#1\endcsname};
   \end{tikzpicture}}



%%%%%%%%%%%%%%%%%%%%%%%%%%%%%%%%%%%%%%%%%%%%%%%%%%%%%%%%%%%%%%%%%%%%%%%%%%%%%%%%
\else
% THUMB INDEXES
% new counter to hold the current number of the letter to determine the vertical position
\newcounter{letternum}
% newcounter for the sum of all letters to get the right height of a box
\newcounter{lettersum}
\setcounter{lettersum}{34}
% some margin settings
\newlength{\thumbtopmargin}
\setlength{\thumbtopmargin}{\ifshowcrop 3cm\else 1.5cm\fi} %1cm
\newlength{\thumbbottommargin}
\setlength{\thumbbottommargin}{\ifshowcrop 6cm\else 2cm\fi} %2.5cm
% calculate the box height by dividing the page height
\newlength{\thumbheight}
\pgfmathsetlength{\thumbheight}{%
(\paperheight-\thumbtopmargin-\thumbbottommargin)%
/%
\value{lettersum}
}
% box width
\newlength{\thumbwidth}
\setlength{\thumbwidth}{\ifshowcrop 2cm\else 0.5cm\fi} %0.5cm
% style the boxes
\tikzset{
thumb/.style={
   text=white,
   minimum height=\thumbheight,
   text width=\thumbwidth,
   outer sep=0pt,
   font=\sffamily\bfseries,
 }
 }

\def\thumbnew{}
\def\thumbold{}
\usepackage{everypage}
\AddEverypageHook{\if\relax\thumbnew\relax\xdef\thumbnew{\thumbold}\fi}

\def\ethumbs#1,#2\relax{\if\relax#1\relax\else\eventhumb{#1}\fi%
                        \if\relax#2\relax\else\ethumbs#2\relax\fi%
                        \gdef\thumbnew{}%
                        \gdef\thumbold{#1,}%
}

\def\othumbs#1,#2\relax{\if\relax#1\relax\else\oddthumb{#1}\fi%
                        \if\relax#2\relax\else\othumbs#2\relax\fi%
                        \gdef\thumbold{#1,}%
                        \gdef\thumbnew{}%
}

\newcommand{\drawthumb}[2]{%
  % see pgfmanual.pdf for more information about this part
  \begin{tikzpicture}[remember picture, overlay]
    \node [thumb, fill=\BoxColor{#1}, text centered, anchor=north #2]
       at ($(current page.north #2)-%
%            (0,\thumbtopmargin+\value{letternum}*\thumbheight)%
            (0,\thumbtopmargin+#1*\thumbheight)$) {\csname Let#1\endcsname};
   \end{tikzpicture}}
\fi

\newcommand{\oddthumb} [1]{\drawthumb{#1}{west}}
\newcommand{\eventhumb}[1]{\drawthumb{#1}{east}}

\newcommand{\lettergroup}[1]%
  {\refstepcounter{letternum}%
   \expandafter\gdef\csname Let\theletternum\endcsname{#1}%
   \xdef\thumbnew{\theletternum,\thumbnew}%
   \fancyhead[LO]{\phvfamily\bfseries\rightmark%
     \expandafter\ethumbs\thumbnew\relax\relax\relax}%
   \fancyhead[RE]{\phvfamily\bfseries\leftmark%
     \expandafter\othumbs\thumbnew\relax\relax\relax}}
     
% ############################################################## LOCALIZATION #

% renames the index name
%\addto\captionsczech{%
%  \renewcommand\indexname{{Seznam autorů fotografií a ilustrací}}}

% renames the contents name
\addto\captionsczech{%
  \renewcommand\contentsname{Obsah}}

\renewcommand \contentsname {Obsah}
\renewcommand \chaptername  {Kapitola}


% ################################################################ SECTIONING #
\usepackage[explicit]{titlesec}

\setlength{\beforetitleunit}{\baselineskip}
\setlength{\aftertitleunit} {\baselineskip}

\titlespacing*{\chapter}   {0em}{*0}{*3} % Remembere there's a topskip too
\titlespacing*{\section}   {0em}{*1}{*2}
\titlespacing*{\subsection}{0em}{*1}{*2}

\newcommand\ghostdrop[2]{%
  \raisebox{-#1\baselineskip}[0pt][0pt]{#2}}

\titleformat {\chapter} {} {} {0pt} % {#1}
  {\ghostdrop{1.0}{\parbox[t]{\columnwidth}
     {\LARGE\bfseries\centering\MakeUppercase{#1}}}}

\def\longchapterskip{\vspace{2\baselineskip}}
\def\longsectionskip{\vspace{2\baselineskip}}

\titleformat {\section} {\bfseries} {} {0pt}
  {\ghostdrop{1.0}{\parbox[t]{\columnwidth}
     {\Large\bfseries\centering#1}}}

\titleformat {\subsection} {\bfseries} {} {0pt}
  {\ghostdrop{1.0}{\hspace{1em}\parbox[t]{\columnwidth-1em}
     {\large\bfseries#1}}}


\setlength\multicolsep{0pt}
%\BeforeBeginEnvironment{multicols}{\vspace{-\dp\strutbox}}


% ######################################################## TABLE OF CONTENTS #
\usepackage{titletoc}

\titlecontents {chapter} [0pt] {\addvspace\baselineskip\bfseries}
  {} {} {\titlerule*[2\baselineskip]{}\contentspage}

\titlecontents {section} [\baselineskip] {} {} {}
  {\titlerule*[1\baselineskip]{.}\contentspage}


% #################################################### ASSORTED CUSTOM MACROS #

\def\textIS#1{\foreignlanguage{icelandic}{#1}}
\def\textCS#1{\foreignlanguage{czech}{#1}}
\def\textLA#1{\foreignlanguage{latin}{#1}}

\usepackage{accsupp}

\makeatletter                                                                                                          
  \def\ACCSUPP@bdc{\special{pdf:code \ACCSUPP@span BDC}}                                                               
  \def\ACCSUPP@emc{\special{pdf:code EMC}}                                                                             
\makeatother

\usepackage{pgffor}

\newcommand\blspace[1][1]{\vspace{#1\baselineskip}}

\usepackage{xparse}


\newcommand \dicsymFrequent  {\raisebox{-.2ex}{\color{darkgreen}\ding{167}}}
\newcommand \dicsymSee       {$\shortrightarrow$}
\newcommand \dicsymCompare   {$\shortuparrow$}
\newcommand \dicsymIdiom     {$\diamondsuit$}
\newcommand \dicsymExampleIS {$\triangleright$}
\newcommand \dicsymExampleCS {\guilsinglright}
\newcommand \dicsymProverb   {{\color{newdev}\ding{96}}}
\newcommand \dicsymPhoto
  {\includegraphics[keepaspectratio,height=0.65em]{photo_icon}}

%\NewDocumentCommand\dicTerm{mo}
 % {{\phvfamily\bfseries\textIS{#1}}\IfValueT{#2}{, \phvfamily\bfseries\textIS{#2}}}

\NewDocumentCommand\dicEntry{o}
  {\par\hangpara{0.25\baselineskip}{1}%
   \IfValueT{#1}{\markboth{#1}{#1}}%
    \ifPDFsearchvalue
    \settowidth{\eyja}{#1}%
  \makebox[\eyja]{\color{white}#1}\hspace{-\eyja}%
  \fi
   \ignorespaces}
   
   \NewDocumentCommand\dicEntryGuide{o}
  {\par\hangpara{0.29\baselineskip}{1}%
   \IfValueT{#1}{\markboth{#1}{#1}}%
    \ifPDFsearchvalue
    \settowidth{\eyja}{#1}%
  \makebox[\eyja]{\color{white}#1}\hspace{-\eyja}%
  \fi
   \ignorespaces}
   
\newcommand\dicSubEntry
  {\\\hspace*{-\hangindent}}

\NewDocumentCommand\textAlt{mm}
  {\BeginAccSupp{method=pdfstringdef,ActualText={#2}}#1\EndAccSupp{}}

\newcounter{terms}

\NewDocumentCommand\dicStyleIntTerm{m}{\textIS{{\phvfamily\bfseries#1}}}
\NewDocumentCommand\dicStyleExtTerm{m}{\textIS{\textit{#1}}}

   \NewDocumentCommand\dicTerm{mo}
  {{\closeupdotbar \phvfamily\bfseries\textIS{#1}}\IfValueT{#2}{, \phvfamily\bfseries\textIS{#2}}}

  \newcommand\dicTermList[1]
  {\foreach \t [count=\n] in {#1}{\ifnum\n=1\else, \fi\dicTerm{\t}}}
  
\def\dicIPA#1{\textipa{[#1]}}

\NewDocumentCommand\dicPos{smo}
  {{\color{darkgreen}\textbf{\small#2}%
   \IfValueT{#3}{\nobreak\textsubscript{#3}}}}

%\NewDocumentCommand\dicFlx{smo}
 % {{\color{darkgreen}\footnotesize#1\IfValueT{#2}{\textsubscript{#2}}}}

  \NewDocumentCommand\dicFlx{smo}
  {{\color{darkgreen}\footnotesize#2%
   \IfValueT{#3}{\nobreak\textsubscript{#3}}}}
  
\newcommand\dicExampleCS[1]{\dicsymExampleCS~\textit{#1}}
\newcommand\dicExampleIS[1]{\dicsymExampleIS~\textIS{\textit{#1}}}

\NewDocumentCommand\dicLink{sm}
  {\IfBooleanF#1{\dicsymSee\nobreak\hspace{.08667em plus .03333em}}\nobreak\hspace{.03667em}\dicTerm{#2}}
\NewDocumentCommand\dicSynonym{sm}
  {(\IfBooleanF#1{\dicsymSee\nobreak\hspace{.08667em plus .03333em}}\nobreak\hspace{.03667em}\textIS{\textit{#2}})}
\NewDocumentCommand\dicAntonym{sm}
  {(\IfBooleanF#1{\dicsymCompare\nobreak\hspace{.08667em plus .03333em}}\nobreak\hspace{.03667em}\textIS{\textit{#2}})}
\NewDocumentCommand\dicCompare{sm}
  {(\IfBooleanF#1{\dicsymCompare\nobreak\hspace{.08667em plus .03333em}}\nobreak\hspace{.03667em}\textIS{\textit{#2}})}


\newcommand\dicPhraseIS[1]{\textIS{\textbf{#1}}}

\NewDocumentCommand\dicIdiom{mo}
  {\dicSubEntry{\color{newdev}\dicStyleIntTerm{#1} %
     \IfValueTF{#2}{+ \dicStyleIntTerm{#2} \markboth{#1 + #2}{#1 + #2}}{\markboth{#1}{#1}}}\dicsymIdiom}
     
\newcommand\dicProverb
  {\dicSubEntry\dicsymProverb}



\newcommand\dicLangCat[1]{{\footnotesize#1}}
\newcommand\dicFieldCat[1]{{\footnotesize#1}}


\newcommand\dicDirectTranslationCS[1]{#1}
\newcommand\dicIndirectTranslationCS[1]{{\footnotesize#1}}

\usepackage{placeins}

\newcommand*{\dicLetter}[2]{%
  \FloatBarrier\newpage\phantomsection\pdfbookmark[1]{#1}{#2}%
  \noindent\parbox[b][9\baselineskip][c]{\columnwidth}
    {\centering\HUGE\strut\smash{\MakeUppercase{#1}~\MakeLowercase{#1}}}%
      %\par\xdef\tpd{\the\prevdepth}
      %\par\prevdepth\tpd%
  \expandafter\lettergroup{\MakeLowercase{#1}}}

  
%tabulky
\newcommand{\specialcell}[2][l]{%
\begin{tabular}[#1]{@{}l@{}}#2\end{tabular}}


%\newcommand{\devisionguide}[1]{\hspace*{0.3cm}{{{{\foreignlanguage{icelandic}{\color{newdev}{\phvfamily{\textbf{#1}}}}}}}}}
%\newcommand{\devisionguideindent}[1]{\hspace*{-0.1cm}{{{{\foreignlanguage{icelandic}{\color{newdev}{\phvfamily{\textbf{#1}}}}}}}}}

%\newcommand\n{$n$\nobreakdash-\hspace{0pt}} %nonbreakable dash in grammar endings
%\def \nobreakseq {\nobreak \hskip 0pt \hbox} %nolinebreak in case like (-s, -)

% rule line
\newcommand{\HRule}{\rule{\linewidth}{0.1mm}}

% width of pdf found box
\newlength\eyja

%\newcommand\strutsmash[1]{\strut\smash{#1}}


% Dictionary symbols


% Command synonyms
%\let\freqstar\dicsymFrequent
%\let\photoicon\dicsymPhoto

\newcommand{\vstretch}[1]{\vspace*{\stretch{#1}}}

\newcommand*{\addthin}{\hskip0.16667em\relax}
\newcommand*{\addthinS}{\hskip0.06667em\relax}
\newcommand*{\addthinM}{\hskip0.10007em\relax}
\newcommand*{\addthinSS}{\hskip0.00007em\relax}
% #################################################################### COVERS #

\def\dictnameCZ{islandsko-český studijní slovník}
\def\dictnameIS{íslensk-tékknesk stúdentaorðabók}

\newcommand*{\frontcoverimages}{%
  ds_image_blaberjalyng_0_2.jpg,
  ds_image_lundi_0_1.jpg,
  ds_image_hestur_0_1.jpg,
  ds_image_hrafnafifa_0_2.jpg,
  ds_image_baldursbra_0_1.jpg,
  20948.jpg,
  ds_image_hvalur_0_1.jpg,
  ds_image_blodberg_0_1.jpg}

\newcommand*{\backcoverimages}{%
  ds_image_spoi_0_1.jpg,
  ds_image_mosalyng_0_1.jpg,
  ds_image_blaklukkulyng_0_1.jpg,
  20787.jpg,
  ds_image_ufsi_0_2.jpg,
  %ds_image_gullbra_0_1.jpg,
  ds_image_vatnsberi_0_1.jpg,
  ds_image_islandsfifill_0_2.jpg,
  21137.jpg}

\newcommand\makecoverwith[1]{
 \pagecolor{title}\afterpage\nopagecolor
 \begin{center}
 \vspace*{0.3cm}%{\fill}
 {\Huge\scshape\color{white}%
   \dictnameCZ\\\dictnameIS}
 \vspace*{1em}                        % I think this is necessary
 \begin{figure}[ht]\centering%
   \setlength\fboxsep{0pt}\setlength\fboxrule{0.5pt}%
   \foreach \file in #1{
     \fbox{\includegraphics[width=5.5cm]{\file}}}
 \end{figure}
 \vspace*{\fill}
 \end{center}
}


% ###################################################### HYPERREF & PDF INFO #

% unicode neccessary so that national characters in hypersetup appear ok
\usepackage[pdftex, unicode, hyperfootnotes=false]{hyperref}

% hyperlinks in black
\makeatletter
\let\Hy@linktoc\Hy@linktoc@none
\makeatother

\hypersetup
  { pdftitle={Islandsko-český studijní slovník}
  , pdfauthor={Aleš Chejn, Jón Gíslason, Ján Zaťko}
  , pdfsubject={Islandsko-český studijní slovník}
  , pdfkeywords={Islandsko-český studijní slovník%
               , islandština, čeština, slovník}
  , bookmarks=true
  , colorlinks=true
  , citecolor=black
  , urlcolor=\ifscreen darkgreen\else black\fi
  , linkcolor=black}


% ################################################################# MICROTYPE #

\usepackage[final]{microtype}

% ############################################################### HYPHENATION #

% this allows to include brackets inside hyphenation rules f.e (pro)-dis-ku-to-vat
\lccode`\(`\(
\lccode`\)`\)
\lccode`\[`\[
\lccode`\]`\]

\begin{hyphenrules}{icelandic}
\hyphenation{
  frum-semja
  burm-neskur
  ó-breytan-le-gur
  pla-tó-nskt
  skatta-fram-tal
  skrúf-gang-ur
  strand-rauð-viður
  prent-sver-ta
  grun-nur
  uppi-skrop-pa
  ger-vi-mennska
  sö-lu-stað-ur
  sý-na-gó-ga
  mark-mið
  trau-st-ur
  ís-kyggi-leg-ur
  brjó-sta-höld
  kné-krjú-pa
  nel-li-ka
  kaffi-kvörn
  margs-konar
  skrúf-gang-ur
  uppi-skrop-pa
  veð-rað-ur
  vest-ur-átt
  mis-mun-andi
  Bar-ba-do-si
  tung-u-mál
  au-ga-brún
  há-sker-pu-sjón-varp
  rúnn-stykki
  ger-vi-tungl
  yfir-borðs-menn-ska
  gæ-fu-mun-ur
  hlí-fi-skil-di
  á-vís-un
  ó-um-flýj-an-leg-ur
  san-mar-ín-sk-ur
  gæ-ða-stað-all
}
\end{hyphenrules}
\begin{hyphenrules}{czech}
\hyphenation{
  teo-rie
  mi-liardy
  (z)-od-po-věd-nost
  tan-tié-ma
  alian-ce
  pe-dia-tr(ička)
  pa-cienty
  zá-chra-nář(ka)
  dia-man-to-vá
  poe-zie
  ma-te-riál-ní
  egoi-s-ta
  egoist-ka
  nai-vi-ta
  ne-upřím-nost
  pě-ti-úhel-ník
  še-s-ti-úhel-ník
  za-šmo-dr-chat
  ar-cheo-lo-gic-ký
  ne-omlu-vi-tel-ný
  (vy)-pro-du-ko-vat
  dia-gno-s-ti-ko-vat
  teo-lož-ka
  ateist-ka
  oceá-nu
  im-pe-ria-li-s-ta
  (vy)-zdvih-nout
  ne-úměr-ný
  kon-ti-nuál-ní
  neu-t-rum
  ne-opráv-ně-ně
  vig-vam
  pro-zaic-ký
  geo-lož-ka
  seiz-mic-ká
  teo-rém
  olean-dr
  poe-ta
  (vy)-fo-to-gra-fo-vat
  (vy)-ven-ti-lo-vat
  pneu-mo-nie
  Lu-cem-bur-čan-(ka)
  ma-te-riály
  am-bi-cióz-ní
  Myan-ma(r)
  Myan-ma-řan-(ka)
  mol-dav-šti-na
  si-tua-cí
  ne-ukáz-ně-ný
  he-te-ro-trof-ní
  ne-omlu-vi-tel-ná
  ne-upřím-ný
  dvou-bliz-ný
  ne-úmy-sl-ný
  ne-ustá-lý-mi
  ne-ohle-du-pl-ný
  ne-osob-ní
  ne-oso-le-ný
  ne-uvě-ři-tel-ný
  ne-ohro-že-ný
  ne-oby-čej-ný
  ne-útul-ný
  ne-únav-ný
  ne-ože-nit
  po-ukáz-ka
  za-úto-čit
  aero-dy-na-mic-ký
  hru-bián-ství
  (pro)-di-s-ku-to-vat
  (vy)-ča-lou-nit
  (na)-oč-ko-vá-ní
  Bhú-tán-ka
  (za)-šně-ro-vat
  soud-ní
  re-vo-lu-cio-nář(ka)
  á-zer-báj-d-žán
  čtyř-hran-ný
  ze-m-dle-lý
  klou-zač-ka
  Vest-mann-ské
  tří-stran-ný
  (za)-bu-rá-ce-ní
  os-tud-ný
  pro-du-cent-(ka)
  chvá-s-tal-(ka)
  apar-theid
  asis-ten-ce
  asis-tent-(ka)
  (u)-či-nit
  (u)-za-mknout
  past-vi-na
  (za)-zá-lo-ho-vat
  Af-ri-čan-(ka)
  (za)-an-ga-žo-vá-ní
  se-rióz-ní
  opo-nent-(ka)
  pro-tiv-ník
  zvlášt-ní
  (po)-div-ný
  ho-mo-se-xuál-ní
  ko-alič-ní
  si-tua-ce
  si-tua-cí
  egoi-s-mus
  spe-ciál-ní
  spe-cia-li-zo-va-ný
  spe-cia-li-zo-va-né
  sobě-stač-ný
  sa-mo-zvo-le-ný
  sa-mo-čin-ně
 chvá-s-ta-vost
 klient-(ka)
 ku-rióz-ní
 (na)-šrou-bo-vat
 (za)-šrou-bo-vat
 krea-ti-vi-ta
 pneu-ma-ti-ka
 toa-le-ta
 so-cia-list-ka
 (vy)-dez-in-fi-ko-va-ný
 vše-uměl
 vše-uměl-(kyně)
 (pře)-škrt-nout
 (vy)-svlék-nout
 vy-učo-va-cí
 vy-učo-va-cích
 (vy)-chla-s-tat
 (vy)-te-to-vat
 pa-cient-(ka)
 ne-opod-s-tat-ně-ný
 de-se-ti-úhel-ník
 pe-rio-di-kum
 (za)-mum-lá-ní
 ne-smě-lá
 ne-dů-vě-ři-vá
 tee-na-ger-(ka)
 (z)-de-ge-ne-ro-vá-ní
 mo-ci-chti-vý
 bez-úhon-nost
 zne-uctít
 vo-dě-odol-ný
 we-bo-vý
 web-de-sign
 in-že-nýr-(ka)
 zá-pa-do-is-land-ský
 klient-(ka)
 za-uja-tý
 vy-uží-va-né
 vy-uží-vat
 (za)-re-gi-s-t-ro-vat
 (z)-kom-pri-mo-va-ný
 na-cio-na-li-s-mus
 ma-k-ro-eko-no-mic-ká
 bez-bo-lest-ný
 viol-ka
 troj-úhel-ník
 (za)-bu-rá-ce-ní
 ne-ústup-nost
 (vy)-cvi-če-ný
 zprá-va
 ge-ria-tr-(ička)
 osud-ný
 ná-pra-vy
 vstup-ní
 obec-ně
 re-pub-li-ka
 (po)-div-ný
 pa-slun-ce
 vliv-ný
 ba-ham-ský
 pe-dia-tr-(ička)
 ha-sič-(ka)
 vnuč-ka
 bro-s-kvoň
 lo-di-vod-(ka)
 mi-s-tr-(yně)
 in-že-nýr-(ka)
 dr-ži-tel-(ka)
 to-pič-(ka)
 pro-mi-nent-ní
 takt-ní
 frag-men-tár-ní
 po-zo-ru-hod-ný
 úspě-š-ný
 trans-fu-ze
 bur-žoa-zie
 pře-stro-je-ní
 oh-ni-vzdor-ný
 dis-tan-ční
 de-frau-da-ce
 fi-s-kál-ní
 zběh-lý
 kvap-ný
 střed-ní
 sa-mo-hlá-s-ka
 pře-vo-di-tel-ný
 zban-k-ro-to-va-ný
 zá-clo-na
 zná-mý
 zá-klad-ní
 na-uka
 ro-š-ťa-č-ka
 lí-bán-ky
 sou-du
 zkou-š-ka
 mlý-nek
 prá-šiv-ka
 kul-ma
 la-hod-ný
 zá-meč-ni-ce
 ne-vý-znam-ný
 ovzdu-ší
 mno-ho-úhel-ní-ko-vý
 žra-vý
 zcest-ný
 esen-ce
 (při)-nu-ce-ný
 skr-b-lic-ký
 skle-s-lý
 bez-bož-ný
 ne-opráv-ně-ný
 ošu-mě-lý
 pru-der-ní
 chvá-s-ta-vý
 sou-drž-ný
 van-dal-(ka)
 sa-bo-tér-(ka)
 (po)-div-ně
 mo-k-řad
 plyt-ký
 na-spo-du
 izo-la-ce
 le-toun
 za-han-be-ný
 ří-ze-ní
 kom-plex-ní
 pr-kén-ko
 ba-vič-(ka)
 plyt-ký
 ne-mo-tor-ně
 šéf-(ka)
 přes-čas
 zji-s-ti-tel-ný
 ta-ko-vý
 kňou-ral-(ka)
 se-me-s-tr
 po-pra-s-ka-ný
 han-ba
 vzpří-me-ný
 chys-ta-ný
 lik-vi-da-ce
 or-ga-ni-za-ce
 ri-s-kant-ní
 (za)-pís-kat
 (za)-hvíz-dat
 vrou-cí
  }
\end{hyphenrules}


% ############################################################# CZECH TYPOGRAPHY #
% https://math.feld.cvut.cz/tkadlec/ftp/text/textypo.pdf



% ################################################################# DOCUMENT #

\begin{document}
\pagenumbering{gobble}

% -------------------------------------------------------------- Front cover -

\ifmakecovers
\covergeometry
  \makecoverwith\frontcoverimages
\restoregeometry  
\fi


% ============================================================= FRONT MATTER =

\ifmakefrontmatter

% ------------------------------------------------------------------ Patitul -

\cleardoublepage

\begin{center}
  %\vspace*{2cm}{\fontsize{20}{20}\selectfont\scshape\dictnameCZ\\\dictnameIS}
  \vspace*{\stretch 1}
  {\Huge\scshape\dictnameCZ\\\dictnameIS}
  \vspace{\stretch 4}
\end{center}

% -------------------------------------------------------------------- Titul -

\cleardoublepage

\begin{center}
  %\vspace*{2cm}{\fontsize{24}{24}\selectfont\scshape\dictnameCZ\\\dictnameIS}
  \vspace*{\stretch 1}
  {\HUGE\scshape\dictnameCZ\\\dictnameIS}
  \vspace{\stretch 2}
\end{center}

% ----------------------------------------------------------------- Impresum -

\clearpage

\noindent\begin{minipage}[t][\textheight][t]{\textwidth}
Autoři \textit{Islandsko-českého studijního slovníku}:\\
  \textbf{Aleš Chejn, Ján Zaťko, Jón Gíslason}\\[\baselineskip]
\ifPDF
\else
Autor \textit{Islandské gramatiky}:\\
  \textbf{Vojtěch Kupča}\\[\baselineskip]
\fi
Poděkování:\\
  \textbf{Dorotě Nierychlewské-Chejn, Jiřině Hanzlové, Jiřině Chejnové,
          Simoně Petrásové,\\ Petru Mikešovi, Vojtěchu Kupčovi, 
          Renatě Peškové Emilsson, Amiru Mulahumici, \\ Kristýně Antonové, 
          Martině Kašparové, Lucii Brabcové, Samoriele, Adamu Kožouškovi}\\[\baselineskip]
Finanční příspěvek na výtisk exemplářů určených knihovnám\\ v České republice,
Slovenské republice a na Islandu poskytli:\\
  \textbf{Petr Mikeš, Dorota Nierychlewska-Chejn, Jiřina Chejnová, Otta Chejn}
  \ifPDF
  \\[39\baselineskip]
  \else
  \\[36\baselineskip]
  \fi
\textbf{Hvalur.org - 2016}\\
Licence: \textbf{GNU Free Documentation License}\\
Sazba: \textbf{Aleš Chejn} v prostředí  \textbf\LaTeX\\
  Zvláštní poděkování patří \textbf{Paolu al-Imk\=an\=\i\ Brasolin} a uživatelům  \textbf{\TeX\ - \LaTeX\ Stack Exchange}\\[\baselineskip]
\ifPDF
Verze: \textbf{PDF verze} \textit{Islandsko-českého studijního slovníku} vhodná pro \textbf{\ifscreen offline prohlížení a vyhledávání\else vlastní tisk\fi}\\ volně dostupná na \url{www.hvalur.org}\\[\baselineskip]
\else
\textbf\ISBN\\[\baselineskip]
\fi
\ifPDF
\else
Tisk: \textbf{Tiskárna Melmen, Pardubice}\\
\fi
Heslová část vysázena patkovým fontem \TeX\ Gyre Pagella,
hesla bezpatkovým fontem {\phvfamily Helvetica}.
\end{minipage}

% -------------------------------------------------------- Table of contents -

\cleardoublepage 

\pdfbookmark\contentsname{toc}    
\tableofcontents
\thispagestyle{empty}
\cleardoublepage 

\fi

\pagenumbering{arabic}
\setcounter{page}{3}
% ============================================================= INTRODUCTION =

\ifinputintroduction

% ------------------------------------- The study Icelandic-Czech dictionary -

\chapter{O Islandsko-českém studijním slovníku}
Cílem vytvoření\textit{ Islandsko-českého studijního slovníku} je zpřístupnění islandského jazyka všem zájemcům o~tento jazyk. 

\textit{Islandsko-český studijní slovník} je určen především českým a~slovenským studentům islandštiny, začátečníkům i~pokročilým, kteří zde naleznou ucelené fonetické, morfologické a~syntaktické informace, jakož
i~slovní zásobu potřebnou k~aktivnímu i~pasivnímu užití s~rozsahem vhodným ke studiu islandštiny na vysoké škole.
Slovník je rovněž určen všem zájemcům o~islandštinu a~Island, jimž slovník nabízí encyklopedické informace, úvod do fonetiky a~gramatiky a~velké množství fotografií a~ilustrací. 
Třebaže se jedná o~slovník primárně zaměřený na potřeby českých a~slovenských uživatelů, obsahuje také řadu informací, které usnadní práci islandskému uživateli (např. použití více jak 10 000 synonym a~antonym, překlady příkladů ap.).

Slovní zásoba původně vychází z~\textit{Concise Icelandic-English Dictionary} (\cite {ic_en}), značně však byla rozšířena díky přístupu k~webovému projektu \textit{ISLEX} (\cite {int1}). Velké množství hesel bylo
přidáno na podnět uživatelů online verze \textit{Islandsko-českého studijního slovníku}.

Slovník obsahuje 30~575 hesel s~46~725 významy nebo slovními spojeními. Z~celkového počtu 30~575 je 28~097 svébytných hesel obsahujících alespoň jeden význam nebo slovní spojení. 
Zbylá hesla (v~počtu 2~478) jsou varianty, dublety nebo nepravidelné tvary.

Slovník je primárně online elektronický slovník, volně přístupný na internetových stránkách \url{www.hvalur.org}  (\cite {int14}). 
Tištěná verze slovníku je publikována na vlastní náklady v~počtu 30 kusů, přičemž většina exemplářů je určena knihovnám a~institucím v~České republice, Slovenské republice a~na Islandu
a~několik exemplářů je určeno autorům nebo zájemcům z~řad veřejnosti. Dvě PDF verze (vyhledávání\,/\addthin prohlížení offline a~vlastní tisk) jsou dostupné zdarma na internetových stránkách slovníku (viz výše).

\textit{Islandsko-český studijní slovník} je otevřený svobodný projekt, jehož se v~průběhu tvorby zúčastnilo velké množství spoluautorů (viz Historie tvorby IČSS), kterým tímto jménem autorů velmi děkuji.
Děkuji také všem uživatelům online verze slovníku za trpělivost, dobrou vůli a~veškeré připomínky a~návrhy, jmenovitě především Petru Mikešovi, Adamu Kožouškovi a~Samoriele za neutuchající víru v~náš projekt.
Poděkování patří také našim rodinám a~blízkým, díky kterým jsme se mohli soustavně věnovat tvorbě slovníku.
Srdečné poděkování patří mé ženě Dorotce za podporu a~nadšení potřebné k~tvorbě slovníku.

\blspace[5]

{\centering Aleš Chejn, Brwinów, 24. prosince 2015\par}


% ------------ The history of the formation Icelandic-Czech dictionary study -

\chapter{Historie tvorby Islandsko-českého studijního slovníku}
\longchapterskip
Projekt \textit{Islandsko-českého studijního slovníku} (IČSS) vznikl v~roce 2001 na nepřímý podnět doc. PhDr. Heleny Kadečkové, CSc. z~Univerzity Karlovy,
kdy byl vypracován seznam slov. 
Tento seznam obsahoval islandská slova a~jejich překlady, později byly přidány základní morfologické informace o~slovním druhu a~deklinační\,/\addthin konjugační koncovky podle vzoru z~\textit{Concise Icelandic-English Dictionary} (\cite {ic_en}). 
Seznam v~této době obsahoval okolo 1 200 slov. Byl zveřejněn pomocí statických webových stránek a~později byl umístěn na webu Univerzity Karlovy. 

V~roce 2006 byl slovník rozšířen o~nových 1 500 heslových slov a~stávající hesla byla rozšířena o~nové významy, zvláště významy frázových sloves. 
Ke konci roku 2006  byl slovník převeden z programu Excel do databáze MySQL a~byly zahájeny práce na vytvoření aplikace \textit{Dictionary System} (DS) (\cite {int19}), která by podporovala tvorbu IČSS. 
Obsah slovníku byl zveřejněn pod svobodnou licencí na webových stránkách \url{www.hvalur.org} (\cite {int14}), kde se nachází dodnes. 

V~roce 2007 byla vytvořena první verze aplikace DS. V~této době se k~projektu přidala Renata Pešková Emilsson, která se nyní věnuje podpoře výuky češtiny na Islandu. 
Ta slovník obohatila rozsáhlým seznamem slov z~praktického života na Islandu. Dále se slovník rozšiřoval o~heslová slova podle \textit{Concise Icelandic-English Dictionary} (\cite {ic_en}). 
V průběhu následujících třech let bylo doplněno přibližně 15~000 slov. V této době byla hesla rozšířena o~synonyma, příklady, frekvenci, oborovou a~stylovou charakteristiku. K projektu se na Islandu připojila na krátkou dobu Kristýna Antonová, studentka islandštiny na Islandské univerzitě, která 
nás upozornila na nesrovnalosti v~uživatelském rozhraní aplikace DS.

Během léta 2009 byl vytvořen skript na generování skloňování a~časování heslových slov. Výsledky generování byly porovnány s~\textit{Beygingarlýsing íslensks nútímamáls} (\cite {DalvikVM}) a~během následujících let konzultovány
s~Jónem Gíslasonem. Kristín Bjarnadóttir, vedoucí projektu \textit{Beygingarlýsing íslensks nútímamáls}, zodpověděla řadu našich dotazů týkajících se deklinací a~konjugací.

V~roce 2009 začala spolupráce s~\textit{Biolib.cz} (\cite {int5}), jmenovitě s~Ondřejem Zichou. Obdrželi jsme  databázi s~náhledy obrázků, autory a~licencemi a~latinskými názvy botanických a~zoologických druhů.

V~zimním semestru 2009 jsme požádali s~Amirem Mulahumicem, bosenským studentem islandštiny na Islandské univerzitě, Studentský fond Islandské univerzity o~grant na tvorbu zvukové databáze islandských slov. 
Studentský fond návrh přijal. Peníze z~grantu jsme určili pro islandského rodilého mluvčího, Jóna Gíslasona, lektora islandského jazyka na Islandské univerzitě v~Reykjavíku, který nahrál více jak 22~000 islandských slov do elektronické formy. 
S~Dorotou Nierychlewskou-Chejn jsme připravili dva seznamy heslových slov (jeden pro nahrávajícího a~druhý pro zpracování nahrávek). Na podporu nahrávek jsme vytvořili webové stránky nazvané \textit{Databáze s~výslovností islandských slov} (\cite {int13}) a~tyto stránky jsme přeložili do 6~jazyků (Jón Gíslason do islandštiny, Ján Zaťko do slovenštiny, francouzštiny a~angličtiny, Amir Mulamuhic do srbochorvatštiny a~Dorota Nierychlewska-Chejn do polštiny).

V~roce 2010 jsme přidali do databáze pravidla islandské výslovnosti, která jsou uvedena v~knize \textit{Handbók um íslenskan framburð} (\cite {is2}). Z~těchto pravidel jsme automaticky vygenerovali fonetický zápis výslovnosti všech hesel ve slovníku. 
Jón Gíslason ověřil správnost fonetického zápisu u~všech heslových slov.

Na konci roku 2011 jsme spojili webové stránky slovníku se stránkami aplikace a~sjednotili jsme a~vylepšili vzhled webových stránek. Cennými připomínkami se na vytvoření aplikace podílela
RNDr. Ingrid Nagyová, PhD., vysokoškolská pedagožka Ostravské Univerzity. Ján Zaťko přeložil nové webové stránky, Průvodce po slovníku do slovenštiny, francouzštiny a~angličtiny a~celý manuál aplikace do angličtiny. 
Dorota Nierychlewska-Chejn přeložila stránky do polštiny, Jón Gíslason do islandštiny a~Marcos Helena do portugalštiny.

Od prosince 2011 do března 2014 jsme realizovali první čtení, ve kterém jsme doplnili chybějící informace, přidali nová hesla a~opravili pravopisné chyby. 
Obsah slovníku jsme porovnávali především s~\textit{Íslensk orðabók} (\mbox{\cite {is}}) a~\textit{Islandsko-českým slovníkem} (\cite {is7}). S~Vojtěchem Kupčou jsme konzultovali mnohé nesnáze při překladech.

V~letech 2013 až 2015 byl vytvořen morfologický klíč, který rozšiřuje klasický systém zobrazení morfologické informace v~islandských slovnících o~přesné zařazení daného hesla do určité deklinační\,/\addthin konjugační skupiny, jejíž
skloňování\,/\addthin časování je ukázáno na příkladu. V~roce 2014 jsme obdrželi roční grant od Ostravské univerzity v~Ostravě na tvorbu morfologického klíče. S~realizací grantu nám velmi pomohla doc. PhDr. Diana Svobodová, Ph.D. a~Ing. Zuzana Havrlantová.

Od března 2014 do srpna 2015 jsme realizovali druhé finální čtení, ve kterém jsme porovnávali informace s~webovým slovníkem \textit{ISLEX} (\cite {int1}). 
Na požádání jsme obdrželi od Þórdís Úlfarsdóttir (vedoucí projektu \textit{ISLEX}) seznam víceslovných příslovečných frází. Takto poskytnutý seznam nám usnadnil opravy slovníku. 
S~Þórdís Úlfarsdóttir jsme konzultovali různé gramatické jevy a~záležitosti týkající se deklinací\,/\addthin konjugací.
Martina Kašparová, studentka islandštiny na Islandské univerzitě, prošla celý slovník, přičemž se soustředila mimo jiné na stylistické a~pravopisné korektury překladů. 
Lucie Brabcová udělala jazykové a~typografické korektury celého slovníku.

Vojtěch Kupča vydal v~roce 2014 \textit{Islandskou gramatiku} (\cite {is77}) a~souhlasil s~připojením této publikace ke~knižní verzi IČSS.
V~roce 2015 jsme s~přispěním Paola al-Imk\=an\=\i\ Brasolin vysázeli slovník v~typografickém prostředí {\LaTeX}. V~únoru 2016 byl slovník vytištěn v~Tiskárně MELMEN, s.r.o v~Pardubicích a svázán 
v~Knihařství Vojtíšková Hrubá v~Pardubicích. 

\blspace[5]

{\centering Všem spoluautorům děkujeme za spolupráci.\par}


% -------------------------------------------------  Guide to the dictionary -

\chapter{Průvodce po slovníku}
\begin{multicols}{2}
\def\tableA{\noindent\begin{tabularx}{\columnwidth}{>{\bfseries}lll}
  \toprule
  \textit{přípona} & \textit{přípona s~koncovkou} & \textit{příklad} \\
  \midrule
  -ing-   & -ingur   & rugl··ingur \\
          & -ing     & hreyf··ing \\
          & -ingi    & erf··ingi \\
          & -ingar   & lækn··ingar \\
  -ling-  & -lingur  & vett··lingur \\
  -ning-  & -ningur  & flut··ningur \\
          & -ning    & set··ning \\
          & -ningar  & kos··ningar \\
  -að-    & -aður    & mark··aður \\
  -nað-   & -naður   & skil··naður \\
  -ar-    & -ari     & dóm··ari \\
  -skap-  & -skapur  & kveð··skapur \\
  -dóm-   & -dómur   & lær··dómur \\
  -leik-  & -leikur  & sann··leikur \\
          & -leiki   & mögu··leiki \\
  -lát-   & -látur   & laus··látur \\
  -átt-   & -átta    & víð··átta \\
  -sam-   & -samur   & starf··samur \\
  -sem-   & -semi    & hamingju··semi \\
  -nesk-  & -neskur  & tékk··neskur \\
          & -neska   & tékk··neska \\
  -leg-   & -legur   & blíð··legur \\
          & -lega    & blíð··lega \\
  -ung-   & -ungur   & gler··ungur \\
          & -unga    & snopp··unga \\
          & -ungi    & sveit··ungi \\
          & -ung     & nýj··ung \\
  -ótt-   & -óttur   & blett··óttur \\
  -ræn-   & -rænn    & suð··rænn \\
          & -ræna    & nor··ræna \\
  -ern-   & -erni    & fað··erni \\
  -ind-   & -indi    & vís··indi \\
  -and-   & -andi    & gef··andi \\
  -ug-    & -ugur    & göf··ugur \\
          & -ugt     & stöð··ugt \\
  -ust-   & -usta    & þjón··usta \\
  -ald-   & -aldur   & far··aldur \\
  -uð-    & -uður    & söfn··uður \\
  -úð-    & -úð      & dul··úð \\
  -versk- & -verskur & ung··verskur \\
          & -verska  & ung··verska \\
  -ist-   & -isti    & hass··isti \\
  -ism-   & -ismi    & húman··ismi \\
  -ang-   & -angur   & far··angur \\
  \bottomrule
  \end{tabularx}
  \captionof{table}{Seznam přípon označených ve slovníku}
  \label{table1}}

\def\tableB{\noindent\begin{tabularx}{\columnwidth}{>{\bfseries}lX}
  \toprule
  \textbf{symbol} & \textbf{vysvětlivka} \\
  \midrule
  m                   & označuje slovní druh \\
  \textsubscript{1}   & označuje třídu \\
  (-a, -ar)           & označuje deklinační / konjugační koncovky hesla \\
  \textsubscript{8}   & označuje podtřídu\\
  \bottomrule
  \end{tabularx}
  \captionof{table}{Morfologická informace}
  \label{table2}}

\def\tableC{\noindent\begin{tabularx}{\columnwidth}
    {>{\bfseries}l>{\footnotesize\itshape}lllll}
  \toprule
  \multirow{3}{*}{{{\textbf{f{\textsubscript{1}}} \Large{\textbf{7}}}}}
    & nom & kisa & kisan    & kisur & kisurnar \\
    & acc & kisu & kisuna   & kisur & kisurnar \\
    & dat & kisu & kisunni  & kisum & kisunum \\
    & gen & kisu & kisunnar & kisa  & kisanna \\
  \bottomrule
  \end{tabularx}
  \captionof{table}{Skloňování hesla \textit{kisa}}
  \label{table3}}

\def\tableD{\noindent\begin{tabularx}{\columnwidth}{>{\bfseries}lllll}
  \toprule
  \textit{zkratka} & \textit{slovo} & \textit{pád} & \textit{číslo}
    & \specialcell{\textit{předmět}\\\textit{podmět}} \\
  \midrule
  e-að  & eitthvað   & 1. & sg & neživotný \\
  e-ð   & eitthvað   & 4. & sg & neživotný \\
  e-ir  & einhverjir & 1. & pl & životný \\
  e-ja  & einhverja  & 4. & pl & životný \\
  e-jum & einhverjum & 3. & pl & životný \\
  e-m   & einhverjum & 3. & sg & životný \\
  e-n   & einhvern   & 4. & sg & životný \\
  e-r   & einhver    & 1. & sg & životný \\
  e-rra & einhverra  & 2. & pl & životný \\
  e-rs  & einhvers   & 2. & sg & životný \\
  e-s   & einhvers  & 2. & sg & neživotný \\
  e-u   & einhverju  & 3. & sg & neživotný \\
  \bottomrule
  \end{tabularx}
  \captionof{table}{Seznam použitých zkratek neurčitých zájmen}
  \label{table4}}

%#############################################################################

\section{Heslo}

\subsection*{Řazení slov}

Heslová slova jsou zobrazena tučně a~jsou seřazena podle islandské abecedy:

\blspace
{\bfseries\centering
  a~á b c d ð e é f g h i~í j k~l m n\\
  o~ó p q r s~t u~ú v~w x y ý z~þ æ ö\par}
\blspace

K~snadnější orientaci v~heslové části je použit u~každého písmene tzv. palcový index, který se nachází na hřbetě slovníku. Každé písmeno je v~palcovém indexu barevně odlišeno.

\subsection*{Homonyma}

Homonyma (heslová slova stejného reprezentativního tvaru) jsou označena povýšenými arabskými číslicemi uváděnými za heslovým slovem.

\blspace
  \dicEntry \dicTerm{vor\textsuperscript{1}} \dicPos{n} \dicFlx{(-s, -)}
  \dicEntry \dicTerm{vor\textsuperscript{2}} \dicPos{pron} \dicFlx{pers gen pl}
\blspace

Homonyma jsou seřazena podle českého ustáleného řazení slovních druhů. Poté následují předpony a~zkratky.
V~případě dvou podstatných jmen je pořadí mužský rod, ženský rod, střední rod. V~případě dvou podstatných jmen stejného rodu má přednost slovo se slabým skloňováním před slovem se silným skloňováním. V~případě dvou sloves stejného reprezentativního tvaru má přednost sloveso se slabým časováním před slovesem se silným časováním.

\subsection*{Varianty a~dublety}

U~některých heslových slov jsou v~záhlaví hesla uvedeny a) varianty (obměny)

\blspace
  \dicEntry \dicTerm{ærsla··fenginn}, \dicTerm{ærslafullur} \dicPos{adj}
\blspace

nebo b) dublety (dvojtvary).

\blspace
  \dicEntry \dicTerm{{ævin··týri}, \dicTerm{ævintýr}} \dicPos{n}
\blspace

Varianty a~dublety jsou méně frekventované než samotné heslo. Lze je vyhledat ve slovníku jako hesla obsahující fonetické a~morfologické informace s~odkazem na původní heslové slovo.

\blspace
  \dicEntry \dicTerm{ærsla··|fullur} \dicIPA{{a}{i}{\textsubring{r}}{s}{\textsubring{d}}{l}{a}{f}{\textscy}{\textsubring{d}}{l}{\textscy}{\textsubring{r}}} \dicPos{adj} \dicFlx{(comp -fyllri, sup -fyllstur)} \dicLink{ærslafenginn}
  \dicEntry \dicTerm{ævin··týr} \dicIPA{{a}{i}{\textlengthmark}{v}{\textsci}{n}{t\textsuperscript{h}}{i}{\textsubring{r}}} \dicPos{n} \dicFlx{(-s, -)} \dicLink{ævintýri}
\blspace

\subsection*{Vybrané tvary jako hesla}

Ve slovníku jsou uvedeny vybrané tvary heslových slov, u~kterých by mohl mít uživatel potíže s~identifikací. Jedná se o~a) genitiv jednotného čísla ({\textbf{sg}}) a~nominativ množného čísla ({\textbf{pl}}) podstatných jmen, b) ženský a~střední rod přídavných jmen a~2. a~3. stupeň rodu mužského přídavných jmen, c) oznamovací způsob ({\textbf{ind}}) přítomného času ({\textbf{praes}}) a~minulého času ({\textbf{pf}}) 1. osoby ({\textbf{pers}}) sg a~pl a~spojovací způsob ({\textbf{con}}) minulého času 1. osoby sg a~jiné nepravidelné tvary.

\blspace
  \dicEntry \dicTerm{öttum} \dicPos{v} \dicFlx{ind pf pl 1 pers} \dicLink{etja}
\blspace

\subsection*{Dělení heslových slov}

K~dělení heslových slov je použita střední tečka \textbf{'·'}, která odděluje části složeného slova, přípony a~předpony. Platí následující pravidla:

\renewcommand{\labelenumi}{\alph{enumi})}
\begin{enumerate}
\item jediné dělení heslového slova je označeno dvěma středními tečkami \textbf{'··'}. Může se jednat o~dělení částí složeného slova
\blspace
\dicEntry \dicTerm{ráð··hús}
\blspace

nebo o~dělení mezi kořenem a~předponou / příponou.

\blspace
\dicEntry \dicTerm{þjón··usta}
\blspace
\item V~případě, že se v~hesle vyskytuje více dělení, jsou použity dvě střední tečky \textbf{'··'} pro významově poslední dělení a~jedna střední tečka \textbf{'·'} pro zbylé / zbylá dělení.
\blspace
\dicEntry \dicTerm{við·bótar··líf·eyris·sparn·aður}
\blspace
\item V~hesle s~předponou je označena jedna nebo i~více předpon (každá zvlášť).
\blspace
\dicEntry \dicTerm{ó··að·finnan·legur}
\blspace
\item V~hesle s~příponou je označena jedna nebo i~více přípon (každá zvlášť).
\blspace
\dicEntry \dicTerm{létt·úð··ugur}
\blspace
\item V~hesle s~příponou je označena pouze přípona / přípony poslední části složeného slova. Důvodem pro neoznačení přípon jiných částí složených slov je lepší čitelnost hesla. Příklad označení pouze přípony poslední části složeniny:
\blspace
\dicEntry \dicTerm{flutninga··miðl·ari}
\end{enumerate}

Následuje seznam přípon, které jsou ve slovníku označeny:

\tableA

\section{Frekvence slov}

Symbolem \dicsymFrequent\ je označeno 2~825 hesel patřících mezi velmi frekventovaná islandská slova. Tento seznam vychází z~korpusu \textit{Íslenskur Orðasjóður} (\cite{int7}) a~knihy \textit{Íslensk orðtíðnibók} (\cite{is2}).

\blspace
  \dicEntry \dicTerm{koma}\textsuperscript{2} \dicsymFrequent
\blspace

\section{Fonetický zápis}

V~hranatých závorkách v~záhlaví hesla je uveden fonetický zápis výslovnosti hesla v~mezinárodní fonetické abecedě (IPA, angl. International Phonetic Alphabet).

\blspace
  \dicEntry \dicTerm{öldu··rót} \dicIPA{{\oe}{l}{\textsubring{d}}{\textscy}{r}{ou}{\textsubring{d}}}
\blspace

Znak '{\textlengthmark}' označuje dlouhou samohlásku nebo souhlásku.

\blspace
  \dicEntry \dicTerm{ör}\textsuperscript{2} \dicIPA{{\oe}{\textlengthmark}{\textsubring{r}}}
\blspace

\section{Slovní druhy, předpony a~zkratky}

V~záhlaví hesla je uvedena morfologická informace popisující a) slovní druh hesla, b) předponu nebo c) zkratku. Pokud se heslo řadí do více slovních druhů, jsou jednotlivé zkratky pro slovní druhy odděleny lomítkem.

\blspace
  \dicEntry \dicTerm{af} \dicPos{prep / adv}
\blspace

Informace o~slovním druhu, předponě či zkratce je vysázena barevně a~tučně.

\subsection*{Podstatná jména}

Podstatná jména jsou uvedena v~1. pádě (\textbf{nom}) jednotného čísla (\textbf{sg}) nebo v~1. pádě množného čísla (\textbf{pl}), pokud se heslo v~jednotném čísle nevyskytuje. Podstatná jména jsou označena následovně: \textbf{m} (l. masculinum) - podstatné jméno rodu mužského,
\textbf{f} (l. femininum) - podstatné jméno rodu ženského, \textbf{n} (l. neutrum) - podstatné jméno rodu středního, \textbf{subs} (l. substantivum) - podstatné jméno bez jednoznačného mluvnického rodu.

\blspace
  \dicEntry \dicTerm{hestur} \dicPos{m}
  \dicEntry \dicTerm{Hanoí} \dicPos{subs}
\blspace

\subsection*{Přídavná jména}

Přídavná jména jsou uvedena v~1. pádě (\textbf{nom}) rodu mužského (\textbf{m}) jednotného čísla (\textbf{sg}) 1. stupně (\textbf{pos}) (l. positivus) a~jsou označena zkratkou \textbf{adj} (l. adjectivum).

\blspace
  \dicEntry \dicTerm{fallegur} \dicPos{adj}
\blspace

V~případě, že u~přídavného jména dochází v~rodě ženském (\textbf{f}) nebo středním (\textbf{n}) k~přehlásce anebo je tvar druhého stupně (\textbf{comp} (l. comparativus)) či třetího stupně (\textbf{sup} (l. superlativus)) nepravidelný, pak jsou tyto tvary uvedeny v~nominativu jako heslová slova s~odkazem na hlavní heslový tvar

\blspace
  \dicEntry \dicTerm{yngri} \dicPos{adj} \dicFlx{comp m} \dicLink{ungur}
\blspace

nebo jako svébytná hesla s~významy.

\blspace
  \dicEntry \dicTerm{fremri} \dicPos{adj} \dicFlx{comp (sup fremstur)}
\blspace

\subsection*{Zájmena}

Zájmena jsou označena zkratkou \textbf{pron} (l. pronomen) a~jsou tříděna na zájmena ukazovací \textbf{dem} (angl. demonstrative), přivlastňovací \textbf{poss} (angl. possessive), osobní \textbf{pers} (angl. personal), neurčitá \textbf{indef} (angl. indefinite), zvratná \textbf{refl} (angl. reflexive) a~vztažná \textbf{rel} (angl. relative).

\blspace
  \dicEntry \dicTerm{þessi} \dicPos{pron} \dicFlx{dem}
\blspace

U~zájmen přivlastňovacích je v~3. osobě označen rod.

\blspace
  \dicEntry \dicTerm{hans} \dicPos{pron} \dicFlx{poss m}
\blspace

U~zájmen osobních je označen pád a~číslo.

\blspace
  \dicEntry \dicTerm{mig} \dicPos{pron} \dicFlx{pers acc sg}
\blspace

\subsection*{Číslovky}

Číslovky jsou označeny zkratkou \textbf{num} (l. numerale). V~případě vybraných číslovek (1 až 4) jsou ve slovníku uvedeny jako heslová slova také tvary ženského i~středního rodu. V~takovém případě jsou označeny zkratkou \textbf{m} pro mužský rod, zkratkou \textbf{f} pro ženský rod a~zkratkou \textbf{n} pro střední rod. Zkratkou \textbf{ord} (angl. ordinal number) jsou označeny číslovky řadové.

\blspace
  \dicEntry \dicTerm{tut··tugu} \dicPos{num}
  \dicEntry \dicTerm{tveir} \dicPos{num} \dicFlx{pl m}
  \dicEntry \dicTerm{fyrsti} \dicPos{num} \dicFlx{ord}
\blspace

\subsection*{Slovesa}

Slovesa jsou uvedena v~infinitivu a~jsou označena zkratkou \textbf{v} (l. verbum).

\blspace
  \dicEntry \dicTerm{fara} \dicPos{v}
\blspace

Slovesa ve středním rodě jsou označena jako \textbf{refl} (angl. reflexive).

\blspace
  \dicEntry \dicTerm{nálgast} \dicPos{v} \dicFlx{refl}
\blspace

Slovesa, která jsou vždy neosobní, jsou označena \textbf{impers} (angl. impersonal).

\blspace
  \dicEntry \dicTerm{svima} \dicPos{v} \dicFlx{impers}
\blspace

U~sloves je dále uvedeno, s~jakým pádem se pojí. Zkratka \textbf{acc} označuje 4. pád, \textbf{dat} 3. pád, \textbf{gen} 2. pád a~\textbf{nom} 1. pád. V~případě, že se sloveso pojí s~různými pády, je použito lomítko.

\blspace
  \dicEntry \dicTerm{klóra} \dicPos{v} \dicFlx{acc / dat}
\blspace

V~případě, že se sloveso pojí s~více pády, je použito znaménko plus.

\blspace
  \dicEntry \dicTerm{gefa} \dicPos{v} \dicFlx{dat + acc}
\blspace

\subsection*{Příslovce}

Příslovce jsou označena zkratkou \textbf{adv} (l. adverbium). Pokud chybí informace o~stupni, jedná se o~1. stupeň stupňování (\textbf{pos}). 2. stupeň je vždy označen zkratkou \textbf{comp} a~3. stupeň zkratkou \textbf{sup}.

\blspace
  \dicEntry \dicTerm{þar} \dicPos{adv}
  \dicEntry \dicTerm{ofar} \dicPos{adv} \dicFlx{comp (pos uppi, sup efst)}
  \dicEntry \dicTerm{síðast} \dicPos{adv} \dicFlx{sup (comp síðar)}
\blspace

\subsection*{Předložky}

Předložky jsou označeny zkratkou \textbf{prep} (l. praepositio). Pokud se předložka pojí výhradně s~jedním pádem, jsou použity zkratky pro 4. pád (\textbf{acc}), 3. pád (\textbf{dat}), 2. pád (\textbf{gen}).

\blspace
  \dicEntry \dicTerm{af} \dicPos{prep / adv} \dicFlx{dat}
\blspace

V~případě, že se předložka pojí s~různými pády, je použito lomítko.

\blspace
  \dicEntry \dicTerm{fyrir} \dicPos{prep / adv} \dicFlx{acc / dat}
\blspace

\subsection*{Spojky}

Spojky jsou označeny zkratkou \textbf{conj} (l. conjunctio).

\subsection*{Částice}

Částice jsou označeny zkratkou \textbf{part} (l. particula).

\subsection*{Citoslovce}

Citoslovce jsou označeny zkratkou \textbf{inter} (l. interjectio).

\subsection*{Předpony}

Předpony jsou označeny zkratkou \textbf{predp}.

\subsection*{Zkratky}

Zkratky jsou označeny zkratkou \textbf{zkr}.

\section{Skloňování, časování a~stupňování}

Každé heslo obsahuje zápis morfologické informace, která přesně informuje, jak se heslo skloňuje / časuje nebo stupňuje.

V~slovníku jsou použity dva odlišné systémy zápisu morfologické informace, které se vzájemně doplňují:

\begin{enumerate}
\item tradiční systém - do závorky za slovním druhem se zapisují deklinační / konjugační koncovky. Počet koncovek závisí na slovním druhu (viz dále).
\item numerický systém (vychází ze systému Helgi Haraldssona (\cite {is6})) - číslice odkazující na morfologický klíč, kde jsou podrobně vypsány deklinační / konjugační třídy.
Další informace o~numerickém systému lze nalézt na straně \pageref{sec:dec}.
\end{enumerate}

Deklinační / konjugační tvary nebo koncovky podstatných jmen, přídavných jmen, zájmen, číslovek, sloves a~příslovcí jsou uvedeny v~závorkách za slovním druhem.

Hnízdovací šev \textbf{'|'} označuje místo ve slově, kde se deklinační koncovka připojuje ke slovu.

\blspace
  \dicEntry \dicTerm{heils|a{\textsuperscript{1}}} \dicPos{f} \dicFlx{(-u)}
\blspace

Pokud není hnízdovací šev \textbf{'|'} uveden, tak se koncovka připojuje k~samotnému heslu.

\blspace
  \dicEntry \dicTerm{mynd} \dicPos{f} \dicFlx{(-ar, -ir)}
\blspace

Pokud dochází v~ohýbání k~přehlásce nebo je ohýbání nepravidelné, pak jsou uvedeny celé tvary hesla.

\blspace
  \dicEntry \dicTerm{maður} \dicPos{m} \dicFlx{(manns, menn)}
\blspace

V~případě, že heslové slovo má více tvarů, jsou varianty koncovky odděleny lomítkem.

\blspace
  \dicEntry \dicTerm{beð|ur} \dicPos{m} \dicFlx{(-s / -jar, -ir)}
\blspace

\subsection*{Podstatná jména}

U~podstatných jmen je uveden tvar hesla pro 2.~pád jednotného čísla a~1.~pád množného čísla.

\blspace
  \dicEntry \dicTerm{hval|ur} \dicPos{m} \dicFlx{(-s, -ir)}
\blspace

kde '\textit{hvals}' je 2.~pád jednotného čísla a~'\textit{hvalir}' 1.~pád množného čísla.
Pokud je uveden pouze jeden tvar, jedná se o~2.~pád jednotného čísla a~znamená to rovněž, že heslové slovo se neohýbá v~množném čísle.

\blspace
  \dicEntry \dicTerm{heils|a{\textsuperscript{1}}} \dicPos{f} \dicFlx{(-u)}
\blspace

Pokud je použita zkratka pro množné číslo (\textbf{pl}), znamená to, že se heslo ohýbá pouze v~množném čísle.

\blspace
  \dicEntry \dicTerm{afar··kostir} \dicPos{m} \dicFlx{pl}
\blspace

Pokud je použita zkratka \textbf{indecl}, znamená to, že podstatné jméno je nesklonné.

\blspace
  \dicEntry \dicTerm{október} \dicPos{m} \dicFlx{indecl}
\blspace

\subsection*{Přídavná jména}

U~přídavných jmen jsou uvedeny v~závorkách tvary, které poukazují na nepravidelnosti ve skloňování.

Pokud přídavné jméno v~ženském rodě neztrácí koncové '\textit{-r}', je uveden tvar hesla v~ženském rodě.

\blspace
  \dicEntry \dicTerm{hýr} \dicPos{adj} \dicFlx{(f -)}
\blspace

Dále jsou uvedeny tvary ženského rodu, ve~kterých dochází k~přehlásce:

\blspace
  \dicEntry \dicTerm{reglu··|samur} \dicPos{adj} \dicFlx{(f -söm)}
\blspace

Pokud je přídavné jméno nesklonné, je uvedena zkratka \textbf{indecl}.

\blspace
  \dicEntry \dicTerm{tví··mála} \dicPos{adj} \dicFlx{indecl}
\blspace

Pokud má přídavné jméno nepravidelné tvary 2.~a~3.~stupně stupňování, jsou tyto tvary uvedeny v~závorce spolu se zkratkou pro 2.~stupeň (\textbf{comp}) a~3.~stupeň (\textbf{sup}).

\blspace
  \dicEntry \dicTerm{þunnur} \dicPos{adj} \dicFlx{(comp þynnri, sup þynnstur)}
\blspace

V~případě, že je uveden pouze jeden ze dvou zbývajících stupňů, znamená to, že heslo se stupňuje pouze v~uvedených stupních.

\blspace
  \dicEntry \dicTerm{syðri} \dicPos{adj} \dicFlx{comp (sup syðstur)}
\blspace

\subsection*{Slovesa}

Islandská slovesa se dělí na slabá, silná a~nepravidelná. Vzhledem k~tomuto dělení se liší počet slovních tvarů uvedených v~závorce za zkratkou \textbf{v}. Princip uvádění slovních tvarů spočívá v~tom, že uživatel
je schopen (více méně) z~těchto tvarů odvodit časování slovesa. V~tomto slovníku jsme
záměrně umístili některá slabá slovesa mezi silná (místo dvou slovesných tvarů je uvedeno pět tvarů) kvůli tomu, aby byl uživatel přesněji informován o~časování slovesa.

Slabá slovesa, která tvoří 1.~osobu jednotného čísla minulého času oznamovacího způsobu koncovkou '\textit{-aði}', mají v~závorce uvedenu jen jednu koncovku, jmenovitě '\textit{(-aði)}'.

\blspace
  \dicEntry \dicTerm{ætl|a} \dicPos{v} \dicFlx{(-aði)}
\blspace

U~zbývajících skupin slabých sloves se v~závorce nachází koncovka 1.~osoby jednotného čísla minulého času oznamovacího způsobu a~supinum.

\blspace
  \dicEntry \dicTerm{kenn|a} \dicPos{v} \dicFlx{(-di, -t)}
\blspace

Silná nebo nepravidelná slovesa mají v~závorce vždy pět tvarů, jmenovitě 1.~osobu jednotného čísla přítomného času oznamovacího způsobu, 1.~osobu jednotného čísla minulého času oznamovacího způsobu, 1.~osobu množného čísla minulého času oznamovacího způsobu, 1.~osobu jednotného čísla minulého času spojovacího způsobu a~supinum.

\blspace
  \dicEntry \dicTerm{grípa} \dicPos{v} \dicFlx{(gríp, greip, gripum, gripi, gripið)}
\blspace

Silná nebo nepravidelná slovesa vyskytující se pouze v~neosobním užití mají v~závorce 4 tvary, jmenovitě 3.~osobu jednotného čísla přítomného času oznamovacího způsobu, 3.~osobu jednotného čísla minulého času oznamovacího způsobu, 3.~osobu jednotného čísla minulého času spojovacího způsobu a~supinum.

\blspace
  \dicEntry \dicTerm{kross··|bregða} \dicPos{v} \dicFlx{(-bregður, -brá, -brygði, -brugðið)}
\blspace

U~některých významů v~heslovém slově se vyskytují gramatické informace, které popisují chování daného významu.

Zkratka \textbf{impers} označuje, že se daný význam slovesa vyskytuje v~neosobním užití. Pomocí slovního spojení je vyjádřeno užití slovesa.

\blspace
  \dicEntry \dicTerm{batn|a} \dicPos{v} \dicFlx{(-aði)} \dicTerm{e-m batnar} \dicFlx{impers} (kdo) se uzdravuje, (kdo) se zotavuje, (komu) je lépe
\blspace

Zkratka \textbf{refl} označuje mediopasivní tvary slovesa.

\blspace
  \dicEntry \dicTerm{að··lag|a} \dicPos{v} \dicFlx{(-aði)} \dicIdiom{aðlagast} {\dicPhraseIS{aðlagast e-u}} \dicFlx{refl} {přizpůsobit se (čemu)}
\blspace

\subsection*{Příslovce}

U~příslovcí se v~záhlaví nachází tvary stupňování. Tyto tvary jsou uvedeny v~závorce spolu se zkratkou pro 2.~stupeň (\textbf{comp}) a~3.~stupeň (\textbf{sup}).

\blspace
  \dicEntry \dicTerm{vel\textsuperscript{2}} \dicPos{adv} \dicFlx{(comp betur, sup best)}
\blspace

Nepravidelné tvary 2.~a~3.~stupně jsou uvedeny buď jako hesla odkazující na základní tvar v~1.~stupni (\textbf{pos}), anebo jako svébytná hesla. V~obou případech jsou uvedeny v~závorce zbývající stupně.

\blspace
  \dicEntry \dicTerm{síst} \dicPos{adv} \dicFlx{sup (pos varla, comp síður)}
\blspace

Pokud některý ze stupňů uveden není, znamená to, že pro dané slovo neexistuje.

\blspace
  \dicEntry \dicTerm{síðar} \dicPos{adv} \dicFlx{comp (sup síðast)}
\blspace

\section{Morfologická informace} \label{sec:dec}

V~záhlaví hesla se kromě deklinačních / konjugačních koncovek nachází také morfologické numerické informace, které odkazují na morfologický klíč (viz strana \pageref{sec:morpho}), kde jsou podrobně vypsány kompletní deklinační / konjugační tabulky.

\blspace
  \dicEntry \dicTerm{ap|i} \dicPos{m}[1] \dicFlx{(-a, -ar)}[8]
\blspace

Morfologická informace je zapsána podle následujícího schématu: \textbf{m{\textsubscript{1}} (-a, -ar){\textsubscript{8}}}

\blspace
\tableB
\blspace

Deklinační / konjugační třídy a~jejich počet nejsou v~islandštině ustálené a~slouží pouze pro rozlišení odlišných skloňování / časování.

Názvy podtříd slouží pro rozlišení odlišných skloňování / časování.

\subsection*{Práce s~morfologickým klíčem}

Pro vyhledání kompletních deklinačních / konjugačních tabulek slouží následující postup, který demonstrujeme na ukázce:

\blspace
  \dicEntry \dicTerm{kis|a} \dicPos{f}[1] \dicFlx{(-u, -ur)}[7]
\blspace

Nejdříve určíme slovní druh (v~případě podstatných jmen také rod), v~tomto případě podstatné jméno rodu ženského, a~navštívíme Morfologický klíč (str. \pageref{sec:morpho}) pro podstatná jména rodu ženského (str. \pageref{sec:morpho_f}).
Následně dohledáme třídu (v~tomto případě číslo {\textsubscript{\textbf{1}}}) a~následně dohledáme dodatečné numerické informace za závorkou (v~tomto případě číslo {\textsubscript{\textbf{7}}}).

\tableC

\section{Definice}

\subsection*{Základní tvar}

Základním tvarem definice je, že jednomu islandskému slovu odpovídá jeden nebo více českých významů.

\blspace
  \dicEntry \dicTerm{landa··fræð|i} \dicPos{f} \dicFlx{(-i)} \textCS{zeměpis, geografie}
\blspace

V~některých případech je v~české části použita závorka. V~prvním případě je závorka užívaná ke~zkrácení zápisu dvou slov ('\textit{pomoc}', '\textit{výpomoc}').

\blspace
  \dicEntry \dicTerm{að··stoð} \dicPos{f} \dicFlx{(-ar)} \textCS{(vý)pomoc, asistence}
\blspace

V~druhém případě se v~závorce nachází upřesňující informace, která zabraňuje dvojznačnosti.

\blspace
  \dicEntry \dicTerm{fíkj|a} \dicPos{f} \dicFlx{(-u, -ur)} \textCS{fík} {\footnotesize {\textCS{(plod)}}}
\blspace

V~mnoha případech je závorka v~české části použita pro příklad užití českého slova a~tím vymezení jeho významu.

\blspace
  \dicEntry \dicTerm{að·gengi··legur} \dicPos{adj} {přístupný, dostupný (informace ap.)}
\blspace

Lomítko \textbf{'/'} slouží v~české části významu k~oddělení více významů pojících se s~jedním nebo více slovy a~je použito v~české části následujícím způsobem:

\blspace
  \dicEntry \dicTerm{aðal··|gata} \dicPos{f} \dicFlx{(-götu, -götur)} \textCS{hlavní třída / ulice}
\blspace

Význam přečteme jako '\textit{hlavní třída, hlavní ulice}'.

\subsection*{Opis heslového slova}

V~případech, kdy islandské slovo nemá ekvivalent v~českém jazyce, jsou významy islandských slov opsány delším opisem a~zobrazeny menším písmem. Příkladem jsou slova z~lidových pověstí, kulinářské speciality nebo botanické pojmy.

\blspace
  \dicEntry \dicTerm{til··ber|i} \dicPos{m} \dicFlx{(-a, -ar)} {\footnotesize{pov.}} {\footnotesize{stvoření, které dojí krávy a~ovce hospodářů}}
\blspace

\subsection*{Slovní spojení}

V~definici hesla jsou kromě samotných významů uváděna také slovní spojení, která jsou sázena tučně. Po slovním spojení následuje překlad slovního spojení.

\blspace
  \dicEntry \dicTerm{lær|a} \dicPos{v} \dicFlx{(-ði, -t)} \textbf{læra e-ð utanbókar} {učit se (co) zpaměti, memorovat (co)}
\blspace

Závorky jsou použity na tu část slovního spojení, kterou je možné vypustit.

\blspace
  \dicEntry \dicTerm{ár\textsuperscript{1}} \dicPos{f} \dicFlx{(-ar, -ar)} \textbf{taka (of) djúpt í árinni} {\footnotesize{přen.}} {přehnat (co v~tvrzení)}
\blspace

Příklad můžeme přečíst jako '\textit{taka djúpt í árinni}' nebo '\textit{taka of djúpt í árinni}'.

Mezi slovní spojení v~tomto slovníku řadíme a) tvary heslového slova, b) příkladová spojení, c) ustálená slovní spojení a~d) rčení nebo přísloví.

\begin{enumerate}
\item Tvary heslového slova\\
Do této skupiny řadíme tvar množného čísla hesla (zkratka \textbf{pl}).
\blspace
\dicEntry \dicTerm{bót} \dicPos{f} \dicFlx{(-ar, bætur)} \textbf{bætur} \dicFlx{pl} {dávky, přídavky, příspěvky}\\

Dále do této skupiny řadíme tvar slovesa v~středním slovesném rodě (zkratka \textbf{refl}).\\

\dicEntry \dicTerm{finna} \dicPos{v} \dicFlx{(finn, fann, fundum, fyndi, fundið)} \dicIdiom{finnast} {\dicPhraseIS{finnast}} \dicFlx{refl} {být nalezen, najít se, nalézt se, nacházet se}
\blspace
\item Příkladová spojení\\
Příkladová spojení ilustrují komunikační chování heslových slov v~příslušných významech a~v~různých syntagmatech (např. sloveso s~určitou předložkou ap.). Jsou označena tučně modrou barvou a~odsazená pro snadnou orientaci jako svébytné heslo.
\blspace
\dicEntry \dicTerm{koma\textsuperscript{2}} \dicPos{v} \dicFlx{(kem, kom, komum, kæmi, komið)} \dicIdiom{koma}[að] {\dicPhraseIS{koma að e-m}} {přiblížit se ke (komu), přijít ke (komu)}
\blspace
\item Ustálená slovní spojení\\
Tento druh slovních spojení je charakterizován většinou přeneseným významem. Tato spojení jsou uváděna za posledním významem nebo příkladovým spojením a~jsou označená tučně modrou barvou a~odsazená pro snadnou orientaci jako svébytné heslo.
\blspace
\dicEntry \dicTerm{kast} \dicPos{n} \dicFlx{(-s, köst)} \dicIdiom{kast} {\dicPhraseIS{þegar til kastanna kemur}} {\footnotesize{přen.}} {když jde do tuhého}
\blspace
\item Rčení a~přísloví\\
Rčení a~přísloví se nacházejí vždy na konci definice a~jsou odlišená symbolem \dicsymProverb\ a~odsazená pro snadnou orientaci.
\blspace
\dicEntry \dicTerm{æfing} \dicPos{f} \dicFlx{(-ar, -ar)} \dicProverb{} \dicPhraseIS{Æfing skapar meistara.} {\footnotesize{přís.}} \textCS{Cvičení dělá mistra.}
\end{enumerate}

\section{Vícevýznamovost}

\subsection*{Vícevýznamová hesla}

Značné množství hesel má více než jeden význam. Významy jsou odlišeny arabskými číslicemi a~v~některých případech i~islandskými synonymy.

\blspace
  \dicEntry \dicTerm{land} \dicPos{n} \dicFlx{(-s, lönd)} {\textbf{1.}} \dicSynonym{þurrlendi} {souš, pevnina, země} {\textbf{2.}} \dicSynonym{árbakki} {břeh} {\textbf{3.}} \dicSynonym{ríki} {země, stát, vlast} {\textbf{4.}} \dicSynonym{landareign} {země, půda, pozemek}
\blspace

\subsection*{Vícevýznamová slovní spojení}

V~případě, že slovní spojení uvnitř definice má více významů, jsou jednotlivé významy odlišeny písmeny {\textbf{a., b., c.}} atd.

\blspace
  \dicEntry \dicTerm{drag} {\dicPos{n} \dicFlx{(-s, drög)} {\textbf{1.}} {mokřina, podmoklý terén;} {\textbf{2.}} \textbf{drög} {\footnotesize{pl}} {\textbf{a.}} \dicSynonym{uppsprettur} {prameny (řeky ap.)} {\textbf{b.}} \dicSynonym{undirbúningur} {náčrt, návrh.}}
\blspace

\section{Řazení významů a~slovních spojení v~definici}
\longsectionskip

V~definici jsou nejdříve seřazeny významy slova a~odlišeny arabskými číslicemi.

\subsection*{V rámci významu}

Slovní spojení jsou uváděna v~rámci jednotlivého významu, pokud se s~ním významově vážou.

\blspace
  \dicEntry \dicTerm{tím|i} \dicPos{m} \dicFlx{(-a, -ar)} {\textbf{1.}} \dicSynonym{tíð} {čas, doba}; \textbf{í þann tíma} \dicFlx{adv} \dicSynonym{þá\textsuperscript{2}} {pak} {\textbf{2.}} \dicSynonym{rétt stund} {(správný) čas / moment}
\blspace

\subsection*{Za posledním očíslovaným významem}

Pokud se slovní spojení neváže k~žádnému významu, je řazeno za posledním očíslovaným významem. Ustálená slovní spojení jsou uvedena symbolem \dicsymIdiom.

\blspace
  \dicEntry \dicTerm{finna} \dicPos{v} \dicFlx{(finn, fann, fundum, fyndi, fundið)} {\textbf{3.}} \dicSynonym{skynja} {cítit, vnímat} \dicIdiom{finna}[að] {\dicPhraseIS{finna að e-u}} {kritizovat (co), nacházet chyby na (čem)}
\blspace

\subsection*{Způsob řazení slovních spojení}

Slovní spojení, která se nachází za posledním očíslovaným významem, jsou řazena podle předložky / příslovce / zájmena, se kterou se pojí. Pro snadnou orientaci v~rozsáhlých definicích je uvedeno heslové slovo nebo tvar heslového slova plus předložka / příslovce / zájmeno a~barevně odlišeno modrou barvou.

\blspace
  \dicEntry \dicTerm{koma\textsuperscript{2}} \dicPos{v} \dicFlx{(kem, kom, komum, kæmi, komið)} \dicIdiom{koma}[út] {\dicPhraseIS{koma út}} {vyjít (kniha ap.), objevit se (na trhu ap.)}
\blspace

Poté následují slovní spojení v~slovesném středním rodě řazená opět podle předložky / příslovce / zájmena, se kterou se pojí:

\blspace
  \dicEntry \dicTerm{koma\textsuperscript{2}} \dicPos{v} \dicFlx{(kem, kom, komum, kæmi, komið)} \dicIdiom{komast} {\dicPhraseIS{komast}} \dicFlx{refl} {dostat se, dostávat se}
\blspace

Poté následují ustálená slovní spojení, která nebylo možné zařadit do předchozích položek:

\blspace
  \dicEntry \dicTerm{koma\textsuperscript{2}} \dicPos{v} \dicFlx{(kem, kom, komum, kæmi, komið)} \dicIdiom{koma} {\dicPhraseIS{(komdu) sæll (og blessaður)}} {ahoj, nazdar, dobrý den}
\blspace

Nakonec jsou uvedena přísloví a~rčení označená symbolem \dicsymProverb\:

\blspace
  \dicEntry \dicTerm{eik} \dicPos{n} \dicFlx{(-ar / -ur, -ur)} \dicProverb{} \dicPhraseIS{Ekki / eigi fellur eik við fyrsta högg.} {\footnotesize{přís.}} {Jednou ranou dub nespadne.}
\blspace

\section{Synonyma a~antonyma / hesla ke~srovnání}
\longsectionskip

\subsection*{Synonyma}

Synonyma nebo synonymní slovní spojení jsou uvedeny před opisem významu hesla. Jsou umístěna v~závorkách. Islandská synonyma slouží a) islandskému uživateli k~rozlišení českých významů b) českému uživateli jako dodatečná informace o~významu slova.

\blspace
  \dicEntry \dicTerm{staða} \dicPos{f} \dicFlx{(stöðu, stöður)} {\textbf{1.}} \dicSynonym{ástand} {situace, stav, poměry} {\textbf{2.}} \dicSynonym{embætti} {post, (pracovní) místo / pozice (ve firmě ap.)}
\blspace

V~případě, že synonymum lze nalézt jako heslo ve slovníku, je použita šipka '\dicsymSee\ '.

\blspace
  \dicEntry \dicTerm{bygging} {\dicPos{f} \dicFlx{(-ar, -ar)} {\textbf{1.}} \dicSynonym{hús} {budova, stavba, stavení} {\textbf{2.}} \textit{(það að byggja)} {(vý)stavba} }
\blspace

\subsection*{Antonyma / Hesla ke~srovnání}

Antonyma či hesla ke~srovnání jsou uvedena za opisem významu heslového slova a~jsou označena šipkou '$\shortuparrow$'.

\blspace
  \dicEntry \dicTerm{kaldur} \dicPos{adj} \dicFlx{(f köld)} {\textbf{1.}} {studený, chladný} \textit{($\shortuparrow$ heitur)}
\blspace

\section{Oborové a~stylové charakteristiky}

K~upřesnění významu heslového slova slouží oborové a~stylové charakteristiky.

\subsection*{Oborové charakteristiky}

Oborové charakteristiky zařazují islandské slovo do určitého oboru a~naznačují použití daného slova v~jazyce

\blspace
  \dicEntry \dicTerm{berg} \dicPos{n} \dicFlx{(-s, -)} {\textbf{1.}} {\footnotesize{geol.}} {hornina, kámen} {\textbf{2.}} \dicSynonym{klöpp\textsuperscript{1}} {skalní stěna, skalisko}
\blspace

a~vymezují význam českého slova.

\blspace
  \dicEntry \dicTerm{mús} \dicPos{f} \dicFlx{(-ar, mýs)} {\textbf{1.}} {\footnotesize{zool.}} {myš} (\textit{l. Mus}) {\textbf{2.}} {\footnotesize{poč.}} {myš}
\blspace

\subsection*{Stylové charakteristiky}

Stylové charakteristiky jsou uváděny v~případě, kdy význam heslového slova je možné použít pouze při určité příležitosti nebo je význam jazykově zabarven.

\blspace
  \dicEntry \dicTerm{gaur} \dicPos{m} \dicFlx{(-s, -ar)} {\textbf{1.}} {\footnotesize{hovor.}} {chlápek, maník, člověk}
\blspace

\section{Syntax}

Syntaktické informace popisují:

\begin{enumerate}
\item zdali je sloveso tranzitivní nebo intranzitivní a~s~jakým(i) pádem (pády) se pojí, pokud je tranzitivní,
\item zdali je podmět / předmět životný či neživotný,
\item zdali je sloveso osobní či neosobní a~v~jakém pádu je podmět, pokud je sloveso neosobní,
\item zdali je podmět gramatický,
\item zdali je sloveso použito ve středním rodě,
\item nejčastější příklady, které se pojí se slovesem,
\item se kterými příslovci nebo předložkami se sloveso pojí.
\end{enumerate}

Syntaktické informace se ve slovníku vyskytují ve třech formách: a) informace uvedené zkratkou, b) teoretické příklady a~c) praktické příklady.

\subsection*{Syntaktické informace uvedené zkratkou}

Zde máme na mysli informace o~pádě (\textbf{nom, acc, dat, gen}), slovesu v~středním slovesném rodě (\textbf{refl}), neosobním slovesu (\textbf{impers}), čísle (\textbf{sg, pl}), rozkazovacím způsobu (\textbf{imper}) a~spojovacím způsobu (\textbf{con}).

\subsection*{Teoretické příklady}

Teoretickým příkladem nazýváme příklad, kdy podmět nebo předmět je zastoupen zájmenem neurčitým v~pádě, se~kterým se slovo pojí a~sloveso je zpravidla v~infinitivu. V~případě, že podmět není v~1.~pádě (např. u~sloves neosobních), je sloveso uvedeno v~3.~osobě jednotného čísla přítomného času.
Ve slovních spojeních používáme zkratky pro tvary neurčitého zájmena \textbf{einhver} (někdo) a~\textbf{eitthvað} (něco). Tyto zkratky označují a) pád, b) číslo a~c) životnost nebo neživotnost podmětu / předmětu.\\
Následuje seznam zkratek neurčitých zájmen použitých ve slovníku: \\

\tableD
\blspace
Příklad použití zkratek s~neurčitými zájmeny:

\blspace
  \dicEntry \dicTerm{leit|a} \dicPos{v} \dicFlx{(-aði)} \textbf{leita að e-u / e-m} {hledat (co / koho)} \dicExampleIS{leita að lyklunum} \dicExampleCS{hledat klíče}
\blspace

Příklad tedy přečteme jako '\textit{leita að einhverju}' (\textit{hledat co}) a~'\textit{leita að einhverjum}' (\textit{hledat koho}). V~tomto případě se sloveso pojí s~3.~pádem a~předmět může být životný i~neživotný.

\subsection*{Praktické příklady}

Ve slovníku se vyskytuje velké množství praktických příkladů a~jejich překladů. Praktický příklad se liší od příkladu teoretického tím, že místo zájmen neurčitých jsou uvedena frekventovaná slova z~běžného projevu. Praktické příklady slouží k~ilustraci použití
heslového slova.\\
Praktickým příkladem může být a) (složené) slovo, b) slovní spojení nebo c) věta.

\begin{enumerate}
\item Slovo jako praktický příklad:
\blspace
\dicEntry \dicTerm{hett|a} \dicPos{f} \dicFlx{(-u, -ur)} \textbf{1.} {kapuce, kápě} \dicExampleIS{hettuúlpa} \dicExampleCS{bunda s~kapucí}
\blspace
\item Slovní spojení jako praktický příklad:
\blspace
\dicEntry \dicTerm{tím|i} \dicPos{m} \dicFlx{(-a, -ar)} \textbf{3.} \dicSynonym{klukkustund} {hodina} {\footnotesize{(šedesát minut)}} \dicExampleIS{tveggja tíma gangur} \dicExampleCS{dvouhodinový pochod}
\blspace
\item Věta jako praktický příklad:
\blspace
\dicEntry \dicTerm{mæt|a} \dicPos{v} \dicFlx{(-ti, -t)} \dicIdiom{mætast} \dicPhraseIS{mætast} \dicSynonym*{ná saman} \dicFlx{refl} {setkat se, potkat se} \dicExampleIS{Þau mættust á gangstéttinni.} \dicExampleCS{Potkali se na chodníku.}
\end{enumerate}

\section{Fotografie a~ilustrace}

Ve slovníku je publikováno 630 fotografií rostlin, živočichů a~fotografií, která dokumentují místopisně heslová slova (např. kulinářské speciality).
Fotografie jsou publikovány pod veřejnými licencemi. Jména autorů a~názvy licencí lze nalézt v~Seznamu autorů fotografií a~ilustrací na straně \pageref{sec:photo}. U~hesla a~konkrétně u~daného významu, ke kterému se fotografie vztahuje, je použit symbol fotoaparátu \dicsymPhoto. Fotografie je umístěna buď pod heslem, anebo v~horní části stránky, kde se heslo nachází, anebo na stránce (stránkách) následující(ch).

%\par\begin{center}\setlength\fboxsep{0pt}\setlength\fboxrule{0.5pt}
%\fbox{\includegraphics[width=6cm]{ds_image_posthus_0_2.jpg}}\end{center}
%\par\begin{center}\footnotesize {Pósthús}\end{center}

\end{multicols}

% ---------------------------------------- List of abbreviations and symbols -

\chapter{Seznam použitých zkratek a symbolů}
\begin{multicols}{2}
\begin{description}[leftmargin=!,labelwidth=1.5cm]
%[\breaklabel\setleftmargin{40pt}\setlabelstyle{\bfseries}
%\addtolength{\itemsep}{-0.5\baselineskip}]
\item[{abb}] {zkrácený}
\item[{acc}] {4. pád, akuzativ}
\item[{adj}] {přídavné jméno}
\item[{adv}] {příslovce}
\item[{akt}] {rod činný, aktivum}
\item[{anat.}] {anatomie}
\item[{angl.}] {anglicky}
\item[{astro.}] {astronomie, výzkum vesmíru}
\item[{ap.}] {a podobně}
\item[{básn.}] {básnický, poetický výraz}
\item[{biol.}] {biologie}
\item[{bot.}] {botanika}

\item[{comp}] {2. stupeň stupňování, komparativ}
\item[{con}] {spojovací způsob, konjunktiv}
\item[{conj}] {spojka}
\item[{dat}] {3. pád, dativ}
\item[{def}] {určitý tvar}
\item[{dem}] {zájmeno ukazovací}
\item[{dět.}] {dětsky}

\item[{e-að}] {co (sg nom neživotný) (\textit{eitthvað})}
\item[{e-ð}] {co (sg acc neživotný) (\textit{eitthvað})}
\item[{e-ir}] {kdo (pl nom životný) (\textit{einhverjir})}
\item[{e-ja}] {koho (pl acc životný) (\textit{einhverja})}
\item[{e-jum}] {komu (pl dat životný) (\textit{einhverjum})}
\item[{e-m}] {komu (sg dat životný) (\textit{einhverjum})}
\item[{e-n}] {koho (sg acc životný) (\textit{einhvern})}
\item[{e-r}] {kdo (sg nom životný) (\textit{einhver})}
\item[{e-rra}] {koho (pl gen životný) (\textit{einhverra})}
\item[{e-rs}] {koho (sg gen životný) (\textit{einhvers})}
\item[{e-s}] {čeho (sg gen neživotný) (\textit{einhvers})}
\item[{e-u}] {čemu (sg dat neživotný) (\textit{einhverju})}

\item[{ekon.}] {ekonomika, obchod}
\item[{elek.}] {elektřina}
\item[{f}] {rod ženský}
\item[{filos.}] {filosofie, logika}
\item[{form.}] {formálně}
\item[{fyz.}] {fyzika}
\item[{gen}] {2. pád, genitiv}
\item[{geog.}] {zeměpis, geografie}
\item[{geol.}] {geologie}
\item[{han.}] {hanlivě, pejorativně}
\item[{hist.}] {historie}
\item[{hovor.}] {hovorově}
\item[{hrub.}] {hrubě}
\item[{hud.}] {hudba}
\item[{imper}] {rozkazovací způsob, imperativ}
\item[{impers}] {sloveso neosobní}
\item[{ind}] {oznamovací způsob, indikativ}
\item[{indecl}] {nesklonný}
\item[{indef}] {neurčitý tvar, (pron +) zájmeno neurčité}
\item[{inf}] {infinitiv}
\item[{int}] {tázací zájmeno}
\item[{inter}] {citoslovce}
\item[{chem.}] {chemie}
\item[{jaz.}] {jazykověda}
\item[{kulin.}] {kulinářství, vaření}
\item[{l.}] {latinsky}
\item[{let.}] {letectví}
\item[{lit.}] {literatura, vydavatelství}
\item[{m}] {rod mužský}
\item[{mat.}] {matematika}
\item[{med. }] {lékařství}
\item[{meteo.}] {meteorologie}
\item[{myt.}] {mytologie}
\item[{n}] {rod střední}
\item[{náb.}] {náboženství}
\item[{nám.}] {námořnictví, rybolov}
\item[{nom}] {1. pád, nominativ}
\item[{num}] {číslovka}
\item[{ord}] {řadová číslovka}
\item[{p}] {osoba}
\item[{part}] {částice}
\item[{pers}] {(v +) osoba, (pron +) zájmeno osobní}
\item[{pf}] {minulý čas}
\item[{pl}] {množné číslo, plurál}
\item[{poč.}] {informatika}
\item[{pol.}] {politika, politologie}
\item[{poss}] {zájmeno přivlastňovací}
\item[{pos}] {1. stupeň stupňování, pozitiv}
\item[{pov.}] {lidové pověsti, folkloristika}
\item[{pp}] {příčestí minulé}
\item[{praes}] {přítomný čas}
\item[{práv.}] {právnictví, soudnictví}
\item[{predp}] {předpona}
\item[{prep}] {předložka}
\item[{presp}] {příčestí přítomné}
\item[{pron}] {zájmeno}
\item[{prop}] {vlastní jméno, proprium}
\item[{přen.}] {přeneseně}
\item[{přís.}] {přísloví}
\item[{psych.}] {psychologie}
\item[{refl}] {(v +) médium, (pron +) zvratné zájmeno}
\item[{rel}] {vztažné zájmeno}
\item[{sg}] {jednotné číslo, singulár}
\item[{slang.}] {slang}
\item[{sport.}] {sport}
\item[{stav.}] {stavebnictví, architektura}
\item[{subs}] {podstatné jméno (bez rodu)}
\item[{sup}] {3. stupeň stupňování, superlativ}
\item[{supin}] {supinum}
\item[{škol.}] {školství}
\item[{techn.}] {technika, mechanika}
\item[{v}] {sloveso}
\item[{voj.}] {vojenství}
\item[{zast.}] {zastarale}
\item[{zkr}] {zkratka}
\item[{zool.}] {zoologie}

\item[{;}] {středník uprostřed definice odděluje věty; středník na konci definice odděluje významy}

\item[{(...)}] {v kulatých závorkách jsou uvedeny doplňkové informace}
\item[{[...]}] {v hranatých závorkách je uveden fonetický zápis výslovnosti}\footnote{Vysvětlení fonetických symbolů je uvedeno v kapitole Fonetika \ref{sec:phon_phonems} na straně \pageref{sec:phon_phonems}}
\item[{|}] {naznačuje místo, kde se deklinační koncovka připojuje ke slovnímu základu}
\item[{·}] {vedlejší dělení složeného slova}
\item[{··}] {hlavní dělení složeného slova}
\item[{+}] {plus}
\item[{/}] {lomítko označuje více možností}
\item[{\dicsymSee}] {šipka odkazuje na heslo ve slovníku}
\item[{\dicsymCompare}] {šipka odkazuje na heslo ve slovníku ke srovnání}
\item[{\dicsymExampleIS}] {uvozuje islandský příklad}
\item[{\dicsymExampleCS}] {uvozuje český překlad příkladu}
\item[{\dicsymIdiom}] {uvozuje slovní spojení}
\item[{\dicsymFrequent}] {heslo patří mezi frekventované výrazy}

\item[{\dicsymPhoto}] {fotografie, ilustrace}
\item[{\dicsymProverb}] {přísloví, rčení}

\end{description}
\end{multicols}

\chapter{Skrá yfir skammstafanir og tákn}
\begin{multicols}{2}
\begin{description}[leftmargin=!,labelwidth=1.5cm]
%[\breaklabel\setleftmargin{40pt}\setlabelstyle{\bfseries}
%\addtolength{\itemsep}{-0.5\baselineskip}]
\item[{abb}] {stýfður}
\item[{acc}] {þolfall}
\item[{adj}] {lýsingarorð}
\item[{adv}] {atviksorð}
\item[{akt}] {germynd}
\item[{anat.}] {líffærafræði}
\item[{angl.}] {enska}
\item[{astro.}] {geimvísindi, stjörnufræði}
\item[{ap.}] {e.þ.h.}
\item[{básn.}] {ljóðrænn}
\item[{biol.}] {líffræði}
\item[{bot.}] {grasafræði}

\item[{comp}] {miðstig}
\item[{con}] {viðtengingarháttur}
\item[{conj}] {samtenging}
\item[{dat}] {þágufall}
\item[{def}] {með greini}
\item[{dem}] {ábendingarfornafn}
\item[{dět.}] {barnamál}

\item[{e-að}] {(\textit{eitthvað})}
\item[{e-ð}] {(\textit{eitthvað})}
\item[{e-ir}] {(\textit{einhverjir})}
\item[{e-ja}] {(\textit{einhverja})}
\item[{e-jum}] {(\textit{einhverjum})}
\item[{e-m}] {(\textit{einhverjum})}
\item[{e-n}] {(\textit{einhvern})}
\item[{e-r}] { (\textit{einhver})}
\item[{e-rra}] {(\textit{einhverra})}
\item[{e-rs}] {(\textit{einhvers})}
\item[{e-s}] {(\textit{einhvers})}
\item[{e-u}] { (\textit{einhverju})}

\item[{ekon.}] {hagfræði, viðskipti}
\item[{elek.}] {rafmagn}
\item[{f}] {kvenkyn}
\item[{filos.}] {heimspeki, rökfræði}
\item[{form.}] {formlegt mál}
\item[{fyz.}] {eðlisfræði}
\item[{gen}] {eignarfall}
\item[{geog.}] {landafræði}
\item[{geol.}] {jarðfræði}
\item[{han.}] {niðrandi}
\item[{hist.}] {sagnfræði}
\item[{hovor.}] {talmál}
\item[{hrub.}] {gróft}
\item[{hud.}] {tónlist}
\item[{imper}] {boðháttur}
\item[{impers}] {ópersónuleg sögn}
\item[{ind}] {framsöguháttur}
\item[{indecl}] {óbeygjanlegur}
\item[{indef}] {óákveðinn, (pron +) óákveðið fornafn}
\item[{inf}] {nafnháttur}
\item[{int}] {spurnigarfornafn}
\item[{inter}] {upphrópun}
\item[{chem.}] {efnafræði}
\item[{jaz.}] {málfræði, málvísindi}
\item[{kulin.}] {eldhús}
\item[{l.}] {latína}
\item[{let.}] {flug, flugmál}
\item[{lit.}] {skáldskapur}
\item[{m}] {nafnorð, karlkyn}
\item[{mat.}] {stærðfræði}
\item[{med. }] {læknisfræði}
\item[{meteo.}] {veðurfræði}
\item[{myt.}] {goðafræði}
\item[{n}] {nafnorð, hvorugkyn}
\item[{náb.}] {trúfræði}
\item[{nám.}] {sjómennska}
\item[{nom}] {nefnifall}
\item[{num}] {töluorð}
\item[{ord}] {raðtala}
\item[{p}] {persóna}
\item[{part}] {ögn}
\item[{pers}] {(v +) persóna, (pron +) persónufornafn}
\item[{pf}] {þátíð}
\item[{pl}] {fleirtala}
\item[{poč.}] {tölvufræði, tölvur}
\item[{pol.}] {stjórnmál, stjórnmálafræði}
\item[{poss}] {eignarfornafn}
\item[{pos}] {frumstig}
\item[{pov.}] {þjóðtrú}
\item[{pp}] {lýsingarháttur þátíðar}
\item[{praes}] {nútíð}
\item[{práv.}] {lögfræði, dómsmál}
\item[{predp}] {forskeyti}
\item[{prep}] {forsetning}
\item[{presp}] {lýsingarháttur nútíðar}
\item[{pron}] {fornafn}
\item[{prop}] {sérnafn}
\item[{přen.}] {afleidd merking}
\item[{přís.}] {málsháttur}
\item[{psych.}] {sálfræði}
\item[{refl}] {(v +) miðmynd, (pron +) afturbeygt fornafn}
\item[{rel}] {tilvísunarfornafn}
\item[{sg}] {eintala}
\item[{slang.}] {slangur}
\item[{sport.}] {íþróttir}
\item[{stav.}] {húsasmíði, byggingarlist}
\item[{subs}] {nafnorð (án kyns)}
\item[{sup}] {efstastig}
\item[{supin}] {sagnbót}
\item[{škol.}] {skólamál}
\item[{techn.}] {tækni, verkfræði}
\item[{v}] {sagnorð}
\item[{voj.}] {her}
\item[{zast.}] {úrelt mál}
\item[{zkr}] {skammstöfun}
\item[{zool.}] {dýrafræði}

\item[{;}] {semíkomma í miðri skýringu aðgreinir setningu; semíkomma í enda skýringar aðgreinir mismunandi merkingarsvið}

\item[{(...)}] {innan sviga eru viðbótarupplýsingar}
\item[{[...]}] {innan hornklofa er hljóðritaður framburður}
\item[{|}] {merkir stað, þar sem beygingarending bætist við stofn}
\item[{·}] {aukasamskeyti samsetts orðs}
\item[{··}] {aðalsamskeyti samsetts orðs}
\item[{+}] {plús}
\item[{/}] {skástrik táknar fleiri möguleika}
\item[{\dicsymSee}] {ör vísar til uppflettiorðs í orðabókinni}
\item[{\dicsymCompare}] {ör vísar til uppflettiorðs í bókinni til samanburðar}
\item[{\dicsymExampleIS}] {kemur á undan dæmi á íslensku}
\item[{\dicsymExampleCS}] {kemur á undan tékkneskri þýðingu á dæminu}
\item[{\dicsymIdiom}] {orð með háa tíðni}

\item[{\dicsymPhoto}] {ljósmynd}
\item[{\dicsymProverb}] {málsháttur}

\end{description}
\end{multicols}

% ------------------------------------------------- Headword part and alphabet

\chapter{Heslová část}
{\renewcommand{\arraystretch}{2}%
%\topskip0pt
\vstretch{1}
{\setlength{\tabcolsep}{1.5em}
\begin{table}[htpb]
\centering

\begin{tabular}{ccccccc} \hline
A a &  & G g &  & O o &  & V v \\
Á á  &  & H h &  & Ó ó &  & W w \\
B b &  & I i &  & P p &  & X x \\
C c &  & Í í &  & Q q &  & Y y \\
D d &  & J j &  & R r &  & Ý ý \\
Ð ð &  & K k &  & S s &  & Z z \\
E e &  & L l &  & T t &  & Þ þ \\
É é  &  & M m &  & U u &  & Æ æ \\
F f &  & N n &  & Ú ú &  & Ö ö \\
\hline
\end{tabular}

%\caption{\textbf{Islandská abeceda}}
\label{my-islandska-abeceda}
\end{table}
}
\vstretch{1}
{\renewcommand{\arraystretch}{1}%

\clearpage

\fi

% =========================================================== THE DICTIONARY =

\ifinputletters

\cleardoublepage

\dictionarygeometry
\pagestyle{myheadings}

%\newcommand*{\alphabet}{%
%  a,aa,b,c,d,dd,e,ee,f,g,h,i,ii,j,k,l,m,
%  %n,o,oo,p,r,s,t,u,uu,v,w,x,y,yy,th,ae,oe
%  }

%\foreach \l in \alphabet{
%  \dicLetter{l}{letter15}
\dicEntry[l] \dicTerm{l} \dicPos{zkr} \dicPhraseIS{lítri} \dicDirectTranslationCS{litr}
\dicEntry[l.] \dicTerm{l.} \dicPos{zkr} \dicPhraseIS{lýsingarorð} \dicFieldCat{jaz.} \dicDirectTranslationCS{přídavné jméno}
\dicEntry[labb] \dicTerm{labb} \dicIPA{{l}{a}{\textsubring{b}}{\textlengthmark}} \dicPos{n}[2] \dicFlx{(‑s)}[2] \dicDirectTranslationCS{procházka, špacír}
\dicEntry[labba] \dicTerm{labb|a} \dicsymFrequent\  \dicIPA{{l}{a}{\textsubring{b}}{\textlengthmark}{a}} \dicPos{v}[1] \dicFlx{(‑aði)}[13] \dicDirectTranslationCS{projít se, procházet se, jít na procházku, špacírovat} \dicExampleIS{Ég labbaði í bæinn.} \dicExampleCS{Prošel jsem se do města.}
\dicEntry[labbakútur] \dicTerm{labba··kút|ur} \dicIPA{{l}{a}{\textsubring{b}}{\textlengthmark}{a}{k\smash{\textsuperscript{h}}}{u}{\textsubring{d}}{\textscy}{\textsubring{r}}} \dicPos{m}[6] \dicFlx{(‑s, ‑ar)}[22] \textbf{1.} \dicSynonym{barn} \dicDirectTranslationCS{klouček, prcek}  \textbf{2.} \dicSynonym*{mannleysa} \dicDirectTranslationCS{trdlo, janek}
\dicEntry[laða] \dicTerm{lað|a} \dicIPA{{l}{a}{\textlengthmark}{ð}{a}} \dicPos{v}[1] \dicFlx{(‑aði)}[13] \dicFlx{acc} \dicDirectTranslationCS{přitahovat, (při)lákat, (při)vábit} \dicExampleIS{laða e‑n að sér} \dicExampleCS{přilákat (koho) k~sobě};  \dicIdiom{laða}[fram]{ \dicPhraseIS{laða e‑u fram}} \dicDirectTranslationCS{vyvolat (co), přivolat (co), přivodit (co)};  \dicIdiom{laðast}{ \dicPhraseIS{laðast}} \dicFlx{refl} \dicDirectTranslationCS{být přitahován} \dicExampleIS{laðast að e‑u} \dicExampleCS{být přitahován (čím)}
\dicEntry[lafa] \dicTerm{laf|a} \dicIPA{{l}{a}{\textlengthmark}{v}{a}} \dicPos{v}[2] \dicFlx{(‑ði, ‑að)}[125] \textbf{1.} \dicSynonym{hanga} \dicDirectTranslationCS{viset} \dicExampleIS{láta fæturna lafa} \dicExampleCS{nechat nohy viset}  \textbf{2.} \dicDirectTranslationCS{prodlévat, zdržovat se} \dicExampleIS{lafa í þorpinu} \dicExampleCS{zdržovat se na vsi}
\dicEntry[lag] \dicTerm{lag} \dicsymFrequent\  \dicIPA{{l}{a}{\textlengthmark}{x}} \dicPos{n}[2] \dicFlx{(‑s, lög)}[8] \textbf{1.} \dicSynonym{samfella} \dicDirectTranslationCS{vrstva} \dicExampleIS{lag af kolum} \dicExampleCS{vrstva uhlí}  \textbf{2.} \dicDirectTranslationCS{(pro)bodnutí (zbraní ap.)} \dicExampleIS{Lagið nam í hjartastað.} \dicExampleCS{Bodnutí zasáhlo srdce.}  \textbf{3.} \dicSynonym{aðferð} \dicDirectTranslationCS{způsob, metoda};  \dicPhraseIS{finna lagið á e‑u} \dicDirectTranslationCS{přijít na to, jak se (co) dělá};  \dicPhraseIS{hafa lag á e‑u} \dicDirectTranslationCS{mít talent na (co)};  \dicPhraseIS{koma e‑m (upp) á lagið með e‑ð} \dicDirectTranslationCS{naučit (koho), jak se (co) dělá};  \dicPhraseIS{með þessu lagi} \dicFlx{adv} \dicDirectTranslationCS{takto, tímto způsobem}  \textbf{4.} \dicSynonym{magn} \dicDirectTranslationCS{množství};  \dicPhraseIS{í einu lagi} \dicFlx{adv} \dicSynonym*{í heild} \dicDirectTranslationCS{najednou, naráz};  \dicPhraseIS{í fyrsta lagi} \dicFlx{adv} {\textbf{a.}} \dicDirectTranslationCS{ne dříve než} \dicExampleIS{Ég kem í fyrsta lagi kl. þrjú.} \dicExampleCS{Nepřijdu dřív než ve tři.};  {\textbf{b.}} \dicDirectTranslationCS{za prvé, v~první řadě};  \dicPhraseIS{í hæsta lagi} \dicFlx{adv} \dicDirectTranslationCS{nanejvýš, maximálně, nejvýše} \dicExampleIS{Þetta tekur í hæsta lagi fimm mínútur.} \dicExampleCS{Zabere to nejvýše pět minut.};  \dicPhraseIS{í meira lagi} \dicFlx{adv} \dicDirectTranslationCS{dosti, značně};  \dicPhraseIS{í mesta lagi} \dicFlx{adv} \dicDirectTranslationCS{ne více než};  \dicPhraseIS{í öðru lagi} \dicFlx{adv} \dicDirectTranslationCS{za druhé}  \textbf{5.} \dicSynonym{regla} \dicDirectTranslationCS{řád, pořádek};  \dicPhraseIS{allt í lagi!} \dicDirectTranslationCS{ok!, všechno v~pořádku!};  \dicPhraseIS{e‑að er fjarri lagi} \dicDirectTranslationCS{(co) není zdaleka v~pořádku};  \dicPhraseIS{koma e‑u í lag} \dicSynonym{lagfæra} \dicDirectTranslationCS{dát (co) do pořádku};  \dicPhraseIS{nærri lagi} \dicFlx{adv} \dicDirectTranslationCS{téměř v~pořádku}  \textbf{6.} \dicSynonym{tónverk} \dicDirectTranslationCS{píseň, písnička, melodie} \dicExampleIS{lag eftir e‑n} \dicExampleCS{(čí) písnička};  \dicPhraseIS{fara út af laginu} \dicDirectTranslationCS{vypadnout z~rytmu (v~písni), zazpívat falešný tón};  \dicPhraseIS{taka lagið} \dicDirectTranslationCS{(začít) zpívat}  \textbf{7.} \dicSynonym{verðlag} \dicDirectTranslationCS{cena, hodnota};  \dicPhraseIS{leggja lag á e‑ð} \dicDirectTranslationCS{ocenit (co), ohodnotit (co), určit cenu (čeho)}  \textbf{8.} \dicSynonym{samneyti} \dicDirectTranslationCS{společnost, přítomnost} \dicIndirectTranslationCS{(společenský styk)};  \dicPhraseIS{leggja lag sitt við e‑n} \dicDirectTranslationCS{dělat (komu) společnost}  \textbf{9.} \dicPhraseIS{lög} \dicFlx{pl} \dicLink{lög};  \dicIdiom{lag}{ \dicPhraseIS{slá e‑n út af laginu}} \dicDirectTranslationCS{vyvést (koho) z~míry, vytočit (koho)}; { \dicPhraseIS{sæta lagi}} \dicDirectTranslationCS{počkat si (na příležitost)}
\dicEntry[laga] \dicTerm{lag|a} \dicsymFrequent\  \dicIPA{{l}{a}{\textlengthmark}{\textbabygamma}{a}} \dicPos{v}[1] \dicFlx{(‑aði)}[13] \dicFlx{acc} \textbf{1.} \dicSynonym{prýða} \dicDirectTranslationCS{upravit, upravovat, urovnat, urovnávat} \dicExampleIS{laga á sér hárið} \dicExampleCS{upravit si vlasy}  \textbf{2.} \dicSynonym*{koma e‑u í lag} \dicDirectTranslationCS{opravit, opravovat, spravit, spravovat, dát\,/\addthin dávat do pořádku} \dicExampleIS{laga vélina} \dicExampleCS{opravit stroj}  \textbf{3.} \dicDirectTranslationCS{připravit, připravovat (kávu ap.)} \dicExampleIS{laga kaffi og te} \dicExampleCS{připravit kávu a~čaj}  \textbf{4.} \dicSynonym{streyma} \dicDirectTranslationCS{proudit, prýštit (krev z~rány ap.)};  \dicIdiom{laga}[að]{ \dicPhraseIS{laga sig að e‑u}} \dicSynonym{aðhæfa} \dicDirectTranslationCS{přizpůsobit se (čemu)};  \dicIdiom{laga}[til]{ \dicPhraseIS{laga til}} \dicDirectTranslationCS{uklízet, poklízet, uklidit} \dicExampleIS{laga til í íbúðinni} \dicExampleCS{uklidit v~bytě}; { \dicPhraseIS{laga sig til}} \dicDirectTranslationCS{poupravit se};  \dicIdiom{lagast}{ \dicPhraseIS{lagast}} \dicFlx{refl} {\textbf{a.}} \dicDirectTranslationCS{(z)lepšit se, vylepšit se} \dicExampleIS{Ástandið lagast af sjálfu sér.} \dicExampleCS{Situace se zlepšuje sama od sebe.};  {\textbf{b.}} \dicDirectTranslationCS{ustoupit, povolit (bolest ap.)}
\dicEntry[lagaákvæði] \dicTerm{laga··á·kvæði} \dicIPA{{l}{a}{\textlengthmark}{\textbabygamma}{a}{au}{k\smash{\textsuperscript{h}}}{v}{a}{i}{ð}{\textsci}} \dicPos{n}[2] \dicFlx{(‑s, ‑)}[14] \dicFieldCat{práv.} \dicDirectTranslationCS{právní předpis\,/\addthin ustanovení}
\dicEntry[lagabálkur] \dicTerm{laga··bálk|ur} \dicIPA{{l}{a}{\textlengthmark}{\textbabygamma}{a}{\textsubring{b}}{au}{\textsubring{l}}{\r{g}}{\textscy}{\textsubring{r}}} \dicPos{m}[6] \dicFlx{(‑s, ‑ar)}[24] \dicFieldCat{práv.} \dicDirectTranslationCS{kodex, zákoník}
\dicEntry[lagabreyting] \dicTerm{laga··breyt·ing} \dicIPA{{l}{a}{\textlengthmark}{\textbabygamma}{a}{\textsubring{b}}{r}{ei}{\textsubring{d}}{i}{\ng}{\r{g}}} \dicPos{f}[4] \dicFlx{(‑ar, ‑ar)}[5] \dicFieldCat{práv.} \dicDirectTranslationCS{novela zákona}
\dicEntry[lagabrot] \dicTerm{laga··brot} \dicIPA{{l}{a}{\textlengthmark}{\textbabygamma}{a}{\textsubring{b}}{r}{\textopeno}{\textsubring{d}}} \dicPos{n}[2] \dicFlx{(‑s, ‑)}[5] \dicDirectTranslationCS{porušení zákona}
\dicEntry[lagadeild] \dicTerm{laga··deild} \dicIPA{{l}{a}{\textlengthmark}{\textbabygamma}{a}{\textsubring{d}}{ei}{l}{\textsubring{d}}} \dicPos{f}[7] \dicFlx{(‑ar, ‑ir)}[1] \dicDirectTranslationCS{právnická fakulta}
\dicEntry[lagafrumvarp] \dicTerm{laga··frum·|varp} \dicIPA{{l}{a}{\textlengthmark}{\textbabygamma}{a}{f}{r}{\textscy}{m}{v}{a}{\textsubring{r}}{\textsubring{b}}} \dicPos{n}[2] \dicFlx{(‑varps, ‑vörp)}[8] \dicFieldCat{práv.} \dicDirectTranslationCS{návrh zákona}
\dicEntry[lagagrein] \dicTerm{laga··grein} \dicIPA{{l}{a}{\textlengthmark}{\textbabygamma}{a}{\r{g}}{r}{ei}{\textsubring{n}}} \dicPos{f}[4] \dicFlx{(‑ar, ‑ar)}[1] \dicFieldCat{práv.} \dicSynonym{grein} \dicDirectTranslationCS{odstavec zákona, paragraf}
\dicEntry[lagaheimild] \dicTerm{laga··heimild} \dicIPA{{l}{a}{\textlengthmark}{\textbabygamma}{a}{h}{ei}{m}{\textsci}{l}{\textsubring{d}}} \dicPos{f}[7] \dicFlx{(‑ar, ‑ir)}[1] \dicDirectTranslationCS{moc zákona, zákonnost}
\dicEntry[lagahöfundur] \dicTerm{laga··höf·und|ur} \dicIPA{{l}{a}{\textlengthmark}{\textbabygamma}{a}{h}{\oe}{v}{\textscy}{n}{\textsubring{d}}{\textscy}{\textsubring{r}}} \dicPos{m}[6] \dicFlx{(‑ar, ‑ar)}[56] \dicDirectTranslationCS{písničkář(ka), skladatel(ka)}
\dicEntry[lagalegur] \dicTerm{laga··legur} \dicIPA{{l}{a}{\textlengthmark}{\textbabygamma}{a}{l}{\textepsilon}{\textbabygamma}{\textscy}{\textsubring{r}}} \dicPos{adj}[1]\dicFlx{}[-8] \dicDirectTranslationCS{zákonný, právní}
\dicEntry[laganir] \dicTerm{laganir} \dicIPA{{l}{a}{\textlengthmark}{\textbabygamma}{a}{n}{\textsci}{\textsubring{r}}} \dicPos{f} \dicFlx{pl nom} \dicLink{lögun}
\dicEntry[lagar] \dicTerm{lagar} \dicIPA{{l}{a}{\textlengthmark}{\textbabygamma}{a}{\textsubring{r}}} \dicPos{m} \dicFlx{sg gen} \dicLink{lögur}
\dicEntry[lagaregla] \dicTerm{laga··regl|a} \dicIPA{{l}{a}{\textlengthmark}{\textbabygamma}{a}{r}{\textepsilon}{\r{g}}{l}{a}} \dicPos{f}[1] \dicFlx{(‑u, ‑ur)}[19] \dicFieldCat{práv.} \dicDirectTranslationCS{právní předpis}
\dicEntry[lagasafn] \dicTerm{laga··|safn} \dicIPA{{l}{a}{\textlengthmark}{\textbabygamma}{a}{s}{a}{\textsubring{b}}{\textsubring{n}}} \dicPos{n}[2] \dicFlx{(‑safns, ‑söfn)}[8] \dicFieldCat{práv.} \dicDirectTranslationCS{sbírka zákonů}
\dicEntry[lager] \dicTerm{lager} \dicIPA{{l}{a}{\textlengthmark}{\textbabygamma}{\textepsilon}{\textsubring{r}}} \dicPos{m}[4] \dicFlx{(‑s, ‑ar)}[14] \textbf{1.} \dicSynonym{vörubirgðir} \dicDirectTranslationCS{sklad, zásoba} \dicExampleIS{Vörur eru sendar samdægurs séu þær til á lager.} \dicExampleCS{Zboží je odesíláno ve stejný den, pokud je na skladě.};  \dicPhraseIS{hafa e‑ð á lager} \dicDirectTranslationCS{mít (co) na skladě}  \textbf{2.} \dicSynonym{vörugeymsla} \dicDirectTranslationCS{sklad, skladiště} \dicExampleIS{vinna á lager} \dicExampleCS{pracovat ve skladu}
\dicEntry[lagermaður] \dicTerm{lager··|maður} \dicIPA{{l}{a}{\textlengthmark}{\textbabygamma}{\textepsilon}{r}{m}{a}{ð}{\textscy}{\textsubring{r}}} \dicPos{m}[13] \dicFlx{(‑manns, ‑menn)}[2] \dicDirectTranslationCS{skladník, skladnice}
\dicEntry[lagfæra] \dicTerm{lag··fær|a} \dicIPA{{l}{a}{x}{f}{a}{i}{r}{a}} \dicPos{v}[2] \dicFlx{(‑ði, ‑t)}[105] \dicFlx{acc} \textbf{1.} \dicDirectTranslationCS{dát\,/\addthin uvést do pořádku, opravit} \dicExampleIS{lagfæra skemmdir á húsi} \dicExampleCS{dát do pořádku poškození na domě}  \textbf{2.} \dicSynonym{endurbæta} \dicDirectTranslationCS{zlepšit, vylepšit, uzpůsobit}
\dicEntry[lagfæring] \dicTerm{lag··fær·ing} \dicIPA{{l}{a}{x}{f}{a}{i}{r}{i}{\ng}{\r{g}}} \dicPos{f}[4] \dicFlx{(‑ar, ‑ar)}[5] \dicSynonym{endurbót} \dicDirectTranslationCS{oprava, náprava, korekce}
\dicEntry[laggar] \dicTerm{laggar} \dicIPA{{l}{a}{\r{g}}{\textlengthmark}{a}{\textsubring{r}}} \dicPos{f} \dicFlx{sg gen} \dicLink{lögg}
\dicEntry[laggir] \dicTerm{laggir} \dicIPA{{l}{a}{\r{\textObardotlessj}}{\textlengthmark}{\textsci}{\textsubring{r}}} \dicPos{f} \dicFlx{pl nom} \dicLink{lögg}
\dicEntry[laggóður] \textls[15]{\dicTerm{lag··|góður} \dicIPA{{l}{a}{\textbabygamma}{\r{g}}{ou}{ð}{\textscy}{\textsubring{r}}} \dicPos{adj}[11] \dicFlx{(comp ‑betri, sup ‑bestur)}[1] \dicSynonym{hnyttinn} \dicDirectTranslationCS{výstižný, trefný} \dicExampleIS{stutt og laggott svar} \dicExampleCS{krátká a~výstižná odpověď}}
\dicEntry[laghendur] \dicTerm{lag··hendur}\dicTerm{, laghentur} \dicIPA{{l}\-{a}\-{\textlengthmark}\-{x}\-{h}\-{\textepsilon}\-{n}\-{\textsubring{d}}\-{\textscy}\-{\textsubring{r}}\-} \dicPos{adj}[2]\dicFlx{}[-14] \dicSynonym{handlaginn} \dicDirectTranslationCS{obratný, zručný}
\dicEntry[laghentur] \dicTerm{lag··hentur} \dicIPA{{l}{a}{\textlengthmark}{x}{h}{\textepsilon}{\textsubring{n}}{\textsubring{d}}{\textscy}{\textsubring{r}}} \dicPos{adj}[1]\dicFlx{}[-10] \dicLink{laghendur}
\dicEntry[laginn] \dicTerm{laginn} \dicIPA{{l}{a}{i}{j}{\textlengthmark}{\textsci}{\textsubring{n}}} \dicPos{adj}[6]\dicFlx{}[-3] \dicDirectTranslationCS{zručný, šikovný};  \dicPhraseIS{e‑r er laginn í höndunum} \dicDirectTranslationCS{(kdo) má šikovné ruce};  \dicPhraseIS{e‑r er laginn við e‑n} \dicDirectTranslationCS{(kdo) má dobrý přístup ke (komu) (zvířatům ap.)}
\dicEntry[lagkaka] \dicTerm{lag··|kaka} \dicIPA{{l}{a}{\textlengthmark}{\textbabygamma}{k\smash{\textsuperscript{h}}}{a}{\r{g}}{a}} \dicPos{f}[1] \dicFlx{(‑köku, ‑kökur)}[20] \dicFieldCat{kulin.} \dicDirectTranslationCS{řezy (čokoládové ap.), vrstvený koláč, dortík, zákusek}
\dicEntry[laglaus] \dicTerm{lag··laus} \dicIPA{{l}{a}{\textbabygamma}{l}{\oe i}{s}} \dicPos{adj}[5]\dicFlx{}[-1] \dicDirectTranslationCS{(jsoucí) bez hudebního sluchu, hluchý (k~hudbě ap.)}
\dicEntry[laglegur] \dicTerm{lag··legur} \dicsymFrequent\  \dicIPA{{l}{a}{\textbabygamma}{l}{\textepsilon}{\textbabygamma}{\textscy}{\textsubring{r}}} \dicPos{adj}[1]\dicFlx{}[-8] \textbf{1.} \dicSynonym{fríður} \dicDirectTranslationCS{pohledný, půvabný} \dicExampleIS{lagleg stúlka} \dicExampleCS{pohledná dívka}  \textbf{2.} \dicSynonym{snotur} \dicDirectTranslationCS{krásný, hezký}  \textbf{3.} \dicDirectTranslationCS{pěkný, pořádný} \dicIndirectTranslationCS{(k~zdůraznění)}
\dicEntry[laglína] \dicTerm{lag··lín|a} \dicIPA{{l}{a}{\textbabygamma}{l}{i}{n}{a}} \dicPos{f}[1] \dicFlx{(‑u, ‑ur)}[7] \dicDirectTranslationCS{melodie}
\dicEntry[lagnar] \dicTerm{lagnar} \dicIPA{{l}{a}{\r{g}}{n}{a}{\textsubring{r}}} \dicPos{f} \dicFlx{sg gen} \dicLink{lögn}
\dicEntry[lagni] \dicTerm{lagn|i} \dicIPA{{l}{a}{\r{g}}{n}{\textsci}} \dicPos{f}[3] \dicFlx{(‑i)}[3] \dicSynonym{leikni} \dicDirectTranslationCS{šikovnost, zručnost, obratnost}
\dicEntry[lagning] \dicTerm{lag··ning} \dicIPA{{l}{a}{\r{g}}{n}{i}{\ng}{\r{g}}} \dicPos{f}[4] \dicFlx{(‑ar, ‑ar)}[5] \textbf{1.} \dicDirectTranslationCS{kladení, pokládání, instalace} \dicExampleIS{lagning vegar} \dicExampleCS{kladení cesty}  \textbf{2.} \dicSynonym*{hárlagning} \dicDirectTranslationCS{úprava\,/\addthin upravení vlasů}
\dicEntry[lagnir] \dicTerm{lagnir} \dicIPA{{l}{a}{\r{g}}{n}{\textsci}{\textsubring{r}}} \dicPos{f} \dicFlx{pl nom} \dicLink{lögn}
\dicEntry[lagskipting] \dicTerm{lag··skipt·ing} \dicIPA{{l}{a}{x}{s}{\r{\textObardotlessj}}{\textsci}{f}{\textsubring{d}}{i}{\ng}{\r{g}}} \dicPos{f}[4] \dicFlx{(‑ar)}[7] \dicDirectTranslationCS{(roz)vrstvení}
\dicEntry[lagsmaður] \dicTerm{lags··|maður} \dicIPA{{l}{a}{x}{s}{m}{a}{ð}{\textscy}{\textsubring{r}}} \dicPos{m}[13] \dicFlx{(‑manns, ‑menn)}[2] \dicSynonym{félagi} \dicDirectTranslationCS{kamarád(ka), přítel(kyně), společník, společnice}
\dicEntry[lagviss] \dicTerm{lag··viss} \dicIPA{{l}{a}{\textbabygamma}{v}{\textsci}{s}} \dicPos{adj}[5]\dicFlx{}[-1] \dicDirectTranslationCS{zpívající čistě, muzikální}
\dicEntry[lak] \dicTerm{lak\smash{\textsuperscript{1}}} \dicIPA{{l}{a}{\textlengthmark}{\r{g}}} \dicPos{n}[2] \dicFlx{(‑s, lök)}[8] \dicDirectTranslationCS{prostěradlo}
\dicEntry[lak] \dicTerm{lak\smash{\textsuperscript{2}}} \dicIPA{{l}{a}{\textlengthmark}{\r{g}}} \dicPos{v} \dicFlx{ind pf sg 1 pers} \dicLink{leka}
\dicEntry[lakk] \dicTerm{lakk} \dicIPA{{l}{a}{h}{\r{g}}} \dicPos{n}[2] \dicFlx{(‑s, lökk)}[8] \textbf{1.} \dicSynonym*{lakkmálning} \dicDirectTranslationCS{lak, nátěr}  \textbf{2.} \dicSynonym{innsigli} \dicDirectTranslationCS{pečeť}  \textbf{3.} \dicSynonym*{innsiglislakk} \dicDirectTranslationCS{pečetní vosk}
\dicEntry[lakka] \dicTerm{lakk|a} \dicIPA{{l}{a}{h}{\r{g}}{a}} \dicPos{v}[1] \dicFlx{(‑aði)}[13] \dicFlx{acc} \textbf{1.} \dicDirectTranslationCS{(na)lakovat} \dicExampleIS{lakka gluggana} \dicExampleCS{nalakovat okna}  \textbf{2.} \dicDirectTranslationCS{zapečetit} \dicExampleIS{lakka bréf} \dicExampleCS{zapečetit dopis}
\dicEntry[lakkrís] \dicTerm{lakk··rís} \dicIPA{{l}{a}{h}{\r{g}}{r}{i}{s}} \dicPos{m}[4] \dicFlx{(‑s)}[18] \dicDirectTranslationCS{lékořice} \dicExampleIS{fá einn poka af lakkrís} \dicExampleCS{dostat jeden pytlík lékořice}
\dicEntry[lakkrísrót] \dicTerm{lakk·rís··|rót} \dicIPA{{l}{a}{h}{\r{g}}{r}{i}{s}{r}{ou}{\textsubring{d}}} \dicPos{f}[8] \dicFlx{(‑rótar, ‑rætur)}[5] \dicFieldCat{bot.} \dicDirectTranslationCS{kořen lékořice lysé} \textit{(l.~{\textLA{Glycyrrhiza glabra}})}
\dicEntry[laktósi] \dicTerm{laktós|i} \dicIPA{{l}{a}{x}{\textsubring{d}}{ou}{s}{\textsci}} \dicPos{m}[1] \dicFlx{(‑a)}[3] \dicFieldCat{chem.} \dicSynonym*{mjólkursykur} \dicDirectTranslationCS{laktóza, mléčný cukr}
\dicEntry[lakur] \dicTerm{lakur} \dicIPA{{l}{a}{\textlengthmark}{\r{g}}{\textscy}{\textsubring{r}}} \dicPos{adj}[1] \dicFlx{(f lök)}[2] \textbf{1.} \dicSynonym{lélegur} \dicDirectTranslationCS{chabý, mizerný (výsledek ap.)} \dicExampleIS{lakur árangur} \dicExampleCS{mizerný výsledek}  \textbf{2.} \dicSynonym{lasinn} \dicDirectTranslationCS{chorý, neduživý}  \textbf{3.} \dicSynonym{tæpur} \dicDirectTranslationCS{necelý, pouhý}
\dicEntry[lama] \dicTerm{lam|a} \dicIPA{{l}{a}{\textlengthmark}{m}{a}} \dicPos{v}[1] \dicFlx{(‑aði)}[13] \dicDirectTranslationCS{ochromit, paralyzovat} \dicExampleIS{Sjúkdómur lamar líkamann.} \dicExampleCS{Nemoc ochromuje tělo.};  \dicIdiom{lamast}{ \dicPhraseIS{lamast}} \dicFlx{refl} {\textbf{a.}} \dicDirectTranslationCS{ochrnout (o~člověku ap.)} \dicExampleIS{lamast við slysið} \dicExampleCS{ochrnout při nehodě};  {\textbf{b.}} \dicDirectTranslationCS{být paralyzován\,/\addthin ochromen (dopravní ruch ap.)}
\dicEntry[lamadýr] \dicTerm{lama··dýr} \dicIPA{{l}{a}{\textlengthmark}{m}{a}{\textsubring{d}}{i}{\textsubring{r}}} \dicPos{n}[2] \dicFlx{(‑s, ‑)}[5] \dicFieldCat{zool.} \dicDirectTranslationCS{lama} \textit{(l.~{\textLA{Lama}})}  \dicsymPhoto\ 
\dicFigure{ds_image_lamadyr_0_1.jpg}{Lamadýr}{Lamadýr - Luca Galuzzi, CC BY-SA 2.5}
\dicEntry[lamaður] \dicTerm{lam··aður} \dicIPA{{l}{a}{\textlengthmark}{m}{a}{ð}{\textscy}{\textsubring{r}}} \dicPos{adj}[3] \dicFlx{(f lömuð)}[1] \dicDirectTranslationCS{chromý, ochrnutý}
\dicEntry[lamasess] \dicTerm{lama··sess} \dicIPA{{l}{a}{\textlengthmark}{m}{a}{s}{\textepsilon}{s}} \dicPos{m}[4] \dicFlx{(‑)}[34] \dicPhraseIS{liggja\,/\addthin vera í lamasessi} \dicDirectTranslationCS{být mimo provoz, nefungovat (o~věcech)}
\dicEntry[lamb] \dicTerm{lamb} \dicsymFrequent\  \dicIPA{{l}{a}{m}{\textsubring{b}}} \dicPos{n}[2] \dicFlx{(‑s, lömb)}[8] \textbf{1.} \dicDirectTranslationCS{jehně} \dicExampleIS{nýfædd lömb} \dicExampleCS{novorozená jehňata}  \textbf{2.} \dicDirectTranslationCS{drahoušek, miláček};  \dicPhraseIS{launa e‑m lambið gráa} \dicLangCat{přen.} \dicDirectTranslationCS{oplatit (komu) stejnou mincí}
\dicEntry[lambagras] \dicTerm{lamba··|gras} \dicIPA{{l}{a}{m}{\textsubring{b}}{a}{\r{g}}{r}{a}{s}} \dicPos{n}[2] \dicFlx{(‑grass, ‑grös)}[8] \dicFieldCat{bot.} \dicDirectTranslationCS{silenka bezlodyžná} \textit{(l.~{\textLA{Silene acaulis}})}  \dicsymPhoto\ 
\dicFigure{9260.jpg}{Lambagras}{Lambagras - Daněk Pavel, Biolib, Copyright/CC-BY-SA}
\dicEntry[lambakjöt] \dicTerm{lamba··kjöt} \dicIPA{{l}{a}{m}{\textsubring{b}}{a}{c\smash{\textsuperscript{h}}}{\oe}{\textsubring{d}}} \dicPos{n}[2] \dicFlx{(‑s)}[2] \dicDirectTranslationCS{jehněčí (maso)}
\dicEntry[lambalæri] \dicTerm{lamba··læri} \dicIPA{{l}{a}{m}{\textsubring{b}}{a}{l}{a}{i}{r}{\textsci}} \dicPos{n}[2] \dicFlx{(‑s, ‑)}[14] \dicDirectTranslationCS{jehněčí kýta}
\dicEntry[lamdi] \dicTerm{lamdi} \dicIPA{{l}{a}{m}{\textsubring{d}}{\textsci}} \dicPos{v} \dicFlx{ind pf sg 1 pers} \dicLink{lemja}
\dicEntry[lamið] \dicTerm{lamið} \dicIPA{{l}{a}{\textlengthmark}{m}{\textsci}{\texttheta}} \dicPos{v} \dicFlx{supin} \dicLink{lemja}
\dicEntry[lampi] \dicTerm{lamp|i} \dicsymFrequent\  \dicIPA{{l}{a}{\textsubring{m}}{\textsubring{b}}{\textsci}} \dicPos{m}[1] \dicFlx{(‑a, ‑ar)}[8] \textbf{1.} \dicDirectTranslationCS{lampa, svítilna} \dicExampleIS{kveikja á lampa} \dicExampleCS{zapálit lampu}  \textbf{2.} \dicFieldCat{elek.} \dicDirectTranslationCS{výbojka}
\dicEntry[land] \dicTerm{land} \dicsymFrequent\  \dicIPA{{l}{a}{n}{\textsubring{d}}} \dicPos{n}[2] \dicFlx{(‑s, lönd)}[8] \textbf{1.} \dicSynonym{þurrlendi} \dicDirectTranslationCS{souš, pevnina, země} \dicExampleIS{á sjó og landi} \dicExampleCS{na moři a~na pevnině}  \textbf{2.} \dicSynonym{árbakki} \dicDirectTranslationCS{břeh} \dicExampleIS{á öðru landinu} \dicExampleCS{na druhém břehu}  \textbf{3.} \dicSynonym{ríki} \dicDirectTranslationCS{země, stát, vlast} \dicExampleIS{höfuðborg landsins} \dicExampleCS{hlavní město země}  \textbf{4.} \dicSynonym{landareign} \dicDirectTranslationCS{země, půda, pozemek};  \dicIdiom{land}{ \dicPhraseIS{e‑að á langt í land}} \dicLangCat{přen.} \dicDirectTranslationCS{(co) je ještě daleko (cíl ap.)}; { \dicPhraseIS{eiga langt í land}} \dicLangCat{přen.} \dicDirectTranslationCS{mít ještě hodně před sebou (větší část práce ap.)}; { \dicPhraseIS{draga í land}} \dicLangCat{přen.} \dicDirectTranslationCS{ustoupit (ze zásad ap.), rezignovat, upustit (od plánů ap.)}; { \dicPhraseIS{draga e‑n að landi}} \dicDirectTranslationCS{dojíst po (kom)}; { \dicPhraseIS{fara með löndum}} \dicLangCat{přen.} \dicDirectTranslationCS{postupovat\,/\addthin jít opatrně}; { \dicPhraseIS{láta e‑ð lönd og leið}} \dicLangCat{přen.} \dicDirectTranslationCS{nestarat se o~(co), nechat (co) být}; { \dicPhraseIS{leggja land undir fót}} \dicDirectTranslationCS{jít pěšky, pochodovat, kráčet}
\dicEntry[landa] \dicTerm{land|a} \dicIPA{{l}{a}{n}{\textsubring{d}}{a}} \dicPos{v}[1] \dicFlx{(‑aði)}[13] \dicFlx{dat} \dicDirectTranslationCS{vyložit, vylodit} \dicExampleIS{landa aflanum} \dicExampleCS{vyložit úlovek}
\dicEntry[landabréf] \dicTerm{landa··bréf} \dicIPA{{l}{a}{n}{\textsubring{d}}{a}{\textsubring{b}}{r}{j}{\textepsilon}{f}} \dicPos{n}[2] \dicFlx{(‑s, ‑)}[5] \dicSynonym{landakort} \dicDirectTranslationCS{mapa (zeměpisná ap.)}
\dicEntry[landafræði] \dicTerm{landa··fræð|i} \dicIPA{{l}{a}{n}{\textsubring{d}}{a}{f}{r}{a}{i}{ð}{\textsci}} \dicPos{f}[3] \dicFlx{(‑i)}[3] \dicDirectTranslationCS{zeměpis, geografie}
\dicEntry[landakort] \dicTerm{landa··kort} \dicIPA{{l}{a}{n}{\textsubring{d}}{a}{k\smash{\textsuperscript{h}}}{\textopeno}{\textsubring{r}}{\textsubring{d}}} \dicPos{n}[2] \dicFlx{(‑s, ‑)}[5] \dicDirectTranslationCS{mapa}
\dicEntry[landakortabók] \dicTerm{landa·korta··|bók} \dicIPA{{l}{a}{n}{\textsubring{d}}{a}{k\smash{\textsuperscript{h}}}{\textopeno}{\textsubring{r}}{\textsubring{d}}{a}{\textsubring{b}}{ou}{\r{g}}} \dicPos{f}[8] \dicFlx{(‑bókar, ‑bækur)}[5] \dicSynonym{atlas} \dicDirectTranslationCS{atlas}
\dicEntry[landalda] \dicTerm{land··|alda} \dicIPA{{l}{a}{n}{\textsubring{d}}{a}{l}{\textsubring{d}}{a}} \dicPos{f}[1] \dicFlx{(‑öldu, ‑öldur)}[20] \dicDirectTranslationCS{pobřežní vlna}
\dicEntry[landamerki] \dicTerm{landa··merki} \dicIPA{{l}{a}{n}{\textsubring{d}}{a}{m}{\textepsilon}{\textsubring{r}}{\r{\textObardotlessj}}{\textsci}} \dicPos{n}[2] \dicFlx{(‑s, ‑)}[16] \dicSynonym{landamæri} \dicDirectTranslationCS{pomezí, hranice (státní ap.)}
\dicEntry[landamæri] \dicTerm{landa··mæri} \dicIPA{{l}{a}{n}{\textsubring{d}}{a}{m}{a}{i}{r}{\textsci}} \dicPos{n}[2] \dicFlx{pl}[19] \dicDirectTranslationCS{(státní) hranice} \dicExampleIS{vera kominn yfir landamærin} \dicExampleCS{překročit státní hranici}
\dicEntry[landareign] \dicTerm{landar··eign} \dicIPA{{l}{a}{n}{\textsubring{d}}{a}{r}{ei}{\r{g}}{\textsubring{n}}} \dicPos{f}[7] \dicFlx{(‑ar, ‑ir)}[1] \dicDirectTranslationCS{pozemek, půda, pozemkové vlastnictví}
\dicEntry[landbúnaðarframleiðsla] \dicTerm{land·búnaðar··fram·leiðsl|a} \dicIPA{{l}\-{a}\-{n}\-{\textsubring{d}}\-{\textsubring{b}}\-{u}\-{n}\-{a}\-{ð}\-{a}\-{\textsubring{r}}\-{f}\-{r}\-{a}\-{m}\-{l}\-{ei}\-{ð}\-{s}\-{\textsubring{d}}\-{l}\-{a}\-} \dicPos{f}[1] \dicFlx{(‑u)}[5] \dicDirectTranslationCS{zemědělská produkce\,/\addthin výroba}
\dicEntry[landbúnaðarráðherra] \dicTerm{land·búnaðar··ráð·herr|a} \dicIPA{{l}\-{a}\-{n}\-{\textsubring{d}}\-{\textsubring{b}}\-{u}\-{n}\-{a}\-{ð}\-{a}\-{r}\-{au}\-{\texttheta}\-{h}\-{\textepsilon}\-{r}\-{a}\-} \dicPos{m}[1] \dicFlx{(‑a, ‑ar)}[17] \dicDirectTranslationCS{ministr(yně) zemědělství}
\dicEntry[landbúnaðarráðuneyti] \textls[15]{\dicTerm{land·búnaðar··ráðu·neyti} \dicIPA{{l}\-{a}\-{n}\-{\textsubring{d}}\-{\textsubring{b}}\-{u}\-{n}\-{a}\-{ð}\-{a}\-{r}\-{au}\-{ð}\-{\textscy}\-{n}\-{ei}\-{\textsubring{d}}\-{\textsci}}} \dicPos{n}[2] \dicFlx{(‑s, ‑)}[14] \dicDirectTranslationCS{ministerstvo zemědělství}
\dicEntry[landbúnaður] \dicTerm{land··bú·nað|ur} \dicIPA{{l}{a}{n}{\textsubring{d}}{\textsubring{b}}{u}{n}{a}{ð}{\textscy}{\textsubring{r}}} \dicPos{m}[10] \dicFlx{(‑ar)}[9] \dicDirectTranslationCS{zemědělství}
\dicEntry[landeigandi] \dicTerm{land··eig·|andi} \dicIPA{{l}{a}{n}{\textsubring{d}}{ei}{\textbabygamma}{a}{n}{\textsubring{d}}{\textsci}} \dicPos{m}[2] \dicFlx{(‑anda, ‑endur)}[1] \dicDirectTranslationCS{majitel(ka) půdy\,/\addthin pozemku}
\dicEntry[landfastur] \dicTerm{land··|fastur} \dicIPA{{l}{a}{n}{\textsubring{d}}{f}{a}{s}{\textsubring{d}}{\textscy}{\textsubring{r}}} \dicPos{adj}[1] \dicFlx{(f ‑föst)}[11] \dicDirectTranslationCS{(jsoucí) při pobřeží (led ap.)}
\dicEntry[landflótta] \dicTerm{land··flótta} \dicIPA{{l}{a}{n}{\textsubring{d}}{f}{l}{ou}{h}{\textsubring{d}}{a}} \dicPos{adj}[13] \dicFlx{indecl}[1] \dicDirectTranslationCS{(jsoucí) v~exilu, (jsoucí) ve vyhnanství}
\dicEntry[landfræðilegur] \dicTerm{land·fræði··legur} \dicIPA{{l}{a}{n}{\textsubring{d}}{f}{r}{a}{i}{ð}{\textsci}{l}{\textepsilon}{\textbabygamma}{\textscy}{\textsubring{r}}} \dicPos{adj}[1]\dicFlx{}[-8] \dicDirectTranslationCS{geografický}
\dicEntry[landfræðingur] \dicTerm{land·fræð··ing|ur} \dicIPA{{l}{a}{n}{\textsubring{d}}{f}{r}{a}{i}{ð}{i}{\ng}{\r{g}}{\textscy}{\textsubring{r}}} \dicPos{m}[6] \dicFlx{(‑s, ‑ar)}[8] \dicDirectTranslationCS{geograf(ka)}
\dicEntry[landgrunn] \dicTerm{land··grunn} \dicIPA{{l}{a}{n}{\textsubring{d}}{\r{g}}{r}{\textscy}{\textsubring{n}}} \dicPos{n}[2] \dicFlx{(‑s, ‑)}[5] \dicFieldCat{geol.} \dicDirectTranslationCS{kontinentální\,/\addthin pevninský šelf, pevninský práh}
\dicEntry[landgræðsla] \dicTerm{land··græðsl|a} \dicIPA{{l}{a}{n}{\textsubring{d}}{\r{g}}{r}{a}{i}{ð}{s}{\textsubring{d}}{l}{a}} \dicPos{f}[1] \dicFlx{(‑u)}[5] \dicDirectTranslationCS{meliorace}
\dicEntry[landgöngubrú] \dicTerm{land··göngu·|brú} \dicIPA{{l}{a}{n}{\textsubring{d}}{\r{g}}{\oe i}{\ng}{\r{g}}{\textscy}{\textsubring{b}}{r}{u}} \dicPos{f}[9] \dicFlx{(‑brúar, ‑brýr)}[4] \dicDirectTranslationCS{lodní lávka}
\dicEntry[landhelgi] \dicTerm{land··helg|i} \dicsymFrequent\  \dicIPA{{l}{a}{n}{\textsubring{d}}{h}{\textepsilon}{l}{\r{\textObardotlessj}}{\textsci}} \dicPos{f}[3] \dicFlx{(‑i)}[3] \dicDirectTranslationCS{výsostné\,/\addthin teritoriální vody} \dicExampleIS{í landhelgi Íslands} \dicExampleCS{v~islandských výsostných vodách}
\dicEntry[landhelgisgæsla] \dicTerm{land·helgis··gæsl|a} \dicIPA{{l}{a}{n}{\textsubring{d}}{h}{\textepsilon}{l}{\r{\textObardotlessj}}{\textsci}{s}{\r{\textObardotlessj}}{a}{i}{s}{\textsubring{d}}{l}{a}} \dicPos{f}[1] \dicFlx{(‑u)}[5] \dicDirectTranslationCS{pobřežní stráž\,/\addthin hlídka}
\dicEntry[landher] \dicTerm{land··her} \dicIPA{{l}{a}{n}{\textsubring{d}}{h}{\textepsilon}{\textsubring{r}}} \dicPos{m}[9] \dicFlx{(‑s, ‑ir)}[21] \dicDirectTranslationCS{pozemní vojsko}
\dicEntry[landhreinsun] \dicTerm{land··hreins|un} \dicIPA{{l}{a}{n}{\textsubring{d}}{\textsubring{r}}{ei}{n}{s}{\textscy}{\textsubring{n}}} \dicPos{f}[7] \dicFlx{(‑unar)}[9] \dicPhraseIS{það er landhreinsun að e‑m} \dicLangCat{přen.} \dicDirectTranslationCS{no sláva, že už je (kdo) pryč}
\dicEntry[landi] \dicTerm{land|i} \dicIPA{{l}{a}{n}{\textsubring{d}}{\textsci}} \dicPos{m}[1] \dicFlx{(‑a, ‑ar)}[8] \textbf{1.} \dicSynonym{samlandi} \dicDirectTranslationCS{krajan(ka)}  \textbf{2.} \dicSynonym{brugg} \dicDirectTranslationCS{(domácí) pálenka}
\dicEntry[landkrabbi] \dicTerm{land··krabb|i} \dicIPA{{l}{a}{n}{\textsubring{d}}{k\smash{\textsuperscript{h}}}{r}{a}{\textsubring{b}}{\textsci}} \dicPos{m}[1] \dicFlx{(‑a, ‑ar)}[8] \textbf{1.} \dicFieldCat{zool.} \dicDirectTranslationCS{krab} \textit{(l.~{\textLA{Gecarcinus}})}  \dicsymPhoto\   \textbf{2.} \dicDirectTranslationCS{suchozemská krysa} \dicIndirectTranslationCS{(o~člověku, který ví málo o~práci na moři)}
\dicFigure{ds_image_landkrabbi_0_1.jpg}{Landkrabbi}{Landkrabbi - Dozenist, CC BY-SA 3.0}
\dicEntry[landkynning] \dicTerm{land··kynn·ing} \dicIPA{{l}{a}{n}{\textsubring{d}}{c\smash{\textsuperscript{h}}}{\textsci}{n}{i}{\ng}{\r{g}}} \dicPos{f}[4] \dicFlx{(‑ar, ‑ar)}[5] \dicDirectTranslationCS{propagace země}
\dicEntry[landkönnuður] \dicTerm{land··könn·uð|ur} \dicIPA{{l}{a}{n}{\textsubring{d}}{k\smash{\textsuperscript{h}}}{\oe}{n}{\textscy}{ð}{\textscy}{\textsubring{r}}} \dicPos{m}[10] \dicFlx{(‑ar, ‑ir)}[4] \dicDirectTranslationCS{cestovatel(ka), průzkumník, průzkumnice, objevitel(ka) (neznámých končin ap.)}
\dicEntry[landluktur] \dicTerm{land··luktur} \dicIPA{{l}{a}{n}{\textsubring{d}}{l}{\textscy}{x}{\textsubring{d}}{\textscy}{\textsubring{r}}} \dicPos{adj}[1]\dicFlx{}[-10] \dicDirectTranslationCS{vnitrozemský}
\dicEntry[landlægur] \dicTerm{land··lægur} \dicIPA{{l}{a}{n}{\textsubring{d}}{l}{a}{i}{\textbabygamma}{\textscy}{\textsubring{r}}} \dicPos{adj}[1]\dicFlx{}[-1] \textbf{1.} \dicFieldCat{med.} \dicDirectTranslationCS{endemický} \dicExampleIS{landlægur sjúkdómur} \dicExampleCS{endemické onemocnění}  \textbf{2.} \dicSynonym{fastur} \dicDirectTranslationCS{zakořeněný} \dicExampleIS{landlægur siður} \dicExampleCS{zakořeněný zvyk}
\dicEntry[landlæknir] \dicTerm{land··lækn|ir} \dicIPA{{l}{a}{n}{\textsubring{d}}{l}{a}{i}{h}{\r{g}}{n}{\textsci}{\textsubring{r}}} \dicPos{m}[7] \dicFlx{(‑is, ‑ar)}[1] \dicDirectTranslationCS{ředitel(ka) zdravotnictví}
\dicEntry[Landmannalaugar] \dicTerm{Land·manna··laugar} \dicIPA{{l}{a}{n}{\textsubring{d}}{m}{a}{n}{a}{l}{\oe i}{\textbabygamma}{a}{\textsubring{r}}} \dicPos{f}[4] \dicFlx{pl}[2] \dicFieldCat{geog.} \dicDirectTranslationCS{Landmannalaugar} \dicIndirectTranslationCS{(oblast s~horkými prameny a~barevnými skalami na jihu Islandu, oblíbený turistický cíl)}
\dicEntry[landmótunarfræði] \dicTerm{land·mótunar··fræð|i} \dicIPA{{l}{a}{n}{\textsubring{d}}{m}{ou}{\textsubring{d}}{\textscy}{n}{a}{\textsubring{r}}{f}{r}{a}{i}{ð}{\textsci}} \dicPos{f}[3] \dicFlx{(‑i)}[3] \dicDirectTranslationCS{geomorfologie}
\dicEntry[landmælingafræði] \dicTerm{land·mælinga··fræð|i} \dicIPA{{l}{a}{n}{\textsubring{d}}{m}{a}{i}{l}{i}{\ng}{\r{g}}{a}{f}{r}{a}{i}{ð}{\textsci}} \dicPos{f}[3] \dicFlx{(‑i)}[3] \dicDirectTranslationCS{geodézie, zeměměřičství}
\dicEntry[landmælingar] \dicTerm{land··mæl·ingar} \dicIPA{{l}{a}{n}{\textsubring{d}}{m}{a}{i}{l}{i}{\ng}{\r{g}}{a}{\textsubring{r}}} \dicPos{f}[4] \dicFlx{pl}[6] \dicDirectTranslationCS{geodetické měření}
\dicEntry[landnám] \dicTerm{land··nám} \dicIPA{{l}{a}{n}{\textsubring{d}}{n}{au}{\textsubring{m}}} \dicPos{n}[2] \dicFlx{(‑s)}[2] \dicDirectTranslationCS{osídlení, osídlování, kolonizace, kolonizování}

\dicEntry[landnámsmaður] \textls[19]{\dicTerm{land·náms··|maður} \dicIPA{{l}{a}{n}{\textsubring{d}}{n}{au}{m}{s}{m}{a}{ð}{\textscy}{\textsubring{r}}} \dicPos{m}[13] \dicFlx{(‑manns, ‑menn)}[2]} \dicDirectTranslationCS{osadník, osadnice, kolonista, kolonistka}

\dicEntry[landnámsöld] \dicTerm{land·náms··|öld} \dicIPA{{l}{a}{n}{\textsubring{d}}{n}{au}{m}{s}{\oe}{l}{\textsubring{d}}} \dicPos{f}[7] \dicFlx{(‑aldar)}[19] \dicFieldCat{hist.} \dicDirectTranslationCS{doba osídlování}
\dicEntry[landnemi] \dicTerm{land··nem|i} \dicIPA{{l}{a}{n}{\textsubring{d}}{n}{\textepsilon}{m}{\textsci}} \dicPos{m}[1] \dicFlx{(‑a, ‑ar)}[1] \dicDirectTranslationCS{osadník, osadnice, kolonista, kolonistka} \dicExampleIS{íslensku landnemarnir í Vesturheimi} \dicExampleCS{islandští osadníci v~Severní Americe}
\dicEntry[landráð] \dicTerm{land··ráð} \dicIPA{{l}{a}{n}{\textsubring{d}}{r}{au}{\texttheta}} \dicPos{n}[2] \dicFlx{pl}[1] \dicFieldCat{práv.} \dicDirectTranslationCS{velezrada, vlastizrada}
\dicEntry[landrek] \dicTerm{land··rek} \dicIPA{{l}{a}{n}{\textsubring{d}}{r}{\textepsilon}{\r{g}}} \dicPos{n}[2] \dicFlx{(‑s)}[2] \dicFieldCat{geol.} \dicDirectTranslationCS{kontinentální drift}
\dicEntry[landsbanki] \dicTerm{lands··bank|i} \dicIPA{{l}{a}{n}{\textsubring{d}}{s}{\textsubring{b}}{au}{\r{\textltailn}}{\r{\textObardotlessj}}{\textsci}} \dicPos{m}[1] \dicFlx{(‑a, ‑ar)}[8] \dicDirectTranslationCS{národní banka}
\dicEntry[landsbyggð] \dicTerm{lands··byggð} \dicIPA{{l}{a}{n}{\textsubring{d}}{s}{\textsubring{b}}{\textsci}{\textbabygamma}{\texttheta}} \dicPos{f}[7] \dicFlx{(‑ar)}[3] \dicDirectTranslationCS{venkov, venkovský kraj} \dicExampleIS{úti á landsbyggðinni} \dicExampleCS{na venkově}
\dicEntry[landsfrægur] \dicTerm{lands··frægur} \dicIPA{{l}{a}{n}{\textsubring{d}}{s}{f}{r}{a}{i}{\textbabygamma}{\textscy}{\textsubring{r}}} \dicPos{adj}[1]\dicFlx{}[-1] \dicDirectTranslationCS{známý\,/\addthin proslavený po celé zemi}
\dicEntry[landshluti] \dicTerm{lands··hlut|i} \dicIPA{{l}{a}{n}{\textsubring{d}}{s}{\textsubring{l}}{\textscy}{\textsubring{d}}{\textsci}} \dicPos{m}[1] \dicFlx{(‑a, ‑ar)}[1] \dicDirectTranslationCS{část země, region}
\dicEntry[landsími] \dicTerm{land··sím|i} \dicIPA{{l}{a}{n}{\textsubring{d}}{s}{i}{m}{\textsci}} \dicPos{m}[1] \dicFlx{(‑a, ‑ar)}[1] \dicFieldCat{hist.} \dicDirectTranslationCS{národní telefonní společnost}
\dicEntry[landskjálfti] \dicTerm{land··skjálft|i} \dicIPA{{l}{a}{n}{\textsubring{d}}{s}{\r{\textObardotlessj}}{au}{\textsubring{l}}{\textsubring{d}}{\textsci}} \dicPos{m}[1] \dicFlx{(‑a, ‑ar)}[1] \dicFieldCat{geol.} \dicLangCat{zast.} \dicSynonym{jarðskjálfti} \dicDirectTranslationCS{zemětřesení}
\dicEntry[landslag] \dicTerm{lands··lag} \dicsymFrequent\  \dicIPA{{l}{a}{n}{\textsubring{d}}{s}{l}{a}{x}} \dicPos{n}[2] \dicFlx{(‑s)}[2] \dicDirectTranslationCS{krajina, terén} \dicExampleIS{skoða landslag í firðinum} \dicExampleCS{zkoumat krajinu ve fjordu}
\dicEntry[landslagsarkitekt] \dicTerm{lands·lags··arkitekt} \dicIPA{{l}{a}{n}{\textsubring{d}}{s}{l}{a}{x}{s}{a}{\textsubring{r}}{\r{\textObardotlessj}}{\textsci}{\textsubring{d}}{\textepsilon}{x}{\textsubring{d}}} \dicPos{m}[4] \dicFlx{(‑s, ‑ar)}[10] \dicDirectTranslationCS{krajinný architekt, krajinná architektka}
\dicEntry[landslagsmynd] \dicTerm{lands·lags··mynd} \dicIPA{{l}{a}{n}{\textsubring{d}}{s}{l}{a}{x}{s}{m}{\textsci}{n}{\textsubring{d}}} \dicPos{f}[7] \dicFlx{(‑ar, ‑ir)}[1] \dicDirectTranslationCS{krajina (obraz), krajinomalba}
\dicEntry[landsleikur] \dicTerm{lands··leik|ur} \dicIPA{{l}{a}{n}{\textsubring{d}}{s}{l}{ei}{\r{g}}{\textscy}{\textsubring{r}}} \dicPos{m}[9] \dicFlx{(‑s, ‑ir)}[15] \dicDirectTranslationCS{mezistátní\,/\addthin mezinárodní zápas\,/\addthin utkání (ve fotbale ap.)}
\dicEntry[landslið] \dicTerm{lands··lið} \dicIPA{{l}{a}{n}{\textsubring{d}}{s}{l}{\textsci}{\texttheta}} \dicPos{n}[2] \dicFlx{(‑s, ‑)}[5] \dicDirectTranslationCS{národní tým\,/\addthin mužstvo}
\dicEntry[landsliðsþjálfari] \dicTerm{lands·liðs··þjálf·ar|i} \dicIPA{{l}\-{a}\-{n}\-{\textsubring{d}}\-{s}\-{l}\-{\textsci}\-{ð}\-{s}\-{\texttheta}\-{j}\-{au}\-{l}\-{v}\-{a}\-{r}\-{\textsci}\-} \dicPos{m}[1] \dicFlx{(‑a, ‑ar)}[13] \dicDirectTranslationCS{trenér(ka) národního týmu}
\dicEntry[landslýður] \dicTerm{lands··lýð|ur} \dicIPA{{l}{a}{n}{\textsubring{d}}{s}{l}{i}{ð}{\textscy}{\textsubring{r}}} \dicPos{m}[9] \dicFlx{(‑s, ‑ir)}[8] \dicDirectTranslationCS{obyvatelstvo}
\dicEntry[landsmaður] \dicTerm{lands··|maður} \dicsymFrequent\  \dicIPA{{l}{a}{n}{\textsubring{d}}{s}{m}{a}{ð}{\textscy}{\textsubring{r}}} \dicPos{m}[13] \dicFlx{(‑manns, ‑menn)}[2] \dicDirectTranslationCS{obyvatel(ka)} \dicExampleIS{allir landsmenn undir fimmtugu} \dicExampleCS{všichni obyvatelé pod padesát let}
\dicEntry[landsmót] \dicTerm{lands··mót} \dicIPA{{l}{a}{n}{\textsubring{d}}{s}{m}{ou}{\textsubring{d}}} \dicPos{n}[2] \dicFlx{(‑s, ‑)}[5] \textbf{1.} \dicDirectTranslationCS{celorepublikový\,/\addthin národní sjezd\,/\addthin setkání}  \textbf{2.} \dicDirectTranslationCS{celorepublikový\,/\addthin národní turnaj (v~házené ap.)}
\dicEntry[landspítali] \dicTerm{land··spítal|i} \dicIPA{{l}{a}{n}{\textsubring{d}}{s}{\textsubring{b}}{i}{t\smash{\textsuperscript{h}}}{a}{l}{\textsci}} \dicPos{m}[1] \dicFlx{(‑a, ‑ar)}[8] \dicDirectTranslationCS{národní univerzitní nemocnice}
\begin{xtolerant}{}{1pt}
\dicEntry[landssamband] \dicTerm{lands··sam·|band} \dicIPA{{l}{a}{n}{\textsubring{d}}{s}{a}{m}{\textsubring{b}}{a}{n}{\textsubring{d}}} \dicPos{n}[2] \dicFlx{(‑bands, ‑bönd)}[8] \dicDirectTranslationCS{národní\,/\addthin celorepublikový svaz}
\end{xtolerant}
\dicEntry[landstjóri] \dicTerm{land··stjór|i} \dicIPA{{l}{a}{n}{\textsubring{d}}{s}{\textsubring{d}}{j}{ou}{r}{\textsci}} \dicPos{m}[1] \dicFlx{(‑a, ‑ar)}[1] \dicDirectTranslationCS{místodržitel(ka), guvernér(ka), zemský správce, zemská správkyně}
\dicEntry[landsvala] \dicTerm{land··|svala} \dicIPA{{l}{a}{n}{\textsubring{d}}{s}{v}{a}{l}{a}} \dicPos{f}[1] \dicFlx{(‑svölu, ‑svölur)}[20] \dicFieldCat{zool.} \dicDirectTranslationCS{vlaštovka, vlaštovka obecná} \textit{(l.~{\textLA{Hirundo rustica}})}  \dicsymPhoto\ 
\dicFigure{20731.jpg}{Landsvala}{Landsvala - Menke, Dave, Biolib, PD}
\dicEntry[Landsvirkjun] \dicTerm{Lands··virkj|un} \dicIPA{{l}{a}{n}{\textsubring{d}}{s}{v}{\textsci}{\textsubring{r}}{\r{\textObardotlessj}}{\textscy}{\textsubring{n}}} \dicPos{f}[7] \dicFlx{(‑unar)}[10] \dicIndirectTranslationCS{islandská společnost provozující elektrárny}
\dicEntry[landsvæði] \dicTerm{land··svæði} \dicIPA{{l}{a}{n}{\textsubring{d}}{s}{v}{a}{i}{ð}{\textsci}} \dicPos{n}[2] \dicFlx{(‑s, ‑)}[14] \dicDirectTranslationCS{území, oblast, teritorium}
\dicEntry[langa] \dicTerm{langa\smash{\textsuperscript{1}}} \dicIPA{{l}{au}{\ng}{\r{g}}{a}} \dicPos{f}[1] \dicFlx{(löngu, löngur)}[20] \dicFieldCat{zool.} \dicDirectTranslationCS{mník mořský} \textit{(l.~{\textLA{Molva molva}})}
\dicEntry[langa] \dicTerm{lang|a\smash{\textsuperscript{2}}} \dicsymFrequent\  \dicIPA{{l}{au}{\ng}{\r{g}}{a}} \dicPos{v}[1] \dicFlx{(‑aði)}[79] \dicFlx{impers} \dicPhraseIS{e‑n langar} \dicDirectTranslationCS{(kdo) touží, (kdo) chce, (kdo) má chuť} \dicExampleIS{Mig langar að fara.} \dicExampleCS{Chce se mi jít.};  \dicPhraseIS{e‑n langar í e‑ð} \dicDirectTranslationCS{(kdo) touží po (čem), (kdo) chce (co), (kdo) má chuť na (co)} \dicExampleIS{Mig langar í sælgæti.} \dicExampleCS{Mám chuť na něco sladkého.};  \dicPhraseIS{e‑n langar til e‑s} \dicDirectTranslationCS{(kdo) touží po (čem), (kdo) chce (co)} \dicExampleIS{Mig langar til að lesa bókina.} \dicExampleCS{Mám chuť si číst knížku.};  \dicPhraseIS{e‑n langar eftir e‑u} \dicDirectTranslationCS{(kdo) touží po (čem), (kdo) prahne po (čem)} \dicExampleIS{Hana langaði eftir friði.} \dicExampleCS{Toužila po klidu.}
\dicEntry[langafasta] \dicTerm{langa··fasta} \dicIPA{{l}{au}{\ng}{\r{g}}{a}{f}{a}{s}{\textsubring{d}}{a}} \dicPos{f}[1] \dicFlx{(lönguföstu, lönguföstur)}[21] \dicFieldCat{náb.} \dicDirectTranslationCS{půst} \dicIndirectTranslationCS{(trvající 7 týdnů od Popeleční středy do Velikonoc)}
\dicEntry[langafi] \dicTerm{lang··af|i} \dicIPA{{l}{au}{\ng}{\r{g}}{a}{v}{\textsci}} \dicPos{m}[1] \dicFlx{(‑a, ‑ar)}[8] \dicDirectTranslationCS{pradědeček} \dicIndirectTranslationCS{(otec dědečka nebo babičky)}
\dicEntry[langamma] \dicTerm{lang··|amma} \dicIPA{{l}{au}{\ng}{\r{g}}{a}{m}{a}} \dicPos{f}[1] \dicFlx{(‑ömmu, ‑ömmur)}[8] \dicDirectTranslationCS{prababička} \dicIndirectTranslationCS{(matka dědečka nebo babičky)}
\dicEntry[langanir] \dicTerm{langanir} \dicIPA{{l}{au}{\ng}{\r{g}}{a}{n}{\textsci}{\textsubring{r}}} \dicPos{f} \dicFlx{pl nom} \dicLink{löngun}
\dicEntry[langatöng] \dicTerm{langa··töng}\dicTerm{, löngutöng} \dicIPA{{l}\-{au}\-{\ng}\-{\r{g}}\-{a}\-{t\smash{\textsuperscript{h}}}\-{\oe i}\-{\ng}\-{\r{g}}\-} \dicPos{f}[8] \dicFlx{(löngutangar, löngutengur\,/\addthin löngutangir)}[10] \dicDirectTranslationCS{prostředník (prst), prostředníček}
\dicEntry[langbestur] \dicTerm{lang··bestur} \dicIPA{{l}{au}{\ng}{\r{g}}{\textsubring{b}}{\textepsilon}{s}{\textsubring{d}}{\textscy}{\textsubring{r}}} \dicPos{adj}[11]\dicFlx{}[-2] \dicFlx{m sg nom sup} \dicDirectTranslationCS{úplně\,/\addthin zdaleka nejlepší}
\dicEntry[langbylgja] \dicTerm{lang··bylgj|a} \dicIPA{{l}{au}{\ng}{\r{g}}{\textsubring{b}}{\textsci}{l}{\r{\textObardotlessj}}{a}} \dicPos{f}[1] \dicFlx{(‑u, ‑ur)}[25] \dicFieldCat{fyz.} \dicDirectTranslationCS{dlouhá vlna} \dicAntonym{stuttbylgja}
\dicEntry[langdreginn] \dicTerm{lang··dreginn} \dicIPA{{l}{au}{\ng}{\r{g}}{\textsubring{d}}{r}{ei}{\textsci}{\textsubring{n}}} \dicPos{adj}[6]\dicFlx{}[-2] \textbf{1.} \dicDirectTranslationCS{zdlouhavý, rozvláč\-ný} \dicExampleIS{langdreginn ræða} \dicExampleCS{zdlouhavý proslov}  \textbf{2.} \dicDirectTranslationCS{dlouhotrvající} \dicExampleIS{langdreginn sjúkdómur} \dicExampleCS{dlouhotrvající onemocnění}
\dicEntry[langdvöl] \dicTerm{lang··|dvöl} \dicIPA{{l}{au}{\ng}{\r{g}}{\textsubring{d}}{v}{\oe}{\textsubring{l}}} \dicPos{f}[7] \dicFlx{(‑dvalar, ‑dvalir)}[16] \dicDirectTranslationCS{dlouhotrvající pobyt};  \dicPhraseIS{dvelja\,/\addthin vera e‑s staðar langdvölum} \dicDirectTranslationCS{pobývat (kde) dlouho}
\dicEntry[langeygur] \dicTerm{lang··eygur} \dicIPA{{l}{au}{\ng}{\r{g}}{ei}{\textbabygamma}{\textscy}{\textsubring{r}}} \dicPos{adj}[1]\dicFlx{}[-1] \dicPhraseIS{vera langeygur eftir e‑u} \dicDirectTranslationCS{být nedočkavý (čeho), nemoct se dočkat (čeho)}
\dicEntry[langframi] \dicTerm{lang··fram|i} \dicIPA{{l}{au}{\ng}{\r{g}}{f}{r}{a}{m}{\textsci}} \dicPos{m}[1] \dicFlx{(‑a)}[3] \dicPhraseIS{til langframa} \dicFlx{adv} \dicDirectTranslationCS{dlouhodobě}
\dicEntry[langhlaup] \dicTerm{lang··hlaup} \dicIPA{{l}{au}{\ng}{\r{g}}{\textsubring{l}}{\oe i}{\textsubring{b}}} \dicPos{n}[2] \dicFlx{(‑s, ‑)}[5] \dicFieldCat{sport.} \dicDirectTranslationCS{dálkový běh}
\dicEntry[langlífi] \dicTerm{lang··lífi} \dicIPA{{l}{au}{\ng}{\r{g}}{l}{i}{v}{\textsci}} \dicPos{n}[2] \dicFlx{(‑s)}[20] \dicDirectTranslationCS{dlouhověkost}
\dicEntry[langlífur] \dicTerm{lang··lífur} \dicIPA{{l}{au}{\ng}{\r{g}}{l}{i}{v}{\textscy}{\textsubring{r}}} \dicPos{adj}[1]\dicFlx{}[-1] \dicDirectTranslationCS{dlouhověký}
\dicEntry[langloka] \dicTerm{lang··lok|a} \dicIPA{{l}{au}{\ng}{\r{g}}{l}{\textopeno}{\r{g}}{a}} \dicPos{f}[1] \dicFlx{(‑u, ‑ur)}[19] \textbf{1.} \dicSynonym{romsa\smash{\textsuperscript{1}}} \dicDirectTranslationCS{litanie, nudná a~zdlouhavá řeč, nudný a~zdlouhavý text}  \textbf{2.} \dicSynonym{samloka} \dicDirectTranslationCS{bageta}
\dicEntry[langorður] \dicTerm{lang··orður} \dicIPA{{l}{au}{\ng}{\r{g}}{\textopeno}{r}{ð}{\textscy}{\textsubring{r}}} \dicPos{adj}[2]\dicFlx{}[-1] \dicDirectTranslationCS{rozvláčný, upovídaný}
\dicEntry[langreyður] \dicTerm{lang··reyð|ur} \dicIPA{{l}{au}{\ng}{\r{g}}{r}{ei}{ð}{\textscy}{\textsubring{r}}} \dicPos{f}[5] \dicFlx{(‑ar, ‑ar)}[1] \dicFieldCat{zool.} \dicDirectTranslationCS{plejtvák myšok} \textit{(l.~{\textLA{Balaenoptera physalus}})}
\dicEntry[langrækinn] \dicTerm{lang··rækinn} \dicIPA{{l}{au}{\ng}{\r{g}}{r}{a}{i}{\r{\textObardotlessj}}{\textsci}{\textsubring{n}}} \dicPos{adj}[6]\dicFlx{}[-2] \dicDirectTranslationCS{neodpouštějící, nesmiřitelný}
\dicEntry[langskip] \dicTerm{lang··skip} \dicIPA{{l}{au}{\ng}{\r{g}}{s}{\r{\textObardotlessj}}{\textsci}{\textsubring{b}}} \dicPos{n}[2] \dicFlx{(‑s, ‑)}[5] \dicFieldCat{nám.} \dicDirectTranslationCS{vikinská válečná loď}
\dicEntry[langsnið] \dicTerm{lang··snið} \dicIPA{{l}{au}{\ng}{\r{g}}{s}{\textsubring{d}}{n}{\textsci}{\texttheta}} \dicPos{n}[2] \dicFlx{(‑s, ‑)}[5] \dicFieldCat{poč.} \dicDirectTranslationCS{formát na délku}
\dicEntry[langsóttur] \dicTerm{lang··sóttur} \dicIPA{{l}{au}{\ng}{\r{g}}{s}{ou}{h}{\textsubring{d}}{\textscy}{\textsubring{r}}} \dicPos{adj}[1]\dicFlx{}[-10] \dicDirectTranslationCS{nevěrohodný, přitažený za vlasy} \dicExampleIS{langsótt sönnun} \dicExampleCS{nevěrohodný důkaz}
\dicEntry[langspil] \dicTerm{lang··spil} \dicIPA{{l}{au}{\ng}{\r{g}}{s}{\textsubring{b}}{\textsci}{\textsubring{l}}} \dicPos{n}[2] \dicFlx{(‑s, ‑)}[5] \dicFieldCat{hud.} \dicIndirectTranslationCS{islandský strunný hudební nástroj}
\dicEntry[langstökk] \dicTerm{lang··stökk} \dicIPA{{l}{au}{\ng}{\r{g}}{s}{\textsubring{d}}{\oe}{h}{\r{g}}} \dicPos{n}[2] \dicFlx{(‑s, ‑)}[5] \dicFieldCat{sport.} \dicDirectTranslationCS{skok do dálky, skok daleký}
\dicEntry[langt] \dicTerm{langt} \dicsymFrequent\  \dicIPA{{l}{au}{\r{\ng}}{\textsubring{d}}} \dicPos{adv} \dicFlx{(comp lengra, sup lengst)} \dicDirectTranslationCS{daleko} \dicExampleIS{Þetta liggur langt inni í fortíðinni.} \dicExampleCS{Má to původ v~daleké minulosti.};  \dicPhraseIS{langt í burtu} \dicFlx{adv} \dicDirectTranslationCS{daleko};  \dicPhraseIS{langt síðan} \dicFlx{adv} \dicDirectTranslationCS{dávno}
\dicEntry[langtíma] \dicTerm{lang··tíma-} \dicIPA{{l}{au}{\ng}{\r{g}}{t\smash{\textsuperscript{h}}}{i}{m}{a}} \dicPos{predp} \dicDirectTranslationCS{dlouhodobý}
\dicEntry[langtímalán] \dicTerm{lang·tíma··lán} \dicIPA{{l}{au}{\ng}{\r{g}}{t\smash{\textsuperscript{h}}}{i}{m}{a}{l}{au}{\textsubring{n}}} \dicPos{n}[2] \dicFlx{(‑s, ‑)}[5] \dicFieldCat{ekon.} \dicDirectTranslationCS{dlouhodobá půjčka}
\dicEntry[langtímasamningur] \dicTerm{lang·tíma··samn·ing|ur} \dicIPA{{l}\-{au}\-{\ng}\-{\r{g}}\-{t\smash{\textsuperscript{h}}}\-{i}\-{m}\-{a}\-{s}\-{a}\-{m}\-{n}\-{i}\-{\ng}\-{\r{g}}\-{\textscy}\-{\textsubring{r}}\-} \dicPos{m}[6] \dicFlx{(‑s, ‑ar)}[8] \dicDirectTranslationCS{dlouhodobá smlouva}
\dicEntry[langur] \dicTerm{lang|ur\smash{\textsuperscript{1}}} \dicIPA{{l}{au}{\ng}{\r{g}}{\textscy}{\textsubring{r}}} \dicPos{m}[6] \dicFlx{(‑s, ‑ar)}[16] \dicPhraseIS{draga e‑ð á langinn} \dicDirectTranslationCS{odložit (co), odkládat (co)}
\dicEntry[langur] \dicTerm{langur\smash{\textsuperscript{2}}} \dicsymFrequent\  \dicIPA{{l}{au}{\ng}{\r{g}}{\textscy}{\textsubring{r}}} \dicPos{adj}[10] \dicFlx{(f löng, comp lengri, sup lengstur)}[2] \mbox{\textbf{1.} \dicDirectTranslationCS{dlouhý}} \dicIndirectTranslationCS{(o~vzdálenosti)} \dicExampleIS{löng leið} \dicExampleCS{dlouhá cesta} \dicAntonym{stuttur}  \textbf{2.} \dicDirectTranslationCS{vysoký} \dicExampleIS{langur maður} \dicExampleCS{vysoký člověk}  \textbf{3.} \dicDirectTranslationCS{dlouhý} \dicIndirectTranslationCS{(o~čase)} \dicExampleIS{langur tími} \dicExampleCS{dlouhý čas};  \dicPhraseIS{fyrir löngu} \dicFlx{adv} \dicDirectTranslationCS{dávno, před dlouhou dobou};  \dicPhraseIS{löngum stundum} \dicFlx{adv} \dicDirectTranslationCS{dlouho, dlouhé hodiny};  \dicPhraseIS{það er langt til e‑s} \dicDirectTranslationCS{je daleko do (čeho) (Vánoc ap.)}
\dicEntry[langvarandi] \dicTerm{lang··var·andi} \dicsymFrequent\  \dicIPA{{l}{au}{\ng}{\r{g}}{v}{a}{r}{a}{n}{\textsubring{d}}{\textsci}} \dicPos{adj}[13] \dicFlx{indecl}[1] \dicDirectTranslationCS{dlouhotrvající, prodlužovaný} \dicExampleIS{langvarandi erfiðleikar} \dicExampleCS{dlouhotrvající potíže}
\dicEntry[langvinnur] \textls[15]{\dicTerm{lang··vinnur} \dicIPA{{l}{au}{\ng}{\r{g}}{v}{\textsci}{n}{\textscy}{\textsubring{r}}} \dicPos{adj}[1]\dicFlx{}[-1] \dicDirectTranslationCS{přetrvávající, táhlý, dlouhotrvající} \dicExampleIS{langvinnur sjúkdómur} \dicExampleCS{přetrvávající onemocnění}}
\dicEntry[langvía] \dicTerm{lang··ví|a} \dicIPA{{l}{au}{\ng}{\r{g}}{v}{i}{j}{a}} \dicPos{f}[1] \dicFlx{(‑u, ‑ur)}[7] \dicFieldCat{zool.} \dicDirectTranslationCS{alkoun úzkozobý} \textit{(l.~{\textLA{Uria aalge}})}  \dicsymPhoto\ 
\dicFigure{20877.jpg}{Langvía}{Langvía - Sowls, Art, Biolib, PD}
\dicEntry[langþráður] \dicTerm{lang··þráður} \dicIPA{{l}{au}{\ng}{\r{g}}{\texttheta}{r}{au}{ð}{\textscy}{\textsubring{r}}} \dicPos{adj}[2]\dicFlx{}[-12] \dicDirectTranslationCS{vytoužený, (dlouho) očekávaný}
\dicEntry[langþreyttur] \dicTerm{lang··þreyttur} \dicIPA{{l}{au}{\ng}{\r{g}}{\texttheta}{r}{ei}{h}{\textsubring{d}}{\textscy}{\textsubring{r}}} \dicPos{adj}[1]\dicFlx{}[-13] \dicDirectTranslationCS{znavený, (velmi) unavený}
\dicEntry[lantan] \dicTerm{lantan} \dicIPA{{l}{a}{\textsubring{n}}{\textsubring{d}}{a}{\textsubring{n}}} \dicPos{n}[2] \dicFlx{(‑s)}[2] \dicFieldCat{chem.} \dicDirectTranslationCS{lanthan} \textit{(l.~{\textLA{La, Lanthanum}})}
\dicEntry[Laos] \dicTerm{Laos} \dicIPA{{l}{a}{\textlengthmark}{\textopeno}{s}} \dicPos{n}[4] \dicFlx{indecl}[2] \dicFieldCat{geog.} \dicDirectTranslationCS{Laos}
\dicEntry[Laosi] \dicTerm{Laos|i} \dicIPA{{l}{a}{\textlengthmark}{\textopeno}{s}{\textsci}} \dicPos{m}[1] \dicFlx{(‑a, ‑ar)}[1] \dicDirectTranslationCS{Laosan(ka)}
\dicEntry[laoska] \dicTerm{laosk|a} \dicIPA{{l}{a}{\textlengthmark}{\textopeno}{s}{\r{g}}{a}} \dicPos{f}[1] \dicFlx{(‑u)}[5] \dicDirectTranslationCS{laoština}
\dicEntry[laoskur] \dicTerm{laoskur} \dicIPA{{l}{a}{\textlengthmark}{\textopeno}{s}{\r{g}}{\textscy}{\textsubring{r}}} \dicPos{adj}[1]\dicFlx{}[-6] \dicDirectTranslationCS{laoský}
\dicEntry[lapið] \dicTerm{lapið} \dicIPA{{l}{a}{\textlengthmark}{\textsubring{b}}{\textsci}{\texttheta}} \dicPos{v} \dicFlx{supin} \dicLink{lepja}
\dicEntry[lappa] \dicTerm{lapp|a} \dicIPA{{l}{a}{h}{\textsubring{b}}{a}} \dicPos{v}[1] \dicFlx{(‑aði)}[13] \dicDirectTranslationCS{ochodit, prochodit (boty ap.)} \dicExampleIS{lappa skóna} \dicExampleCS{prochodit boty};  \dicPhraseIS{e‑að lappar e‑n} \dicDirectTranslationCS{(co) je (komu) volné (boty ap.)} \dicExampleIS{Skórinn lappar mig.} \dicExampleCS{Boty mi jsou volné.};  \dicIdiom{lappa}[upp á]{ \dicPhraseIS{lappa upp á e‑ð}} \dicDirectTranslationCS{osvěžit (co), vyspravit (co)}
\dicEntry[lappajaðrakan] \dicTerm{lappa··jað·rakan} \dicIPA{{l}{a}{h}{\textsubring{b}}{a}{j}{a}{ð}{r}{a}{\r{g}}{a}{\textsubring{n}}} \dicPos{m}[4] \dicFlx{(‑s, ‑ar)}[12] \dicFieldCat{zool.} \dicDirectTranslationCS{břehouš rudý} \textit{(l.~{\textLA{Limosa lapponica}})}  \dicsymPhoto\ 
\dicFigure{ds_image_lappajadrakan_0_1.jpg}{Lappajaðrakan}{Lappajaðrakan - Andreas Trepte, CC BY-SA 2.5}
\dicEntry[lappalaus] \dicTerm{lappa··laus} \dicIPA{{l}{a}{h}{\textsubring{b}}{a}{l}{\oe i}{s}} \dicPos{adj}[5]\dicFlx{}[-1] \dicPhraseIS{vera orðinn lappalaus} \dicLangCat{přen.} \dicDirectTranslationCS{být k~smrti unavený}
\dicEntry[lappar] \dicTerm{lappar} \dicIPA{{l}{a}{h}{\textsubring{b}}{a}{\textsubring{r}}} \dicPos{f} \dicFlx{sg gen} \dicLink{löpp}
\dicEntry[Lappi] \dicTerm{Lapp|i} \dicIPA{{l}{a}{h}{\textsubring{b}}{\textsci}} \dicPos{m}[1] \dicFlx{(‑a, ‑ar)}[8] \dicDirectTranslationCS{Laponec, Laponka}
\dicEntry[lappir] \dicTerm{lappir} \dicIPA{{l}{a}{h}{\textsubring{b}}{\textsci}{\textsubring{r}}} \dicPos{f} \dicFlx{pl nom} \dicLink{löpp}
\dicEntry[Lappland] \dicTerm{Lapp··land} \dicIPA{{l}{a}{h}{\textsubring{b}}{l}{a}{n}{\textsubring{d}}} \dicPos{n}[2] \dicFlx{(‑s)}[4] \dicFieldCat{geog.} \dicSynonym{Samaland} \dicDirectTranslationCS{Laponsko}
\dicEntry[lappneska] \dicTerm{lapp··nesk|a} \dicIPA{{l}{a}{h}{\textsubring{b}}{n}{\textepsilon}{s}{\r{g}}{a}} \dicPos{f}[1] \dicFlx{(‑u)}[5] \dicDirectTranslationCS{laponština}
\dicEntry[lappneskur] \dicTerm{lapp··neskur} \dicIPA{{l}{a}{h}{\textsubring{b}}{n}{\textepsilon}{s}{\r{g}}{\textscy}{\textsubring{r}}} \dicPos{adj}[1]\dicFlx{}[-1] \dicDirectTranslationCS{laponský}
\dicEntry[lapti] \dicTerm{lapti} \dicIPA{{l}{a}{f}{\textsubring{d}}{\textsci}} \dicPos{v} \dicFlx{ind pf sg 1 pers} \dicLink{lepja}
\dicEntry[las] \dicTerm{las} \dicIPA{{l}{a}{\textlengthmark}{s}} \dicPos{v} \dicFlx{ind pf sg 1 pers} \dicLink{lesa}
\dicEntry[lasanja] \dicTerm{lasanja} \dicIPA{{l}{a}{\textlengthmark}{s}{a}{n}{j}{a}} \dicPos{n}[4] \dicFlx{indecl}[1] \dicFieldCat{kulin.} \dicDirectTranslationCS{lasagne}
\dicEntry[lasburða] \dicTerm{las··burða} \dicIPA{{l}{a}{\textlengthmark}{s}{\textsubring{b}}{\textscy}{r}{ð}{a}} \dicPos{adj}[13] \dicFlx{indecl}[1] \dicDirectTranslationCS{churavý, neduživý}
\dicEntry[lasinn] \dicTerm{lasinn} \dicIPA{{l}{a}{\textlengthmark}{s}{\textsci}{\textsubring{n}}} \dicPos{adj}[6]\dicFlx{}[-3] \textbf{1.} \dicSynonym{veikur} \dicDirectTranslationCS{nemocný, indisponovaný}  \textbf{2.} \dicSynonym{lélegur} \dicDirectTranslationCS{opotřebovaný}
\dicEntry[laska] \dicTerm{lask|a} \dicIPA{{l}{a}{s}{\r{g}}{a}} \dicPos{v}[1] \dicFlx{(‑aði)}[13] \dicFlx{acc} \dicSynonym{skadda} \dicDirectTranslationCS{poškodit, způsobit škodu}
\dicEntry[lasleiki] \dicTerm{las··leik|i} \dicIPA{{l}{a}{\textlengthmark}{s}{l}{ei}{\r{\textObardotlessj}}{\textsci}} \dicPos{m}[1] \dicFlx{(‑a)}[3] \dicDirectTranslationCS{indispozice, churavost}
\dicEntry[last] \dicTerm{last} \dicIPA{{l}{a}{s}{\textsubring{d}}} \dicPos{n}[2] \dicFlx{(‑s)}[2] \dicSynonym*{baktal} \dicDirectTranslationCS{pomluva, osočení};  \dicPhraseIS{leggja e‑m e‑ð til lasts} \dicDirectTranslationCS{dávat (co komu) za vinu}
\dicEntry[lasta] \dicTerm{last|a} \dicIPA{{l}{a}{s}{\textsubring{d}}{a}} \dicPos{v}[1] \dicFlx{(‑aði)}[13] \dicFlx{acc} \dicSynonym{gagnrýna} \dicDirectTranslationCS{očernit, očerňovat, (po)hanit, pohanět} \dicExampleIS{lasta þýðingu á bók} \dicExampleCS{očerňovat překlad knížky}
\dicEntry[lastar] \dicTerm{lastar} \dicIPA{{l}{a}{s}{\textsubring{d}}{a}{\textsubring{r}}} \dicPos{m} \dicFlx{sg gen} \dicLink{löstur}
\dicEntry[lastyrði] \dicTerm{last··yrði} \dicIPA{{l}{a}{s}{\textsubring{d}}{\textsci}{r}{ð}{\textsci}} \dicPos{n}[2] \dicFlx{(‑s, ‑)}[14] \dicSynonym{last} \dicDirectTranslationCS{hana, hanobení}
\dicEntry[latína] \dicTerm{latín|a} \dicIPA{{l}{a}{\textlengthmark}{t\smash{\textsuperscript{h}}}{i}{n}{a}} \dicPos{f}[1] \dicFlx{(‑u)}[5] \dicDirectTranslationCS{latina}
\dicEntry[latínusegl] \dicTerm{latínu··segl} \dicIPA{{l}{a}{\textlengthmark}{t\smash{\textsuperscript{h}}}{i}{n}{\textscy}{s}{\textepsilon}{\r{g}}{\textsubring{l}}} \dicPos{n}[2] \dicFlx{(‑s, ‑)}[5] \dicFieldCat{nám.} \dicDirectTranslationCS{latinská plachta}
\dicEntry[latneskur] \dicTerm{latn··eskur} \dicIPA{{l}{a}{h}{\textsubring{d}}{n}{\textepsilon}{s}{\r{g}}{\textscy}{\textsubring{r}}} \dicPos{adj}[1]\dicFlx{}[-1] \dicDirectTranslationCS{latinský}
\dicEntry[latt] \dicTerm{latt} \dicIPA{{l}{a}{h}{\textsubring{d}}} \dicPos{v} \dicFlx{supin} \dicLink{letja}
\dicEntry[latti] \dicTerm{latti} \dicIPA{{l}{a}{h}{\textsubring{d}}{\textsci}} \dicPos{v} \dicFlx{ind pf sg 1 pers} \dicLink{letja}
\dicEntry[latur] \dicTerm{latur} \dicIPA{{l}{a}{\textlengthmark}{\textsubring{d}}{\textscy}{\textsubring{r}}} \dicPos{adj}[1] \dicFlx{(f löt)}[2] \dicDirectTranslationCS{líný, lenivý}
\dicEntry[lauf] \dicTerm{lauf} \dicsymFrequent\  \dicIPA{{l}{\oe i}{\textlengthmark}{f}} \dicPos{n}[2] \dicFlx{(‑s, ‑)}[5] \textbf{1.} \dicSynonym{laufblað} \dicDirectTranslationCS{list (stromu ap.)} \dicExampleIS{gulnuð lauf að hausti} \dicExampleCS{zežloutlé listy na podzim}  \textbf{2.} \dicSynonym*{litur í spilum} \dicDirectTranslationCS{trefy, kříže} \dicIndirectTranslationCS{(barva v~kartách)}
\dicEntry[laufabrauð] \dicTerm{laufa··brauð} \dicIPA{{l}{\oe i}{\textlengthmark}{v}{a}{\textsubring{b}}{r}{\oe i}{\texttheta}} \dicPos{n}[2] \dicFlx{(‑s, ‑)}[5] \dicFieldCat{kulin.} \dicIndirectTranslationCS{islandský vánoční chléb ručně zdobený různými motivy} \dicsymPhoto\ 
\dicFigure{ds_image_laufabraud_0_2.jpg}{Laufabrauð}{Laufabrauð - Fany Larota Catunta, PD}
\dicEntry[laufblað] \dicTerm{lauf··|blað} \dicIPA{{l}{\oe i}{v}{\textsubring{b}}{l}{a}{\texttheta}} \dicPos{n}[2] \dicFlx{(‑blaðs, ‑blöð)}[8] \dicSynonym{lauf} \dicDirectTranslationCS{list (stromu ap.)}
\dicEntry[laufgast] \dicTerm{laufg|ast} \dicIPA{{l}{\oe i}{v}{\r{g}}{a}{s}{\textsubring{d}}} \dicPos{v}[1] \dicFlx{(‑aðist)}[95] \dicFlx{refl} \dicDirectTranslationCS{nasadit listy, zazelenat se}
\dicEntry[laufglói] \dicTerm{lauf··gló|i} \dicIPA{{l}{\oe i}{v}{\r{g}}{l}{ou}{\textsci}} \dicPos{m}[1] \dicFlx{(‑a, ‑ar)}[1] \dicFieldCat{zool.} \dicSynonym{gullþröstur} \dicDirectTranslationCS{žluva, žluva hajní} \textit{(l.~{\textLA{Oriolus oriolus}})}  \dicsymPhoto\ 
\dicFigure{ds_image_laufgloi_0_1.jpg}{Laufglói}{Laufglói - Dixi, GFDL}
\dicEntry[laufgræna] \dicTerm{lauf··græn|a} \dicIPA{{l}{\oe i}{v}{\r{g}}{r}{a}{i}{n}{a}} \dicPos{f}[1] \dicFlx{(‑u)}[5] \dicFieldCat{bot.} \dicSynonym{blaðgræna} \dicDirectTranslationCS{chlorofyl}
\dicEntry[laufskógur] \dicTerm{lauf··skóg|ur} \dicIPA{{l}{\oe i}{f}{s}{\r{g}}{ou}{\textscy}{\textsubring{r}}} \dicPos{m}[6] \dicFlx{(‑ar, ‑ar)}[56] \dicDirectTranslationCS{listnatý les}
\dicEntry[laufsöngvari] \dicTerm{lauf··söngv·ar|i} \dicIPA{{l}{\oe i}{f}{s}{\oe i}{\ng}{\r{g}}{v}{a}{r}{\textsci}} \dicPos{m}[1] \dicFlx{(‑a, ‑ar)}[13] \dicFieldCat{zool.} \dicDirectTranslationCS{budníček větší} \textit{(l.~{\textLA{Phylloscopus trochilus}})}  \dicsymPhoto\ 
\dicFigure{ds_image_laufsongvari_0_1.jpg}{Laufsöngvari}{Laufsöngvari - Vogelartinfo, CC BY-SA 3.0}
\dicEntry[lauftré] \dicTerm{lauf··tré} \dicIPA{{l}{\oe i}{v}{t\smash{\textsuperscript{h}}}{r}{j}{\textepsilon}} \dicPos{n}[2] \dicFlx{(‑s, ‑)}[36] \dicDirectTranslationCS{listnatý strom, listnáč}
\dicEntry[laug] \dicTerm{laug\smash{\textsuperscript{1}}} \dicIPA{{l}{\oe i}{\textlengthmark}{x}} \dicPos{f}[4] \dicFlx{(‑ar, ‑ar)}[1] \textbf{1.} \dicSynonym{sundlaug} \dicDirectTranslationCS{(plavecký) bazén, koupaliště, plovárna}  \textbf{2.} \dicSynonym*{heit lind} \dicDirectTranslationCS{vřídlo, horký pramen}
\dicEntry[laug] \dicTerm{laug\smash{\textsuperscript{2}}} \dicIPA{{l}{\oe i}{\textlengthmark}{x}} \dicPos{v} \dicFlx{ind pf sg 1 pers} \dicLink{ljúga}
\dicEntry[lauga] \dicTerm{laug|a} \dicIPA{{l}{\oe i}{\textlengthmark}{\textbabygamma}{a}} \dicPos{v}[1] \dicFlx{(‑aði)}[1] \dicFlx{acc} \dicLangCat{zast.} \dicSynonym{þvo} \dicDirectTranslationCS{(vy)koupat} \dicExampleIS{lauga barn} \dicExampleCS{umýt dítě}
\dicEntry[laugardagskvöld] \dicTerm{laugar·dags··kvöld} \dicIPA{{l}{\oe i}{\textlengthmark}{\textbabygamma}{a}{r}{\textsubring{d}}{a}{x}{s}{k\smash{\textsuperscript{h}}}{v}{\oe}{l}{\textsubring{d}}} \dicPos{n}[2] \dicFlx{(‑s, ‑)}[5] \dicDirectTranslationCS{sobotní večer, sobota večer}
\dicEntry[laugardagur] \dicTerm{laugar··dag|ur} \dicIPA{{l}{\oe i}{\textlengthmark}{\textbabygamma}{a}{r}{\textsubring{d}}{a}{\textbabygamma}{\textscy}{\textsubring{r}}} \dicPos{m}[6] \dicFlx{(‑s, ‑ar)}[62] \dicDirectTranslationCS{sobota}
\dicEntry[Laugavegur] \dicTerm{Lauga··veg|ur} \dicIPA{{l}{\oe i}{\textlengthmark}{\textbabygamma}{a}{v}{\textepsilon}{\textbabygamma}{\textscy}{\textsubring{r}}} \dicPos{m}[10] \dicFlx{(‑ar\,/\addthin ‑s)}[28] \dicFlx{prop} \dicDirectTranslationCS{Laugavegur} \dicIndirectTranslationCS{(hlavní nákupní třída v~Reykjavíku)}
\dicEntry[lauk] \dicTerm{lauk} \dicIPA{{l}{\oe i}{\textlengthmark}{\r{g}}} \dicPos{v} \dicFlx{ind pf sg 1 pers} \dicLink{ljúka}
\dicEntry[laukur] \dicTerm{lauk|ur} \dicIPA{{l}{\oe i}{\textlengthmark}{\r{g}}{\textscy}{\textsubring{r}}} \dicPos{m}[6] \dicFlx{(‑s, ‑ar)}[22] \dicFieldCat{bot.} \dicDirectTranslationCS{cibule, cibule kuchyňská} \textit{(l.~{\textLA{Allium cepa}})};  \dicPhraseIS{laukur ættar} \dicLangCat{přen.} \dicDirectTranslationCS{pýcha rodiny}
\dicEntry[laum] \dicTerm{laum} \dicIPA{{l}{\oe i}{\textlengthmark}{\textsubring{m}}} \dicPos{n}[2] \dicFlx{(‑s)}[2] \dicPhraseIS{í laumi} \dicFlx{adv} \dicDirectTranslationCS{potají, v~tajnosti, tajně}
\dicEntry[lauma] \dicTerm{laum|a} \dicsymFrequent\  \dicIPA{{l}{\oe i}{\textlengthmark}{m}{a}} \dicPos{v}[1] \dicFlx{(‑aði)}[1] \dicFlx{dat} \dicDirectTranslationCS{(nenápadně) dát, vsunout, podstrčit, propašovat} \dicExampleIS{lauma bréfi að e‑m} \dicExampleCS{nenápadně (komu) podstrčit dopis};  \dicIdiom{laumast}{ \dicPhraseIS{laumast}} \dicFlx{refl} \dicSynonym*{læðast} \dicDirectTranslationCS{(v)plížit se, (v)krást se} \dicExampleIS{laumast úr bænum} \dicExampleCS{vykrást se z~města}
\dicEntry[laumufarþegi] \dicTerm{laumu··far·þeg|i} \dicIPA{{l}{\oe i}{\textlengthmark}{m}{\textscy}{f}{a}{\textsubring{r}}{\texttheta}{ei}{\textsci}} \dicPos{m}[1] \dicFlx{(‑a, ‑ar)}[1] \dicDirectTranslationCS{černý pasažér, černá pasažérka}
\dicEntry[laumulegur] \dicTerm{laumu··legur} \dicIPA{{l}{\oe i}{\textlengthmark}{m}{\textscy}{l}{\textepsilon}{\textbabygamma}{\textscy}{\textsubring{r}}} \dicPos{adj}[1]\dicFlx{}[-8] \dicDirectTranslationCS{tajný, skrytý, kradmý}
\dicEntry[laumuspil] \dicTerm{laumu··spil} \dicIPA{{l}{\oe i}{\textlengthmark}{m}{\textscy}{s}{\textsubring{b}}{\textsci}{\textsubring{l}}} \dicPos{n}[2] \dicFlx{(‑s, ‑)}[5] \dicDirectTranslationCS{tajnosti, tajnůstkaření}
\dicEntry[laun] \dicTerm{laun\smash{\textsuperscript{1}}} \dicIPA{{l}{\oe i}{\textlengthmark}{\textsubring{n}}} \dicPos{f}[4] \dicFlx{(‑ar)}[3] \dicSynonym{leynd} \dicDirectTranslationCS{tajemství};  \dicPhraseIS{á laun} \dicFlx{adv} \dicDirectTranslationCS{potají, v~tajnosti}
\dicEntry[laun] \dicTerm{laun\smash{\textsuperscript{2}}} \dicsymFrequent\  \dicIPA{{l}{\oe i}{\textlengthmark}{\textsubring{n}}} \dicPos{n}[2] \dicFlx{pl}[1] \textbf{1.} \dicSynonym{kaup} \dicDirectTranslationCS{plat, výplata, mzda} \dicExampleIS{há laun} \dicExampleCS{vysoký plat}  \textbf{2.} \dicSynonym{endurgjald} \dicDirectTranslationCS{vděk, odměna}  \textbf{3.} \dicSynonym{borgun} \dicDirectTranslationCS{(peněžitá) odměna, honorář}
\dicEntry[launa] \dicTerm{laun|a} \dicIPA{{l}{\oe i}{\textlengthmark}{n}{a}} \dicPos{v}[1] \dicFlx{(‑aði)}[1] \dicFlx{dat + acc} \textbf{1.} \dicSynonym{endurgjalda} \dicDirectTranslationCS{odplatit, oplatit, odvděčit se};  \dicPhraseIS{launa e‑m e‑ð} \dicDirectTranslationCS{odvděčit se (komu čím), oplatit (komu čím)}  \textbf{2.} \dicDirectTranslationCS{ohodnotit, honorovat, platit} \dicExampleIS{illa launað starf} \dicExampleCS{špatně honorovaná práce}
\dicEntry[launaauki] \dicTerm{launa··auk|i} \dicIPA{{l}{\oe i}{\textlengthmark}{n}{a}{\oe i}{\r{\textObardotlessj}}{\textsci}} \dicPos{m}[1] \dicFlx{(‑a)}[3] \dicDirectTranslationCS{příplatek\,/\addthin přídavek k~platu, prémie, bonus}
\dicEntry[launadeila] \dicTerm{launa··deil|a} \dicIPA{{l}{\oe i}{\textlengthmark}{n}{a}{\textsubring{d}}{ei}{l}{a}} \dicPos{f}[1] \dicFlx{(‑u, ‑ur)}[13] \dicDirectTranslationCS{mzdový spor}
\dicEntry[launaflokkur] \dicTerm{launa··flokk|ur} \dicIPA{{l}{\oe i}{\textlengthmark}{n}{a}{f}{l}{\textopeno}{h}{\r{g}}{\textscy}{\textsubring{r}}} \dicPos{m}[6] \dicFlx{(‑s, ‑ar)}[8] \dicDirectTranslationCS{platové zařazení, platová skupina\,/\addthin třída}
\dicEntry[launagreiðsla] \dicTerm{launa··greiðsl|a} \dicIPA{{l}{\oe i}{\textlengthmark}{n}{a}{\r{g}}{r}{ei}{ð}{s}{\textsubring{d}}{l}{a}} \dicPos{f}[1] \dicFlx{(‑u, ‑ur)}[13] \dicDirectTranslationCS{výplata mzdy}
\dicEntry[launahækkun] \dicTerm{launa··hækk|un} \dicIPA{{l}{\oe i}{\textlengthmark}{n}{a}{h}{a}{i}{h}{\r{g}}{\textscy}{\textsubring{n}}} \dicPos{f}[7] \dicFlx{(‑unar, ‑anir)}[8] \dicDirectTranslationCS{zvýšení platu, mzdový růst}
\dicEntry[launakjör] \dicTerm{launa··kjör} \dicIPA{{l}{\oe i}{\textlengthmark}{n}{a}{c\smash{\textsuperscript{h}}}{\oe}{\textsubring{r}}} \dicPos{n}[2] \dicFlx{pl}[9] \dicDirectTranslationCS{platové\,/\addthin mzdové podmínky}
\begin{xtolerant}{}{1pt}
\dicEntry[launakostnaður] \dicTerm{launa··kost·nað|ur}\addthinS\dicIPA{{l}{\oe i}{\textlengthmark}{n}{a}{k\smash{\textsuperscript{h}}}{\textopeno}{s}{\textsubring{d}}{n}{a}{ð}{\textscy}{\textsubring{r}}}\addthinS\dicPos{m}[10]\addthinS\dicFlx{(‑ar)}[9] \dicDirectTranslationCS{mzdové náklady}
\end{xtolerant}
\dicEntry[launalækkun] \dicTerm{launa··lækk|un} \dicIPA{{l}{\oe i}{\textlengthmark}{n}{a}{l}{a}{i}{h}{\r{g}}{\textscy}{\textsubring{n}}} \dicPos{f}[7] \dicFlx{(‑unar)}[9] \dicDirectTranslationCS{snížení platu, mzdový pokles}
\dicEntry[launamaður] \dicTerm{launa··|maður} \dicIPA{{l}{\oe i}{\textlengthmark}{n}{a}{m}{a}{ð}{\textscy}{\textsubring{r}}} \dicPos{m}[13] \dicFlx{(‑manns, ‑menn)}[2] \dicDirectTranslationCS{zaměstnanec, zaměstnankyně}
\dicEntry[launaskrá] \dicTerm{launa··skrá} \dicIPA{{l}{\oe i}{\textlengthmark}{n}{a}{s}{\r{g}}{r}{au}} \dicPos{f}[4] \dicFlx{(‑r\,/\addthin ‑ar, ‑r)}[21] \dicDirectTranslationCS{výplatní listina}
\dicEntry[launaskrið] \dicTerm{launa··skrið} \dicIPA{{l}{\oe i}{\textlengthmark}{n}{a}{s}{\r{g}}{r}{\textsci}{\texttheta}} \dicPos{n}[2] \dicFlx{(‑s)}[2] \dicFieldCat{ekon.} \dicDirectTranslationCS{stoupání mezd}
\dicEntry[launastigi] \dicTerm{launa··stig|i} \dicIPA{{l}{\oe i}{\textlengthmark}{n}{a}{s}{\textsubring{d}}{i}{j}{\textsci}} \dicPos{m}[1] \dicFlx{(‑a)}[3] \dicDirectTranslationCS{mzdový tarif}
\dicEntry[launatengdur] \dicTerm{launa··tengdur} \dicIPA{{l}{\oe i}{\textlengthmark}{n}{a}{t\smash{\textsuperscript{h}}}{\textepsilon}{\ng}{\textsubring{d}}{\textscy}{\textsubring{r}}} \dicPos{adj}[2]\dicFlx{}[-17] \dicDirectTranslationCS{mzdový, spojený se mzdou (náklady ap.)};  \dicPhraseIS{launatengd gjöld} \dicDirectTranslationCS{příspěvky zaměstnavatele}
\dicEntry[launauppbót] \dicTerm{launa··upp·|bót} \dicIPA{{l}{\oe i}{\textlengthmark}{n}{a}{\textscy}{h}{\textsubring{b}}{ou}{\textsubring{d}}} \dicPos{f}[8] \dicFlx{(‑bótar, ‑bætur)}[5] \dicDirectTranslationCS{platový bonus}
\dicEntry[launráð] \dicTerm{laun··ráð} \dicIPA{{l}{\oe i}{n}{r}{au}{\texttheta}} \dicPos{n}[2] \dicFlx{pl}[1] \dicSynonym{ráðabrugg} \dicDirectTranslationCS{pikle, intriky, komplot}
\dicEntry[launsátur] \dicTerm{laun··sátur} \dicIPA{{l}{\oe i}{n}{s}{au}{\textsubring{d}}{\textscy}{\textsubring{r}}} \dicPos{n}[2] \dicFlx{(‑s, ‑)}[25] \dicSynonym{fyrirsát} \dicDirectTranslationCS{léčka, nástraha}
\dicEntry[launung] \dicTerm{laun··ung} \dicIPA{{l}{\oe i}{\textlengthmark}{n}{u}{\ng}{\r{g}}} \dicPos{f}[4] \dicFlx{(‑ar)}[3] \dicSynonym{leynd} \dicDirectTranslationCS{tajemství, tajnost};  \dicPhraseIS{með launung} \dicFlx{adv} \dicSynonym*{leynilega} \dicDirectTranslationCS{tajně};  \dicPhraseIS{e‑m er engin launung á e‑u} \dicFlx{impers} \dicDirectTranslationCS{(kdo) se netají (čím)}
\dicEntry[launþegi] \dicTerm{laun··þeg|i} \dicIPA{{l}{\oe i}{n}{\texttheta}{ei}{\textsci}} \dicPos{m}[1] \dicFlx{(‑a, ‑ar)}[1] \dicDirectTranslationCS{zaměstnanec, zaměstnankyně, pracovník, pracovnice}
\dicEntry[laus] \dicTerm{laus} \dicsymFrequent\  \dicIPA{{l}{\oe i}{\textlengthmark}{s}} \dicPos{adj}[5] \dicFlx{(f ‑)}[1] \textbf{1.} \dicSynonym{óbundinn} \dicDirectTranslationCS{volný, uvolněný, povolený} \dicExampleIS{Steinninn er orðinn laus frá veggnum.} \dicExampleCS{Kámen se už uvolnil ze zdi.}  \textbf{2.} \dicSynonym{frjáls} \dicDirectTranslationCS{volný, svobodný} \dicExampleIS{Ertu laus í kvöld?} \dicExampleCS{Máš večer čas?};  \dicPhraseIS{ganga laus} \dicDirectTranslationCS{jít volně (pes ap.)};  \dicPhraseIS{láta e‑n lausan} \dicDirectTranslationCS{pustit (koho) na svobodu};  \dicPhraseIS{vera laus við e‑ð} \dicDirectTranslationCS{zbavit se (čeho), oprostit se od (čeho)}  \textbf{3.} \dicSynonym*{auður} \dicDirectTranslationCS{volný, prázdný, neobsazený} \dicExampleIS{Er þetta sæti laust?} \dicExampleCS{Je toto místo volné?} \dicAntonym{upptekinn}  \textbf{4.} \dicSynonym{léttur} \dicDirectTranslationCS{lehký} \dicExampleIS{laust högg} \dicExampleCS{lehká rána}  \textbf{5.} \dicSynonym{óáreiðanlegur} \dicDirectTranslationCS{nespolehlivý, nestabilní};  \dicPhraseIS{vera á lausu} \dicDirectTranslationCS{být nezadaný};  \dicPhraseIS{vera laus á kostunum} \dicSynonym{lauslátur} \dicDirectTranslationCS{mít volné mravy} \dicIndirectTranslationCS{(zvláště o~ženě)}  \textbf{6.} \dicSynonym{örlátur} \dicDirectTranslationCS{štědrý};  \dicPhraseIS{vera laus á fé} \dicDirectTranslationCS{být štědrý}
\dicEntry[lausafé] \dicTerm{lausa··|fé} \dicIPA{{l}{\oe i}{\textlengthmark}{s}{a}{f}{j}{\textepsilon}} \dicPos{n}[3] \dicFlx{(‑fjár)}[5] \dicDirectTranslationCS{hotovost}
\dicEntry[lausafjárstaða] \dicTerm{lausa··fjár·|staða} \dicIPA{{l}{\oe i}{\textlengthmark}{s}{a}{f}{j}{au}{\textsubring{r}}{s}{\textsubring{d}}{a}{ð}{a}} \dicPos{f}[1] \dicFlx{(‑stöðu)}[2] \dicFieldCat{ekon.} \dicDirectTranslationCS{likvidita}
\dicEntry[lausaleikur] \dicTerm{lausa··leik|ur} \dicIPA{{l}{\oe i}{\textlengthmark}{s}{a}{l}{ei}{\r{g}}{\textscy}{\textsubring{r}}} \dicPos{m}[6] \dicFlx{(‑s)}[17] \dicPhraseIS{eiga\,/\addthin eignast barn í lausaleik} \dicDirectTranslationCS{mít nemanželské dítě}
\dicEntry[lausamennska] \dicTerm{lausa··mennsk|a} \dicIPA{{l}{\oe i}{\textlengthmark}{s}{a}{m}{\textepsilon}{n}{s}{\r{g}}{a}} \dicPos{f}[1] \dicFlx{(‑u)}[5] \dicDirectTranslationCS{práce na volné noze}
\dicEntry[lausamjöll] \dicTerm{lausa··|mjöll} \dicIPA{{l}{\oe i}{\textlengthmark}{s}{a}{m}{j}{\oe}{\textsubring{d}}{\textsubring{l}}} \dicPos{f}[7] \dicFlx{(‑mjallar)}[19] \dicDirectTranslationCS{čerstvě napadlý sníh, prašan}
\dicEntry[lausasala] \dicTerm{lausa··|sala} \dicIPA{{l}{\oe i}{\textlengthmark}{s}{a}{s}{a}{l}{a}} \dicPos{f}[1] \dicFlx{(‑sölu, ‑sölur)}[20] \dicDirectTranslationCS{volný prodej}
\dicEntry[lausavísa] \dicTerm{lausa··vís|a} \dicIPA{{l}{\oe i}{\textlengthmark}{s}{a}{v}{i}{s}{a}} \dicPos{f}[1] \dicFlx{(‑u, ‑ur)}[13] \dicFieldCat{lit.} \dicDirectTranslationCS{epigram} \dicIndirectTranslationCS{(často od neznámého autora, velmi populární na Islandu)}
\dicEntry[lausblaðabók] \dicTerm{laus·blaða··|bók} \dicIPA{{l}{\oe i}{s}{\textsubring{b}}{l}{a}{ð}{a}{\textsubring{b}}{ou}{\r{g}}} \dicPos{f}[8] \dicFlx{(‑bókar, ‑bækur)}[5] \dicDirectTranslationCS{zápisník\,/\addthin sešit s~volnými listy}

\dicEntry[lausblaðamappa] \textls[15]{\dicTerm{laus·blaða··|mappa} \dicIPA{{l}{\oe i}{s}{\textsubring{b}}{l}{a}{ð}{a}{m}{a}{h}{\textsubring{b}}{a}} \dicPos{f}[1]} \dicFlx{(‑möppu, ‑möpp\-ur)}[8] \dicDirectTranslationCS{pořadač, šanon}

\dicEntry[lauslátur] \dicTerm{laus··látur} \dicIPA{{l}{\oe i}{s}{l}{au}{\textsubring{d}}{\textscy}{\textsubring{r}}} \dicPos{adj}[1]\dicFlx{}[-1] \dicDirectTranslationCS{nezávazný (flirt ap.), promiskuitní}
\dicEntry[lauslegur] \dicTerm{laus··legur} \dicIPA{{l}{\oe i}{s}{l}{\textepsilon}{\textbabygamma}{\textscy}{\textsubring{r}}} \dicPos{adj}[1]\dicFlx{}[-8] \textbf{1.} \dicSynonym{yfirborðslegur} \dicDirectTranslationCS{zběžný, letmý, hrubý (odhad ap.)} \dicExampleIS{lauslegur uppdráttur} \dicExampleCS{zběžný nákres}  \textbf{2.} \dicSynonym{óstöðugur} \dicDirectTranslationCS{volný, nepřipevněný}
\dicEntry[lauslæti] \dicTerm{laus··læti} \dicIPA{{l}{\oe i}{s}{l}{a}{i}{\textsubring{d}}{\textsci}} \dicPos{n}[2] \dicFlx{(‑s)}[20] \dicDirectTranslationCS{nezávaznost, promiskuita}
\dicEntry[lausmáll] \dicTerm{laus··máll} \dicIPA{{l}{\oe i}{s}{m}{au}{\textsubring{d}}{\textsubring{l}}} \dicPos{adj}[8]\dicFlx{}[-1] \textbf{1.} \dicSynonym{málgefinn} \dicDirectTranslationCS{upovídaný, užvaněný}  \textbf{2.} \dicSynonym*{óorðheldinn} \dicDirectTranslationCS{indiskrétní, netaktní}
\dicEntry[lausn] \dicTerm{lausn} \dicsymFrequent\  \dicIPA{{l}{\oe i}{s}{\textsubring{d}}{\textsubring{n}}} \dicPos{f}[7] \dicFlx{(‑ar, ‑ir)}[1] \textbf{1.} \dicSynonym{frelsun} \dicDirectTranslationCS{osvobození, vysvobození} \dicExampleIS{lausn úr þrældómi} \dicExampleCS{vysvobození z~otroctví}  \textbf{2.} \dicDirectTranslationCS{uvolnění, propuštění (z~práce ap.)} \dicExampleIS{lausn frá embætti} \dicExampleCS{uvolnění z~úřadu};  \dicPhraseIS{biðjast lausnar} \dicFlx{refl} \dicDirectTranslationCS{žádat o~uvolnění (z~pozice ap.)}  \textbf{3.} \dicSynonym{úrslit} \dicDirectTranslationCS{(vy)řešení, rozřešení, rozuzlení} \dicExampleIS{lausn á vandanum} \dicExampleCS{vyřešení problému}  \textbf{4.} \dicSynonym{endurlausn} \dicDirectTranslationCS{vykoupení, spasení}  \textbf{5.} \dicFieldCat{chem.} \dicDirectTranslationCS{roztok};  \dicPhraseIS{mettuð lausn} \dicFieldCat{chem.} \dicDirectTranslationCS{nasycený roztok}
\dicEntry[lausnarbeiðni] \dicTerm{lausnar··beiðn|i} \dicIPA{{l}{\oe i}{s}{\textsubring{d}}{n}{a}{r}{\textsubring{b}}{ei}{ð}{n}{\textsci}} \dicPos{f}[3] \dicFlx{(‑i, ‑ir)}[1] \dicDirectTranslationCS{výpověď, žádost o~uvolnění (z~práce ap.)};  \dicPhraseIS{senda inn lausnarbeiðni} \dicDirectTranslationCS{podat výpověď (z~práce ap.)}
\dicEntry[lausnargjald] \dicTerm{lausnar··|gjald} \dicIPA{{l}{\oe i}{s}{\textsubring{d}}{n}{a}{r}{\r{\textObardotlessj}}{a}{l}{\textsubring{d}}} \dicPos{n}[2] \dicFlx{(‑gjalds, ‑gjöld)}[8] \dicDirectTranslationCS{výkupné}
\dicEntry[lausnarhnappur] \dicTerm{lausnar··hnapp|ur} \dicIPA{{l}{\oe i}{s}{\textsubring{d}}{n}{a}{\textsubring{r}}{\textsubring{n}}{a}{h}{\textsubring{b}}{\textscy}{\textsubring{r}}} \dicPos{m}[6] \dicFlx{(‑s, ‑ar)}[25] \dicFieldCat{poč.} \dicDirectTranslationCS{klávesa escape}
\dicEntry[lausnari] \dicTerm{lausn··ar|i} \dicIPA{{l}{\oe i}{s}{\textsubring{d}}{n}{a}{r}{\textsci}} \dicPos{m}[1] \dicFlx{(‑a, ‑ar)}[13] \dicFieldCat{náb.} \dicSynonym{frelsari} \dicDirectTranslationCS{spasitel, vykupitel}
\dicEntry[laust] \dicTerm{laust\smash{\textsuperscript{1}}} \dicIPA{{l}{\oe i}{s}{\textsubring{d}}} \dicPos{v} \dicFlx{ind pf sg 1 pers} \dicLink{ljósta}
\dicEntry[laust] \dicTerm{laus|t\smash{\textsuperscript{2}}} \dicIPA{{l}{\oe i}{s}{\textsubring{d}}} \dicPos{adv} \dicFlx{(comp ‑ar, sup ‑ast)} \textbf{1.} \dicDirectTranslationCS{lehce, slabě, mírně} \dicExampleIS{berja laust að dyrum} \dicExampleCS{zaklepat lehce na dveře}  \textbf{2.} \dicDirectTranslationCS{krátce, těsně} \dicExampleIS{laust fyrir klukkan fimm} \dicExampleCS{krátce před pátou hodinou}
\dicEntry[lausung] \dicTerm{laus··ung} \dicIPA{{l}{\oe i}{\textlengthmark}{s}{u}{\ng}{\r{g}}} \dicPos{f}[4] \dicFlx{(‑ar)}[3] \textbf{1.} \dicSynonym*{staðfestuleysi} \dicDirectTranslationCS{nestabilita, nestálost}  \textbf{2.} \dicSynonym{lauslæti} \dicDirectTranslationCS{promiskuita};  \dicIdiom{lausung}{ \dicPhraseIS{gjalda lausung við lygi}} \dicLangCat{přen.} \dicDirectTranslationCS{oplatit faleš lží}
\dicEntry[laut] \dicTerm{laut\smash{\textsuperscript{1}}} \dicIPA{{l}{\oe i}{\textlengthmark}{\textsubring{d}}} \dicPos{f}[7] \dicFlx{(‑ar, ‑ir)}[6] \dicSynonym{lægð} \dicDirectTranslationCS{prohlubeň, dolík}
\dicEntry[laut] \dicTerm{laut\smash{\textsuperscript{2}}} \dicIPA{{l}{\oe i}{\textlengthmark}{\textsubring{d}}} \dicPos{v} \dicFlx{ind pf sg 1 pers} \dicLink{lúta}
\dicEntry[lautinant] \dicTerm{lautinant} \dicIPA{{l}{\oe i}{\textlengthmark}{\textsubring{d}}{\textsci}{n}{a}{\textsubring{n}}{\textsubring{d}}} \dicPos{m}[4] \dicFlx{(‑s, ‑ar)}[11] \dicFieldCat{voj.} \dicDirectTranslationCS{poručík, poručice}
\dicEntry[lax] \dicTerm{lax}\dicTerm{, laxfiskur} \dicIPA{{l}\-{a}\-{x}\-{s}\-} \dicPos{m}[4] \dicFlx{(‑, ‑ar)}[29] \dicFieldCat{zool.} \dicDirectTranslationCS{losos, losos obecný} \textit{(l.~{\textLA{Salmo salar}})}
\dicEntry[laxá] \dicTerm{lax··á} \dicIPA{{l}{a}{x}{s}{au}} \dicPos{f}[4] \dicFlx{(‑r, ‑r)}[18] \dicDirectTranslationCS{lososí řeka}
\dicEntry[laxerolía] \dicTerm{laxer··olí|a} \dicIPA{{l}{a}{x}{s}{\textepsilon}{r}{\textopeno}{l}{i}{j}{a}} \dicPos{f}[1] \dicFlx{(‑u)}[5] \dicDirectTranslationCS{ricinový olej}
\dicEntry[laxfiskur] \dicTerm{lax··fisk|ur} \dicIPA{{l}{a}{x}{s}{f}{\textsci}{s}{\r{g}}{\textscy}{\textsubring{r}}} \dicPos{m}[6] \dicFlx{(‑s, ‑ar)}[10] \dicLink{lax}
\dicEntry[laxveiði] \dicTerm{lax··veið|i} \dicIPA{{l}{a}{x}{s}{v}{ei}{ð}{\textsci}} \dicPos{f}[2] \dicFlx{(‑i\,/\addthin ‑ar, ‑ar)}[3] \textbf{1.} \dicDirectTranslationCS{lososí úlovek}  \textbf{2.} \dicDirectTranslationCS{chytání\,/\addthin lov lososů}
\dicEntry[lá] \dicTerm{lá\smash{\textsuperscript{1}}} \dicIPA{{l}{au}{\textlengthmark}} \dicPos{v}[2] \dicFlx{(‑ði, ‑ð)}[117] \dicFlx{dat} \dicSynonym{ásaka} \dicDirectTranslationCS{vyčíst, vyčítat, dávat za vinu} \dicIndirectTranslationCS{(používané v~otázce nebo záporu)} \dicExampleIS{Ég lái þér það ekki.} \dicExampleCS{Já ti to nevyčítám.};  \dicIdiom{lást}{ \dicPhraseIS{e‑m láist að (gera e‑ð)}} \dicFlx{refl impers} \dicDirectTranslationCS{(kdo) zapomíná (udělat (co)), (komu) vypadá z~paměti (udělat (co))} \dicExampleIS{Mér láðist að senda bréfið.} \dicExampleCS{Vypadlo mi (z~paměti), že mám poslat dopis.}
\dicEntry[lá] \dicTerm{lá\smash{\textsuperscript{2}}} \dicIPA{{l}{au}{\textlengthmark}} \dicPos{v} \dicFlx{ind pf sg 1 pers} \dicLink{liggja}
\dicEntry[láð] \dicTerm{láð} \dicIPA{{l}{au}{\textlengthmark}{\texttheta}} \dicPos{n}[2] \dicFlx{(‑s, ‑)}[5] \dicSynonym{land} \dicDirectTranslationCS{souš, země} \dicExampleIS{um láð og lög} \dicExampleCS{po souši a~po vodě}
\dicEntry[lág] \dicTerm{lág} \dicIPA{{l}{au}{\textlengthmark}} \dicPos{f}[4] \dicFlx{(‑ar, ‑ar)}[1] \dicSynonym{dæld} \dicDirectTranslationCS{dolina, kotlina};  \dicPhraseIS{láta e‑ð liggja í láginni} \dicLangCat{přen.} \dicDirectTranslationCS{nechat (co) ležet u~ledu}
\dicEntry[lágfiðla] \dicTerm{lág··fiðl|a} \dicIPA{{l}{au}{\textlengthmark}{f}{\textsci}{ð}{l}{a}} \dicPos{f}[1] \dicFlx{(‑u, ‑ur)}[7] \dicFieldCat{hud.} \dicDirectTranslationCS{viola}
\dicEntry[lágkúra] \dicTerm{lág··kúr|a} \dicIPA{{l}{au}{\textlengthmark}{k\smash{\textsuperscript{h}}}{u}{r}{a}} \dicPos{f}[1] \dicFlx{(‑u)}[5] \dicSynonym*{svipleysa} \dicDirectTranslationCS{banalita, přízemnost}
\dicEntry[lágkúrulegur] \textls[15]{\dicTerm{lág·kúru··legur} \dicIPA{{l}{au}{\textlengthmark}{k\smash{\textsuperscript{h}}}{u}{r}{\textscy}{l}{\textepsilon}{\textbabygamma}{\textscy}{\textsubring{r}}} \dicPos{adj}[1]\dicFlx{}[-8] \textbf{1.} \dicSynonym*{flatneskjulegur} \dicDirectTranslationCS{banální, všední, nezajímavý}} \dicExampleIS{lágkúrulegt hús} \dicExampleCS{všední dům}  \textbf{2.} \dicSynonym{lítilmótlegur} \dicDirectTranslationCS{přízemní, nízký, mrzký} \dicExampleIS{lágkúrulegur maður} \dicExampleCS{přízemní člověk}
\dicEntry[lágkveða] \dicTerm{lág··kveð|a} \dicIPA{{l}{au}{\textlengthmark}{k\smash{\textsuperscript{h}}}{v}{\textepsilon}{ð}{a}} \dicPos{f}[1] \dicFlx{(‑u, ‑ur)}[19] \dicFieldCat{lit.} \dicDirectTranslationCS{sudá stopa se slabším přízvukem}
\dicEntry[láglaunafólk] \dicTerm{lág·launa··fólk} \dicIPA{{l}{au}{\textlengthmark}{l}{\oe i}{n}{a}{f}{ou}{\textsubring{l}}{\r{g}}} \dicPos{n}[2] \dicFlx{(‑s)}[2] \dicDirectTranslationCS{osoby s~nízkým příjmem}
\dicEntry[láglendi] \dicTerm{lág··lendi} \dicIPA{{l}{au}{\textlengthmark}{l}{\textepsilon}{n}{\textsubring{d}}{\textsci}} \dicPos{n}[2] \dicFlx{(‑s)}[20] \dicDirectTranslationCS{nížina} \dicAntonym{hálendi}
\dicEntry[lágmark] \dicTerm{lág··|mark} \dicIPA{{l}{au}{\textlengthmark}{m}{a}{\textsubring{r}}{\r{g}}} \dicPos{n}[2] \dicFlx{(‑marks, ‑mörk)}[8] \dicDirectTranslationCS{minimum} \dicExampleIS{Kuldinn er í algjöru lágmarki.} \dicExampleCS{Zima je na úplném minimu.} \dicAntonym{hámark}
\dicEntry[lágmarksframfærsla] \dicTerm{lág·marks··fram·færsl|a} \dicIPA{{l}\-{au}\-{\textlengthmark}\-{m}\-{a}\-{\textsubring{r}}\-{x}\-{s}\-{f}\-{r}\-{a}\-{m}\-{f}\-{a}\-{i}\-{\textsubring{r}}\-{s}\-{\textsubring{d}}\-{l}\-{a}\-} \dicPos{f}[1] \dicFlx{(‑u)}[5] \dicFieldCat{ekon.} \dicDirectTranslationCS{životní minimum}
\dicEntry[lágmarksgjald] \dicTerm{lág·marks··|gjald} \dicIPA{{l}{au}{\textlengthmark}{m}{a}{\textsubring{r}}{x}{s}{\r{\textObardotlessj}}{a}{l}{\textsubring{d}}} \dicPos{n}[2] \dicFlx{(‑gjalds, ‑gjöld)}[8] \dicDirectTranslationCS{minimální sazba}
\begin{xtolerant}{}{1pt}
\dicEntry[lágmarkslaun] \dicTerm{lág·marks··laun}\addthinS\dicIPA{{l}{au}{\textlengthmark}{m}{a}{\textsubring{r}}{x}{s}{l}{\oe i}{\textsubring{n}}}\addthinS\dicPos{n}[2]\addthinS\dicFlx{pl}[1]\addthinS\dicDirectTranslationCS{minimální mzda\hspace{-0.06667em}\,/\addthin plat}
\end{xtolerant}
\dicEntry[lágmarkstímasókn] \dicTerm{lág·marks··tíma·sókn} \dicIPA{{l}{au}{\textlengthmark}{m}{a}{\textsubring{r}}{x}{s}{t\smash{\textsuperscript{h}}}{i}{m}{a}{s}{ou}{h}{\r{g}}{\textsubring{n}}} \dicPos{f}[7] \dicFlx{(‑ar)}[3] \dicDirectTranslationCS{minimální počet hodin školní docházky} \dicIndirectTranslationCS{(pro ukončení kurzu ap.)}
\dicEntry[lágmarksverð] \dicTerm{lág·marks··verð} \dicIPA{{l}{au}{\textlengthmark}{m}{a}{\textsubring{r}}{x}{s}{v}{\textepsilon}{r}{\texttheta}} \dicPos{n}[2] \dicFlx{(‑s)}[2] \dicDirectTranslationCS{minimální cena}
\dicEntry[lágmynd] \dicTerm{lág··mynd} \dicIPA{{l}{au}{\textlengthmark}{m}{\textsci}{n}{\textsubring{d}}} \dicPos{f}[7] \dicFlx{(‑ar, ‑ir)}[1] \dicDirectTranslationCS{reliéf}
\dicEntry[lágnætti] \dicTerm{lág··nætti} \dicIPA{{l}{au}{\textlengthmark}{n}{a}{i}{h}{\textsubring{d}}{\textsci}} \dicPos{n}[2] \dicFlx{(‑s, ‑)}[14] \dicDirectTranslationCS{půlnoc} \dicIndirectTranslationCS{(zvláště v~létě)}
\dicEntry[lágspenna] \dicTerm{lág··spenn|a} \dicIPA{{l}{au}{\textlengthmark}{s}{\textsubring{b}}{\textepsilon}{n}{a}} \dicPos{f}[1] \dicFlx{(‑u)}[5] \dicFieldCat{elek.} \dicDirectTranslationCS{nízké napětí} \dicAntonym{háspenna}
\dicEntry[lágstétt] \dicTerm{lág··stétt} \dicIPA{{l}{au}{\textlengthmark}{s}{\textsubring{d}}{j}{\textepsilon}{h}{\textsubring{d}}} \dicPos{f}[7] \dicFlx{(‑ar, ‑ir)}[1] \dicDirectTranslationCS{nižší (společenská) třída, proletariát} \dicAntonym{hástétt}
\dicEntry[lágt] \dicTerm{lágt} \dicsymFrequent\  \dicIPA{{l}{au}{x}{\textsubring{d}}} \dicPos{adv} \dicFlx{(comp lægra, sup lægst)} \textbf{1.} \dicSynonym*{lítill hljóðstyrkur} \dicDirectTranslationCS{tiše, potichu, tlumeně} \dicExampleIS{tala lágt} \dicExampleCS{mluvit potichu}  \textbf{2.} \dicSynonym{niðri} \dicDirectTranslationCS{nízko, dole}
\dicEntry[lágum] \dicTerm{lágum} \dicIPA{{l}{au}{\textlengthmark}{\textbabygamma}{\textscy}{\textsubring{m}}} \dicPos{v} \dicFlx{ind pf pl 1 pers} \dicLink{liggja}
\dicEntry[lágur] \dicTerm{lágur} \dicsymFrequent\  \dicIPA{{l}{au}{\textlengthmark}{\textscy}{\textsubring{r}}} \dicPos{adj}[10] \dicFlx{(comp lægri, sup lægstur)}[16] \textbf{1.} \dicSynonym{stuttur} \dicDirectTranslationCS{nízký, malý} \dicIndirectTranslationCS{(o~výšce)};  \dicPhraseIS{vera lágur vexti} \dicDirectTranslationCS{být malého vzrůstu}  \textbf{2.} \dicDirectTranslationCS{nízký} \dicIndirectTranslationCS{(o~čísle, ceně ap.)}  \textbf{3.} \dicDirectTranslationCS{tichý, tlumený} \dicExampleIS{lágur söngur} \dicExampleCS{tichý zpěv};  \dicPhraseIS{í lágum hljóðum} \dicFlx{adv} \dicDirectTranslationCS{tlumeně, tiše}  \textbf{4.} \dicDirectTranslationCS{nízký (tón ap.)} \dicExampleIS{lágur tónn} \dicExampleCS{nízký tón};  \dicIdiom{lágur}{ \dicPhraseIS{bera\,/\addthin bíða lægri hlut}} \dicLangCat{přen.} \dicDirectTranslationCS{tahat za kratší konec provazu}
\dicEntry[lágvaxinn] \dicTerm{lág··vaxinn} \dicIPA{{l}{au}{\textlengthmark}{v}{a}{x}{s}{\textsci}{\textsubring{n}}} \dicPos{adj}[6]\dicFlx{}[-3] \dicDirectTranslationCS{malý (vzrůstem), (jsoucí) malého vzrůstu} \dicAntonym{hávaxinn}
\dicEntry[lágþrýstingur] \dicTerm{lág··þrýst·ing|ur} \dicIPA{{l}{au}{\textlengthmark}{\texttheta}{r}{i}{s}{\textsubring{d}}{i}{\ng}{\r{g}}{\textscy}{\textsubring{r}}} \dicPos{m}[6] \dicFlx{(‑s)}[9] \textbf{1.} \dicFieldCat{med.} \dicDirectTranslationCS{nízký tlak}  \textbf{2.} \dicFieldCat{meteo.} \dicDirectTranslationCS{nízký tlak}
\dicEntry[lágþrýstisvæði] \dicTerm{lág·þrýsti··svæði} \dicIPA{{l}{au}{\textlengthmark}{\texttheta}{r}{i}{s}{\textsubring{d}}{\textsci}{s}{v}{a}{i}{ð}{\textsci}} \dicPos{n}[2] \dicFlx{(‑s, ‑)}[14] \dicFieldCat{meteo.} \dicDirectTranslationCS{oblast nízkého tlaku} \dicAntonym{háþrýstisvæði}
\dicEntry[lágþýska] \dicTerm{lág··þýsk|a} \dicIPA{{l}{au}{\textlengthmark}{\texttheta}{i}{s}{\r{g}}{a}} \dicPos{f}[1] \dicFlx{(‑u)}[5] \dicDirectTranslationCS{dolní němčina} \dicIndirectTranslationCS{(dialekt němčiny používaný v~severních částech Německa)} \dicAntonym{háþýska}
\dicEntry[lákum] \dicTerm{lákum} \dicIPA{{l}{au}{\textlengthmark}{\r{g}}{\textscy}{\textsubring{m}}} \dicPos{v} \dicFlx{ind pf pl 1 pers} \dicLink{leka}
\dicEntry[lán] \dicTerm{lán} \dicsymFrequent\  \dicIPA{{l}{au}{\textlengthmark}{\textsubring{n}}} \dicPos{n}[2] \dicFlx{(‑s, ‑)}[5] \textbf{1.} \dicSynonym{lánsfé} \dicDirectTranslationCS{(vý)půjčka, úvěr} \dicExampleIS{borga lán} \dicExampleCS{splatit půjčku};  \dicPhraseIS{taka lán} \dicDirectTranslationCS{vzít si půjčku\,/\addthin úvěr}  \textbf{2.} \dicSynonym*{það að lána} \dicDirectTranslationCS{(vy)půjčení};  \dicPhraseIS{fá e‑ð að láni} \dicDirectTranslationCS{(vy)půjčit si (co)};  \dicPhraseIS{hafa e‑ð að láni, vera með e‑ð í láni} \dicDirectTranslationCS{(vy)půjčit si (co), mít (co) půjčené}  \textbf{3.} \dicSynonym{gæfa} \dicDirectTranslationCS{štěstí, štěstěna};  \dicPhraseIS{lánið leikur við e‑n} \dicLangCat{přen.} \dicDirectTranslationCS{(kdo) má štěstí, na (koho) se usmívá štěstí};  \dicPhraseIS{það er lán í óláni að} \dicLangCat{přen.} \dicDirectTranslationCS{je to štěstí v~neštěstí, že}
\dicEntry[lána] \dicTerm{lán|a} \dicsymFrequent\  \dicIPA{{l}{au}{\textlengthmark}{n}{a}} \dicPos{v}[1] \dicFlx{(‑aði)}[1] \dicFlx{dat + acc} \dicDirectTranslationCS{(vy)půjčit, zapůjčit};  \dicPhraseIS{lána e‑m e‑ð} \dicDirectTranslationCS{půjčit (komu co)} \dicExampleIS{lána e‑m bókina} \dicExampleCS{půjčit (komu) knížku};  \dicPhraseIS{fá e‑ð lánað} \dicDirectTranslationCS{(vy)půjčit si (co)};  \dicIdiom{lána}[út]{ \dicPhraseIS{lána út e‑ð}} \dicDirectTranslationCS{(za)půjčovat (co)}
\dicEntry[lánardrottinn] \dicTerm{lánar··drott|inn} \dicIPA{{l}{au}{\textlengthmark}{n}{a}{r}{\textsubring{d}}{r}{\textopeno}{h}{\textsubring{d}}{\textsci}{\textsubring{n}}} \dicPos{m}[6] \dicFlx{(‑ins, ‑nar)}[39] \dicDirectTranslationCS{věřitel(ka)}
\dicEntry[lánasjóður] \dicTerm{lána··sjóð|ur} \dicIPA{{l}{au}{\textlengthmark}{n}{a}{s}{j}{ou}{ð}{\textscy}{\textsubring{r}}} \dicPos{m}[9] \dicFlx{(‑s, ‑ir)}[6] \dicFieldCat{ekon.} \dicDirectTranslationCS{úvěrový fond}
\dicEntry[lánastofnun] \dicTerm{lána··stofn|un} \dicIPA{{l}{au}{\textlengthmark}{n}{a}{s}{\textsubring{d}}{\textopeno}{\textsubring{b}}{n}{\textscy}{\textsubring{n}}} \dicPos{f}[7] \dicFlx{(‑unar, ‑anir)}[8] \dicFieldCat{ekon.} \dicDirectTranslationCS{úvěrová společnost}
\dicEntry[lánaveitandi] \dicTerm{lána··veit·|andi} \dicIPA{{l}{au}{\textlengthmark}{n}{a}{v}{ei}{\textsubring{d}}{a}{n}{\textsubring{d}}{\textsci}} \dicPos{m}[2] \dicFlx{(‑anda, ‑endur)}[1] \dicDirectTranslationCS{věřitel(ka)}
\dicEntry[lánlaus] \dicTerm{lán··laus} \dicIPA{{l}{au}{n}{l}{\oe i}{s}} \dicPos{adj}[5]\dicFlx{}[-1] \dicDirectTranslationCS{smolařský, nemající štěstí}
\dicEntry[lánsamur] \dicTerm{lán··|samur} \dicIPA{{l}{au}{n}{s}{a}{m}{\textscy}{\textsubring{r}}} \dicPos{adj}[1] \dicFlx{(f ‑söm)}[2] \dicSynonym{gæfusamur} \dicDirectTranslationCS{šťastný, mající štěstí}
\dicEntry[lánsfé] \dicTerm{láns··|fé} \dicIPA{{l}{au}{n}{s}{f}{j}{\textepsilon}} \dicPos{n}[3] \dicFlx{(‑fjár)}[5] \dicDirectTranslationCS{půjčka, půjčené peníze}
\dicEntry[lánsfjármögnun] \dicTerm{láns·fjár··mögn|un} \dicIPA{{l}{au}{n}{s}{f}{j}{au}{r}{m}{\oe}{\r{g}}{n}{\textscy}{\textsubring{n}}} \dicPos{f}[7] \dicFlx{(‑unar)}[12] \dicFieldCat{ekon.} \dicDirectTranslationCS{financování dluhu}
\begin{xtolerant}{}{1pt}
\dicEntry[lánskjaravísitala] \dicTerm{láns·kjara··vísi·|tala}\addthinS\dicIPA{{l}\-{au}\-{n}\-{s}\-{c\smash{\textsuperscript{h}}}\-{a}\-{r}\-{a}\-{v}\-{i}\-{s}\-{\textsci}\-{t\smash{\textsuperscript{h}}}\-{a}\-{l}\-{a}\-}\addthinS\dicPos{f}[1]\addthinS\dicFlx{(‑tölu, ‑tölur)}[14] \dicFieldCat{ekon.} \dicDirectTranslationCS{index úvěrových podmínek}
\end{xtolerant}
\dicEntry[lánskjör] \dicTerm{láns··kjör} \dicIPA{{l}{au}{n}{s}{c\smash{\textsuperscript{h}}}{\oe}{\textsubring{r}}} \dicPos{n}[2] \dicFlx{pl}[9] \dicFieldCat{ekon.} \dicDirectTranslationCS{úvěrové podmínky}
\dicEntry[lánstími] \dicTerm{láns··tím|i} \dicIPA{{l}{au}{n}{s}{t\smash{\textsuperscript{h}}}{i}{m}{\textsci}} \dicPos{m}[1] \dicFlx{(‑a)}[3] \textbf{1.} \dicFieldCat{ekon.} \dicDirectTranslationCS{doba splatnosti úvěru}  \textbf{2.} \dicDirectTranslationCS{výpůjční doba (v~knihovně ap.)}
\dicEntry[lánstraust] \dicTerm{láns··traust} \dicIPA{{l}{au}{n}{s}{t\smash{\textsuperscript{h}}}{r}{\oe i}{s}{\textsubring{d}}} \dicPos{n}[2] \dicFlx{(‑s)}[2] \dicFieldCat{ekon.} \dicDirectTranslationCS{úvěruschopnost}
\dicEntry[lántakandi] \dicTerm{lán··tak·|andi} \dicIPA{{l}{au}{n}{t\smash{\textsuperscript{h}}}{a}{\r{g}}{a}{n}{\textsubring{d}}{\textsci}} \dicPos{m}[2] \dicFlx{(‑anda, ‑endur)}[1] \dicSynonym{lánþegi} \dicDirectTranslationCS{vypůjčovatel(ka), příjemce\,/\addthin příjemkyně půjčky}
\dicEntry[lántaki] \dicTerm{lán··tak|i} \dicIPA{{l}{au}{n}{t\smash{\textsuperscript{h}}}{a}{\r{\textObardotlessj}}{\textsci}} \dicPos{m}[1] \dicFlx{(‑a, ‑ar)}[8] \dicSynonym{lánþegi} \dicDirectTranslationCS{vypůjčovatel(ka), příjemce\,/\addthin příjemkyně půjčky}
\dicEntry[lánveiting] \dicTerm{lán··veit·ing} \dicIPA{{l}{au}{n}{v}{ei}{\textsubring{d}}{i}{\ng}{\r{g}}} \dicPos{f}[4] \dicFlx{(‑ar, ‑ar)}[5] \dicFieldCat{ekon.} \dicDirectTranslationCS{poskytnutí\,/\addthin poskytování půjčky\,/\addthin úvěru}
\dicEntry[lánþegi] \dicTerm{lán··þeg|i} \dicIPA{{l}{au}{n}{\texttheta}{ei}{\textsci}} \dicPos{m}[1] \dicFlx{(‑a, ‑ar)}[1] \dicDirectTranslationCS{vypůjčovatel(ka), příjemce\,/\addthin příjemkyně půjčky}
\dicEntry[láréttur] \dicTerm{lá··réttur} \dicIPA{{l}{au}{\textlengthmark}{r}{j}{\textepsilon}{h}{\textsubring{d}}{\textscy}{\textsubring{r}}} \dicPos{adj}[1]\dicFlx{}[-10] \dicDirectTranslationCS{vodorovný, horizontální}
\dicEntry[lárpera] \dicTerm{lár··per|a} \dicIPA{{l}{au}{r}{p\smash{\textsuperscript{h}}}{\textepsilon}{r}{a}} \dicPos{f}[1] \dicFlx{(‑u, ‑ur)}[7] \dicDirectTranslationCS{avokádo} \dicIndirectTranslationCS{(plod)}
\dicEntry[lárperutré] \dicTerm{lár·peru··tré} \dicIPA{{l}{au}{r}{p\smash{\textsuperscript{h}}}{\textepsilon}{r}{\textscy}{t\smash{\textsuperscript{h}}}{r}{j}{\textepsilon}} \dicPos{n}[2] \dicFlx{(‑s, ‑)}[36] \dicFieldCat{bot.} \dicDirectTranslationCS{hruškovec přelahodný} \textit{(l.~{\textLA{Persea americana}})}  \dicsymPhoto\ 
\dicFigure{54128.jpg}{Lárperutré}{Lárperutré - Forest \& Kim [[p:2684;Starr]], Biolib, CC-BY}
\dicEntry[lárviðarlauf] \dicTerm{lár·viðar··lauf} \dicIPA{{l}{au}{r}{v}{\textsci}{ð}{a}{r}{l}{\oe i}{f}} \dicPos{n}[2] \dicFlx{(‑s, ‑)}[5] \dicDirectTranslationCS{bobkový list}
\dicEntry[lárviðarrós] \dicTerm{lár·viðar··rós} \dicIPA{{l}{au}{r}{v}{\textsci}{ð}{a}{r}{ou}{s}} \dicPos{f}[7] \dicFlx{(‑ar, ‑ir)}[1] \dicFieldCat{bot.} \dicSynonym{nería} \dicDirectTranslationCS{oleandr, oleandr obecný} \textit{(l.~{\textLA{Nerium oleander}})}  \dicsymPhoto\ 
\dicFigure{ds_image_larvidarros_0_1.jpg}{Lárviðarrós}{Lárviðarrós - Alvesgaspar, CC BY-SA 3.0}
\dicEntry[lárviður] \dicTerm{lár··við|ur} \dicIPA{{l}{au}{r}{v}{\textsci}{ð}{\textscy}{\textsubring{r}}} \dicPos{m}[10] \dicFlx{(‑ar, ‑ir)}[19] \dicFieldCat{bot.} \dicDirectTranslationCS{vavřín, vavřín vznešený} \textit{(l.~{\textLA{Laurus nobilis}})}  \dicsymPhoto\ 
\dicFigure{55905.jpg}{Lárviður}{Lárviður - Forest \& Kim [[p:2684;Starr]], Biolib, CC-BY}
\dicEntry[lás] \dicTerm{lás} \dicIPA{{l}{au}{\textlengthmark}{s}} \dicPos{m}[4] \dicFlx{(‑s, ‑ar)}[25] \textbf{1.} \dicDirectTranslationCS{zámek (u~dveří ap.)} \dicExampleIS{lásinn á hjólinu} \dicExampleCS{zámek na kolo};  \dicPhraseIS{skella e‑u í lás} \dicDirectTranslationCS{zabouchnout (si) (co) (dveře ap.)}  \textbf{2.} \dicDirectTranslationCS{pojistka} \dicIndirectTranslationCS{(součást střelné zbraně)} \dicExampleIS{lás á byssu} \dicExampleCS{pojistka pistole}
\dicEntry[lásasmiður] \dicTerm{lása··smið|ur} \dicIPA{{l}{au}{\textlengthmark}{s}{a}{s}{m}{\textsci}{ð}{\textscy}{\textsubring{r}}} \dicPos{m}[9] \dicFlx{(‑s, ‑ir)}[8] \dicDirectTranslationCS{zámečník, zámečnice}
\dicEntry[lásbogi] \dicTerm{lás··bog|i} \dicIPA{{l}{au}{\textlengthmark}{s}{\textsubring{b}}{\textopeno i}{j}{\textsci}} \dicPos{m}[1] \dicFlx{(‑a, ‑ar)}[1] \dicDirectTranslationCS{kuše, samostříl}
\dicEntry[lásum] \dicTerm{lásum} \dicIPA{{l}{au}{\textlengthmark}{s}{\textscy}{\textsubring{m}}} \dicPos{v} \dicFlx{ind pf pl 1 pers} \dicLink{lesa}
\dicEntry[lát] \dicTerm{lát} \dicIPA{{l}{au}{\textlengthmark}{\textsubring{d}}} \dicPos{n}[2] \dicFlx{(‑s, ‑)}[5] \textbf{1.} \dicSynonym{andlát} \dicDirectTranslationCS{skon, úmrtí} \dicExampleIS{auglýsa lát e‑rs} \dicExampleCS{oznámit (čí) úmrtí}  \textbf{2.} \dicSynonym{hlé} \dicDirectTranslationCS{přestávka, chvíle oddechu} \dicExampleIS{Það var ekkert lát á storminum alla nóttina.} \dicExampleCS{Bouře zuřila bez přestávky celou noc.}  \textbf{3.} \dicSynonym*{undanlát} \dicDirectTranslationCS{poddajnost} \dicExampleIS{Það er ekkert lát á honum.} \dicExampleCS{Není v~něm kousek poddajnosti.}
\dicEntry[láta] \dicTerm{láta} \dicsymFrequent\  \dicIPA{{l}{au}{\textlengthmark}{\textsubring{d}}{a}} \dicPos{v}[6] \dicFlx{(læt, lét, létum, léti, látið)}[84] \dicFlx{acc} \textbf{1.} \dicSynonym*{leyfa að e‑að gerist} \dicDirectTranslationCS{(po)nechat, (po)nechávat} \dicExampleIS{láta e‑n lesa} \dicExampleCS{nechat (koho) číst};  \dicPhraseIS{láta e‑n í friði} \dicDirectTranslationCS{nechat (koho) na pokoji};  \dicPhraseIS{láta e‑n vera} \dicDirectTranslationCS{nechat (koho) být}  \textbf{2.} \dicDirectTranslationCS{(po)nechat, (po)nechávat, zanechat, zanechávat} \dicExampleIS{láta bókina þar} \dicExampleCS{nechat tam knížku};  \dicPhraseIS{láta e‑ð af hendi} \dicDirectTranslationCS{dát (co) z~ruky};  \dicPhraseIS{láta e‑m e‑ð í té} \dicDirectTranslationCS{poskytnout (komu co), předat (komu co)}  \textbf{3.} \dicSynonym*{hegða sér} \dicDirectTranslationCS{chovat se, jednat} \dicExampleIS{láta eins og fífl} \dicExampleCS{chovat se jako hlupák};  \dicPhraseIS{láta sem ekkert sé} \dicDirectTranslationCS{dělat jakoby nic; dělat, jako by se nic nestalo};  \dicPhraseIS{láta öllum illum látum} \dicDirectTranslationCS{dělat problémy\,/\addthin potíže};  \dicPhraseIS{láttu ekki svona} \dicFlx{imper} \dicDirectTranslationCS{nech toho}  \textbf{4.} \dicPhraseIS{láta lífið} \dicDirectTranslationCS{přijít o~život}  \textbf{5.} \dicPhraseIS{láta í haf} \dicDirectTranslationCS{vyplout na moře};  \dicIdiom{láta}[að]{ \dicPhraseIS{láta vel að e‑n}} \dicDirectTranslationCS{být milý\,/\addthin laskavý ke (komu)};  \dicIdiom{láta}[af]{ \dicPhraseIS{láta af e‑u}} \dicDirectTranslationCS{(za)nechat (čeho) (práce ap.)};  \dicIdiom{láta}[aftur]{ \dicPhraseIS{láta aftur e‑ð}} \dicDirectTranslationCS{zavřít\,/\addthin zavírat (co)} \dicExampleIS{láta aftur hurðina\,/\addthin dyrnar} \dicExampleCS{zavřít dveře};  \dicIdiom{láta}[eftir]{ \dicPhraseIS{láta allt eftir sér}} \dicDirectTranslationCS{dopřávat si, nic si neodepřít}; { \dicPhraseIS{láta e‑ð eftir e‑m}} \dicDirectTranslationCS{nechat (co) na (kom), neodepřít (komu co)}; { \dicPhraseIS{láta eftir}} \dicDirectTranslationCS{podřídit se, podvolit se} \dicExampleIS{Hann varð að láta eftir.} \dicExampleCS{Musel se podřídit.};  \dicIdiom{láta}[nærri]{ \dicPhraseIS{það lætur nærri}} \dicFlx{impers} \dicDirectTranslationCS{moc tomu nechybí, téměř};  \dicIdiom{láta}[sér]{ \dicPhraseIS{láta sér e‑ð vel líka}} \dicDirectTranslationCS{být s~(čím) spokojený}; { \dicPhraseIS{láta sér ekki bregða}} \dicDirectTranslationCS{nenechat se překvapit\,/\addthin zaskočit}; { \dicPhraseIS{láta sér fátt um finnast}} \dicDirectTranslationCS{nedat najevo příliš velký zájem};  \dicIdiom{láta}[sig]{ \dicPhraseIS{láta sig ekki}} \dicDirectTranslationCS{nevzdat se, nevzdávat se};  \dicIdiom{láta}[undan]{ \dicPhraseIS{láta undan e‑m}} \dicSynonym{beygja\smash{\textsuperscript{2}}} \dicDirectTranslationCS{podvolit se (komu), podrobit se (komu)}; { \dicPhraseIS{láta undan}} \dicDirectTranslationCS{povolit, spadnout (střecha ap.)};  \dicIdiom{láta}[uppi]{ \dicPhraseIS{láta e‑ð uppi}} \dicDirectTranslationCS{vyjevit (co), dát (co) najevo};  \dicIdiom{láta}[út úr]{ \dicPhraseIS{láta e‑ð út úr sér}} \dicDirectTranslationCS{(dovolit si) říct (co)};  \dicIdiom{láta}[yfir]{ \dicPhraseIS{láta illa yfir e‑u}} \dicDirectTranslationCS{neschvalovat (co), odsuzovat (co)} \dicExampleIS{láta illa yfir ástandinu} \dicExampleCS{neschvalovat takovýto stav};  \dicIdiom{látast}{ \dicPhraseIS{látast}} \dicFlx{refl} {\textbf{a.}} \dicDirectTranslationCS{předstírat} \dicExampleIS{látast vera reiður} \dicExampleCS{předstírat rozčilení};  {\textbf{b.}} \dicDirectTranslationCS{skonat, zemřít} \dicExampleIS{látast úr elli} \dicExampleCS{zemřít v~důsledku vysokého věku};  \dicIdiom{láta}{ \dicPhraseIS{láta e‑ð eiga sig}} \dicDirectTranslationCS{odpustit si (co)}; { \dicPhraseIS{láta e‑ð ekki á sig fá}} \dicDirectTranslationCS{nepoddat se (čemu)}; { \dicPhraseIS{láta e‑ð í ljós}} \dicDirectTranslationCS{dát (co) najevo, vyjádřit (co)} \dicExampleIS{láta skoðun sína í ljós} \dicExampleCS{dát najevo své názory}; { \dicPhraseIS{láta fara vel um sig}} \dicDirectTranslationCS{udělat si pohodlí} \dicExampleIS{Afi lætur fara vel um sig í hægindastólnum.} \dicExampleCS{Děda si dělá pohodlí v~křesle.}; { \dicPhraseIS{láta ekki hugfallast}} \dicDirectTranslationCS{neztratit kuráž}; { \dicPhraseIS{láta ekki e‑ð hjá líða}} \dicDirectTranslationCS{nenechat si ujít (co)}; { \dicPhraseIS{láta til skarar skríða}} \dicLangCat{přen.} \dicDirectTranslationCS{jít na věc}
\dicEntry[látalæti] \dicTerm{láta··læti} \dicIPA{{l}{au}{\textlengthmark}{\textsubring{d}}{a}{l}{a}{i}{\textsubring{d}}{\textsci}} \dicPos{n}[2] \dicFlx{pl}[19] \dicSynonym{uppgerð} \dicDirectTranslationCS{předstírání, přetvářka, pokrytectví}
\dicEntry[látbragð] \dicTerm{lát··bragð} \dicIPA{{l}{au}{\textlengthmark}{\textsubring{d}}{\textsubring{b}}{r}{a}{\textbabygamma}{\texttheta}} \dicPos{n}[2] \dicFlx{(‑s)}[2] \textbf{1.} \dicDirectTranslationCS{posunek, gesto}  \textbf{2.} \dicDirectTranslationCS{pantomima}
\dicEntry[látbragðsleikur] \dicTerm{lát·bragðs··leik|ur} \dicIPA{{l}{au}{\textlengthmark}{\textsubring{d}}{\textsubring{b}}{r}{a}{\textbabygamma}{ð}{s}{l}{ei}{\r{g}}{\textscy}{\textsubring{r}}} \dicPos{m}[6] \dicFlx{(‑s)}[17] \dicDirectTranslationCS{němohra, pantomima}
\dicEntry[látið] \dicTerm{látið} \dicIPA{{l}{au}{\textlengthmark}{\textsubring{d}}{\textsci}{\texttheta}} \dicPos{v} \dicFlx{supin} \dicLink{láta}
\dicEntry[látinn] \dicTerm{látinn} \dicsymFrequent\  \dicIPA{{l}{au}{\textlengthmark}{\textsubring{d}}{\textsci}{\textsubring{n}}} \dicPos{adj}[6]\dicFlx{}[-6] \dicSynonym{dáinn} \dicDirectTranslationCS{zesnulý, zemřelý} \dicExampleIS{níu látnir og fjölmargir slasaðir} \dicExampleCS{devět mrtvých a~mnoho zraněných};  \dicPhraseIS{vera vant við látinn} \dicDirectTranslationCS{mít napilno, být zaneprázdněný};  \dicPhraseIS{vera vel látinn} \dicDirectTranslationCS{být populární\,/\addthin uznávaný}
\dicEntry[látlaus] \dicTerm{lát··laus} \dicIPA{{l}{au}{\textlengthmark}{\textsubring{d}}{l}{\oe i}{s}} \dicPos{adj}[5]\dicFlx{}[-1] \textbf{1.} \dicSynonym{tilgerðarlaus} \dicDirectTranslationCS{prostý, skromný, nenáročný}  \textbf{2.} \dicSynonym{linnulaus} \dicDirectTranslationCS{vytrvalý, soustavný, nepřetržitý} \dicExampleIS{látlaus rigning} \dicExampleCS{soustavný déšť}
\dicEntry[látleysi] \dicTerm{lát··leysi} \dicIPA{{l}{au}{\textlengthmark}{\textsubring{d}}{l}{ei}{s}{\textsci}} \dicPos{n}[2] \dicFlx{(‑s)}[20] \dicSynonym*{tilgerðarleysi} \dicDirectTranslationCS{prostota, přirozenost, jednoduchost}
\dicEntry[látún] \dicTerm{látún} \dicIPA{{l}{au}{\textlengthmark}{t\smash{\textsuperscript{h}}}{u}{\textsubring{n}}} \dicPos{n}[2] \dicFlx{(‑s)}[2] \dicDirectTranslationCS{mosaz}
\dicEntry[látæði] \dicTerm{lát··æði} \dicIPA{{l}{au}{\textlengthmark}{\textsubring{d}}{a}{i}{ð}{\textsci}} \dicPos{n}[2] \dicFlx{(‑s)}[20] \textbf{1.} \dicSynonym{hegðun} \dicDirectTranslationCS{chování, jednání}  \textbf{2.} \dicSynonym{látbragð} \dicDirectTranslationCS{posunek, gesto}
\dicEntry[lávarður] \dicTerm{lávarð|ur} \dicIPA{{l}{au}{\textlengthmark}{v}{a}{r}{ð}{\textscy}{\textsubring{r}}} \dicPos{m}[6] \dicFlx{(‑s\,/\addthin ‑ar, ‑ar)}[64] \dicDirectTranslationCS{lord}
\dicEntry[Lbs.] \dicTerm{Lbs.} \dicPos{zkr} \dicPhraseIS{Landsbókasafn} \dicDirectTranslationCS{Národní knihovna}
\dicEntry[ld.] \dicTerm{ld.} \dicPos{zkr} \dicPhraseIS{laugardagur} \dicDirectTranslationCS{sobota}
\dicEntry[leðja] \dicTerm{leðj|a} \dicIPA{{l}{\textepsilon}{ð}{j}{a}} \dicPos{f}[1] \dicFlx{(‑u)}[5] \dicSynonym{for\smash{\textsuperscript{1}}} \dicDirectTranslationCS{bláto, bahno} \dicExampleIS{sökkva ofan í leðjuna} \dicExampleCS{zapadnout do bláta}
\dicEntry[leður] \dicTerm{leður} \dicsymFrequent\  \dicIPA{{l}{\textepsilon}{\textlengthmark}{ð}{\textscy}{\textsubring{r}}} \dicPos{n}[2] \dicFlx{(‑s, ‑)}[25] \dicDirectTranslationCS{kůže} \dicIndirectTranslationCS{(materiál)} \dicExampleIS{vörur úr leðri} \dicExampleCS{výrobky z~kůže}
\dicEntry[leðurblaka] \dicTerm{leður··|blaka} \dicIPA{{l}{\textepsilon}{\textlengthmark}{ð}{\textscy}{r}{\textsubring{b}}{l}{a}{\r{g}}{a}} \dicPos{f}[1] \dicFlx{(‑blöku, ‑blökur)}[20] \textbf{1.} \dicPhraseIS{leðurblökur} \dicFlx{pl} \dicFieldCat{zool.} \dicDirectTranslationCS{letouni} \textit{(l.~{\textLA{Chiroptera}})}  \dicsymPhoto\   \textbf{2.} \dicFieldCat{zool.} \dicDirectTranslationCS{netopýr} \textit{(l.~{\textLA{Microchiroptera}})}
\dicFigure{62444.jpg}{Leðurblaka}{Leðurblaka - Krejčík Stanislav, Biolib, Copyright/CC-BY}
\dicEntry[leðurjakki] \dicTerm{leður··jakk|i} \dicIPA{{l}{\textepsilon}{\textlengthmark}{ð}{\textscy}{r}{j}{a}{h}{\r{\textObardotlessj}}{\textsci}} \dicPos{m}[1] \dicFlx{(‑a, ‑ar)}[8] \dicDirectTranslationCS{kožená bunda}
\dicEntry[leðurlíki] \dicTerm{leður··líki} \dicIPA{{l}{\textepsilon}{\textlengthmark}{ð}{\textscy}{r}{l}{i}{\r{\textObardotlessj}}{\textsci}} \dicPos{n}[2] \dicFlx{(‑s, ‑)}[16] \dicDirectTranslationCS{koženka}
\dicEntry[leg] \dicTerm{leg} \dicIPA{{l}{\textepsilon}{\textlengthmark}{x}} \dicPos{n}[2] \dicFlx{(‑s, ‑)}[5] \textbf{1.} \dicDirectTranslationCS{umístění, položení}  \textbf{2.} \dicFieldCat{anat.} \dicDirectTranslationCS{děloha}
\dicEntry[lega] \dicTerm{leg|a} \dicsymFrequent\  \dicIPA{{l}{\textepsilon}{\textlengthmark}{\textbabygamma}{a}} \dicPos{f}[1] \dicFlx{(‑u, ‑ur)}[19] \textbf{1.} \dicSynonym{staðsetning} \dicDirectTranslationCS{uložení, pozice} \dicExampleIS{landfræðileg lega Íslands} \dicExampleCS{geografická pozice Islandu}  \textbf{2.} \dicSynonym{rúmlega\smash{\textsuperscript{1}}} \dicDirectTranslationCS{ležení (na lůžku)}  \textbf{3.} \dicFieldCat{techn.} \dicDirectTranslationCS{ložisko}
\dicEntry[leggja] \dicTerm{leggja} \dicsymFrequent\  \dicIPA{{l}{\textepsilon}{\r{\textObardotlessj}}{\textlengthmark}{a}} \dicPos{v}[4] \dicFlx{(legg, lagði, lögðum, legði, lagt)}[20] \dicFlx{acc\,/\addthin dat} \textbf{1.} \dicFlx{acc} \dicDirectTranslationCS{položit, pokládat, umístit, umísťovat, uložit, ukládat} \dicExampleIS{leggja hlutinn á borðið} \dicExampleCS{položit věc na stůl};  \dicPhraseIS{leggja sig} \dicDirectTranslationCS{položit se, lehnout si, natáhnout se}  \textbf{2.} \dicFlx{acc} \dicDirectTranslationCS{(na)instalovat (potrubí ap.), (na)táhnout (elek. vedení ap.), položit, pokládat (dlaždice ap.)} \dicExampleIS{leggja nýjan veg} \dicExampleCS{položit novou silnici}  \textbf{3.} \dicDirectTranslationCS{vyjít, vyjet, vyrazit};  \dicPhraseIS{leggja af\,/\addthin á stað} \dicDirectTranslationCS{vyrazit (na cestu ap.), vydat se} \dicExampleIS{Hún leggur af stað til Akureyrar.} \dicExampleCS{Vyráží do Akureyri.}  \textbf{4.} \dicFlx{dat} \dicDirectTranslationCS{(za)parkovat} \dicExampleIS{leggja bílnum hjá búðinni} \dicExampleCS{zaparkovat auto u~obchodu}  \textbf{5.} \dicFlx{dat} \dicDirectTranslationCS{vyřadit, vyřazovat} \dicExampleIS{leggja gömlu kaffivélinni} \dicExampleCS{vyřadit starý kávovar}  \textbf{6.} \dicFlx{acc} \dicDirectTranslationCS{porazit, porážet (protivníka ap.)}  \textbf{7.} \dicFlx{acc} \dicDirectTranslationCS{(za)útočit, napadnout, napadat}  \textbf{8.} \dicPhraseIS{e‑ð leggur} \dicFlx{impers} {\textbf{a.}} \dicDirectTranslationCS{(co) se šíří, (co) se rozšiřuje (vůně ap.)};  {\textbf{b.}} \dicDirectTranslationCS{(co) zamrzá};  \dicIdiom{leggja}[að]{ \dicPhraseIS{leggja hart að sér}} \dicDirectTranslationCS{tvrdě pracovat, vyvinout veškeré úsilí}; { \dicPhraseIS{leggja fast að e‑m}} \dicDirectTranslationCS{vyvíjet tlak na (koho)}; { \dicPhraseIS{leggja e‑ð að veði}} \dicDirectTranslationCS{dát (co) v~sázku, vsadit (co)};  \dicIdiom{leggja}[af]{ \dicPhraseIS{leggja af}} \dicDirectTranslationCS{(z)hubnout} \dicExampleIS{Hún lagði af vegna þess að hún var veik.} \dicExampleCS{Zhubla, protože byla nemocná.};  \dicIdiom{leggja}[aftur]{ \dicPhraseIS{leggja aftur e‑ð}} \dicSynonym{loka\smash{\textsuperscript{2}}} \dicDirectTranslationCS{zavřít (co), zamknout (co)} \dicExampleIS{leggja aftur augun} \dicExampleCS{zavřít oči};  \dicIdiom{leggja}[á]{ \dicPhraseIS{leggja á borð}} \dicDirectTranslationCS{prostřít (stůl)}; { \dicPhraseIS{leggja ást á e‑n}} \dicDirectTranslationCS{zamilovat se do (koho), milovat (koho)}; { \dicPhraseIS{leggja áherslu á e‑ð}} \dicDirectTranslationCS{zdůraznit\,/\addthin zdůrazňovat (co), klást na (co) důraz}; { \dicPhraseIS{leggja mikið á sig}} \dicSynonym*{beita sig hörku} \dicDirectTranslationCS{opřít se do toho, vynasnažit se}; { \dicPhraseIS{leggja á}} \dicDirectTranslationCS{zavěsit, položit (telefon)} \dicExampleIS{Hann lagði símann á.} \dicExampleCS{Položil telefon.}; { \dicPhraseIS{leggja e‑ð á minnið, leggja e‑ð í minni}} \dicDirectTranslationCS{uložit (co) do paměti, zapamatovat si (co)}; { \dicPhraseIS{leggja skatt á e‑ð}} \dicDirectTranslationCS{(z)danit (co), uvalit daň na (co)}; { \dicPhraseIS{leggja hatur á e‑n}} \dicDirectTranslationCS{nenávidět (koho)}; { \dicPhraseIS{leggja dul á e‑ð}} \dicDirectTranslationCS{(u)tajit (co)}; { \dicPhraseIS{leggja stund á e‑ð}} \dicDirectTranslationCS{věnovat se (čemu)};  \dicIdiom{leggja}[fram]{ \dicPhraseIS{leggja sig allan fram}} \dicDirectTranslationCS{vynasnažit se, vyvinout úsilí}; { \dicPhraseIS{leggja fram e‑ð}} \dicDirectTranslationCS{předložit\,/\addthin předkládat (co) (k~diskuzi ap.)};  \dicIdiom{leggja}[fyrir]{ \dicPhraseIS{leggja e‑ð fyrir e‑n}} \dicDirectTranslationCS{předložit\,/\addthin předkládat (komu co) (otázku ap.)}; { \dicPhraseIS{leggja fyrir sig e‑ð}} \dicDirectTranslationCS{zabývat se (čím), být zapojen do (čeho)}; { \dicPhraseIS{leggja fyrir}} \dicDirectTranslationCS{odkládat, šetřit};  \dicIdiom{leggja}[inn]{ \dicPhraseIS{leggja e‑ð inn}} \dicDirectTranslationCS{uložit\,/\addthin ukládat (co)} \dicExampleIS{Ég legg peninga inn í bankann.} \dicExampleCS{Ukládám si peníze do banky.};  \dicIdiom{leggja}[í]{ \dicPhraseIS{leggja í e‑ð}} \dicDirectTranslationCS{odvážit se na (co)}; { \dicPhraseIS{leggja í hann}} \dicDirectTranslationCS{vydat se na cestu, vyrazit};  \dicIdiom{leggja}[niður]{ \dicPhraseIS{leggja niður e‑ð}} \dicDirectTranslationCS{skončit s~(čím), skoncovat s~(čím)} \dicExampleIS{leggja niður starf} \dicExampleCS{skončit s~prací}; { \dicPhraseIS{leggja niður vopn}} \dicDirectTranslationCS{složit zbraně}; \textls[15]{{ \dicPhraseIS{leggja sig niður við e‑ð}} \dicDirectTranslationCS{sní\nobreak žit se k~(čemu)};}  \dicIdiom{leggja}[saman]{ \dicPhraseIS{leggja saman e‑ð}} \dicFieldCat{mat.} \dicDirectTranslationCS{sčítat (co), přičítat (co)} \dicExampleIS{Barnið leggur saman tvo og þrjá og fær fimm.} \dicExampleCS{Dítě sečte dva plus tři a~vyjde mu pět.}; { \dicPhraseIS{leggja saman tvo og tvo}} \dicLangCat{přen.} \dicDirectTranslationCS{dát si dvě a~dvě dohromady};  \dicIdiom{leggja}[til]{ \dicPhraseIS{leggja e‑m til e‑ð}} \dicDirectTranslationCS{poskytnout\,/\addthin poskytovat (komu co)}; { \dicPhraseIS{leggja e‑ð til}} \dicDirectTranslationCS{navrhnout\,/\addthin navrhovat (co), předložit návrh o~(čem)} \dicExampleIS{Ég legg til að þessi plaköt verði lagfærð.} \dicExampleCS{Navrhuji, aby tyto plakáty byly upraveny.};  \dicPhraseIS{leggja sér e‑ð til munns} \dicDirectTranslationCS{jíst (co), pojídat (co)};  \dicIdiom{leggja}[undir]{ \dicPhraseIS{leggja undir sig e‑ð}} \dicDirectTranslationCS{podmanit\,/\addthin podmaňovat si (co), podrobit\,/\addthin podrobovat si (co)};  \dicIdiom{leggja}[upp]{ \dicPhraseIS{leggja upp}} \dicDirectTranslationCS{vyrazit, vydat se (na cestu ap.)} \dicExampleIS{leggja upp í ferðina} \dicExampleCS{vydat se na cestu};  \dicIdiom{leggja}[upp úr]{ \dicPhraseIS{leggja mikið\,/\addthin lítið upp úr e‑u}} \dicDirectTranslationCS{pokládat (co) za hodně\,/\addthin málo důležité};  \dicIdiom{leggja}[út]{ \dicPhraseIS{leggja út e‑ð}} {\textbf{a.}} \dicSynonym{borga} \dicDirectTranslationCS{vyplatit (co), přispět (čím) (částkou ap.)};  {\textbf{b.}} \dicSynonym{túlka} \dicDirectTranslationCS{přeložit (co), tlumočit (co)};  \dicIdiom{leggja}[við]{ \dicPhraseIS{leggja rækt við e‑ð}} \dicDirectTranslationCS{kultivovat (co), rozvíjet (co)}; { \dicPhraseIS{leggja e‑ð við e‑ð}} \dicDirectTranslationCS{přidat (co) k~(čemu)} \dicExampleIS{leggja tvo við þrjá} \dicExampleCS{přidat dvě ke třem};  \dicIdiom{leggjast}{ \dicPhraseIS{leggjast}} \dicFlx{refl} \dicDirectTranslationCS{lehnout si, natáhnout se};  \dicIdiom{leggjast}[að]{ \dicPhraseIS{leggjast að bryggju}} \dicFlx{refl} \dicDirectTranslationCS{přistát u~mola};  \dicIdiom{leggjast}[af]{ \dicPhraseIS{leggjast af}} \dicFlx{refl} \dicDirectTranslationCS{vymizet, zaniknout};  \dicIdiom{leggjast}[á]{ \dicPhraseIS{leggjast á eitt}} \dicFlx{refl} \dicDirectTranslationCS{spojit síly, táhnout za jeden provaz}; { \dicPhraseIS{e‑að leggst þungt á e‑n}} \dicFlx{refl} \dicLangCat{přen.} \dicDirectTranslationCS{(co koho) tíží (obavy ap.)};  \dicIdiom{leggjast}[gegn]{ \dicPhraseIS{leggjast gegn e‑u}} \dicFlx{refl} \dicDirectTranslationCS{postavit\,/\addthin stavět se proti (čemu)};  \dicIdiom{leggjast}[í]{ \dicPhraseIS{leggjast í drykkju}} \dicFlx{refl} \dicDirectTranslationCS{oddávat se pití (alkoholu)}; { \dicPhraseIS{leggjast í þunglyndi}} \dicFlx{refl} \dicDirectTranslationCS{upadnout do deprese};  \dicIdiom{leggjast}[með]{ \dicPhraseIS{leggjast með e‑m}} \dicFlx{refl} \dicDirectTranslationCS{vyspat se s~(kým), spát s~(kým)};  \dicIdiom{leggjast}[niður]{ \dicPhraseIS{leggjast niður}} \dicFlx{refl} \dicDirectTranslationCS{být přerušen, přestat};  \dicIdiom{leggjast}[yfir]{ \dicPhraseIS{e‑að leggst yfir e‑ð}} \dicFlx{refl} \dicDirectTranslationCS{(co) se usazuje nad (čím) (mlha nad městem ap.)}; { \dicPhraseIS{leggjast yfir e‑ð}} \dicFlx{refl} \dicDirectTranslationCS{pročíst (co), (důkladně) přečíst (co)}
\dicEntry[leggur] \dicTerm{legg|ur} \dicsymFrequent\  \dicIPA{{l}{\textepsilon}{\r{g}}{\textlengthmark}{\textscy}{\textsubring{r}}} \dicPos{m}[9] \dicFlx{(‑jar\,/\addthin ‑s, ‑ir)}[26] \textbf{1.} \dicSynonym{bein} \dicDirectTranslationCS{kost}  \textbf{2.} \dicSynonym{fótleggur} \dicDirectTranslationCS{noha, dolní končetina} \dicExampleIS{Ég er tilfinningalaus í leggnum.} \dicExampleCS{Nemám v~noze žádný cit.}  \textbf{3.} \dicFieldCat{bot.} \dicSynonym{stilkur} \dicDirectTranslationCS{stopka, stonek};  \dicIdiom{leggur}{ \dicPhraseIS{komast á legg}} \dicFlx{refl} \dicLangCat{přen.} \dicDirectTranslationCS{postavit se na vlastní nohy}
\dicEntry[leggöng] \dicTerm{leg··göng} \dicIPA{{l}{\textepsilon}{\textbabygamma}{\r{g}}{\oe i}{\ng}{\r{g}}} \dicPos{n}[2] \dicFlx{pl}[9] \dicFieldCat{anat.} \dicDirectTranslationCS{vagina, pochva}
\dicEntry[legið] \dicTerm{legið} \dicIPA{{l}{ei}{\textsci}{\texttheta}} \dicPos{v} \dicFlx{supin} \dicLink{liggja}
\dicEntry[legir] \dicTerm{legir} \dicIPA{{l}{ei}{\textsci}{\textsubring{r}}} \dicPos{m} \dicFlx{pl nom} \dicLink{lögur}
\dicEntry[legkaka] \dicTerm{leg··|kaka} \dicIPA{{l}{\textepsilon}{\textbabygamma}{k\smash{\textsuperscript{h}}}{a}{\r{g}}{a}} \dicPos{f}[1] \dicFlx{(‑köku, ‑kökur)}[20] \dicFieldCat{anat.} \dicDirectTranslationCS{placenta, plodový obal}
\dicEntry[legstaður] \dicTerm{leg··stað|ur} \dicIPA{{l}{\textepsilon}{x}{s}{\textsubring{d}}{a}{ð}{\textscy}{\textsubring{r}}} \dicPos{m}[10] \dicFlx{(‑ar, ‑ir)}[14] \dicSynonym{gröf} \dicDirectTranslationCS{hrob(ka)}
\dicEntry[legsteinn] \dicTerm{leg··stein|n} \dicIPA{{l}{\textepsilon}{x}{s}{\textsubring{d}}{ei}{\textsubring{d}}{\textsubring{n}}} \dicPos{m}[6] \dicFlx{(‑s, ‑ar)}[42] \dicDirectTranslationCS{náhrobní kámen}
\dicEntry[legubekkur] \dicTerm{legu··bekk|ur} \dicIPA{{l}{\textepsilon}{\textlengthmark}{\textbabygamma}{\textscy}{\textsubring{b}}{\textepsilon}{h}{\r{g}}{\textscy}{\textsubring{r}}} \dicPos{m}[9] \dicFlx{(‑s\,/\addthin ‑jar, ‑ir)}[26] \dicDirectTranslationCS{pohovka, lehátko}
\dicEntry[legudagur] \dicTerm{legu··dag|ur} \dicIPA{{l}{\textepsilon}{\textlengthmark}{\textbabygamma}{\textscy}{\textsubring{d}}{a}{\textbabygamma}{\textscy}{\textsubring{r}}} \dicPos{m}[6] \dicFlx{(‑s, ‑ar)}[62] \textbf{1.} \dicDirectTranslationCS{den na lůžku (v~nemocnici ap.)}  \textbf{2.} \dicDirectTranslationCS{den (lodi ap.) v~přístavu}
\dicEntry[legugjald] \dicTerm{legu··|gjald} \dicIPA{{l}{\textepsilon}{\textlengthmark}{\textbabygamma}{\textscy}{\r{\textObardotlessj}}{a}{l}{\textsubring{d}}} \dicPos{n}[2] \dicFlx{(‑gjalds, ‑gjöld)}[8] \dicFieldCat{nám.} \dicDirectTranslationCS{kotevné} \dicIndirectTranslationCS{(poplatek za kotvení)}
\dicEntry[legusár] \dicTerm{legu··sár} \dicIPA{{l}{\textepsilon}{\textlengthmark}{\textbabygamma}{\textscy}{s}{au}{\textsubring{r}}} \dicPos{n}[2] \dicFlx{(‑s, ‑)}[5] \dicFieldCat{med.} \dicDirectTranslationCS{proleženina}
\dicEntry[leið] \dicTerm{leið\smash{\textsuperscript{1}}} \dicsymFrequent\  \dicIPA{{l}{ei}{\textlengthmark}{\texttheta}} \dicPos{f}[7] \dicFlx{(‑ar, ‑ir)}[1] \textbf{1.} \dicSynonym{vegur} \dicDirectTranslationCS{cesta, stezka} \dicExampleIS{bein leið} \dicExampleCS{přímá cesta};  \dicPhraseIS{eiga leið hjá} \dicDirectTranslationCS{mít cestu kolem};  \dicPhraseIS{vera á leið(inni) (þangað)} \dicDirectTranslationCS{být na cestě (tam)}  \textbf{2.} \dicSynonym{háttur\smash{\textsuperscript{1}}} \dicDirectTranslationCS{způsob, postup};  \dicPhraseIS{það er engin leið að \dots{}} \dicDirectTranslationCS{není žádný způsob, jak\dots{}}  \textbf{3.} \dicSynonym{vegalengd} \dicDirectTranslationCS{cesta, vzdálenost} \dicExampleIS{dagleið} \dicExampleCS{jednodenní cesta}  \textbf{4.} \dicSynonym{áætlunarleið} \dicDirectTranslationCS{(dopravní) linka};  \dicIdiom{leið}{ \dicPhraseIS{á miðri leið}} \dicFlx{adv} \dicDirectTranslationCS{v~půli cesty}; { \dicPhraseIS{heim á leið}} \dicFlx{adv} \dicDirectTranslationCS{domů}; { \dicPhraseIS{leggja leið sína, leggja leiðir sínar}} \dicDirectTranslationCS{vydat se (do cizí země ap.)} \dicExampleIS{leggja leið sína til Íslands} \dicExampleCS{vydat se na Island};  \dicPhraseIS{um leið} \dicFlx{adv} \dicDirectTranslationCS{zároveň, současně, naráz}; { \dicPhraseIS{um leið og}} \dicFlx{conj} \dicSynonym*{strax og} \dicDirectTranslationCS{jakmile, hned jak} \dicExampleIS{Hann fór að sofa um leið og hann kom heim.} \dicExampleCS{Jakmile přišel domů, šel si lehnout.}; { \dicPhraseIS{vera komin (fimm) mánuði á leið}} \dicDirectTranslationCS{být v~(pátém) měsíci (těhotenství)}
\dicEntry[leið] \dicTerm{leið\smash{\textsuperscript{2}}} \dicIPA{{l}{ei}{\textlengthmark}{\texttheta}} \dicPos{v} \dicFlx{ind pf sg 1 pers} \dicLink{líða}
\dicEntry[leiða] \dicTerm{lei|ða} \dicsymFrequent\  \dicIPA{{l}{ei}{\textlengthmark}{ð}{a}} \dicPos{v}[2] \dicFlx{(‑ddi, ‑tt)}[169] \dicFlx{acc} \textbf{1.} \dicDirectTranslationCS{(při)vést, přivádět, vodit} \dicExampleIS{leiða e‑n þangað} \dicExampleCS{přivést (koho) tam}  \textbf{2.} \dicDirectTranslationCS{(při)vést, přivádět (elektřinu ap.)} \dicExampleIS{leiða rafmagn í hús} \dicExampleCS{přivádět elektřinu do domu};  \dicIdiom{leiða}[að]{ \dicPhraseIS{leiða hugann að e‑u}} \dicDirectTranslationCS{zavést myšlenku na (co)}; { \dicPhraseIS{leiða tal að e‑u}} \dicDirectTranslationCS{zavést řeč na (co), nadhodit (co)};  \dicIdiom{leiða}[af]{ \dicPhraseIS{leiða e‑ð af sér}} \dicDirectTranslationCS{mít (co) za následek, vést k~(čemu), skončit se (čím)};  \dicIdiom{leiða}[fyrir]{ \dicPhraseIS{leiða e‑ð e‑m fyrir sjónir}} \dicDirectTranslationCS{ukázat (komu co), vysvětlit (komu co)};  \dicIdiom{leiða}[hjá]{ \dicPhraseIS{leiða e‑ð hjá sér}} \dicDirectTranslationCS{nevšímat si (čeho), nebrat na vědomí (co)} \dicExampleIS{leiða orð hans hjá sér} \dicExampleCS{nevšímat si jeho slov};  \dicIdiom{leiða}[í]{ \dicPhraseIS{leiða e‑u í ljós}} \dicDirectTranslationCS{(po)odhalit (co)};  \dicIdiom{leiða}[til]{ \dicPhraseIS{leiða til e‑s}} \dicDirectTranslationCS{vést k~(čemu), vyústit v~(co)};  \dicIdiom{leiðast}{ \dicPhraseIS{leiðast}} \dicFlx{refl} \dicDirectTranslationCS{vést se za ruce} \dicExampleIS{Þau leiddust niður götuna.} \dicExampleCS{Šli po ulici a~vedli se za ruce.}; { \dicPhraseIS{leiðast út í drykkju}} \dicFlx{refl} \dicDirectTranslationCS{začít pít (alkohol)}; { \dicPhraseIS{láta til leiðast}} \dicDirectTranslationCS{nechat se přesvědčit};  \dicPhraseIS{e‑m leiðist} \dicFlx{refl impers} \dicDirectTranslationCS{(kdo) se nudí, (komu) je dlouhá chvíle} \dicExampleIS{Mér leiðist.} \dicExampleCS{Nudím se.}; { \dicPhraseIS{e‑m leiðist e‑að}} \dicFlx{refl impers} \dicDirectTranslationCS{(koho) nudí (co), (komu) přijde nudné (co)}
\dicEntry[leiðabók] \dicTerm{leiða··|bók} \dicIPA{{l}{ei}{\textlengthmark}{ð}{a}{\textsubring{b}}{ou}{\r{g}}} \dicPos{f}[8] \dicFlx{(‑bókar, ‑bækur)}[5] \dicDirectTranslationCS{jízdní řád}
\dicEntry[leiðangur] \dicTerm{leið··ang|ur} \dicsymFrequent\  \dicIPA{{l}{ei}{\textlengthmark}{ð}{au}{\ng}{\r{g}}{\textscy}{\textsubring{r}}} \dicPos{m}[5] \dicFlx{(‑urs, ‑rar)}[3] \dicSynonym{könnunarferð} \dicDirectTranslationCS{výprava, expedice, mise} \dicExampleIS{Leiðangurinn tekur tvær vikur.} \dicExampleCS{Expedice potrvá dva týdny.}
\dicEntry[leiðarendi] \dicTerm{leiðar··end|i} \dicIPA{{l}{ei}{\textlengthmark}{ð}{a}{r}{\textepsilon}{n}{\textsubring{d}}{\textsci}} \dicPos{m}[1] \dicFlx{(‑a, ‑ar)}[1] \dicSynonym{ákvörðunarstaður} \dicDirectTranslationCS{cíl cesty, destinace} \dicExampleIS{komast á leiðarenda} \dicExampleCS{dorazit do cíle cesty}
\dicEntry[leiðari] \dicTerm{leið··ar|i} \dicIPA{{l}{ei}{\textlengthmark}{ð}{a}{r}{\textsci}} \dicPos{m}[1] \dicFlx{(‑a, ‑ar)}[13] \textbf{1.} \dicSynonym{grein} \dicDirectTranslationCS{úvodník} \dicIndirectTranslationCS{(článek umístěný na první stránce novin)}  \textbf{2.} \dicFieldCat{fyz.} \dicDirectTranslationCS{vodič (elektřiny ap.)}
\dicEntry[leiðarljós] \dicTerm{leiðar··ljós} \dicIPA{{l}{ei}{\textlengthmark}{ð}{a}{r}{l}{j}{ou}{s}} \dicPos{n}[2] \dicFlx{(‑s, ‑)}[5] \dicLangCat{přen.} \dicDirectTranslationCS{záchytný bod, vodítko};  \dicPhraseIS{hafa e‑ð að leiðarljósi} \dicLangCat{přen.} \dicDirectTranslationCS{mít (co) jako záchytný bod}
\dicEntry[leiðarlok] \dicTerm{leiðar··lok} \dicIPA{{l}{ei}{\textlengthmark}{ð}{a}{r}{l}{\textopeno}{\r{g}}} \dicPos{n}[2] \dicFlx{pl}[1] \dicDirectTranslationCS{závěr\,/\addthin konec cesty (silnice i~putování)}
\dicEntry[leiðarlýsing] \dicTerm{leiðar··lýs·ing} \dicIPA{{l}{ei}{\textlengthmark}{ð}{a}{r}{l}{i}{s}{i}{\ng}{\r{g}}} \dicPos{f}[4] \dicFlx{(‑ar, ‑ar)}[5] \dicDirectTranslationCS{itinerář, cestovní plán}
\dicEntry[leiðarminni] \dicTerm{leiðar··minni} \dicIPA{{l}{ei}{\textlengthmark}{ð}{a}{r}{m}{\textsci}{n}{\textsci}} \dicPos{n}[2] \dicFlx{(‑s, ‑)}[14] \dicDirectTranslationCS{příznačný motiv, leitmotiv}
\dicEntry[leiðarvísir] \dicTerm{leiðar··vís|ir} \dicIPA{{l}{ei}{\textlengthmark}{ð}{a}{r}{v}{i}{s}{\textsci}{\textsubring{r}}} \dicPos{m}[7] \dicFlx{(‑is, ‑ar)}[1] \dicDirectTranslationCS{příručka, manuál, návod} \dicExampleIS{leiðarvísir um vél} \dicExampleCS{manuál k~zařízení}
\dicEntry[leiðbeina] \dicTerm{leið··bein|a} \dicIPA{{l}{ei}{ð}{\textsubring{b}}{ei}{n}{a}} \dicPos{v}[2] \dicFlx{(‑di, ‑t)}[152] \dicFlx{dat} \dicSynonym{kenna} \dicDirectTranslationCS{učit, vyučovat, cvičit, školit} \dicExampleIS{leiðbeina e‑m í e‑u} \dicExampleCS{školit (koho) v~(čem)}
\dicEntry[leiðbeinandi] \dicTerm{leið··bein·|andi} \dicIPA{{l}{ei}{ð}{\textsubring{b}}{ei}{n}{a}{n}{\textsubring{d}}{\textsci}} \dicPos{m}[2] \dicFlx{(‑anda, ‑endur)}[1] \textbf{1.} \dicDirectTranslationCS{školitel(ka), instruktor(ka), cvičitel(ka)}  \textbf{2.} \dicFieldCat{škol.} \dicDirectTranslationCS{nekvalifikovaný učitel, nekvalifikovaná učitelka}
\dicEntry[leiðbeining] \dicTerm{leið··bein·ing} \dicIPA{{l}{ei}{ð}{\textsubring{b}}{ei}{n}{i}{\ng}{\r{g}}} \dicPos{f}[4] \dicFlx{(‑ar, ‑ar)}[5] \textbf{1.} \dicSynonym{leiðsögn} \dicDirectTranslationCS{pokyn, instrukce}  \textbf{2.} \dicPhraseIS{leiðbeiningar} \dicFlx{pl} \dicSynonym{leiðarvísir} \dicDirectTranslationCS{pokyny, návod, manuál, instrukce} \dicExampleIS{leiðbeiningar um notkun} \dicExampleCS{návod k~použití}
\dicEntry[leiði] \dicTerm{leið|i\smash{\textsuperscript{1}}} \dicIPA{{l}{ei}{\textlengthmark}{ð}{\textsci}} \dicPos{m}[1] \dicFlx{(‑a)}[3] \dicSynonym{leiðindi} \dicDirectTranslationCS{nuda, dlouhá chvíle} \dicExampleIS{fá leiða á leikjum} \dicExampleCS{nudit se u~hry}
\dicEntry[leiði] \dicTerm{leiði\smash{\textsuperscript{2}}} \dicIPA{{l}{ei}{\textlengthmark}{ð}{\textsci}} \dicPos{n}[2] \dicFlx{(‑s, ‑)}[14] \textbf{1.} \dicSynonym{gröf} \dicDirectTranslationCS{náhrobek, hrobka} \dicExampleIS{setja kross á leiðið} \dicExampleCS{umístit kříž na náhrobek}  \textbf{2.} \dicSynonym{meðvindur} \dicDirectTranslationCS{příznivý vítr}  \textbf{3.} \dicSynonym{færð} \dicDirectTranslationCS{sjízdnost}
\dicEntry[leiðindamál] \dicTerm{leiðinda··mál} \dicIPA{{l}{ei}{\textlengthmark}{ð}{\textsci}{n}{\textsubring{d}}{a}{m}{au}{\textsubring{l}}} \dicPos{n}[2] \dicFlx{(‑s, ‑)}[5] \dicDirectTranslationCS{nepříjemnost, nepříjemná\,/\addthin otravná věc}
\dicEntry[leiðindi] \dicTerm{leið··indi} \dicIPA{{l}{ei}{\textlengthmark}{ð}{\textsci}{n}{\textsubring{d}}{\textsci}} \dicPos{n}[2] \dicFlx{pl}[19] \textbf{1.} \dicSynonym{leiði\smash{\textsuperscript{1}}} \dicDirectTranslationCS{nuda, nudnost, dlouhá chvíle}  \textbf{2.} \dicSynonym{óánægja} \dicDirectTranslationCS{protivnost, špatná nálada}
\dicEntry[leiðinlegur] \dicTerm{leiðin··legur} \dicsymFrequent\  \dicIPA{{l}{ei}{\textlengthmark}{ð}{\textsci}{n}{l}{\textepsilon}{\textbabygamma}{\textscy}{\textsubring{r}}} \dicPos{adj}[1]\dicFlx{}[-8] \textbf{1.} \dicDirectTranslationCS{nudný, fádní} \dicExampleIS{leiðinlegur kennari} \dicExampleCS{nudný učitel}  \textbf{2.} \dicSynonym{dapurlegur} \dicDirectTranslationCS{smutný, tristní}  \textbf{3.} \dicDirectTranslationCS{protivný, nepříjemný (počasí ap.)}
\dicEntry[leiðitamur] \textls[15]{\dicTerm{leiði··|tamur} \dicIPA{{l}{ei}{\textlengthmark}{ð}{\textsci}{t\smash{\textsuperscript{h}}}{a}{m}{\textscy}{\textsubring{r}}} \dicPos{adj}[1] \dicFlx{(f ‑töm)}[2] \dicSynonym{eftirlátur}} \dicDirectTranslationCS{svolný, povolný}
\dicEntry[leiðni] \dicTerm{leiðn|i} \dicIPA{{l}{ei}{ð}{n}{\textsci}} \dicPos{f}[3] \dicFlx{(‑i)}[3] \dicFieldCat{fyz.} \dicDirectTranslationCS{vodivost}
\dicEntry[leiðrétta] \dicTerm{leið··rétt|a} \dicIPA{{l}{ei}{ð}{r}{j}{\textepsilon}{h}{\textsubring{d}}{a}} \dicPos{v}[2] \dicFlx{(‑i, ‑)}[1] \dicFlx{acc} \dicDirectTranslationCS{opravit, opravovat, (z)korigovat (chyby ap.)} \dicExampleIS{leiðrétta verkefnið} \dicExampleCS{opravit cvičení}
\dicEntry[leiðrétting] \dicTerm{leið··rétt·ing} \dicIPA{{l}{ei}{ð}{r}{j}{\textepsilon}{h}{\textsubring{d}}{i}{\ng}{\r{g}}} \dicPos{f}[4] \dicFlx{(‑ar, ‑ar)}[5] \dicDirectTranslationCS{oprava, opravení, korekce, korektura} \dicExampleIS{leiðrétting  á mistökum} \dicExampleCS{oprava chyb}
\dicEntry[leiðsla] \dicTerm{leiðsl|a} \dicIPA{{l}{ei}{ð}{s}{\textsubring{d}}{l}{a}} \dicPos{f}[1] \dicFlx{(‑u, ‑ur)}[19] \textbf{1.} \dicSynonym{leiðsögn} \dicDirectTranslationCS{pokyn, instrukce}  \textbf{2.} \dicSynonym{rör} \dicDirectTranslationCS{roura, trubka, potrubí}  \textbf{3.} \dicSynonym{rafleiðsla} \dicDirectTranslationCS{(elektrické) vedení}  \textbf{4.} \dicSynonym*{ofurhrifning} \dicDirectTranslationCS{vytržení, trans} \dicExampleIS{vera í leiðslu} \dicExampleCS{být v~transu}
\dicEntry[leiðsögn] \dicTerm{leið··|sögn} \dicIPA{{l}{ei}{ð}{s}{\oe}{\r{g}}{\textsubring{n}}} \dicPos{f}[7] \dicFlx{(‑sagnar)}[19] \dicDirectTranslationCS{vedení, navádění, provádění}
\dicEntry[leiðsögumaður] \dicTerm{leið·sögu··|maður} \dicIPA{{l}{ei}{ð}{s}{\oe}{\textbabygamma}{\textscy}{m}{a}{ð}{\textscy}{\textsubring{r}}} \dicPos{m}[13] \dicFlx{(‑manns, ‑menn)}[2] \dicDirectTranslationCS{průvodce, průvodkyně} \dicExampleIS{leiðsögumaður ferðamanna} \dicExampleCS{průvodce turistů}
\dicEntry[leiðtogi] \dicTerm{leið··tog|i} \dicIPA{{l}{ei}{ð}{t\smash{\textsuperscript{h}}}{\textopeno i}{j}{\textsci}} \dicPos{m}[1] \dicFlx{(‑a, ‑ar)}[1] \dicSynonym{foringi} \dicDirectTranslationCS{vedoucí představitel(ka), lídr(yně)} \dicExampleIS{Flokkurinn hefur kosið sér leiðtoga.} \dicExampleCS{Strana si zvolila vedoucího představitele.}
\dicEntry[leiður] \dicTerm{leiður} \dicsymFrequent\  \dicIPA{{l}{ei}{\textlengthmark}{ð}{\textscy}{\textsubring{r}}} \dicPos{adj}[2]\dicFlx{}[-6] \textbf{1.} \dicSynonym{sorgmæddur} \dicDirectTranslationCS{zarmoucený, cítící lítost} \dicExampleIS{leiður út af e‑u} \dicExampleCS{zarmoucený kvůli čemu}  \textbf{2.} \dicSynonym{hvimleiður} \dicDirectTranslationCS{znuděný, unavený (z~poslouchání nářků ap.)} \dicExampleIS{vera orðinn leiður á verkinu} \dicExampleCS{být už unavený z~práce}
\dicEntry[leif] \dicTerm{leif} \dicsymFrequent\  \dicIPA{{l}{ei}{\textlengthmark}{f}} \dicPos{f}[4] \dicFlx{(‑ar, ‑ar)}[1] \textbf{1.} \dicLangCat{básn.} \dicSynonym{arfur} \dicDirectTranslationCS{dědictví, pozůstalost}  \textbf{2.} \dicPhraseIS{leifar} \dicFlx{pl} \dicDirectTranslationCS{zbytky, ostatky} \dicExampleIS{matarleifar} \dicExampleCS{zbytky jídla}  \textbf{3.} \dicFieldCat{mat.} \dicSynonym{afgangur} \dicDirectTranslationCS{zbytek}
\dicEntry[leifa] \dicTerm{leif|a} \dicIPA{{l}{ei}{\textlengthmark}{v}{a}} \dicPos{v}[2] \dicFlx{(‑ði, ‑t)}[99] \dicFlx{dat} \dicDirectTranslationCS{(za)nechat (jídlo nedojedené ap.)} \dicExampleIS{leifa matnum} \dicExampleCS{nechat jídlo}
\dicEntry[leiftra] \dicTerm{leiftr|a} \dicIPA{{l}{ei}{f}{\textsubring{d}}{r}{a}} \dicPos{v}[1] \dicFlx{(‑aði)}[1] \dicSynonym{glitra} \dicDirectTranslationCS{blýskat se, blyštět se, třpytit se} \dicExampleIS{Stjörnurnar leiftra á himninum.} \dicExampleCS{Hvězdy se blyští na obloze.}
\dicEntry[leiftur] \dicTerm{leift|ur\smash{\textsuperscript{1}}} \dicIPA{{l}{ei}{f}{\textsubring{d}}{\textscy}{\textsubring{r}}} \dicPos{m}[5] \dicFlx{(‑urs, ‑rar)}[1] \dicFieldCat{zool.} \dicDirectTranslationCS{plískavice běloboká} \textit{(l.~{\textLA{Lagenorhynchus acutus}})}  \dicsymPhoto\ 
\dicFigure{ds_image_leiftur_1_1.jpg}{Leiftur}{Leiftur - Twp, CC BY-SA 3.0}
\dicEntry[leiftur] \dicTerm{leiftur\smash{\textsuperscript{2}}} \dicIPA{{l}{ei}{f}{\textsubring{d}}{\textscy}{\textsubring{r}}} \dicPos{n}[2] \dicFlx{(‑s, ‑)}[25] \dicSynonym{blossi} \dicDirectTranslationCS{(za)blýsknutí, zá\-blesk}
\dicEntry[leifturhraði] \dicTerm{leiftur··hrað|i} \dicIPA{{l}{ei}{f}{\textsubring{d}}{\textscy}{\textsubring{r}}{a}{ð}{\textsci}} \dicPos{m}[1] \dicFlx{(‑a)}[3] \dicLangCat{přen.} \dicDirectTranslationCS{rychlost blesku, blesková rychlost}
\dicEntry[leiftursókn] \dicTerm{leiftur··sókn} \dicIPA{{l}{ei}{f}{\textsubring{d}}{\textscy}{\textsubring{r}}{s}{ou}{h}{\r{g}}{\textsubring{n}}} \dicPos{f}[7] \dicFlx{(‑ar, ‑ir)}[1] \dicDirectTranslationCS{bleskový útok}
\dicEntry[leiga] \dicTerm{leig|a} \dicsymFrequent\  \dicIPA{{l}{ei}{\textlengthmark}{\textbabygamma}{a}} \dicPos{f}[1] \dicFlx{(‑u, ‑ur)}[19] \textbf{1.} \dicSynonym*{það að leigja} \dicDirectTranslationCS{(pro)nájem} \dicExampleIS{leiga á íbúð} \dicExampleCS{pronájem bytu}  \textbf{2.} \dicSynonym*{leigugjald} \dicDirectTranslationCS{nájem, nájemné} \dicExampleIS{borga leiguna} \dicExampleCS{zaplatit nájem}
\dicEntry[leigja] \dicTerm{leig|ja} \dicsymFrequent\  \dicIPA{{l}{ei}{j}{\textlengthmark}{a}} \dicPos{v}[2] \dicFlx{(‑ði, ‑t)}[89] \dicFlx{acc} \dicDirectTranslationCS{(pro)najmout, (pro)najímat} \dicExampleIS{leigja íbúð af e‑m} \dicExampleCS{pronajmout si od (koho) byt}
\dicEntry[leigjandi] \dicTerm{leigj··|andi} \dicIPA{{l}{ei}{j}{\textlengthmark}{a}{n}{\textsubring{d}}{\textsci}} \dicPos{m}[2] \dicFlx{(‑anda, ‑endur)}[1] \dicDirectTranslationCS{(pod)nájemník, (pod)nájemnice}
\dicEntry[leigubifreið] \dicTerm{leigu··bif·reið} \dicIPA{{l}{ei}{\textlengthmark}{\textbabygamma}{\textscy}{\textsubring{b}}{\textsci}{v}{r}{ei}{\texttheta}} \dicPos{f}[4] \dicFlx{(‑ar, ‑ar\,/\addthin ‑ir)}[30] \dicLangCat{form.} \dicDirectTranslationCS{taxi, taxislužba}
\dicEntry[leigubílastöð] \textls[10]{\dicTerm{leigu·bíla··stöð} \dicIPA{{l}{ei}{\textlengthmark}{\textbabygamma}{\textscy}{\textsubring{b}}{i}{l}{a}{s}{\textsubring{d}}{\oe}{\texttheta}} \dicPos{f}[6] \dicFlx{(‑var, ‑var)}[1] \textbf{1.} \dicDirectTranslationCS{taxi\-služba}} \dicIndirectTranslationCS{(firma)}  \textbf{2.} \dicDirectTranslationCS{stanoviště taxi}
\dicEntry[leigubíll] \dicTerm{leigu··bíl|l} \dicIPA{{l}{ei}{\textlengthmark}{\textbabygamma}{\textscy}{\textsubring{b}}{i}{\textsubring{d}}{\textsubring{l}}} \dicPos{m}[6] \dicFlx{(‑s, ‑ar)}[48] \dicDirectTranslationCS{taxi, taxík}
\dicEntry[leigubílstjóri] \dicTerm{leigu··bíl·stjór|i} \dicIPA{{l}{ei}{\textlengthmark}{\textbabygamma}{\textscy}{\textsubring{b}}{i}{l}{s}{\textsubring{d}}{j}{ou}{r}{\textsci}} \dicPos{m}[1] \dicFlx{(‑a, ‑ar)}[1] \dicDirectTranslationCS{taxikář(ka), řidič(ka) taxi}
\dicEntry[leiguflug] \dicTerm{leigu··flug} \dicIPA{{l}{ei}{\textlengthmark}{\textbabygamma}{\textscy}{f}{l}{\textscy}{x}} \dicPos{n}[2] \dicFlx{(‑s, ‑)}[5] \dicDirectTranslationCS{charterový let}
\dicEntry[leiguíbúð] \dicTerm{leigu··í·búð} \dicIPA{{l}{ei}{\textlengthmark}{\textbabygamma}{\textscy}{i}{\textsubring{b}}{u}{\texttheta}} \dicPos{f}[7] \dicFlx{(‑ar, ‑ir)}[1] \dicDirectTranslationCS{pronajatý byt}
\dicEntry[leiguliði] \dicTerm{leigu··lið|i} \dicIPA{{l}{ei}{\textlengthmark}{\textbabygamma}{\textscy}{l}{\textsci}{ð}{\textsci}} \dicPos{m}[1] \dicFlx{(‑a, ‑ar)}[1] \dicDirectTranslationCS{nájemce, nájemkyně}
\dicEntry[leigumiðlari] \dicTerm{leigu··miðl·ar|i} \dicIPA{{l}{ei}{\textlengthmark}{\textbabygamma}{\textscy}{m}{\textsci}{ð}{l}{a}{r}{\textsci}} \dicPos{m}[1] \dicFlx{(‑a, ‑ar)}[13] \dicDirectTranslationCS{zprostředkovatel(ka) pronájmu}
\dicEntry[leigumorð] \dicTerm{leigu··morð} \dicIPA{{l}{ei}{\textlengthmark}{\textbabygamma}{\textscy}{m}{\textopeno}{r}{\texttheta}} \dicPos{n}[2] \dicFlx{(‑s, ‑)}[5] \dicDirectTranslationCS{nájemná vražda}
\dicEntry[leigumorðingi] \dicTerm{leigu··morð·ing|i} \dicIPA{{l}{ei}{\textlengthmark}{\textbabygamma}{\textscy}{m}{\textopeno}{r}{ð}{i}{\textltailn}{\r{\textObardotlessj}}{\textsci}} \dicPos{m}[1] \dicFlx{(‑ja, ‑jar)}[14] \dicDirectTranslationCS{nájemný vrah, nájemná vražedkyně}
\dicEntry[leigusali] \dicTerm{leigu··sal|i} \dicIPA{{l}{ei}{\textlengthmark}{\textbabygamma}{\textscy}{s}{a}{l}{\textsci}} \dicPos{m}[1] \dicFlx{(‑a, ‑ar)}[8] \dicDirectTranslationCS{pronajímatel(ka)}
\dicEntry[leigusamningur] \dicTerm{leigu··samn·ing|ur} \dicIPA{{l}{ei}{\textlengthmark}{\textbabygamma}{\textscy}{s}{a}{m}{n}{i}{\ng}{\r{g}}{\textscy}{\textsubring{r}}} \dicPos{m}[6] \dicFlx{(‑s, ‑ar)}[8] \dicDirectTranslationCS{nájem\-ní smlouva}
\dicEntry[leigutaki] \dicTerm{leigu··tak|i} \dicIPA{{l}{ei}{\textlengthmark}{\textbabygamma}{\textscy}{t\smash{\textsuperscript{h}}}{a}{\r{\textObardotlessj}}{\textsci}} \dicPos{m}[1] \dicFlx{(‑a, ‑ar)}[8] \dicDirectTranslationCS{nájemce, nájemkyně}
\dicEntry[leik] \dicTerm{leik} \dicIPA{{l}{ei}{\textlengthmark}{\r{g}}} \dicPos{v} \dicFlx{ind praes sg 1 pers} \dicLink{leika}
\dicEntry[leika] \dicTerm{leika} \dicsymFrequent\  \dicIPA{{l}{ei}{\textlengthmark}{\r{g}}{a}} \dicPos{v}[6] \dicFlx{(leik, lék, lékum, léki, leikið)}[63] \dicFlx{acc} \textbf{1.} \dicDirectTranslationCS{(za)hrát} \dicIndirectTranslationCS{(bavit se pro zábavu)};  \dicPhraseIS{leika sér} \dicDirectTranslationCS{hrát si} \dicExampleIS{Börnin leika sér saman.} \dicExampleCS{Děti si spolu hrají.};  \dicPhraseIS{leika sér að e‑u} \dicDirectTranslationCS{hrát si s~(čím)} \dicExampleIS{leika sér að dúkkum} \dicExampleCS{hrát si s~panenkami}  \textbf{2.} \dicDirectTranslationCS{(za)hrát (si)} \dicIndirectTranslationCS{(provozovat jako sport)} \dicExampleIS{leika knattspyrnu} \dicExampleCS{hrát fotbal}  \textbf{3.} \dicDirectTranslationCS{(za)hrát} \dicIndirectTranslationCS{(provozovat hudbu)};  \dicPhraseIS{leika á hljóðfæri} \dicDirectTranslationCS{hrát na hudební nástroj} \dicExampleIS{Hún lék á fiðlu sína.} \dicExampleCS{Hrála na housle.}  \textbf{4.} \dicDirectTranslationCS{(za)hrát} \dicIndirectTranslationCS{(předvádět divadelní hru nebo osobu v~ní)} \dicExampleIS{leika Hamlet} \dicExampleCS{hrát Hamleta}  \textbf{5.} \dicPhraseIS{leika e‑n grátt\,/\addthin hart\,/\addthin illa} \dicDirectTranslationCS{zatočit s~(kým), nemazlit se s~(kým)};  \dicIdiom{leika}[af]{ \dicPhraseIS{leika af sér}} \dicDirectTranslationCS{udělat špatný tah (v~šachu ap.)};  \dicIdiom{leika}[á]{ \dicPhraseIS{leika á e‑n}} \dicDirectTranslationCS{vyzrát na (koho), přelstít (koho)}; { \dicPhraseIS{leika á als oddi, leika á alsoddi}} \dicLangCat{přen.} \dicDirectTranslationCS{mít výbornou náladu, být v~dobrém rozmaru}; { \dicPhraseIS{það leikur ekki á tveim tungum}} \dicFlx{impers} \dicLangCat{přen.} \dicDirectTranslationCS{o~tom není pochyb};  \dicIdiom{leika}[eftir]{ \dicPhraseIS{leika e‑ð eftir e‑m}} \dicDirectTranslationCS{zopakovat (co) po (kom)};  \dicIdiom{leika}[í]{ \dicPhraseIS{leika í e‑u}} \dicDirectTranslationCS{hrát v~(čem) (ve filmu ap.)};  \dicIdiom{leika}[undir]{ \dicPhraseIS{leika undir (e‑ð)}} \dicFieldCat{hud.} \dicDirectTranslationCS{doprovodit\,/\addthin doprovázet ((co))};  \dicIdiom{leika}[við]{ \dicPhraseIS{leika við e‑n}} \dicDirectTranslationCS{hrát si s~(kým)} \dicExampleIS{Lilja leikur við mömmu sína.} \dicExampleCS{Lilja si hraje s~maminkou.}; { \dicPhraseIS{leika við hvern sinn fingur}} \dicLangCat{přen.} \dicDirectTranslationCS{být dobře naladěn, mít dobrou náladu}
\dicEntry[leikandi] \dicTerm{leik··|andi\smash{\textsuperscript{1}}} \dicIPA{{l}{ei}{\textlengthmark}{\r{g}}{a}{n}{\textsubring{d}}{\textsci}} \dicPos{m}[2] \dicFlx{(‑anda, ‑endur)}[1] \textbf{1.} \dicDirectTranslationCS{hráč(ka) (ve fotbale ap.)}  \textbf{2.} \dicDirectTranslationCS{herec, herečka}
\dicEntry[leikandi] \dicTerm{leik··andi\smash{\textsuperscript{2}}} \dicIPA{{l}{ei}{\textlengthmark}{\r{g}}{a}{n}{\textsubring{d}}{\textsci}} \dicPos{adv} \dicDirectTranslationCS{hravě, s~lehkostí} \dicExampleIS{stökkva leikandi yfir skurðinn} \dicExampleCS{přeskočit hravě (přes) příkop}
\dicEntry[leikaraskapur] \dicTerm{leikara··skap|ur} \dicIPA{{l}{ei}{\textlengthmark}{\r{g}}{a}{r}{a}{s}{\r{g}}{a}{\textsubring{b}}{\textscy}{\textsubring{r}}} \dicPos{m}[10] \dicFlx{(‑ar)}[15] \textbf{1.} \dicSynonym{fíflalæti} \dicDirectTranslationCS{šaškoviny, komediantství}  \textbf{2.} \dicSynonym{látalæti} \dicDirectTranslationCS{předstírání, hraní}
\dicEntry[leikari] \dicTerm{leik··ar|i} \dicsymFrequent\  \dicIPA{{l}{ei}{\textlengthmark}{\r{g}}{a}{r}{\textsci}} \dicPos{m}[1] \dicFlx{(‑a, ‑ar)}[13] \dicDirectTranslationCS{herec, herečka} \dicExampleIS{Leikarinn hlaut verðlaun.} \dicExampleCS{Herec obdržel ocenění.}
\dicEntry[leikbrúða] \dicTerm{leik··brúð|a} \dicIPA{{l}{ei}{\textlengthmark}{\r{g}}{\textsubring{b}}{r}{u}{ð}{a}} \dicPos{f}[1] \dicFlx{(‑u, ‑ur)}[19] \textbf{1.} \dicSynonym{brúða} \dicDirectTranslationCS{loutka} \dicIndirectTranslationCS{(v~loutkovém divadle)}  \textbf{2.} \dicSynonym{handbendi} \dicDirectTranslationCS{loutka} \dicIndirectTranslationCS{(zmanipulovaný člověk)}
\dicEntry[leikdómur] \dicTerm{leik··dóm|ur} \dicIPA{{l}{ei}{\textlengthmark}{\r{g}}{\textsubring{d}}{ou}{m}{\textscy}{\textsubring{r}}} \dicPos{m}[6] \dicFlx{(‑s, ‑ar)}[10] \dicDirectTranslationCS{divadelní kritika}
\dicEntry[leikfang] \dicTerm{leik··|fang} \dicIPA{{l}{ei}{\textlengthmark}{\r{g}}{f}{au}{\ng}{\r{g}}} \dicPos{n}[2] \dicFlx{(‑fangs, ‑föng)}[8] \dicDirectTranslationCS{hračka (na hraní)}
\dicEntry[leikfangabúð] \dicTerm{leik·fanga··búð} \dicIPA{{l}{ei}{\textlengthmark}{\r{g}}{f}{au}{\ng}{\r{g}}{a}{\textsubring{b}}{u}{\texttheta}} \dicPos{f}[7] \dicFlx{(‑ar, ‑ir)}[1] \dicDirectTranslationCS{hračkářství}
\dicEntry[leikfélag] \dicTerm{leik··fé·|lag} \dicIPA{{l}{ei}{\textlengthmark}{\r{g}}{f}{j}{\textepsilon}{l}{a}{x}} \dicPos{n}[2] \dicFlx{(‑lags, ‑lög)}[8] \dicDirectTranslationCS{divadelní společnost}
\dicEntry[leikfélagi] \dicTerm{leik··fé·lag|i} \dicIPA{{l}{ei}{\textlengthmark}{\r{g}}{f}{j}{\textepsilon}{l}{a}{i}{j}{\textsci}} \dicPos{m}[1] \dicFlx{(‑a, ‑ar)}[8] \dicDirectTranslationCS{kamarád(ka) z~dětství}
\dicEntry[leikfimi] \dicTerm{leik··fim|i} \dicIPA{{l}{ei}{\textlengthmark}{\r{g}}{f}{\textsci}{m}{\textsci}} \dicPos{f}[3] \dicFlx{(‑i)}[3] \dicDirectTranslationCS{tělocvik, (tělesné) cvičení}
\dicEntry[leikfimisalur] \dicTerm{leik·fimi··sal|ur} \dicIPA{{l}{ei}{\textlengthmark}{\r{g}}{f}{\textsci}{m}{\textsci}{s}{a}{l}{\textscy}{\textsubring{r}}} \dicPos{m}[10] \dicFlx{(‑ar, ‑ir)}[14] \dicDirectTranslationCS{tělocvična}
\dicEntry[leikfimiskennari] \dicTerm{leik·fimis··kenn·ar|i} \dicIPA{{l}\-{ei}\-{\textlengthmark}\-{\r{g}}\-{f}\-{\textsci}\-{m}\-{\textsci}\-{s}\-{c\smash{\textsuperscript{h}}}\-{\textepsilon}\-{n}\-{a}\-{r}\-{\textsci}\-} \dicPos{m}[1] \dicFlx{(‑a, ‑ar)}[13] \dicDirectTranslationCS{učitel(ka) tělesné výchovy, tělocvikář(ka)}
\dicEntry[leikgagnrýnandi] \dicTerm{leik··gagn·rýn·|andi} \dicIPA{{l}{ei}{\r{g}}{\textlengthmark}{a}{\r{g}}{n}{r}{i}{n}{a}{n}{\textsubring{d}}{\textsci}} \dicPos{m}[2] \dicFlx{(‑anda, ‑endur)}[1] \dicDirectTranslationCS{divadelní kritik\,/\addthin kritička}
\dicEntry[leikgerð] \dicTerm{leik··gerð} \dicIPA{{l}{ei}{\textlengthmark}{\r{g}}{\r{\textObardotlessj}}{\textepsilon}{r}{\texttheta}} \dicPos{f}[7] \dicFlx{(‑ar, ‑ir)}[1] \dicFieldCat{lit.} \dicDirectTranslationCS{dramatizace} \dicExampleIS{leikgerð Sölku Völku} \dicExampleCS{dramatizace Salky Valky}
\dicEntry[leikgrind] \dicTerm{leik··grind} \dicIPA{{l}{ei}{\r{g}}{\textlengthmark}{r}{\textsci}{n}{\textsubring{d}}} \dicPos{f}[8] \dicFlx{(‑ar, ‑ur)}[1] \dicDirectTranslationCS{(dětská) ohrádka}
\dicEntry[leikhús] \dicTerm{leik··hús} \dicIPA{{l}{ei}{\textlengthmark}{\r{g}}{h}{u}{s}} \dicPos{n}[2] \dicFlx{(‑s, ‑)}[5] \dicDirectTranslationCS{divadlo} \dicExampleIS{sýning í leikhúsinu} \dicExampleCS{představení v~divadle}
\dicEntry[leikhússtjóri] \dicTerm{leik·hús··stjór|i} \dicIPA{{l}{ei}{\textlengthmark}{\r{g}}{h}{u}{s}{\textsubring{d}}{j}{ou}{r}{\textsci}} \dicPos{m}[1] \dicFlx{(‑a, ‑ar)}[1] \dicDirectTranslationCS{divadelní ředitel(ka), ředitel(ka) divadla}
\dicEntry[leikið] \dicTerm{leikið} \dicIPA{{l}{ei}{\textlengthmark}{\r{\textObardotlessj}}{\textsci}{\texttheta}} \dicPos{v} \dicFlx{supin} \dicLink{leika}
\dicEntry[leikinn] \dicTerm{leikinn} \dicIPA{{l}{ei}{\textlengthmark}{\r{\textObardotlessj}}{\textsci}{\textsubring{n}}} \dicPos{adj}[6]\dicFlx{}[-2] \dicSynonym{snjall} \dicDirectTranslationCS{obratný, šikovný, dovedný} \dicExampleIS{vera leikinn í e‑u} \dicExampleCS{být šikovný v~(čem)};  \dicPhraseIS{vera grátt\,/\addthin illa leikinn} \dicDirectTranslationCS{nést známky špatného zacházení}
\dicEntry[leikjanámskeið] \dicTerm{leikja··nám·skeið} \dicIPA{{l}{ei}{\textlengthmark}{\r{\textObardotlessj}}{a}{n}{au}{m}{s}{\r{\textObardotlessj}}{ei}{\texttheta}} \dicPos{n}[2] \dicFlx{(‑s, ‑)}[5] \dicDirectTranslationCS{tábor, (dětský) prázdninový kurz}
\dicEntry[leikjatölva] \dicTerm{leikja··tölv|a} \dicIPA{{l}{ei}{\textlengthmark}{\r{\textObardotlessj}}{a}{t\smash{\textsuperscript{h}}}{\oe}{l}{v}{a}} \dicPos{f}[1] \dicFlx{(‑u, ‑ur)}[7] \dicFieldCat{poč.} \dicDirectTranslationCS{(hrací) konzole}
\dicEntry[leikkerahjól] \dicTerm{leik·kera··hjól} \dicIPA{{l}{ei}{\textlengthmark}{\r{g}}{c\smash{\textsuperscript{h}}}{\textepsilon}{r}{a}{\c{c}}{ou}{\textsubring{l}}} \dicPos{n}[2] \dicFlx{(‑s, ‑)}[5] \dicDirectTranslationCS{hrnčířské kolo}
\dicEntry[leikkerasmiður] \dicTerm{leik·kera··smið|ur} \dicIPA{{l}{ei}{\textlengthmark}{\r{g}}{c\smash{\textsuperscript{h}}}{\textepsilon}{r}{a}{s}{m}{\textsci}{ð}{\textscy}{\textsubring{r}}} \dicPos{m}[9] \dicFlx{(‑s, ‑ir)}[11] \dicDirectTranslationCS{hrnčíř(ka), keramik, keramička}
\dicEntry[leikkona] \dicTerm{leik··kon|a} \dicIPA{{l}{ei}{\textlengthmark}{\r{g}}{k\smash{\textsuperscript{h}}}{\textopeno}{n}{a}} \dicPos{f}[1] \dicFlx{(‑u, ‑ur)}[27] \dicDirectTranslationCS{herečka}
\dicEntry[leiklist] \dicTerm{leik··list} \dicIPA{{l}{ei}{\textlengthmark}{\r{g}}{l}{\textsci}{s}{\textsubring{d}}} \dicPos{f}[7] \dicFlx{(‑ar)}[3] \dicDirectTranslationCS{divadlo, divadelní umění}
\dicEntry[leiklistarskóli] \dicTerm{leik·listar··skól|i} \dicIPA{{l}{ei}{\textlengthmark}{\r{g}}{l}{\textsci}{s}{\textsubring{d}}{a}{\textsubring{r}}{s}{\r{g}}{ou}{l}{\textsci}} \dicPos{m}[1] \dicFlx{(‑a, ‑ar)}[1] \dicFieldCat{škol.} \dicDirectTranslationCS{divadelní škola}
\dicEntry[leikmaður] \dicTerm{leik··|maður} \dicIPA{{l}{ei}{\textlengthmark}{\r{g}}{m}{a}{ð}{\textscy}{\textsubring{r}}} \dicPos{m}[13] \dicFlx{(‑manns, ‑menn)}[2] \textbf{1.} \dicDirectTranslationCS{laik, laička, neodborník, neodbornice}  \textbf{2.} \dicDirectTranslationCS{hráč(ka)}  \textbf{3.} \dicSynonym{leikari} \dicDirectTranslationCS{herec, herečka}
\dicEntry[leikmunur] \dicTerm{leik··mun|ur} \dicIPA{{l}{ei}{\textlengthmark}{\r{g}}{m}{\textscy}{n}{\textscy}{\textsubring{r}}} \dicPos{m}[10] \dicFlx{(‑ar, ‑ir)}[11] \dicDirectTranslationCS{rekvizita}
\dicEntry[leikmynd] \dicTerm{leik··mynd} \dicIPA{{l}{ei}{\textlengthmark}{\r{g}}{m}{\textsci}{n}{\textsubring{d}}} \dicPos{f}[7] \dicFlx{(‑ar, ‑ir)}[1] \dicDirectTranslationCS{scénografie, scéna}
\dicEntry[leikmyndahönnuður] \dicTerm{leik·mynda··hönn·uð|ur} \dicIPA{{l}\-{ei}\-{\textlengthmark}\-{\r{g}}\-{m}\-{\textsci}\-{n}\-{\textsubring{d}}\-{a}\-{h}\-{\oe}\-{n}\-{\textscy}\-{ð}\-{\textscy}\-{\textsubring{r}}\-} \dicPos{m}[10] \dicFlx{(‑ar, ‑ir)}[4] \dicDirectTranslationCS{jevištní výtvarník\,/\addthin výtvarnice, scénograf(ka)}
\dicEntry[leikni] \dicTerm{leikn|i} \dicIPA{{l}{ei}{h}{\r{g}}{n}{\textsci}} \dicPos{f}[3] \dicFlx{(‑i)}[3] \dicDirectTranslationCS{šikovnost, obratnost, zručnost, dovednost} \dicExampleIS{leikni í íþróttum} \dicExampleCS{dovednost ve sportu}
\dicEntry[leikregla] \dicTerm{leik··regl|a} \dicIPA{{l}{ei}{\textlengthmark}{\r{g}}{r}{\textepsilon}{\r{g}}{l}{a}} \dicPos{f}[1] \dicFlx{(‑u, ‑ur)}[19] \dicDirectTranslationCS{pravidlo hry}
\dicEntry[leikrit] \dicTerm{leik··rit} \dicIPA{{l}{ei}{\textlengthmark}{\r{g}}{r}{\textsci}{\textsubring{d}}} \dicPos{n}[2] \dicFlx{(‑s, ‑)}[5] \dicDirectTranslationCS{divadelní hra}
\dicEntry[leikritahöfundur] \dicTerm{leik·rita··höf·und|ur} \dicIPA{{l}\-{ei}\-{\textlengthmark}\-{\r{g}}\-{r}\-{\textsci}\-{\textsubring{d}}\-{a}\-{h}\-{\oe}\-{v}\-{\textscy}\-{n}\-{\textsubring{d}}\-{\textscy}\-{\textsubring{r}}\-} \dicPos{m}[6] \dicFlx{(‑ar, ‑ar)}[56] \dicDirectTranslationCS{(divadelní) dramatik, dramatička} \dicExampleIS{frægur tékkneskur leikritahöfundur} \dicExampleCS{slavný český dramatik}
\dicEntry[leikrænn] \dicTerm{leik··rænn} \dicIPA{{l}{ei}{\textlengthmark}{\r{g}}{r}{a}{i}{\textsubring{d}}{\textsubring{n}}} \dicPos{adj}[7]\dicFlx{}[-1] \dicDirectTranslationCS{divadelní}
\dicEntry[leikskáld] \dicTerm{leik··skáld} \dicIPA{{l}{ei}{\textlengthmark}{\r{g}}{s}{\r{g}}{au}{l}{\textsubring{d}}} \dicPos{n}[2] \dicFlx{(‑s, ‑)}[5] \dicDirectTranslationCS{dramatik, dramatička, autor(ka) divadelních her}
\dicEntry[leikskóladeild] \dicTerm{leik·skóla··deild} \dicIPA{{l}{ei}{\textlengthmark}{\r{g}}{s}{\r{g}}{ou}{l}{a}{\textsubring{d}}{ei}{l}{\textsubring{d}}} \dicPos{f}[7] \dicFlx{(‑ar, ‑ir)}[1] \dicDirectTranslationCS{oddělení mateřské školky} \dicIndirectTranslationCS{(v~rámci školy ap.)}
\dicEntry[leikskóladvöl] \dicTerm{leik·skóla··|dvöl} \dicIPA{{l}{ei}{\textlengthmark}{\r{g}}{s}{\r{g}}{ou}{l}{a}{\textsubring{d}}{v}{\oe}{\textsubring{l}}} \dicPos{f}[7] \dicFlx{(‑dvalar, ‑dvalir)}[16] \dicDirectTranslationCS{pobyt ve školce}
\dicEntry[leikskólakennari] \dicTerm{leik·skóla··kenn·ar|i} \dicIPA{{l}\-{ei}\-{\textlengthmark}\-{\r{g}}\-{s}\-{\r{g}}\-{ou}\-{l}\-{a}\-{c\smash{\textsuperscript{h}}}\-{\textepsilon}\-{n}\-{a}\-{r}\-{\textsci}\-} \dicPos{m}[1] \dicFlx{(‑a, ‑ar)}[13] \dicDirectTranslationCS{učitel(ka) v~mateřské školce}
\dicEntry[leikskóli] \dicTerm{leik··skól|i} \dicIPA{{l}{ei}{\textlengthmark}{\r{g}}{s}{\r{g}}{ou}{l}{\textsci}} \dicPos{m}[1] \dicFlx{(‑a, ‑ar)}[1] \dicDirectTranslationCS{(mateřská) školka}
\dicEntry[leikskrá] \dicTerm{leik··skrá} \dicIPA{{l}{ei}{\textlengthmark}{\r{g}}{s}{\r{g}}{r}{au}} \dicPos{f}[4] \dicFlx{(‑r\,/\addthin ‑ar, ‑r)}[21] \dicDirectTranslationCS{divadelní program}
\dicEntry[leikslok] \dicTerm{leiks··lok} \dicIPA{{l}{ei}{\textlengthmark}{\r{g}}{s}{l}{\textopeno}{\r{g}}} \dicPos{n}[2] \dicFlx{pl}[1] \textbf{1.} \dicSynonym{niðurstaða} \dicDirectTranslationCS{závěr, rozhodnutí}  \textbf{2.} \dicDirectTranslationCS{výsledek (zápasu, utkání ap.)}
\dicEntry[leiksoppur] \dicTerm{leik··sopp|ur} \dicIPA{{l}{ei}{\textlengthmark}{\r{g}}{s}{\textopeno}{h}{\textsubring{b}}{\textscy}{\textsubring{r}}} \dicPos{m}[6] \dicFlx{(‑s, ‑ar)}[22] \dicDirectTranslationCS{loutka} \dicIndirectTranslationCS{(člověk vnitřně prázdný)}
\dicEntry[leikstjóri] \dicTerm{leik··stjór|i} \dicIPA{{l}{ei}{\textlengthmark}{\r{g}}{s}{\textsubring{d}}{j}{ou}{r}{\textsci}} \dicPos{m}[1] \dicFlx{(‑a, ‑ar)}[1] \textbf{1.} \dicDirectTranslationCS{režisér(ka)}  \textbf{2.} \dicFieldCat{sport.} \dicDirectTranslationCS{rozhodčí}
\dicEntry[leikstjórn] \dicTerm{leik··stjórn} \dicIPA{{l}{ei}{\textlengthmark}{\r{g}}{s}{\textsubring{d}}{j}{ou}{r}{\textsubring{d}}{\textsubring{n}}} \dicPos{f}[7] \dicFlx{(‑ar)}[3] \dicDirectTranslationCS{režie, režírování}
\dicEntry[leikstýra] \dicTerm{leik··stýr|a} \dicIPA{{l}{ei}{\textlengthmark}{\r{g}}{s}{\textsubring{d}}{i}{r}{a}} \dicPos{v}[2] \dicFlx{(‑ði, ‑t)}[105] \dicFlx{dat} \dicDirectTranslationCS{režírovat}
\dicEntry[leiksvið] \dicTerm{leik··svið} \dicIPA{{l}{ei}{\textlengthmark}{\r{g}}{s}{v}{\textsci}{\texttheta}} \dicPos{n}[2] \dicFlx{(‑s, ‑)}[5] \dicDirectTranslationCS{jeviště, scéna}
\dicEntry[leiksýning] \dicTerm{leik··sýn·ing} \dicIPA{{l}{ei}{\textlengthmark}{\r{g}}{s}{i}{n}{i}{\ng}{\r{g}}} \dicPos{f}[4] \dicFlx{(‑ar, ‑ar)}[5] \dicDirectTranslationCS{(divadelní) představení}
\dicEntry[leiktíð] \dicTerm{leik··tíð} \dicIPA{{l}{ei}{\textlengthmark}{\r{g}}{t\smash{\textsuperscript{h}}}{i}{\texttheta}} \dicPos{f}[7] \dicFlx{(‑ar, ‑ir)}[1] \dicFieldCat{sport.} \dicDirectTranslationCS{sezóna (fotbalová ap.)}
\dicEntry[leikur] \dicTerm{leik|ur} \dicsymFrequent\  \dicIPA{{l}{ei}{\textlengthmark}{\r{g}}{\textscy}{\textsubring{r}}} \dicPos{m}[6] \dicFlx{(‑s, ‑ir\,/\addthin ‑ar)}[67] \textbf{1.} \dicFlx{pl (-ir)} \dicSynonym{barnaleikur} \dicDirectTranslationCS{hra, hraní} \dicExampleIS{leikur barna} \dicExampleCS{hra dětí};  \dicPhraseIS{bregða á leik} \dicDirectTranslationCS{hrát si, skotačit}  \textbf{2.} \dicFlx{pl (-ar)} \dicSynonym{viðureign} \dicDirectTranslationCS{boj, střetnutí, utkání}  \textbf{3.} \dicFlx{pl (-ar)} \dicSynonym*{íþróttakeppni} \dicDirectTranslationCS{zápas, utkání, hra} \dicExampleIS{Ólympíuleikar} \dicExampleCS{olympijské hry};  \dicPhraseIS{fara á leikinn} \dicDirectTranslationCS{jít na zápas}  \textbf{4.} \dicSynonym{leikrit} \dicDirectTranslationCS{hra (divadelní ap.), představení}  \textbf{5.} \dicDirectTranslationCS{tah} \dicIndirectTranslationCS{(posunutí šachové figury ap.)};  \dicPhraseIS{eiga leik} \dicDirectTranslationCS{být na tahu};  \dicIdiom{leikur}{ \dicPhraseIS{á nýjan leik}} \dicFlx{adv} \dicDirectTranslationCS{znovu, opět}; { \dicPhraseIS{e‑að er leikur einn}} \dicLangCat{přen.} \dicDirectTranslationCS{(co) je hračka}; { \dicPhraseIS{með leik}} \dicFlx{adv} \dicDirectTranslationCS{snadno, lehce}; { \dicPhraseIS{vera úr leik}} \dicLangCat{přen.} \dicDirectTranslationCS{být mimo hru}; { \dicPhraseIS{við illan leik}} \dicFlx{adv} \dicDirectTranslationCS{horko těžko}
\dicEntry[leikvangur] \dicTerm{leik··vang|ur} \dicIPA{{l}{ei}{\textlengthmark}{\r{g}}{v}{au}{\ng}{\r{g}}{\textscy}{\textsubring{r}}} \dicPos{m}[6] \dicFlx{(‑s, ‑ar)}[25] \dicSynonym{íþróttavöllur} \dicDirectTranslationCS{stadión}
\dicEntry[leikvöllur] \dicTerm{leik··|völlur} \dicIPA{{l}{ei}{\textlengthmark}{\r{g}}{v}{\oe}{\textsubring{d}}{l}{\textscy}{\textsubring{r}}} \dicPos{m}[11] \dicFlx{(‑vallar, ‑vellir)}[5] \textbf{1.} \dicFieldCat{sport.} \dicSynonym{í\-þrótta\-völlur} \dicDirectTranslationCS{(sportovní) hřiště}  \textbf{2.} \dicSynonym*{barnaleikvöllur} \dicDirectTranslationCS{dětské hřiště}
\dicEntry[leir] \dicTerm{leir} \dicsymFrequent\  \dicIPA{{l}{ei}{\textlengthmark}{\textsubring{r}}} \dicPos{m}[4] \dicFlx{(‑s)}[18] \dicDirectTranslationCS{jíl, hlína} \dicExampleIS{hvít englabörn úr leir} \dicExampleCS{bílí andělíčci z~hlíny}
\dicEntry[leira] \dicTerm{leir|a\smash{\textsuperscript{1}}} \dicIPA{{l}{ei}{\textlengthmark}{r}{a}} \dicPos{f}[1] \dicFlx{(‑u, ‑ur)}[19] \dicDirectTranslationCS{jílovitá půda}
\dicEntry[Leira] \dicTerm{Leir|a\smash{\textsuperscript{2}}} \dicIPA{{l}{ei}{\textlengthmark}{r}{a}} \dicPos{f}[1] \dicFlx{(‑u)}[6] \dicFieldCat{geog.} \dicDirectTranslationCS{Loira} \dicIndirectTranslationCS{(řeka ve Francii)}
\dicEntry[leirburður] \dicTerm{leir··burð|ur} \dicIPA{{l}{ei}{r}{\textsubring{b}}{\textscy}{r}{ð}{\textscy}{\textsubring{r}}} \dicPos{m}[10] \dicFlx{(‑ar)}[7] \dicDirectTranslationCS{špatná poezie}
\dicEntry[leirker] \dicTerm{leir··ker} \dicIPA{{l}{ei}{\textsubring{r}}{c\smash{\textsuperscript{h}}}{\textepsilon}{\textsubring{r}}} \dicPos{n}[2] \dicFlx{(‑s, ‑)}[6] \dicDirectTranslationCS{keramika} \dicIndirectTranslationCS{(především z~archeologických nálezů)}
\dicEntry[leirlist] \dicTerm{leir··list} \dicIPA{{l}{ei}{r}{l}{\textsci}{s}{\textsubring{d}}} \dicPos{f}[7] \dicFlx{(‑ar)}[3] \dicDirectTranslationCS{keramika, keramické umění}
\dicEntry[leirmunur] \dicTerm{leir··mun|ur} \dicIPA{{l}{ei}{r}{m}{\textscy}{n}{\textscy}{\textsubring{r}}} \dicPos{m}[10] \dicFlx{(‑ar, ‑ir)}[11] \dicDirectTranslationCS{keramika, keramický výrobek}
\dicEntry[leit] \dicTerm{leit\smash{\textsuperscript{1}}} \dicsymFrequent\  \dicIPA{{l}{ei}{\textlengthmark}{\textsubring{d}}} \dicPos{f}[7] \dicFlx{(‑ar, ‑ir)}[1] \textbf{1.} \dicDirectTranslationCS{hledání, pátrání, vyhledávání} \dicExampleIS{Leitin er mjög umfangsmikil.} \dicExampleCS{Pátrání je velmi rozsáhlé.}  \textbf{2.} \dicPhraseIS{leitir} \dicFlx{pl} \dicSynonym{smölun} \dicDirectTranslationCS{shánění ovcí};  \dicPhraseIS{fara í leitir} \dicDirectTranslationCS{jet shánět ovce}
\dicEntry[leit] \dicTerm{leit\smash{\textsuperscript{2}}} \dicIPA{{l}{ei}{\textlengthmark}{\textsubring{d}}} \dicPos{v} \dicFlx{ind pf sg 1 pers} \dicLink{líta}
\dicEntry[leita] \dicTerm{leit|a} \dicsymFrequent\  \dicIPA{{l}{ei}{\textlengthmark}{\textsubring{d}}{a}} \dicPos{v}[1] \dicFlx{(‑aði)}[1] \dicFlx{gen} \textbf{1.} \dicDirectTranslationCS{(vy)hledat, (vy)pátrat} \dicIndirectTranslationCS{(snažit se nalézt)};  \dicPhraseIS{leita e‑s} \dicDirectTranslationCS{hledat (co)} \dicExampleIS{leita bátsins} \dicExampleCS{hledat člun};  \dicPhraseIS{leita að e‑u\,/\addthin e‑m} \dicDirectTranslationCS{hledat (co\,/\addthin koho)} \dicExampleIS{leita að lyklunum} \dicExampleCS{hledat klíče}  \textbf{2.} \dicDirectTranslationCS{vyhledat, navštívit (lékaře ap.)};  \dicPhraseIS{leita læknis} \dicDirectTranslationCS{vyhledat doktora};  \dicIdiom{leita}[á]{ \dicPhraseIS{leita á e‑m}} \dicDirectTranslationCS{prohledat (koho), prošacovat (koho)};  \dicIdiom{leita}[eftir]{ \dicPhraseIS{leita eftir e‑u við e‑n}} \dicDirectTranslationCS{požádat (koho) o~(co)};  \dicIdiom{leita}[fyrir]{ \dicPhraseIS{leita fyrir sér um e‑ð}} \dicDirectTranslationCS{shánět se po (čem), shánět (co) (práci ap.)};  \dicIdiom{leita}[til]{ \dicPhraseIS{leita til e‑rs}} \dicDirectTranslationCS{obrátit se na (koho)} \dicExampleIS{leita til hennar með vandamál sín} \dicExampleCS{obrátit se na ni se svými problémy};  \dicIdiom{leita}[uppi]{ \dicPhraseIS{leita e‑n uppi}} \dicDirectTranslationCS{vyhledat (koho), vypátrat (koho)};  \dicIdiom{leitast}[við]{ \dicPhraseIS{leitast við}} \dicFlx{refl} \dicDirectTranslationCS{pokoušet se, snažit se}
\dicEntry[leitandi] \dicTerm{leit··andi} \dicIPA{{l}{ei}{\textlengthmark}{\textsubring{d}}{a}{n}{\textsubring{d}}{\textsci}} \dicPos{adj}[13] \dicFlx{indecl}[1] \dicDirectTranslationCS{hledající, pátrající}
\dicEntry[leitarflokkur] \dicTerm{leitar··flokk|ur} \dicIPA{{l}{ei}{\textlengthmark}{\textsubring{d}}{a}{\textsubring{r}}{f}{l}{\textopeno}{h}{\r{g}}{\textscy}{\textsubring{r}}} \dicPos{m}[6] \dicFlx{(‑s, ‑ar)}[8] \dicDirectTranslationCS{pátrací skupina\,/\addthin četa}
\dicEntry[leitarvél] \dicTerm{leitar··vél} \dicIPA{{l}{ei}{\textlengthmark}{\textsubring{d}}{a}{r}{v}{j}{\textepsilon}{\textsubring{l}}} \dicPos{f}[4] \dicFlx{(‑ar, ‑ar)}[1] \dicFieldCat{poč.} \dicDirectTranslationCS{(internetový) vyhledávač}
\dicEntry[leiti] \dicTerm{leiti} \dicIPA{{l}{ei}{\textlengthmark}{\textsubring{d}}{\textsci}} \dicPos{n}[2] \dicFlx{(‑s, ‑)}[14] \dicSynonym{hæð} \dicDirectTranslationCS{vrch, vyvýšenina, kopec};  \dicPhraseIS{á næsta leiti} \dicFlx{adv} \dicDirectTranslationCS{blízko, nedaleko (časově), za dveřmi\,/\addthin rohem} \dicExampleIS{Hátíðin er á næsta leiti.} \dicExampleCS{Svátky jsou za dveřmi.}
\dicEntry[leitun] \dicTerm{leit|un} \dicIPA{{l}{ei}{\textlengthmark}{\textsubring{d}}{\textscy}{\textsubring{n}}} \dicPos{f}[7] \dicFlx{(‑unar)}[9] \dicDirectTranslationCS{hledání, vyhledávání, pátrání};  \dicPhraseIS{leitun að e‑u} \dicDirectTranslationCS{hledání (čeho)};  \dicPhraseIS{það er leitun á\,/\addthin að e‑u} \dicDirectTranslationCS{je vzácné\,/\addthin těžké nalézt (co)}
\dicEntry[lek] \dicTerm{lek} \dicIPA{{l}{\textepsilon}{\textlengthmark}{\r{g}}} \dicPos{v} \dicFlx{ind praes sg 1 pers} \dicLink{leka}
\dicEntry[leka] \dicTerm{leka} \dicsymFrequent\  \dicIPA{{l}{\textepsilon}{\textlengthmark}{\r{g}}{a}} \dicPos{v}[6] \dicFlx{(lek, lak, lákum, læki, lekið)}[17] \textbf{1.} \dicSynonym*{falla í dropum} \dicDirectTranslationCS{téct, stékat, kapat} \dicExampleIS{Vatnið lekur niður á gólfið.} \dicExampleCS{Voda stéká dolů na podlahu.}  \textbf{2.} \dicSynonym*{vera lekur} \dicDirectTranslationCS{téct, prosakovat} \dicIndirectTranslationCS{(propouštět kapalinu)}  \textbf{3.} \dicLangCat{přen.} \dicDirectTranslationCS{uniknout (na veřejnost), prosáknout (na veřejnost) (zpráva ap.)}
\dicEntry[lekandi] \dicTerm{lek··|andi} \dicIPA{{l}{\textepsilon}{\textlengthmark}{\r{g}}{a}{n}{\textsubring{d}}{\textsci}} \dicPos{m}[2] \dicFlx{(‑anda)}[2] \dicFieldCat{med.} \dicDirectTranslationCS{kapavka}
\dicEntry[leki] \dicTerm{lek|i} \dicIPA{{l}{\textepsilon}{\textlengthmark}{\r{\textObardotlessj}}{\textsci}} \dicPos{m}[1] \dicFlx{(‑a, ‑ar)}[1] \textbf{1.} \dicDirectTranslationCS{kapání, odkapávání} \dicExampleIS{leki úr þaki} \dicExampleCS{kapání ze střechy}  \textbf{2.} \dicDirectTranslationCS{prosakování, zatékání} \dicExampleIS{Það kom leki að bátnum.} \dicExampleCS{Loď začala zatékat.}  \textbf{3.} \dicDirectTranslationCS{únik, unikání (informací ap.)}
\dicEntry[lekið] \dicTerm{lekið} \dicIPA{{l}{\textepsilon}{\textlengthmark}{\r{\textObardotlessj}}{\textsci}{\texttheta}} \dicPos{v} \dicFlx{supin} \dicLink{leka}
\dicEntry[lektor] \dicTerm{lektor} \dicIPA{{l}{\textepsilon}{x}{\textsubring{d}}{\textopeno}{\textsubring{r}}} \dicPos{m}[4] \dicFlx{(‑s, ‑ar)}[14] \dicFieldCat{škol.} \dicDirectTranslationCS{odborný asistent, odborná asistentka}
\dicEntry[lekur] \dicTerm{lekur} \dicIPA{{l}{\textepsilon}{\textlengthmark}{\r{g}}{\textscy}{\textsubring{r}}} \dicPos{adj}[1]\dicFlx{}[-1] \dicDirectTranslationCS{zatékající, prosakující} \dicExampleIS{lekur bátur} \dicExampleCS{zatékající člun}
\dicEntry[lem] \dicTerm{lem} \dicIPA{{l}{\textepsilon}{\textlengthmark}{\textsubring{m}}} \dicPos{v} \dicFlx{ind praes sg 1 pers} \dicLink{lemja}
\dicEntry[lemdi] \dicTerm{lemdi} \dicIPA{{l}{\textepsilon}{m}{\textsubring{d}}{\textsci}} \dicPos{v} \dicFlx{con pf sg 1 pers} \dicLink{lemja}
\dicEntry[lemja] \dicTerm{lemja} \dicsymFrequent\  \dicIPA{{l}{\textepsilon}{m}{j}{a}} \dicPos{v}[4] \dicFlx{(lem, lamdi, lömdum, lemdi, lamið)}[3] \dicFlx{acc} \dicSynonym{slá\smash{\textsuperscript{3}}} \dicDirectTranslationCS{uhodit, udeřit, tlouct, mlátit} \dicExampleIS{lemja e‑n með barefli} \dicExampleCS{uhodit (koho) pálkou}
\dicEntry[lemstra] \dicTerm{lemstr|a} \dicIPA{{l}{\textepsilon}{m}{s}{\textsubring{d}}{r}{a}} \dicPos{v}[1] \dicFlx{(‑aði)}[1] \dicFlx{acc} \dicSynonym{limlesta} \dicDirectTranslationCS{zmrzačit}
\dicEntry[lend] \dicTerm{lend} \dicIPA{{l}{\textepsilon}{n}{\textsubring{d}}} \dicPos{f}[4] \dicFlx{(‑ar, ‑ar)}[1] \textbf{1.} \dicDirectTranslationCS{zadek (koně ap.)}  \textbf{2.} \dicPhraseIS{lendar} \dicFlx{pl} \dicSynonym{mjaðmir} \dicDirectTranslationCS{boky, bedra}
\dicEntry[lenda] \dicTerm{len|da} \dicsymFrequent\  \dicIPA{{l}{\textepsilon}{n}{\textsubring{d}}{a}} \dicPos{v}[2] \dicFlx{(‑ti, ‑t)}[41] \dicFlx{dat} \textbf{1.} \dicSynonym*{koma að landi} \dicDirectTranslationCS{přistát, přistávat, dosednout, dosedat} \dicExampleIS{lenda flugvélinni} \dicExampleCS{přistát s~letadlem}  \textbf{2.} \dicDirectTranslationCS{přistát, dopadnout (na zadek ap.)};  \dicIdiom{lenda}[í]{ \dicPhraseIS{lenda í e‑u}} \dicSynonym*{rata í e‑ð} \dicDirectTranslationCS{spadnout do (čeho), dostat se do (čeho)} \dicExampleIS{lenda í vandræðum} \dicExampleCS{dostat se do potíží};  \dicIdiom{lenda}[saman]{ \dicPhraseIS{e‑jum lendir saman}} \dicFlx{impers} \dicLangCat{přen.} \dicDirectTranslationCS{(kdo) se dostává do křížku} \dicExampleIS{Þeim lenti saman.} \dicExampleCS{Dostali se do křížku.}
\dicEntry[lending] \dicTerm{lend··ing} \dicIPA{{l}{\textepsilon}{n}{\textsubring{d}}{i}{\ng}{\r{g}}} \dicPos{f}[4] \dicFlx{(‑ar, ‑ar)}[5] \dicDirectTranslationCS{přistání, přistávání, dosednutí} \dicExampleIS{mjúk lending} \dicExampleCS{měkké přistání}
\dicEntry[lengd] \dicTerm{lengd} \dicsymFrequent\  \dicIPA{{l}{\textepsilon}{\ng}{\textsubring{d}}} \dicPos{f}[7] \dicFlx{(‑ar, ‑ir)}[1] \textbf{1.} \dicDirectTranslationCS{délka} \dicExampleIS{í fullri lengd} \dicExampleCS{v~plné délce}  \textbf{2.} \dicFieldCat{geog.} \dicDirectTranslationCS{zeměpisná délka};  \dicIdiom{lengd}{ \dicPhraseIS{til lengdar}} \dicFlx{adv} \dicDirectTranslationCS{dlouhodobě, v~dlouhém časovém období}
\dicEntry[lengdarbaugur] \dicTerm{lengdar··baug|ur} \dicIPA{{l}{\textepsilon}{\ng}{\textsubring{d}}{a}{r}{\textsubring{b}}{\oe i}{\textbabygamma}{\textscy}{\textsubring{r}}} \dicPos{m}[6] \dicFlx{(‑s, ‑ar)}[24] \dicFieldCat{geog.} \dicDirectTranslationCS{poledník}
\dicEntry[lengi] \dicTerm{lengi} \dicsymFrequent\  \dicIPA{{l}{ei}{\textltailn}{\r{\textObardotlessj}}{\textsci}} \dicPos{adv} \dicFlx{(comp lengur, sup lengst)} \dicSynonym*{langan tíma} \dicDirectTranslationCS{dlouho, dlouze} \dicExampleIS{lengi vetrar} \dicExampleCS{po většinu zimy};  \dicPhraseIS{lengi vel} \dicFlx{adv} \dicDirectTranslationCS{dlouho, dlouhou dobu}
\dicEntry[lenging] \dicTerm{leng··ing} \dicIPA{{l}{ei}{\textltailn}{\r{\textObardotlessj}}{i}{\ng}{\r{g}}} \dicPos{f}[4] \dicFlx{(‑ar, ‑ar)}[5] \textbf{1.} \dicSynonym*{það að lengja e‑ð} \dicDirectTranslationCS{prodloužení, prodlužování} \dicExampleIS{lenging á jólafríinu} \dicExampleCS{prodloužení vánočních prázdnin}  \textbf{2.} \dicSynonym{viðbót} \dicDirectTranslationCS{vylepšení, vylepšování} \dicExampleIS{lenging á bátnum} \dicExampleCS{vylepšení člunu}
\dicEntry[lengja] \dicTerm{lengj|a\smash{\textsuperscript{1}}} \dicIPA{{l}{ei}{\textltailn}{\r{\textObardotlessj}}{a}} \dicPos{f}[1] \dicFlx{(‑u, ‑ur)}[7] \textbf{1.} \dicSynonym{ræma} \dicDirectTranslationCS{pruh, proužek}  \textbf{2.} \dicDirectTranslationCS{karton (cigaret)}
\dicEntry[lengja] \dicTerm{leng|ja\smash{\textsuperscript{2}}} \dicIPA{{l}{ei}{\textltailn}{\r{\textObardotlessj}}{a}} \dicPos{v}[2] \dicFlx{(‑di, ‑t)}[130] \dicFlx{acc} \textbf{1.} \dicDirectTranslationCS{prodloužit, prodlužovat} \dicExampleIS{lengja vinnutímann} \dicExampleCS{prodloužit pracovní dobu}  \textbf{2.} \dicPhraseIS{e‑ð lengir} \dicFlx{impers} \dicDirectTranslationCS{(co) se prodlužuje} \dicExampleIS{Daginn lengir.} \dicExampleCS{Den se prodlužuje.};  \dicIdiom{lengja}[eftir]{ \dicPhraseIS{e‑n lengir eftir e‑m\,/\addthin e‑u}} \dicFlx{impers} \dicDirectTranslationCS{(kdo) touží po (kom\,/\addthin čem)} \dicExampleIS{Mig lengir eftir henni.} \dicExampleCS{Toužím po ní.};  \dicIdiom{lengjast}{ \dicPhraseIS{lengjast}} \dicFlx{refl} \dicDirectTranslationCS{stávat se delší, prodlužovat se} \dicExampleIS{Dagurinn lengist.} \dicExampleCS{Den se prodlužuje.}
\dicEntry[lengra] \dicTerm{lengra} \dicIPA{{l}{\textepsilon}{\ng}{\r{g}}{r}{a}} \dicPos{adv} \dicFlx{comp} \dicLink{langt}
\dicEntry[lengri] \dicTerm{lengri} \dicIPA{{l}{\textepsilon}{\ng}{\r{g}}{r}{\textsci}} \dicPos{adj} \dicFlx{comp m} \dicLink{langur\smash{\textsuperscript{2}}}
\dicEntry[lengst] \dicTerm{lengst} \dicIPA{{l}{\textepsilon}{\ng}{s}{\textsubring{d}}} \dicPos{adv} \dicFlx{sup} \textbf{1.} \dicLink{langt}  \textbf{2.} \dicLink{lengi}
\dicEntry[lengstur] \dicTerm{lengstur} \dicIPA{{l}{\textepsilon}{\ng}{s}{\textsubring{d}}{\textscy}{\textsubring{r}}} \dicPos{adj} \dicFlx{m sg nom sup} \dicLink{langur\smash{\textsuperscript{2}}}
\dicEntry[lengur] \dicTerm{lengur} \dicIPA{{l}{ei}{\ng}{\r{g}}{\textscy}{\textsubring{r}}} \dicPos{adv} \dicFlx{comp} \dicLink{lengi}
\dicEntry[lens] \dicTerm{lens} \dicIPA{{l}{\textepsilon}{n}{s}} \dicPos{n}[2] \dicFlx{(‑)}[24] \dicFieldCat{nám.} \dicSynonym{byr} \dicDirectTranslationCS{zadní vítr}
\dicEntry[lenska] \dicTerm{lensk|a} \dicIPA{{l}{\textepsilon}{n}{s}{\r{g}}{a}} \dicPos{f}[1] \dicFlx{(‑u)}[5] \dicSynonym*{landsvenja} \dicDirectTranslationCS{zvyk, obyčej, zvyklost}
\dicEntry[lenskur] \dicTerm{lenskur} \dicIPA{{l}{\textepsilon}{n}{s}{\r{g}}{\textscy}{\textsubring{r}}} \dicPos{adj}[1]\dicFlx{}[-1] \dicPhraseIS{hvers lenskur} \dicDirectTranslationCS{jaké národnosti} \dicExampleIS{Ég vissi ekki hvers lenskur hann var.} \dicExampleCS{Nevěděl jsem, jaké je národnosti.}
\dicEntry[lep] \dicTerm{lep} \dicIPA{{l}{\textepsilon}{\textlengthmark}{\textsubring{b}}} \dicPos{v} \dicFlx{ind praes sg 1 pers} \dicLink{lepja}
\dicEntry[lepja] \dicTerm{lepja} \dicIPA{{l}{\textepsilon}{\textlengthmark}{\textsubring{b}}{j}{a}} \dicPos{v}[4] \dicFlx{(lep, lapti, löptum, lepti, lapið)}[47] \dicFlx{acc} \textbf{1.} \dicDirectTranslationCS{(s)lízat, vylízat} \dicIndirectTranslationCS{(o~zvířatech)} \dicExampleIS{lepja úr skálinni} \dicExampleCS{lízat z~misky}  \textbf{2.} \dicDirectTranslationCS{papouškovat, opakovat (bez porozumění ap.)}
\dicEntry[leppa] \dicTerm{lepp|a} \dicIPA{{l}{\textepsilon}{h}{\textsubring{b}}{a}} \dicPos{v}[1] \dicFlx{(‑aði)}[1] \dicFlx{acc} \textbf{1.} \dicPhraseIS{leppa sig} \dicSynonym*{klæða sig} \dicDirectTranslationCS{obléct se, oblékat se}  \textbf{2.} \dicSynonym{bæta} \dicDirectTranslationCS{opravit, vyspravit} \dicExampleIS{leppa upp á e‑ð} \dicExampleCS{vyspravit (co)}  \textbf{3.} \dicDirectTranslationCS{vázat} \dicIndirectTranslationCS{(v~šachách)}
\dicEntry[Leppalúði] \dicTerm{Leppa··lúð|i} \dicIPA{{l}{\textepsilon}{h}{\textsubring{b}}{a}{l}{u}{ð}{\textsci}} \dicPos{m}[1] \dicFlx{(‑a)}[5] \dicFlx{prop} \dicFieldCat{pov.} \dicDirectTranslationCS{Leppalúði} \dicIndirectTranslationCS{(nesmírně líný obr, otec 13 vánočních skřítků)}
\dicEntry[leppríki] \dicTerm{lepp··ríki} \dicIPA{{l}{\textepsilon}{h}{\textsubring{b}}{r}{i}{\r{\textObardotlessj}}{\textsci}} \dicPos{n}[2] \dicFlx{(‑s, ‑)}[16] \dicFieldCat{pol.} \dicDirectTranslationCS{loutkový\,/\addthin satelitní stát}
\dicEntry[leppstjórn] \dicTerm{lepp··stjórn} \dicIPA{{l}{\textepsilon}{h}{\textsubring{b}}{s}{\textsubring{d}}{j}{ou}{r}{\textsubring{d}}{\textsubring{n}}} \dicPos{f}[7] \dicFlx{(‑ar, ‑ir)}[1] \dicDirectTranslationCS{loutková vláda}
\dicEntry[leppur] \dicTerm{lepp|ur} \dicIPA{{l}{\textepsilon}{h}{\textsubring{b}}{\textscy}{\textsubring{r}}} \dicPos{m}[6] \dicFlx{(‑s, ‑ar)}[22] \textbf{1.} \dicSynonym{bót} \dicDirectTranslationCS{záplata}  \textbf{2.} \dicDirectTranslationCS{vložka do bot}  \textbf{3.} \dicPhraseIS{leppar} \dicFlx{pl} \dicSynonym*{tötrar} \dicDirectTranslationCS{hadry} \dicIndirectTranslationCS{(o~šatech)}  \textbf{4.} \dicSynonym{handbendi} \dicDirectTranslationCS{loutka, nohsled}  \textbf{5.} \dicDirectTranslationCS{klapka (na oku ap.)}  \textbf{6.} \dicDirectTranslationCS{vazba} \dicIndirectTranslationCS{(v~šachách)}
\dicEntry[lepti] \dicTerm{lepti} \dicIPA{{l}{\textepsilon}{f}{\textsubring{d}}{\textsci}} \dicPos{v} \dicFlx{con pf sg 1 pers} \dicLink{lepja}
\dicEntry[lerki] \dicTerm{lerki} \dicIPA{{l}{\textepsilon}{\textsubring{r}}{\r{\textObardotlessj}}{\textsci}} \dicPos{n}[2] \dicFlx{(‑s)}[20] \dicFieldCat{bot.} \dicDirectTranslationCS{modřín} \textit{(l.~{\textLA{Larix}})}  \dicsymPhoto\ 
\dicFigure{286.jpg}{Lerki}{Lerki - Zicha Ondřej, Biolib, Copyright/CC-BY-NC}
\dicEntry[les] \dicTerm{les} \dicIPA{{l}{\textepsilon}{\textlengthmark}{s}} \dicPos{v} \dicFlx{ind praes sg 1 pers} \dicLink{lesa}
\dicEntry[lesa] \dicTerm{lesa} \dicsymFrequent\  \dicIPA{{l}{\textepsilon}{\textlengthmark}{s}{a}} \dicPos{v}[6] \dicFlx{(les, las, lásum, læsi, lesið)}[16] \dicFlx{acc} \textbf{1.} \dicDirectTranslationCS{(pře)číst} \dicExampleIS{lesa söguna upphátt} \dicExampleCS{číst příběh nahlas}  \textbf{2.} \dicSynonym{læra} \dicDirectTranslationCS{studovat, učit se} \dicExampleIS{Hann les sagnfræði í háskóla.} \dicExampleCS{Studuje historii na vysoké škole.}  \textbf{3.} \dicLangCat{zast.} \dicDirectTranslationCS{sbírat (květy ap.)} \dicExampleIS{lesa aldin af trjánum} \dicExampleCS{sbírat plody ze stromů};  \dicIdiom{lesa}[fyrir]{ \dicPhraseIS{lesa fyrir}} \dicDirectTranslationCS{diktovat}; { \dicPhraseIS{lesa fyrir e‑m e‑ð}} \dicDirectTranslationCS{(na)diktovat (komu co)} \dicExampleIS{lesa henni bréfið fyrir} \dicExampleCS{diktovat jí dopis};  \dicIdiom{lesa}[í]{ \dicPhraseIS{lesa í lófa}} \dicLangCat{přen.} \dicDirectTranslationCS{číst z~dlaně};  \dicIdiom{lesa}[saman]{ \dicPhraseIS{lesa saman e‑ð og e‑ð}} \dicDirectTranslationCS{srovnávat\,/\addthin porovnávat (co) s~(čím)} \dicExampleIS{lesa frumritið og afritið saman} \dicExampleCS{srovnávat původní text s~opisem};  \dicIdiom{lesa}[til]{ \dicPhraseIS{lesa sér til e‑ð}} \dicDirectTranslationCS{seznámit se s~(čím), dozvědět se (co) nového (čtením)};  \dicIdiom{lesa}[undir]{ \dicPhraseIS{lesa undir próf}} \dicDirectTranslationCS{učit se na zkoušku, studovat na zkoušku} \dicExampleIS{lesa undir skóla} \dicExampleCS{učit se do školy};  \dicIdiom{lesa}[upp]{ \dicPhraseIS{lesa upp e‑ð}} \dicDirectTranslationCS{recitovat (co), přednášet (co)} \dicExampleIS{lesa upp ljóð} \dicExampleCS{recitovat báseň}; { \dicPhraseIS{lesa sig upp}} \dicSynonym{klifra} \dicDirectTranslationCS{šplhat, vyšplhat (se)};  \dicIdiom{lesa}[yfir]{ \dicPhraseIS{lesa e‑ð yfir}} \dicDirectTranslationCS{pročíst si (co)}
\dicEntry[lesandi] \dicTerm{les··|andi} \dicsymFrequent\  \dicIPA{{l}{\textepsilon}{\textlengthmark}{s}{a}{n}{\textsubring{d}}{\textsci}} \dicPos{m}[2] \dicFlx{(‑anda, ‑endur)}[1] \dicDirectTranslationCS{čtenář(ka)} \dicExampleIS{Lesendurnir eru ánægðir með bókina.} \dicExampleCS{Čtenářům se knížka líbí.}
\dicEntry[lesbía] \dicTerm{lesbí|a} \dicIPA{{l}{\textepsilon}{s}{\textsubring{b}}{i}{j}{a}} \dicPos{f}[1] \dicFlx{(‑u, ‑ur)}[7] \dicDirectTranslationCS{lesba, lesbička}
\dicEntry[lesblinda] \dicTerm{les··blind|a} \dicIPA{{l}{\textepsilon}{\textlengthmark}{s}{\textsubring{b}}{l}{\textsci}{n}{\textsubring{d}}{a}} \dicPos{f}[1] \dicFlx{(‑u)}[5] \dicFieldCat{med.} \dicDirectTranslationCS{dyslexie}
\dicEntry[lesbók] \dicTerm{les··|bók} \dicIPA{{l}{\textepsilon}{\textlengthmark}{s}{\textsubring{b}}{ou}{\r{g}}} \dicPos{f}[8] \dicFlx{(‑bókar, ‑bækur)}[5] \dicDirectTranslationCS{čítanka}
\dicEntry[lesefni] \dicTerm{les··efni} \dicIPA{{l}{\textepsilon}{\textlengthmark}{s}{\textepsilon}{\textsubring{b}}{n}{\textsci}} \dicPos{n}[2] \dicFlx{(‑s, ‑)}[14] \dicDirectTranslationCS{čtení, četba, čtivo}
\dicEntry[lesgleraugu] \dicTerm{les··gler·augu} \dicIPA{{l}{\textepsilon}{\textlengthmark}{s}{\r{g}}{l}{\textepsilon}{r}{\oe i}{\textbabygamma}{\textscy}} \dicPos{n}[1] \dicFlx{pl}[3] \dicDirectTranslationCS{brýle na čtení}
\dicEntry[lesið] \dicTerm{lesið} \dicIPA{{l}{\textepsilon}{\textlengthmark}{s}{\textsci}{\texttheta}} \dicPos{v} \dicFlx{supin} \dicLink{lesa}
\dicEntry[lesinn] \dicTerm{lesinn} \dicIPA{{l}{\textepsilon}{\textlengthmark}{s}{\textsci}{\textsubring{n}}} \dicPos{adj}[6]\dicFlx{}[-2] \dicDirectTranslationCS{sečtělý}
\dicEntry[lesminni] \dicTerm{les··minni} \dicIPA{{l}{\textepsilon}{\textlengthmark}{s}{m}{\textsci}{n}{\textsci}} \dicPos{n}[2] \dicFlx{(‑s, ‑)}[14] \dicFieldCat{poč.} \dicDirectTranslationCS{ROM} \dicIndirectTranslationCS{(typ elektronické paměti používaný pro uložení firmware v~elektronických přístrojích)} \dicAntonym{vinnsluminni}
\dicEntry[lesning] \dicTerm{les··ning} \dicIPA{{l}{\textepsilon}{s}{\textsubring{d}}{n}{i}{\ng}{\r{g}}} \dicPos{f}[4] \dicFlx{(‑ar, ‑ar)}[5] \dicDirectTranslationCS{četba, čtení}
\dicEntry[Lesótó] \dicTerm{Lesótó} \dicIPA{{l}{\textepsilon}{\textlengthmark}{s}{ou}{\textsubring{d}}{ou}} \dicPos{n}[4] \dicFlx{indecl}[2] \dicFieldCat{geog.} \dicDirectTranslationCS{Lesotho}
\dicEntry[Lesótói] \dicTerm{Lesótó|i} \dicIPA{{l}{\textepsilon}{\textlengthmark}{s}{ou}{\textsubring{d}}{ou}{\textsci}} \dicPos{m}[1] \dicFlx{(‑a, ‑ar)}[1] \dicLink{Lesótómaður}
\dicEntry[Lesótómaður] \dicTerm{Lesótó··|maður}\dicTerm{, Lesótói} \dicIPA{{l}\-{\textepsilon}\-{\textlengthmark}\-{s}\-{ou}\-{\textsubring{d}}\-{ou}\-{m}\-{a}\-{ð}\-{\textscy}\-{\textsubring{r}}\-} \dicPos{m}[13] \dicFlx{(‑manns, ‑menn)}[2] \dicDirectTranslationCS{Lesothan(ka)}
\dicEntry[lesótóskur] \dicTerm{lesótóskur} \dicIPA{{l}{\textepsilon}{\textlengthmark}{s}{ou}{\textsubring{d}}{ou}{s}{\r{g}}{\textscy}{\textsubring{r}}} \dicPos{adj}[1]\dicFlx{}[-6] \dicDirectTranslationCS{lesothský}
\dicEntry[lesskilningur] \dicTerm{les··skiln·ing|ur} \dicIPA{{l}{\textepsilon}{\textlengthmark}{s}{\r{\textObardotlessj}}{\textsci}{l}{n}{i}{\ng}{\r{g}}{\textscy}{\textsubring{r}}} \dicPos{m}[6] \dicFlx{(‑s)}[9] \dicDirectTranslationCS{porozumění (psanému) textu}
\dicEntry[lest] \dicTerm{lest} \dicsymFrequent\  \dicIPA{{l}{\textepsilon}{s}{\textsubring{d}}} \dicPos{f}[7] \dicFlx{(‑ar, ‑ir)}[1] \textbf{1.} \dicDirectTranslationCS{vlak, vláček} \dicExampleIS{fara með lest} \dicExampleCS{jet vlakem};  \dicPhraseIS{missa af lestinni} {\textbf{a.}} \dicDirectTranslationCS{zmeškat vlak};  {\textbf{b.}} \dicLangCat{přen.} \dicDirectTranslationCS{nechat si ujet vlak} \dicIndirectTranslationCS{(promeškat příležitost)}  \textbf{2.} \dicFieldCat{nám.} \dicSynonym{farmur} \dicDirectTranslationCS{náklad}  \textbf{3.} \dicSynonym{tonn} \dicDirectTranslationCS{(lodní) tuna}  \textbf{4.} \dicFieldCat{nám.} \dicSynonym*{lestarrúm} \dicDirectTranslationCS{nákladní prostor}  \textbf{5.} \dicSynonym{röð} \dicDirectTranslationCS{řada, karavana} \dicExampleIS{löng lest af bílum} \dicExampleCS{dlouhá řada aut}
\dicEntry[lesta] \dicTerm{lest|a} \dicIPA{{l}{\textepsilon}{s}{\textsubring{d}}{a}} \dicPos{v}[1] \dicFlx{(‑aði)}[1] \dicSynonym{hlaða\smash{\textsuperscript{2}}} \dicDirectTranslationCS{naložit, nakládat} \dicExampleIS{lesta skipið} \dicExampleCS{nakládat loď}
\dicEntry[lestarstjóri] \dicTerm{lestar··stjór|i} \dicIPA{{l}{\textepsilon}{s}{\textsubring{d}}{a}{\textsubring{r}}{s}{\textsubring{d}}{j}{ou}{r}{\textsci}} \dicPos{m}[1] \dicFlx{(‑a, ‑ar)}[1] \dicDirectTranslationCS{strojvedoucí, strojvůdce, strojvůdkyně}
\dicEntry[lestir] \dicTerm{lestir} \dicIPA{{l}{\textepsilon}{s}{\textsubring{d}}{\textsci}{\textsubring{r}}} \dicPos{m} \dicFlx{pl nom} \dicLink{löstur}
\dicEntry[lestrarbók] \dicTerm{lestrar··|bók} \dicIPA{{l}{\textepsilon}{s}{\textsubring{d}}{r}{a}{r}{\textsubring{b}}{ou}{\r{g}}} \dicPos{f}[8] \dicFlx{(‑bókar, ‑bækur)}[5] \dicDirectTranslationCS{čítanka}
\dicEntry[lestrarsalur] \dicTerm{lestrar··sal|ur} \dicIPA{{l}{\textepsilon}{s}{\textsubring{d}}{r}{a}{\textsubring{r}}{s}{a}{l}{\textscy}{\textsubring{r}}} \dicPos{m}[10] \dicFlx{(‑ar, ‑ir)}[14] \dicDirectTranslationCS{čítárna}
\dicEntry[lestur] \dicTerm{lest|ur} \dicsymFrequent\  \dicIPA{{l}{\textepsilon}{s}{\textsubring{d}}{\textscy}{\textsubring{r}}} \dicPos{m}[5] \dicFlx{(‑rar\,/\addthin ‑urs, ‑rar)}[9] \textbf{1.} \dicSynonym*{það að lesa} \dicDirectTranslationCS{čtení, četba} \dicExampleIS{skemmtilegur lestur} \dicExampleCS{zábavné čtení}  \textbf{2.} \dicSynonym{nám} \dicDirectTranslationCS{studium, učení se}  \textbf{3.} \dicSynonym{klifur} \dicDirectTranslationCS{šplhání}
\dicEntry[let] \dicTerm{let} \dicIPA{{l}{\textepsilon}{\textlengthmark}{\textsubring{d}}} \dicPos{v} \dicFlx{ind praes sg 1 pers} \dicLink{letja}
\dicEntry[leti] \dicTerm{let|i} \dicIPA{{l}{\textepsilon}{\textlengthmark}{\textsubring{d}}{\textsci}} \dicPos{f}[3] \dicFlx{(‑i)}[3] \dicDirectTranslationCS{lenost, lenivost, zahálčivost}
\dicEntry[letidýr] \dicTerm{leti··dýr} \dicIPA{{l}{\textepsilon}{\textlengthmark}{\textsubring{d}}{\textsci}{\textsubring{d}}{i}{\textsubring{r}}} \dicPos{n}[2] \dicFlx{(‑s, ‑)}[5] \dicFieldCat{zool.} \dicDirectTranslationCS{lenochod} \textit{(l.~{\textLA{Bradypus}})}  \dicsymPhoto\ 
\dicFigure{ds_image_letidyr_0_1.jpg}{Letidýr}{Letidýr - Stefan Laube, PD}
\dicEntry[letihaugur] \dicTerm{leti··haug|ur} \dicIPA{{l}{\textepsilon}{\textlengthmark}{\textsubring{d}}{\textsci}{h}{\oe i}{\textbabygamma}{\textscy}{\textsubring{r}}} \dicPos{m}[6] \dicFlx{(‑s, ‑ar)}[22] \dicDirectTranslationCS{lenoch, lenoška, povaleč(ka)}
\dicEntry[letilegur] \dicTerm{leti··legur} \dicIPA{{l}{\textepsilon}{\textlengthmark}{\textsubring{d}}{\textsci}{l}{\textepsilon}{\textbabygamma}{\textscy}{\textsubring{r}}} \dicPos{adj}[1]\dicFlx{}[-8] \dicDirectTranslationCS{líný, lenivý}
\dicEntry[letingi] \dicTerm{let··ing|i} \dicIPA{{l}{\textepsilon}{\textlengthmark}{\textsubring{d}}{i}{\textltailn}{\r{\textObardotlessj}}{\textsci}} \dicPos{m}[1] \dicFlx{(‑ja, ‑jar)}[14] \dicDirectTranslationCS{lenoch, lenoška, povaleč(ka)}
\dicEntry[letja] \dicTerm{letja} \dicIPA{{l}{\textepsilon}{\textlengthmark}{\textsubring{d}}{j}{a}} \dicPos{v}[4] \dicFlx{(let, latti, löttum, letti, latt)}[48] \dicFlx{acc} \dicSynonym*{telja úr} \dicDirectTranslationCS{odradit, odrazovat};  \dicPhraseIS{letja e‑n e‑s} \dicDirectTranslationCS{odrazovat (koho) od (čeho)} \dicExampleIS{letja e‑n fararinnar} \dicExampleCS{odrazovat (koho) od cesty}
\dicEntry[letra] \dicTerm{letr|a} \dicIPA{{l}{\textepsilon}{\textlengthmark}{\textsubring{d}}{r}{a}} \dicPos{v}[1] \dicFlx{(‑aði)}[1] \dicFlx{acc} \textbf{1.} \dicSynonym{skrifa} \dicDirectTranslationCS{psát, napsat}  \textbf{2.} \dicSynonym*{höggva stafi í stein} \dicDirectTranslationCS{vysázet\,/\addthin vyrýt písmem}
\dicEntry[Letti] \dicTerm{Lett|i\smash{\textsuperscript{1}}} \dicIPA{{l}{\textepsilon}{h}{\textsubring{d}}{\textsci}} \dicPos{m}[1] \dicFlx{(‑a, ‑ar)}[1] \dicDirectTranslationCS{Lotyš(ka)}
\dicEntry[letti] \dicTerm{letti\smash{\textsuperscript{2}}} \dicIPA{{l}{\textepsilon}{h}{\textsubring{d}}{\textsci}} \dicPos{v} \dicFlx{con pf sg 1 pers} \dicLink{letja}
\dicEntry[Lettland] \dicTerm{Lett··land} \dicIPA{{l}{\textepsilon}{h}{\textsubring{d}}{l}{a}{n}{\textsubring{d}}} \dicPos{n}[2] \dicFlx{(‑s)}[4] \dicFieldCat{geog.} \dicDirectTranslationCS{Lotyšsko}
\dicEntry[lettneska] \dicTerm{lett··nesk|a} \dicIPA{{l}{\textepsilon}{h}{\textsubring{d}}{n}{\textepsilon}{s}{\r{g}}{a}} \dicPos{f}[1] \dicFlx{(‑u)}[5] \dicDirectTranslationCS{lotyština}
\dicEntry[lettneskur] \dicTerm{lett··neskur} \dicIPA{{l}{\textepsilon}{h}{\textsubring{d}}{n}{\textepsilon}{s}{\r{g}}{\textscy}{\textsubring{r}}} \dicPos{adj}[1]\dicFlx{}[-6] \dicDirectTranslationCS{lotyšský}
\dicEntry[letur] \dicTerm{letur} \dicIPA{{l}{\textepsilon}{\textlengthmark}{\textsubring{d}}{\textscy}{\textsubring{r}}} \dicPos{n}[2] \dicFlx{(‑s, ‑)}[25] \textbf{1.} \dicSynonym*{stafategund} \dicDirectTranslationCS{písmo, druh\,/\addthin typ písma, font}  \textbf{2.} \dicSynonym{áletrun} \dicDirectTranslationCS{nápis, epigraf} \dicExampleIS{letur á steini} \dicExampleCS{nápis na kameni}
\dicEntry[leturbreyting] \dicTerm{letur··breyt·ing} \dicIPA{{l}{\textepsilon}{\textlengthmark}{\textsubring{d}}{\textscy}{r}{\textsubring{b}}{r}{ei}{\textsubring{d}}{i}{\ng}{\r{g}}} \dicPos{f}[4] \dicFlx{(‑ar, ‑ar)}[5] \dicDirectTranslationCS{změna písma} \dicIndirectTranslationCS{(z~tučného písma na normální ap.)}
\dicEntry[leturgerð] \dicTerm{letur··gerð} \dicIPA{{l}{\textepsilon}{\textlengthmark}{\textsubring{d}}{\textscy}{r}{\r{\textObardotlessj}}{\textepsilon}{r}{\texttheta}} \dicPos{f}[7] \dicFlx{(‑ar, ‑ir)}[1] \dicSynonym*{leturtegund} \dicDirectTranslationCS{font, znaková sada}
\dicEntry[leturgrafari] \dicTerm{letur··graf·ar|i} \dicIPA{{l}{\textepsilon}{\textlengthmark}{\textsubring{d}}{\textscy}{r}{\r{g}}{r}{a}{v}{a}{r}{\textsci}} \dicPos{m}[1] \dicFlx{(‑a, ‑ar)}[10] \dicDirectTranslationCS{rytec\,/\addthin rytečka (písma)}
\dicEntry[leturstærð] \dicTerm{letur··stærð} \dicIPA{{l}{\textepsilon}{\textlengthmark}{\textsubring{d}}{\textscy}{\textsubring{r}}{s}{\textsubring{d}}{a}{i}{r}{\texttheta}} \dicPos{f}[7] \dicFlx{(‑ar, ‑ir)}[1] \dicDirectTranslationCS{velikost písma}
\dicEntry[lexía] \dicTerm{lexí|a} \dicIPA{{l}{\textepsilon}{x}{s}{i}{j}{a}} \dicPos{f}[1] \dicFlx{(‑u, ‑ur)}[7] \textbf{1.} \dicSynonym{námsefni} \dicDirectTranslationCS{lekce} \dicIndirectTranslationCS{(úsek učební látky)};  \dicPhraseIS{lesa lexíuna sína} \dicDirectTranslationCS{naučit se, udělat domácí úkoly}  \textbf{2.} \dicSynonym{áminning} \dicDirectTranslationCS{lekce, (nepříjemné) poučení}
\dicEntry[leyfa] \dicTerm{leyf|a} \dicsymFrequent\  \dicIPA{{l}{ei}{\textlengthmark}{v}{a}} \dicPos{v}[2] \dicFlx{(‑ði, ‑t)}[106] \dicFlx{dat + acc} \dicDirectTranslationCS{dovolit, dovolovat, povolit, povolovat, nechat, nechávat};  \dicPhraseIS{leyfa e‑m e‑ð} \dicDirectTranslationCS{dovolit (co komu)} \dicExampleIS{leyfa henni að halda á hestinum} \dicExampleCS{dovolit jí vést koně};  \dicIdiom{leyfast}{ \dicPhraseIS{e‑m leyfist að (gera e‑ð)}} \dicFlx{refl impers} \dicDirectTranslationCS{(komu) je dovoleno (udělat (co))} \dicExampleIS{ef mér leyfist að spyrja} \dicExampleCS{pokud se můžu zeptat}
\dicEntry[leyfður] \dicTerm{leyfður} \dicIPA{{l}{ei}{v}{ð}{\textscy}{\textsubring{r}}} \dicPos{adj}[2]\dicFlx{}[-1] \dicDirectTranslationCS{dovolený, povolený} \dicExampleIS{leyfður hraði} \dicExampleCS{povolená rychlost}
\dicEntry[leyfi] \dicTerm{leyfi} \dicsymFrequent\  \dicIPA{{l}{ei}{\textlengthmark}{v}{\textsci}} \dicPos{n}[2] \dicFlx{(‑s, ‑)}[14] \textbf{1.} \dicSynonym{heimild} \dicDirectTranslationCS{licence, oprávnění}  \textbf{2.} \dicDirectTranslationCS{povolení, dovolení, souhlas, svolení} \dicExampleIS{án leyfis} \dicExampleCS{bez povolení};  \dicPhraseIS{með leyfi} \dicFlx{adv} \dicDirectTranslationCS{s~dovolením}  \textbf{3.} \dicSynonym{frí} \dicDirectTranslationCS{dovolená, volno}
\dicEntry[leyfilegur] \dicTerm{leyfi··legur} \dicIPA{{l}{ei}{\textlengthmark}{v}{\textsci}{l}{\textepsilon}{\textbabygamma}{\textscy}{\textsubring{r}}} \dicPos{adj}[1]\dicFlx{}[-8] \dicSynonym{heimill} \dicDirectTranslationCS{povolený, dovolený, přípustný}
\dicEntry[leyfisleysi] \dicTerm{leyfis··leysi} \dicIPA{{l}{ei}{\textlengthmark}{v}{\textsci}{s}{l}{ei}{s}{\textsci}} \dicPos{n}[2] \dicFlx{(‑s)}[20] \dicPhraseIS{gera e‑ð í leyfisleysi} \dicDirectTranslationCS{udělat (co) bez povolení\,/\addthin dovolení\,/\addthin svolení}
\dicEntry[leyfisveiting] \dicTerm{leyfis··veit·ing} \dicIPA{{l}{ei}{\textlengthmark}{v}{\textsci}{s}{v}{ei}{\textsubring{d}}{i}{\ng}{\r{g}}} \dicPos{f}[4] \dicFlx{(‑ar, ‑ar)}[5] \dicDirectTranslationCS{udělení\,/\addthin vydání oprávnění\,/\addthin povolení}
\dicEntry[leyna] \dicTerm{leyn|a} \dicsymFrequent\  \dicIPA{{l}{ei}{\textlengthmark}{n}{a}} \dicPos{v}[2] \dicFlx{(‑di, ‑t)}[141] \dicFlx{dat} \dicDirectTranslationCS{ukrýt, ukrývat, skrýt, skrývat, (u)tajit, zatajovat};  \dicPhraseIS{leyna e‑n e‑u, leyna e‑u fyrir e‑m} \dicDirectTranslationCS{skrýt (co) před (kým), zatajit (co) před (kým), zatajit (komu co)} \dicExampleIS{leyna hana vitneskju um málið} \dicExampleCS{zatajit jí informace o~té záležitosti};  \dicIdiom{leyna}[á]{ \dicPhraseIS{leyna á sér}} \dicDirectTranslationCS{nezdát se} \dicIndirectTranslationCS{(neukazovat pravou skutečnost)};  \dicIdiom{leynast}{ \dicPhraseIS{leynast}} \dicFlx{refl} \dicDirectTranslationCS{skrýt se, skrývat se}
\dicEntry[leynd] \dicTerm{leynd} \dicIPA{{l}{ei}{n}{\textsubring{d}}} \dicPos{f}[7] \dicFlx{(‑ar, ‑ir)}[1] \textbf{1.} \dicSynonym{launung} \dicDirectTranslationCS{tajnost};  \dicPhraseIS{með leynd} \dicFlx{adv} \dicDirectTranslationCS{v~tajnosti, potají}  \textbf{2.} \dicSynonym*{leynistaður} \dicDirectTranslationCS{(místo) utajení}
\begin{xtolerant}{}{1pt}
\dicEntry[leyndardómsfullur] \dicTerm{leyndar·dóms··|fullur} \dicsymFrequent\addthinS\textls[-10]{\dicIPA{{l}\-{ei}\-{n}\-{\textsubring{d}}\-{a}\-{r}\-{\textsubring{d}}\-{ou}\-{m}\-{s}\-{f}\-{\textscy}\-{\textsubring{d}}\-{l}\-{\textscy}\-{\textsubring{r}}}} \dicPos{adj}[10]\addthinS\dicFlx{(comp ‑fyllri, sup ‑fyllstur)}[7] \dicDirectTranslationCS{tajuplný, tajemný, záhadný} \dicExampleIS{leyndardómsfullur á svip} \dicExampleCS{s~tajuplným výrazem ve tváři}
\end{xtolerant}
\dicEntry[leyndardómur] \dicTerm{leyndar··dóm|ur} \dicIPA{{l}{ei}{n}{\textsubring{d}}{a}{r}{\textsubring{d}}{ou}{m}{\textscy}{\textsubring{r}}} \dicPos{m}[6] \dicFlx{(‑s, ‑ar)}[10] \dicDirectTranslationCS{tajemství, tajemnost, záhada} \dicExampleIS{leyndardómur lífsins} \dicExampleCS{tajemství života}
\dicEntry[leyndarmál] \dicTerm{leyndar··mál} \dicsymFrequent\  \dicIPA{{l}{ei}{n}{\textsubring{d}}{a}{r}{m}{au}{\textsubring{l}}} \dicPos{n}[2] \dicFlx{(‑s, ‑)}[5] \dicDirectTranslationCS{tajemství} \dicExampleIS{opinbert leyndarmál} \dicExampleCS{veřejné tajemství}
\dicEntry[leyndur] \dicTerm{leyndur} \dicsymFrequent\  \dicIPA{{l}{ei}{n}{\textsubring{d}}{\textscy}{\textsubring{r}}} \dicPos{adj}[2]\dicFlx{}[-14] \textbf{1.} \dicDirectTranslationCS{skrytý, schovaný} \dicExampleIS{leyndir hæfileikar} \dicExampleCS{skryté schopnosti}  \textbf{2.} \dicDirectTranslationCS{tajný, utajený};  \dicPhraseIS{fara leynt með e‑ð} \dicDirectTranslationCS{tajit (co), držet (co) v~tajnosti};  \dicPhraseIS{halda e‑ð leyndum} \dicDirectTranslationCS{tajit (co), držet (co) v~tajnosti}
\dicEntry[leyni] \dicTerm{leyni} \dicIPA{{l}{ei}{\textlengthmark}{n}{\textsci}} \dicPos{n}[2] \dicFlx{(‑s, ‑)}[14] \dicDirectTranslationCS{skrýš, úkryt};  \dicPhraseIS{í leyni} \dicFlx{adv} \dicDirectTranslationCS{tajně, potají};  \dicPhraseIS{liggja í leyni} \dicDirectTranslationCS{skrývat se, ukrývat se}
\dicEntry[leynifélag] \dicTerm{leyni··fé·|lag} \dicIPA{{l}{ei}{\textlengthmark}{n}{\textsci}{f}{j}{\textepsilon}{l}{a}{x}} \dicPos{n}[2] \dicFlx{(‑lags, ‑lög)}[8] \dicDirectTranslationCS{tajné společenství, tajná společnost}
\dicEntry[leynilegur] \dicTerm{leyni··legur} \dicIPA{{l}{ei}{\textlengthmark}{n}{\textsci}{l}{\textepsilon}{\textbabygamma}{\textscy}{\textsubring{r}}} \dicPos{adj}[1]\dicFlx{}[-8] \dicSynonym{dulinn} \dicDirectTranslationCS{tajný, utajený}
\dicEntry[leynilögregla] \dicTerm{leyni··lög·regl|a} \dicIPA{{l}{ei}{\textlengthmark}{n}{\textsci}{l}{\oe}{\textbabygamma}{r}{\textepsilon}{\r{g}}{l}{a}} \dicPos{f}[1] \dicFlx{(‑u)}[5] \dicDirectTranslationCS{tajná policie}

\dicEntry[leynilögreglumaður] \textls[23]{\dicTerm{leyni··lög·reglu·|maður} \dicIPA{{l}\-{ei}\-{\textlengthmark}\-{n}\-{\textsci}\-{l}\-{\oe}\-{\textbabygamma}\-{r}\-{\textepsilon}\-{\r{g}}\-{l}\-{\textscy}\-{m}\-{a}\-{ð}\-{\textscy}\-{\textsubring{r}}} \dicPos{m}[13] \dicFlx{(‑manns, ‑menn)}[2]} \dicDirectTranslationCS{detektiv, tajný policista, tajná policistka}

\dicEntry[leynilögreglusaga] \dicTerm{leyni·lög·reglu··|saga} \dicIPA{{l}\-{ei}\-{\textlengthmark}\-{n}\-{\textsci}\-{l}\-{\oe}\-{\textbabygamma}\-{r}\-{\textepsilon}\-{\r{g}}\-{l}\-{\textscy}\-{s}\-{a}\-{\textbabygamma}\-{a}\-} \dicPos{f}[1] \dicFlx{(‑sögu, ‑sög\-ur)}[14] \dicDirectTranslationCS{detektivka, detektivní příběh}
\dicEntry[leyniskytta] \dicTerm{leyni··skytt|a} \dicIPA{{l}{ei}{\textlengthmark}{n}{\textsci}{s}{\r{\textObardotlessj}}{\textsci}{h}{\textsubring{d}}{a}} \dicPos{f}[1] \dicFlx{(‑u, ‑ur)}[19] \dicDirectTranslationCS{odstřelovač(ka), snajper(ka)}
\dicEntry[leyniþjónusta] \textls[15]{\dicTerm{leyni··þjón·ust|a} \dicIPA{{l}{ei}{\textlengthmark}{n}{\textsci}{\texttheta}{j}{ou}{n}{\textscy}{s}{\textsubring{d}}{a}} \dicPos{f}[1] \dicFlx{(‑u)}[5] \dicDirectTranslationCS{tajná\,/\addthin vý\-zvěd\-ná služba, rozvědka}}
\dicEntry[leysa] \dicTerm{leys|a} \dicsymFrequent\  \dicIPA{{l}{ei}{\textlengthmark}{s}{a}} \dicPos{v}[2] \dicFlx{(‑ti, ‑t)}[64] \dicFlx{acc} \textbf{1.} \dicSynonym{losa} \dicDirectTranslationCS{rozvázat, rozvazovat, rozmotat, rozmotávat, uvolnit, uvolňovat} \dicExampleIS{leysa bandið} \dicExampleCS{rozvázat provázek}  \textbf{2.} \dicSynonym*{finna lausn} \dicDirectTranslationCS{vyřešit, rozřešit, rozuzlit}  \textbf{3.} \dicPhraseIS{e‑ð leysir} \dicFlx{impers} \dicDirectTranslationCS{(co) taje, (co) roztává} \dicExampleIS{Ísa leysir.} \dicExampleCS{Ledy tají.};  \dicIdiom{leysa}[af]{ \dicPhraseIS{leysa e‑n af}} \dicDirectTranslationCS{zastoupit (koho)}; { \dicPhraseIS{leysa e‑ð af hólmi}} \dicDirectTranslationCS{nahradit (co), zastoupit (co)};  \dicIdiom{leysa}[niður]{ \dicPhraseIS{leysa niður um sig}} \dicDirectTranslationCS{stáhnout si kalhoty};  \dicIdiom{leysa}[undan]{ \dicPhraseIS{leysa e‑n undan e‑u}} \dicDirectTranslationCS{oprostit (koho) od (čeho), osvobodit (koho) od (čeho)};  \dicIdiom{leysa}[upp]{ \dicPhraseIS{leysa e‑ð upp}} \dicDirectTranslationCS{rozpustit\,/\addthin rozpouštět (co) (sůl ve vodě ap.)};  \dicIdiom{leysa}[úr]{ \dicPhraseIS{leysa úr e‑u}} \dicDirectTranslationCS{vyřešit (co), nalézt řešení (čeho)} \dicExampleIS{leysa úr vanda} \dicExampleCS{vyřešit problém};  \dicIdiom{leysa}[út]{ \dicPhraseIS{leysa e‑ð út}} \dicDirectTranslationCS{zaplatit za (co) (za objednané zboží)};  \dicPhraseIS{leysa e‑n út með e‑ð} \dicDirectTranslationCS{obdarovat (koho čím) na odchodu};  \dicIdiom{leysast}{ \dicPhraseIS{leysast}} \dicFlx{refl} {\textbf{a.}} \dicDirectTranslationCS{rozvázat se, rozmotat se};  {\textbf{b.}} \dicDirectTranslationCS{vyřešit se, rozuzlit se} \dicExampleIS{Deilan leystist með samkomulagi.} \dicExampleCS{Spor se vyřešil dohodou.}
\dicEntry[leysanlegur] \dicTerm{leysan··legur} \dicIPA{{l}{ei}{\textlengthmark}{s}{a}{n}{l}{\textepsilon}{\textbabygamma}{\textscy}{\textsubring{r}}} \dicPos{adj}[1]\dicFlx{}[-8] \textbf{1.} \dicFieldCat{chem.} \dicDirectTranslationCS{rozpustný}  \textbf{2.} \dicDirectTranslationCS{(vy)řešitelný}
\dicEntry[leysiefni] \dicTerm{leysi··efni} \dicIPA{{l}{ei}{\textlengthmark}{s}{\textsci}{\textepsilon}{\textsubring{b}}{n}{\textsci}} \dicPos{n}[2] \dicFlx{(‑s, ‑)}[14] \dicFieldCat{chem.} \dicDirectTranslationCS{rozpouštědlo}
\dicEntry[leysigeisli] \dicTerm{leysi··geisl|i} \dicIPA{{l}{ei}{\textlengthmark}{s}{\textsci}{\r{\textObardotlessj}}{ei}{s}{\textsubring{d}}{l}{\textsci}} \dicPos{m}[1] \dicFlx{(‑a, ‑ar)}[1] \dicFieldCat{fyz.} \dicDirectTranslationCS{laserový paprsek}
\dicEntry[leysing] \dicTerm{leys··ing} \dicIPA{{l}{ei}{\textlengthmark}{s}{i}{\ng}{\r{g}}} \dicPos{f}[4] \dicFlx{(‑ar, ‑ar)}[5] \textbf{1.} \dicSynonym{lausn} \dicDirectTranslationCS{(vy)řešení, rozuzlení}  \textbf{2.} \dicSynonym*{snjóbráð} \dicDirectTranslationCS{tání, obleva}  \textbf{3.} \dicSynonym{upplausn} \dicDirectTranslationCS{rozpuštění, rozpouštění (látky ve vodě ap.)} \dicExampleIS{leysing sykurs í vatni} \dicExampleCS{rozpuštění cukru ve vodě}
\dicEntry[leysingi] \dicTerm{leys··ing|i} \dicIPA{{l}{ei}{\textlengthmark}{s}{i}{\textltailn}{\r{\textObardotlessj}}{\textsci}} \dicPos{m}[1] \dicFlx{(‑ja, ‑jar)}[14] \dicFieldCat{hist.} \dicDirectTranslationCS{osvobozený otrok, osvobozená otrokyně}
\dicEntry[leysinn] \dicTerm{leysinn} \dicIPA{{l}{ei}{\textlengthmark}{s}{\textsci}{\textsubring{n}}} \dicPos{adj}[6]\dicFlx{}[-6] \dicDirectTranslationCS{rozpustný}
\dicEntry[leysiprentari] \dicTerm{leysi··prent·ar|i} \dicIPA{{l}{ei}{\textlengthmark}{s}{\textsci}{p\smash{\textsuperscript{h}}}{r}{\textepsilon}{\textsubring{n}}{\textsubring{d}}{a}{r}{\textsci}} \dicPos{m}[1] \dicFlx{(‑a, ‑ar)}[13] \dicFieldCat{poč.} \dicDirectTranslationCS{laserová tiskárna}
\dicEntry[leysir] \dicTerm{leys|ir} \dicIPA{{l}{ei}{\textlengthmark}{s}{\textsci}{\textsubring{r}}} \dicPos{m}[7] \dicFlx{(‑is, ‑ar)}[1] \dicFieldCat{fyz.} \dicDirectTranslationCS{laser}
\dicEntry[leysni] \dicTerm{leysn|i} \dicIPA{{l}{ei}{s}{\textsubring{d}}{n}{\textsci}} \dicPos{f}[3] \dicFlx{(‑i)}[3] \dicFieldCat{chem.} \dicDirectTranslationCS{rozpustnost (materiálu ap.)}
\dicEntry[leyti] \dicTerm{leyti} \dicsymFrequent\  \dicIPA{{l}{ei}{\textlengthmark}{\textsubring{d}}{\textsci}} \dicPos{n}[2] \dicFlx{(‑s)}[20] \textbf{1.} \dicSynonym{tillit} \dicDirectTranslationCS{ohled, zřetel, míra};  \dicPhraseIS{að mestu leyti} \dicFlx{adv} \dicDirectTranslationCS{převážně, většinou};  \dicPhraseIS{að miklu leyti} \dicFlx{adv} \dicDirectTranslationCS{do velké míry, z~velké části};  \dicPhraseIS{að mörgu leyti} \dicFlx{adv} \dicDirectTranslationCS{do značné míry};  \dicPhraseIS{að nokkru leyti} \dicFlx{adv} \dicDirectTranslationCS{do určité míry, částečně};  \dicPhraseIS{að vissu leyti} \dicFlx{adv} \dicDirectTranslationCS{do jisté míry};  \dicPhraseIS{að öðru leyti} \dicFlx{adv} \dicDirectTranslationCS{jinak};  \dicPhraseIS{að þessu leyti} \dicFlx{adv} \dicDirectTranslationCS{v~tomto ohledu\,/\addthin směru}  \textbf{2.} \dicSynonym{tími} \dicDirectTranslationCS{doba, čas} \dicExampleIS{um sjöleytið} \dicExampleCS{kolem sedmé hodiny};  \dicPhraseIS{um hvaða leyti} \dicDirectTranslationCS{v~kolik (hodin)};  \dicPhraseIS{um sama leyti} \dicFlx{adv} \dicDirectTranslationCS{(přibližně) ve stejný čas};  \dicPhraseIS{um það leyti} \dicFlx{adv} \dicDirectTranslationCS{tehdy, tenkrát}
\dicEntry[léð] \dicTerm{léð} \dicIPA{{l}{j}{\textepsilon}{\textlengthmark}{\texttheta}} \dicPos{v} \dicFlx{supin} \dicLink{ljá}
\dicEntry[léði] \dicTerm{léði} \dicIPA{{l}{j}{\textepsilon}{\textlengthmark}{ð}{\textsci}} \dicPos{v} \dicFlx{ind\,/\addthin con pf sg 1 pers} \dicLink{ljá}
\dicEntry[léðum] \dicTerm{léðum} \dicIPA{{l}{j}{\textepsilon}{\textlengthmark}{ð}{\textscy}{\textsubring{m}}} \dicPos{v} \dicFlx{ind pf pl 1 pers} \dicLink{ljá}
\dicEntry[lék] \dicTerm{lék} \dicIPA{{l}{j}{\textepsilon}{\textlengthmark}{\r{g}}} \dicPos{v} \dicFlx{ind pf sg 1 pers} \dicLink{leika}
\dicEntry[léki] \dicTerm{léki} \dicIPA{{l}{j}{\textepsilon}{\textlengthmark}{\r{\textObardotlessj}}{\textsci}} \dicPos{v} \dicFlx{con pf sg 1 pers} \dicLink{leika}
\dicEntry[lékum] \dicTerm{lékum} \dicIPA{{l}{j}{\textepsilon}{\textlengthmark}{\r{g}}{\textscy}{\textsubring{m}}} \dicPos{v} \dicFlx{ind pf pl 1 pers} \dicLink{leika}
\dicEntry[lélega] \dicTerm{lé··lega} \dicIPA{{l}{j}{\textepsilon}{\textlengthmark}{l}{\textepsilon}{\textbabygamma}{a}} \dicPos{adv} \dicDirectTranslationCS{špatně, bídně, mizerně}
\dicEntry[lélegur] \dicTerm{lé··legur} \dicsymFrequent\  \dicIPA{{l}{j}{\textepsilon}{\textlengthmark}{l}{\textepsilon}{\textbabygamma}{\textscy}{\textsubring{r}}} \dicPos{adj}[1]\dicFlx{}[-8] \dicSynonym{slæmur} \dicDirectTranslationCS{špatný, bídný, mizerný, chabý} \dicExampleIS{benda á lélegan árangur} \dicExampleCS{upozornit na špatný výsledek}
\dicEntry[lén] \dicTerm{lén} \dicIPA{{l}{j}{\textepsilon}{\textlengthmark}{\textsubring{n}}} \dicPos{n}[2] \dicFlx{(‑s, ‑)}[5] \textbf{1.} \dicSynonym{ríki} \dicDirectTranslationCS{léno}  \textbf{2.} \dicFieldCat{poč.} \dicSynonym{umdæmi} \dicDirectTranslationCS{doména}  \textbf{3.} \dicFieldCat{biol.} \dicDirectTranslationCS{doména}
\dicEntry[lénskur] \dicTerm{lénskur} \dicIPA{{l}{j}{\textepsilon}{n}{s}{\r{g}}{\textscy}{\textsubring{r}}} \dicPos{adj}[1]\dicFlx{}[-1] \dicDirectTranslationCS{feudální}
\dicEntry[lénsmaður] \dicTerm{léns··|maður} \dicIPA{{l}{j}{\textepsilon}{n}{s}{m}{a}{ð}{\textscy}{\textsubring{r}}} \dicPos{m}[13] \dicFlx{(‑manns, ‑menn)}[2] \dicDirectTranslationCS{leník, vazal(ka)}
\dicEntry[léreft] \dicTerm{léreft} \dicIPA{{l}{j}{\textepsilon}{\textlengthmark}{r}{\textepsilon}{f}{\textsubring{d}}} \dicPos{n}[2] \dicFlx{(‑s, ‑)}[5] \dicDirectTranslationCS{(lněné) plátno}
\dicEntry[lét] \dicTerm{lét} \dicIPA{{l}{j}{\textepsilon}{\textlengthmark}{\textsubring{d}}} \dicPos{v} \dicFlx{ind pf sg 1 pers} \dicLink{láta}
\dicEntry[léti] \dicTerm{léti} \dicIPA{{l}{j}{\textepsilon}{\textlengthmark}{\textsubring{d}}{\textsci}} \dicPos{v} \dicFlx{con pf sg 1 pers} \dicLink{láta}
\dicEntry[létt] \dicTerm{létt} \dicsymFrequent\  \dicIPA{{l}{j}{\textepsilon}{h}{\textsubring{d}}} \dicPos{adv} \dicFlx{(comp ‑ar, sup ‑ast)} \dicDirectTranslationCS{lehce, mírně, jemně, zlehka} \dicExampleIS{Hann klappaði létt á bakið á mér.} \dicExampleCS{Lehce mě poklepal po zádech.}
\dicEntry[létta] \dicTerm{létt|a} \dicsymFrequent\  \dicIPA{{l}{j}{\textepsilon}{h}{\textsubring{d}}{a}} \dicPos{v}[2] \dicFlx{(‑i, ‑)}[2] \dicFlx{acc\,/\addthin dat} \textbf{1.} \dicDirectTranslationCS{odlehčit, odlehčovat, ulehčit, ulehčovat, (u)činit lehčím} \dicExampleIS{létta byrðina} \dicExampleCS{odlehčit břemeno}  \textbf{2.} \dicPhraseIS{e‑m léttir} \dicFlx{impers} \dicDirectTranslationCS{(komu) se ulehčuje, (komu) se ulevuje}  \textbf{3.} \dicPhraseIS{e‑u léttir} \dicFlx{impers} \dicDirectTranslationCS{(čeho) ubývá (mlhy ap.), (co) mizí};  \dicIdiom{létta}[af]{ \dicPhraseIS{létta af e‑m e‑u}} \dicDirectTranslationCS{odlehčit (komu) od (čeho)}; { \dicPhraseIS{það léttir af e‑m e‑u}} \dicFlx{impers} \dicDirectTranslationCS{(komu) se ulevuje od (čeho)};  \dicIdiom{létta}[á]{ \dicPhraseIS{létta á sér}} \dicDirectTranslationCS{ulevit si} \dicIndirectTranslationCS{(vykonat tělesnou potřebu)};  \dicIdiom{létta}[sér]{ \dicPhraseIS{létta sér lífið}} \dicDirectTranslationCS{pobavit se, povyrazit si};  \dicIdiom{létta}[til]{ \dicPhraseIS{létta til}} \dicDirectTranslationCS{vyčasit se, vyjasnit se, vyjasňovat se} \dicExampleIS{Það er að létta til.} \dicExampleCS{Začíná se vyjasňovat.};  \dicIdiom{létta}[undir]{ \dicPhraseIS{létta undir með e‑m}} \dicDirectTranslationCS{ulehčit (komu), učinit to (komu) lehčím};  \dicIdiom{léttast}{ \dicPhraseIS{léttast}} \dicFlx{refl} \dicDirectTranslationCS{(z)hubnout};  \dicIdiom{létta}{ \dicPhraseIS{létta akkerum}} \dicFieldCat{nám.} \dicDirectTranslationCS{zvednout kotvu}
\dicEntry[léttasótt] \dicTerm{létta··sótt} \dicIPA{{l}{j}{\textepsilon}{h}{\textsubring{d}}{a}{s}{ou}{h}{\textsubring{d}}} \dicPos{f}[7] \dicFlx{(‑ar)}[3] \dicDirectTranslationCS{porodní bolesti, kontrakce} \dicExampleIS{taka léttasóttina} \dicExampleCS{mít porodní bolesti}
\dicEntry[létteind] \dicTerm{létt··eind} \dicIPA{{l}{j}{\textepsilon}{h}{\textsubring{d}}{ei}{n}{\textsubring{d}}} \dicPos{f}[7] \dicFlx{(‑ar, ‑ir)}[1] \dicFieldCat{fyz.} \dicDirectTranslationCS{lepton}
\dicEntry[létti] \dicTerm{létt|i} \dicIPA{{l}{j}{\textepsilon}{h}{\textsubring{d}}{\textsci}} \dicPos{m}[1] \dicFlx{(‑a, ‑ar)}[1] \textbf{1.} \dicSynonym{léttir} \dicDirectTranslationCS{ulehčení, úleva}  \textbf{2.} \dicSynonym{axlaband} \dicDirectTranslationCS{popruh, šle}  \textbf{3.} \dicSynonym*{fúsleiki} \dicDirectTranslationCS{ochota};  \dicPhraseIS{segja allt af létta} \dicDirectTranslationCS{ochotně všechno povědět}
\dicEntry[léttir] \dicTerm{létt|ir} \dicIPA{{l}{j}{\textepsilon}{h}{\textsubring{d}}{\textsci}{\textsubring{r}}} \dicPos{m}[7] \dicFlx{(‑is, ‑ar)}[1] \textbf{1.} \dicSynonym{huggun} \dicDirectTranslationCS{útěcha, úleva}  \textbf{2.} \dicSynonym{hjálp} \dicDirectTranslationCS{pomoc};  \dicPhraseIS{vera e‑m til léttis} \dicDirectTranslationCS{pomoci (komu), ulehčit (komu)}  \textbf{3.} \dicFieldCat{zool.} \dicDirectTranslationCS{delfín obecný} \textit{(l.~{\textLA{Delphinus delphis}})}  \dicsymPhoto\ 
\dicFigure{26462.jpg}{Léttir}{Léttir - Shane Anderson, Biolib, PD}
\dicEntry[léttklæddur] \dicTerm{létt··klæddur} \dicIPA{{l}{j}{\textepsilon}{h}{\textsubring{d}}{k\smash{\textsuperscript{h}}}{l}{a}{i}{\textsubring{d}}{\textscy}{\textsubring{r}}} \dicPos{adj}[2]\dicFlx{}[-18] \dicDirectTranslationCS{lehce oblečený}
\dicEntry[léttleiki] \dicTerm{létt··leik|i} \dicIPA{{l}{j}{\textepsilon}{h}{\textsubring{d}}{l}{ei}{\r{\textObardotlessj}}{\textsci}} \dicPos{m}[1] \dicFlx{(‑a)}[3] \dicDirectTranslationCS{lehkost} \dicExampleIS{óbærilegur léttleiki tilverunnar} \dicExampleCS{nesnesitelná lehkost bytí}
\dicEntry[léttlyndi] \dicTerm{létt··lyndi} \dicIPA{{l}{j}{\textepsilon}{h}{\textsubring{d}}{l}{\textsci}{n}{\textsubring{d}}{\textsci}} \dicPos{n}[2] \dicFlx{(‑s)}[20] \dicDirectTranslationCS{dobrá nálada, uvolněnost, bezstarostnost}
\dicEntry[léttlyndur] \dicTerm{létt··lyndur} \dicIPA{{l}{j}{\textepsilon}{h}{\textsubring{d}}{l}{\textsci}{n}{\textsubring{d}}{\textscy}{\textsubring{r}}} \dicPos{adj}[2]\dicFlx{}[-14] \dicDirectTranslationCS{uvolněný, bezstarostný}
\dicEntry[léttmjólk] \dicTerm{létt··mjólk} \dicIPA{{l}{j}{\textepsilon}{h}{\textsubring{d}}{m}{j}{ou}{\textsubring{l}}{\r{g}}} \dicPos{f}[10] \dicFlx{(‑ur)}[2] \dicDirectTranslationCS{nízkotučné mléko}
\dicEntry[léttskýjaður] \dicTerm{létt··skýj·|aður} \dicIPA{{l}{j}{\textepsilon}{h}{\textsubring{d}}{s}{\r{\textObardotlessj}}{i}{j}{a}{ð}{\textscy}{\textsubring{r}}} \dicPos{adj}[3] \dicFlx{(f ‑uð)}[3] \dicPhraseIS{það er léttskýjað} \dicDirectTranslationCS{je polojasno}
\dicEntry[léttur] \dicTerm{léttur} \dicsymFrequent\  \dicIPA{{l}{j}{\textepsilon}{h}{\textsubring{d}}{\textscy}{\textsubring{r}}} \dicPos{adj}[1]\dicFlx{}[-10] \textbf{1.} \dicDirectTranslationCS{lehký, snadný} \dicExampleIS{létt verk} \dicExampleCS{lehká práce} \dicAntonym{erfiður}  \textbf{2.} \dicDirectTranslationCS{lehký} \dicIndirectTranslationCS{(vážící málo)} \dicAntonym{þungur}  \textbf{3.} \dicSynonym{glaður} \dicDirectTranslationCS{uvolněný, klidný, bezstarostný} \dicExampleIS{Hún er létt í lund.} \dicExampleCS{Je bezstarostná.}  \textbf{4.} \dicDirectTranslationCS{dietní, lehký (strava ap.)} \dicExampleIS{léttar veitingar} \dicExampleCS{lehké občerstvení}  \textbf{5.} \dicDirectTranslationCS{lehký, slabý (alkohol ap.)} \dicExampleIS{létt vín} \dicExampleCS{slabý alkohol}
\dicEntry[léttúð] \dicTerm{létt··úð} \dicIPA{{l}{j}{\textepsilon}{h}{\textsubring{d}}{u}{\texttheta}} \dicPos{f}[4] \dicFlx{(‑ar)}[3] \textbf{1.} \dicSynonym{alvöruleysi} \dicDirectTranslationCS{lehkovážnost}  \textbf{2.} \dicSynonym{lauslæti} \dicDirectTranslationCS{rozpustilost, frivolnost}
\dicEntry[léttúðugur] \dicTerm{létt·úð··ugur} \dicIPA{{l}{j}{\textepsilon}{h}{\textsubring{d}}{u}{ð}{\textscy}{\textbabygamma}{\textscy}{\textsubring{r}}} \dicPos{adj}[1]\dicFlx{}[-8] \dicSynonym{lauslátur} \dicDirectTranslationCS{frivolní, prostopášný}
\dicEntry[léttvín] \dicTerm{létt··vín} \dicIPA{{l}{j}{\textepsilon}{h}{\textsubring{d}}{v}{i}{\textsubring{n}}} \dicPos{n}[2] \dicFlx{(‑s, ‑)}[5] \dicDirectTranslationCS{víno}
\dicEntry[léttvægur] \dicTerm{létt··vægur} \dicIPA{{l}{j}{\textepsilon}{h}{\textsubring{d}}{v}{a}{i}{\textbabygamma}{\textscy}{\textsubring{r}}} \dicPos{adj}[1]\dicFlx{}[-1] \dicDirectTranslationCS{nedůležitý, nepodstat\-ný, nevýznamný}
\dicEntry[létum] \dicTerm{létum} \dicIPA{{l}{j}{\textepsilon}{\textlengthmark}{\textsubring{d}}{\textscy}{\textsubring{m}}} \dicPos{v} \dicFlx{ind pf pl 1 pers} \dicLink{láta}
\dicEntry[lh. nt.] \dicTerm{lh. nt.} \dicPos{zkr} \dicPhraseIS{lýsingarháttur nútíðar} \dicFieldCat{jaz.} \dicDirectTranslationCS{příčestí přítomné}
\dicEntry[lh. þt.] \dicTerm{lh. þt.} \dicPos{zkr} \dicPhraseIS{lýsingarháttur þátíðar} \dicFieldCat{jaz.} \dicDirectTranslationCS{příčestí minulé}
\dicEntry[lið] \dicTerm{lið} \dicsymFrequent\  \dicIPA{{l}{\textsci}{\textlengthmark}{\texttheta}} \dicPos{n}[2] \dicFlx{(‑s, ‑)}[5] \textbf{1.} \dicSynonym{hópur} \dicDirectTranslationCS{družstvo, oddíl, četa};  \dicPhraseIS{ganga í lið við e‑n} \dicDirectTranslationCS{spojit síly s~(kým)};  \dicPhraseIS{safna liði} \dicDirectTranslationCS{shromažďovat podporu}  \textbf{2.} \dicFieldCat{sport.} \dicDirectTranslationCS{družina, tým, mužstvo}  \textbf{3.} \dicSynonym{aðstoð} \dicDirectTranslationCS{výpomoc, pomocná ruka};  \dicPhraseIS{koma til liðs við e‑n} \dicDirectTranslationCS{vypomoct (komu), přijít (komu) na pomoc};  \dicPhraseIS{verða e‑m að liði} \dicDirectTranslationCS{vypomoct (komu), asistovat (komu)}  \textbf{4.} \dicLangCat{hovor.} \dicSynonym{fólk} \dicDirectTranslationCS{lidi}
\dicEntry[liða] \dicTerm{lið|a} \dicIPA{{l}{\textsci}{\textlengthmark}{ð}{a}} \dicPos{v}[1] \dicFlx{(‑aði)}[1] \dicFlx{acc} \textbf{1.} \dicSynonym{skipta} \dicDirectTranslationCS{(roz)dělit (do skupin)}  \textbf{2.} \dicSynonym*{gera hrokkið} \dicDirectTranslationCS{natočit, natáčet, (na)kadeřit (vlasy ap.)} \dicExampleIS{liða hár} \dicExampleCS{natáčet vlasy};  \dicIdiom{liðast}{ \dicPhraseIS{liðast}} \dicFlx{refl} \dicDirectTranslationCS{vlnit se, kroutit se, kadeřit se}
\dicEntry[liðaður] \dicTerm{lið··|aður} \dicIPA{{l}{\textsci}{\textlengthmark}{ð}{a}{ð}{\textscy}{\textsubring{r}}} \dicPos{adj}[3] \dicFlx{(f ‑uð)}[3] \dicSynonym*{bylgjaður} \dicDirectTranslationCS{natočený, nakadeřený, navlněný}
\dicEntry[liðagigt] \dicTerm{liða··gigt} \dicIPA{{l}{\textsci}{\textlengthmark}{ð}{a}{\r{\textObardotlessj}}{\textsci}{x}{\textsubring{d}}} \dicPos{f}[4] \dicFlx{(‑ar)}[3] \dicFieldCat{med.} \dicDirectTranslationCS{revmatoidní artritida}
\dicEntry[liðamót] \dicTerm{liða··mót} \dicIPA{{l}{\textsci}{\textlengthmark}{ð}{a}{m}{ou}{\textsubring{d}}} \dicPos{n}[2] \dicFlx{pl}[1] \dicFieldCat{anat.} \dicDirectTranslationCS{kloub}
\dicEntry[liðband] \dicTerm{lið··|band} \dicIPA{{l}{\textsci}{ð}{\textsubring{b}}{a}{n}{\textsubring{d}}} \dicPos{n}[2] \dicFlx{(‑bands, ‑bönd)}[8] \dicFieldCat{anat.} \dicDirectTranslationCS{vazivo}
\dicEntry[liðdýr] \dicTerm{lið··dýr} \dicIPA{{l}{\textsci}{ð}{\textsubring{d}}{i}{\textsubring{r}}} \dicPos{n}[2] \dicFlx{(‑s, ‑)}[5] \dicFieldCat{zool.} \dicDirectTranslationCS{členovec} \textit{(l.~{\textLA{Arthropoda}})}
\dicEntry[liðfall] \dicTerm{lið··|fall} \dicIPA{{l}{\textsci}{ð}{f}{a}{\textsubring{d}}{\textsubring{l}}} \dicPos{n}[2] \dicFlx{(‑falls, ‑föll)}[8] \dicFieldCat{jaz.} \dicSynonym{brottfall} \dicDirectTranslationCS{elipsa, výpustka}
\dicEntry[liðfelldur] \dicTerm{lið··felldur} \dicIPA{{l}{\textsci}{ð}{f}{\textepsilon}{l}{\textsubring{d}}{\textscy}{\textsubring{r}}} \dicPos{adj}[2]\dicFlx{}[-14] \dicFieldCat{jaz.} \dicDirectTranslationCS{eliptický} \dicExampleIS{liðfellt orðtak} \dicExampleCS{eliptický výraz}
\dicEntry[liðhlaupi] \dicTerm{lið··hlaup|i} \dicIPA{{l}{\textsci}{ð}{\textsubring{l}}{\oe i}{\textsubring{b}}{\textsci}} \dicPos{m}[1] \dicFlx{(‑a, ‑ar)}[1] \dicFieldCat{voj.} \dicDirectTranslationCS{zběh, dezertér(ka)}
\dicEntry[liði] \dicTerm{lið|i\smash{\textsuperscript{1}}} \dicIPA{{l}{\textsci}{\textlengthmark}{ð}{\textsci}} \dicPos{m}[1] \dicFlx{(‑a, ‑ar)}[1] \dicDirectTranslationCS{člen(ka) týmu (vojenského, zdravotního ap.)} \dicIndirectTranslationCS{(vyskytuje se jako druhá část složeného slova)} \dicExampleIS{skæruliði} \dicExampleCS{partyzán}
\dicEntry[liði] \dicTerm{liði\smash{\textsuperscript{2}}} \dicIPA{{l}{\textsci}{\textlengthmark}{ð}{\textsci}} \dicPos{v} \dicFlx{con pf sg 1 pers} \dicLink{líða}
\dicEntry[liðið] \dicTerm{liðið} \dicIPA{{l}{\textsci}{\textlengthmark}{ð}{\textsci}{\texttheta}} \dicPos{v} \dicFlx{supin} \dicLink{líða}
\dicEntry[liðinn] \dicTerm{liðinn} \dicsymFrequent\  \dicIPA{{l}{\textsci}{\textlengthmark}{ð}{\textsci}{\textsubring{n}}} \dicPos{adj}[6]\dicFlx{}[-6] \textbf{1.} \dicDirectTranslationCS{minulý, uplynulý} \dicExampleIS{í liðnum mánuði} \dicExampleCS{v~minulém měsíci}  \textbf{2.} \dicDirectTranslationCS{zesnulý};  \dicPhraseIS{vera vel liðinn} \dicDirectTranslationCS{být populární\,/\addthin oblíbený}
\dicEntry[liðka] \dicTerm{liðk|a} \dicIPA{{l}{\textsci}{\texttheta}{\r{g}}{a}} \dicPos{v}[1] \dicFlx{(‑aði)}[1] \dicFlx{acc} \dicDirectTranslationCS{rozhýbat, rozpohybovat} \dicExampleIS{liðka handlegginn} \dicExampleCS{rozhýbat ruku};  \dicPhraseIS{liðka sig} \dicDirectTranslationCS{rozhýbat se};  \dicIdiom{liðka}[til]{ \dicPhraseIS{liðka e‑ð til}} \dicDirectTranslationCS{usnadnit (co)};  \dicIdiom{liðkast}{ \dicPhraseIS{liðkast}} \dicFlx{refl} \dicDirectTranslationCS{rozhýbat se}
\dicEntry[liðlangur] \dicTerm{lið··|langur} \dicIPA{{l}{\textsci}{ð}{l}{au}{\ng}{\r{g}}{\textscy}{\textsubring{r}}} \dicPos{adj}[10] \dicFlx{(f ‑löng, comp ‑lengri, sup ‑leng\-st\-ur)}[2] \dicPhraseIS{allan liðlangan daginn} \dicFlx{adv} \dicDirectTranslationCS{(po) celý celičký den}
\dicEntry[liðlegur] \dicTerm{lið··legur} \dicIPA{{l}{\textsci}{ð}{l}{\textepsilon}{\textbabygamma}{\textscy}{\textsubring{r}}} \dicPos{adj}[1]\dicFlx{}[-8] \textbf{1.} \dicSynonym{fimur} \dicDirectTranslationCS{obratný}  \textbf{2.} \dicSynonym{hjálpsamur} \dicDirectTranslationCS{vstřícný, úslužný, ochotný pomoci}
\dicEntry[liðleiki] \dicTerm{lið··leik|i} \dicIPA{{l}{\textsci}{ð}{l}{ei}{\r{\textObardotlessj}}{\textsci}} \dicPos{m}[1] \dicFlx{(‑a)}[3] \textbf{1.} \dicDirectTranslationCS{pohyblivost, obratnost}  \textbf{2.} \dicDirectTranslationCS{vstřícnost, ochota pomoci}
\dicEntry[liðsauki] \dicTerm{liðs··auk|i} \dicIPA{{l}{\textsci}{ð}{s}{\oe i}{\r{\textObardotlessj}}{\textsci}} \dicPos{m}[1] \dicFlx{(‑a)}[3] \dicDirectTranslationCS{posila, posílení\,/\addthin rozšíření řad}
\dicEntry[liðsemd] \textls[30]{\dicTerm{lið··semd} \dicIPA{{l}{\textsci}{ð}{s}{\textepsilon}{m}{\textsubring{d}}} \dicPos{f}[7] \dicFlx{(‑ar)}[3] \dicSynonym{aðstoð} \dicDirectTranslationCS{výpomoc, asis\nobreak tence}}
\dicEntry[liðsforingi] \dicTerm{liðs··for·ing|i} \dicIPA{{l}{\textsci}{ð}{s}{f}{\textopeno}{r}{i}{\textltailn}{\r{\textObardotlessj}}{\textsci}} \dicPos{m}[1] \dicFlx{(‑ja, ‑jar)}[14] \dicFieldCat{voj.} \dicDirectTranslationCS{poručík, poručice}
\dicEntry[liðsinna] \dicTerm{lið··sinn|a} \dicIPA{{l}{\textsci}{ð}{s}{\textsci}{n}{a}} \dicPos{v}[2] \dicFlx{(‑ti, ‑t)}[78] \dicFlx{dat} \dicSynonym{hjálpa} \dicDirectTranslationCS{asistovat, (vy)pomáhat} \dicExampleIS{liðsinna e‑m við e‑ð} \dicExampleCS{asistovat (komu) při (čem)}
\dicEntry[liðsinni] \dicTerm{lið··sinni} \dicIPA{{l}{\textsci}{ð}{s}{\textsci}{n}{\textsci}} \dicPos{n}[2] \dicFlx{(‑s)}[20] \dicSynonym{aðstoð} \dicDirectTranslationCS{asistence, (vý)pomoc}
\dicEntry[liðsmaður] \dicTerm{liðs··|maður} \dicIPA{{l}{\textsci}{ð}{s}{m}{a}{ð}{\textscy}{\textsubring{r}}} \dicPos{m}[13] \dicFlx{(‑manns, ‑menn)}[2] \textbf{1.} \dicSynonym{hermaður} \dicDirectTranslationCS{voják(yně)}  \textbf{2.} \dicSynonym{stuðningsmaður} \dicDirectTranslationCS{přívrženec, přívrženkyně, stoupenec, stoupenkyně}  \textbf{3.} \dicDirectTranslationCS{(spolu)hráč(ka), člen(ka) družiny\,/\addthin týmu}
\dicEntry[liðtækur] \dicTerm{lið··tækur} \dicIPA{{l}{\textsci}{ð}{t\smash{\textsuperscript{h}}}{a}{i}{\r{g}}{\textscy}{\textsubring{r}}} \dicPos{adj}[1]\dicFlx{}[-6] \dicSynonym{hæfur} \dicDirectTranslationCS{kompetentní, způsobilý}
\dicEntry[liðugur] \dicTerm{lið··ugur} \dicIPA{{l}{\textsci}{\textlengthmark}{ð}{\textscy}{\textbabygamma}{\textscy}{\textsubring{r}}} \dicPos{adj}[1]\dicFlx{}[-8] \dicSynonym{fimur} \dicDirectTranslationCS{hbitý, mrštný} \dicExampleIS{vera liðugur í hreyfingum} \dicExampleCS{pohybovat se hbitě};  \dicPhraseIS{vera laus og liðugur} \dicDirectTranslationCS{být volný, užívat si svobody}
\dicEntry[liðum] \dicTerm{liðum} \dicIPA{{l}{\textsci}{\textlengthmark}{ð}{\textscy}{\textsubring{m}}} \dicPos{v} \dicFlx{ind pf pl 1 pers} \dicLink{líða}
\dicEntry[liður] \dicTerm{lið|ur} \dicsymFrequent\  \dicIPA{{l}{\textsci}{\textlengthmark}{ð}{\textscy}{\textsubring{r}}} \dicPos{m}[10] \dicFlx{(‑ar\,/\addthin ‑s, ‑ir)}[37] \textbf{1.} \dicFieldCat{anat.} \dicSynonym{liðamót} \dicDirectTranslationCS{kloub} \dicExampleIS{axlarliður} \dicExampleCS{ramenní kloub};  \dicPhraseIS{geta hvorki hrært legg né lið} \dicLangCat{přen.} \dicDirectTranslationCS{nemoci se vůbec pohnout}  \textbf{2.} \dicSynonym{bylgja} \dicDirectTranslationCS{kudrna, vlnka (ve vlasech)} \dicExampleIS{liður í hári} \dicExampleCS{kudrna ve vlasech}  \textbf{3.} \dicSynonym{atriði} \dicDirectTranslationCS{položka, bod} \dicExampleIS{Þessi liður er rangur.} \dicExampleCS{Tato položka je špatně.}  \textbf{4.} \dicSynonym{ættliður} \dicDirectTranslationCS{koleno, pokolení, generace}  \textbf{5.} \dicFieldCat{mat.} \dicDirectTranslationCS{člen}
\dicEntry[liðveisla] \dicTerm{lið··veisl|a} \dicIPA{{l}{\textsci}{ð}{v}{ei}{s}{\textsubring{d}}{l}{a}} \dicPos{f}[1] \dicFlx{(‑u, ‑ur)}[13] \dicDirectTranslationCS{asistence, (vý)pomoc}
\dicEntry[liðþjálfi] \dicTerm{lið··þjálf|i} \dicIPA{{l}{\textsci}{ð}{\texttheta}{j}{au}{l}{v}{\textsci}} \dicPos{m}[1] \dicFlx{(‑a, ‑ar)}[1] \dicFieldCat{voj.} \dicDirectTranslationCS{seržant(ka)}
\dicEntry[Liechtenstein] \dicTerm{Liechtenstein} \dicIPA{{l}{i}{x}{\textsubring{d}}{\textepsilon}{n}{s}{\textsubring{d}}{a}{i}{\textsubring{n}}} \dicPos{n}[4] \dicFlx{indecl}[2] \dicFieldCat{geog.} \dicDirectTranslationCS{Lichtenštejnsko}
\dicEntry[Liechtensteini] \dicTerm{Liechtenstein|i} \dicIPA{{l}{i}{x}{\textsubring{d}}{\textepsilon}{n}{s}{\textsubring{d}}{a}{i}{n}{\textsci}} \dicPos{m}[1] \dicFlx{(‑a, ‑ar)}[1] \dicDirectTranslationCS{Lichtenštejňan(ka), Lichtenštejnec, Lichtenštejnka}
\dicEntry[liechtensteinskur] \dicTerm{liechtensteinskur} \dicIPA{{l}{i}{x}{\textsubring{d}}{\textepsilon}{n}{s}{\textsubring{d}}{a}{i}{n}{s}{\r{g}}{\textscy}{\textsubring{r}}} \dicPos{adj}[1]\dicFlx{}[-6] \dicDirectTranslationCS{lichtenštejnský}
\dicEntry[lifa] \dicTerm{lif|a} \dicsymFrequent\  \dicIPA{{l}{\textsci}{\textlengthmark}{v}{a}} \dicPos{v}[2] \dicFlx{(‑ði, ‑að)}[124] \dicFlx{acc} \textbf{1.} \dicDirectTranslationCS{žít} \dicExampleIS{lifa langa ævi} \dicExampleCS{žít dlouhý život}  \textbf{2.} \dicDirectTranslationCS{existovat}  \textbf{3.} \dicDirectTranslationCS{zažít, prožít} \dicExampleIS{lifa þessa atburði} \dicExampleCS{zažít tyto události}  \textbf{4.} \dicSynonym*{vera logandi} \dicDirectTranslationCS{planout, svítit};  \dicIdiom{lifa}[af]{ \dicPhraseIS{lifa af}} \dicSynonym{endast} \dicDirectTranslationCS{přežít}; { \dicPhraseIS{lifa af e‑u}} \dicDirectTranslationCS{vyžít z~(čeho)};  \dicIdiom{lifa}[á]{ \dicPhraseIS{lifa á e‑u}} \dicDirectTranslationCS{živit se (čím)};  \dicIdiom{lifa}[eftir]{ \dicPhraseIS{lifa eftir e‑u}} \dicDirectTranslationCS{žít podle (čeho), žít v~souladu s~(čím)};  \dicIdiom{lifa}[fyrir]{ \dicPhraseIS{lifa fyrir e‑n}} \dicDirectTranslationCS{žít pro (koho)};  \dicIdiom{lifa}[inn í]{ \dicPhraseIS{lifa sig inn í e‑ð}} \dicDirectTranslationCS{vžít se do (čeho)} \dicExampleIS{lifa sig inn í starfið} \dicExampleCS{vžít se do práce};  \dicIdiom{lifa}[við]{ \dicPhraseIS{lifa við e‑ð}} \dicDirectTranslationCS{žít s~(čím) (pocitem viny ap.)}
\dicEntry[lifandi] \dicTerm{lif··|andi\smash{\textsuperscript{1}}} \dicIPA{{l}{\textsci}{\textlengthmark}{v}{a}{n}{\textsubring{d}}{\textsci}} \dicPos{m}[2] \dicFlx{(‑anda, ‑endur)}[1] \dicDirectTranslationCS{živý člověk, živá bytost} \dicIndirectTranslationCS{(zvláště v~pl)}
\dicEntry[lifandi] \dicTerm{lif··andi\smash{\textsuperscript{2}}} \dicsymFrequent\  \dicIPA{{l}{\textsci}{\textlengthmark}{v}{a}{n}{\textsubring{d}}{\textsci}} \dicPos{adj}[13] \dicFlx{indecl}[1] \textbf{1.} \dicSynonym*{á lífi} \dicDirectTranslationCS{živý, (jsoucí) naživu} \dicExampleIS{ná birninum lifandi} \dicExampleCS{chytit medvěda živého}  \textbf{2.} \dicSynonym{líflegur} \dicDirectTranslationCS{živý, čilý};  \dicIdiom{lifandi}{ \dicPhraseIS{ekki nokkur\,/\addthin enginn lifandi maður}} \dicLangCat{přen.} \dicDirectTranslationCS{ani živáček};  \dicPhraseIS{lifandi tónlist} \dicDirectTranslationCS{živá hudba}
\dicEntry[lifna] \dicTerm{lifn|a} \dicsymFrequent\  \dicIPA{{l}{\textsci}{\textsubring{b}}{n}{a}} \dicPos{v}[1] \dicFlx{(‑aði)}[1] \textbf{1.} \dicDirectTranslationCS{oživit se, rozhořet se}  \textbf{2.} \dicPhraseIS{lifna við} \dicSynonym*{vakna til lífs} \dicDirectTranslationCS{oživnout, ožít} \dicExampleIS{Fjöllin lifnuðu við eftir fimmtíu ár.} \dicExampleCS{Hory ožily po padesáti letech.}  \textbf{3.} \dicPhraseIS{það lifnar yfir e‑m} \dicFlx{impers} \dicLangCat{přen.} \dicDirectTranslationCS{(kdo) ožívá, (kdo) se stává živějším}
\dicEntry[lifnaður] \dicTerm{lif··nað|ur} \dicIPA{{l}{\textsci}{\textsubring{b}}{n}{a}{ð}{\textscy}{\textsubring{r}}} \dicPos{m}[10] \dicFlx{(‑ar)}[7] \dicSynonym{líferni} \dicDirectTranslationCS{způsob života, životní styl}
\dicEntry[lifrarkæfa] \dicTerm{lifrar··kæf|a} \dicIPA{{l}{\textsci}{v}{r}{a}{\textsubring{r}}{c\smash{\textsuperscript{h}}}{a}{i}{v}{a}} \dicPos{f}[1] \dicFlx{(‑u)}[5] \dicFieldCat{kulin.} \dicDirectTranslationCS{játrová paštika, játrovka}
\dicEntry[lifrarpylsa] \dicTerm{lifrar··pyls|a} \dicIPA{{l}{\textsci}{v}{r}{a}{\textsubring{r}}{p\smash{\textsuperscript{h}}}{\textsci}{l}{s}{a}} \dicPos{f}[1] \dicFlx{(‑u, ‑ur)}[13] \dicFieldCat{kulin.} \dicDirectTranslationCS{jitrnice}
\dicEntry[lifur] \dicTerm{lif|ur} \dicIPA{{l}{\textsci}{\textlengthmark}{v}{\textscy}{\textsubring{r}}} \dicPos{f}[6] \dicFlx{(‑rar, ‑rar)}[4] \dicFieldCat{anat.} \dicDirectTranslationCS{játra}
\dicEntry[ligg] \dicTerm{ligg} \dicIPA{{l}{\textsci}{\r{g}}{\textlengthmark}} \dicPos{v} \dicFlx{ind praes sg 1 pers} \dicLink{liggja}
\dicEntry[liggja] \dicTerm{liggja} \dicsymFrequent\  \dicIPA{{l}{\textsci}{\r{\textObardotlessj}}{\textlengthmark}{a}} \dicPos{v}[7] \dicFlx{(ligg, lá, lágum, lægi, legið)}[9] \textbf{1.} \dicDirectTranslationCS{ležet} \dicIndirectTranslationCS{(spočívat na něčem podélně)} \dicExampleIS{Blaðið liggur á gólfinu.} \dicExampleCS{Noviny leží na podlaze.}  \textbf{2.} \dicDirectTranslationCS{ležet} \dicIndirectTranslationCS{(být uložen ke spánku)}  \textbf{3.} \dicDirectTranslationCS{ležet, rozkládat se, prostírat se} \dicExampleIS{Vegurinn liggur meðfram ánni.} \dicExampleCS{Cesta vede podél řeky.};  \dicIdiom{liggja}[að]{ \dicPhraseIS{liggja að e‑u}} \dicDirectTranslationCS{sousedit s~(čím)};  \dicIdiom{liggja}[á]{ \dicPhraseIS{ekkert liggur á}} \dicDirectTranslationCS{všechno popořadě, vše ve správný čas}; { \dicPhraseIS{e‑m liggur á}} \dicFlx{impers} \dicDirectTranslationCS{(kdo) má naspěch, (kdo) pospíchá, (kdo) chvátá} \dicExampleIS{Honum lá á.} \dicExampleCS{Měl naspěch.}; { \dicPhraseIS{e‑m liggur á e‑u}} \dicFlx{impers} \dicDirectTranslationCS{(kdo) má naspěch s~(čím), (komu) záleží na (čem)} \dicExampleIS{Honum liggur mikið á sendingunni.} \dicExampleCS{Na té zásilce mu moc záleží.}; { \dicPhraseIS{liggja á eggjum}} \dicDirectTranslationCS{sedět na vejcích}; { \dicPhraseIS{það liggur vel á e‑m}} \dicFlx{impers} \dicDirectTranslationCS{(kdo) se dobře cítí, (kdo) je v~dobrém rozpoložení};  \dicIdiom{liggja}[fyrir]{ \dicPhraseIS{e‑að á fyrir e‑m að liggja}} \dicDirectTranslationCS{(co) na (koho) čeká, (co) má (kdo) před sebou}; { \dicPhraseIS{e‑að liggur fyrir e‑m}} \dicDirectTranslationCS{(co) na (koho) čeká, (co) leží před (kým) (rozhodnutí ap.)}; { \dicPhraseIS{liggja fyrir}} \dicDirectTranslationCS{ležet v~posteli};  \dicIdiom{liggja}[í]{ \dicPhraseIS{liggja í e‑u}} \dicDirectTranslationCS{ležet s~(čím) (chřipkou ap.)} \dicExampleIS{liggja í flensu} \dicExampleCS{ležet s~chřipkou}; { \dicPhraseIS{e‑að liggur í loftinu}} \dicLangCat{přen.} \dicDirectTranslationCS{(co) visí ve vzduchu, (co) se vznáší ve vzduchu (napětí ap.)}; { \dicPhraseIS{liggja í því}} \dicDirectTranslationCS{prohrát, dostat nakládačku};  \dicIdiom{liggja}[með]{ \dicPhraseIS{liggja með e‑m}} \dicDirectTranslationCS{spát s~(kým)};  \dicIdiom{liggja}[nærri]{ \dicPhraseIS{það lá nærri e‑u}} \dicFlx{impers} \dicDirectTranslationCS{schylovalo se k~(čemu), (co) bylo na spadnutí};  \dicIdiom{liggja}[undir]{ \dicPhraseIS{liggja undir grun}} \dicDirectTranslationCS{čelit podezření, být v~podezření};  \dicIdiom{liggja}[úti]{ \dicPhraseIS{liggja úti}} \dicDirectTranslationCS{spát venku, spát pod širým nebem};  \dicIdiom{liggja}[við]{ \dicPhraseIS{e‑m liggur við e‑u}} \dicFlx{impers} \dicDirectTranslationCS{(kdo) málem (co) (upadl ap.)}; { \dicPhraseIS{e‑að liggur beint við}} \dicDirectTranslationCS{(co) je přirozené}; { \dicPhraseIS{ef\,/\addthin þegar mikið liggur við}} \dicDirectTranslationCS{pokud to je nutné}; { \dicPhraseIS{það lá við e‑u}} \dicFlx{impers} \dicDirectTranslationCS{schylovalo se k~(čemu)};  \dicIdiom{liggja}[yfir]{ \dicPhraseIS{liggja yfir e‑u}} \dicDirectTranslationCS{ležet nad (čím), ležet v~(čem) (v~učení ap.)};  \dicIdiom{liggja}{ \dicPhraseIS{liggja í augum uppi}} \dicDirectTranslationCS{být jasný\,/\addthin zřejmý\,/\addthin očividný, bít do očí}
\dicEntry[lilja] \dicTerm{lilj|a} \dicIPA{{l}{\textsci}{l}{j}{a}} \dicPos{f}[1] \dicFlx{(‑u, ‑ur)}[7] \dicFieldCat{bot.} \dicDirectTranslationCS{lilie} \textit{(l.~{\textLA{Lilium}})}  \dicsymPhoto\ 
\dicFigure{1249.jpg}{Lilja}{Lilja - Zicha Ondřej, Biolib, Copyright/CC-BY-NC}
\dicEntry[lim] \dicTerm{lim} \dicIPA{{l}{\textsci}{\textlengthmark}{\textsubring{m}}} \dicPos{f}[4] \dicFlx{(‑ar, ‑ar)}[1] \dicDirectTranslationCS{listoví}
\dicEntry[limaburður] \dicTerm{lima··burð|ur} \dicIPA{{l}{\textsci}{\textlengthmark}{m}{a}{\textsubring{b}}{\textscy}{r}{ð}{\textscy}{\textsubring{r}}} \dicPos{m}[10] \dicFlx{(‑ar)}[7] \dicDirectTranslationCS{postoj, póza} \dicIndirectTranslationCS{(držení těla při stání)}
\dicEntry[limgerði] \dicTerm{lim··gerði} \dicIPA{{l}{\textsci}{m}{\r{\textObardotlessj}}{\textepsilon}{r}{ð}{\textsci}} \dicPos{n}[2] \dicFlx{(‑s, ‑)}[14] \dicDirectTranslationCS{živý plot}
\dicEntry[limlesta] \dicTerm{lim··lest|a} \dicIPA{{l}{\textsci}{m}{l}{\textepsilon}{s}{\textsubring{d}}{a}} \dicPos{v}[2] \dicFlx{(‑i, ‑)}[12] \dicFlx{acc} \dicDirectTranslationCS{zmrzačit, polámat kosti}
\dicEntry[limra] \dicTerm{limr|a} \dicIPA{{l}{\textsci}{m}{r}{a}} \dicPos{f}[1] \dicFlx{(‑u, ‑ur)}[7] \dicFieldCat{lit.} \dicDirectTranslationCS{limerik} \dicIndirectTranslationCS{(zábavný poetický útvar)}
\dicEntry[limur] \dicTerm{lim|ur} \dicIPA{{l}{\textsci}{\textlengthmark}{m}{\textscy}{\textsubring{r}}} \dicPos{m}[9] \dicFlx{(‑s, ‑ir)}[8] \textbf{1.} \dicSynonym{útlimur} \dicDirectTranslationCS{úd, končetina} \dicExampleIS{limir líkamans} \dicExampleCS{údy těla}  \textbf{2.} \dicSynonym{getnaðarlimur} \dicDirectTranslationCS{pyj}
\dicEntry[lina] \dicTerm{lin|a} \dicIPA{{l}{\textsci}{\textlengthmark}{n}{a}} \dicPos{v}[1] \dicFlx{(‑aði)}[1] \dicFlx{acc} \dicDirectTranslationCS{(z)mírnit, (z)tlumit, (z)měkčit} \dicExampleIS{Kærleikurinn linar sorg hennar.} \dicExampleCS{Láska mírní její žal.};  \dicIdiom{linast}{ \dicPhraseIS{linast}} \dicFlx{refl} {\textbf{a.}} \dicSynonym*{mýkjast} \dicDirectTranslationCS{(z)měknout} \dicIndirectTranslationCS{(stát\,/\addthin stávat se měkkým)};  {\textbf{b.}} \dicSynonym*{slá af fyrri skoðun} \dicDirectTranslationCS{obměkčit se, ustoupit};  {\textbf{c.}} \dicSynonym*{missa krafta} \dicDirectTranslationCS{(ze)slábnout, ztratit\,/\addthin ztrácet sílu}
\dicEntry[lind] \dicTerm{lind} \dicsymFrequent\  \dicIPA{{l}{\textsci}{n}{\textsubring{d}}} \dicPos{f}[7] \dicFlx{(‑ar, ‑ir)}[1] \textbf{1.} \dicSynonym{uppspretta} \dicDirectTranslationCS{zřídlo, pramen} \dicExampleIS{Hún vísaði þyrstum hermönnum á lind.} \dicExampleCS{Ukázala žíznivým vojákům cestu k~pramenu.}  \textbf{2.} \dicLangCat{přen.} \dicSynonym{upptök} \dicDirectTranslationCS{zřídlo, původ}  \textbf{3.} \dicFieldCat{bot.} \dicSynonym{linditré} \dicDirectTranslationCS{lípa} \textit{(l.~{\textLA{Tilia}})}  \dicsymPhoto\ 
\dicFigure{16367.jpg}{Lind}{Lind - Hanzlík Václav, Biolib, Copyright/CC-BY-NC}
\dicEntry[lindá] \dicTerm{lind··á} \dicIPA{{l}{\textsci}{n}{\textsubring{d}}{au}} \dicPos{f}[4] \dicFlx{(‑r, ‑r)}[18] \dicFieldCat{geol.} \dicDirectTranslationCS{řeka napájená zřídly}
\dicEntry[lindi] \dicTerm{lind|i} \dicIPA{{l}{\textsci}{n}{\textsubring{d}}{\textsci}} \dicPos{m}[1] \dicFlx{(‑a, ‑ar)}[1] \dicSynonym{band} \dicDirectTranslationCS{opasek, pás}
\dicEntry[linditré] \dicTerm{lindi··tré} \dicIPA{{l}{\textsci}{n}{\textsubring{d}}{\textsci}{t\smash{\textsuperscript{h}}}{r}{j}{\textepsilon}} \dicPos{n}[2] \dicFlx{(‑s, ‑)}[36] \dicFieldCat{bot.} \dicSynonym{lind} \dicDirectTranslationCS{lípa} \textit{(l.~{\textLA{Tilia}})}
\dicEntry[lindýr] \dicTerm{lin··dýr} \dicIPA{{l}{\textsci}{n}{\textsubring{d}}{i}{\textsubring{r}}} \dicPos{n}[2] \dicFlx{(‑s, ‑)}[5] \dicFieldCat{zool.} \dicDirectTranslationCS{měkkýš} \textit{(l.~{\textLA{Mollusca}})}
\dicEntry[lingerður] \dicTerm{lin··gerður} \dicIPA{{l}{\textsci}{n}{\r{\textObardotlessj}}{\textepsilon}{r}{ð}{\textscy}{\textsubring{r}}} \dicPos{adj}[2]\dicFlx{}[-1] \textbf{1.} \dicSynonym*{kveifarlegur} \dicDirectTranslationCS{změkčilý}  \textbf{2.} \dicSynonym{veikburða} \dicDirectTranslationCS{zesláblý, neduživý}
\dicEntry[linka] \dicTerm{link|a} \dicIPA{{l}{i}{\r{\ng}}{\r{g}}{a}} \dicPos{f}[1] \dicFlx{(‑u)}[5] \textbf{1.} \dicSynonym{lasleiki} \dicDirectTranslationCS{churavost, neduživost}  \textbf{2.} \dicSynonym{dugleysi} \dicDirectTranslationCS{neschopnost, nezpůsobilost}
\dicEntry[linmæli] \dicTerm{lin··mæli} \dicIPA{{l}{\textsci}{n}{m}{a}{i}{l}{\textsci}} \dicPos{n}[2] \dicFlx{(‑s)}[20] \dicFieldCat{jaz.} \dicIndirectTranslationCS{druh výslovnosti -- používání neaspirované plozivy ([\textsubring{b}], [\textsubring{d}], [\r{g}], [\r{\textObardotlessj}]) uprostřed slova po dlouhé samohlásce (standardní výslovnost)} \dicAntonym{harðmæli}
\dicEntry[linna] \dicTerm{linn|a} \dicIPA{{l}{\textsci}{n}{\textlengthmark}{a}} \dicPos{v}[2] \dicFlx{(‑ti, ‑t)}[80] \dicPhraseIS{e‑u linnir} \dicFlx{impers} \dicDirectTranslationCS{(co) přestává\,/\addthin ustává\,/\addthin končí} \dicExampleIS{Óveðrinu er tekið að linna.} \dicExampleCS{Bouře pomalu ustává.};  \dicPhraseIS{linna ekki látum} \dicDirectTranslationCS{nepřestat, neustat (s~prosbami ap.)}
\dicEntry[linnulaus] \dicTerm{linnu··laus} \dicIPA{{l}{\textsci}{n}{\textlengthmark}{\textscy}{l}{\oe i}{s}} \dicPos{adj}[5]\dicFlx{}[-1] \dicSynonym{sífelldur} \dicDirectTranslationCS{nepřetržitý, neustálý, trvalý} \dicExampleIS{linnulausar árásir} \dicExampleCS{nepřetržité útoky}
\dicEntry[linsa] \dicTerm{lins|a} \dicIPA{{l}{\textsci}{n}{s}{a}} \dicPos{f}[1] \dicFlx{(‑u, ‑ur)}[7] \textbf{1.} \dicFieldCat{fyz.} \dicDirectTranslationCS{čočka}  \textbf{2.} \dicSynonym*{augnlinsa} \dicDirectTranslationCS{(kontaktní) čočka}
\dicEntry[linsoðinn] \dicTerm{lin··soðinn} \dicIPA{{l}{\textsci}{n}{s}{\textopeno}{ð}{\textsci}{\textsubring{n}}} \dicPos{adj}[6]\dicFlx{}[-2] \dicDirectTranslationCS{(jsoucí) naměkko} \dicExampleIS{linsoðið egg} \dicExampleCS{vejce uvařené naměkko}
\dicEntry[linsubaun] \dicTerm{linsu··baun} \dicIPA{{l}{\textsci}{n}{s}{\textscy}{\textsubring{b}}{\oe i}{\textsubring{n}}} \dicPos{f}[7] \dicFlx{(‑ar, ‑ir)}[1] \dicFieldCat{bot.} \dicDirectTranslationCS{čočka} \textit{(l.~{\textLA{Lens culinaris}})}
\dicEntry[linun] \dicTerm{lin|un} \dicIPA{{l}{\textsci}{\textlengthmark}{n}{\textscy}{\textsubring{n}}} \dicPos{f}[7] \dicFlx{(‑unar)}[9] \textbf{1.} \dicSynonym{fró} \dicDirectTranslationCS{uklidnění, úleva}  \textbf{2.} \dicSynonym{mildun} \dicDirectTranslationCS{zmírnění} \dicExampleIS{linun refsingar} \dicExampleCS{zmírnění trestu}
\dicEntry[linur] \dicTerm{linur} \dicIPA{{l}{\textsci}{\textlengthmark}{n}{\textscy}{\textsubring{r}}} \dicPos{adj}[1]\dicFlx{}[-1] \textbf{1.} \dicSynonym{mjúkur} \dicDirectTranslationCS{hebký, měkký, jemný}  \textbf{2.} \dicSynonym{slakur} \dicDirectTranslationCS{slabý, chabý} \dicExampleIS{linur í fótbolta} \dicExampleCS{slabý ve fotbale}  \textbf{3.} \dicSynonym{slappur} \dicDirectTranslationCS{indisponovaný, slabý}
\dicEntry[lipur] \dicTerm{lipur} \dicIPA{{l}{\textsci}{\textlengthmark}{\textsubring{b}}{\textscy}{\textsubring{r}}} \dicPos{adj}[9] \dicFlx{(f ‑)}[1] \textbf{1.} \dicSynonym*{flinkur} \dicDirectTranslationCS{zručný, šikovný} \dicExampleIS{Hún er lipur í höndunum.} \dicExampleCS{Má šikovné ruce.}  \textbf{2.} \dicSynonym{þægilegur} \dicDirectTranslationCS{vstřícný, přívětivý}  \textbf{3.} \dicSynonym{meðfærilegur} \dicDirectTranslationCS{(uživatelsky) přívětivý}
\dicEntry[lipurð] \dicTerm{lipurð} \dicIPA{{l}{\textsci}{\textlengthmark}{\textsubring{b}}{\textscy}{r}{\texttheta}} \dicPos{f}[4] \dicFlx{(‑ar)}[3] \textbf{1.} \dicSynonym*{fimleiki} \dicDirectTranslationCS{obratnost, zručnost}  \textbf{2.} \dicSynonym*{lipur í viðmóti} \dicDirectTranslationCS{vstřícnost, přívětivost}
\dicEntry[lirfa] \dicTerm{lirf|a} \dicIPA{{l}{\textsci}{r}{v}{a}} \dicPos{f}[1] \dicFlx{(‑u, ‑ur)}[19] \dicDirectTranslationCS{larva}
\dicEntry[list] \dicTerm{list} \dicsymFrequent\  \dicIPA{{l}{\textsci}{s}{\textsubring{d}}} \dicPos{f}[7] \dicFlx{(‑ar, ‑ir)}[1] \textbf{1.} \dicDirectTranslationCS{umění} \dicIndirectTranslationCS{(odvětví lidské tvorby)} \dicExampleIS{unna listum} \dicExampleCS{zbožňovat umění}  \textbf{2.} \dicSynonym{færni} \dicDirectTranslationCS{umění, um, dovednost}
\dicEntry[listaháskóli] \dicTerm{lista··há·skól|i} \dicIPA{{l}{\textsci}{s}{\textsubring{d}}{a}{h}{au}{s}{\r{g}}{ou}{l}{\textsci}} \dicPos{m}[1] \dicFlx{(‑a, ‑ar)}[1] \dicFieldCat{škol.} \dicDirectTranslationCS{akademie umění}
\dicEntry[listahátíð] \dicTerm{lista··há·tíð} \dicIPA{{l}{\textsci}{s}{\textsubring{d}}{a}{h}{au}{t\smash{\textsuperscript{h}}}{i}{\texttheta}} \dicPos{f}[7] \dicFlx{(‑ar, ‑ir)}[1] \dicDirectTranslationCS{umělecký festival}
\dicEntry[listakona] \dicTerm{lista··kon|a} \dicIPA{{l}{\textsci}{s}{\textsubring{d}}{a}{k\smash{\textsuperscript{h}}}{\textopeno}{n}{a}} \dicPos{f}[1] \dicFlx{(‑u, ‑ur)}[27] \dicDirectTranslationCS{umělkyně}
\dicEntry[listamaður] \dicTerm{lista··|maður} \dicsymFrequent\  \dicIPA{{l}{\textsci}{s}{\textsubring{d}}{a}{m}{a}{ð}{\textscy}{\textsubring{r}}} \dicPos{m}[13] \dicFlx{(‑manns, ‑menn)}[2] \dicDirectTranslationCS{umělec, umělkyně} \dicExampleIS{Það er eitt dáðasta verk listamannsins.} \dicExampleCS{Je to jedno z~nejvíce obdivovaných děl umělce.}
\dicEntry[listasafn] \dicTerm{lista··|safn} \dicIPA{{l}{\textsci}{s}{\textsubring{d}}{a}{s}{a}{\textsubring{b}}{\textsubring{n}}} \dicPos{n}[2] \dicFlx{(‑safns, ‑söfn)}[8] \dicDirectTranslationCS{galerie (umění)}
\dicEntry[listasaga] \dicTerm{lista··|saga} \dicIPA{{l}{\textsci}{s}{\textsubring{d}}{a}{s}{a}{\textbabygamma}{a}} \dicPos{f}[1] \dicFlx{(‑sögu)}[2] \dicDirectTranslationCS{dějiny umění}
\dicEntry[listaskóli] \dicTerm{lista··skól|i} \dicIPA{{l}{\textsci}{s}{\textsubring{d}}{a}{s}{\r{g}}{ou}{l}{\textsci}} \dicPos{m}[1] \dicFlx{(‑a, ‑ar)}[1] \dicFieldCat{škol.} \dicDirectTranslationCS{umělecká škola}
\dicEntry[listaverk] \dicTerm{lista··verk} \dicIPA{{l}{\textsci}{s}{\textsubring{d}}{a}{v}{\textepsilon}{\textsubring{r}}{\r{g}}} \dicPos{n}[2] \dicFlx{(‑s, ‑)}[5] \dicDirectTranslationCS{umělecké dílo}
\dicEntry[listdans] \dicTerm{list··dans} \dicIPA{{l}{\textsci}{s}{\textsubring{d}}{a}{n}{s}} \dicPos{m}[4] \dicFlx{(‑, ‑ar)}[29] \dicSynonym{ballett} \dicDirectTranslationCS{umělecký tanec, balet}
\dicEntry[listfengur] \dicTerm{list··fengur} \dicIPA{{l}{\textsci}{s}{\textsubring{d}}{f}{ei}{\ng}{\r{g}}{\textscy}{\textsubring{r}}} \dicPos{adj}[1]\dicFlx{}[-1] \dicDirectTranslationCS{umělecky nadaný}
\dicEntry[listfræðingur] \dicTerm{list·fræð··ing|ur} \dicIPA{{l}{\textsci}{s}{\textsubring{d}}{f}{r}{a}{i}{ð}{i}{\ng}{\r{g}}{\textscy}{\textsubring{r}}} \dicPos{m}[6] \dicFlx{(‑s, ‑ar)}[8] \dicDirectTranslationCS{historik\,/\addthin historička umění, kunsthistorik, kunsthistorička}
\dicEntry[listgáfa] \dicTerm{list··gáf|a} \dicIPA{{l}{\textsci}{s}{\textsubring{d}}{\r{g}}{au}{v}{a}} \dicPos{f}[1] \dicFlx{(‑u, ‑ur)}[13] \dicDirectTranslationCS{umělecký talent}
\dicEntry[listi] \dicTerm{list|i} \dicsymFrequent\  \dicIPA{{l}{\textsci}{s}{\textsubring{d}}{\textsci}} \dicPos{m}[1] \dicFlx{(‑a, ‑ar)}[1] \textbf{1.} \dicSynonym{skrá\smash{\textsuperscript{1}}} \dicDirectTranslationCS{seznam, soupis, listina, rejstřík, žebříček} \dicExampleIS{óskalisti} \dicExampleCS{seznam přání};  \dicPhraseIS{listi yfir e‑n\,/\addthin e‑ð} \dicDirectTranslationCS{seznam (koho\,/\addthin čeho)}  \textbf{2.} \dicSynonym{ræma} \dicDirectTranslationCS{lišta, laťka}
\dicEntry[listiðn] \dicTerm{list··iðn} \dicIPA{{l}{\textsci}{s}{\textsubring{d}}{\textsci}{ð}{\textsubring{n}}} \dicPos{f}[7] \dicFlx{(‑ar, ‑ir)}[1] \dicDirectTranslationCS{umělecké řemeslo}
\dicEntry[listiðnaður] \dicTerm{list··iðn·að|ur} \dicIPA{{l}{\textsci}{s}{\textsubring{d}}{\textsci}{ð}{n}{a}{ð}{\textscy}{\textsubring{r}}} \dicPos{m}[10] \dicFlx{(‑ar)}[9] \dicDirectTranslationCS{užité umění}
\dicEntry[listmálari] \dicTerm{list··mál·ar|i} \dicIPA{{l}{\textsci}{s}{\textsubring{d}}{m}{au}{l}{a}{r}{\textsci}} \dicPos{m}[1] \dicFlx{(‑a, ‑ar)}[13] \dicDirectTranslationCS{malíř(ka)} \dicIndirectTranslationCS{(umělec malující obrazy)}
\dicEntry[listnám] \dicTerm{list··nám} \dicIPA{{l}{\textsci}{s}{\textsubring{d}}{n}{au}{\textsubring{m}}} \dicPos{n}[2] \dicFlx{(‑s)}[2] \dicDirectTranslationCS{studium umění}
\dicEntry[listrænn] \dicTerm{list··rænn} \dicIPA{{l}{\textsci}{s}{\textsubring{d}}{r}{a}{i}{\textsubring{d}}{\textsubring{n}}} \dicPos{adj}[7]\dicFlx{}[-1] \dicDirectTranslationCS{umělecký}
\dicEntry[listskautar] \dicTerm{list··skautar} \dicIPA{{l}{\textsci}{s}{\textsubring{d}}{s}{\r{g}}{\oe i}{\textsubring{d}}{a}{\textsubring{r}}} \dicPos{m}[1] \dicFlx{pl}[2] \dicFieldCat{sport.} \dicDirectTranslationCS{krasobruslení}
\dicEntry[listunnandi] \dicTerm{list··unn·|andi} \dicIPA{{l}{\textsci}{s}{\textsubring{d}}{\textscy}{n}{a}{n}{\textsubring{d}}{\textsci}} \dicPos{m}[2] \dicFlx{(‑anda, ‑endur)}[1] \dicDirectTranslationCS{milovník\,/\addthin milovnice umění}
\dicEntry[lita] \dicTerm{lit|a} \dicsymFrequent\  \dicIPA{{l}{\textsci}{\textlengthmark}{\textsubring{d}}{a}} \dicPos{v}[1] \dicFlx{(‑aði)}[1] \dicFlx{acc} \textbf{1.} \dicSynonym{mála} \dicDirectTranslationCS{(na)barvit, nabarvovat, obarvit, obarvovat} \dicExampleIS{lita fötin} \dicExampleCS{nabarvit oblečení}  \textbf{2.} \dicDirectTranslationCS{zabarvit, zkreslit (skutečnost ap.)} \dicExampleIS{Skáldskapurinn litar frá sér.} \dicExampleCS{Literární dílo zkresluje.};  \dicIdiom{litast}{ \dicPhraseIS{litast}} \dicFlx{refl} \dicDirectTranslationCS{zbarvit se, zabarvit se, obarvit se};  \dicIdiom{litast}[um]{ \dicPhraseIS{litast um}} \dicFlx{refl} \dicDirectTranslationCS{(po)rozhlédnout se};  \dicIdiom{litast}[upp]{ \dicPhraseIS{litast upp}} \dicFlx{refl} \dicSynonym*{tapa lit} \dicDirectTranslationCS{(vy)blednout, ztratit\,/\addthin ztrácet barvu}
\dicEntry[litaraft] \dicTerm{litar··aft} \dicIPA{{l}{\textsci}{\textlengthmark}{\textsubring{d}}{a}{r}{a}{f}{\textsubring{d}}} \dicPos{n}[2] \dicFlx{(‑s)}[2] \dicDirectTranslationCS{barva\,/\addthin zbarvení obličeje}
\dicEntry[litarháttur] \dicTerm{litar··hátt|ur} \dicIPA{{l}{\textsci}{\textlengthmark}{\textsubring{d}}{a}{\textsubring{r}}{h}{au}{h}{\textsubring{d}}{\textscy}{\textsubring{r}}} \dicPos{m}[12] \dicFlx{(‑ar)}[8] \dicDirectTranslationCS{barva pleti} \dicExampleIS{fordómar vegna litarháttar} \dicExampleCS{předsudky k~barvě pleti}
\dicEntry[litasamsetning] \dicTerm{lita··sam·set·ning} \dicIPA{{l}{\textsci}{\textlengthmark}{\textsubring{d}}{a}{s}{a}{m}{s}{\textepsilon}{h}{\textsubring{d}}{n}{i}{\ng}{\r{g}}} \dicPos{f}[4] \dicFlx{(‑ar, ‑ar)}[5] \dicDirectTranslationCS{barevná kombinace, barevné schéma}
\dicEntry[litaspjald] \dicTerm{lita··|spjald} \dicIPA{{l}{\textsci}{\textlengthmark}{\textsubring{d}}{a}{s}{\textsubring{b}}{j}{a}{l}{\textsubring{d}}} \dicPos{n}[2] \dicFlx{(‑spjalds, ‑spjöld)}[8] \textbf{1.} \dicDirectTranslationCS{paleta} \dicIndirectTranslationCS{(malířská deska k~míchání a~roztírání barev)}  \textbf{2.} \dicFieldCat{poč.} \dicDirectTranslationCS{paleta} \dicIndirectTranslationCS{(na výběr barev)}
\dicEntry[litblinda] \dicTerm{lit··blind|a} \dicIPA{{l}{\textsci}{\textlengthmark}{\textsubring{d}}{\textsubring{b}}{l}{\textsci}{n}{\textsubring{d}}{a}} \dicPos{f}[1] \dicFlx{(‑u)}[5] \dicDirectTranslationCS{barvoslepost}
\dicEntry[litblindur] \dicTerm{lit··blindur} \dicIPA{{l}{\textsci}{\textlengthmark}{\textsubring{d}}{\textsubring{b}}{l}{\textsci}{n}{\textsubring{d}}{\textscy}{\textsubring{r}}} \dicPos{adj}[2]\dicFlx{}[-14] \dicFieldCat{med.} \dicDirectTranslationCS{barvoslepý}
\dicEntry[litbrigði] \dicTerm{lit··brigði} \dicIPA{{l}{\textsci}{\textlengthmark}{\textsubring{d}}{\textsubring{b}}{r}{\textsci}{\textbabygamma}{ð}{\textsci}} \dicPos{n}[2] \dicFlx{(‑s, ‑)}[14] \dicDirectTranslationCS{odstín (barvy)}
\dicEntry[litgreining] \dicTerm{lit··grein·ing} \dicIPA{{l}{\textsci}{\textlengthmark}{\textsubring{d}}{\r{g}}{r}{ei}{n}{i}{\ng}{\r{g}}} \dicPos{f}[4] \dicFlx{(‑ar)}[7] \textbf{1.} \dicFieldCat{polygr.} \dicDirectTranslationCS{separace barev}  \textbf{2.} \dicDirectTranslationCS{barevná analýza} \dicIndirectTranslationCS{(osobní kombinace barev)}
\dicEntry[Litháen] \dicTerm{Litháen} \dicIPA{{l}{\textsci}{\textlengthmark}{t\smash{\textsuperscript{h}}}{au}{\textepsilon}{\textsubring{n}}} \dicPos{n}[2] \dicFlx{(‑s)}[32] \dicFieldCat{geog.} \dicDirectTranslationCS{Litva}
\dicEntry[Lithái] \dicTerm{Lithá|i} \dicIPA{{l}{\textsci}{\textlengthmark}{t\smash{\textsuperscript{h}}}{au}{\textsci}} \dicPos{m}[1] \dicFlx{(‑a, ‑ar)}[1] \dicDirectTranslationCS{Litevec, Litevka}
\dicEntry[litháíska] \dicTerm{litháísk|a} \dicIPA{{l}{\textsci}{\textlengthmark}{t\smash{\textsuperscript{h}}}{au}{i}{s}{\r{g}}{a}} \dicPos{f}[1] \dicFlx{(‑u)}[5] \dicLink{litháska}
\dicEntry[litháískur] \dicTerm{litháískur} \dicIPA{{l}{\textsci}{\textlengthmark}{t\smash{\textsuperscript{h}}}{au}{i}{s}{\r{g}}{\textscy}{\textsubring{r}}} \dicPos{adj}[1]\dicFlx{}[-6] \dicLink{litháskur}
\dicEntry[litháska] \dicTerm{lithásk|a}\dicTerm{, litháíska} \dicIPA{{l}\-{\textsci}\-{\textlengthmark}\-{t\smash{\textsuperscript{h}}}\-{au}\-{s}\-{\r{g}}\-{a}\-} \dicPos{f}[1] \dicFlx{(‑u)}[5] \dicDirectTranslationCS{litevština}
\dicEntry[litháskur] \dicTerm{litháskur}\dicTerm{, litháískur} \dicIPA{{l}\-{\textsci}\-{\textlengthmark}\-{t\smash{\textsuperscript{h}}}\-{au}\-{s}\-{\r{g}}\-{\textscy}\-{\textsubring{r}}\-} \dicPos{adj}[1]\dicFlx{}[-6] \dicDirectTranslationCS{litevský}
\dicEntry[lithimna] \dicTerm{lit··himn|a} \dicIPA{{l}{\textsci}{\textlengthmark}{\textsubring{d}}{h}{\textsci}{m}{n}{a}} \dicPos{f}[1] \dicFlx{(‑u, ‑ur)}[7] \dicFieldCat{anat.} \dicDirectTranslationCS{duhovka}
\dicEntry[liti] \dicTerm{liti} \dicIPA{{l}{\textsci}{\textlengthmark}{\textsubring{d}}{\textsci}} \dicPos{v} \dicFlx{con pf sg 1 pers} \dicLink{líta}
\dicEntry[litið] \dicTerm{litið} \dicIPA{{l}{\textsci}{\textlengthmark}{\textsubring{d}}{\textsci}{\texttheta}} \dicPos{v} \dicFlx{supin} \dicLink{líta}
\dicEntry[litlafingurs] \dicTerm{litla··fingurs} \dicIPA{{l}{\textsci}{h}{\textsubring{d}}{l}{a}{f}{\textsci}{\ng}{\r{g}}{\textscy}{\textsubring{r}}{s}} \dicPos{m} \dicFlx{sg gen} \dicLink{litlifingur}
\dicEntry[litlatá] \dicTerm{litla··tá} \dicIPA{{l}{\textsci}{h}{\textsubring{d}}{l}{a}{t\smash{\textsuperscript{h}}}{au}} \dicPos{f}[9] \dicFlx{(litlutáar, litlutær)}[2] \dicDirectTranslationCS{malík\,/\addthin malíček (u~nohy)}
\dicEntry[litlaus] \dicTerm{lit··laus} \dicIPA{{l}{\textsci}{\textlengthmark}{\textsubring{d}}{l}{\oe i}{s}} \dicPos{adj}[5]\dicFlx{}[-1] \textbf{1.} \dicSynonym*{ekki í lit} \dicDirectTranslationCS{bezbarvý}  \textbf{2.} \dicSynonym{dauflegur} \dicDirectTranslationCS{mdlý, fádní, bezbarvý} \dicExampleIS{litlaus frásögn} \dicExampleCS{bezbarvé vyprávění}
\dicEntry[litlifingur] \dicTerm{litli··fingur} \dicIPA{{l}{\textsci}{h}{\textsubring{d}}{l}{\textsci}{f}{i}{\ng}{\r{g}}{\textscy}{\textsubring{r}}} \dicPos{m}[5] \dicFlx{(litlafingurs, litlufingur)}[14] \dicDirectTranslationCS{malík\,/\addthin malíček (u~ruky)}
\dicEntry[litlufingur] \dicTerm{litlu··fingur} \dicIPA{{l}{\textsci}{h}{\textsubring{d}}{l}{\textscy}{f}{\textsci}{\ng}{\r{g}}{\textscy}{\textsubring{r}}} \dicPos{m} \dicFlx{pl nom} \dicLink{litlifingur}
\dicEntry[litmynd] \dicTerm{lit··mynd} \dicIPA{{l}{\textsci}{\textlengthmark}{\textsubring{d}}{m}{\textsci}{n}{\textsubring{d}}} \dicPos{f}[7] \dicFlx{(‑ar, ‑ir)}[1] \dicDirectTranslationCS{barevný obraz, barevná fotografie}
\dicEntry[litningur] \dicTerm{litn··ing|ur} \dicIPA{{l}{\textsci}{h}{\textsubring{d}}{n}{i}{\ng}{\r{g}}{\textscy}{\textsubring{r}}} \dicPos{m}[6] \dicFlx{(‑s, ‑ar)}[8] \dicFieldCat{biol.} \dicDirectTranslationCS{chromozom}
\dicEntry[litprentun] \dicTerm{lit··prent|un} \dicIPA{{l}{\textsci}{\textlengthmark}{\textsubring{d}}{p\smash{\textsuperscript{h}}}{r}{\textepsilon}{\textsubring{n}}{\textsubring{d}}{\textscy}{\textsubring{n}}} \dicPos{f}[7] \dicFlx{(‑unar)}[9] \dicFieldCat{polygr.} \dicDirectTranslationCS{barevný tisk}
\dicEntry[litríkur] \dicTerm{lit··ríkur} \dicIPA{{l}{\textsci}{\textlengthmark}{\textsubring{d}}{r}{i}{\r{g}}{\textscy}{\textsubring{r}}} \dicPos{adj}[1]\dicFlx{}[-1] \textbf{1.} \dicDirectTranslationCS{pestrobarevný, pestrý} \dicExampleIS{litrík föt} \dicExampleCS{pestrobarevné oblečení}  \textbf{2.} \dicDirectTranslationCS{pestrý, rozmanitý} \dicExampleIS{Söngkonan á litríkan feril að baki.} \dicExampleCS{Zpěvačka má za sebou pestrou kariéru.}
\dicEntry[litróf] \dicTerm{lit··róf} \dicIPA{{l}{\textsci}{\textlengthmark}{\textsubring{d}}{r}{ou}{f}} \dicPos{n}[2] \dicFlx{(‑s, ‑)}[5] \textbf{1.} \dicFieldCat{fyz.} \dicDirectTranslationCS{barevné spektrum}  \textbf{2.} \dicLangCat{přen.} \dicDirectTranslationCS{spektrum (názorové ap.)}
\dicEntry[litskrúð] \dicTerm{lit··skrúð} \dicIPA{{l}{\textsci}{\textlengthmark}{\textsubring{d}}{s}{\r{g}}{r}{u}{\texttheta}} \dicPos{n}[2] \dicFlx{(‑s)}[2] \dicDirectTranslationCS{pestrobarevnost}
\dicEntry[litskyggna] \dicTerm{lit··skyggn|a} \dicIPA{{l}{\textsci}{\textlengthmark}{\textsubring{d}}{s}{\r{\textObardotlessj}}{\textsci}{\r{g}}{n}{a}} \dicPos{f}[1] \dicFlx{(‑u, ‑ur)}[7] \dicDirectTranslationCS{barevný diapozitiv}
\dicEntry[litum] \dicTerm{litum} \dicIPA{{l}{\textsci}{\textlengthmark}{\textsubring{d}}{\textscy}{\textsubring{m}}} \dicPos{v} \dicFlx{ind pf pl 1 pers} \dicLink{líta}
\dicEntry[litun] \dicTerm{lit|un} \dicIPA{{l}{\textsci}{\textlengthmark}{\textsubring{d}}{\textscy}{\textsubring{n}}} \dicPos{f}[7] \dicFlx{(‑unar)}[9] \dicDirectTranslationCS{(na)barvení}
\dicEntry[litur] \dicTerm{lit|ur} \dicsymFrequent\  \dicIPA{{l}{\textsci}{\textlengthmark}{\textsubring{d}}{\textscy}{\textsubring{r}}} \dicPos{m}[10] \dicFlx{(‑ar, ‑ir)}[13] \textbf{1.} \dicDirectTranslationCS{barva, barevný odstín} \dicExampleIS{Hvernig er það á litinn?} \dicExampleCS{Jakou to má barvu?}  \textbf{2.} \dicSynonym*{litarefni} \dicDirectTranslationCS{barva, nátěr}  \textbf{3.} \dicSynonym*{tegund spila} \dicDirectTranslationCS{barva} \dicIndirectTranslationCS{(v~kartách)}
\dicEntry[Líbani] \dicTerm{Líban|i} \dicIPA{{l}{i}{\textlengthmark}{\textsubring{b}}{a}{n}{\textsci}} \dicPos{m}[3] \dicFlx{(‑a, ‑ar\,/\addthin ‑ir)}[3] \dicDirectTranslationCS{Libanonec, Libanonka}
\dicEntry[Líbanon] \dicTerm{Líbanon} \dicIPA{{l}{i}{\textlengthmark}{\textsubring{b}}{a}{n}{\textopeno}{\textsubring{n}}} \dicPos{n}[2] \dicFlx{(‑s)}[32] \dicFieldCat{geog.} \dicDirectTranslationCS{Libanon}
\dicEntry[líbanskur] \dicTerm{líbanskur} \dicIPA{{l}{i}{\textlengthmark}{\textsubring{b}}{a}{n}{s}{\r{g}}{\textscy}{\textsubring{r}}} \dicPos{adj}[1] \dicFlx{(f líbönsk)}[3] \dicDirectTranslationCS{libanonský}
\dicEntry[Líbería] \dicTerm{Líberí|a} \dicIPA{{l}{i}{\textlengthmark}{\textsubring{b}}{\textepsilon}{r}{i}{j}{a}} \dicPos{f}[1] \dicFlx{(‑u)}[6] \dicFieldCat{geog.} \dicDirectTranslationCS{Libérie}
\dicEntry[líberískur] \dicTerm{líberískur} \dicIPA{{l}{i}{\textlengthmark}{\textsubring{b}}{\textepsilon}{r}{i}{s}{\r{g}}{\textscy}{\textsubring{r}}} \dicPos{adj}[1]\dicFlx{}[-6] \dicDirectTranslationCS{liberijský}
\dicEntry[Líberíumaður] \dicTerm{Líberíu··|maður} \dicIPA{{l}{i}{\textlengthmark}{\textsubring{b}}{\textepsilon}{r}{i}{j}{\textscy}{m}{a}{ð}{\textscy}{\textsubring{r}}} \dicPos{m}[13] \dicFlx{(‑manns, ‑menn)}[2] \dicDirectTranslationCS{Liberijec, Liberijka}
\dicEntry[Líbía] \dicTerm{Líbí|a}\dicTerm{, Líbýa} \dicIPA{{l}\-{i}\-{\textlengthmark}\-{\textsubring{b}}\-{i}\-{j}\-{a}\-} \dicPos{f}[1] \dicFlx{(‑u)}[6] \dicFieldCat{geog.} \dicDirectTranslationCS{Libye}
\dicEntry[líbískur] \dicTerm{líbískur}\dicTerm{, líbýskur} \dicIPA{{l}\-{i}\-{\textlengthmark}\-{\textsubring{b}}\-{i}\-{s}\-{\r{g}}\-{\textscy}\-{\textsubring{r}}\-} \dicPos{adj}[1]\dicFlx{}[-6] \dicDirectTranslationCS{libyjský}
\dicEntry[Líbíumaður] \dicTerm{Líbíu··|maður} \dicIPA{{l}{i}{\textlengthmark}{\textsubring{b}}{i}{j}{\textscy}{m}{a}{ð}{\textscy}{\textsubring{r}}} \dicPos{m}[13] \dicFlx{(‑manns, ‑menn)}[2] \dicDirectTranslationCS{Libyjec, Libyjka}
\dicEntry[Líbýa] \dicTerm{Líbý|a} \dicIPA{{l}{i}{\textlengthmark}{\textsubring{b}}{i}{j}{a}} \dicPos{f}[1] \dicFlx{(‑u)}[6] \dicLink{Líbía}
\dicEntry[líbýskur] \dicTerm{líbýskur} \dicIPA{{l}{i}{\textlengthmark}{\textsubring{b}}{i}{s}{\r{g}}{\textscy}{\textsubring{r}}} \dicPos{adj}[1]\dicFlx{}[-6] \dicLink{líbískur}
\dicEntry[líð] \dicTerm{líð} \dicIPA{{l}{i}{\textlengthmark}{\texttheta}} \dicPos{v} \dicFlx{ind praes sg 1 pers} \dicLink{líða}
\dicEntry[líða] \dicTerm{líða} \dicsymFrequent\  \dicIPA{{l}{i}{\textlengthmark}{ð}{a}} \dicPos{v}[6] \dicFlx{(líð, leið, liðum, liði, liðið)}[66] \dicFlx{acc\,/\addthin dat} \textbf{1.} \dicSynonym*{um tíma} \dicDirectTranslationCS{plynout, ubíhat, utíkat} \dicIndirectTranslationCS{(o~čase)} \dicExampleIS{Árið líður hratt.} \dicExampleCS{Rok rychle ubíhá.}  \textbf{2.} \dicFlx{dat} \dicDirectTranslationCS{pokračovat, postupovat} \dicExampleIS{Hvernig líður verkinu?} \dicExampleCS{Jak pokračuje práce?}  \textbf{3.} \dicFlx{acc} \dicSynonym{þola} \dicDirectTranslationCS{snést, snášet, tolerovat} \dicIndirectTranslationCS{(často se záporem)} \dicExampleIS{líða þetta ekki} \dicExampleCS{nesnést to}  \textbf{4.} \dicPhraseIS{e‑m líður} \dicFlx{impers} \dicDirectTranslationCS{(kdo) se cítí (dobře ap.)} \dicExampleIS{Mér líður vel.} \dicExampleCS{Cítím se dobře.\,/\addthin Je mi dobře.};  \dicPhraseIS{hvernig líður þér?} \dicDirectTranslationCS{jak se máš?\,/\addthin jak se cítíš?};  \dicIdiom{líða}[að]{ \dicPhraseIS{það líður að e‑u}} \dicFlx{impers} \dicDirectTranslationCS{(co) se blíží (svátky ap.)};  \dicIdiom{líða}[hjá]{ \dicPhraseIS{e‑að líður hjá}} \dicDirectTranslationCS{(co) míjí, (co) pomine};  \dicIdiom{líða}[í]{ \dicPhraseIS{líða í ómegin}} \dicDirectTranslationCS{omdlít, upadnout do mdlob};  \dicIdiom{líða}[undir]{ \dicPhraseIS{e‑að líður undir lok}} \dicDirectTranslationCS{(co) se schyluje\,/\addthin chýlí ke konci};  \dicIdiom{líða}[út af]{ \dicPhraseIS{líða út af}} {\textbf{a.}} \dicDirectTranslationCS{omdlít};  {\textbf{b.}} \dicDirectTranslationCS{usnout} \dicExampleIS{líða út af í hægindastólnum} \dicExampleCS{usnout v~křesle};  \dicIdiom{líða}[yfir]{ \dicPhraseIS{það líður yfir e‑n}} \dicFlx{impers} \dicDirectTranslationCS{(kdo) omdlívá, (kdo) ztrácí vědomí} \dicExampleIS{Það leið yfir mig.} \dicExampleCS{Omdlel jsem.};  \dicIdiom{líðast}{ \dicPhraseIS{e‑m líðst e‑að}} \dicFlx{refl impers} \dicDirectTranslationCS{(co) se (komu) toleruje} \dicIndirectTranslationCS{(často se záporem)};  \dicIdiom{líða}{ \dicPhraseIS{hvað sem öðru líður}} \dicDirectTranslationCS{beztak, když se to vezme kolem a~kolem}
\dicEntry[líðan] \dicTerm{líðan} \dicIPA{{l}{i}{\textlengthmark}{ð}{a}{\textsubring{n}}} \dicPos{f}[4] \dicFlx{(‑ar)}[3] \dicDirectTranslationCS{(zdravotní) stav} \dicExampleIS{Líðan sjúklingsins fer batnandi.} \dicExampleCS{Zdravotní stav pacienta se zlepšuje.}
\dicEntry[líðandi] \dicTerm{líð··andi} \dicIPA{{l}{i}{\textlengthmark}{ð}{a}{n}{\textsubring{d}}{\textsci}} \dicPos{adj}[13] \dicFlx{indecl}[1] \dicPhraseIS{á líðandi stund} \dicFlx{adv} \dicSynonym{núna} \dicDirectTranslationCS{teď, v~současné chvíli}
\dicEntry[líf] \dicTerm{líf} \dicsymFrequent\  \dicIPA{{l}{i}{\textlengthmark}{f}} \dicPos{n}[2] \dicFlx{(‑s, ‑)}[5] \textbf{1.} \dicSynonym*{það að lifa} \dicDirectTranslationCS{život, žití, bytí};  \dicPhraseIS{vera á lífi} \dicDirectTranslationCS{být naživu}  \textbf{2.} \dicSynonym*{lífverur} \dicDirectTranslationCS{život, existence (organismu)} \dicExampleIS{lífið á jörðinni} \dicExampleCS{život na Zemi}  \textbf{3.} \dicSynonym*{dagleg tilvera} \dicDirectTranslationCS{život (všední ap.)}  \textbf{4.} \dicSynonym{fjör} \dicDirectTranslationCS{živo, rušno};  \dicIdiom{líf}{ \dicPhraseIS{af lífi og sál}} \dicFlx{adv} \dicDirectTranslationCS{z~celých sil}; { \dicPhraseIS{blása lífið í e‑ð}} \dicLangCat{přen.} \dicDirectTranslationCS{vdechnout život (čemu)}; { \dicPhraseIS{elska e‑n út af lífinu}} \dicDirectTranslationCS{milovat (koho) z~celé duše};  \dicPhraseIS{láta lífið} \dicDirectTranslationCS{přijít o~život}; { \dicPhraseIS{svipta sig lífi}} \dicDirectTranslationCS{vzít si život}; { \dicPhraseIS{taka e‑n af lífi}} \dicDirectTranslationCS{vzít (komu) život}; { \dicPhraseIS{vekja til lífsins}} \dicLangCat{přen.} \dicDirectTranslationCS{probudit k~životu}
\dicEntry[lífbátur] \dicTerm{líf··bát|ur} \dicIPA{{l}{i}{v}{\textsubring{b}}{au}{\textsubring{d}}{\textscy}{\textsubring{r}}} \dicPos{m}[6] \dicFlx{(‑s, ‑ar)}[22] \dicDirectTranslationCS{záchranný člun}
\dicEntry[lífbein] \dicTerm{líf··bein} \dicIPA{{l}{i}{v}{\textsubring{b}}{ei}{\textsubring{n}}} \dicPos{n}[2] \dicFlx{(‑s, ‑)}[5] \dicFieldCat{anat.} \dicDirectTranslationCS{stydká kost}
\dicEntry[lífeðlisfræði] \dicTerm{líf··eðlis·fræð|i} \dicIPA{{l}{i}{\textlengthmark}{v}{\textepsilon}{ð}{l}{\textsci}{s}{f}{r}{a}{i}{ð}{\textsci}} \dicPos{f}[3] \dicFlx{(‑i)}[3] \dicDirectTranslationCS{fyziologie}
\dicEntry[lífeðlisfræðingur] \dicTerm{líf·eðlis·fræð··ing|ur} \dicIPA{{l}\-{i}\-{\textlengthmark}\-{v}\-{\textepsilon}\-{ð}\-{l}\-{\textsci}\-{s}\-{f}\-{r}\-{a}\-{i}\-{ð}\-{i}\-{\ng}\-{\r{g}}\-{\textscy}\-{\textsubring{r}}\-} \dicPos{m}[6] \dicFlx{(‑s, ‑ar)}[8] \dicDirectTranslationCS{fyziolog, fyzioložka}
\dicEntry[lífefnafræði] \dicTerm{líf··efna·fræð|i} \dicIPA{{l}{i}{\textlengthmark}{v}{\textepsilon}{\textsubring{b}}{n}{a}{f}{r}{a}{i}{ð}{\textsci}} \dicPos{f}[3] \dicFlx{(‑i)}[3] \dicDirectTranslationCS{biochemie}
\dicEntry[lífefnafræðingur] \dicTerm{líf·efna·fræð··ing|ur} \dicIPA{{l}\-{i}\-{\textlengthmark}\-{v}\-{\textepsilon}\-{\textsubring{b}}\-{n}\-{a}\-{f}\-{r}\-{a}\-{i}\-{ð}\-{i}\-{\ng}\-{\r{g}}\-{\textscy}\-{\textsubring{r}}\-} \dicPos{m}[6] \dicFlx{(‑s, ‑ar)}[8] \dicDirectTranslationCS{biochemik, biochemička}
\dicEntry[líferni] \dicTerm{líf··erni} \dicIPA{{l}{i}{\textlengthmark}{v}{\textepsilon}{r}{\textsubring{d}}{n}{\textsci}} \dicPos{n}[2] \dicFlx{(‑s)}[20] \dicDirectTranslationCS{životospráva, životní styl}
\dicEntry[lífeyrir] \dicTerm{líf··eyr|ir} \dicIPA{{l}{i}{\textlengthmark}{v}{ei}{r}{\textsci}{\textsubring{r}}} \dicPos{m}[7] \dicFlx{(‑is)}[2] \dicDirectTranslationCS{důchod, penze} \dicIndirectTranslationCS{(pravidelný příjem osob plynoucí z~pojištění)} \dicExampleIS{taka lífeyris} \dicExampleCS{příjem důchodu}
\dicEntry[lífeyrisgreiðsla] \dicTerm{líf·eyris··greiðsl|a} \dicIPA{{l}{i}{\textlengthmark}{v}{ei}{r}{\textsci}{s}{\r{g}}{r}{ei}{ð}{s}{\textsubring{d}}{l}{a}} \dicPos{f}[1] \dicFlx{(‑u, ‑ur)}[13] \dicDirectTranslationCS{výplata důchodu}
\dicEntry[lífeyrisréttindi] \dicTerm{líf·eyris··rétt·indi} \dicIPA{{l}{i}{\textlengthmark}{v}{ei}{r}{\textsci}{s}{r}{j}{\textepsilon}{h}{\textsubring{d}}{\textsci}{n}{\textsubring{d}}{\textsci}} \dicPos{n}[2] \dicFlx{pl}[19] \dicDirectTranslationCS{nárok na důchod}
\dicEntry[lífeyrissjóðstekjur] \dicTerm{líf·eyris·sjóðs··tekjur} \dicIPA{{l}\-{i}\-{\textlengthmark}\-{v}\-{ei}\-{r}\-{\textsci}\-{s}\-{j}\-{ou}\-{ð}\-{s}\-{t\smash{\textsuperscript{h}}}\-{\textepsilon}\-{\r{\textObardotlessj}}\-{\textscy}\-{\textsubring{r}}\-} \dicPos{f}[1] \dicFlx{pl}[18] \dicDirectTranslationCS{příjmy z~penzijního fondu}
\dicEntry[lífeyrissjóður] \dicTerm{líf·eyris··sjóð|ur} \dicIPA{{l}{i}{\textlengthmark}{v}{ei}{r}{\textsci}{s}{j}{ou}{ð}{\textscy}{\textsubring{r}}} \dicPos{m}[9] \dicFlx{(‑s, ‑ir)}[6] \dicFieldCat{ekon.} \dicDirectTranslationCS{penzijní fond}
\dicEntry[lífeyrissparnaður] \dicTerm{líf·eyris··spar·nað|ur} \dicIPA{{l}\-{i}\-{\textlengthmark}\-{v}\-{ei}\-{r}\-{\textsci}\-{s}\-{\textsubring{b}}\-{a}\-{r}\-{\textsubring{d}}\-{n}\-{a}\-{ð}\-{\textscy}\-{\textsubring{r}}\-} \dicPos{m}[10] \dicFlx{(‑ar)}[9] \dicDirectTranslationCS{důchodové spoření}
\dicEntry[lífeyristryggingar] \dicTerm{líf··eyris··trygg·ingar} \dicIPA{{l}\-{i}\-{\textlengthmark}\-{v}\-{ei}\-{r}\-{\textsci}\-{s}\-{t\smash{\textsuperscript{h}}}\-{r}\-{\textsci}\-{\r{\textObardotlessj}}\-{i}\-{\ng}\-{\r{g}}\-{a}\-{\textsubring{r}}\-} \dicPos{f}[4] \dicFlx{pl}[6] \dicDirectTranslationCS{důchodové pojištění}
\dicEntry[lífeyrisþegi] \dicTerm{líf·eyris··þeg|i} \dicIPA{{l}{i}{\textlengthmark}{v}{ei}{r}{\textsci}{s}{\texttheta}{ei}{\textsci}} \dicPos{m}[1] \dicFlx{(‑a, ‑ar)}[1] \dicDirectTranslationCS{důchodce, důchodkyně, penzista, penzistka}
\dicEntry[líffræði] \dicTerm{líf··fræð|i} \dicIPA{{l}{i}{f}{\textlengthmark}{r}{a}{i}{ð}{\textsci}} \dicPos{f}[3] \dicFlx{(‑i)}[3] \dicDirectTranslationCS{biologie}
\dicEntry[líffræðingur] \dicTerm{líf·fræð··ing|ur} \dicIPA{{l}{i}{f}{\textlengthmark}{r}{a}{i}{ð}{i}{\ng}{\r{g}}{\textscy}{\textsubring{r}}} \dicPos{m}[6] \dicFlx{(‑s, ‑ar)}[8] \dicDirectTranslationCS{biolog, bioložka}
\dicEntry[líffærafræði] \dicTerm{líf·færa··fræð|i} \dicIPA{{l}{i}{f}{\textlengthmark}{a}{i}{r}{a}{f}{r}{a}{i}{ð}{\textsci}} \dicPos{f}[3] \dicFlx{(‑i)}[3] \dicDirectTranslationCS{anatomie}
\dicEntry[líffæri] \dicTerm{líf··færi} \dicIPA{{l}{i}{f}{\textlengthmark}{a}{i}{r}{\textsci}} \dicPos{n}[2] \dicFlx{(‑s, ‑)}[14] \dicFieldCat{anat.} \dicDirectTranslationCS{orgán, ústrojí}
\dicEntry[lífga] \dicTerm{lífg|a} \dicIPA{{l}{i}{v}{\r{g}}{a}} \dicPos{v}[1] \dicFlx{(‑aði)}[1] \dicFlx{acc} \textbf{1.} \dicSynonym{fjörga} \dicDirectTranslationCS{oživit, oživovat, vzkřísit} \dicIndirectTranslationCS{(činit živějším)} \dicExampleIS{lífga eldinn} \dicExampleCS{oživit oheň};  \dicPhraseIS{lífga upp á e‑ð} \dicDirectTranslationCS{oživit\,/\addthin osvěžit (co) (pokoj barvami ap.)}  \textbf{2.} \dicSynonym*{vekja til lífs} \dicDirectTranslationCS{oživit, oživovat, vzkřísit, resuscitovat} \dicIndirectTranslationCS{(probudit člověka k~životu)};  \dicPhraseIS{lífga e‑n úr dauðadái} \dicDirectTranslationCS{probudit (koho) z~kómatu};  \dicPhraseIS{lífga e‑n við} \dicDirectTranslationCS{oživit (koho), probudit (koho) k~životu}
\dicEntry[lífgjafi] \dicTerm{líf··gjaf|i} \dicIPA{{l}{i}{v}{\r{\textObardotlessj}}{a}{v}{\textsci}} \dicPos{m}[1] \dicFlx{(‑a, ‑ar)}[8] \dicDirectTranslationCS{zachránce\,/\addthin zachránkyně (života)}
\dicEntry[lífgjöf] \dicTerm{líf··|gjöf} \dicIPA{{l}{i}{v}{\r{\textObardotlessj}}{\oe}{f}} \dicPos{f}[7] \dicFlx{(‑gjafar)}[19] \dicDirectTranslationCS{záchrana života}
\dicEntry[lífgun] \dicTerm{lífg|un} \dicIPA{{l}{i}{v}{\r{g}}{\textscy}{\textsubring{n}}} \dicPos{f}[7] \dicFlx{(‑unar, ‑anir)}[8] \dicDirectTranslationCS{resuscitace, reanimace, oživování}
\dicEntry[lífhimna] \dicTerm{líf··himn|a} \dicIPA{{l}{i}{\textlengthmark}{f}{h}{\textsci}{m}{n}{a}} \dicPos{f}[1] \dicFlx{(‑u, ‑ur)}[7] \dicFieldCat{anat.} \dicDirectTranslationCS{pobřišnice}
\dicEntry[lífklukka] \dicTerm{líf··klukk|a} \dicIPA{{l}{i}{v}{k\smash{\textsuperscript{h}}}{l}{\textscy}{h}{\r{g}}{a}} \dicPos{f}[1] \dicFlx{(‑u, ‑ur)}[13] \dicDirectTranslationCS{biologické hodiny}
\dicEntry[líflaus] \dicTerm{líf··laus} \dicIPA{{l}{i}{v}{l}{\oe i}{s}} \dicPos{adj}[5]\dicFlx{}[-1] \textbf{1.} \dicDirectTranslationCS{(jsoucí) bez života, mrtvý}  \textbf{2.} \dicFieldCat{biol.} \dicLangCat{zast.} \dicSynonym{ólífrænn} \dicDirectTranslationCS{neživý}
\dicEntry[líflát] \dicTerm{líf··lát} \dicIPA{{l}{i}{v}{l}{au}{\textsubring{d}}} \dicPos{n}[2] \dicFlx{(‑s, ‑)}[5] \dicSynonym{aftaka\smash{\textsuperscript{1}}} \dicDirectTranslationCS{poprava} \dicExampleIS{hóta e‑m  líflát} \dicExampleCS{hrozit (komu) popravou}
\dicEntry[lífláta] \dicTerm{líf··|láta} \dicIPA{{l}{i}{v}{l}{au}{\textsubring{d}}{a}} \dicPos{v}[6] \dicFlx{(‑læt, ‑lét, ‑létum, ‑léti, ‑látið)}[87] \dicFlx{acc} \dicDirectTranslationCS{popravit} \dicExampleIS{lífláta e‑n} \dicExampleCS{popravit (koho)}
\dicEntry[líflátið] \dicTerm{líf··látið} \dicIPA{{l}{i}{v}{l}{au}{\textsubring{d}}{\textsci}{\texttheta}} \dicPos{v} \dicFlx{supin} \dicLink{lífláta}
\dicEntry[líflegur] \dicTerm{líf··legur} \dicIPA{{l}{i}{v}{l}{\textepsilon}{\textbabygamma}{\textscy}{\textsubring{r}}} \dicPos{adj}[1]\dicFlx{}[-8] \dicSynonym{fjörugur} \dicDirectTranslationCS{živý, svěží, plný života}
\dicEntry[líflét] \dicTerm{líf··lét} \dicIPA{{l}{i}{v}{l}{j}{\textepsilon}{\textsubring{d}}} \dicPos{v} \dicFlx{ind pf sg 1 pers} \dicLink{lífláta}
\dicEntry[lífléti] \dicTerm{líf··léti} \dicIPA{{l}{i}{v}{l}{j}{\textepsilon}{\textsubring{d}}{\textsci}} \dicPos{v} \dicFlx{con pf sg 1 pers} \dicLink{lífláta}
\dicEntry[líflétum] \dicTerm{líf··létum} \dicIPA{{l}{i}{v}{l}{j}{\textepsilon}{\textsubring{d}}{\textscy}{\textsubring{m}}} \dicPos{v} \dicFlx{ind pf pl 1 pers} \dicLink{lífláta}
\dicEntry[líflæt] \dicTerm{líf··læt} \dicIPA{{l}{i}{v}{l}{a}{i}{\textsubring{d}}} \dicPos{v} \dicFlx{ind praes sg 1 pers} \dicLink{lífláta}
\dicEntry[lífríki] \dicTerm{líf··ríki} \dicIPA{{l}{i}{v}{r}{i}{\r{\textObardotlessj}}{\textsci}} \dicPos{n}[2] \dicFlx{(‑s, ‑)}[16] \dicDirectTranslationCS{biosystém, biologický systém}
\dicEntry[lífrænn] \dicTerm{líf··rænn} \dicIPA{{l}{i}{v}{r}{a}{i}{\textsubring{d}}{\textsubring{n}}} \dicPos{adj}[7]\dicFlx{}[-1] \textbf{1.} \dicDirectTranslationCS{živý, živoucí, rušný} \dicExampleIS{lífrænt starf} \dicExampleCS{živá práce}  \textbf{2.} \dicSynonym{náttúrulegur} \dicDirectTranslationCS{organický, ekologický} \dicExampleIS{lífrænn búskapur} \dicExampleCS{organické zemědělství} \dicAntonym{ólífrænn}
\dicEntry[lífsbarátta] \dicTerm{lífs··bar·átt|a} \dicIPA{{l}{i}{f}{s}{\textsubring{b}}{a}{r}{au}{h}{\textsubring{d}}{a}} \dicPos{f}[1] \dicFlx{(‑u)}[5] \dicDirectTranslationCS{existenční\,/\addthin životní boj}
\dicEntry[lífseigur] \dicTerm{líf··seigur} \dicIPA{{l}{i}{f}{s}{ei}{\textbabygamma}{\textscy}{\textsubring{r}}} \dicPos{adj}[1]\dicFlx{}[-1] \dicDirectTranslationCS{houževnatý, (jsoucí) s~tuhým kořínkem}
\dicEntry[lífsflótti] \dicTerm{lífs··flótt|i} \dicIPA{{l}{i}{f}{s}{f}{l}{ou}{h}{\textsubring{d}}{\textsci}} \dicPos{m}[1] \dicFlx{(‑a)}[3] \dicDirectTranslationCS{útěk před realitou, eskapismus}
\dicEntry[lífsglaður] \dicTerm{lífs··|glaður} \dicIPA{{l}{i}{f}{s}{\r{g}}{l}{a}{ð}{\textscy}{\textsubring{r}}} \dicPos{adj}[2] \dicFlx{(f ‑glöð)}[7] \dicSynonym{glaðlyndur} \dicDirectTranslationCS{radostný, šťastný, radující se ze života}
\dicEntry[lífsgæði] \dicTerm{lífs··gæði} \dicIPA{{l}{i}{f}{s}{\r{\textObardotlessj}}{a}{i}{ð}{\textsci}} \dicPos{n}[2] \dicFlx{pl}[19] \dicDirectTranslationCS{životní úroveň, kvalita života}
\dicEntry[lífsháski] \dicTerm{lífs··hásk|i} \dicIPA{{l}{i}{f}{s}{h}{au}{s}{\r{\textObardotlessj}}{\textsci}} \dicPos{m}[1] \dicFlx{(‑a)}[3] \dicDirectTranslationCS{ohrožení života}
\dicEntry[lífshætta] \dicTerm{lífs··hætt|a} \dicIPA{{l}{i}{f}{s}{h}{a}{i}{h}{\textsubring{d}}{a}} \dicPos{f}[1] \dicFlx{(‑u)}[5] \dicDirectTranslationCS{nebezpečí života}
\dicEntry[lífshættulegur] \dicTerm{lífs··hættu·legur} \dicIPA{{l}{i}{f}{s}{h}{a}{i}{h}{\textsubring{d}}{\textscy}{l}{\textepsilon}{\textbabygamma}{\textscy}{\textsubring{r}}} \dicPos{adj}[1]\dicFlx{}[-8] \dicDirectTranslationCS{životu nebezpečný}
\dicEntry[lífskjör] \dicTerm{lífs··kjör} \dicIPA{{l}{i}{f}{s}{c\smash{\textsuperscript{h}}}{\oe}{\textsubring{r}}} \dicPos{n}[2] \dicFlx{pl}[9] \dicDirectTranslationCS{životní podmínky}
\dicEntry[lífsleiði] \dicTerm{lífs··leið|i} \dicIPA{{l}{i}{f}{s}{l}{ei}{ð}{\textsci}} \dicPos{m}[1] \dicFlx{(‑a)}[3] \dicDirectTranslationCS{životní nuda, omrzelost životem}
\dicEntry[lífsleikni] \dicTerm{lífs··leikn|i} \dicIPA{{l}{i}{f}{s}{l}{ei}{h}{\r{g}}{n}{\textsci}} \dicPos{f}[3] \dicFlx{(‑i)}[3] \dicFieldCat{škol.} \dicDirectTranslationCS{praktická výchova} \dicIndirectTranslationCS{(předmět ve škole)}
\dicEntry[lífsnauðsyn] \dicTerm{lífs··nauð·syn} \dicIPA{{l}{i}{f}{s}{n}{\oe i}{ð}{s}{\textsci}{\textsubring{n}}} \dicPos{f}[4] \dicFlx{(‑jar, ‑jar)}[9] \dicDirectTranslationCS{životní nutnost}
\dicEntry[lífsorka] \dicTerm{lífs··ork|a} \dicIPA{{l}{i}{f}{s}{\textopeno}{\textsubring{r}}{\r{g}}{a}} \dicPos{f}[1] \dicFlx{(‑u)}[5] \dicDirectTranslationCS{vitalita, životní energie}
\dicEntry[lífsregla] \dicTerm{lífs··regl|a} \dicIPA{{l}{i}{f}{s}{r}{\textepsilon}{\r{g}}{l}{a}} \dicPos{f}[1] \dicFlx{(‑u, ‑ur)}[19] \dicDirectTranslationCS{životní zásada}
\dicEntry[lífsreynsla] \dicTerm{lífs··reynsl|a} \dicIPA{{l}{i}{f}{s}{r}{ei}{n}{s}{\textsubring{d}}{l}{a}} \dicPos{f}[1] \dicFlx{(‑u)}[5] \dicDirectTranslationCS{životní zkušenost}
\dicEntry[lífsskilyrði] \dicTerm{lífs··skil·yrði} \dicIPA{{l}{i}{f}{s}{\r{\textObardotlessj}}{\textsci}{l}{\textsci}{r}{ð}{\textsci}} \dicPos{n}[2] \dicFlx{pl}[19] \dicDirectTranslationCS{životní podmínky}
\dicEntry[lífsskoðun] \dicTerm{lífs··skoð|un} \dicIPA{{l}{i}{f}{s}{\r{g}}{\textopeno}{ð}{\textscy}{\textsubring{n}}} \dicPos{f}[7] \dicFlx{(‑unar, ‑anir)}[8] \dicDirectTranslationCS{životní postoj\,/\addthin názor}
\dicEntry[lífsspeki] \dicTerm{lífs··spek|i} \dicIPA{{l}{i}{f}{s}{\textsubring{b}}{\textepsilon}{\r{\textObardotlessj}}{\textsci}} \dicPos{f}[3] \dicFlx{(‑i)}[3] \dicDirectTranslationCS{životní moudrost}
\dicEntry[lífsstarf] \dicTerm{lífs··|starf} \dicIPA{{l}{i}{f}{s}{\textsubring{d}}{a}{r}{f}} \dicPos{n}[2] \dicFlx{(‑starfs, ‑störf)}[8] \dicDirectTranslationCS{životní poslání}
\dicEntry[lífsstíll] \dicTerm{lífs··stíl|l} \dicIPA{{l}{i}{f}{s}{\textsubring{d}}{i}{\textsubring{d}}{\textsubring{l}}} \dicPos{m}[6] \dicFlx{(‑s)}[50] \dicDirectTranslationCS{životní styl}
\dicEntry[lífstykki] \dicTerm{líf··stykki} \dicIPA{{l}{i}{f}{s}{\textsubring{d}}{\textsci}{h}{\r{\textObardotlessj}}{\textsci}} \dicPos{n}[2] \dicFlx{(‑s, ‑)}[16] \dicDirectTranslationCS{korzet, šněrovačka}
\dicEntry[lífsveig] \dicTerm{lífs··veig} \dicIPA{{l}{i}{f}{s}{v}{ei}{x}} \dicPos{f}[4] \dicFlx{(‑ar, ‑ar)}[1] \dicDirectTranslationCS{elixír života}
\dicEntry[lífsviðurværi] \dicTerm{lífs··viður·væri} \dicIPA{{l}{i}{f}{s}{v}{\textsci}{ð}{\textscy}{r}{v}{a}{i}{r}{\textsci}} \dicPos{n}[2] \dicFlx{(‑s)}[20] \dicDirectTranslationCS{obživa (jídlo ap.)}
\dicEntry[lífsþægindi] \dicTerm{lífs··þæg·indi} \dicIPA{{l}{i}{f}{s}{\texttheta}{a}{i}{j}{\textsci}{n}{\textsubring{d}}{\textsci}} \dicPos{n}[2] \dicFlx{pl}[19] \dicDirectTranslationCS{životní komfort\,/\addthin pohodlí}
\dicEntry[líftóra] \dicTerm{líf··tór|a} \dicIPA{{l}{i}{f}{t\smash{\textsuperscript{h}}}{ou}{r}{a}} \dicPos{f}[1] \dicFlx{(‑u)}[5] \dicDirectTranslationCS{jiskra života};  \dicPhraseIS{hræða úr e‑m líftóruna} \dicLangCat{přen.} \dicDirectTranslationCS{vyděsit (koho) k~smrti}
\dicEntry[líftrygging] \dicTerm{líf··trygg·ing} \dicIPA{{l}{i}{f}{t\smash{\textsuperscript{h}}}{r}{\textsci}{\r{\textObardotlessj}}{i}{\ng}{\r{g}}} \dicPos{f}[4] \dicFlx{(‑ar, ‑ar)}[5] \dicDirectTranslationCS{životní pojištění}
\dicEntry[líftryggja] \textls[15]{\dicTerm{líf··trygg|ja} \dicIPA{{l}{i}{f}{t\smash{\textsuperscript{h}}}{r}{\textsci}{\r{\textObardotlessj}}{a}} \dicPos{v}[2] \dicFlx{(‑ði, ‑t)}[92] \dicFlx{acc} \dicDirectTranslationCS{uzavřít životní pojiš\nobreak tění}}
\dicEntry[lífvera] \dicTerm{líf··ver|a} \dicIPA{{l}{i}{v}{\textlengthmark}{\textepsilon}{r}{a}} \dicPos{f}[1] \dicFlx{(‑u, ‑ur)}[7] \dicDirectTranslationCS{živá bytost, živý organismus}
\dicEntry[lífvænlegur] \dicTerm{líf··væn·legur} \dicIPA{{l}{i}{v}{\textlengthmark}{a}{i}{n}{l}{\textepsilon}{\textbabygamma}{\textscy}{\textsubring{r}}} \dicPos{adj}[1]\dicFlx{}[-8] \textbf{1.} \dicDirectTranslationCS{rentabilní} \dicExampleIS{lífvænleg atvinna} \dicExampleCS{rentabilní zaměstnání}  \textbf{2.} \dicDirectTranslationCS{příznivý pro život} \dicExampleIS{lífvænlegur staður} \dicExampleCS{místo příznivé pro život}
\begin{xtolerant}{}{1pt}
\dicEntry[lífvörður] \dicTerm{líf··|vörður}\addthinS\dicIPA{{l}{i}{v}{\textlengthmark}{\oe}{r}{ð}{\textscy}{\textsubring{r}}}\addthinS\dicPos{m}[11]\addthinS\dicFlx{(‑varðar, ‑verðir)}[5] \dicDirectTranslationCS{osobní strážce\,/\addthin strážkyně, tělesná stráž, bodyguard}
\end{xtolerant}
\dicEntry[lík] \dicTerm{lík} \dicsymFrequent\  \dicIPA{{l}{i}{\textlengthmark}{\r{g}}} \dicPos{n}[2] \dicFlx{(‑s, ‑)}[5] \textbf{1.} \dicSynonym*{dauður líkami} \dicDirectTranslationCS{mrtvola} \dicExampleIS{líkið af e‑m} \dicExampleCS{(čí) mrtvola}  \textbf{2.} \dicLangCat{hovor.} \dicSynonym*{tóm vínflaska} \dicDirectTranslationCS{prázdná láhev (po alkoholu)}
\dicEntry[líka] \dicTerm{lík|a\smash{\textsuperscript{1}}} \dicsymFrequent\  \dicIPA{{l}{i}{\textlengthmark}{\r{g}}{a}} \dicPos{v}[1] \dicFlx{(‑aði)}[82] \dicFlx{impers} \dicPhraseIS{e‑m líkar e‑að} \dicDirectTranslationCS{(co) se (komu) líbí, (co) má (kdo) rád} \dicExampleIS{Mér líkar það ekki.} \dicExampleCS{Mně se to nelíbí.};  \dicPhraseIS{e‑m líkar vel við e‑n} \dicDirectTranslationCS{(kdo) má (koho) rád}
\dicEntry[líka] \dicTerm{líka\smash{\textsuperscript{2}}} \dicsymFrequent\  \dicIPA{{l}{i}{\textlengthmark}{\r{g}}{a}} \dicPos{adv} \dicSynonym{einnig} \dicDirectTranslationCS{také, taky, rovněž} \dicExampleIS{Dagblaðið er líka á netinu.} \dicExampleCS{Deník je také na internetu.}
\dicEntry[líkami] \dicTerm{líkam|i} \dicsymFrequent\  \dicIPA{{l}{i}{\textlengthmark}{\r{g}}{a}{m}{\textsci}} \dicPos{m}[1] \dicFlx{(‑a, ‑ar)}[8] \dicDirectTranslationCS{tělo} \dicExampleIS{æfingar fyrir líkama} \dicExampleCS{cvičení pro tělo}
\dicEntry[líkamlegur] \dicTerm{líkam··legur} \dicsymFrequent\  \dicIPA{{l}{i}{\r{g}}{a}{m}{l}{\textepsilon}{\textbabygamma}{\textscy}{\textsubring{r}}} \dicPos{adj}[1]\dicFlx{}[-8] \dicDirectTranslationCS{tělesný, fyzický} \dicExampleIS{líkamleg atvinna} \dicExampleCS{fyzická práce}
\dicEntry[líkamsárás] \dicTerm{líkams··á·rás} \dicIPA{{l}{i}{\textlengthmark}{\r{g}}{a}{m}{s}{au}{r}{au}{s}} \dicPos{f}[7] \dicFlx{(‑ar, ‑ir)}[1] \dicDirectTranslationCS{tělesný útok, fyzické napadení}
\dicEntry[líkamsástand] \dicTerm{líkams··á·stand} \dicIPA{{l}{i}{\textlengthmark}{\r{g}}{a}{m}{s}{au}{s}{\textsubring{d}}{a}{n}{\textsubring{d}}} \dicPos{n}[2] \dicFlx{(‑s)}[2] \dicDirectTranslationCS{tělesný stav}
\dicEntry[líkamsburðir] \dicTerm{líkams··burðir} \dicIPA{{l}{i}{\textlengthmark}{\r{g}}{a}{m}{s}{\textsubring{b}}{\textscy}{r}{ð}{\textsci}{\textsubring{r}}} \dicPos{m}[10] \dicFlx{pl}[1] \dicDirectTranslationCS{tělesná\,/\addthin fyzická síla}
\dicEntry[líkamshiti] \dicTerm{líkams··hit|i} \dicIPA{{l}{i}{\textlengthmark}{\r{g}}{a}{m}{s}{h}{\textsci}{\textsubring{d}}{\textsci}} \dicPos{m}[1] \dicFlx{(‑a)}[3] \dicDirectTranslationCS{tělesná teplota}
\dicEntry[líkamsrækt] \dicTerm{líkams··rækt} \dicIPA{{l}{i}{\textlengthmark}{\r{g}}{a}{m}{s}{r}{a}{i}{x}{\textsubring{d}}} \dicPos{f}[4] \dicFlx{(‑ar)}[3] \dicDirectTranslationCS{tělesné cvičení, posilování, fitness}
\dicEntry[líkamsræktarstöð] \dicTerm{líkams·ræktar··stöð} \dicIPA{{l}{i}{\textlengthmark}{\r{g}}{a}{m}{s}{r}{a}{i}{x}{\textsubring{d}}{a}{\textsubring{r}}{s}{\textsubring{d}}{\oe}{\texttheta}} \dicPos{f}[6] \dicFlx{(‑var, ‑var)}[1] \dicDirectTranslationCS{posilovna, fitness centrum}
\dicEntry[líkamsstærð] \dicTerm{líkams··stærð} \dicIPA{{l}{i}{\textlengthmark}{\r{g}}{a}{m}{s}{\textsubring{d}}{a}{i}{r}{\texttheta}} \dicPos{f}[7] \dicFlx{(‑ar)}[3] \dicDirectTranslationCS{životní velikost} \dicExampleIS{í fullri líkamsstærð} \dicExampleCS{v~životní velikosti}
\dicEntry[líkan] \dicTerm{lík|an} \dicIPA{{l}{i}{\textlengthmark}{\r{g}}{a}{\textsubring{n}}} \dicPos{n}[2] \dicFlx{(‑ans, ‑ön)}[8] \dicDirectTranslationCS{model, maketa} \dicIndirectTranslationCS{(zmenšené provedení předmětu)} \dicExampleIS{líkan af húsi} \dicExampleCS{model domu}
\dicEntry[líkbrennsla] \dicTerm{lík··brennsl|a} \dicIPA{{l}{i}{\textlengthmark}{\r{g}}{\textsubring{b}}{r}{\textepsilon}{n}{s}{\textsubring{d}}{l}{a}} \dicPos{f}[1] \dicFlx{(‑u, ‑ur)}[19] \dicSynonym{bálför} \dicDirectTranslationCS{kremace}
\dicEntry[líkbörur] \dicTerm{lík··börur} \dicIPA{{l}{i}{\textlengthmark}{\r{g}}{\textsubring{b}}{\oe}{r}{\textscy}{\textsubring{r}}} \dicPos{f}[12] \dicFlx{pl}[4] \dicDirectTranslationCS{máry}
\dicEntry[líkfylgd] \dicTerm{lík··fylgd} \dicIPA{{l}{i}{\textlengthmark}{\r{g}}{f}{\textsci}{l}{\textsubring{d}}} \dicPos{f}[7] \dicFlx{(‑ar, ‑ir)}[1] \dicDirectTranslationCS{pohřební průvod}
\dicEntry[líkhús] \dicTerm{lík··hús} \dicIPA{{l}{i}{\textlengthmark}{\r{g}}{h}{u}{s}} \dicPos{n}[2] \dicFlx{(‑s, ‑)}[5] \dicDirectTranslationCS{márnice}
\dicEntry[líki] \dicTerm{lík|i\smash{\textsuperscript{1}}} \dicIPA{{l}{i}{\textlengthmark}{\r{\textObardotlessj}}{\textsci}} \dicPos{m}[1] \dicFlx{(‑a, ‑ar)}[1] \dicSynonym{jafningi} \dicDirectTranslationCS{sobě rovný, podobný (člověk)} \dicExampleIS{hann og hans líkar} \dicExampleCS{on a~jemu podobní};  \dicPhraseIS{eiga engan sinn líka} \dicDirectTranslationCS{nemít sobě rovného}
\dicEntry[líki] \dicTerm{líki\smash{\textsuperscript{2}}} \dicIPA{{l}{i}{\textlengthmark}{\r{\textObardotlessj}}{\textsci}} \dicPos{n}[2] \dicFlx{(‑s, ‑)}[16] \dicSynonym{sköpulag} \dicDirectTranslationCS{podoba, forma};  \dicPhraseIS{í heilu líki} \dicFlx{adj} \dicDirectTranslationCS{celý, netknutý}
\dicEntry[líkindi] \dicTerm{lík··indi} \dicsymFrequent\  \dicIPA{{l}{i}{\textlengthmark}{\r{\textObardotlessj}}{\textsci}{n}{\textsubring{d}}{\textsci}} \dicPos{n}[2] \dicFlx{pl}[19] \dicDirectTranslationCS{pravděpodobnost} \dicExampleIS{Það eru mikil líkindi til þess að \dots{}} \dicExampleCS{Je velká pravděpodobnost, že\dots{}};  \dicPhraseIS{að öllum líkindum} \dicFlx{adv} \dicDirectTranslationCS{s~největší pravděpodobností}
\dicEntry[líking] \dicTerm{lík··ing} \dicIPA{{l}{i}{\textlengthmark}{\r{\textObardotlessj}}{i}{\ng}{\r{g}}} \dicPos{f}[4] \dicFlx{(‑ar, ‑ar)}[5] \textbf{1.} \dicSynonym{líki\smash{\textsuperscript{2}}} \dicDirectTranslationCS{podoba, podobnost, obdoba, analogie}  \textbf{2.} \dicDirectTranslationCS{přirovnání, podobenství} \dicExampleIS{tala í líkingum} \dicExampleCS{mluvit v~přirovnáních}  \textbf{3.} \dicFieldCat{lit.} \dicSynonym{samlíking} \dicDirectTranslationCS{metafora} \dicExampleIS{líkingar í skáldskap} \dicExampleCS{metafora v~literatuře}  \textbf{4.} \dicFieldCat{mat.} \dicSynonym{jafna\smash{\textsuperscript{1}}} \dicDirectTranslationCS{rovnost}
\dicEntry[líkja] \dicTerm{lík|ja} \dicsymFrequent\  \dicIPA{{l}{i}{\textlengthmark}{\r{\textObardotlessj}}{a}} \dicPos{v}[2] \dicFlx{(‑ti, ‑t)}[25] \dicFlx{dat} \textbf{1.} \dicDirectTranslationCS{přirovnat, přirovnávat};  \dicPhraseIS{líkja e‑u saman} \dicDirectTranslationCS{srovnávat (co)};  \dicPhraseIS{líkja e‑u við e‑ð} \dicDirectTranslationCS{přirovnávat (co) k~(čemu)} \dicExampleIS{líkja ævi mannsins við blóm} \dicExampleCS{přirovnávat život člověka ke květině}  \textbf{2.} \dicDirectTranslationCS{napodobit, napodobovat};  \dicPhraseIS{líkja eftir e‑u} \dicDirectTranslationCS{napodobit (co), imitovat (co)} \dicExampleIS{líkja eftir hljóði fuglsins} \dicExampleCS{napodobovat zvuk ptáčka};  \dicIdiom{líkjast}{ \dicPhraseIS{líkjast e‑m}} \dicFlx{refl} \dicDirectTranslationCS{připomínat (koho), podobat se (komu)} \dicExampleIS{Hverjum líkist ég?} \dicExampleCS{Komu se podobám?}
\dicEntry[líkjör] \dicTerm{líkjör} \dicIPA{{l}{i}{\textlengthmark}{\r{\textObardotlessj}}{\oe}{\textsubring{r}}} \dicPos{m}[4] \dicFlx{(‑s, ‑ar)}[14] \dicDirectTranslationCS{likér}
\dicEntry[líkkista] \dicTerm{lík··kist|a} \dicIPA{{l}{i}{\textlengthmark}{\r{g}}{c\smash{\textsuperscript{h}}}{\textsci}{s}{\textsubring{d}}{a}} \dicPos{f}[1] \dicFlx{(‑u, ‑ur)}[19] \dicDirectTranslationCS{rakev}
\dicEntry[líkklæði] \dicTerm{lík··klæði} \dicIPA{{l}{i}{\textlengthmark}{\r{g}}{k\smash{\textsuperscript{h}}}{l}{a}{i}{ð}{\textsci}} \dicPos{n}[2] \dicFlx{pl}[19] \dicDirectTranslationCS{pohřební roucho}
\dicEntry[líkl.] \dicTerm{líkl.} \dicPos{zkr} \dicPhraseIS{líklega} \dicFlx{adv} \dicDirectTranslationCS{pravděpodobně}
\dicEntry[líklega] \dicTerm{lík··leg|a} \dicsymFrequent\  \dicIPA{{l}{i}{\textlengthmark}{\r{g}}{l}{\textepsilon}{\textbabygamma}{a}} \dicPos{adv} \dicFlx{(comp ‑ar, sup ‑ast)} \dicSynonym{sennilega} \dicDirectTranslationCS{pravděpodobně, nejspíš, asi} \dicExampleIS{Þau segja líklega nei.} \dicExampleCS{Pravděpodobně řeknou ne.}
\dicEntry[líklegur] \dicTerm{lík··legur} \dicsymFrequent\  \dicIPA{{l}{i}{\textlengthmark}{\r{g}}{l}{\textepsilon}{\textbabygamma}{\textscy}{\textsubring{r}}} \dicPos{adj}[1]\dicFlx{}[-8] \textbf{1.} \dicSynonym{efnilegur} \dicDirectTranslationCS{slibný, nadějný}  \textbf{2.} \dicSynonym{sennilegur} \dicDirectTranslationCS{pravděpodobný, předpokládaný} \dicExampleIS{Þeir eru líklegir til að feta í sömu fótspor.} \dicExampleCS{Pravděpodobně budou kráčet ve stejných stopách.}
\dicEntry[líkn] \dicTerm{líkn} \dicIPA{{l}{i}{h}{\r{g}}{\textsubring{n}}} \dicPos{f}[7] \dicFlx{(‑ar, ‑ir)}[1] \textbf{1.} \dicSynonym{miskunn} \dicDirectTranslationCS{milost, smilování}  \textbf{2.} \dicSynonym{hjúkrun} \dicDirectTranslationCS{opatrování, péče}
\dicEntry[líkna] \dicTerm{líkn|a} \dicIPA{{l}{i}{h}{\r{g}}{n}{a}} \dicPos{v}[1] \dicFlx{(‑aði)}[1] \dicFlx{dat} \dicSynonym{hjúkra} \dicDirectTranslationCS{pečovat, poskytnout péči, opatrovat} \dicExampleIS{líkna e‑m} \dicExampleCS{pečovat o~(koho)}
\dicEntry[líknarbelgur] \dicTerm{líknar··belg|ur} \dicIPA{{l}{i}{h}{\r{g}}{n}{a}{r}{\textsubring{b}}{\textepsilon}{l}{\r{g}}{\textscy}{\textsubring{r}}} \dicPos{m}[9] \dicFlx{(‑s\,/\addthin ‑jar, ‑ir)}[26] \textbf{1.} \dicFieldCat{anat.} \dicSynonym*{fósturfylgja} \dicDirectTranslationCS{amnion, vnitřní plodový obal}  \textbf{2.} \dicSynonym*{loftpúði í bíl} \dicDirectTranslationCS{airbag} \dicIndirectTranslationCS{(zařízení pasivní bezpečnosti v~autě)}
\dicEntry[líknardeild] \dicTerm{líknar··deild} \dicIPA{{l}{i}{h}{\r{g}}{n}{a}{r}{\textsubring{d}}{ei}{l}{\textsubring{d}}} \dicPos{f}[7] \dicFlx{(‑ar, ‑ir)}[1] \dicDirectTranslationCS{hospic}
\dicEntry[líknardráp] \dicTerm{líknar··dráp} \dicIPA{{l}{i}{h}{\r{g}}{n}{a}{r}{\textsubring{d}}{r}{au}{\textsubring{b}}} \dicPos{n}[2] \dicFlx{(‑s, ‑)}[5] \dicDirectTranslationCS{eutanazie}
\dicEntry[líknarmál] \dicTerm{líknar··mál} \dicIPA{{l}{i}{h}{\r{g}}{n}{a}{r}{m}{au}{\textsubring{l}}} \dicPos{n}[2] \dicFlx{(‑s, ‑)}[5] \dicDirectTranslationCS{charita, dobročinnost}
\dicEntry[líknarstofnun] \dicTerm{líknar··stofn|un} \dicIPA{{l}{i}{h}{\r{g}}{n}{a}{\textsubring{r}}{s}{\textsubring{d}}{\textopeno}{\textsubring{b}}{n}{\textscy}{\textsubring{n}}} \dicPos{f}[7] \dicFlx{(‑unar, ‑anir)}[8] \dicSynonym{góðgerðastofnun} \dicDirectTranslationCS{dobročinná\,/\addthin charitativní organizace}
\dicEntry[líkneski] \dicTerm{lík··neski} \dicIPA{{l}{i}{h}{\r{g}}{n}{\textepsilon}{s}{\r{\textObardotlessj}}{\textsci}} \dicPos{n}[2] \dicFlx{(‑s, ‑)}[16] \dicSynonym{höggmynd} \dicDirectTranslationCS{podobizna}
\dicEntry[líknsamur] \dicTerm{líkn··|samur} \dicIPA{{l}{i}{h}{\r{g}}{\textsubring{n}}{s}{a}{m}{\textscy}{\textsubring{r}}} \dicPos{adj}[1] \dicFlx{(f ‑söm)}[2] \dicDirectTranslationCS{milosrdný, milostivý, soucitný}
\dicEntry[líkræða] \dicTerm{lík··ræð|a} \dicIPA{{l}{i}{\textlengthmark}{\r{g}}{r}{a}{i}{ð}{a}} \dicPos{f}[1] \dicFlx{(‑u, ‑ur)}[19] \dicDirectTranslationCS{smuteční proslov}
\dicEntry[líkskoðun] \dicTerm{lík··skoð|un} \dicIPA{{l}{i}{\textlengthmark}{\r{g}}{s}{\r{g}}{\textopeno}{ð}{\textscy}{\textsubring{n}}} \dicPos{f}[7] \dicFlx{(‑unar)}[9] \dicDirectTranslationCS{ohledání mrtvoly}
\dicEntry[líkskurður] \dicTerm{lík··skurð|ur} \dicIPA{{l}{i}{\textlengthmark}{\r{g}}{s}{\r{g}}{\textscy}{r}{ð}{\textscy}{\textsubring{r}}} \dicPos{m}[10] \dicFlx{(‑ar, ‑ir)}[4] \dicDirectTranslationCS{pitva}
\dicEntry[líkur] \dicTerm{líkur\smash{\textsuperscript{1}}} \dicsymFrequent\  \dicIPA{{l}{i}{\textlengthmark}{\r{g}}{\textscy}{\textsubring{r}}} \dicPos{f}[12] \dicFlx{pl}[5] \dicSynonym{líkindi} \dicDirectTranslationCS{šance, pravděpodobnost} \dicExampleIS{líkur á e‑u\,/\addthin líkur til e‑s} \dicExampleCS{šance na (co)};  \dicPhraseIS{ef að líkum lætur} \dicFlx{adv} \dicDirectTranslationCS{pravděpodobně, patrně}
\dicEntry[líkur] \dicTerm{líkur\smash{\textsuperscript{2}}} \dicsymFrequent\  \dicIPA{{l}{i}{\textlengthmark}{\r{g}}{\textscy}{\textsubring{r}}} \dicPos{adj}[1]\dicFlx{}[-1] \textbf{1.} \dicDirectTranslationCS{podobný} \dicExampleIS{lík vandamál} \dicExampleCS{podobné problémy};  \dicPhraseIS{vera líkur e‑m} \dicDirectTranslationCS{být (komu) podobný, podobat se (komu)}  \textbf{2.} \dicDirectTranslationCS{obdobný, analogický};  \dicIdiom{líkur}{ \dicPhraseIS{líkt og}} \dicFlx{conj} \dicDirectTranslationCS{jak, jako} \dicExampleIS{líta út líkt og Íslendingur} \dicExampleCS{vypadat jako Islanďan}
\dicEntry[líkþorn] \dicTerm{lík··þorn} \dicIPA{{l}{i}{\textlengthmark}{\r{g}}{\texttheta}{\textopeno}{r}{\textsubring{d}}{\textsubring{n}}} \dicPos{n}[2] \dicFlx{(‑s, ‑)}[5] \dicDirectTranslationCS{kuří oko}
\dicEntry[líkþrá] \dicTerm{lík··þrá} \dicIPA{{l}{i}{\textlengthmark}{\r{g}}{\texttheta}{r}{au}} \dicPos{f}[4] \dicFlx{(‑r)}[19] \dicFieldCat{med.} \dicDirectTranslationCS{lepra, malomocenství}
\dicEntry[líkþrár] \dicTerm{lík··þrár} \dicIPA{{l}{i}{\textlengthmark}{\r{g}}{\texttheta}{r}{au}{\textsubring{r}}} \dicPos{adj}[4]\dicFlx{}[-3] \dicDirectTranslationCS{malomocný}
\dicEntry[lím] \dicTerm{lím} \dicIPA{{l}{i}{\textlengthmark}{\textsubring{m}}} \dicPos{n}[2] \dicFlx{(‑s, ‑)}[5] \dicDirectTranslationCS{lepidlo, lep, klih}
\dicEntry[líma] \dicTerm{lím|a} \dicsymFrequent\  \dicIPA{{l}{i}{\textlengthmark}{m}{a}} \dicPos{v}[2] \dicFlx{(‑di, ‑t)}[141] \dicFlx{acc} \dicDirectTranslationCS{(na)lepit, nalepovat, přilepit, přilepovat} \dicExampleIS{líma mynd á blað} \dicExampleCS{nalepit obrázek na papír};  \dicIdiom{límast}{ \dicPhraseIS{límast}} \dicFlx{refl} \dicDirectTranslationCS{(při)lepit se}
\dicEntry[límband] \dicTerm{lím··|band} \dicIPA{{l}{i}{m}{\textsubring{b}}{a}{n}{\textsubring{d}}} \dicPos{n}[2] \dicFlx{(‑bands, ‑bönd)}[8] \dicDirectTranslationCS{lepicí páska, izolepa}
\dicEntry[límmiði] \dicTerm{lím··mið|i} \dicIPA{{l}{i}{m}{\textlengthmark}{\textsci}{ð}{\textsci}} \dicPos{m}[1] \dicFlx{(‑a, ‑ar)}[1] \dicDirectTranslationCS{nálepka, štítek, etiketa}
\dicEntry[lín] \dicTerm{lín} \dicIPA{{l}{i}{\textlengthmark}{\textsubring{n}}} \dicPos{n}[2] \dicFlx{(‑s)}[2] \textbf{1.} \dicFieldCat{bot.} \dicSynonym{hör} \dicDirectTranslationCS{len setý} \textit{(l.~{\textLA{Linum usitatissimum}})}  \dicsymPhoto\   \textbf{2.} \dicSynonym{léreft} \dicDirectTranslationCS{len, lněné plátno} \dicExampleIS{spinna lín} \dicExampleCS{tkát len}
\dicFigure{ds_image_lin_0_2.jpg}{Lín}{Lín - Franz Eugen Kohler, PD}
\dicEntry[lína] \dicTerm{lín|a} \dicsymFrequent\  \dicIPA{{l}{i}{\textlengthmark}{n}{a}} \dicPos{f}[1] \dicFlx{(‑u, ‑ur)}[7] \textbf{1.} \dicSynonym{strik} \dicDirectTranslationCS{čára, linka, linie}  \textbf{2.} \dicSynonym*{textalína} \dicDirectTranslationCS{řádek, řádka} \dicExampleIS{nokkrar línur úr textanum} \dicExampleCS{několik řádek z~textu}  \textbf{3.} \dicSynonym{band} \dicDirectTranslationCS{lano, šňůra, provázek}  \textbf{4.} \dicSynonym{lóð\smash{\textsuperscript{2}}} \dicDirectTranslationCS{lanová síť}  \textbf{5.} \dicSynonym{stefna\smash{\textsuperscript{1}}} \dicDirectTranslationCS{politický směr, politická linie}  \textbf{6.} \dicFieldCat{mat.} \dicDirectTranslationCS{přímka};  \dicIdiom{lína}{ \dicPhraseIS{lesa milli línanna}} \dicLangCat{přen.} \dicDirectTranslationCS{číst mezi řádky} \dicIndirectTranslationCS{(porozumět více, než je napsáno)}; { \dicPhraseIS{passa upp á línurnar}} \dicLangCat{přen.} \dicDirectTranslationCS{dávat pozor na linii} \dicIndirectTranslationCS{(snažit se neztloustnout)}
\dicEntry[línolía] \dicTerm{lín··olí|a} \dicIPA{{l}{i}{\textlengthmark}{n}{\textopeno}{l}{i}{j}{a}} \dicPos{f}[1] \dicFlx{(‑u)}[5] \dicDirectTranslationCS{lněný olej}
\dicEntry[línubil] \dicTerm{línu··bil} \dicIPA{{l}{i}{\textlengthmark}{n}{\textscy}{\textsubring{b}}{\textsci}{\textsubring{l}}} \dicPos{n}[2] \dicFlx{(‑s, ‑)}[5] \dicFieldCat{poč.} \dicDirectTranslationCS{řádkování}
\dicEntry[línumaður] \dicTerm{línu··|maður} \dicIPA{{l}{i}{\textlengthmark}{n}{\textscy}{m}{a}{ð}{\textscy}{\textsubring{r}}} \dicPos{m}[13] \dicFlx{(‑manns, ‑menn)}[2] \dicFieldCat{sport.} \dicDirectTranslationCS{krajní útočník\,/\addthin útočnice, křídlo}
\dicEntry[línurit] \dicTerm{línu··rit} \dicIPA{{l}{i}{\textlengthmark}{n}{\textscy}{r}{\textsci}{\textsubring{d}}} \dicPos{n}[2] \dicFlx{(‑s, ‑)}[5] \dicDirectTranslationCS{graf, diagram} \dicExampleIS{töflur og línurit} \dicExampleCS{tabulky a~grafy}
\dicEntry[línuskauti] \dicTerm{línu··skaut|i} \dicIPA{{l}{i}{\textlengthmark}{n}{\textscy}{s}{\r{g}}{\oe i}{\textsubring{d}}{\textsci}} \dicPos{m}[1] \dicFlx{(‑a, ‑ar)}[1] \dicPhraseIS{línuskautar} \dicFlx{pl} \dicFieldCat{sport.} \dicDirectTranslationCS{kolečkové brusle}
\dicEntry[línuskil] \dicTerm{línu··skil} \dicIPA{{l}{i}{\textlengthmark}{n}{\textscy}{s}{\r{\textObardotlessj}}{\textsci}{\textsubring{l}}} \dicPos{n}[2] \dicFlx{pl}[1] \dicFieldCat{poč.} \dicDirectTranslationCS{zalomení řádku}
\dicEntry[líra] \dicTerm{lír|a} \dicIPA{{l}{i}{\textlengthmark}{r}{a}} \dicPos{f}[1] \dicFlx{(‑u, ‑ur)}[7] \textbf{1.} \dicDirectTranslationCS{lira} \dicIndirectTranslationCS{(bývalá italská měna)}  \textbf{2.} \dicFieldCat{hud.} \dicDirectTranslationCS{lyra}
\dicEntry[lít] \dicTerm{lít} \dicIPA{{l}{i}{\textlengthmark}{\textsubring{d}}} \dicPos{v} \dicFlx{ind praes sg 1 pers} \dicLink{líta}
\dicEntry[líta] \dicTerm{líta} \dicsymFrequent\  \dicIPA{{l}{i}{\textlengthmark}{\textsubring{d}}{a}} \dicPos{v}[6] \dicFlx{(lít, leit, litum, liti, litið)}[76] \dicFlx{acc} \dicDirectTranslationCS{(po)dívat se, hledět, pohlédnout, pohlížet} \dicExampleIS{líta á klukkuna} \dicExampleCS{podívat se na hodiny};  \dicPhraseIS{líta e‑n} \dicDirectTranslationCS{pohlédnout na (koho)};  \dicPhraseIS{líta á e‑n} \dicDirectTranslationCS{podívat se na (koho)} \dicExampleIS{líta spyrjandi á hana} \dicExampleCS{podívat se na ni tázavě};  \dicIdiom{líta}[af]{ \dicPhraseIS{líta af e‑u\,/\addthin e‑m}} \dicDirectTranslationCS{moct spustit oči z~(čeho\,/\addthin koho)} \dicIndirectTranslationCS{(většinou v~záporu)};  \dicIdiom{líta}[á]{ \dicPhraseIS{líta á e‑ð}} \dicDirectTranslationCS{podívat se na (co), kouknout na (co), prozkoumat (co)} \dicExampleIS{líta á vélina} \dicExampleCS{kouknout na motor};  \dicPhraseIS{líta stórt á sig} \dicDirectTranslationCS{mít o~sobě vysoké mínění};  \dicIdiom{líta}[eftir]{ \dicPhraseIS{líta eftir e‑m}} \dicDirectTranslationCS{dohlédnout\,/\addthin dohlížet na (koho)} \dicExampleIS{líta eftir barninu} \dicExampleCS{dohlížet na dítě};  \dicIdiom{líta}[fram hjá]{ \dicPhraseIS{líta fram hjá e‑u}} \dicDirectTranslationCS{přehlédnout\,/\addthin přehlížet (co)} \dicExampleIS{líta fram hjá mikilvægum staðreyndum} \dicExampleCS{přehlížet důležitá fakta};  \dicIdiom{líta}[inn]{ \dicPhraseIS{líta inn hjá e‑m}} \dicDirectTranslationCS{(na)kouknout ke (komu), (za)stavit se u~(koho) (na návštěvu ap.)};  \dicIdiom{líta}[í]{ \dicPhraseIS{líta í e‑ð}} \dicDirectTranslationCS{podívat se do (čeho), nahlédnout do (čeho) (novin ap.)};  \dicIdiom{líta}[niður]{ \dicPhraseIS{líta niður á e‑n}} \dicDirectTranslationCS{dívat se na (koho) spatra, shlížet na (koho) povýšeně};  \dicIdiom{líta}[til]{ \dicPhraseIS{líta til e‑rs}} \dicDirectTranslationCS{zastavit se u~(koho), podívat se ke (komu)};  \dicIdiom{líta}[undan]{ \dicPhraseIS{líta undan}} \dicDirectTranslationCS{sklopit zrak, podívat se dolů};  \dicIdiom{líta}[upp]{ \dicPhraseIS{líta upp til e‑rs}} \dicDirectTranslationCS{obdivovat (koho), vzhlížet ke (komu) (s~úctou ap.)};  \dicIdiom{líta}[út]{ \dicPhraseIS{líta út}} \dicDirectTranslationCS{vypadat, vyhlížet} \dicExampleIS{líta vel út} \dicExampleCS{vypadat dobře}; { \dicPhraseIS{líta út fyrir að (vera ríkur)}} \dicDirectTranslationCS{vypadat na (boháče), vypadat jako (bohatý)};  \dicIdiom{líta}[við]{ \dicPhraseIS{líta ekki við e‑u}} \dicDirectTranslationCS{netknout se (čeho), nejevit zájem o~(co)}; { \dicPhraseIS{líta við (þar)}} \dicDirectTranslationCS{stavit se (tam) na skok, zaskočit (tam)};  \dicIdiom{líta}[yfir]{ \dicPhraseIS{líta yfir e‑ð}} \dicDirectTranslationCS{projít (co), prohlédnout (co)} \dicExampleIS{líta yfir textann} \dicExampleCS{projít text};  \dicIdiom{lítast}[á]{ \dicPhraseIS{e‑m líst vel á e‑ð}} \dicFlx{refl impers} \dicDirectTranslationCS{(co) se (komu) líbí}
\dicEntry[lítið] \dicTerm{lítið} \dicsymFrequent\  \dicIPA{{l}{i}{\textlengthmark}{\textsubring{d}}{\textsci}{\texttheta}} \dicPos{adv} \dicFlx{(comp minna, sup minnst)} \dicDirectTranslationCS{málo, trochu} \dicExampleIS{vera lítið áberandi} \dicExampleCS{být trochu nápadný}
\dicEntry[lítilfjörlegur] \dicTerm{lítil··fjör·legur} \dicIPA{{l}{i}{\textlengthmark}{\textsubring{d}}{\textsci}{l}{f}{j}{\oe}{r}{l}{\textepsilon}{\textbabygamma}{\textscy}{\textsubring{r}}} \dicPos{adj}[1]\dicFlx{}[-8] \textbf{1.} \dicSynonym{ómerkilegur} \dicDirectTranslationCS{nezajímavý, nezáživný, všední}  \textbf{2.} \dicSynonym{lasinn} \dicDirectTranslationCS{indisponovaný, slabý}
\dicEntry[lítill] \dicTerm{lítill} \dicsymFrequent\  \dicIPA{{l}{i}{\textlengthmark}{\textsubring{d}}{\textsci}{\textsubring{d}}{\textsubring{l}}} \dicPos{adj}[11] \dicFlx{(comp minni, sup minnstur)}[10] \textbf{1.} \dicSynonym{smár} \dicDirectTranslationCS{malý, drobný} \dicExampleIS{lítill maður} \dicExampleCS{malý člověk} \dicAntonym{stór}  \textbf{2.} \dicSynonym*{þann yngri til aðgreiningar frá þeim eldri} \dicDirectTranslationCS{malý, mladší} \dicExampleIS{Sölvi litli} \dicExampleCS{malý Sölvi}  \textbf{3.} \dicSynonym{óverulegur} \dicDirectTranslationCS{malý, zanedbatelný, nepatrný}  \textbf{4.} \dicSynonym{óduglegur} \dicDirectTranslationCS{zahálčivý, málo snaživý} \dicExampleIS{lítill íþróttamaður} \dicExampleCS{malý sportovec};  \dicIdiom{lítill}{ \dicPhraseIS{vera lítill í sér}} \dicDirectTranslationCS{být zbabělý}
\dicEntry[lítillátur] \dicTerm{lítil··látur} \dicIPA{{l}{i}{\textlengthmark}{\textsubring{d}}{\textsci}{l}{au}{\textsubring{d}}{\textscy}{\textsubring{r}}} \dicPos{adj}[1]\dicFlx{}[-1] \textbf{1.} \dicSynonym*{hrokalaus} \dicDirectTranslationCS{pokorný}  \textbf{2.} \dicSynonym{auðmjúkur} \dicDirectTranslationCS{nenáročný, skromný}
\dicEntry[lítillækka] \dicTerm{lítil··lækk|a} \dicIPA{{l}{i}{\textlengthmark}{\textsubring{d}}{\textsci}{l}{a}{i}{h}{\r{g}}{a}} \dicPos{v}[1] \dicFlx{(‑aði)}[1] \dicFlx{acc} \dicDirectTranslationCS{pokořit, pokořovat, ponížit, ponižovat} \dicExampleIS{lítillækka sig} \dicExampleCS{ponížit se}
\dicEntry[lítillækkun] \dicTerm{lítil··lækk|un} \dicIPA{{l}{i}{\textlengthmark}{\textsubring{d}}{\textsci}{l}{a}{i}{h}{\r{g}}{\textscy}{\textsubring{n}}} \dicPos{f}[7] \dicFlx{(‑unar)}[9] \dicDirectTranslationCS{pokoření, ponížení}
\dicEntry[lítillæti] \dicTerm{lítil··læti} \dicIPA{{l}{i}{\textlengthmark}{\textsubring{d}}{\textsci}{l}{a}{i}{\textsubring{d}}{\textsci}} \dicPos{n}[2] \dicFlx{(‑s)}[20] \textbf{1.} \dicSynonym{yfirlætisleysi} \dicDirectTranslationCS{pokora}  \textbf{2.} \dicSynonym{auðmýkt} \dicDirectTranslationCS{skromnost}
\dicEntry[lítilmagni] \dicTerm{lítil··magn|i} \dicIPA{{l}{i}{\textsubring{d}}{\textsci}{l}{m}{a}{\r{g}}{n}{\textsci}} \dicPos{m}[1] \dicFlx{(‑a, ‑ar)}[8] \dicDirectTranslationCS{padavka, outsider(ka)}
\dicEntry[lítilmenni] \dicTerm{lítil··menni} \dicIPA{{l}{i}{\textlengthmark}{\textsubring{d}}{\textsci}{l}{m}{\textepsilon}{n}{\textsci}} \dicPos{n}[2] \dicFlx{(‑s, ‑)}[14] \textbf{1.} \dicDirectTranslationCS{malý člověk (vzrůstem)}  \textbf{2.} \dicDirectTranslationCS{malicherný\,/\addthin nanicovatý člověk}
\dicEntry[lítilmótlegur] \dicTerm{lítil··mót·legur} \dicIPA{{l}{i}{\textlengthmark}{\textsubring{d}}{\textsci}{l}{m}{ou}{\textsubring{d}}{l}{\textepsilon}{\textbabygamma}{\textscy}{\textsubring{r}}} \dicPos{adj}[1]\dicFlx{}[-8] \dicSynonym{ómerkilegur} \dicDirectTranslationCS{ubohý, mizerný, mrzký}
\dicEntry[lítilræði] \dicTerm{lítil··ræði} \dicIPA{{l}{i}{\textlengthmark}{\textsubring{d}}{\textsci}{l}{r}{a}{i}{ð}{\textsci}} \dicPos{n}[2] \dicFlx{(‑s)}[20] \dicDirectTranslationCS{maličkost, maličko}
\dicEntry[lítilsháttar] \dicTerm{lítils··háttar\smash{\textsuperscript{1}}} \dicIPA{{l}{i}{\textlengthmark}{\textsubring{d}}{\textsci}{l}{s}{h}{au}{h}{\textsubring{d}}{a}{\textsubring{r}}} \dicPos{adj}[13] \dicFlx{indecl}[1] \dicDirectTranslationCS{malý, drobný}
\dicEntry[lítilsháttar] \dicTerm{lítils··háttar\smash{\textsuperscript{2}}} \dicIPA{{l}{i}{\textlengthmark}{\textsubring{d}}{\textsci}{l}{s}{h}{au}{h}{\textsubring{d}}{a}{\textsubring{r}}} \dicPos{adv} \dicDirectTranslationCS{málo, trochu}
\dicEntry[lítilsigldur] \dicTerm{lítil··sigldur} \dicIPA{{l}{i}{\textlengthmark}{\textsubring{d}}{\textsci}{l}{s}{\textsci}{l}{\textsubring{d}}{\textscy}{\textsubring{r}}} \dicPos{adj}[2]\dicFlx{}[-14] \dicSynonym{lítilfjörlegur} \dicDirectTranslationCS{nevýznamný, ubohý, nicotný}
\dicEntry[lítilsvirða] \dicTerm{lítils··vir|ða} \dicIPA{{l}{i}{\textlengthmark}{\textsubring{d}}{\textsci}{l}{s}{v}{\textsci}{r}{ð}{a}} \dicPos{v}[2] \dicFlx{(‑ti, ‑t)}[52] \dicFlx{acc} \dicDirectTranslationCS{pohrdat, opovrhovat}
\dicEntry[lítilsvirðing] \dicTerm{lítils··virð·ing} \dicIPA{{l}{i}{\textlengthmark}{\textsubring{d}}{\textsci}{l}{s}{v}{\textsci}{r}{ð}{i}{\ng}{\r{g}}} \dicPos{f}[4] \dicFlx{(‑ar)}[7] \dicSynonym{óvirðing} \dicDirectTranslationCS{opovržení, pohrdání}
\dicEntry[lítilvægur] \dicTerm{lítil··vægur} \dicIPA{{l}{i}{\textlengthmark}{\textsubring{d}}{\textsci}{l}{v}{a}{i}{\textbabygamma}{\textscy}{\textsubring{r}}} \dicPos{adj}[1]\dicFlx{}[-1] \dicSynonym*{þýðingarlítill} \dicDirectTranslationCS{nedůležitý, nepodstatný, bezvýznamný}
\dicEntry[lítilþægur] \dicTerm{lítil··þægur} \dicIPA{{l}{i}{\textlengthmark}{\textsubring{d}}{\textsci}{l}{\texttheta}{a}{i}{\textbabygamma}{\textscy}{\textsubring{r}}} \dicPos{adj}[1]\dicFlx{}[-1] \dicDirectTranslationCS{skromný, nenáročný} \dicIndirectTranslationCS{(často používané v~záporu)}
\dicEntry[lítri] \dicTerm{lítr|i} \dicIPA{{l}{i}{\textlengthmark}{\textsubring{d}}{r}{\textsci}} \dicPos{m}[1] \dicFlx{(‑a, ‑ar)}[1] \dicDirectTranslationCS{litr}
\dicEntry[lítt] \dicTerm{lítt} \dicsymFrequent\  \dicIPA{{l}{i}{h}{\textsubring{d}}} \dicPos{adv} \dicFlx{(comp miður, sup minnst)} \textbf{1.} \dicSynonym{lítið} \dicDirectTranslationCS{trochu, málo} \dicExampleIS{lítt þekktur ráðherra} \dicExampleCS{málo známý ministr}  \textbf{2.} \dicSynonym{illa} \dicDirectTranslationCS{špatně}
\dicEntry[ljá] \dicTerm{ljá} \dicIPA{{l}{j}{au}{\textlengthmark}} \dicPos{v}[5] \dicFlx{(ljæ, léði, léðum, léði, léð)}[15] \dicFlx{dat + acc} \textbf{1.} \dicLangCat{zast.} \dicDirectTranslationCS{(vy)půjčit} \dicExampleIS{ljá e‑m e‑ð} \dicExampleCS{půjčit (komu co)}  \textbf{2.} \dicPhraseIS{ljá e‑m eyra} \dicLangCat{přen.} \dicDirectTranslationCS{vyslyšet (koho), dopřát (komu) sluchu}  \textbf{3.} \dicDirectTranslationCS{vnést, vnášet, přidat, přidávat} \dicIndirectTranslationCS{(způsobit uplatnění, začlenění něčeho)}  \textbf{4.} \dicPhraseIS{ljá máls á e‑u} \dicDirectTranslationCS{zavést téma na (co)}
\dicEntry[ljár] \dicTerm{ljá|r} \dicIPA{{l}{j}{au}{\textlengthmark}{\textsubring{r}}} \dicPos{m}[9] \dicFlx{(‑s, ‑ir)}[12] \dicDirectTranslationCS{kosa} \dicIndirectTranslationCS{(ruční nástroj)};  \dicPhraseIS{slá með orfi og ljá} \dicDirectTranslationCS{sekat kosou};  \dicPhraseIS{vera e‑m óþægur ljár í þúfu} \dicLangCat{přen.} \dicDirectTranslationCS{být (komu) trnem v~oku}
\dicEntry[ljóð] \dicTerm{ljóð} \dicsymFrequent\  \dicIPA{{l}{j}{ou}{\textlengthmark}{\texttheta}} \dicPos{n}[2] \dicFlx{(‑s, ‑)}[5] \dicDirectTranslationCS{báseň} \dicExampleIS{ljóð eftir e‑n} \dicExampleCS{(čí) báseň}
\dicEntry[ljóðabók] \dicTerm{ljóða··|bók} \dicIPA{{l}{j}{ou}{\textlengthmark}{ð}{a}{\textsubring{b}}{ou}{\r{g}}} \dicPos{f}[8] \dicFlx{(‑bókar, ‑bækur)}[5] \dicDirectTranslationCS{sbírka\,/\addthin kniha poezie}
\dicEntry[ljóðaháttur] \dicTerm{ljóða··|háttur} \dicIPA{{l}{j}{ou}{\textlengthmark}{ð}{a}{h}{au}{h}{\textsubring{d}}{\textscy}{\textsubring{r}}} \dicPos{m}[12] \dicFlx{(‑háttar, ‑hættir)}[7] \dicFieldCat{lit.} \dicDirectTranslationCS{ed\nobreak dické metrum (šestiverší)}
\dicEntry[ljóðlist] \dicTerm{ljóð··list} \dicIPA{{l}{j}{ou}{ð}{l}{\textsci}{s}{\textsubring{d}}} \dicPos{f}[7] \dicFlx{(‑ar)}[3] \dicDirectTranslationCS{básnictví, básnické umění}
\dicEntry[ljóðmæli] \dicTerm{ljóð··mæli} \dicIPA{{l}{j}{ou}{ð}{m}{a}{i}{l}{\textsci}} \dicPos{n}[2] \dicFlx{pl}[19] \dicDirectTranslationCS{sbírka poezie}
\dicEntry[ljóðræna] \dicTerm{ljóð··ræn|a} \dicIPA{{l}{j}{ou}{ð}{r}{a}{i}{n}{a}} \dicPos{f}[1] \dicFlx{(‑u)}[5] \dicFieldCat{lit.} \dicDirectTranslationCS{lyrika}
\dicEntry[ljóðrænn] \dicTerm{ljóð··rænn} \dicIPA{{l}{j}{ou}{ð}{r}{a}{i}{\textsubring{d}}{\textsubring{n}}} \dicPos{adj}[7]\dicFlx{}[-1] \dicDirectTranslationCS{poetický, lyrický}
\dicEntry[ljóðskáld] \dicTerm{ljóð··skáld} \dicIPA{{l}{j}{ou}{ð}{s}{\r{g}}{au}{l}{\textsubring{d}}} \dicPos{n}[2] \dicFlx{(‑s, ‑)}[5] \dicDirectTranslationCS{básník, básnířka, poeta, poetka}
\dicEntry[ljóður] \dicTerm{ljóð|ur} \dicIPA{{l}{j}{ou}{\textlengthmark}{ð}{\textscy}{\textsubring{r}}} \dicPos{m}[6] \dicFlx{(‑s)}[26] \dicSynonym{lýti} \dicDirectTranslationCS{kaz, vada}
\dicEntry[ljóma] \dicTerm{ljóm|a} \dicIPA{{l}{j}{ou}{\textlengthmark}{m}{a}} \dicPos{v}[1] \dicFlx{(‑aði)}[1] \dicDirectTranslationCS{zářit, vyzařovat} \dicExampleIS{Sólin ljómar á himninum.} \dicExampleCS{Slunce září na obloze.};  \dicPhraseIS{ljóma af gleði} \dicLangCat{přen.} \dicDirectTranslationCS{zářit štěstím}
\dicEntry[ljómandi] \dicTerm{ljóm··andi\smash{\textsuperscript{1}}} \dicIPA{{l}{j}{ou}{\textlengthmark}{m}{a}{n}{\textsubring{d}}{\textsci}} \dicPos{adj}[13] \dicFlx{indecl}[1] \dicDirectTranslationCS{zářivý, zářící} \dicExampleIS{ljómandi sól} \dicExampleCS{zářící slunce}
\dicEntry[ljómandi] \dicTerm{ljóm··andi\smash{\textsuperscript{2}}} \dicIPA{{l}{j}{ou}{\textlengthmark}{m}{a}{n}{\textsubring{d}}{\textsci}} \dicPos{adv} \dicDirectTranslationCS{úžasně, fantasticky} \dicExampleIS{ljómandi vel} \dicExampleCS{fantasticky dobře}
\dicEntry[ljómi] \dicTerm{ljóm|i} \dicIPA{{l}{j}{ou}{\textlengthmark}{m}{\textsci}} \dicPos{m}[1] \dicFlx{(‑a)}[3] \textbf{1.} \dicSynonym{birta\smash{\textsuperscript{1}}} \dicDirectTranslationCS{zář(e), svit}  \textbf{2.} \dicFieldCat{fyz.} \dicDirectTranslationCS{jas}
\dicEntry[ljón] \dicTerm{ljón} \dicsymFrequent\  \dicIPA{{l}{j}{ou}{\textlengthmark}{\textsubring{n}}} \dicPos{n}[2] \dicFlx{(‑s, ‑)}[5] \textbf{1.} \dicFieldCat{zool.} \dicDirectTranslationCS{lev} \textit{(l.~{\textLA{Panthera leo}})}  \dicsymPhoto\  \dicExampleIS{Ljón hleypur í átt að manni.} \dicExampleCS{Směrem k~člověku běží lev.}  \textbf{2.} \dicDirectTranslationCS{Lev} \dicIndirectTranslationCS{(znamení zvěrokruhu)} \textit{(l.~{\textLA{Leo}})}
\dicFigure{20786.jpg}{Ljón}{Ljón - Stansell, Ken, Biolib, PD}
\dicEntry[ljónatemjari] \dicTerm{ljóna··temj·ar|i} \dicIPA{{l}{j}{ou}{\textlengthmark}{n}{a}{t\smash{\textsuperscript{h}}}{\textepsilon}{m}{j}{a}{r}{\textsci}} \dicPos{m}[1] \dicFlx{(‑a, ‑ar)}[13] \dicDirectTranslationCS{krotitel(ka) lvů\,/\addthin šelem\,/\addthin dravé zvěře}
\dicEntry[ljónslappi] \dicTerm{ljóns··lapp|i} \dicIPA{{l}{j}{ou}{n}{s}{l}{a}{h}{\textsubring{b}}{\textsci}} \dicPos{m}[1] \dicFlx{(‑a)}[3] \dicFieldCat{bot.} \dicDirectTranslationCS{kontryhel alpinský} \textit{(l.~{\textLA{Alchemilla alpina}})}  \dicsymPhoto\ 
\dicFigure{88632.jpg}{Ljónslappi}{Ljónslappi - Pleva František, Biolib, PD}
\dicEntry[ljónynja] \dicTerm{ljón··ynj|a} \dicIPA{{l}{j}{ou}{\textlengthmark}{n}{\textsci}{n}{j}{a}} \dicPos{f}[1] \dicFlx{(‑u, ‑ur)}[7] \dicDirectTranslationCS{lvice}
\dicEntry[ljóri] \dicTerm{ljór|i} \dicIPA{{l}{j}{ou}{\textlengthmark}{r}{\textsci}} \dicPos{m}[1] \dicFlx{(‑a, ‑ar)}[1] \dicDirectTranslationCS{světlík}
\dicEntry[ljós] \dicTerm{ljós\smash{\textsuperscript{1}}} \dicsymFrequent\  \dicIPA{{l}{j}{ou}{\textlengthmark}{s}} \dicPos{n}[2] \dicFlx{(‑s, ‑)}[5] \textbf{1.} \dicSynonym{birta\smash{\textsuperscript{1}}} \dicDirectTranslationCS{světlo, záře, jas, svit} \dicExampleIS{ljósið frá lampanum} \dicExampleCS{světlo lampy};  \dicPhraseIS{kveikja ljósið} \dicDirectTranslationCS{rozsvítit, zapnout\,/\addthin zapálit světlo};  \dicPhraseIS{slökkva ljósið} \dicDirectTranslationCS{zhasnout, vypnout světlo}  \textbf{2.} \dicSynonym*{þægt barn} \dicDirectTranslationCS{zlatíčko, andílek} \dicIndirectTranslationCS{(o~poslušném dítěti)};  \dicIdiom{ljós}{ \dicPhraseIS{í nýju ljósi}} \dicFlx{adv} \dicLangCat{přen.} \dicDirectTranslationCS{v~novém světle}; { \dicPhraseIS{það kom í ljós að}} \dicDirectTranslationCS{ukázalo se, že; vyšlo najevo, že}
\dicEntry[ljós] \dicTerm{ljós\smash{\textsuperscript{2}}} \dicsymFrequent\  \dicIPA{{l}{j}{ou}{\textlengthmark}{s}} \dicPos{adj}[5]\dicFlx{}[-1] \dicDirectTranslationCS{jasný, světlý} \dicExampleIS{vera ljós á húð} \dicExampleCS{mít světlou pokožku};  \dicPhraseIS{gera sér e‑ð ljóst} \dicDirectTranslationCS{uvědomit si (co)};  \dicPhraseIS{e‑m er ljóst e‑að} \dicFlx{impers} \dicDirectTranslationCS{(co) je (komu) jasné};  \dicPhraseIS{e‑að er deginum ljósara} \dicLangCat{přen.} \dicDirectTranslationCS{(co) je nad slunce jasnější, (co) je jasné jak facka}
\dicEntry[ljósakróna] \dicTerm{ljósa··krón|a} \dicIPA{{l}{j}{ou}{\textlengthmark}{s}{a}{k\smash{\textsuperscript{h}}}{r}{ou}{n}{a}} \dicPos{f}[1] \dicFlx{(‑u, ‑ur)}[7] \dicDirectTranslationCS{(mnohoramenný) lustr}
\dicEntry[ljósapera] \dicTerm{ljósa··per|a} \dicIPA{{l}{j}{ou}{\textlengthmark}{s}{a}{p\smash{\textsuperscript{h}}}{\textepsilon}{r}{a}} \dicPos{f}[1] \dicFlx{(‑u, ‑ur)}[7] \dicFieldCat{elek.} \dicDirectTranslationCS{žárovka}
\dicEntry[ljósaskilti] \dicTerm{ljósa··skilti} \dicIPA{{l}{j}{ou}{\textlengthmark}{s}{a}{s}{\r{\textObardotlessj}}{\textsci}{\textsubring{l}}{\textsubring{d}}{\textsci}} \dicPos{n}[2] \dicFlx{(‑s, ‑)}[14] \dicDirectTranslationCS{neonová reklama}
\dicEntry[ljósaskipti] \dicTerm{ljósa··skipti} \dicIPA{{l}{j}{ou}{\textlengthmark}{s}{a}{s}{\r{\textObardotlessj}}{\textsci}{f}{\textsubring{d}}{\textsci}} \dicPos{n}[2] \dicFlx{pl}[19] \dicSynonym{rökkur} \dicDirectTranslationCS{soumrak}
\dicEntry[ljósastaur] \dicTerm{ljósa··staur} \dicIPA{{l}{j}{ou}{\textlengthmark}{s}{a}{s}{\textsubring{d}}{\oe i}{\textsubring{r}}} \dicPos{m}[4] \dicFlx{(‑s, ‑ar)}[14] \dicDirectTranslationCS{kandelábr, stojan veřejného osvětlení}
\dicEntry[ljósastæði] \dicTerm{ljósa··stæði} \dicIPA{{l}{j}{ou}{\textlengthmark}{s}{a}{s}{\textsubring{d}}{a}{i}{ð}{\textsci}} \dicPos{n}[2] \dicFlx{(‑s, ‑)}[14] \dicDirectTranslationCS{přípojka na světlo}
\dicEntry[Ljósálfheimur] \dicTerm{Ljós··álf·heim|ur}\dicTerm{, Álfheimur} \dicIPA{{l}\-{j}\-{ou}\-{\textlengthmark}\-{s}\-{au}\-{l}\-{f}\-{h}\-{ei}\-{m}\-{\textscy}\-{\textsubring{r}}\-} \dicPos{m}[6] \dicFlx{(‑s)}[3] \dicFieldCat{myt.} \dicDirectTranslationCS{Álfheim} \dicIndirectTranslationCS{(říše světla, světlých elfů)} \dicAntonym{Svartálfheimur}
\dicEntry[ljósár] \dicTerm{ljós··ár} \dicIPA{{l}{j}{ou}{\textlengthmark}{s}{au}{\textsubring{r}}} \dicPos{n}[2] \dicFlx{(‑s, ‑)}[5] \dicFieldCat{astro.} \dicDirectTranslationCS{světelný rok}
\dicEntry[ljósberi] \dicTerm{ljós··ber|i} \dicIPA{{l}{j}{ou}{s}{\textsubring{b}}{\textepsilon}{r}{\textsci}} \dicPos{m}[1] \dicFlx{(‑a, ‑ar)}[1] \dicFieldCat{bot.} \dicDirectTranslationCS{kohoutek alpský} \textit{(l.~{\textLA{Lychnis alpina}})}  \dicsymPhoto\ 
\dicFigure{72890.jpg}{Ljósberi}{Ljósberi - Kesl Michael, Biolib, Copyright/CC-BY-NC}
\dicEntry[ljóseind] \dicTerm{ljós··eind} \dicIPA{{l}{j}{ou}{\textlengthmark}{s}{ei}{n}{\textsubring{d}}} \dicPos{f}[7] \dicFlx{(‑ar, ‑ir)}[1] \dicFieldCat{fyz.} \dicDirectTranslationCS{foton}
\dicEntry[ljósfræði] \dicTerm{ljós··fræð|i} \dicIPA{{l}{j}{ou}{s}{f}{r}{a}{i}{ð}{\textsci}} \dicPos{f}[3] \dicFlx{(‑i)}[3] \dicFieldCat{fyz.} \dicDirectTranslationCS{optika}
\dicEntry[ljósgeisli] \dicTerm{ljós··geisl|i} \dicIPA{{l}{j}{ou}{s}{\r{\textObardotlessj}}{ei}{s}{\textsubring{d}}{l}{\textsci}} \dicPos{m}[1] \dicFlx{(‑a, ‑ar)}[1] \dicDirectTranslationCS{světelný paprsek, paprsek světla}
\dicEntry[ljósgulur] \dicTerm{ljós··gulur} \dicIPA{{l}{j}{ou}{s}{\r{g}}{\textscy}{l}{\textscy}{\textsubring{r}}} \dicPos{adj}[1]\dicFlx{}[-1] \dicDirectTranslationCS{světle žlutý, nažloutlý}
\dicEntry[ljóshraði] \dicTerm{ljós··hrað|i} \dicIPA{{l}{j}{ou}{\textlengthmark}{s}{\textsubring{r}}{a}{ð}{\textsci}} \dicPos{m}[1] \dicFlx{(‑a)}[3] \dicFieldCat{fyz.} \dicDirectTranslationCS{rychlost světla}
\dicEntry[ljóshærður] \dicTerm{ljós··hærður} \dicsymFrequent\  \dicIPA{{l}{j}{ou}{\textlengthmark}{s}{h}{a}{i}{r}{ð}{\textscy}{\textsubring{r}}} \dicPos{adj}[2]\dicFlx{}[-1] \dicDirectTranslationCS{blonďatý, světlovlasý} \dicExampleIS{Hann er ljóshærður.} \dicExampleCS{Má světlé vlasy.}
\dicEntry[ljóshöfðaönd] \dicTerm{ljós·höfða··|önd} \dicIPA{{l}{j}{ou}{\textlengthmark}{s}{h}{\oe}{v}{ð}{a}{\oe}{n}{\textsubring{d}}} \dicPos{f}[8] \dicFlx{(‑andar, ‑endur)}[2] \dicFieldCat{zool.} \dicDirectTranslationCS{hvízdák americký} \textit{(l.~{\textLA{Anas americana}})}  \dicsymPhoto\ 
\dicFigure{ds_image_ljoshofdaond_0_2.jpg}{Ljóshöfðaönd}{Ljóshöfðaönd - Donna Dewhurst, USFWS, PD}
\dicEntry[ljóska] \dicTerm{ljósk|a} \dicIPA{{l}{j}{ou}{s}{\r{g}}{a}} \dicPos{f}[1] \dicFlx{(‑u, ‑ur)}[19] \dicDirectTranslationCS{blondýna, blondýnka}
\dicEntry[ljóskastari] \dicTerm{ljós··kast·ar|i} \dicIPA{{l}{j}{ou}{s}{k\smash{\textsuperscript{h}}}{a}{s}{\textsubring{d}}{a}{r}{\textsci}} \dicPos{m}[1] \dicFlx{(‑a, ‑ar)}[10] \dicDirectTranslationCS{světlomet, reflektor}
\dicEntry[ljósker] \dicTerm{ljós··ker} \dicIPA{{l}{j}{ou}{s}{c\smash{\textsuperscript{h}}}{\textepsilon}{\textsubring{r}}} \dicPos{n}[2] \dicFlx{(‑s, ‑)}[6] \dicDirectTranslationCS{lampa, svítilna, lucerna}
\dicEntry[ljóslifandi] \dicTerm{ljós··lif·andi} \dicIPA{{l}{j}{ou}{s}{l}{\textsci}{v}{a}{n}{\textsubring{d}}{\textsci}} \dicPos{adj}[13] \dicFlx{indecl}[1] \dicDirectTranslationCS{živý, barvitý} \dicExampleIS{e‑að stendur e‑m ljóslifandi fyrir hugskotssjónum} \dicExampleCS{(co) má (kdo) v~živé paměti}
\dicEntry[ljósmerki] \dicTerm{ljós··merki} \dicIPA{{l}{j}{ou}{s}{m}{\textepsilon}{\textsubring{r}}{\r{\textObardotlessj}}{\textsci}} \dicPos{n}[2] \dicFlx{(‑s, ‑)}[16] \dicDirectTranslationCS{světelný signál, světelné znamení}
\dicEntry[ljósmóðir] \dicTerm{ljós··|móðir} \dicIPA{{l}{j}{ou}{s}{m}{ou}{ð}{\textsci}{\textsubring{r}}} \dicPos{f}[11] \dicFlx{(‑móður, ‑mæður)}[1] \dicDirectTranslationCS{porodní asistentka}
\dicEntry[ljósmynd] \dicTerm{ljós··mynd} \dicsymFrequent\  \dicIPA{{l}{j}{ou}{s}{m}{\textsci}{n}{\textsubring{d}}} \dicPos{f}[7] \dicFlx{(‑ar, ‑ir)}[1] \dicDirectTranslationCS{fotografie, fotka} \dicExampleIS{svarthvít ljósmynd af ungri konu} \dicExampleCS{černobílá fotografie mladé ženy}
\dicEntry[ljósmynda] \dicTerm{ljós··mynd|a} \dicIPA{{l}{j}{ou}{s}{m}{\textsci}{n}{\textsubring{d}}{a}} \dicPos{v}[1] \dicFlx{(‑aði)}[1] \dicDirectTranslationCS{(vy)fotit, (vy)fotografovat, (u)dělat fotografii}
\dicEntry[ljósmyndari] \dicTerm{ljós··mynd·ar|i} \dicIPA{{l}{j}{ou}{s}{m}{\textsci}{n}{\textsubring{d}}{a}{r}{\textsci}} \dicPos{m}[1] \dicFlx{(‑a, ‑ar)}[13] \dicDirectTranslationCS{fotograf(ka)}
\dicEntry[ljósmyndavél] \dicTerm{ljós··mynda·vél} \dicIPA{{l}{j}{ou}{s}{m}{\textsci}{n}{\textsubring{d}}{a}{v}{j}{\textepsilon}{\textsubring{l}}} \dicPos{f}[4] \dicFlx{(‑ar, ‑ar)}[1] \dicDirectTranslationCS{fotoaparát}
\dicEntry[ljósmyndun] \dicTerm{ljós··mynd|un} \dicIPA{{l}{j}{ou}{s}{m}{\textsci}{n}{\textsubring{d}}{\textscy}{\textsubring{n}}} \dicPos{f}[7] \dicFlx{(‑unar, ‑anir)}[8] \textbf{1.} \dicDirectTranslationCS{fotografování}  \textbf{2.} \dicDirectTranslationCS{fotografie} \dicIndirectTranslationCS{(umělecký obor)}
\dicEntry[ljósmælir] \dicTerm{ljós··mæl|ir} \dicIPA{{l}{j}{ou}{s}{m}{a}{i}{l}{\textsci}{\textsubring{r}}} \dicPos{m}[7] \dicFlx{(‑is, ‑ar)}[1] \dicFieldCat{fyz.} \dicDirectTranslationCS{expozimetr}
\dicEntry[ljósnemi] \dicTerm{ljós··nem|i} \dicIPA{{l}{j}{ou}{s}{n}{\textepsilon}{m}{\textsci}} \dicPos{m}[1] \dicFlx{(‑a, ‑ar)}[1] \dicFieldCat{fyz.} \dicDirectTranslationCS{fotočlánek}
\dicEntry[ljósnæmur] \dicTerm{ljós··næmur} \dicIPA{{l}{j}{ou}{s}{n}{a}{i}{m}{\textscy}{\textsubring{r}}} \dicPos{adj}[1]\dicFlx{}[-1] \dicDirectTranslationCS{citlivý na světlo, fotosenzitivní}
\dicEntry[ljósrit] \dicTerm{ljós··rit} \dicIPA{{l}{j}{ou}{s}{r}{\textsci}{\textsubring{d}}} \dicPos{n}[2] \dicFlx{(‑s, ‑)}[5] \dicDirectTranslationCS{fotokopie}
\dicEntry[ljósrita] \dicTerm{ljós··rit|a} \dicIPA{{l}{j}{ou}{s}{r}{\textsci}{\textsubring{d}}{a}} \dicPos{v}[1] \dicFlx{(‑aði)}[1] \dicFlx{acc} \dicDirectTranslationCS{(o)kopírovat (pomocí fotokopírky)}
\dicEntry[ljósritunarvél] \dicTerm{ljós·ritunar··vél} \dicIPA{{l}{j}{ou}{s}{r}{\textsci}{\textsubring{d}}{\textscy}{n}{a}{r}{v}{j}{\textepsilon}{\textsubring{l}}} \dicPos{f}[4] \dicFlx{(‑ar, ‑ar)}[1] \dicDirectTranslationCS{(foto)kopírka}
\dicEntry[ljósta] \dicTerm{ljósta} \dicsymFrequent\  \dicIPA{{l}{j}{ou}{s}{\textsubring{d}}{a}} \dicPos{v}[6] \dicFlx{(lýst, laust, lustum, lysti, lostið)}[97] \dicFlx{acc} \dicLangCat{zast.} \dicDirectTranslationCS{udeřit, uhodit} \dicExampleIS{ljósta e‑n í andlitið með e‑u} \dicExampleCS{uhodit (koho) do tváře (čím)};  \dicIdiom{ljósta}[niður]{ \dicPhraseIS{e‑u lýstur niður}} \dicFlx{impers} \dicDirectTranslationCS{(co) udeří, (co) zasáhne (blesk ap.)} \dicExampleIS{Eldingu laust niður í húsið.} \dicExampleCS{Blesk udeřil do domu.};  \dicIdiom{ljósta}[saman]{ \dicPhraseIS{e‑m lýstur saman}} \dicFlx{impers} \dicDirectTranslationCS{(kdo) se střetne} \dicExampleIS{Herjunum laust saman.} \dicExampleCS{Armády se střetly.};  \dicIdiom{ljósta}[upp]{ \dicPhraseIS{ljósta upp ópi}} \dicDirectTranslationCS{vykřiknout}
\dicEntry[ljóstillífun] \dicTerm{ljós··til·líf|un} \dicIPA{{l}{j}{ou}{s}{t\smash{\textsuperscript{h}}}{\textsci}{l}{i}{v}{\textscy}{\textsubring{n}}} \dicPos{f}[7] \dicFlx{(‑unar)}[9] \dicFieldCat{biol.} \dicDirectTranslationCS{fotosyntéza}
\dicEntry[ljóstra] \dicTerm{ljóstr|a} \dicIPA{{l}{j}{ou}{s}{\textsubring{d}}{r}{a}} \dicPos{v}[1] \dicFlx{(‑aði)}[1] \dicFlx{dat} \dicPhraseIS{ljóstra e‑u upp} \dicDirectTranslationCS{odhalit\,/\addthin odhalovat (co), odkrýt\,/\addthin odkrývat (co)} \dicExampleIS{ljóstra upp leyndarmáli} \dicExampleCS{odkrýt tajemství}
\dicEntry[ljósvaki] \dicTerm{ljós··vak|i} \dicIPA{{l}{j}{ou}{\textlengthmark}{s}{v}{a}{\r{\textObardotlessj}}{\textsci}} \dicPos{m}[1] \dicFlx{(‑a, ‑ar)}[8] \dicDirectTranslationCS{éter} \dicIndirectTranslationCS{(podle starých představ látka zaplňující vesmírný prostor)};  \dicPhraseIS{á öldum ljósvakans} \dicFlx{adv} \dicDirectTranslationCS{na vlnách éteru}
\dicEntry[ljótleiki] \dicTerm{ljót··leik|i} \dicIPA{{l}{j}{ou}{\textlengthmark}{\textsubring{d}}{l}{ei}{\r{\textObardotlessj}}{\textsci}} \dicPos{m}[1] \dicFlx{(‑a)}[3] \dicSynonym{lýti} \dicDirectTranslationCS{šerednost, ohyzdnost (vzhled ap.)}
\dicEntry[ljótur] \dicTerm{ljótur} \dicsymFrequent\  \dicIPA{{l}{j}{ou}{\textlengthmark}{\textsubring{d}}{\textscy}{\textsubring{r}}} \dicPos{adj}[1]\dicFlx{}[-1] \dicDirectTranslationCS{šeredný, ošklivý, škaredý} \dicExampleIS{ljótar athugasemdir} \dicExampleCS{ošklivé poznámky}
\dicEntry[Ljubljana] \dicTerm{Ljubljana} \dicIPA{{l}{j}{\textscy}{\textsubring{b}}{l}{j}{a}{n}{a}} \dicPos{subs} \dicFlx{indecl} \dicFieldCat{geog.} \dicDirectTranslationCS{Lublaň} \dicIndirectTranslationCS{(hlavní město Slovinska)}
\dicEntry[ljúffengur] \dicTerm{ljúf··fengur} \dicIPA{{l}{j}{u}{f}{\textlengthmark}{ei}{\ng}{\r{g}}{\textscy}{\textsubring{r}}} \dicPos{adj}[1]\dicFlx{}[-1] \dicDirectTranslationCS{chutný, lahodný, delikátní} \dicIndirectTranslationCS{(o~jídle)}
\dicEntry[ljúflingur] \dicTerm{ljúf··ling|ur} \dicIPA{{l}{j}{u}{v}{l}{i}{\ng}{\r{g}}{\textscy}{\textsubring{r}}} \dicPos{m}[6] \dicFlx{(‑s, ‑ar)}[8] \textbf{1.} \dicSynonym{huldumaður} \dicDirectTranslationCS{elf}  \textbf{2.} \dicSynonym{eftirlæti} \dicDirectTranslationCS{miláček, zlato}
\dicEntry[ljúflyndi] \dicTerm{ljúf··lyndi} \dicIPA{{l}{j}{u}{v}{l}{\textsci}{n}{\textsubring{d}}{\textsci}} \dicPos{n}[2] \dicFlx{(‑s)}[20] \dicDirectTranslationCS{přívětivost, laskavost}
\dicEntry[ljúfmenni] \dicTerm{ljúf··menni} \dicIPA{{l}{j}{u}{v}{m}{\textepsilon}{n}{\textsci}} \dicPos{n}[2] \dicFlx{(‑s, ‑)}[14] \dicSynonym{góðmenni} \dicDirectTranslationCS{milý\,/\addthin přívětivý člověk}
\dicEntry[ljúfmennska] \dicTerm{ljúf··mennsk|a} \dicIPA{{l}{j}{u}{v}{m}{\textepsilon}{n}{s}{\r{g}}{a}} \dicPos{f}[1] \dicFlx{(‑u)}[5] \dicDirectTranslationCS{dobrota, přívětivost, vlídnost}
\dicEntry[ljúfur] \dicTerm{ljúf|ur\smash{\textsuperscript{1}}} \dicIPA{{l}{j}{u}{\textlengthmark}{v}{\textscy}{\textsubring{r}}} \dicPos{m}[6] \dicFlx{(‑s, ‑ar)}[24] \dicDirectTranslationCS{poklad, drahoušek, miláček} \dicExampleIS{ljúfur(inn) (minn)} \dicExampleCS{ty můj poklade}
\dicEntry[ljúfur] \dicTerm{ljúfur\smash{\textsuperscript{2}}} \dicsymFrequent\  \dicIPA{{l}{j}{u}{\textlengthmark}{v}{\textscy}{\textsubring{r}}} \dicPos{adj}[1]\dicFlx{}[-1] \textbf{1.} \dicSynonym{blíður} \dicDirectTranslationCS{laskavý, přívětivý, vlídný}  \textbf{2.} \dicSynonym{þægilegur} \dicDirectTranslationCS{příjemný, milý} \dicExampleIS{ljúfar minningar} \dicExampleCS{příjemné vzpomínky};  \dicPhraseIS{e‑m er ljúft að (gera e‑ð)} \dicFlx{impers} \dicDirectTranslationCS{(komu) je příjemné (udělat (co))}
\dicEntry[ljúga] \dicTerm{ljúga} \dicsymFrequent\  \dicIPA{{l}{j}{u}{\textlengthmark}{a}} \dicPos{v}[6] \dicFlx{(lýg, laug, lugum, lygi, logið)}[52] \dicFlx{dat} \dicDirectTranslationCS{lhát, zalhat} \dicExampleIS{ljúga e‑u} \dicExampleCS{zalhat o~(čem)};  \dicPhraseIS{ljúga að e‑m} \dicDirectTranslationCS{lhát (komu)};  \dicPhraseIS{ljúga til um e‑ð} \dicDirectTranslationCS{lhát o~(čem)} \dicExampleIS{ljúga til um aldur} \dicExampleCS{lhát o~(svém) věku}
\dicEntry[ljúgvitni] \dicTerm{ljúg··vitni} \dicIPA{{l}{j}{u}{\textlengthmark}{v}{\textsci}{h}{\textsubring{d}}{n}{\textsci}} \dicPos{n}[2] \dicFlx{(‑s, ‑)}[14] \dicDirectTranslationCS{křivá přísaha, křivé doznání};  \dicPhraseIS{bera ljúgvitni gegn e‑m} \dicDirectTranslationCS{křivě svědčit proti (komu)}
\dicEntry[ljúka] \dicTerm{ljúka} \dicsymFrequent\  \dicIPA{{l}{j}{u}{\textlengthmark}{\r{g}}{a}} \dicPos{v}[6] \dicFlx{(lýk, lauk, lukum, lyki, lokið)}[47] \dicFlx{dat} \textbf{1.} \dicSynonym{klára} \dicDirectTranslationCS{(u)končit, ukončovat, zakončit, zakončovat, dokončit, dokončovat};  \dicPhraseIS{ljúka e‑u} \dicDirectTranslationCS{ukončit (co)};  \dicPhraseIS{ljúka við e‑ð} \dicDirectTranslationCS{ukončit (co), (s)končit s~(čím)}  \textbf{2.} \dicPhraseIS{e‑u lýkur} \dicFlx{impers} \dicDirectTranslationCS{(co) končí} \dicExampleIS{Skólanum lýkur í júní.} \dicExampleCS{Škola končí v~červnu.};  \dicIdiom{ljúka}[af]{ \dicPhraseIS{ljúka e‑u af}} \dicDirectTranslationCS{ukončit (co), skončit (co)};  \dicIdiom{ljúka}[upp]{ \dicPhraseIS{ljúka upp e‑u}} \dicLangCat{zast.} \dicSynonym{opna\smash{\textsuperscript{2}}} \dicDirectTranslationCS{otevřít (co)} \dicExampleIS{ljúka dyrunum upp} \dicExampleCS{otevřít dveře};  \dicIdiom{ljúka}[yfir]{ \dicPhraseIS{þar til yfir lýkur}} \dicFlx{adv} \dicDirectTranslationCS{až do samého konce}; { \dicPhraseIS{áður en yfir lýkur}} \dicFlx{adv} \dicDirectTranslationCS{dříve než bude po všem};  \dicIdiom{ljúkast}[saman]{ \dicPhraseIS{ljúkast saman}} \dicFlx{refl} \dicLangCat{zast.} \dicDirectTranslationCS{zavřít se};  \dicIdiom{ljúkast}[upp]{ \dicPhraseIS{ljúkast upp}} \dicFlx{refl} \dicLangCat{zast.} \dicDirectTranslationCS{otevřít se, rozevřít se}; { \dicPhraseIS{e‑að lýkst upp fyrir e‑m}} \dicFlx{refl} \dicLangCat{přen.} \dicDirectTranslationCS{(co) se otevírá před (kým) (nová myšlenka ap.)}
\dicEntry[ljæ] \dicTerm{ljæ} \dicIPA{{l}{j}{a}{i}{\textlengthmark}} \dicPos{v} \dicFlx{ind praes sg 1 pers} \dicLink{ljá}
\dicEntry[lo.] \dicTerm{lo.} \dicPos{zkr} \dicPhraseIS{lýsingarorð} \dicFieldCat{jaz.} \dicDirectTranslationCS{přídavné jméno}
\dicEntry[lobbíisti] \dicTerm{lobbí··ist|i} \dicIPA{{l}{\textopeno}{\textsubring{b}}{\textlengthmark}{i}{j}{\textsci}{s}{\textsubring{d}}{\textsci}} \dicPos{m}[1] \dicFlx{(‑a, ‑ar)}[1] \dicDirectTranslationCS{lobbista, lobbistka}
\dicEntry[loddari] \dicTerm{lodd··ar|i} \dicIPA{{l}{\textopeno}{\textsubring{d}}{\textlengthmark}{a}{r}{\textsci}} \dicPos{m}[1] \dicFlx{(‑a, ‑ar)}[13] \textbf{1.} \dicSynonym{svikari} \dicDirectTranslationCS{šarlatán(ka)}  \textbf{2.} \dicSynonym{trúður} \dicDirectTranslationCS{šašek, klaun(ka)}
\dicEntry[loða] \dicTerm{lo|ða} \dicIPA{{l}{\textopeno}{\textlengthmark}{ð}{a}} \dicPos{v}[2] \dicFlx{(‑ddi, ‑ðað)}[182] \dicSynonym{tolla\smash{\textsuperscript{2}}} \dicDirectTranslationCS{přiléhat, lnout, přimknout se, přimykat};  \dicIdiom{loða}[saman]{ \dicPhraseIS{loða saman}} \dicDirectTranslationCS{přilnout k~sobě};  \dicIdiom{loða}[við]{ \dicPhraseIS{loða við e‑ð}} \dicDirectTranslationCS{přilnout k~(čemu)} \dicExampleIS{Málningin loðir við vegginn.} \dicExampleCS{Nátěr přilne ke zdi.}
\dicEntry[loðdýr] \dicTerm{loð··dýr} \dicIPA{{l}{\textopeno}{ð}{\textsubring{d}}{i}{\textsubring{r}}} \dicPos{n}[2] \dicFlx{(‑s, ‑)}[5] \dicDirectTranslationCS{kožešinové zvíře}
\dicEntry[loðfeldur] \dicTerm{loð··feld|ur} \dicIPA{{l}{\textopeno}{ð}{f}{\textepsilon}{l}{\textsubring{d}}{\textscy}{\textsubring{r}}} \dicPos{m}[10] \dicFlx{(‑ar, ‑ir)}[4] \dicSynonym{loðskinn} \dicDirectTranslationCS{kožešina}
\dicEntry[loðfíll] \dicTerm{loð··fíl|l} \dicIPA{{l}{\textopeno}{ð}{f}{i}{\textsubring{d}}{\textsubring{l}}} \dicPos{m}[6] \dicFlx{(‑s, ‑ar)}[48] \dicFieldCat{zool.} \dicDirectTranslationCS{mamut} \textit{(l.~{\textLA{Mammuthus primigenius}})}
\dicEntry[loðhúfa] \dicTerm{loð··húf|a} \dicIPA{{l}{\textopeno}{\textlengthmark}{\texttheta}{h}{u}{v}{a}} \dicPos{f}[1] \dicFlx{(‑u, ‑ur)}[19] \dicDirectTranslationCS{kožešinová čepice, beranice}
\dicEntry[loðinn] \dicTerm{loðinn} \dicIPA{{l}{\textopeno}{\textlengthmark}{ð}{\textsci}{\textsubring{n}}} \dicPos{adj}[6]\dicFlx{}[-2] \textbf{1.} \dicDirectTranslationCS{chlupatý, srstnatý, huňatý}  \textbf{2.} \dicSynonym{óhreinn} \dicDirectTranslationCS{vyhýbavý, mlhavý (odpověď ap.)}
\dicEntry[loðkápa] \dicTerm{loð··káp|a} \dicIPA{{l}{\textopeno}{ð}{k\smash{\textsuperscript{h}}}{au}{\textsubring{b}}{a}} \dicPos{f}[1] \dicFlx{(‑u, ‑ur)}[19] \dicDirectTranslationCS{kožich}
\dicEntry[loðna] \dicTerm{loðn|a} \dicIPA{{l}{\textopeno}{ð}{n}{a}} \dicPos{f}[1] \dicFlx{(‑u, ‑ur)}[7] \textbf{1.} \dicDirectTranslationCS{ochlupení}  \textbf{2.} \dicFieldCat{zool.} \dicDirectTranslationCS{huňáček severní} \textit{(l.~{\textLA{Mallotus villosus}})}
\dicEntry[loðskinn] \dicTerm{loð··skinn} \dicIPA{{l}{\textopeno}{ð}{s}{\r{\textObardotlessj}}{\textsci}{\textsubring{n}}} \dicPos{n}[2] \dicFlx{(‑s, ‑)}[5] \dicDirectTranslationCS{kožešina, kožka}
\dicEntry[loðvíðir] \dicTerm{loð··víð|ir} \dicIPA{{l}{\textopeno}{ð}{v}{i}{ð}{\textsci}{\textsubring{r}}} \dicPos{m}[7] \dicFlx{(‑is)}[2] \dicFieldCat{bot.} \dicDirectTranslationCS{vrba vlnatá} \textit{(l.~{\textLA{Salix lanata}})}  \dicsymPhoto\ 
\dicFigure{66869.jpg}{Loðvíðir}{Loðvíðir - Hroneš Michal, Biolib, Copyright/CC-BY-NC}
\dicEntry[lof] \dicTerm{lof} \dicIPA{{l}{\textopeno}{\textlengthmark}{f}} \dicPos{n}[2] \dicFlx{(‑s)}[2] \dicSynonym{hrós} \dicDirectTranslationCS{chvála, pochvala};  \dicPhraseIS{bera lof á e‑n} \dicDirectTranslationCS{chválit (koho)};  \dicPhraseIS{guði sé lof} \dicDirectTranslationCS{díky bohu};  \dicPhraseIS{klappa lof í lófa} \dicDirectTranslationCS{tleskat, aplaudovat}
\dicEntry[lofa] \dicTerm{lof|a} \dicsymFrequent\  \dicIPA{{l}{\textopeno}{\textlengthmark}{v}{a}} \dicPos{v}[1] \dicFlx{(‑aði)}[1] \dicFlx{acc\,/\addthin dat} \textbf{1.} \dicFlx{acc} \dicSynonym{hrósa} \dicDirectTranslationCS{(po)chválit, vychvalovat};  \dicPhraseIS{lofa e‑n\,/\addthin e‑ð} \dicDirectTranslationCS{chválit (koho\,/\addthin co)} \dicExampleIS{lofa listaverkið} \dicExampleCS{pochválit umělecké dílo}  \textbf{2.} \dicFlx{dat + dat} \dicSynonym{heita} \dicDirectTranslationCS{(při)slíbit, slibovat};  \dicPhraseIS{lofa e‑m e‑u} \dicDirectTranslationCS{slíbit (komu co)}  \textbf{3.} \mbox{\dicFlx{dat + acc}} \dicSynonym{leyfa} \dicDirectTranslationCS{dovolit, dovolovat, povolit, povolovat};  \dicPhraseIS{lofa e‑m e‑ð} \dicDirectTranslationCS{dovolit (komu co)} \dicExampleIS{Ég lofaði henni að dansa.} \dicExampleCS{Dovolil jsem jí tančit.};  \dicIdiom{lofast}{ \dicPhraseIS{lofast e‑m}} \dicFlx{refl} \dicLangCat{zast.} \dicDirectTranslationCS{zaslíbit se (komu)}
\dicEntry[lofgerð] \dicTerm{lof··gerð} \dicIPA{{l}{\textopeno}{v}{\r{\textObardotlessj}}{\textepsilon}{r}{\texttheta}} \dicPos{f}[7] \dicFlx{(‑ar, ‑ir)}[1] \dicSynonym{lof} \dicDirectTranslationCS{chvála, pochvala}
\dicEntry[loforð] \dicTerm{lof··orð} \dicIPA{{l}{\textopeno}{\textlengthmark}{v}{\textopeno}{r}{\texttheta}} \dicPos{n}[2] \dicFlx{(‑s, ‑)}[5] \textbf{1.} \dicDirectTranslationCS{slib, slovo};  \dicPhraseIS{standa við loforð} \dicDirectTranslationCS{držet slovo}  \textbf{2.} \dicDirectTranslationCS{přísaha}
\dicEntry[lofsamlegur] \dicTerm{lof·sam··legur} \dicIPA{{l}{\textopeno}{f}{s}{a}{m}{l}{\textepsilon}{\textbabygamma}{\textscy}{\textsubring{r}}} \dicPos{adj}[1]\dicFlx{}[-8] \dicDirectTranslationCS{pochvalný}
\dicEntry[lofsorð] \dicTerm{lofs··orð} \dicIPA{{l}{\textopeno}{f}{s}{\textopeno}{r}{\texttheta}} \dicPos{n}[2] \dicFlx{(‑s, ‑)}[5] \dicPhraseIS{ljúka lofsorði á e‑ð} \dicDirectTranslationCS{chválit (co), vyjádřit se o~(čem) pochvalně}
\dicEntry[lofsverður] \dicTerm{lofs··verður} \dicIPA{{l}{\textopeno}{f}{s}{v}{\textepsilon}{r}{ð}{\textscy}{\textsubring{r}}} \dicPos{adj}[2]\dicFlx{}[-1] \dicDirectTranslationCS{chvályhodný, záslužný} \dicExampleIS{lofsverð viðleitni} \dicExampleCS{chvályhodné úsilí}
\dicEntry[lofsöngur] \dicTerm{lof··söng|ur} \dicIPA{{l}{\textopeno}{f}{s}{\oe i}{\ng}{\r{g}}{\textscy}{\textsubring{r}}} \dicPos{m}[8] \dicFlx{(‑s, ‑var)}[2] \dicDirectTranslationCS{chvalozpěv}
\dicEntry[loft] \dicTerm{loft} \dicsymFrequent\  \dicIPA{{l}{\textopeno}{f}{\textsubring{d}}} \dicPos{n}[2] \dicFlx{(‑s, ‑)}[5] \textbf{1.} \dicSynonym{andrúmsloft} \dicDirectTranslationCS{vzduch, povětří, ovzduší} \dicExampleIS{ferskt loft} \dicExampleCS{čerstvý vzduch}  \textbf{2.} \dicSynonym{himinn} \dicDirectTranslationCS{vzduch, obloha} \dicExampleIS{kasta e‑u upp í loftið} \dicExampleCS{vyhodit (co) do vzduchu}  \textbf{3.} \dicDirectTranslationCS{poschodí, podlaží};  \dicPhraseIS{uppi á lofti} \dicFlx{adv} \dicDirectTranslationCS{v~horním patře}  \textbf{4.} \dicSynonym{háaloft} \dicDirectTranslationCS{podkroví};  \dicIdiom{loft}{ \dicPhraseIS{það er allt upp í loft}} \dicLangCat{přen.} \dicDirectTranslationCS{všechno je vzhůru nohama};  \dicPhraseIS{e‑að er úr lausu lofti gripið} \dicLangCat{přen.} \dicDirectTranslationCS{(co) je neopodstatněné}
\dicEntry[loftárás] \dicTerm{loft··á·rás} \dicIPA{{l}{\textopeno}{f}{\textsubring{d}}{au}{r}{au}{s}} \dicPos{f}[7] \dicFlx{(‑ar, ‑ir)}[1] \dicDirectTranslationCS{nálet, letecký útok}
\dicEntry[loftbelgur] \dicTerm{loft··belg|ur} \dicIPA{{l}{\textopeno}{f}{\textsubring{d}}{\textsubring{b}}{\textepsilon}{l}{\r{g}}{\textscy}{\textsubring{r}}} \dicPos{m}[9] \dicFlx{(‑s\,/\addthin ‑jar, ‑ir)}[26] \dicFieldCat{let.} \dicDirectTranslationCS{(horkovzdušný) balón}
\dicEntry[loftbor] \dicTerm{loft··bor} \dicIPA{{l}{\textopeno}{f}{\textsubring{d}}{\textsubring{b}}{\textopeno}{\textsubring{r}}} \dicPos{m}[4] \dicFlx{(‑s, ‑ar)}[14] \dicDirectTranslationCS{sbíječka, sbíjecí kladivo}
\dicEntry[loftbréf] \dicTerm{loft··bréf} \dicIPA{{l}{\textopeno}{f}{\textsubring{d}}{\textsubring{b}}{r}{j}{\textepsilon}{f}} \dicPos{n}[2] \dicFlx{(‑s, ‑)}[5] \dicDirectTranslationCS{letecké psaní}
\dicEntry[loftbrú] \dicTerm{loft··|brú} \dicIPA{{l}{\textopeno}{f}{\textsubring{d}}{\textsubring{b}}{r}{u}} \dicPos{f}[9] \dicFlx{(‑brúar, ‑brýr)}[4] \dicDirectTranslationCS{letecký most}
\dicEntry[loftfar] \dicTerm{loft··|far} \dicIPA{{l}{\textopeno}{f}{\textsubring{d}}{f}{a}{\textsubring{r}}} \dicPos{n}[2] \dicFlx{(‑fars, ‑för)}[8] \dicDirectTranslationCS{létající prostředek (letadlo, balón, vzducholoď ap.)}
\dicEntry[loftfimleikar] \dicTerm{loft··fim·leikar} \dicIPA{{l}{\textopeno}{f}{\textsubring{d}}{f}{\textsci}{m}{l}{ei}{\r{g}}{a}{\textsubring{r}}} \dicPos{m}[6] \dicFlx{pl}[2] \dicDirectTranslationCS{letecká akrobacie}
\dicEntry[lofthræddur] \dicTerm{loft··hræddur} \dicIPA{{l}{\textopeno}{f}{\textsubring{d}}{\textsubring{r}}{a}{i}{\textsubring{d}}{\textscy}{\textsubring{r}}} \dicPos{adj}[2]\dicFlx{}[-18] \dicDirectTranslationCS{mající strach z~výšek}
\dicEntry[lofthræðsla] \dicTerm{loft··hræðsl|a} \dicIPA{{l}{\textopeno}{f}{\textsubring{d}}{\textsubring{r}}{a}{i}{ð}{s}{\textsubring{d}}{l}{a}} \dicPos{f}[1] \dicFlx{(‑u)}[5] \dicDirectTranslationCS{strach z~výšek}
\dicEntry[loftkastali] \dicTerm{loft··kastal|i} \dicIPA{{l}{\textopeno}{f}{\textsubring{d}}{k\smash{\textsuperscript{h}}}{a}{s}{\textsubring{d}}{a}{l}{\textsci}} \dicPos{m}[1] \dicFlx{(‑a, ‑ar)}[12] \dicDirectTranslationCS{vzdušný zámek};  \dicPhraseIS{byggja loftkastala} \dicLangCat{přen.} \dicDirectTranslationCS{stavět vzdušné zámky}
\dicEntry[loftkenndur] \dicTerm{loft··kenndur} \dicIPA{{l}{\textopeno}{f}{\textsubring{d}}{c\smash{\textsuperscript{h}}}{\textepsilon}{n}{\textsubring{d}}{\textscy}{\textsubring{r}}} \dicPos{adj}[2]\dicFlx{}[-14] \textbf{1.} \dicDirectTranslationCS{vzdušný}  \textbf{2.} \dicDirectTranslationCS{plynný} \dicExampleIS{loftkennd gosefni} \dicExampleCS{plynné magma}
\dicEntry[loftkældur] \dicTerm{loft··kældur} \dicIPA{{l}{\textopeno}{f}{\textsubring{d}}{c\smash{\textsuperscript{h}}}{a}{i}{l}{\textsubring{d}}{\textscy}{\textsubring{r}}} \dicPos{adj}[2]\dicFlx{}[-17] \dicDirectTranslationCS{klimatizovaný} \dicExampleIS{loftkæld skrifstofa} \dicExampleCS{klimatizovaná kancelář}
\dicEntry[loftkæling] \dicTerm{loft··kæl·ing} \dicIPA{{l}{\textopeno}{f}{\textsubring{d}}{c\smash{\textsuperscript{h}}}{a}{i}{l}{i}{\ng}{\r{g}}} \dicPos{f}[4] \dicFlx{(‑ar, ‑ar)}[5] \dicDirectTranslationCS{klimatizace}
\begin{xtolerant}{}{1pt}
\dicEntry[loftmengun] \dicTerm{loft··meng|un}\addthinS\dicIPA{{l}{\textopeno}{f}{\textsubring{d}}{m}{ei}{\ng}{\r{g}}{\textscy}{\textsubring{n}}}\addthinS\dicPos{f}[7]\addthinS\dicFlx{(‑unar)}[9] \dicDirectTranslationCS{znečištění ovzduší}
\end{xtolerant}
\dicEntry[loftnet] \dicTerm{loft··net} \dicIPA{{l}{\textopeno}{f}{\textsubring{d}}{n}{\textepsilon}{\textsubring{d}}} \dicPos{n}[2] \dicFlx{(‑s, ‑)}[5] \dicDirectTranslationCS{anténa}
\dicEntry[loftræsta] \dicTerm{loft··ræst|a} \dicIPA{{l}{\textopeno}{f}{\textsubring{d}}{r}{a}{i}{s}{\textsubring{d}}{a}} \dicPos{v}[2] \dicFlx{(‑i, ‑)}[16] \dicFlx{acc} \dicDirectTranslationCS{(vy)větrat, (vy)ventilovat} \dicExampleIS{loftræsta sal} \dicExampleCS{vyvětrat sál}
\dicEntry[loftræsting] \dicTerm{loft··ræst·ing} \dicIPA{{l}{\textopeno}{f}{\textsubring{d}}{r}{a}{i}{s}{\textsubring{d}}{i}{\ng}{\r{g}}} \dicPos{f}[4] \dicFlx{(‑ar, ‑ar)}[5] \dicDirectTranslationCS{větrání, ventilace}
\begin{xtolerant}{}{1pt}
\dicEntry[loftskeytamaður] \dicTerm{loft·skeyta··|maður}\addthinSS\dicIPA{{l}{\textopeno}{f}{\textsubring{d}}{s}{c\smash{\textsuperscript{h}}}{ei}{\textsubring{d}}{a}{m}{a}{ð}{\textscy}{\textsubring{r}}}\addthinSS\dicPos{m}[1\hspace{-0.10667em}3]\addthinS\dicFlx{(‑manns,\addthinS‑menn)}[2] \dicDirectTranslationCS{radista, radistka, rádiový operátor, rádiová operátorka}
\end{xtolerant}
\dicEntry[loftskip] \dicTerm{loft··skip} \dicIPA{{l}{\textopeno}{f}{\textsubring{d}}{s}{\r{\textObardotlessj}}{\textsci}{\textsubring{b}}} \dicPos{n}[2] \dicFlx{(‑s, ‑)}[5] \dicFieldCat{let.} \dicDirectTranslationCS{vzducholoď}
\dicEntry[loftslag] \dicTerm{lofts··lag} \dicIPA{{l}{\textopeno}{f}{s}{l}{a}{x}} \dicPos{n}[2] \dicFlx{(‑s)}[2] \dicSynonym{veðurfar} \dicDirectTranslationCS{podnebí, klima}
\dicEntry[loftslagsbelti] \dicTerm{lofts·lags··belti} \dicIPA{{l}{\textopeno}{f}{s}{l}{a}{x}{s}{\textsubring{b}}{\textepsilon}{\textsubring{l}}{\textsubring{d}}{\textsci}} \dicPos{n}[2] \dicFlx{(‑s, ‑)}[14] \dicFieldCat{meteo.} \dicDirectTranslationCS{podnebný pás}
\dicEntry[loftslagsfræði] \dicTerm{lofts·lags··fræð|i} \dicIPA{{l}{\textopeno}{f}{s}{l}{a}{x}{s}{f}{r}{a}{i}{ð}{\textsci}} \dicPos{f}[3] \dicFlx{(‑i)}[3] \dicSynonym{veðurfarsfræði} \dicDirectTranslationCS{klimatologie}
\dicEntry[loftsteinn] \dicTerm{loft··stein|n} \dicIPA{{l}{\textopeno}{f}{\textsubring{d}}{s}{\textsubring{d}}{ei}{\textsubring{d}}{\textsubring{n}}} \dicPos{m}[6] \dicFlx{(‑s, ‑ar)}[42] \dicFieldCat{astro.} \dicDirectTranslationCS{meteorit}
\dicEntry[lofttegund] \dicTerm{loft··tegund} \dicIPA{{l}{\textopeno}{f}{t\smash{\textsuperscript{h}}}{\textepsilon}{\textbabygamma}{\textscy}{n}{\textsubring{d}}} \dicPos{f}[7] \dicFlx{(‑ar, ‑ir)}[1] \dicFieldCat{chem.} \dicDirectTranslationCS{plyn}
\dicEntry[lofttóm] \textls[15]{\dicTerm{loft··tóm} \dicIPA{{l}{\textopeno}{f}{t\smash{\textsuperscript{h}}}{ou}{\textsubring{m}}} \dicPos{n}[2] \dicFlx{(‑s)}[2] \dicFieldCat{fyz.} \dicDirectTranslationCS{vakuum, vzduchoprázdno}}
\dicEntry[loftvarnir] \dicTerm{loft··varnir} \dicIPA{{l}{\textopeno}{f}{\textsubring{d}}{v}{a}{r}{\textsubring{d}}{n}{\textsci}{\textsubring{r}}} \dicPos{f}[7] \dicFlx{pl}[18] \dicFieldCat{voj.} \dicDirectTranslationCS{protiletecká\,/\addthin protiletadlová obrana}
\dicEntry[loftveiki] \dicTerm{loft··veik|i} \dicIPA{{l}{\textopeno}{f}{\textsubring{d}}{v}{ei}{\r{\textObardotlessj}}{\textsci}} \dicPos{f}[3] \dicFlx{(‑i)}[3] \dicFieldCat{med.} \dicDirectTranslationCS{výšková\,/\addthin horská nemoc}
\dicEntry[loftvog] \dicTerm{loft··vog} \dicIPA{{l}{\textopeno}{f}{\textsubring{d}}{v}{\textopeno}{x}} \dicPos{f}[7] \dicFlx{(‑ar, ‑ir)}[1] \dicDirectTranslationCS{barometr}
\dicEntry[loftþéttur] \dicTerm{loft··þéttur} \dicIPA{{l}{\textopeno}{f}{\textsubring{d}}{\texttheta}{j}{\textepsilon}{h}{\textsubring{d}}{\textscy}{\textsubring{r}}} \dicPos{adj}[1]\dicFlx{}[-10] \dicDirectTranslationCS{vzduchotěsný, hermetický}
\dicEntry[loftþrýstingur] \dicTerm{loft··þrýst·ing|ur} \dicIPA{{l}{\textopeno}{f}{\textsubring{d}}{\texttheta}{r}{i}{s}{\textsubring{d}}{i}{\ng}{\r{g}}{\textscy}{\textsubring{r}}} \dicPos{m}[6] \dicFlx{(‑s)}[9] \dicDirectTranslationCS{tlak vzduchu, atmosférický tlak}
\dicEntry[loga] \dicTerm{log|a} \dicsymFrequent\  \dicIPA{{l}{\textopeno}{\textlengthmark}{\textbabygamma}{a}} \dicPos{v}[1] \dicFlx{(‑aði)}[44] \textbf{1.} \dicDirectTranslationCS{hořet, planout};  \dicPhraseIS{það logar á e‑u} \dicFlx{impers} \dicDirectTranslationCS{(co) hoří} \dicExampleIS{Það logar vel á kertinu.} \dicExampleCS{Svíčka dobře hoří.}  \textbf{2.} \dicDirectTranslationCS{hořet, zuřit (město ap.)} \dicIndirectTranslationCS{(o~nepokojích nebo válce)};  \dicPhraseIS{e‑að logar í e‑u} \dicDirectTranslationCS{v~(čem) zuří (co) (ve městě nepokoje ap.)}
\dicEntry[logandi] \dicTerm{log··andi\smash{\textsuperscript{1}}} \dicIPA{{l}{\textopeno}{\textlengthmark}{\textbabygamma}{a}{n}{\textsubring{d}}{\textsci}} \dicPos{adj}[13] \dicFlx{indecl}[1] \dicDirectTranslationCS{hořící, planoucí} \dicExampleIS{logandi kerti} \dicExampleCS{hořící svíce}
\dicEntry[logandi] \dicTerm{log··andi\smash{\textsuperscript{2}}} \dicIPA{{l}{\textopeno}{\textlengthmark}{\textbabygamma}{a}{n}{\textsubring{d}}{\textsci}} \dicPos{adv} \dicDirectTranslationCS{hrozně, strašně} \dicIndirectTranslationCS{(k~zdůraznění)} \dicExampleIS{Hann var logandi hræddur.} \dicExampleCS{Strašně se bál.};  \dicPhraseIS{hvert þó í logandi!} \dicDirectTranslationCS{panebože!, svatá dobroto!}
\dicEntry[logga] \dicTerm{logg|a} \dicIPA{{l}{\textopeno}{\r{g}}{\textlengthmark}{a}} \dicPos{v}[1] \dicFlx{(‑aði)}[1] \dicFieldCat{poč.} \dicDirectTranslationCS{logovat, zaznamenat do souboru};  \dicIdiom{logga}[inn]{ \dicPhraseIS{logga sig inn}} \dicFieldCat{poč.} \dicDirectTranslationCS{(za)logovat se, nalogovat se, přihlásit se};  \dicIdiom{logga}[út]{ \dicPhraseIS{logga sig út}} \dicFieldCat{poč.} \dicDirectTranslationCS{vylogovat se, odlogovat se, odhlásit se}
\dicEntry[logi] \dicTerm{log|i} \dicsymFrequent\  \dicIPA{{l}{\textopeno i}{j}{\textlengthmark}{\textsci}} \dicPos{m}[1] \dicFlx{(‑a, ‑ar)}[1] \dicDirectTranslationCS{plamen, planutí} \dicExampleIS{Loginn hefur dofnað.} \dicExampleCS{Plamen zeslábl.};  \dicPhraseIS{fara eins og logi yfir akur} \dicLangCat{přen.} \dicDirectTranslationCS{přehnat se jako bouře}
\dicEntry[logið] \dicTerm{logið} \dicIPA{{l}{\textopeno i}{j}{\textlengthmark}{\textsci}{\texttheta}} \dicPos{v} \dicFlx{supin} \dicLink{ljúga}
\dicEntry[logn] \dicTerm{logn} \dicsymFrequent\  \dicIPA{{l}{\textopeno}{\r{g}}{\textsubring{n}}} \dicPos{n}[2] \dicFlx{(‑s, ‑)}[5] \dicSynonym*{vindleysa} \dicDirectTranslationCS{bezvětří} \dicIndirectTranslationCS{(Beaufort 0) \mbox{(0--0,2~m/s)}} \dicExampleIS{Það er logn í dag.} \dicExampleCS{Dnes je bezvětří.}
\dicEntry[lognmolla] \dicTerm{logn··moll|a} \dicIPA{{l}{\textopeno}{\r{g}}{n}{m}{\textopeno}{\textsubring{d}}{l}{a}} \dicPos{f}[1] \dicFlx{(‑u)}[5] \textbf{1.} \dicSynonym{andleysi} \dicDirectTranslationCS{nečinnost, neaktivita}  \textbf{2.} \dicDirectTranslationCS{poklidné počasí}
\dicEntry[logsauð] \dicTerm{log··sauð} \dicIPA{{l}{\textopeno}{x}{s}{\oe i}{\texttheta}} \dicPos{v} \dicFlx{ind pf sg 1 pers} \dicLink{logsjóða}
\dicEntry[logsjóða] \dicTerm{log··|sjóða} \dicIPA{{l}{\textopeno}{x}{s}{j}{ou}{ð}{a}} \dicPos{v}[6] \dicFlx{(‑sýð, ‑sauð, ‑suðum, ‑syði, ‑soðið)}[104] \dicFlx{acc} \dicFieldCat{techn.} \dicDirectTranslationCS{svářet, svařovat}
\dicEntry[logsoðið] \dicTerm{log··soðið} \dicIPA{{l}{\textopeno}{x}{s}{\textopeno}{ð}{\textsci}{\texttheta}} \dicPos{v} \dicFlx{supin} \dicLink{logsjóða}
\dicEntry[logsuða] \dicTerm{log··suð|a} \dicIPA{{l}{\textopeno}{x}{s}{\textscy}{ð}{a}} \dicPos{f}[1] \dicFlx{(‑u)}[5] \dicSynonym*{það að logsjóða} \dicDirectTranslationCS{sváření, svařování}
\dicEntry[logsuðum] \dicTerm{log··suðum} \dicIPA{{l}{\textopeno}{x}{s}{\textscy}{ð}{\textscy}{\textsubring{m}}} \dicPos{v} \dicFlx{ind pf pl 1 pers} \dicLink{logsjóða}
\dicEntry[logsyði] \dicTerm{log··syði} \dicIPA{{l}{\textopeno}{x}{s}{\textsci}{ð}{\textsci}} \dicPos{v} \dicFlx{con pf sg 1 pers} \dicLink{logsjóða}
\dicEntry[logsýð] \dicTerm{log··sýð} \dicIPA{{l}{\textopeno}{x}{s}{i}{\texttheta}} \dicPos{v} \dicFlx{ind praes sg 1 pers} \dicLink{logsjóða}
\dicEntry[lok] \dicTerm{lok} \dicsymFrequent\  \dicIPA{{l}{\textopeno}{\textlengthmark}{\r{g}}} \dicPos{n}[2] \dicFlx{(‑s, ‑)}[5] \textbf{1.} \dicDirectTranslationCS{uzávěr, víčko, víko, poklička, poklice} \dicExampleIS{lokið á dósinni} \dicExampleCS{víčko plechovky}  \textbf{2.} \dicPhraseIS{lok} \dicFlx{pl} \dicSynonym{endir} \dicDirectTranslationCS{konec, závěr};  \dicPhraseIS{að lokum} \dicFlx{adv} \dicDirectTranslationCS{nakonec, na závěr};  \dicPhraseIS{e‑að líður undir lok} \dicDirectTranslationCS{(co) se chýlí k~závěru};  \dicPhraseIS{í lokin} \dicFlx{adv} \dicDirectTranslationCS{nakonec}
\dicEntry[loka] \dicTerm{lok|a\smash{\textsuperscript{1}}} \dicIPA{{l}{\textopeno}{\textlengthmark}{\r{g}}{a}} \dicPos{f}[1] \dicFlx{(‑u, ‑ur)}[19] \textbf{1.} \dicSynonym{hespa\smash{\textsuperscript{1}}} \dicDirectTranslationCS{západka, petlice, závora}  \textbf{2.} \dicFieldCat{anat.} \dicDirectTranslationCS{chlopeň}
\dicEntry[loka] \dicTerm{lok|a\smash{\textsuperscript{2}}} \dicsymFrequent\  \dicIPA{{l}{\textopeno}{\textlengthmark}{\r{g}}{a}} \dicPos{v}[1] \dicFlx{(‑aði)}[1] \dicFlx{dat} \dicDirectTranslationCS{(u)zavřít, zavírat} \dicExampleIS{loka glugganum} \dicExampleCS{zavřít okno};  \dicIdiom{loka}[að]{ \dicPhraseIS{loka að sér}} \dicDirectTranslationCS{zavřít za sebou (dveře)};  \dicIdiom{loka}[fyrir]{ \dicPhraseIS{loka fyrir e‑ð}} \dicDirectTranslationCS{zavřít (co), vypnout (co) (elektřinu ap.)};  \dicIdiom{loka}[inni]{ \dicPhraseIS{loka e‑n inni}} \dicDirectTranslationCS{zavřít (koho) (do blázince ap.)};  \dicIdiom{loka}[úti]{ \dicPhraseIS{loka e‑n úti}} \dicDirectTranslationCS{vyloučit (koho) (z~party ap.), nevpustit (koho) dovnitř};  \dicIdiom{lokast}{ \dicPhraseIS{lokast}} \dicFlx{refl} \dicDirectTranslationCS{zavřít se, zavírat se}
\dicEntry[lokaáfangi] \dicTerm{loka··á·fang|i} \dicIPA{{l}{\textopeno}{\textlengthmark}{\r{g}}{a}{au}{f}{au}{\textltailn}{\r{\textObardotlessj}}{\textsci}} \dicPos{m}[1] \dicFlx{(‑a, ‑ar)}[8] \dicDirectTranslationCS{poslední etapa}
\dicEntry[lokaður] \dicTerm{lok··|aður} \dicsymFrequent\  \dicIPA{{l}{\textopeno}{\textlengthmark}{\r{g}}{a}{ð}{\textscy}{\textsubring{r}}} \dicPos{adj}[3] \dicFlx{(f ‑uð)}[3] \textbf{1.} \dicSynonym*{ekki opinn} \dicDirectTranslationCS{zavřený, uzavřený} \dicExampleIS{lokaður vegur} \dicExampleCS{uzavřená cesta}  \textbf{2.} \dicLangCat{přen.} \dicDirectTranslationCS{uzavřený, odměřený, zdrženlivý, rezervovaný}
\dicEntry[lokafrestur] \dicTerm{loka··frest|ur} \dicIPA{{l}{\textopeno}{\textlengthmark}{\r{g}}{a}{f}{r}{\textepsilon}{s}{\textsubring{d}}{\textscy}{\textsubring{r}}} \dicPos{m}[6] \dicFlx{(‑s)}[7] \dicDirectTranslationCS{uzávěrka, (konečný) termín, deadline}
\dicEntry[lokaorð] \dicTerm{loka··orð} \dicIPA{{l}{\textopeno}{\textlengthmark}{\r{g}}{a}{\textopeno}{r}{\texttheta}} \dicPos{n}[2] \dicFlx{(‑s, ‑)}[5] \dicDirectTranslationCS{konečné\,/\addthin závěrečné slovo, doslov};  \dicPhraseIS{eiga lokaorðið} \dicLangCat{přen.} \dicDirectTranslationCS{mít poslední slovo}
\dicEntry[lokapróf] \dicTerm{loka··próf} \dicIPA{{l}{\textopeno}{\textlengthmark}{\r{g}}{a}{p\smash{\textsuperscript{h}}}{r}{ou}{f}} \dicPos{n}[2] \dicFlx{(‑s, ‑)}[5] \dicFieldCat{škol.} \dicDirectTranslationCS{závěrečná zkouška}
\dicEntry[lokaritgerð] \dicTerm{loka··rit·gerð} \dicIPA{{l}{\textopeno}{\textlengthmark}{\r{g}}{a}{r}{\textsci}{\textsubring{d}}{\r{\textObardotlessj}}{\textepsilon}{r}{\texttheta}} \dicPos{f}[7] \dicFlx{(‑ar, ‑ir)}[1] \dicFieldCat{škol.} \dicDirectTranslationCS{závěrečná práce}
\dicEntry[lokasjóður] \dicTerm{loka··sjóð|ur} \dicIPA{{l}{\textopeno}{\textlengthmark}{\r{g}}{a}{s}{j}{ou}{ð}{\textscy}{\textsubring{r}}} \dicPos{m}[9] \dicFlx{(‑s, ‑ir)}[6] \dicFieldCat{bot.} \dicDirectTranslationCS{kokrhel menší} \textit{(l.~{\textLA{Rhinanthus minor}})}  \dicsymPhoto\ 
\dicFigure{71849.jpg}{Lokasjóður}{Lokasjóður - ŠARŽÍK František, Biolib, Copyright/PD}
\dicEntry[lokaæfing] \dicTerm{loka··æf·ing} \dicIPA{{l}{\textopeno}{\textlengthmark}{\r{g}}{a}{a}{i}{v}{i}{\ng}{\r{g}}} \dicPos{f}[4] \dicFlx{(‑ar, ‑ar)}[5] \dicDirectTranslationCS{generální zkouška, generálka}
\dicEntry[lokhljóð] \dicTerm{lok··hljóð} \dicIPA{{l}{\textopeno}{\textlengthmark}{\r{g}}{\textsubring{l}}{j}{ou}{\texttheta}} \dicPos{n}[2] \dicFlx{(‑s, ‑)}[5] \dicFieldCat{jaz.} \dicDirectTranslationCS{ražená souhláska, ploziva}
\dicEntry[loki] \dicTerm{lok|i} \dicIPA{{l}{\textopeno}{\textlengthmark}{\r{\textObardotlessj}}{\textsci}} \dicPos{m}[1] \dicFlx{(‑a, ‑ar)}[1] \textbf{1.} \dicSynonym{hnútur} \dicDirectTranslationCS{uzel, zauzlení}  \textbf{2.} \dicFieldCat{techn.} \dicSynonym*{öryggisloki} \dicDirectTranslationCS{ventil, ventilek}
\dicEntry[lokið] \dicTerm{lokið} \dicIPA{{l}{\textopeno}{\textlengthmark}{\r{\textObardotlessj}}{\textsci}{\texttheta}} \dicPos{v} \dicFlx{supin} \dicLink{ljúka}
\dicEntry[lokinn] \dicTerm{lokinn} \dicsymFrequent\  \dicIPA{{l}{\textopeno}{\textlengthmark}{\r{\textObardotlessj}}{\textsci}{\textsubring{n}}} \dicPos{adj}[6]\dicFlx{}[-6] \dicSynonym{búinn} \dicDirectTranslationCS{skončený, uzavřený} \dicExampleIS{Rannsókn er lokið.} \dicExampleCS{Vyšetřování je uzavřené.};  \dicPhraseIS{e‑m er öllum lokið} \dicFlx{impers} \dicLangCat{přen.} \dicDirectTranslationCS{(kdo) toho má dost, (kdo) je hotový}
\dicEntry[lokka] \dicTerm{lokk|a} \dicIPA{{l}{\textopeno}{h}{\r{g}}{a}} \dicPos{v}[1] \dicFlx{(‑aði)}[1] \dicFlx{acc} \textbf{1.} \dicSynonym{ginna} \dicDirectTranslationCS{(na)lákat, zlákat} \dicExampleIS{lokka e‑n til sín} \dicExampleCS{nalákat (koho) k~sobě}  \textbf{2.} \dicDirectTranslationCS{vylákat, vymámit, vyloudit} \dicExampleIS{lokka e‑ð út úr e‑m} \dicExampleCS{vylákat (co) od (koho)}
\dicEntry[lokkur] \dicTerm{lokk|ur} \dicIPA{{l}{\textopeno}{h}{\r{g}}{\textscy}{\textsubring{r}}} \dicPos{m}[6] \dicFlx{(‑s, ‑ar)}[22] \dicDirectTranslationCS{lokna, kudrna}
\dicEntry[lokræsi] \dicTerm{lok··ræsi} \dicIPA{{l}{\textopeno}{\textlengthmark}{\r{g}}{r}{a}{i}{s}{\textsci}} \dicPos{n}[2] \dicFlx{(‑s, ‑)}[14] \dicLangCat{zast.} \dicDirectTranslationCS{stoka, odpad}
\dicEntry[loks] \dicTerm{loks} \dicsymFrequent\  \dicIPA{{l}{\textopeno}{\r{g}}{s}} \dicPos{adv} \dicSynonym*{að lokum} \dicDirectTranslationCS{nakonec, konečně} \dicExampleIS{Hurðin var loks opnuð.} \dicExampleCS{Dveře se nakonec otevřely.}
\dicEntry[loksins] \dicTerm{loksins} \dicsymFrequent\  \dicIPA{{l}{\textopeno}{\r{g}}{s}{\textsci}{n}{s}} \dicPos{adv} \dicSynonym*{á endanum} \dicDirectTranslationCS{konečně} \dicExampleIS{Fundur er loksins að hefjast.} \dicExampleCS{Schůze konečně začíná.}
\dicEntry[lokun] \dicTerm{lok|un} \dicIPA{{l}{\textopeno}{\textlengthmark}{\r{g}}{\textscy}{\textsubring{n}}} \dicPos{f}[7] \dicFlx{(‑unar, ‑anir)}[8] \textbf{1.} \dicSynonym*{það að loka} \dicDirectTranslationCS{zavírání, (u)zavření} \dicExampleIS{lokun búða} \dicExampleCS{zavření obchodu}  \textbf{2.} \dicSynonym*{það að setja hindrun} \dicDirectTranslationCS{uzavírka (cesty ap.)} \dicExampleIS{lokun vegar} \dicExampleCS{uzavírka cesty}
\dicEntry[lokunartími] \textls[15]{\dicTerm{lokunar··tím|i} \dicIPA{{l}{\textopeno}{\r{g}}{\textscy}{n}{a}{\textsubring{r}}{t\smash{\textsuperscript{h}}}{i}{m}{\textsci}} \dicPos{m}[1] \dicFlx{(‑a)}[3] \dicDirectTranslationCS{zavírací doba}} \dicAntonym{opnunartími}
\dicEntry[lopapeysa] \dicTerm{lopa··peys|a} \dicIPA{{l}{\textopeno}{\textlengthmark}{\textsubring{b}}{a}{p\smash{\textsuperscript{h}}}{ei}{s}{a}} \dicPos{f}[1] \dicFlx{(‑u, ‑ur)}[7] \dicIndirectTranslationCS{pletený vlněný svetr s~typickými islandskými vzory}
\dicEntry[lopi] \dicTerm{lop|i} \dicIPA{{l}{\textopeno}{\textlengthmark}{\textsubring{b}}{\textsci}} \dicPos{m}[1] \dicFlx{(‑a, ‑ar)}[1] \textbf{1.} \dicSynonym*{ullarþráður} \dicDirectTranslationCS{mykaná vlna}  \textbf{2.} \dicFieldCat{med.} \dicSynonym{bjúgur} \dicDirectTranslationCS{edém, otok};  \dicIdiom{lopi}{ \dicPhraseIS{spinna\,/\addthin teygja lopann}} \dicLangCat{přen.} \dicDirectTranslationCS{pustit si pusu na špacír}
\dicEntry[loppa] \dicTerm{lopp|a} \dicIPA{{l}{\textopeno}{h}{\textsubring{b}}{a}} \dicPos{f}[1] \dicFlx{(‑u, ‑ur)}[19] \dicDirectTranslationCS{tlapa, tlapka}
\dicEntry[los] \dicTerm{los} \dicIPA{{l}{\textopeno}{\textlengthmark}{s}} \dicPos{n}[2] \dicFlx{(‑s, ‑)}[5] \dicDirectTranslationCS{nestabilita, nestabilnost, nesoudržnost} \dicExampleIS{það er los á e‑u} \dicExampleCS{(co) je nestabilní}
\dicEntry[losa] \dicTerm{los|a} \dicsymFrequent\  \dicIPA{{l}{\textopeno}{\textlengthmark}{s}{a}} \dicPos{v}[1] \dicFlx{(‑aði)}[1] \dicFlx{acc} \textbf{1.} \dicDirectTranslationCS{uvolnit, uvolňovat, povolit, povolovat (šroubek ap.)} \dicExampleIS{Hann losar hjólið af bílnum.} \dicExampleCS{Uvolňuje kolo u~auta.}  \textbf{2.} \dicDirectTranslationCS{uvolnit, uvolňovat (pokoj ap.)} \dicExampleIS{losa hótelherbergi} \dicExampleCS{uvolnit hotelový pokoj}  \textbf{3.} \dicDirectTranslationCS{vyprázdnit, vyprazdňovat} \dicExampleIS{losa flöskuna} \dicExampleCS{vyprázdnit láhev}  \textbf{4.} \dicPhraseIS{e‑að losar (20 tonn)} \dicDirectTranslationCS{(co) má něco málo přes (20 tun), (co) přesahuje (20 tun)};  \dicIdiom{losa}[um]{ \dicPhraseIS{losa um e‑ð}} \dicDirectTranslationCS{uvolnit (co), zmírnit (co) (o\nobreak mezení ap.)};  \dicIdiom{losa}[úr]{ \dicPhraseIS{losa e‑n úr e‑u}} \dicDirectTranslationCS{uvolnit (koho) z~(čeho), vymanit (koho) z~(čeho)};  \dicIdiom{losa}[við]{ \dicPhraseIS{losa e‑n við e‑ð}} \dicDirectTranslationCS{pomoct (komu) s~(čím)}; { \dicPhraseIS{losa sig við e‑ð}} \dicDirectTranslationCS{zbavit se (čeho)}
\dicEntry[losaralegur] \dicTerm{losara··legur} \dicIPA{{l}{\textopeno}{\textlengthmark}{s}{a}{r}{a}{l}{\textepsilon}{\textbabygamma}{\textscy}{\textsubring{r}}} \dicPos{adj}[1]\dicFlx{}[-8] \dicSynonym{sundurlaus} \dicDirectTranslationCS{náhod\-ný, nahodilý, neplánovaný}
\dicEntry[losna] \dicTerm{losn|a} \dicsymFrequent\  \dicIPA{{l}{\textopeno}{s}{\textsubring{d}}{n}{a}} \dicPos{v}[1] \dicFlx{(‑aði)}[1] \textbf{1.} \dicDirectTranslationCS{uvolnit se, uprázdnit se (místo ap.)} \dicExampleIS{Starfið losnar bráðlega.} \dicExampleCS{Zaměstnání se brzy uvolní.}  \textbf{2.} \dicDirectTranslationCS{uvolnit\,/\addthin uvolňovat se, povolit, povolovat (šroubek ap.)} \dicExampleIS{Það hefur losnað hjól af bílnum.} \dicExampleCS{U~auta se uvolnilo kolo.}  \textbf{3.} \dicDirectTranslationCS{uvolnit\,/\addthin uvolňovat se (z~práce ap.)} \dicExampleIS{losna úr vinnunni} \dicExampleCS{uvolnit se z~práce};  \dicIdiom{losna}[undan]{ \dicPhraseIS{losna undan e‑u}} \dicDirectTranslationCS{vyhnout\,/\addthin vyhýbat se (čemu), uniknout\,/\addthin unikat (čemu)} \dicExampleIS{losna undan mörgum skyldum} \dicExampleCS{vyhýbat se mnoha povinnostem};  \dicIdiom{losna}[við]{ \dicPhraseIS{losna við e‑ð}} \dicSynonym*{koma e‑u af sér} \dicDirectTranslationCS{zbavit se (čeho)}
\dicEntry[lost] \dicTerm{lost} \dicIPA{{l}{\textopeno}{s}{\textsubring{d}}} \dicPos{n}[2] \dicFlx{(‑s, ‑)}[5] \textbf{1.} \dicFieldCat{med.} \dicDirectTranslationCS{šok}  \textbf{2.} \dicLangCat{hovor.} \dicDirectTranslationCS{šok, překvapení} \dicExampleIS{Ég var alveg í losti á eftir.} \dicExampleCS{Byl jsem po tom v~šoku.}
\dicEntry[lostafullur] \dicTerm{losta··|fullur} \dicIPA{{l}{\textopeno}{s}{\textsubring{d}}{a}{f}{\textscy}{\textsubring{d}}{l}{\textscy}{\textsubring{r}}} \dicPos{adj}[10] \dicFlx{(comp ‑fyllri, sup ‑fyllstur)}[7] \dicSynonym*{munaðargjarn} \dicDirectTranslationCS{chlípný, žádostivý, chtivý, vilný}
\dicEntry[losti] \dicTerm{lost|i} \dicIPA{{l}{\textopeno}{s}{\textsubring{d}}{\textsci}} \dicPos{m}[1] \dicFlx{(‑a)}[3] \dicSynonym{girnd} \dicDirectTranslationCS{chlípnost, žádostivost, chtíč, vilnost}
\dicEntry[lostið] \dicTerm{lostið} \dicIPA{{l}{\textopeno}{s}{\textsubring{d}}{\textsci}{\texttheta}} \dicPos{v} \dicFlx{supin} \dicLink{ljósta}
\dicEntry[lostinn] \dicTerm{lostinn} \dicIPA{{l}{\textopeno}{s}{\textsubring{d}}{\textsci}{\textsubring{n}}} \dicPos{adj}[6]\dicFlx{}[-6] \dicPhraseIS{skelfingu lostinn} \dicFlx{adj} \dicDirectTranslationCS{vyděšený, vystrašený};  \dicPhraseIS{vera furðu lostinn yfir e‑u} \dicDirectTranslationCS{oněmět úžasem nad (čím)};  \dicPhraseIS{vera þrumu lostinn yfir e‑u} \dicDirectTranslationCS{být šokován (čím)}
\dicEntry[lostæti] \dicTerm{lost··æti} \dicIPA{{l}{\textopeno}{s}{\textsubring{d}}{a}{i}{\textsubring{d}}{\textsci}} \dicPos{n}[2] \dicFlx{(‑s, ‑)}[14] \dicSynonym{góðgæti} \dicDirectTranslationCS{pochoutka, lahůdka, delikatesa}
\dicEntry[losun] \dicTerm{los|un} \dicIPA{{l}{\textopeno}{\textlengthmark}{s}{\textscy}{\textsubring{n}}} \dicPos{f}[7] \dicFlx{(‑unar)}[9] \textbf{1.} \dicSynonym*{losun sorps o.fl.} \dicDirectTranslationCS{vypouštění, vyprazdňování (odpadu, odpadků ap.)}  \textbf{2.} \dicSynonym*{afferming} \dicDirectTranslationCS{vykládka, vyložení} \dicExampleIS{losun á farminum} \dicExampleCS{vyložení nákladu}
\dicEntry[lota] \dicTerm{lot|a} \dicIPA{{l}{\textopeno}{\textlengthmark}{\textsubring{d}}{a}} \dicPos{f}[1] \dicFlx{(‑u, ‑ur)}[19] \dicDirectTranslationCS{období, cyklus, perioda, (časový) úsek} \dicExampleIS{í einni lotu} \dicExampleCS{na jeden zátah}
\dicEntry[lotið] \dicTerm{lotið} \dicIPA{{l}{\textopeno}{\textlengthmark}{\textsubring{d}}{\textsci}{\texttheta}} \dicPos{v} \dicFlx{supin} \dicLink{lúta}
\dicEntry[lotinn] \dicTerm{lotinn} \dicIPA{{l}{\textopeno}{\textlengthmark}{\textsubring{d}}{\textsci}{\textsubring{n}}} \dicPos{adj}[6]\dicFlx{}[-2] \dicSynonym{boginn} \dicDirectTranslationCS{shrbený} \dicExampleIS{lotinn í herðum} \dicExampleCS{shrbený v~zádech}
\dicEntry[lotning] \dicTerm{lot··ning} \dicIPA{{l}{\textopeno}{h}{\textsubring{d}}{n}{i}{\ng}{\r{g}}} \dicPos{f}[4] \dicFlx{(‑ar)}[7] \dicSynonym{virðing} \dicDirectTranslationCS{úcta, uctívání} \dicExampleIS{sýna e‑m lotningu} \dicExampleCS{projevovat ke (komu) úctu}
\begin{xtolerant}{}{1pt}
\dicEntry[lotningarhreimur] \dicTerm{lotningar··hreim|ur}\addthinS\dicIPA{{l}{\textopeno}{h}{\textsubring{d}}{n}{i}{\ng}{\r{g}}{a}{\textsubring{r}}{ei}{m}{\textscy}{\textsubring{r}}}\addthinS\dicPos{m}[6]\addthinS\dicFlx{(‑s, ‑ar)}[22] \dicDirectTranslationCS{pochlebovačný tón}
\end{xtolerant}
\dicEntry[lottó] \dicTerm{lottó} \dicIPA{{l}{\textopeno}{h}{\textsubring{d}}{ou}} \dicPos{n}[2] \dicFlx{(‑s, ‑)}[34] \dicDirectTranslationCS{loterie, sportka}
\dicEntry[lotugræðgi] \dicTerm{lotu··græðg|i} \dicIPA{{l}{\textopeno}{\textlengthmark}{\textsubring{d}}{\textscy}{\r{g}}{r}{a}{i}{ð}{\r{g}}{\textsci}} \dicPos{f}[3] \dicFlx{(‑i)}[3] \dicFieldCat{med.} \dicDirectTranslationCS{bulimie}
\dicEntry[lotukerfi] \dicTerm{lotu··kerfi} \dicIPA{{l}{\textopeno}{\textlengthmark}{\textsubring{d}}{\textscy}{c\smash{\textsuperscript{h}}}{\textepsilon}{r}{v}{\textsci}} \dicPos{n}[2] \dicFlx{(‑s)}[20] \dicFieldCat{chem.} \dicDirectTranslationCS{periodická tabulka\,/\addthin soustava prvků}
\dicEntry[lotutúlkun] \dicTerm{lotu··túlk|un} \dicIPA{{l}{\textopeno}{\textlengthmark}{\textsubring{d}}{\textscy}{t\smash{\textsuperscript{h}}}{u}{\textsubring{l}}{\r{g}}{\textscy}{\textsubring{n}}} \dicPos{f}[7] \dicFlx{(‑unar, ‑anir)}[8] \dicDirectTranslationCS{konsekutivní tlumočení}
\dicEntry[lóa] \dicTerm{ló|a} \dicIPA{{l}{ou}{\textlengthmark}{a}} \dicPos{f}[1] \dicFlx{(‑u, ‑ur)}[7] \dicFieldCat{zool.} \dicSynonym{heiðlóa} \dicDirectTranslationCS{kulík, kulík zlatý} \textit{(l.~{\textLA{Pluvialis apricaria}})}  \dicsymPhoto\ 
\dicFigure{21284.jpg}{Lóa}{Lóa - Karney, Lee, Biolib, PD}
\dicEntry[lóð] \dicTerm{lóð\smash{\textsuperscript{1}}} \dicIPA{{l}{ou}{\textlengthmark}{\texttheta}} \dicPos{f}[7] \dicFlx{(‑ar, ‑ir)}[1] \dicSynonym*{landspilda} \dicDirectTranslationCS{parcela, pozemek} \dicExampleIS{Húsið stendur á lóðinni.} \dicExampleCS{Dům stojí na pozemku.}
\dicEntry[lóð] \dicTerm{lóð\smash{\textsuperscript{2}}} \dicIPA{{l}{ou}{\textlengthmark}{\texttheta}} \dicPos{f}[7] \dicFlx{(‑ar, ‑ir)}[1] \dicFieldCat{nám.} \dicSynonym{fiskilína} \dicDirectTranslationCS{síť s~mnoha háčky};  \dicPhraseIS{leggja lóðirnar} {\textbf{a.}} \dicDirectTranslationCS{položit síť} \dicIndirectTranslationCS{(síť s~mnoha háčky)};  {\textbf{b.}} \dicSynonym*{selja upp af sjóveiki} \dicDirectTranslationCS{zvracet} \dicIndirectTranslationCS{(kvůli mořské nemoci)}
\dicEntry[lóð] \dicTerm{lóð\smash{\textsuperscript{3}}} \dicIPA{{l}{ou}{\textlengthmark}{\texttheta}} \dicPos{n}[2] \dicFlx{(‑s, ‑)}[5] \dicDirectTranslationCS{závaží} \dicExampleIS{lóð á vog} \dicExampleCS{závaží na váze};  \dicPhraseIS{það er lóðið} \dicLangCat{přen.} \dicDirectTranslationCS{to je jádro věci}
\dicEntry[lóða] \dicTerm{lóð|a} \dicIPA{{l}{ou}{\textlengthmark}{ð}{a}} \dicPos{v}[1] \dicFlx{(‑aði)}[1] \dicFlx{acc} \textbf{1.} \dicFieldCat{techn.} \dicDirectTranslationCS{letovat, pájet} \dicExampleIS{lóða e‑ð saman} \dicExampleCS{sletovat (co)}  \textbf{2.} \dicSynonym*{mæla dýpt} \dicDirectTranslationCS{sondovat (dno nebo moře pro výskyt ryb)};  \dicPhraseIS{lóða á fisk(i)} \dicFieldCat{nám.} \dicDirectTranslationCS{hledat ryby sonarem}
\dicEntry[lóðbolti] \dicTerm{lóð··bolt|i} \dicIPA{{l}{ou}{ð}{\textsubring{b}}{\textopeno}{\textsubring{l}}{\textsubring{d}}{\textsci}} \dicPos{m}[1] \dicFlx{(‑a, ‑ar)}[1] \dicFieldCat{techn.} \dicDirectTranslationCS{páječka, pájedlo}
\dicEntry[lóðréttur] \dicTerm{lóð··réttur} \dicIPA{{l}{ou}{ð}{r}{j}{\textepsilon}{h}{\textsubring{d}}{\textscy}{\textsubring{r}}} \dicPos{adj}[1]\dicFlx{}[-10] \dicDirectTranslationCS{svislý, vertikální, kolmý}
\dicEntry[lóðun] \dicTerm{lóð|un} \dicIPA{{l}{ou}{\textlengthmark}{ð}{\textscy}{\textsubring{n}}} \dicPos{f}[7] \dicFlx{(‑unar)}[9] \dicFieldCat{techn.} \dicSynonym*{það að lóða} \dicDirectTranslationCS{letování, pájení}
\dicEntry[lófaklapp] \dicTerm{lófa··klapp} \dicIPA{{l}{ou}{\textlengthmark}{v}{a}{k\smash{\textsuperscript{h}}}{l}{a}{h}{\textsubring{b}}} \dicPos{n}[2] \dicFlx{(‑s)}[2] \dicDirectTranslationCS{potlesk, aplaus}
\dicEntry[lófatak] \dicTerm{lófa··|tak} \dicIPA{{l}{ou}{\textlengthmark}{v}{a}{t\smash{\textsuperscript{h}}}{a}{\r{g}}} \dicPos{n}[2] \dicFlx{(‑taks, ‑tök)}[8] \dicSynonym{lófaklapp} \dicDirectTranslationCS{potlesk, aplaus}
\dicEntry[lófatölva] \dicTerm{lófa··tölv|a} \dicIPA{{l}{ou}{\textlengthmark}{v}{a}{t\smash{\textsuperscript{h}}}{\oe}{l}{v}{a}} \dicPos{f}[1] \dicFlx{(‑u, ‑ur)}[7] \dicFieldCat{poč.} \dicDirectTranslationCS{palmtop, PDA, kapesní počítač}
\dicEntry[lófi] \dicTerm{lóf|i} \dicsymFrequent\  \dicIPA{{l}{ou}{\textlengthmark}{v}{\textsci}} \dicPos{m}[1] \dicFlx{(‑a, ‑ar)}[1] \dicDirectTranslationCS{dlaň} \dicExampleIS{geyma hring í lófa sínum} \dicExampleCS{držet prsten ve své dlani}
\dicEntry[lófótur] \textls[15]{\dicTerm{ló··|fótur} \dicIPA{{l}{ou}{\textlengthmark}{f}{ou}{\textsubring{d}}{\textscy}{\textsubring{r}}} \dicPos{m}[13] \dicFlx{(‑fótar, ‑fætur)}[5] \dicFieldCat{bot.} \dicDirectTranslationCS{prustka obecná}} \textit{(l.~{\textLA{Hippuris vulgaris}})}  \dicsymPhoto\ 
\dicFigure{3847.jpg}{Lófótur}{Lófótur - Zicha Ondřej, Biolib, Copyright/CC-BY-NC}
\dicEntry[lóga] \dicTerm{lóg|a} \dicIPA{{l}{ou}{\textlengthmark}{a}} \dicPos{v}[1] \dicFlx{(‑aði)}[34] \dicFlx{dat} \textbf{1.} \dicSynonym{slátra} \dicDirectTranslationCS{porazit, porážet (ovci ap.)} \dicExampleIS{lóga fé} \dicExampleCS{porážet ovce}  \textbf{2.} \dicSynonym*{aflífa} \dicDirectTranslationCS{utratit, uspat (psa ap.)} \dicExampleIS{lóga hundi} \dicExampleCS{uspat psa}
\dicEntry[lógaritmi] \dicTerm{lógaritm|i} \dicIPA{{l}{ou}{\textlengthmark}{\r{g}}{a}{r}{\textsci}{h}{\textsubring{d}}{m}{\textsci}} \dicPos{m}[1] \dicFlx{(‑a, ‑ar)}[1] \dicFieldCat{mat.} \dicDirectTranslationCS{logaritmus}
\dicEntry[lómur] \dicTerm{lóm|ur} \dicIPA{{l}{ou}{\textlengthmark}{m}{\textscy}{\textsubring{r}}} \dicPos{m}[6] \dicFlx{(‑s, ‑ar)}[24] \textbf{1.} \dicFieldCat{zool.} \dicDirectTranslationCS{potáplice malá} \textit{(l.~{\textLA{Gavia stellata}})}  \dicsymPhoto\   \textbf{2.} \dicSynonym{kveinstafir} \dicDirectTranslationCS{nářky, stížnosti};  \dicPhraseIS{berja\,/\addthin lemja lóminn} \dicDirectTranslationCS{stěžovat si, naříkat}
\dicFigure{21071.jpg}{Lómur}{Lómur - Bowman, Tim, Biolib, PD}
\dicEntry[lón] \dicTerm{lón} \dicIPA{{l}{ou}{\textlengthmark}{\textsubring{n}}} \dicPos{n}[2] \dicFlx{(‑s, ‑)}[5] \dicDirectTranslationCS{laguna};  \dicPhraseIS{Bláa lónið} \dicFieldCat{geog.} \dicDirectTranslationCS{Modrá laguna} \dicIndirectTranslationCS{(laguna v~blízkosti Keflavíku)}
\dicEntry[lóuþræll] \dicTerm{lóu··þræl|l} \dicIPA{{l}{ou}{\textlengthmark}{\textscy}{\texttheta}{r}{a}{i}{\textsubring{d}}{\textsubring{l}}} \dicPos{m}[6] \dicFlx{(‑s, ‑ar)}[51] \dicFieldCat{zool.} \dicDirectTranslationCS{jespák, jespák obecný} \textit{(l.~{\textLA{Calidris alpina}})}  \dicsymPhoto\ 
\dicFigure{21101.jpg}{Lóuþræll}{Lóuþræll - Bowman, Tim, Biolib, PD}
\dicEntry[luggarasegl] \dicTerm{luggara··segl} \dicIPA{{l}{\textscy}{\r{g}}{\textlengthmark}{a}{r}{a}{s}{\textepsilon}{\r{g}}{\textsubring{l}}} \dicPos{n}[2] \dicFlx{(‑s, ‑)}[5] \dicFieldCat{nám.} \dicDirectTranslationCS{lugrová plachta}
\dicEntry[lugt] \dicTerm{lugt} \dicIPA{{l}{\textscy}{x}{\textsubring{d}}} \dicPos{f}[7] \dicFlx{(‑ar, ‑ir)}[1] \dicLink{lukt}
\dicEntry[lugum] \dicTerm{lugum} \dicIPA{{l}{\textscy}{\textlengthmark}{\textbabygamma}{\textscy}{\textsubring{m}}} \dicPos{v} \dicFlx{ind pf pl 1 pers} \dicLink{ljúga}
\dicEntry[lukið] \dicTerm{lukið} \dicIPA{{l}{\textscy}{\textlengthmark}{\r{\textObardotlessj}}{\textsci}{\texttheta}} \dicPos{v} \dicFlx{supin} \dicLink{lykja}
\dicEntry[lukka] \dicTerm{lukk|a} \dicIPA{{l}{\textscy}{h}{\r{g}}{a}} \dicPos{f}[1] \dicFlx{(‑u)}[5] \dicSynonym{hamingja} \dicDirectTranslationCS{štěstí, zdar} \dicExampleIS{óska e‑m til lukku með afmælið} \dicExampleCS{popřát (komu) štěstí k~narozeninám}
\dicEntry[lukkast] \dicTerm{lukk|ast} \dicIPA{{l}{\textscy}{h}{\r{g}}{a}{s}{\textsubring{d}}} \dicPos{v}[1] \dicFlx{(‑aðist)}[97] \dicFlx{refl} \dicSynonym{heppnast} \dicDirectTranslationCS{(po)dařit se, zdařit se} \dicExampleIS{Hjónabandið lukkast vel.} \dicExampleCS{Manželství se daří.}
\dicEntry[lukkulegur] \dicTerm{lukku··legur} \dicIPA{{l}{\textscy}{h}{\r{g}}{\textscy}{l}{\textepsilon}{\textbabygamma}{\textscy}{\textsubring{r}}} \dicPos{adj}[1]\dicFlx{}[-8] \dicSynonym{ánægður} \dicDirectTranslationCS{šťastný, spokojený, potěšený}
\dicEntry[lukt] \dicTerm{lukt}\dicTerm{, lugt} \dicIPA{{l}\-{\textscy}\-{x}\-{\textsubring{d}}\-} \dicPos{f}[7] \dicFlx{(‑ar, ‑ir)}[1] \textbf{1.} \dicDirectTranslationCS{lucerna, svítilna} \dicExampleIS{vasalukt} \dicExampleCS{kapesní svítilna}  \textbf{2.} \dicDirectTranslationCS{přední světlo (u~auta, kola ap.)}
\dicEntry[lukti] \dicTerm{lukti} \dicIPA{{l}{\textscy}{x}{\textsubring{d}}{\textsci}} \dicPos{v} \dicFlx{ind pf sg 1 pers} \dicLink{lykja}
\dicEntry[luktum] \dicTerm{luktum} \dicIPA{{l}{\textscy}{x}{\textsubring{d}}{\textscy}{\textsubring{m}}} \dicPos{v} \dicFlx{ind pf pl 1 pers} \dicLink{lykja}
\dicEntry[luktur] \dicTerm{luktur} \dicIPA{{l}{\textscy}{x}{\textsubring{d}}{\textscy}{\textsubring{r}}} \dicPos{adj}[1]\dicFlx{}[-10] \dicDirectTranslationCS{(u)zavřený, obklopený}
\dicEntry[lukum] \dicTerm{lukum} \dicIPA{{l}{\textscy}{\textlengthmark}{\r{g}}{\textscy}{\textsubring{m}}} \dicPos{v} \dicFlx{ind pf pl 1 pers} \dicLink{ljúka}
\dicEntry[luma] \dicTerm{lum|a} \dicIPA{{l}{\textscy}{\textlengthmark}{m}{a}} \dicPos{v}[1] \dicFlx{(‑aði)}[1] \dicPhraseIS{luma á e‑u} \dicDirectTranslationCS{ukrýt si (co), schovat si (co), ulít si (co)} \dicExampleIS{Hann lumar á peningum.} \dicExampleCS{Ukrývá si peníze.}
\dicEntry[lumbra] \dicTerm{lumbr|a} \dicIPA{{l}{\textscy}{m}{\textsubring{b}}{r}{a}} \dicPos{v}[1] \dicFlx{(‑aði)}[1] \dicDirectTranslationCS{zašantročit, zastrčit} \dicExampleIS{lumbra e‑u e‑s staðar} \dicExampleCS{zašantročit (co kde)};  \dicIdiom{lumbra}[á]{ \dicPhraseIS{lumbra á e‑m}} \dicSynonym{berja} \dicDirectTranslationCS{(z)bít (koho), (z)mlátit (koho)}
\dicEntry[lumma] \dicTerm{lumm|a} \dicIPA{{l}{\textscy}{m}{\textlengthmark}{a}} \dicPos{f}[1] \dicFlx{(‑u, ‑ur)}[7] \dicFieldCat{kulin.} \dicDirectTranslationCS{lívanec, lívaneček};  \dicPhraseIS{e‑að selst eins og heitar lummur} \dicFlx{refl} \dicLangCat{přen.} \dicDirectTranslationCS{(co) jde na dračku};  \dicPhraseIS{það er gömul lumma} \dicLangCat{přen.} \dicDirectTranslationCS{to je ohraná písnička}
\dicEntry[lummulegur] \dicTerm{lummu··legur} \dicIPA{{l}{\textscy}{m}{\textlengthmark}{\textscy}{l}{\textepsilon}{\textbabygamma}{\textscy}{\textsubring{r}}} \dicPos{adj}[1]\dicFlx{}[-8] \dicLangCat{hovor.} \dicSynonym{tilkomulítill} \dicDirectTranslationCS{nijaký, nezáživný, nudný}
\dicEntry[lund] \dicTerm{lund} \dicsymFrequent\  \dicIPA{{l}{\textscy}{n}{\textsubring{d}}} \dicPos{f}[7] \dicFlx{(‑ar, ‑ir)}[1] \textbf{1.} \dicSynonym{geð} \dicDirectTranslationCS{povaha, temperament} \dicExampleIS{vera léttur í lund} \dicExampleCS{být šťastný}  \textbf{2.} \dicSynonym{háttur\smash{\textsuperscript{1}}} \dicDirectTranslationCS{způsob}
\dicEntry[Lundey] \dicTerm{Lund··ey} \dicIPA{{l}{\textscy}{n}{\textsubring{d}}{ei}} \dicPos{f}[4] \dicFlx{(‑jar)}[17] \dicFieldCat{geog.} \dicDirectTranslationCS{Lundey} \dicIndirectTranslationCS{(ostrov severně od Islandu)}
\dicEntry[lundi] \dicTerm{lund|i} \dicIPA{{l}{\textscy}{n}{\textsubring{d}}{\textsci}} \dicPos{m}[1] \dicFlx{(‑a, ‑ar)}[1] \dicFieldCat{zool.} \dicDirectTranslationCS{papuchálek, papuchalk, papuchalk ploskozobý\,/\addthin bělobradý} \textit{(l.~{\textLA{Fratercula arctica}})}  \dicsymPhoto\ 
\dicFigure{ds_image_lundi_0_1.jpg}{Lundi}{Lundi - Tgo2002, CC BY-SA 2.5}
\dicEntry[lundir] \dicTerm{lundir} \dicIPA{{l}{\textscy}{n}{\textsubring{d}}{\textsci}{\textsubring{r}}} \dicPos{f}[7] \dicFlx{pl}[2] \dicDirectTranslationCS{jemné filety, svíčková}
\dicEntry[lundur] \dicTerm{lund|ur} \dicIPA{{l}{\textscy}{n}{\textsubring{d}}{\textscy}{\textsubring{r}}} \dicPos{m}[10] \dicFlx{(‑ar, ‑ir)}[4] \dicDirectTranslationCS{háj, lesík}
\dicEntry[Lundúnir] \dicTerm{Lundúnir} \dicIPA{{l}{\textscy}{n}{\textsubring{d}}{u}{n}{\textsci}{\textsubring{r}}} \dicPos{f}[12] \dicFlx{pl}[6] \dicFieldCat{geog.} \dicDirectTranslationCS{Londýn} \dicIndirectTranslationCS{(hlavní město Spojeného království)}
\dicEntry[lunga] \dicTerm{lung|a} \dicIPA{{l}{u}{\ng}{\r{g}}{a}} \dicPos{n}[1] \dicFlx{(‑a, ‑u)}[2] \dicFieldCat{anat.} \dicDirectTranslationCS{plíce}
\dicEntry[lungnabólga] \dicTerm{lungna··bólg|a} \dicIPA{{l}{u}{\ng}{n}{a}{\textsubring{b}}{ou}{l}{\r{g}}{a}} \dicPos{f}[1] \dicFlx{(‑u)}[5] \dicFieldCat{med.} \dicDirectTranslationCS{zápal plic, pneumonie}
\dicEntry[lungnakvef] \dicTerm{lungna··kvef} \dicIPA{{l}{u}{\ng}{n}{a}{k\smash{\textsuperscript{h}}}{v}{\textepsilon}{f}} \dicPos{n}[2] \dicFlx{(‑s)}[2] \dicFieldCat{med.} \dicDirectTranslationCS{zánět průdušek, bronchitida}
\dicEntry[lungnapípa] \dicTerm{lungna··píp|a} \dicIPA{{l}{u}{\ng}{n}{a}{p\smash{\textsuperscript{h}}}{i}{\textsubring{b}}{a}} \dicPos{f}[1] \dicFlx{(‑u, ‑ur)}[19] \dicFieldCat{anat.} \dicDirectTranslationCS{průduška}
\dicEntry[luntalegur] \dicTerm{lunta··legur} \dicIPA{{l}{\textscy}{\textsubring{n}}{\textsubring{d}}{a}{l}{\textepsilon}{\textbabygamma}{\textscy}{\textsubring{r}}} \dicPos{adj}[1]\dicFlx{}[-8] \dicSynonym{ólundarlegur} \dicDirectTranslationCS{mrzutý, nevrlý}
\dicEntry[luralegur] \dicTerm{lura··legur} \dicIPA{{l}{\textscy}{\textlengthmark}{r}{a}{l}{\textepsilon}{\textbabygamma}{\textscy}{\textsubring{r}}} \dicPos{adj}[1]\dicFlx{}[-8] \dicSynonym*{þyngslalegur} \dicDirectTranslationCS{nemotorný, neohrabaný}
\dicEntry[lustum] \dicTerm{lustum} \dicIPA{{l}{\textscy}{s}{\textsubring{d}}{\textscy}{\textsubring{m}}} \dicPos{v} \dicFlx{ind pf pl 1 pers} \dicLink{ljósta}
\dicEntry[lutum] \dicTerm{lutum} \dicIPA{{l}{\textscy}{\textlengthmark}{\textsubring{d}}{\textscy}{\textsubring{m}}} \dicPos{v} \dicFlx{ind pf pl 1 pers} \dicLink{lúta}
\dicEntry[lúalegur] \dicTerm{lúa··legur} \dicIPA{{l}{u}{\textlengthmark}{a}{l}{\textepsilon}{\textbabygamma}{\textscy}{\textsubring{r}}} \dicPos{adj}[1]\dicFlx{}[-8] \textbf{1.} \dicSynonym{þreytulegur} \dicDirectTranslationCS{vyčerpaný, unavený}  \textbf{2.} \dicSynonym{ódrengilegur} \dicDirectTranslationCS{hanebný, sprostý} \dicExampleIS{lúalegar getsakir} \dicExampleCS{hanebné narážky}
\dicEntry[lúð] \dicTerm{lúð} \dicIPA{{l}{u}{\textlengthmark}{\texttheta}} \dicPos{v} \dicFlx{supin} \dicLink{lýja}
\dicEntry[lúða] \dicTerm{lúð|a} \dicIPA{{l}{u}{\textlengthmark}{ð}{a}} \dicPos{f}[1] \dicFlx{(‑u, ‑ur)}[19] \dicFieldCat{zool.} \dicSynonym{flyðra} \dicDirectTranslationCS{platýs, platýs obecný} \textit{(l.~{\textLA{Hippoglossus hippoglossus}})}
\dicEntry[lúðalegur] \dicTerm{lúða··legur} \dicIPA{{l}{u}{\textlengthmark}{ð}{a}{l}{\textepsilon}{\textbabygamma}{\textscy}{\textsubring{r}}} \dicPos{adj}[1]\dicFlx{}[-8] \dicDirectTranslationCS{slizký, odporný}
\dicEntry[lúði] \dicTerm{lúð|i\smash{\textsuperscript{1}}} \dicIPA{{l}{u}{\textlengthmark}{ð}{\textsci}} \dicPos{m}[1] \dicFlx{(‑a, ‑ar)}[1] \dicDirectTranslationCS{křupan, balík}
\dicEntry[lúði] \dicTerm{lúði\smash{\textsuperscript{2}}} \dicIPA{{l}{u}{\textlengthmark}{ð}{\textsci}} \dicPos{v} \dicFlx{ind pf sg 1 pers} \dicLink{lýja}
\dicEntry[lúðrasveit] \dicTerm{lúðra··sveit} \dicIPA{{l}{u}{ð}{r}{a}{s}{v}{ei}{\textsubring{d}}} \dicPos{f}[7] \dicFlx{(‑ar, ‑ir)}[1] \dicFieldCat{hud.} \dicDirectTranslationCS{dechová kapela, dechovka}
\dicEntry[lúðum] \dicTerm{lúðum} \dicIPA{{l}{u}{\textlengthmark}{ð}{\textscy}{\textsubring{m}}} \dicPos{v} \dicFlx{ind pf pl 1 pers} \dicLink{lýja}
\dicEntry[lúður] \dicTerm{lúð|ur} \dicIPA{{l}{u}{\textlengthmark}{ð}{\textscy}{\textsubring{r}}} \dicPos{m}[5] \dicFlx{(‑urs, ‑rar)}[1] \dicFieldCat{hud.} \dicDirectTranslationCS{trumpeta} \dicExampleIS{blása í lúður} \dicExampleCS{foukat na trumpetu}
\dicEntry[lúðurþeytari] \dicTerm{lúður··þeyt·ar|i} \dicIPA{{l}{u}{\textlengthmark}{ð}{\textscy}{\textsubring{r}}{\texttheta}{ei}{\textsubring{d}}{a}{r}{\textsci}} \dicPos{m}[1] \dicFlx{(‑a, ‑ar)}[13] \dicDirectTranslationCS{trubač(ka)}
\dicEntry[lúga] \dicTerm{lúg|a} \dicIPA{{l}{u}{\textlengthmark}{a}} \dicPos{f}[1] \dicFlx{(‑u, ‑ur)}[7] \dicDirectTranslationCS{poklop, příklop}
\dicEntry[lúi] \dicTerm{lú|i} \dicIPA{{l}{u}{\textlengthmark}{\textsci}} \dicPos{m}[1] \dicFlx{(‑a)}[3] \dicSynonym{þreyta\smash{\textsuperscript{1}}} \dicDirectTranslationCS{znavení, vyčerpanost, vysílení} \dicExampleIS{finna fyrir lúa eftir ferðina} \dicExampleCS{cítit vysílení po cestě}
\dicEntry[lúið] \dicTerm{lúið} \dicIPA{{l}{u}{\textlengthmark}{\textsci}{\texttheta}} \dicPos{v} \dicFlx{supin} \dicLink{lýja}
\dicEntry[lúinn] \dicTerm{lúinn} \dicIPA{{l}{u}{\textlengthmark}{\textsci}{\textsubring{n}}} \dicPos{adj}[6]\dicFlx{}[-2] \textbf{1.} \dicSynonym{þreyttur} \dicDirectTranslationCS{znavený, vyčerpaný, vysílený}  \textbf{2.} \dicSynonym{slitinn} \dicDirectTranslationCS{opotřebovaný, obnošený}
\dicEntry[lúka] \dicTerm{lúk|a} \dicIPA{{l}{u}{\textlengthmark}{\r{g}}{a}} \dicPos{f}[1] \dicFlx{(‑u, ‑ur)}[19] \textbf{1.} \dicSynonym{lófi} \dicDirectTranslationCS{dlaň};  \dicPhraseIS{vera með lífið í lúkunum} \dicLangCat{přen.} \dicDirectTranslationCS{mít srdce v~kalhotách}  \textbf{2.} \dicSynonym{handfylli} \dicDirectTranslationCS{hrst} \dicExampleIS{lúka af salti} \dicExampleCS{hrst soli}
\dicEntry[lúkar] \dicTerm{lúkar} \dicIPA{{l}{u}{\textlengthmark}{\r{g}}{a}{\textsubring{r}}} \dicPos{m}[4] \dicFlx{(‑s, ‑ar)}[15] \dicFieldCat{nám.} \dicSynonym{káeta} \dicDirectTranslationCS{důstojnická kabina\,/\addthin kajuta}
\dicEntry[lúlla] \dicTerm{lúll|a} \dicIPA{{l}{u}{l}{\textlengthmark}{a}} \dicPos{v}[1] \dicFlx{(‑aði)}[44] \dicLangCat{dět.} \dicSynonym{sofa} \dicDirectTranslationCS{spinkat, hajat} \dicExampleIS{Farðu nú að lúlla.} \dicExampleCS{Jdi spinkat.}
\dicEntry[lúmskur] \dicTerm{lúmskur} \dicIPA{{l}{u}{m}{s}{\r{g}}{\textscy}{\textsubring{r}}} \dicPos{adj}[1]\dicFlx{}[-1] \dicSynonym{undirförull} \dicDirectTranslationCS{podlý, zákeřný};  \dicPhraseIS{hafa lúmskt gaman af að (gera e‑ð)} \dicDirectTranslationCS{mít škodolibou radost z~(dělání (čeho))}
\dicEntry[lúpína] \dicTerm{lúpín|a} \dicIPA{{l}{u}{\textlengthmark}{p\smash{\textsuperscript{h}}}{i}{n}{a}} \dicPos{f}[1] \dicFlx{(‑u, ‑ur)}[7] \dicFieldCat{bot.} \dicDirectTranslationCS{lupina nutkajská, vlčí bob} \textit{(l.~{\textLA{Lupinus nootkatensis}})}  \dicsymPhoto\ 
\dicFigure{ds_image_lupina_0_1.jpg}{Lúpína}{Lúpína - Jutta234, GFDL}
\dicEntry[lúpulegur] \dicTerm{lúpu··legur} \dicIPA{{l}{u}{\textlengthmark}{\textsubring{b}}{\textscy}{l}{\textepsilon}{\textbabygamma}{\textscy}{\textsubring{r}}}\addthinS\dicPos{adj}[1]\dicFlx{}[-8] \dicSynonym{skömmustulegur} \dicDirectTranslationCS{zahanbený, (jsoucí) v~rozpacích}
\dicEntry[lúr] \dicTerm{lúr} \dicIPA{{l}{u}{\textlengthmark}{\textsubring{r}}} \dicPos{m}[4] \dicFlx{(‑s, ‑ar)}[14] \dicSynonym{blundur} \dicDirectTranslationCS{zdřímnutí, šlofík} \dicExampleIS{fá sér svolítinn lúr} \dicExampleCS{trochu si zdřímnout}
\dicEntry[lúra] \dicTerm{lúr|a} \dicIPA{{l}{u}{\textlengthmark}{r}{a}} \dicPos{v}[2] \dicFlx{(‑ði, ‑að\,/\addthin ‑t)}[184] \dicSynonym{sofa} \dicDirectTranslationCS{zdřímnout si, dřímat, podřimovat};  \dicIdiom{lúra}[á]{ \dicPhraseIS{lúra á e‑u}} \dicDirectTranslationCS{schovávat (co), ukrýt (si) (co)} \dicExampleIS{lúra á góðgæti} \dicExampleCS{schovávat dobroty}
\dicEntry[lús] \dicTerm{lús} \dicIPA{{l}{u}{\textlengthmark}{s}} \dicPos{f}[9] \dicFlx{(‑ar, lýs)}[6] \dicFieldCat{zool.} \dicDirectTranslationCS{veš} \textit{(l.~{\textLA{Anoplura}})};  \dicPhraseIS{læðast eins og lús með saum(i)} \dicFlx{refl} \dicLangCat{přen.} \dicDirectTranslationCS{loudat se jako hlemýžď};  \dicPhraseIS{lús kviknaði milli þeirra} \dicLangCat{přen.} \dicDirectTranslationCS{dostali se do křížku};  \dicPhraseIS{nú detta mér allar dauðar lýs úr höfði} \dicLangCat{přen.} \dicDirectTranslationCS{to jsem z~toho jelen}
\dicEntry[lúsarlegur] \dicTerm{lúsar··legur} \dicIPA{{l}{u}{\textlengthmark}{s}{a}{r}{l}{\textepsilon}{\textbabygamma}{\textscy}{\textsubring{r}}} \dicPos{adj}[1]\dicFlx{}[-8] \dicSynonym{smásálarlegur} \dicDirectTranslationCS{ubohý, tristní, mizerný (plat ap.)}
\dicEntry[lúta] \dicTerm{lúta} \dicsymFrequent\  \dicIPA{{l}{u}{\textlengthmark}{\textsubring{d}}{a}} \dicPos{v}[6] \dicFlx{(lýt, laut, lutum, lyti, lotið)}[106] \dicFlx{dat} \textbf{1.} \dicSynonym{beygja\smash{\textsuperscript{2}}} \dicDirectTranslationCS{sklonit, sklánět, sehnout, shýbat};  \dicPhraseIS{lúta höfði} \dicDirectTranslationCS{sklonit hlavu}  \textbf{2.} \dicDirectTranslationCS{sklonit se, podrobit se};  \dicPhraseIS{lúta e‑m} \dicDirectTranslationCS{sklonit se před (kým), podrobit se (komu)};  \dicIdiom{lúta}[að]{ \dicPhraseIS{lúta að e‑u}} \dicDirectTranslationCS{vztahovat se na (co), týkat se (čeho)};  \dicIdiom{lúta}{ \dicPhraseIS{lúta (svo) lágt}} \dicLangCat{přen.} \dicDirectTranslationCS{snížit se, klesnout} \dicIndirectTranslationCS{(propůjčit se k~něčemu špatnému)}; { \dicPhraseIS{lúta í lægra haldi fyrir e‑m}} \dicLangCat{přen.} \dicDirectTranslationCS{sklonit se před (kým)} \dicIndirectTranslationCS{(přijmout prohru)}
\dicEntry[lúterskur] \dicTerm{lúterskur} \dicIPA{{l}{u}{\textlengthmark}{t\smash{\textsuperscript{h}}}{e}{\textsubring{r}}{s}{\r{g}}{\textscy}{\textsubring{r}}} \dicPos{adj}[1]\dicFlx{}[-1] \dicFieldCat{náb.} \dicDirectTranslationCS{luteránský}
\dicEntry[lúterstrú] \dicTerm{lúters··trú} \dicIPA{{l}{u}{\textlengthmark}{t\smash{\textsuperscript{h}}}{\textepsilon}{\textsubring{r}}{s}{t\smash{\textsuperscript{h}}}{r}{u}} \dicPos{f}[4] \dicFlx{(‑ar)}[27] \dicFieldCat{náb.} \dicDirectTranslationCS{luteránství}
\dicEntry[lútur] \dicTerm{lút|ur} \dicIPA{{l}{u}{\textlengthmark}{\textsubring{d}}{\textscy}{\textsubring{r}}} \dicPos{m}[6] \dicFlx{(‑s\,/\addthin ‑ar, ‑ar)}[65] \dicFieldCat{chem.} \dicDirectTranslationCS{louh}
\dicEntry[Lúxemborg] \dicTerm{Lúxem··borg} \dicIPA{{l}{u}{x}{s}{\textepsilon}{m}{\textsubring{b}}{\textopeno}{r}{\r{g}}} \dicPos{f}[7] \dicFlx{(‑ar)}[4] \textbf{1.} \dicFieldCat{geog.} \dicDirectTranslationCS{Lucembursko}  \textbf{2.} \dicFieldCat{geog.} \dicDirectTranslationCS{Lucemburk} \dicIndirectTranslationCS{(hlavní město Lucemburska)}
\dicEntry[Lúxemborgari] \dicTerm{Lúxem··borg·ar|i} \dicIPA{{l}{u}{x}{s}{\textepsilon}{m}{\textsubring{b}}{\textopeno}{r}{\r{g}}{a}{r}{\textsci}} \dicPos{m}[1] \dicFlx{(‑a, ‑ar)}[13] \dicDirectTranslationCS{Lucemburčan(ka)}
\dicEntry[lúxemborgíska] \dicTerm{lúxem··borgísk|a} \dicIPA{{l}{u}{x}{s}{\textepsilon}{m}{\textsubring{b}}{\textopeno}{r}{\r{\textObardotlessj}}{i}{s}{\r{g}}{a}} \dicPos{f}[1] \dicFlx{(‑u)}[5] \dicDirectTranslationCS{lucemburština}
\dicEntry[lúxemborgskur] \dicTerm{lúxem··borgskur} \dicIPA{{l}{u}{x}{s}{\textepsilon}{m}{\textsubring{b}}{\textopeno}{r}{\r{g}}{s}{\r{g}}{\textscy}{\textsubring{r}}} \dicPos{adj}[1]\dicFlx{}[-6] \dicDirectTranslationCS{lucemburský}
\dicEntry[lúxus] \dicTerm{lúxus} \dicIPA{{l}{u}{x}{s}{\textscy}{s}} \dicPos{m}[4] \dicFlx{(‑s)}[21] \dicDirectTranslationCS{luxus, přepych}
\dicEntry[lyf] \dicTerm{lyf} \dicsymFrequent\  \dicIPA{{l}{\textsci}{\textlengthmark}{f}} \dicPos{n}[2] \dicFlx{(‑s, ‑)}[13] \dicDirectTranslationCS{lék} \dicExampleIS{lyf gegn lyfseðli} \dicExampleCS{lék na předpis};  \dicPhraseIS{lyf gegn e‑u} \dicDirectTranslationCS{lék proti (čemu), lék na (co)} \dicExampleIS{lyf gegn kvefi} \dicExampleCS{lék na rýmu};  \dicPhraseIS{lyf við e‑ð} \dicDirectTranslationCS{lék na (co)};  \dicPhraseIS{taka inn lyf} \dicDirectTranslationCS{brát léky};  \dicPhraseIS{verja á lyfjum} \dicDirectTranslationCS{být na prášcích}
\dicEntry[lyfjabúð] \dicTerm{lyfja··búð} \dicIPA{{l}{\textsci}{v}{j}{a}{\textsubring{b}}{u}{\texttheta}} \dicPos{f}[7] \dicFlx{(‑ar, ‑ir)}[1] \dicDirectTranslationCS{lékárna}
\dicEntry[lyfjafræði] \dicTerm{lyfja··fræð|i} \dicIPA{{l}{\textsci}{v}{j}{a}{f}{r}{a}{i}{ð}{\textsci}} \dicPos{f}[3] \dicFlx{(‑i)}[3] \textbf{1.} \dicDirectTranslationCS{lékárnictví, farmacie}  \textbf{2.} \dicDirectTranslationCS{farmakologie}
\dicEntry[lyfjafræðingur] \dicTerm{lyfja·fræð··ing|ur} \dicIPA{{l}{\textsci}{v}{j}{a}{f}{r}{a}{i}{ð}{i}{\ng}{\r{g}}{\textscy}{\textsubring{r}}} \dicPos{m}[6] \dicFlx{(‑s, ‑ar)}[8] \textbf{1.} \dicDirectTranslationCS{lékárník, lékárnice, farmaceut(ka)}  \textbf{2.} \dicDirectTranslationCS{farmakolog, farmakoložka}
\dicEntry[lyfjagras] \dicTerm{lyfja··|gras} \dicIPA{{l}{\textsci}{v}{j}{a}{\r{g}}{r}{a}{s}} \dicPos{n}[2] \dicFlx{(‑grass, ‑grös)}[8] \dicFieldCat{bot.} \dicDirectTranslationCS{tučnice obecná} \textit{(l.~{\textLA{Pinguicula vulgaris}})}  \dicsymPhoto\ 
\dicFigure{31254.jpg}{Lyfjagras}{Lyfjagras - Dvořák Václav, Biolib, Copyright/CC-BY-NC}
\dicEntry[lyfjakort] \dicTerm{lyfja··kort} \dicIPA{{l}{\textsci}{v}{j}{a}{k\smash{\textsuperscript{h}}}{\textopeno}{\textsubring{r}}{\textsubring{d}}} \dicPos{n}[2] \dicFlx{(‑s, ‑)}[5] \dicIndirectTranslationCS{karta pacienta s~informacemi o~užívaných lécích}
\dicEntry[lyfjakostnaður] \dicTerm{lyfja··kost·nað|ur} \dicIPA{{l}{\textsci}{v}{j}{a}{k\smash{\textsuperscript{h}}}{\textopeno}{s}{\textsubring{d}}{n}{a}{ð}{\textscy}{\textsubring{r}}} \dicPos{m}[10] \dicFlx{(‑ar)}[9] \dicDirectTranslationCS{náklady na léky}
\dicEntry[lyfjaskírteini] \dicTerm{lyfja··skír·teini} \dicIPA{{l}{\textsci}{v}{j}{a}{s}{\r{\textObardotlessj}}{i}{\textsubring{r}}{t\smash{\textsuperscript{h}}}{ei}{n}{\textsci}} \dicPos{n}[2] \dicFlx{(‑s, ‑)}[14] \dicIndirectTranslationCS{průkaz opravňující ke slevě na léky}
\dicEntry[lyfjaverð] \dicTerm{lyfja··verð} \dicIPA{{l}{\textsci}{v}{j}{a}{v}{\textepsilon}{r}{\texttheta}} \dicPos{n}[2] \dicFlx{(‑s)}[2] \dicDirectTranslationCS{cena léku}
\dicEntry[lyfjaverslun] \dicTerm{lyfja··versl|un} \dicIPA{{l}{\textsci}{v}{j}{a}{v}{\textepsilon}{\textsubring{r}}{s}{\textsubring{d}}{l}{\textscy}{\textsubring{n}}} \dicPos{f}[7] \dicFlx{(‑unar, ‑anir)}[8] \dicSynonym{lyfjabúð} \dicDirectTranslationCS{lékárna}
\dicEntry[lyflækningar] \dicTerm{lyf··lækn·ingar} \dicIPA{{l}{\textsci}{v}{l}{a}{i}{h}{\r{g}}{n}{i}{\ng}{\r{g}}{a}{\textsubring{r}}} \dicPos{f}[4] \dicFlx{pl}[6] \dicFieldCat{med.} \dicDirectTranslationCS{klinická medicína}
\dicEntry[lyfsali] \dicTerm{lyf··sal|i} \dicIPA{{l}{\textsci}{f}{s}{a}{l}{\textsci}} \dicPos{m}[1] \dicFlx{(‑a, ‑ar)}[8] \dicDirectTranslationCS{lékárník, lékárnice}
\dicEntry[lyfseðill] \dicTerm{lyf··seð|ill} \dicIPA{{l}{\textsci}{f}{s}{\textepsilon}{ð}{\textsci}{\textsubring{d}}{\textsubring{l}}} \dicPos{m}[6] \dicFlx{(‑ils, ‑lar)}[35] \dicDirectTranslationCS{(lékařský) předpis}
\dicEntry[lyfta] \dicTerm{lyft|a\smash{\textsuperscript{1}}} \dicIPA{{l}{\textsci}{f}{\textsubring{d}}{a}} \dicPos{f}[1] \dicFlx{(‑u, ‑ur)}[7] \textbf{1.} \dicDirectTranslationCS{výtah}  \textbf{2.} \dicDirectTranslationCS{zdviž}  \textbf{3.} \dicSynonym{skíðalyfta} \dicDirectTranslationCS{vlek (lyžařský ap.)}
\dicEntry[lyfta] \dicTerm{lyft|a\smash{\textsuperscript{2}}} \dicsymFrequent\  \dicIPA{{l}{\textsci}{f}{\textsubring{d}}{a}} \dicPos{v}[2] \dicFlx{(‑i, ‑)}[14] \dicFlx{dat} \dicSynonym{hefja} \dicDirectTranslationCS{zvednout, zvedat, zdvihat, nadzvednout, nadzdvihnout} \dicExampleIS{lyfta steini} \dicExampleCS{zvednout kámen};  \dicPhraseIS{lyfta brúnum} \dicDirectTranslationCS{zvednout obočí};  \dicIdiom{lyfta}[sér]{ \dicPhraseIS{lyfta sér}} \dicDirectTranslationCS{(na)kynout (těsto ap.)};  \dicIdiom{lyfta}[upp]{ \dicPhraseIS{lyfta sér upp}} \dicSynonym*{skemmta sér} \dicDirectTranslationCS{povyrazit si, odvázat se};  \dicIdiom{lyftast}{ \dicPhraseIS{lyftast}} \dicFlx{refl} \dicDirectTranslationCS{zvednout\,/\addthin zvedat se (letadlo ap.)}
\dicEntry[lyftari] \dicTerm{lyft··ar|i} \dicIPA{{l}{\textsci}{f}{\textsubring{d}}{a}{r}{\textsci}} \dicPos{m}[1] \dicFlx{(‑a, ‑ar)}[13] \dicFieldCat{techn.} \dicDirectTranslationCS{vysokozdvižný vozík}
\dicEntry[lyftiduft] \dicTerm{lyfti··duft} \dicIPA{{l}{\textsci}{f}{\textsubring{d}}{\textsci}{\textsubring{d}}{\textscy}{f}{\textsubring{d}}} \dicPos{n}[2] \dicFlx{(‑s)}[2] \dicDirectTranslationCS{prášek do pečiva}
\dicEntry[lyfting] \dicTerm{lyft··ing} \dicIPA{{l}{\textsci}{f}{\textsubring{d}}{i}{\ng}{\r{g}}} \dicPos{f}[4] \dicFlx{(‑ar, ‑ar)}[5] \textbf{1.} \dicSynonym*{það að lyfta} \dicDirectTranslationCS{zvednutí, zvedání, (vy)zdvižení, nadzvednutí}  \textbf{2.} \dicSynonym*{fjörgun} \dicDirectTranslationCS{pozvednutí, povznesení (na duchu ap.)} \dicExampleIS{lyfting andans} \dicExampleCS{povznesení na duchu}  \textbf{3.} \dicPhraseIS{lyftingar} \dicFlx{pl} \dicFieldCat{sport.} \dicDirectTranslationCS{vzpírání}
\begin{xtolerant}{}{1pt}
\dicEntry[lyftingamaður] \dicTerm{lyftinga··|maður}\addthinS\dicIPA{{l}{\textsci}{f}{\textsubring{d}}{i}{\ng}{\r{g}}{a}{m}{a}{ð}{\textscy}{\textsubring{r}}}\addthinS\dicPos{m}[13]\addthinS\dicFlx{(‑manns, ‑menn)}[2] \dicDirectTranslationCS{vzpě\-rač(ka)}
\end{xtolerant}
\dicEntry[lygamælir] \dicTerm{lyga··mæl|ir} \dicIPA{{l}{\textsci}{\textlengthmark}{\textbabygamma}{a}{m}{a}{i}{l}{\textsci}{\textsubring{r}}} \dicPos{m}[7] \dicFlx{(‑is, ‑ar)}[1] \dicDirectTranslationCS{detektor lži}
\dicEntry[lygar] \dicTerm{lygar} \dicIPA{{l}{\textsci}{\textlengthmark}{\textbabygamma}{a}{\textsubring{r}}} \dicPos{f} \dicFlx{pl nom} \dicLink{lygi\smash{\textsuperscript{1}}}
\dicEntry[lygari] \dicTerm{lyg··ar|i} \dicIPA{{l}{\textsci}{\textlengthmark}{\textbabygamma}{a}{r}{\textsci}} \dicPos{m}[1] \dicFlx{(‑a, ‑ar)}[13] \dicDirectTranslationCS{lhář(ka)}
\dicEntry[lygasaga] \dicTerm{lyga··|saga} \dicIPA{{l}{\textsci}{\textlengthmark}{\textbabygamma}{a}{s}{a}{\textbabygamma}{a}} \dicPos{f}[1] \dicFlx{(‑sögu, ‑sögur)}[14] \dicDirectTranslationCS{báchorka, pohádka}
\dicEntry[lygi] \dicTerm{lyg|i\smash{\textsuperscript{1}}} \dicsymFrequent\  \dicIPA{{l}{i}{j}{\textlengthmark}{\textsci}} \dicPos{f}[2] \dicFlx{(‑i, ‑ar)}[1] \dicDirectTranslationCS{lež, nepravda} \dicExampleIS{Bæjarstjórinn segir lygar.} \dicExampleCS{Starosta lže.};  \dicPhraseIS{fara með lygar} \dicDirectTranslationCS{lhát}
\dicEntry[lygi] \dicTerm{lygi\smash{\textsuperscript{2}}} \dicIPA{{l}{i}{j}{\textlengthmark}{\textsci}} \dicPos{v} \dicFlx{con pf sg 1 pers} \dicLink{ljúga}
\dicEntry[lygilegur] \textls[15]{\dicTerm{lygi··legur} \dicIPA{{l}{i}{j}{\textlengthmark}{\textsci}{l}{\textepsilon}{\textbabygamma}{\textscy}{\textsubring{r}}} \dicPos{adj}[1]\dicFlx{}[-8] \textbf{1.} \dicSynonym{ótrúlegur} \dicDirectTranslationCS{nevěrohodný, nepravděpodobný}}  \textbf{2.} \dicSynonym{ósannur} \dicDirectTranslationCS{nepravdivý, lživý}
\dicEntry[lyginn] \dicTerm{lyginn} \dicIPA{{l}{i}{j}{\textlengthmark}{\textsci}{\textsubring{n}}} \dicPos{adj}[6]\dicFlx{}[-2] \dicDirectTranslationCS{ulhaný, prolhaný}
\dicEntry[lygn] \dicTerm{lygn} \dicIPA{{l}{\textsci}{\r{g}}{\textsubring{n}}} \dicPos{adj}[5]\dicFlx{}[-5] \textbf{1.} \dicSynonym*{vindlaus} \dicDirectTranslationCS{bezvětrný, klidný} \dicIndirectTranslationCS{(o~počasí)}  \textbf{2.} \dicSynonym*{straumlaus} \dicDirectTranslationCS{nehybný, nezčeřený, klidný} \dicIndirectTranslationCS{(o~vodě)}
\dicEntry[lygna] \dicTerm{lygn|a\smash{\textsuperscript{1}}} \dicIPA{{l}{\textsci}{\r{g}}{n}{a}} \dicPos{f}[1] \dicFlx{(‑u, ‑ur)}[7] \textbf{1.} \dicSynonym*{kyrrt vatn} \dicDirectTranslationCS{klidná voda}  \textbf{2.} \dicSynonym*{kyrrt loft} \dicDirectTranslationCS{bezvětří}
\dicEntry[lygna] \dicTerm{lygn|a\smash{\textsuperscript{2}}} \dicIPA{{l}{\textsci}{\r{g}}{n}{a}} \dicPos{v}[2] \dicFlx{(‑di, ‑t)}[150] \dicFlx{dat} \textbf{1.} \dicPhraseIS{það lygnir} \dicFlx{impers} \dicSynonym*{kyrrast} \dicDirectTranslationCS{utišuje se (vítr), přestává foukat}  \textbf{2.} \dicPhraseIS{lygna aftur augunum} \dicDirectTranslationCS{zavřít\,/\addthin přivřít oči}
\dicEntry[lyk] \dicTerm{lyk} \dicIPA{{l}{\textsci}{\textlengthmark}{\r{g}}} \dicPos{v} \dicFlx{ind praes sg 1 pers} \dicLink{lykja}
\dicEntry[lyki] \dicTerm{lyki} \dicIPA{{l}{\textsci}{\textlengthmark}{\r{\textObardotlessj}}{\textsci}} \dicPos{v} \dicFlx{con pf sg 1 pers} \dicLink{ljúka}
\dicEntry[lykilaðstaða] \dicTerm{lykil··að·|staða} \dicIPA{{l}{\textsci}{\textlengthmark}{\r{\textObardotlessj}}{\textsci}{l}{a}{ð}{s}{\textsubring{d}}{a}{ð}{a}} \dicPos{f}[1] \dicFlx{(‑stöðu)}[2] \dicDirectTranslationCS{klíčové postavení}
\dicEntry[lykilatriði] \dicTerm{lykil··at·riði} \dicIPA{{l}{\textsci}{\textlengthmark}{\r{\textObardotlessj}}{\textsci}{l}{a}{\textsubring{d}}{r}{\textsci}{ð}{\textsci}} \dicPos{n}[2] \dicFlx{(‑s, ‑)}[14] \dicDirectTranslationCS{klíčový bod, klíčová otázka}
\dicEntry[lykill] \dicTerm{lyk|ill} \dicsymFrequent\  \dicIPA{{l}{\textsci}{\textlengthmark}{\r{\textObardotlessj}}{\textsci}{\textsubring{d}}{\textsubring{l}}} \dicPos{m}[6] \dicFlx{(‑ils, ‑lar)}[35] \textbf{1.} \dicDirectTranslationCS{klíč (od dveří ap.)};  \dicPhraseIS{lykill að e‑u} \dicDirectTranslationCS{klíč od (čeho)} \dicExampleIS{lykill að húsinu} \dicExampleCS{klíč od domu}  \textbf{2.} \dicFieldCat{techn.} \dicDirectTranslationCS{klíč} \dicIndirectTranslationCS{(ruční nástroj)}  \textbf{3.} \dicFieldCat{hud.} \dicDirectTranslationCS{klíč} \dicIndirectTranslationCS{(hudební značka)} \dicExampleIS{nótnalykill} \dicExampleCS{hudební klíč}  \textbf{4.} \dicSynonym{skýring} \dicDirectTranslationCS{klíč, nápověda} \dicIndirectTranslationCS{(pomůcka podstatná pro vysvětlení něčeho)} \dicExampleIS{orðalykill} \dicExampleCS{konkordance}
\dicEntry[lykilmaður] \dicTerm{lykil··|maður} \dicIPA{{l}{\textsci}{\textlengthmark}{\r{\textObardotlessj}}{\textsci}{l}{m}{a}{ð}{\textscy}{\textsubring{r}}} \dicPos{m}[13] \dicFlx{(‑manns, ‑menn)}[2] \dicDirectTranslationCS{klíčová postava}
\dicEntry[lykilorð] \dicTerm{lykil··orð} \dicIPA{{l}{\textsci}{\textlengthmark}{\r{\textObardotlessj}}{\textsci}{l}{\textopeno}{r}{\texttheta}} \dicPos{n}[2] \dicFlx{(‑s, ‑)}[5] \textbf{1.} \dicSynonym{aðalatriði} \dicDirectTranslationCS{klíčové slovo}  \textbf{2.} \dicFieldCat{poč.} \dicSynonym*{aðgangsorð} \dicDirectTranslationCS{(přístupové) heslo}
\dicEntry[lykja] \dicTerm{lykja} \dicIPA{{l}{\textsci}{\textlengthmark}{\r{\textObardotlessj}}{a}} \dicPos{v}[4] \dicFlx{(lyk, lukti, luktum, lykti, lukið)}[61] \dicDirectTranslationCS{obklopit, obklopovat, obstoupit, obstupovat} \dicExampleIS{Vatn lukti um hana.} \dicExampleCS{Voda ji obklopila.};  \dicIdiom{lykjast}{ \dicPhraseIS{lykjast}} \dicFlx{refl} \dicDirectTranslationCS{obklopit, obklopovat, obstoupit, obstupovat}
\dicEntry[lykkja] \dicTerm{lykkj|a} \dicIPA{{l}{\textsci}{h}{\r{\textObardotlessj}}{a}} \dicPos{f}[1] \dicFlx{(‑u, ‑ur)}[25] \textbf{1.} \dicSynonym{hanki} \dicDirectTranslationCS{smyčka, oko, klička (na provázku ap.)} \dicExampleIS{lykkja á bandi} \dicExampleCS{smyčka na provázku}  \textbf{2.} \dicDirectTranslationCS{oko} \dicIndirectTranslationCS{(v~pletení)}  \textbf{3.} \dicSynonym{getnaðarvörn} \dicDirectTranslationCS{spirála} \dicIndirectTranslationCS{(druh antikoncepce)}
\dicEntry[lykkjufall] \dicTerm{lykkju··|fall} \dicIPA{{l}{\textsci}{h}{\r{\textObardotlessj}}{\textscy}{f}{a}{\textsubring{d}}{\textsubring{l}}} \dicPos{n}[2] \dicFlx{(‑falls, ‑föll)}[8] \dicDirectTranslationCS{puštěné oko} \dicIndirectTranslationCS{(v~pletení)} \dicExampleIS{lykkjufall á sokkabuxurnar} \dicExampleCS{oko na punčoše}
\dicEntry[lyklaborð] \dicTerm{lykla··borð} \dicsymFrequent\  \dicIPA{{l}{\textsci}{h}{\r{g}}{l}{a}{\textsubring{b}}{\textopeno}{r}{\texttheta}} \dicPos{n}[2] \dicFlx{(‑s, ‑)}[5] \dicFieldCat{poč.} \dicDirectTranslationCS{klávesnice} \dicExampleIS{íslenskt lyklaborð} \dicExampleCS{islandská klávesnice}
\dicEntry[lyklakippa] \dicTerm{lykla··kipp|a} \dicIPA{{l}{\textsci}{h}{\r{g}}{l}{a}{c\smash{\textsuperscript{h}}}{\textsci}{h}{\textsubring{b}}{a}} \dicPos{f}[1] \dicFlx{(‑u, ‑ur)}[19] \dicDirectTranslationCS{svazek klíčů}
\dicEntry[lykt] \dicTerm{lykt\smash{\textsuperscript{1}}} \dicIPA{{l}{\textsci}{x}{\textsubring{d}}} \dicPos{f}[7] \dicFlx{(‑ar, ‑ir)}[1] \dicSynonym{endir} \dicDirectTranslationCS{konec, závěr};  \dicPhraseIS{að lyktum} \dicFlx{adv} \dicSynonym*{að lokum} \dicDirectTranslationCS{nakonec, na závěr};  \dicPhraseIS{leiða e‑ð til lykta} \dicSynonym*{ljúka e‑u} \dicDirectTranslationCS{dovést (co) do konce, dokončit (co)};  \dicPhraseIS{gera lyktir á e‑u} \dicSynonym*{ljúka við e‑ð} \dicDirectTranslationCS{dokončit\,/\addthin dokončovat (co)}
\dicEntry[lykt] \dicTerm{lykt\smash{\textsuperscript{2}}} \dicsymFrequent\  \dicIPA{{l}{\textsci}{x}{\textsubring{d}}} \dicPos{f}[7] \dicFlx{(‑ar, ‑ir)}[1] \dicSynonym{ilmur} \dicDirectTranslationCS{(zá)pach, vůně, aroma} \dicExampleIS{lykt af blómunum} \dicExampleCS{vůně květin}
\dicEntry[lykta] \dicTerm{lykt|a\smash{\textsuperscript{1}}} \dicIPA{{l}{\textsci}{x}{\textsubring{d}}{a}} \dicPos{v}[1] \dicFlx{(‑aði)}[5] \dicPhraseIS{e‑u lyktar} \dicFlx{impers} \dicDirectTranslationCS{(co) končí} \dicExampleIS{Leiknum lyktaði með jafntefli.} \dicExampleCS{Zápas skončil remízou.}
\dicEntry[lykta] \dicTerm{lykt|a\smash{\textsuperscript{2}}} \dicIPA{{l}{\textsci}{x}{\textsubring{d}}{a}} \dicPos{v}[1] \dicFlx{(‑aði)}[34] \textbf{1.} \dicSynonym{ilma} \dicDirectTranslationCS{páchnout, být cítit, vonět, vydávat vůni} \dicExampleIS{lykta vel} \dicExampleCS{vonět dobře};  \dicPhraseIS{lykta af e‑u} \dicDirectTranslationCS{páchnout (čím), být cítit (čím), vonět (čím)} \dicExampleIS{lykta af ilmvatni} \dicExampleCS{vonět parfémem}  \textbf{2.} \dicSynonym{þefa} \dicDirectTranslationCS{přivonět (si), přičichnout (si)};  \dicPhraseIS{lykta af e‑u} \dicSynonym{þefa} \dicDirectTranslationCS{přivonět (si) k~(čemu)} \dicExampleIS{lykta af blóminu} \dicExampleCS{přivonět si ke květině}
\dicEntry[lyktandi] \dicTerm{lykt··andi} \dicIPA{{l}{\textsci}{x}{\textsubring{d}}{a}{n}{\textsubring{d}}{\textsci}} \dicPos{adj}[13] \dicFlx{indecl}[1] \dicDirectTranslationCS{vonící, šířící vůni};  \dicPhraseIS{e‑að er illa lyktandi} \dicDirectTranslationCS{(co) smrdí\,/\addthin páchne}
\dicEntry[lykti] \dicTerm{lykti} \dicIPA{{l}{\textsci}{x}{\textsubring{d}}{\textsci}} \dicPos{v} \dicFlx{con pf sg 1 pers} \dicLink{lykja}
\dicEntry[lyktnæmur] \dicTerm{lykt··næmur} \dicIPA{{l}{\textsci}{x}{\textsubring{d}}{n}{a}{i}{m}{\textscy}{\textsubring{r}}} \dicPos{adj}[1]\dicFlx{}[-1] \dicDirectTranslationCS{mající dobrý čich}
\dicEntry[lymska] \dicTerm{lymsk|a} \dicIPA{{l}{\textsci}{m}{s}{\r{g}}{a}} \dicPos{f}[1] \dicFlx{(‑u)}[5] \dicSynonym*{lævísi} \dicDirectTranslationCS{záludnost, zákeřnost, proradnost}
\dicEntry[lymskur] \dicTerm{lymskur} \dicIPA{{l}{\textsci}{m}{s}{\r{g}}{\textscy}{\textsubring{r}}} \dicPos{adj}[1]\dicFlx{}[-1] \dicSynonym{lævís} \dicDirectTranslationCS{záludný, zákeřný, proradný}
\dicEntry[lynda] \dicTerm{lyn|da} \dicIPA{{l}{\textsci}{n}{\textsubring{d}}{a}} \dicPos{v}[2] \dicFlx{(‑ti, ‑t)}[44] \dicPhraseIS{e‑m lyndir við e‑n} \dicFlx{impers} \dicDirectTranslationCS{(kdo) si rozumí s~(kým), (kdo) má pochopení pro (koho)} \dicExampleIS{Mér lyndir ekki við hana.} \dicExampleCS{Nerozumím si s~ní.};  \dicPhraseIS{e‑jum lyndir} \dicFlx{impers} \dicDirectTranslationCS{(kdo) si rozumí, (kdo) má pochopení jeden pro druhého};  \dicPhraseIS{láta sér e‑ð lynda} \dicDirectTranslationCS{spokojit se s~(čím)}
\dicEntry[lyndi] \dicTerm{lyndi} \dicIPA{{l}{\textsci}{n}{\textsubring{d}}{\textsci}} \dicPos{n}[2] \dicFlx{(‑s)}[20] \dicSynonym{skap} \dicDirectTranslationCS{nálada, rozpoložení};  \dicPhraseIS{e‑að leikur í lyndi} \dicDirectTranslationCS{(co) klape, (co) hraje do karet};  \dicPhraseIS{e‑m leikur allt í lyndi} \dicFlx{impers} \dicDirectTranslationCS{(komu) vše vychází, (komu) vše hraje do noty}
\dicEntry[lyndiseinkunn] \dicTerm{lyndis··ein·kunn} \dicIPA{{l}{\textsci}{n}{\textsubring{d}}{\textsci}{s}{ei}{\r{\ng}}{\r{g}}{\textscy}{\textsubring{n}}} \dicPos{f}[7] \dicFlx{(‑ar, ‑ir)}[1] \dicSynonym*{skaphöfn} \dicDirectTranslationCS{charakter}
\dicEntry[lyng] \dicTerm{lyng} \dicIPA{{l}{i}{\ng}{\r{g}}} \dicPos{n}[2] \dicFlx{(‑s)}[2] \dicFieldCat{bot.} \dicDirectTranslationCS{vřesovec} \textit{(l.~{\textLA{Erica}})}  \dicsymPhoto\ 
\dicFigure{74636.jpg}{Lyng}{Lyng - Kesl Michael, Biolib, Copyright/CC-BY-NC}
\dicEntry[lyst] \dicTerm{lyst} \dicIPA{{l}{\textsci}{s}{\textsubring{d}}} \dicPos{f}[7] \dicFlx{(‑ar, ‑ir)}[1] \textbf{1.} \dicSynonym{matarlyst} \dicDirectTranslationCS{chuť (k~jídlu), apetit};  \dicPhraseIS{hafa lyst á e‑u} \dicDirectTranslationCS{mít chuť na (co)} \dicExampleIS{hafa góða lyst á matnum} \dicExampleCS{mít dobrý apetit}  \textbf{2.} \dicSynonym{löngun} \dicDirectTranslationCS{chuť, přání};  \dicPhraseIS{af hjartans lyst} \dicFlx{adv} \dicDirectTranslationCS{z~celého srdce}
\dicEntry[lystarlaus] \dicTerm{lystar··laus} \dicIPA{{l}{\textsci}{s}{\textsubring{d}}{a}{r}{l}{\oe i}{s}} \dicPos{adj}[5]\dicFlx{}[-1] \dicDirectTranslationCS{(jsoucí) bez apetitu, nemající chuť (na jídlo)}
\dicEntry[lystarleysi] \dicTerm{lystar··leysi} \dicIPA{{l}{\textsci}{s}{\textsubring{d}}{a}{r}{l}{ei}{s}{\textsci}} \dicPos{n}[2] \dicFlx{(‑s)}[20] \dicDirectTranslationCS{nechutenství}
\dicEntry[lystarstol] \dicTerm{lystar··stol} \dicIPA{{l}{\textsci}{s}{\textsubring{d}}{a}{\textsubring{r}}{s}{\textsubring{d}}{\textopeno}{\textsubring{l}}} \dicPos{n}[2] \dicFlx{(‑s)}[2] \dicFieldCat{med.} \dicDirectTranslationCS{(mentální) anorexie}
\dicEntry[lysti] \dicTerm{lysti} \dicIPA{{l}{\textsci}{s}{\textsubring{d}}{\textsci}} \dicPos{v} \dicFlx{con pf sg 1 pers} \dicLink{ljósta}
\dicEntry[lystibátur] \dicTerm{lysti··bát|ur} \dicIPA{{l}{\textsci}{s}{\textsubring{d}}{\textsci}{\textsubring{b}}{au}{\textsubring{d}}{\textscy}{\textsubring{r}}} \dicPos{m}[6] \dicFlx{(‑s, ‑ar)}[22] \dicFieldCat{nám.} \dicDirectTranslationCS{jachta}
\dicEntry[lystigarður] \dicTerm{lysti··garð|ur} \dicIPA{{l}{\textsci}{s}{\textsubring{d}}{\textsci}{\r{g}}{a}{r}{ð}{\textscy}{\textsubring{r}}} \dicPos{m}[6] \dicFlx{(‑s, ‑ar)}[5] \dicSynonym{skemmtigarður} \dicDirectTranslationCS{(zábavní) park}
\dicEntry[lystiskip] \dicTerm{lysti··skip} \dicIPA{{l}{\textsci}{s}{\textsubring{d}}{\textsci}{s}{\r{\textObardotlessj}}{\textsci}{\textsubring{b}}} \dicPos{n}[2] \dicFlx{(‑s, ‑)}[5] \dicDirectTranslationCS{jachta, výletní loď}
\dicEntry[lystugur] \dicTerm{lyst··ugur} \dicIPA{{l}{\textsci}{s}{\textsubring{d}}{\textscy}{\textbabygamma}{\textscy}{\textsubring{r}}} \dicPos{adj}[1]\dicFlx{}[-8] \textbf{1.} \dicSynonym{sólginn} \dicDirectTranslationCS{mající dobrý apetit} \dicExampleIS{vera lystugur á matinn} \dicExampleCS{mít dobrý apetit}  \textbf{2.} \dicSynonym*{sem vekur lyst} \dicDirectTranslationCS{chutný, vábný} \dicIndirectTranslationCS{(o~jídle)}
\dicEntry[lyti] \dicTerm{lyti} \dicIPA{{l}{\textsci}{\textlengthmark}{\textsubring{d}}{\textsci}} \dicPos{v} \dicFlx{con pf sg 1 pers} \dicLink{lúta}
\dicEntry[lýðfrjáls] \dicTerm{lýð··frjáls} \dicIPA{{l}{i}{ð}{f}{r}{j}{au}{l}{s}} \dicPos{adj}[5]\dicFlx{}[-1] \dicDirectTranslationCS{demokratický, svobodný} \dicExampleIS{lýðfrjálst land} \dicExampleCS{demokratická země}
\dicEntry[lýðháskóli] \dicTerm{lýð··há·skól|i} \dicIPA{{l}{i}{\textlengthmark}{\texttheta}{h}{au}{s}{\r{g}}{ou}{l}{\textsci}} \dicPos{m}[1] \dicFlx{(‑a, ‑ar)}[1] \dicFieldCat{škol.} \dicDirectTranslationCS{veřejná vysoká škola}
\dicEntry[lýðhylli] \dicTerm{lýð··hyll|i} \dicIPA{{l}{i}{\textlengthmark}{\texttheta}{h}{\textsci}{\textsubring{d}}{l}{\textsci}} \dicPos{f}[3] \dicFlx{(‑i)}[3] \dicDirectTranslationCS{přízeň lidu, popularita}
\dicEntry[lýði] \dicTerm{lýði} \dicIPA{{l}{i}{\textlengthmark}{ð}{\textsci}} \dicPos{v} \dicFlx{con pf sg 1 pers} \dicLink{lýja}
\dicEntry[lýðréttindi] \dicTerm{lýð··rétt·indi} \dicIPA{{l}{i}{ð}{r}{j}{\textepsilon}{h}{\textsubring{d}}{\textsci}{n}{\textsubring{d}}{\textsci}} \dicPos{n}[2] \dicFlx{pl}[19] \dicDirectTranslationCS{občanská práva}
\dicEntry[lýðræði] \dicTerm{lýð··ræði} \dicIPA{{l}{i}{ð}{r}{a}{i}{ð}{\textsci}} \dicPos{n}[2] \dicFlx{(‑s)}[20] \dicDirectTranslationCS{demokracie}
\dicEntry[lýðræðislegur] \dicTerm{lýð·ræðis··legur} \dicIPA{{l}{i}{ð}{r}{a}{i}{ð}{\textsci}{s}{l}{\textepsilon}{\textbabygamma}{\textscy}{\textsubring{r}}} \dicPos{adj}[1]\dicFlx{}[-8] \dicDirectTranslationCS{demokratický}
\dicEntry[lýðræðissinnaður] \dicTerm{lýð·ræðis··sinn·|aður} \dicIPA{{l}\-{i}\-{ð}\-{r}\-{a}\-{i}\-{ð}\-{\textsci}\-{s}\-{\textsci}\-{n}\-{a}\-{ð}\-{\textscy}\-{\textsubring{r}}\-} \dicPos{adj}[3] \dicFlx{(f ‑uð)}[3] \dicDirectTranslationCS{(pro)demokratický}
\dicEntry[lýðræðissinni] \dicTerm{lýð·ræðis··sinn|i} \dicIPA{{l}{i}{ð}{r}{a}{i}{ð}{\textsci}{s}{\textsci}{n}{\textsci}} \dicPos{m}[1] \dicFlx{(‑a, ‑ar)}[1] \dicDirectTranslationCS{demokrat(ka)}
\dicEntry[lýðskrum] \dicTerm{lýð··skrum} \dicIPA{{l}{i}{ð}{s}{\r{g}}{r}{\textscy}{\textsubring{m}}} \dicPos{n}[2] \dicFlx{(‑s)}[2] \dicDirectTranslationCS{demagogie, populismus}
\dicEntry[lýðskrumari] \dicTerm{lýð··skrum·ar|i} \dicIPA{{l}{i}{ð}{s}{\r{g}}{r}{\textscy}{m}{a}{r}{\textsci}} \dicPos{m}[1] \dicFlx{(‑a, ‑ar)}[13] \dicDirectTranslationCS{demagog, demagožka, populista, populistka}
\dicEntry[lýður] \dicTerm{lýð|ur} \dicIPA{{l}{i}{\textlengthmark}{ð}{\textscy}{\textsubring{r}}} \dicPos{m}[9] \dicFlx{(‑s, ‑ir)}[8] \textbf{1.} \dicSynonym{almenningur} \dicDirectTranslationCS{lid(é), občané}  \textbf{2.} \dicSynonym{skríll} \dicDirectTranslationCS{lůza, holota};  \dicPhraseIS{e‑að er enn við lýði} \dicDirectTranslationCS{(co) stále existuje, (co) se stále vyskytuje (zvyk ap.)}
\dicEntry[lýðveldi] \dicTerm{lýð··veldi} \dicIPA{{l}{i}{ð}{v}{\textepsilon}{l}{\textsubring{d}}{\textsci}} \dicPos{n}[2] \dicFlx{(‑s, ‑)}[14] \dicDirectTranslationCS{republika}
\dicEntry[lýg] \dicTerm{lýg} \dicIPA{{l}{i}{\textlengthmark}{x}} \dicPos{v} \dicFlx{ind praes sg 1 pers} \dicLink{ljúga}
\dicEntry[lýi] \dicTerm{lýi} \dicIPA{{l}{i}{j}{\textsci}} \dicPos{v} \dicFlx{ind praes sg 1 pers} \dicLink{lýja}
\dicEntry[lýja] \dicTerm{lýja} \dicIPA{{l}{i}{j}{\textlengthmark}{a}} \dicPos{v}[4] \dicFlx{(lýi, lúði, lúðum, lýði, lúið\,/\addthin lúð)}[54] \dicFlx{acc} \textbf{1.} \dicDirectTranslationCS{kovat, kout} \dicExampleIS{lýja járn} \dicExampleCS{kovat železo}  \textbf{2.} \dicSynonym{þreyta\smash{\textsuperscript{2}}} \dicDirectTranslationCS{unavit, utahat} \dicExampleIS{Þetta lýir mann.} \dicExampleCS{To člověka unaví.};  \dicIdiom{lýjast}{ \dicPhraseIS{lýjast}} \dicFlx{refl} \dicSynonym*{þreytast} \dicDirectTranslationCS{vyčerpat se, unavit se} \dicExampleIS{Strákurinn er farinn að lýjast.} \dicExampleCS{Kluk se začíná unavovat.}
\dicEntry[lýjandi] \dicTerm{lýj··andi} \dicIPA{{l}{i}{j}{\textlengthmark}{a}{n}{\textsubring{d}}{\textsci}} \dicPos{adj}[13] \dicFlx{indecl}[1] \dicDirectTranslationCS{vyčerpávající, vysilující} \dicExampleIS{lýjandi starf} \dicExampleCS{vyčerpávající práce}
\dicEntry[lýk] \dicTerm{lýk} \dicIPA{{l}{i}{\textlengthmark}{\r{g}}} \dicPos{v} \dicFlx{ind praes sg 1 pers} \dicLink{ljúka}
\dicEntry[lýrik] \dicTerm{lýrik} \dicIPA{{l}{i}{\textlengthmark}{r}{\textsci}{\r{g}}} \dicPos{f}[10] \dicFlx{(‑ur)}[2] \dicFieldCat{lit.} \dicSynonym{ljóðræna} \dicDirectTranslationCS{lyrika}
\dicEntry[lýrískur] \dicTerm{lýrískur} \dicIPA{{l}{i}{\textlengthmark}{r}{i}{s}{\r{g}}{\textscy}{\textsubring{r}}} \dicPos{adj}[1]\dicFlx{}[-1] \dicFieldCat{lit.} \dicSynonym{ljóðrænn} \dicDirectTranslationCS{lyrický}
\dicEntry[lýs] \dicTerm{lýs} \dicIPA{{l}{i}{\textlengthmark}{s}} \dicPos{f} \dicFlx{pl nom} \dicLink{lús}
\dicEntry[lýsa] \dicTerm{lýs|a\smash{\textsuperscript{1}}} \dicIPA{{l}{i}{\textlengthmark}{s}{a}} \dicPos{f}[1] \dicFlx{(‑u, ‑ur)}[7] \textbf{1.} \dicSynonym{ljómi} \dicDirectTranslationCS{světlo, paprsek} \dicExampleIS{fyrsta lýsa morgunsins} \dicExampleCS{první ranní paprsek}  \textbf{2.} \dicFieldCat{zool.} \dicDirectTranslationCS{treska bezvousá} \textit{(l.~{\textLA{Merlangius merlangus}})}  \dicsymPhoto\ 
\dicFigure{ds_image_lysa_1_2.jpg}{Lýsa}{Lýsa - Georges Jansoone (JoJan), CC BY 3.0}
\dicEntry[lýsa] \dicTerm{lýs|a\smash{\textsuperscript{2}}} \dicsymFrequent\  \dicIPA{{l}{i}{\textlengthmark}{s}{a}} \dicPos{v}[2] \dicFlx{(‑ti, ‑t)}[67] \dicFlx{dat} \textbf{1.} \dicSynonym*{bera birtu} \dicDirectTranslationCS{osvětlit, osvětlovat, osvítit, ozářit, ozařovat} \dicExampleIS{Tunglið lýsir mér.} \dicExampleCS{Měsíc mě osvětluje.};  \dicPhraseIS{það lýsir (af degi)} \dicFlx{impers} \dicSynonym*{það birtir} \dicDirectTranslationCS{rozednívá se}  \textbf{2.} \dicSynonym*{lýsa mynd} \dicDirectTranslationCS{vyvolat, vyvolávat (film ap.)} \dicExampleIS{lýsa myndina} \dicExampleCS{vyvolat film}  \textbf{3.} \dicSynonym{sýna} \dicDirectTranslationCS{projevit, projevovat, vyjádřit, vyjadřovat}  \textbf{4.} \dicSynonym*{gefa lýsingu} \dicDirectTranslationCS{popsat, popisovat, (vy)líčit, vykreslit, vykreslovat} \dicExampleIS{lýsa reynslu sinni} \dicExampleCS{popsat svou zkušenost}  \textbf{5.} \dicSynonym{tilkynna} \dicDirectTranslationCS{prohlásit, prohlašovat, deklarovat};  \dicIdiom{lýsa}[eftir]{ \dicPhraseIS{lýsa eftir e‑m}} \dicDirectTranslationCS{vyhlásit pátrání po (kom)};  \dicIdiom{lýsa}[upp]{ \dicPhraseIS{lýsa upp e‑ð}} \dicDirectTranslationCS{osvětlit (co) (místnost ap.)};  \dicIdiom{lýsa}[yfir]{ \dicPhraseIS{lýsa yfir e‑u}} \dicDirectTranslationCS{ohlásit (co), oznámit (co)}
\dicEntry[lýsi] \dicTerm{lýsi} \dicIPA{{l}{i}{\textlengthmark}{s}{\textsci}} \dicPos{n}[2] \dicFlx{(‑s, ‑)}[14] \dicDirectTranslationCS{rybí tuk} \dicIndirectTranslationCS{(medicína)}
\dicEntry[lýsing] \dicTerm{lýs··ing} \dicsymFrequent\  \dicIPA{{l}{i}{\textlengthmark}{s}{i}{\ng}{\r{g}}} \dicPos{f}[4] \dicFlx{(‑ar, ‑ar)}[5] \textbf{1.} \dicDirectTranslationCS{osvětlení, nasvícení, iluminace} \dicExampleIS{lýsing í eldhúsinu} \dicExampleCS{osvětlení v~kuchyni}  \textbf{2.} \dicSynonym{frásögn} \dicDirectTranslationCS{popis, vylíčení, vykreslení} \dicExampleIS{lýsing á ferðinni} \dicExampleCS{vylíčení cesty}  \textbf{3.} \dicDirectTranslationCS{osvit, expozice}  \textbf{4.} \dicSynonym{skreyting} \dicDirectTranslationCS{ilustrace, ilustrování (rukopisu, knihy ap.)}
\dicEntry[lýsingarháttur] \dicTerm{lýsingar··|háttur} \dicIPA{{l}{i}{\textlengthmark}{s}{i}{\ng}{\r{g}}{a}{\textsubring{r}}{h}{au}{h}{\textsubring{d}}{\textscy}{\textsubring{r}}} \dicPos{m}[12] \dicFlx{(‑háttar, ‑hættir)}[7] \dicFieldCat{jaz.} \dicDirectTranslationCS{příčestí};  \dicPhraseIS{lýsingarháttur nútíðar} \dicFieldCat{jaz.} \dicDirectTranslationCS{příčestí přítomné};  \dicPhraseIS{lýsingarháttur þátíðar} \dicFieldCat{jaz.} \dicDirectTranslationCS{příčestí minulé}
\dicEntry[lýsingarorð] \dicTerm{lýsingar··orð} \dicIPA{{l}{i}{\textlengthmark}{s}{i}{\ng}{\r{g}}{a}{r}{\textopeno}{r}{\texttheta}} \dicPos{n}[2] \dicFlx{(‑s, ‑)}[5] \dicFieldCat{jaz.} \dicDirectTranslationCS{přídavné jméno, adjektivum} \textit{(l.~{\textLA{nomen adiectivum}})}
\dicEntry[lýst] \dicTerm{lýst} \dicIPA{{l}{i}{s}{\textsubring{d}}} \dicPos{v} \dicFlx{ind praes sg 1 pers} \dicLink{ljósta}
\dicEntry[lýt] \dicTerm{lýt} \dicIPA{{l}{i}{\textlengthmark}{\textsubring{d}}} \dicPos{v} \dicFlx{ind praes sg 1 pers} \dicLink{lúta}
\dicEntry[lýtalaus] \textls[15]{\dicTerm{lýta··laus} \dicIPA{{l}{i}{\textlengthmark}{\textsubring{d}}{a}{l}{\oe i}{s}} \dicPos{adj}[5]\dicFlx{}[-1] \dicSynonym{gallalaus} \dicDirectTranslationCS{bezchybný, (jsoucí) bez poskvrnky}}
\dicEntry[lýtalækningar] \dicTerm{lýta··lækn·ingar} \dicIPA{{l}{i}{\textlengthmark}{\textsubring{d}}{a}{l}{a}{i}{h}{\r{g}}{n}{i}{\ng}{\r{g}}{a}{\textsubring{r}}} \dicPos{f}[4] \dicFlx{pl}[6] \dicFieldCat{med.} \dicDirectTranslationCS{plastická chirurgie}
\dicEntry[lýtalæknir] \dicTerm{lýta··lækn|ir} \dicIPA{{l}{i}{\textlengthmark}{\textsubring{d}}{a}{l}{a}{i}{h}{\r{g}}{n}{\textsci}{\textsubring{r}}} \dicPos{m}[7] \dicFlx{(‑is, ‑ar)}[1] \dicDirectTranslationCS{plastický chirurg, plastická chirurgyně\,/\addthin chiruržka}
\dicEntry[lýti] \dicTerm{lýti} \dicIPA{{l}{i}{\textlengthmark}{\textsubring{d}}{\textsci}} \dicPos{n}[2] \dicFlx{(‑s, ‑)}[14] \textbf{1.} \dicSynonym{ljóður} \dicDirectTranslationCS{poskvrna, kaz, vada} \dicExampleIS{lýti á bókinni} \dicExampleCS{vada na knížce}  \textbf{2.} \dicSynonym{ljótleiki} \dicDirectTranslationCS{šerednost, ošklivost (vzhled člověka ap.)}
\dicEntry[læða] \dicTerm{læð|a\smash{\textsuperscript{1}}} \dicIPA{{l}{a}{i}{\textlengthmark}{ð}{a}} \dicPos{f}[1] \dicFlx{(‑u, ‑ur)}[19] \dicDirectTranslationCS{kočka} \dicIndirectTranslationCS{(samice)}
\dicEntry[læða] \dicTerm{læ|ða\smash{\textsuperscript{2}}} \dicsymFrequent\  \dicIPA{{l}{a}{i}{\textlengthmark}{ð}{a}} \dicPos{v}[2] \dicFlx{(‑ddi, ‑tt)}[172] \dicFlx{dat} \dicPhraseIS{læða e‑u að e‑m} \dicDirectTranslationCS{propašovat (komu co), vstrčit (komu co) skrytě} \dicExampleIS{læða þessari frétt að henni} \dicExampleCS{nenápadně jí sdělit zprávu};  \dicIdiom{læðast}{ \dicPhraseIS{læðast}} \dicFlx{refl} \dicDirectTranslationCS{(v)plížit se, plazit se, krást se} \dicExampleIS{læðast burt} \dicExampleCS{odplížit se pryč}
\dicEntry[læðupokast] \dicTerm{læðu··pok|ast} \dicIPA{{l}{a}{i}{\textlengthmark}{ð}{\textscy}{p\smash{\textsuperscript{h}}}{\textopeno}{\r{g}}{a}{s}{\textsubring{d}}} \dicPos{v}[1] \dicFlx{(‑aðist)}[91] \dicFlx{refl} \dicDirectTranslationCS{(při)krást se, plížit se}
\dicEntry[lægð] \dicTerm{lægð} \dicIPA{{l}{a}{i}{\textbabygamma}{\texttheta}} \dicPos{f}[7] \dicFlx{(‑ar, ‑ir)}[1] \textbf{1.} \dicSynonym{dæld} \dicDirectTranslationCS{prohlubeň, proláklina} \dicExampleIS{Það myndast djúp lægð í götuna.} \dicExampleCS{Na ulici se tvoří hluboká prohlubeň.}  \textbf{2.} \dicFieldCat{meteo.} \dicSynonym{lágþrýstisvæði} \dicDirectTranslationCS{tlaková níže} \dicAntonym{hægð}  \textbf{3.} \dicSynonym{deyfð} \dicDirectTranslationCS{útlum};  \dicPhraseIS{vera í lægð} \dicDirectTranslationCS{být v~útlumu}
\dicEntry[lægi] \dicTerm{lægi\smash{\textsuperscript{1}}} \dicIPA{{l}{a}{i}{j}{\textlengthmark}{\textsci}} \dicPos{n}[2] \dicFlx{(‑s, ‑)}[16] \dicSynonym{skipalægi} \dicDirectTranslationCS{kotviště}
\dicEntry[lægi] \dicTerm{lægi\smash{\textsuperscript{2}}} \dicIPA{{l}{a}{i}{j}{\textlengthmark}{\textsci}} \dicPos{v} \dicFlx{con pf sg 1 pers} \dicLink{liggja}
\dicEntry[lægja] \dicTerm{læg|ja} \dicIPA{{l}{a}{i}{j}{\textlengthmark}{a}} \dicPos{v}[2] \dicFlx{(‑ði, ‑t)}[90] \dicFlx{acc} \dicSynonym{auðmýkja} \dicDirectTranslationCS{zklidnit, zklidňovat, zmírnit, zmírňovat} \dicExampleIS{lægja æsing} \dicExampleCS{zklidnit vztek};  \dicPhraseIS{lægja rostann í e‑m} \dicDirectTranslationCS{zklidnit (koho)};  \dicPhraseIS{e‑ð lægir} \dicFlx{impers} {\textbf{a.}} \dicDirectTranslationCS{(co) klesá, (co) zapadá (slunce ap.)} \dicExampleIS{Sólina lægir.} \dicExampleCS{Slunce klesá.};  {\textbf{b.}} \dicDirectTranslationCS{(co) ustává (o~větru ap.), (co) se uklidňuje} \dicExampleIS{Veðrið lægir.} \dicExampleCS{Počasí se uklidňuje.}
\dicEntry[lægra] \dicTerm{lægra} \dicIPA{{l}{a}{i}{\textbabygamma}{r}{a}} \dicPos{adv} \dicFlx{comp} \dicLink{lágt}
\dicEntry[lægri] \dicTerm{lægri} \dicIPA{{l}{a}{i}{\textbabygamma}{r}{\textsci}} \dicPos{adj} \dicFlx{comp m} \dicLink{lágur}
\dicEntry[lægst] \dicTerm{lægst} \dicIPA{{l}{a}{i}{x}{s}{\textsubring{d}}} \dicPos{adv} \dicFlx{sup} \dicLink{lágt}
\dicEntry[lægstur] \dicTerm{lægstur} \dicIPA{{l}{a}{i}{x}{s}{\textsubring{d}}{\textscy}{\textsubring{r}}} \dicPos{adj} \dicFlx{m sg nom sup} \dicLink{lágur}
\dicEntry[læki] \dicTerm{læki} \dicIPA{{l}{a}{i}{\textlengthmark}{\r{\textObardotlessj}}{\textsci}} \dicPos{v} \dicFlx{con pf sg 1 pers} \dicLink{leka}
\dicEntry[lækka] \dicTerm{lækk|a} \dicsymFrequent\  \dicIPA{{l}{a}{i}{h}{\r{g}}{a}} \dicPos{v}[1] \dicFlx{(‑aði)}[48] \dicFlx{acc} \textbf{1.} \dicSynonym{minnka} \dicDirectTranslationCS{snížit, snižovat, (z)redukovat} \dicExampleIS{lækka skatta} \dicExampleCS{snížit daně};  \dicPhraseIS{lækka í e‑u} \dicDirectTranslationCS{zeslabit (co), ztišit (co), ztlumit (co), zmírnit (co)} \dicExampleIS{lækka í útvarpinu} \dicExampleCS{ztišit rádio};  \dicPhraseIS{lækka röddina\,/\addthin róminn} \dicDirectTranslationCS{snížit hlas}  \textbf{2.} \dicDirectTranslationCS{klesnout, klesat, sestoupit, sestupovat, snášet se} \dicExampleIS{lækka flugið} \dicExampleCS{klesat s~letadlem};  \dicPhraseIS{e‑að lækkar} \dicDirectTranslationCS{(co) klesá, (co) ustupuje (voda ap.)};  \dicPhraseIS{það lækkar í e‑u} \dicFlx{impers} \dicDirectTranslationCS{(co) klesá, (co) opadává} \dicExampleIS{Það lækkar í ánni.} \dicExampleCS{Voda v~řece opadává.}
\dicEntry[lækkun] \dicTerm{lækk|un} \dicIPA{{l}{a}{i}{h}{\r{g}}{\textscy}{\textsubring{n}}} \dicPos{f}[7] \dicFlx{(‑unar)}[9] \textbf{1.} \dicSynonym*{það að lækka} \dicDirectTranslationCS{snížení, snižování, redukce, redukování} \dicExampleIS{lækkun skatta} \dicExampleCS{snížení daní}  \textbf{2.} \dicSynonym*{það að e‑að lækkar} \dicDirectTranslationCS{klesání, opadání, ubývání} \dicExampleIS{lækkun á yfirborði sjávar} \dicExampleCS{klesání hladiny moře}
\dicEntry[lækna] \dicTerm{lækn|a} \dicsymFrequent\  \dicIPA{{l}{a}{i}{h}{\r{g}}{n}{a}} \dicPos{v}[1] \dicFlx{(‑aði)}[1] \dicFlx{acc} \dicDirectTranslationCS{(vy)léčit};  \dicPhraseIS{lækna e‑n af e‑u} \dicDirectTranslationCS{léčit (koho) z~(čeho)} \dicExampleIS{lækna e‑n af berklum} \dicExampleCS{léčit (koho) z~tuberkulózy}
\dicEntry[læknadeild] \dicTerm{lækna··deild} \dicIPA{{l}{a}{i}{h}{\r{g}}{n}{a}{\textsubring{d}}{ei}{l}{\textsubring{d}}} \dicPos{f}[7] \dicFlx{(‑ar, ‑ir)}[1] \dicFieldCat{škol.} \dicDirectTranslationCS{lékařská fakulta}
\dicEntry[læknandi] \dicTerm{lækn··andi} \dicIPA{{l}{a}{i}{h}{\r{g}}{n}{a}{n}{\textsubring{d}}{\textsci}} \dicPos{adj}[13] \dicFlx{indecl}[1] \dicDirectTranslationCS{léčivý, léčebný}
\dicEntry[læknanemi] \dicTerm{lækna··nem|i} \dicIPA{{l}{a}{i}{h}{\r{g}}{n}{a}{n}{\textepsilon}{m}{\textsci}} \dicPos{m}[1] \dicFlx{(‑a, ‑ar)}[1] \dicDirectTranslationCS{student(ka) medicíny\,/\addthin lékařství}
\dicEntry[læknanlegur] \dicTerm{læknan··legur} \dicIPA{{l}{a}{i}{h}{\r{g}}{n}{a}{n}{l}{\textepsilon}{\textbabygamma}{\textscy}{\textsubring{r}}} \dicPos{adj}[1]\dicFlx{}[-8] \dicDirectTranslationCS{(vy)léčitelný}
\dicEntry[læknastofa] \dicTerm{lækna··stof|a} \dicIPA{{l}{a}{i}{h}{\r{g}}{n}{a}{s}{\textsubring{d}}{\textopeno}{v}{a}} \dicPos{f}[1] \dicFlx{(‑u, ‑ur)}[7] \dicDirectTranslationCS{ordinace}
\dicEntry[lækning] \dicTerm{lækn··ing} \dicsymFrequent\  \dicIPA{{l}{a}{i}{h}{\r{g}}{n}{i}{\ng}{\r{g}}} \dicPos{f}[4] \dicFlx{(‑ar, ‑ar)}[5] \dicDirectTranslationCS{léčení, léčba};  \dicPhraseIS{lækning við e‑u} \dicDirectTranslationCS{léčba (čeho)} \dicExampleIS{lækning við sjúkdómi} \dicExampleCS{léčba nemoci}
\dicEntry[lækningastofa] \dicTerm{lækninga··stof|a} \dicIPA{{l}{a}{i}{h}{\r{g}}{n}{i}{\ng}{\r{g}}{a}{s}{\textsubring{d}}{\textopeno}{v}{a}} \dicPos{f}[1] \dicFlx{(‑u, ‑ur)}[7] \dicDirectTranslationCS{ordinace}
\dicEntry[læknir] \dicTerm{lækn|ir} \dicsymFrequent\  \dicIPA{{l}{a}{i}{h}{\r{g}}{n}{\textsci}{\textsubring{r}}} \dicPos{m}[7] \dicFlx{(‑is, ‑ar)}[1] \dicDirectTranslationCS{doktor(ka), lékař(ka)} \dicExampleIS{panta tíma hjá lækni} \dicExampleCS{objednat se k~lékaři}
\dicEntry[læknisfræði] \dicTerm{læknis··fræð|i} \dicIPA{{l}{a}{i}{h}{\r{g}}{n}{\textsci}{s}{f}{r}{a}{i}{ð}{\textsci}} \dicPos{f}[3] \dicFlx{(‑i)}[3] \dicDirectTranslationCS{lékařství, medicína} \dicIndirectTranslationCS{(nauka)}
\dicEntry[læknisskoðun] \dicTerm{læknis··skoð|un} \dicIPA{{l}{a}{i}{h}{\r{g}}{n}{\textsci}{s}{\r{g}}{\textopeno}{ð}{\textscy}{\textsubring{n}}} \dicPos{f}[7] \dicFlx{(‑unar)}[9] \dicDirectTranslationCS{zdravotní\,/\addthin lékařská prohlídka}
\dicEntry[læknisvottorð] \dicTerm{læknis··vott·orð} \dicIPA{{l}{a}{i}{h}{\r{g}}{n}{\textsci}{s}{v}{\textopeno}{h}{\textsubring{d}}{\textopeno}{r}{\texttheta}} \dicPos{n}[2] \dicFlx{(‑s, ‑)}[5] \dicDirectTranslationCS{lékařská zpráva}
\dicEntry[læknisþjónusta] \dicTerm{læknis··þjón·ust|a} \dicIPA{{l}{a}{i}{h}{\r{g}}{n}{\textsci}{s}{\texttheta}{j}{ou}{n}{\textscy}{s}{\textsubring{d}}{a}} \dicPos{f}[1] \dicFlx{(‑u)}[5] \dicDirectTranslationCS{zdravot\-ní\,/\addthin lékařská služba}
\dicEntry[lækur] \dicTerm{læk|ur} \dicsymFrequent\  \dicIPA{{l}{a}{i}{\textlengthmark}{\r{g}}{\textscy}{\textsubring{r}}} \dicPos{m}[10] \dicFlx{(‑jar, ‑ir)}[23] \textbf{1.} \dicDirectTranslationCS{potok, říčka} \dicExampleIS{ár og lækir} \dicExampleCS{řeky a~potoky}  \textbf{2.} \dicFieldCat{anat.} \dicSynonym{naflastrengur} \dicDirectTranslationCS{pupeční šňůra}
\dicEntry[læmingi] \dicTerm{læm··ing|i} \dicIPA{{l}{a}{i}{\textlengthmark}{m}{i}{\textltailn}{\r{\textObardotlessj}}{\textsci}} \dicPos{m}[1] \dicFlx{(‑ja, ‑jar)}[14] \dicFieldCat{zool.} \dicDirectTranslationCS{lumík norský} \textit{(l.~{\textLA{Lemmus lemmus}})}  \dicsymPhoto\ 
\dicFigure{ds_image_laemingi_0_1.jpg}{Læmingi}{Læmingi - Argus fin, PD}
\dicEntry[læra] \dicTerm{lær|a} \dicsymFrequent\  \dicIPA{{l}{a}{i}{\textlengthmark}{r}{a}} \dicPos{v}[2] \dicFlx{(‑ði, ‑t)}[103] \dicFlx{acc} \dicDirectTranslationCS{(na)učit se, studovat} \dicExampleIS{læra íslensku} \dicExampleCS{učit se islandštinu};  \dicPhraseIS{læra e‑ð utanbókar} \dicDirectTranslationCS{učit se (co) zpaměti, memorovat (co)} \dicExampleIS{læra utanbókar ljóð} \dicExampleCS{učit se básničku nazpaměť};  \dicIdiom{læra}[af]{ \dicPhraseIS{læra af e‑m}} \dicDirectTranslationCS{učit se od (koho)};  \dicIdiom{læra}[á]{ \dicPhraseIS{læra á e‑ð}} \dicDirectTranslationCS{učit se na (co) (na hudební nástroj ap.)} \dicExampleIS{Ég lærði á píanó.} \dicExampleCS{Učil jsem se na piano.};  \dicIdiom{læra}[fyrir]{ \dicPhraseIS{læra fyrir e‑ð}} \dicDirectTranslationCS{učit se na (co)} \dicExampleIS{læra fyrir próf} \dicExampleCS{učit se na zkoušku};  \dicIdiom{læra}[undir]{ \dicPhraseIS{læra undir e‑ð}} \dicDirectTranslationCS{učit se na (co) (na zkoušku ap.)} \dicExampleIS{læra undir skóla} \dicExampleCS{učit se do školy};  \dicIdiom{læra}[utan að]{ \dicPhraseIS{læra utan að}} \dicDirectTranslationCS{učit se zpaměti, memorovat};  \dicIdiom{lærast}{ \dicPhraseIS{e‑að lærist vel}} \dicFlx{refl} \dicDirectTranslationCS{(co) se učí dobře}; { \dicPhraseIS{e‑m lærist að (gera e‑ð)}} \dicFlx{refl impers} \dicDirectTranslationCS{(kdo) se učí (dělat (co)), (kdo) si osvojuje (dělat (co))};  \dicProverb\  \dicPhraseIS{Svo lengi lærir sem lifir.} \dicLangCat{přís.} \dicDirectTranslationCS{Člověk se učí po celý život.}
\dicEntry[lærbein] \dicTerm{lær··bein} \dicIPA{{l}{a}{i}{r}{\textsubring{b}}{ei}{\textsubring{n}}} \dicPos{n}[2] \dicFlx{(‑s, ‑)}[5] \dicFieldCat{anat.} \dicDirectTranslationCS{stehenní kost}
\dicEntry[lærdómur] \dicTerm{lær··dóm|ur} \dicIPA{{l}{a}{i}{r}{\textsubring{d}}{ou}{m}{\textscy}{\textsubring{r}}} \dicPos{m}[6] \dicFlx{(‑s, ‑ar)}[10] \dicSynonym{þekking} \dicDirectTranslationCS{učení, učenost, vzdělání};  \dicPhraseIS{draga lærdóma af e‑u} \dicDirectTranslationCS{poučit se z~(čeho)}
\dicEntry[lærður] \dicTerm{lærður} \dicsymFrequent\  \dicIPA{{l}{a}{i}{r}{ð}{\textscy}{\textsubring{r}}} \dicPos{adj}[2]\dicFlx{}[-1] \dicDirectTranslationCS{učený, vzdělaný} \dicExampleIS{mjög lærður maður} \dicExampleCS{velmi vzdělaný člověk}
\dicEntry[læri] \dicTerm{læri} \dicsymFrequent\  \dicIPA{{l}{a}{i}{\textlengthmark}{r}{\textsci}} \dicPos{n}[2] \dicFlx{(‑s, ‑)}[14] \textbf{1.} \dicFieldCat{anat.} \dicDirectTranslationCS{stehno}  \textbf{2.} \dicDirectTranslationCS{kýta} \dicExampleIS{læri af lambi} \dicExampleCS{jehněčí kýta}  \textbf{3.} \dicSynonym{nám} \dicDirectTranslationCS{učení} \dicIndirectTranslationCS{(učňovská léta)};  \dicPhraseIS{vera í læri hjá e‑m} \dicDirectTranslationCS{být u~(koho) v~učení}
\dicEntry[lærifaðir] \dicTerm{læri··|faðir} \dicIPA{{l}{a}{i}{\textlengthmark}{r}{\textsci}{f}{a}{ð}{\textsci}{\textsubring{r}}} \dicPos{m}[13] \dicFlx{(‑föður, ‑feður)}[3] \dicSynonym{kennari} \dicDirectTranslationCS{duchovní otec, mistr, guru}
\dicEntry[lærisveinn] \dicTerm{læri··svein|n} \dicIPA{{l}{a}{i}{\textlengthmark}{r}{\textsci}{s}{v}{ei}{\textsubring{d}}{\textsubring{n}}} \dicPos{m}[6] \dicFlx{(‑s, ‑ar)}[42] \dicSynonym{nemandi} \dicDirectTranslationCS{učeň(ka), učnice, učedník, učednice, žák(yně), žačka} \dicExampleIS{lærisveinar Jesú} \dicExampleCS{učedníci Ježíše}
\dicEntry[lærleggur] \dicTerm{lær··legg|ur} \dicIPA{{l}{a}{i}{r}{l}{\textepsilon}{\r{g}}{\textscy}{\textsubring{r}}} \dicPos{m}[9] \dicFlx{(‑jar\,/\addthin ‑s, ‑ir)}[26] \dicFieldCat{anat.} \dicDirectTranslationCS{stehenní kost}
\dicEntry[lærlingur] \dicTerm{lær··ling|ur} \dicIPA{{l}{a}{i}{r}{l}{i}{\ng}{\r{g}}{\textscy}{\textsubring{r}}} \dicPos{m}[6] \dicFlx{(‑s, ‑ar)}[8] \dicSynonym{nemandi} \dicDirectTranslationCS{učedník, učednice, učeň(ka), učnice}
\dicEntry[læs] \dicTerm{læs} \dicIPA{{l}{a}{i}{\textlengthmark}{s}} \dicPos{adj}[5]\dicFlx{}[-1] \dicDirectTranslationCS{gramotný, umějící číst}
\dicEntry[læsa] \dicTerm{læs|a} \dicsymFrequent\  \dicIPA{{l}{a}{i}{\textlengthmark}{s}{a}} \dicPos{v}[2] \dicFlx{(‑ti, ‑t)}[63] \dicFlx{dat\,/\addthin acc} \textbf{1.} \dicFlx{dat} \dicSynonym*{loka með lás} \dicDirectTranslationCS{(u)zamknout, (u)zamykat, (u)zavřít, (u)zavírat} \dicExampleIS{Dyrunum var læst.} \dicExampleCS{Dveře byly zamčeny.};  \dicPhraseIS{læsa að sér} \dicDirectTranslationCS{zamknout (za sebou) dveře};  \dicPhraseIS{læsa á eftir sér} \dicDirectTranslationCS{zavřít za sebou (dveře)};  \dicPhraseIS{læsa e‑n inni} \dicDirectTranslationCS{zavřít (koho), dát (koho) pod zámek (do vězení ap.)}  \textbf{2.} \dicFlx{dat} \dicSynonym{innsigla} \dicDirectTranslationCS{zalepit, zalepovat, zapečetit, zapečeťovat (obálku ap.)}  \textbf{3.} \dicFlx{acc} \dicSynonym*{breiðast út} \dicDirectTranslationCS{(roz)šířit se, rozšiřovat se};  \dicPhraseIS{læsa sig um e‑ð} \dicDirectTranslationCS{rozšířit se po (čem)}  \textbf{4.} \dicFlx{dat} \dicDirectTranslationCS{sevřít, svírat (v~drápech\,/\addthin zubech)};  \dicIdiom{læsast}{ \dicPhraseIS{læsast}} \dicFlx{refl} {\textbf{a.}} \dicDirectTranslationCS{zavřít se (klíče v~domě ap.)};  {\textbf{b.}} \dicDirectTranslationCS{sevřít, chopit se}
\dicEntry[læsi] \dicTerm{læsi\smash{\textsuperscript{1}}} \dicIPA{{l}{a}{i}{\textlengthmark}{s}{\textsci}} \dicPos{n}[2] \dicFlx{(‑s)}[20] \dicDirectTranslationCS{gramotnost, dovednost čtení}
\dicEntry[læsi] \dicTerm{læsi\smash{\textsuperscript{2}}} \dicIPA{{l}{a}{i}{\textlengthmark}{s}{\textsci}} \dicPos{v} \dicFlx{con pf sg 1 pers} \dicLink{lesa}
\dicEntry[læsilegur] \dicTerm{læsi··legur} \dicIPA{{l}{a}{i}{\textlengthmark}{s}{\textsci}{l}{\textepsilon}{\textbabygamma}{\textscy}{\textsubring{r}}} \dicPos{adj}[1]\dicFlx{}[-8] \textbf{1.} \dicDirectTranslationCS{čitelný} \dicExampleIS{læsileg skrift} \dicExampleCS{čitelné písmo}  \textbf{2.} \dicDirectTranslationCS{čtivý} \dicExampleIS{læsileg bók} \dicExampleCS{čtivá knížka}
\dicEntry[læsing] \dicTerm{læs··ing} \dicIPA{{l}{a}{i}{\textlengthmark}{s}{i}{\ng}{\r{g}}} \dicPos{f}[4] \dicFlx{(‑ar, ‑ar)}[5] \textbf{1.} \dicSynonym{lás} \dicDirectTranslationCS{zámek, zamykání (u~dveří ap.)} \dicExampleIS{læsing á dyrunum} \dicExampleCS{zámek u~dveří}  \textbf{2.} \dicSynonym{lokun} \dicDirectTranslationCS{(u)zamčení, (u)zamknutí}
\dicEntry[læstur] \dicTerm{læstur} \dicsymFrequent\  \dicIPA{{l}{a}{i}{s}{\textsubring{d}}{\textscy}{\textsubring{r}}} \dicPos{adj}[1]\dicFlx{}[-10] \dicSynonym*{aflokaður} \dicDirectTranslationCS{(u)zamčený, (u)zamknutý} \dicExampleIS{læstar dyr} \dicExampleCS{zamknuté dveře}
\dicEntry[læt] \dicTerm{læt} \dicIPA{{l}{a}{i}{\textlengthmark}{\textsubring{d}}} \dicPos{v} \dicFlx{ind praes sg 1 pers} \dicLink{láta}
\dicEntry[læti] \dicTerm{læti} \dicsymFrequent\  \dicIPA{{l}{a}{i}{\textlengthmark}{\textsubring{d}}{\textsci}} \dicPos{n}[2] \dicFlx{pl}[19] \dicDirectTranslationCS{rozruch, povyk} \dicExampleIS{vera með læti} \dicExampleCS{dělat povyk};  \dicPhraseIS{kunna sér ekki læti} \dicLangCat{přen.} \dicDirectTranslationCS{být štěstím bez sebe}
\dicEntry[lævirki] \dicTerm{læ··virk|i} \dicIPA{{l}{a}{i}{\textlengthmark}{v}{\textsci}{\textsubring{r}}{\r{\textObardotlessj}}{\textsci}} \dicPos{m}[1] \dicFlx{(‑ja, ‑jar)}[14] \dicFieldCat{zool.} \dicDirectTranslationCS{skřivan} \textit{(l.~{\textLA{Alauda}})}  \dicsymPhoto\ 
\dicFigure{ds_image_laevirki_0_2.jpg}{Lævirki}{Lævirki - Daniel Pettersson, CC BY-SA 2.5}
\dicEntry[lævís] \dicTerm{læ··vís} \dicIPA{{l}{a}{i}{\textlengthmark}{v}{i}{s}} \dicPos{adj}[5]\dicFlx{}[-1] \dicSynonym{slóttugur} \dicDirectTranslationCS{záludný, zákeřný}
\dicEntry[löðrungur] \dicTerm{löðr··ung|ur} \dicIPA{{l}{\oe}{ð}{r}{u}{\ng}{\r{g}}{\textscy}{\textsubring{r}}} \dicPos{m}[6] \dicFlx{(‑s, ‑ar)}[8] \dicSynonym{kinnhestur} \dicDirectTranslationCS{facka, políček} \dicExampleIS{gefa e‑m löðrungur} \dicExampleCS{dát (komu) facku}
\dicEntry[löður] \dicTerm{löður} \dicIPA{{l}{\oe}{\textlengthmark}{ð}{\textscy}{\textsubring{r}}} \dicPos{n}[2] \dicFlx{(‑s)}[28] \dicSynonym{froða} \dicDirectTranslationCS{(vodní) pěna}
\dicEntry[lög] \dicTerm{lög} \dicIPA{{l}{\oe}{\textlengthmark}{x}} \dicPos{n}[2] \dicFlx{pl}[9] \dicFieldCat{práv.} \dicDirectTranslationCS{zákon};  \dicPhraseIS{leiða e‑ð í lög} \dicDirectTranslationCS{uzákonit (co)};  \dicPhraseIS{setja lög} \dicDirectTranslationCS{vydat zákon}
\dicEntry[lögaldur] \dicTerm{lög··ald|ur} \dicIPA{{l}{\oe}{\textlengthmark}{\textbabygamma}{a}{l}{\textsubring{d}}{\textscy}{\textsubring{r}}} \dicPos{m}[5] \dicFlx{(‑urs, ‑rar)}[3] \dicFieldCat{práv.} \dicDirectTranslationCS{plnoletost, zletilost, zákonný věk};  \dicPhraseIS{ná lögaldri} \dicDirectTranslationCS{dosáhnout plnoletosti}
\dicEntry[lögbann] \dicTerm{lög··bann} \dicIPA{{l}{\oe}{\textbabygamma}{\textsubring{b}}{a}{\textsubring{n}}} \dicPos{n}[2] \dicFlx{(‑s)}[2] \dicFieldCat{práv.} \dicDirectTranslationCS{soudní zákaz}
\dicEntry[lögbrot] \dicTerm{lög··brot} \dicIPA{{l}{\oe}{\textbabygamma}{\textsubring{b}}{r}{\textopeno}{\textsubring{d}}} \dicPos{n}[2] \dicFlx{(‑s, ‑)}[5] \dicFieldCat{práv.} \dicDirectTranslationCS{porušení zákona, delikt}
\dicEntry[lögfesta] \dicTerm{lög··fest|a} \dicIPA{{l}{\oe}{x}{f}{\textepsilon}{s}{\textsubring{d}}{a}} \dicPos{v}[2] \dicFlx{(‑i, ‑)}[12] \dicFlx{acc} \dicFieldCat{práv.} \dicDirectTranslationCS{uzákonit, legalizovat, kodifikovat}
\dicEntry[lögformlegur] \dicTerm{lög··form·legur} \dicIPA{{l}{\oe}{x}{f}{\textopeno}{r}{m}{l}{\textepsilon}{\textbabygamma}{\textscy}{\textsubring{r}}} \dicPos{adj}[1]\dicFlx{}[-8] \dicFieldCat{práv.} \dicDirectTranslationCS{zákonný}
\dicEntry[lögfræði] \dicTerm{lög··fræð|i} \dicIPA{{l}{\oe}{x}{f}{r}{a}{i}{ð}{\textsci}} \dicPos{f}[3] \dicFlx{(‑i)}[3] \dicDirectTranslationCS{právní věda, práva}
\dicEntry[lögfræðideild] \dicTerm{lög·fræði··deild} \dicIPA{{l}{\oe}{x}{f}{r}{a}{i}{ð}{\textsci}{\textsubring{d}}{ei}{l}{\textsubring{d}}} \dicPos{f}[7] \dicFlx{(‑ar, ‑ir)}[1] \dicFieldCat{škol.} \dicDirectTranslationCS{právnická fakulta}
\dicEntry[lögfræðilegur] \dicTerm{lög·fræði··legur} \dicIPA{{l}{\oe}{x}{f}{r}{a}{i}{ð}{\textsci}{l}{\textepsilon}{\textbabygamma}{\textscy}{\textsubring{r}}} \dicPos{adj}[1]\dicFlx{}[-8] \dicDirectTranslationCS{právní} \dicExampleIS{lögfræðilegur ráðgjafi} \dicExampleCS{právní poradce}
\dicEntry[lögfræðingur] \dicTerm{lög·fræð··ing|ur} \dicsymFrequent\  \dicIPA{{l}{\oe}{x}{f}{r}{a}{i}{ð}{i}{\ng}{\r{g}}{\textscy}{\textsubring{r}}} \dicPos{m}[6] \dicFlx{(‑s, ‑ar)}[8] \dicDirectTranslationCS{právník, právnička} \dicExampleIS{aðstoð lögfræðings} \dicExampleCS{pomoc právníka}
\dicEntry[lögg] \dicTerm{lögg} \dicIPA{{l}{\oe}{\r{g}}{\textlengthmark}} \dicPos{f}[7] \dicFlx{(laggar\,/\addthin ‑var, laggir\,/\addthin ‑ir\,/\addthin ‑var)}[30] \dicDirectTranslationCS{kapka, doušek (v~láhvi ap.)} \dicExampleIS{kaffilögg} \dicExampleCS{doušek kávy};  \dicPhraseIS{setja e‑ð á laggirnar} \dicDirectTranslationCS{založit (co), zřídit (co)} \dicExampleIS{setja fyrirtæki á laggirnar} \dicExampleCS{zřídit společnost};  \dicPhraseIS{koma e‑u á laggirnar} \dicDirectTranslationCS{založit (co), zřídit (co)}
\dicEntry[lögga] \dicTerm{lögg|a} \dicsymFrequent\  \dicIPA{{l}{\oe}{\r{g}}{\textlengthmark}{a}} \dicPos{f}[1] \dicFlx{(‑u, ‑ur)}[7] \dicLangCat{hovor.} \dicSynonym{lögreglumaður} \dicDirectTranslationCS{polda, policajt(ka)} \dicExampleIS{löggur í kvikmyndum} \dicExampleCS{poldové ve filmech}
\dicEntry[löggilda] \dicTerm{lög··gil|da} \dicIPA{{l}{\oe}{\textbabygamma}{\r{\textObardotlessj}}{\textsci}{l}{\textsubring{d}}{a}} \dicPos{v}[2] \dicFlx{(‑ti, ‑t)}[37] \dicFlx{acc} \dicFieldCat{práv.} \dicDirectTranslationCS{oprávnit, zmocnit, autorizovat}
\dicEntry[löggilding] \dicTerm{lög··gild·ing} \dicIPA{{l}{\oe}{\textbabygamma}{\r{\textObardotlessj}}{\textsci}{l}{\textsubring{d}}{i}{\ng}{\r{g}}} \dicPos{f}[4] \dicFlx{(‑ar, ‑ar)}[5] \dicFieldCat{práv.} \dicDirectTranslationCS{oprávnění, zmocnění, autorizace}
\dicEntry[löggiltur] \dicTerm{lög··giltur} \dicIPA{{l}{\oe}{\textbabygamma}{\r{\textObardotlessj}}{\textsci}{\textsubring{l}}{\textsubring{d}}{\textscy}{\textsubring{r}}} \dicPos{adj}[1]\dicFlx{}[-13] \dicFieldCat{práv.} \dicDirectTranslationCS{oprávněný, autorizovaný};  \dicPhraseIS{löggiltur skjalaþýðandi} \dicDirectTranslationCS{úřední překladatel(ka)}
\dicEntry[löggjafarvald] \dicTerm{lög·gjafar··vald} \dicIPA{{l}{\oe}{\textbabygamma}{\r{\textObardotlessj}}{a}{v}{a}{r}{v}{a}{l}{\textsubring{d}}} \dicPos{n}[2] \dicFlx{(‑s)}[2] \dicDirectTranslationCS{zákonodárná moc}
\dicEntry[löggjafarþing] \dicTerm{lög·gjafar··þing} \dicIPA{{l}{\oe}{\textbabygamma}{\r{\textObardotlessj}}{a}{f}{a}{\textsubring{r}}{\texttheta}{i}{\ng}{\r{g}}} \dicPos{n}[2] \dicFlx{(‑s, ‑)}[5] \dicFieldCat{práv.} \dicDirectTranslationCS{zákonodárný sbor}
\dicEntry[löggjafi] \dicTerm{lög··gjaf|i} \dicIPA{{l}{\oe}{\textbabygamma}{\r{\textObardotlessj}}{a}{v}{\textsci}} \dicPos{m}[1] \dicFlx{(‑a, ‑ar)}[8] \dicFieldCat{pol.} \dicDirectTranslationCS{zákonodárce, zákonodárkyně}
\dicEntry[löggjöf] \dicTerm{lög··|gjöf} \dicIPA{{l}{\oe}{\textbabygamma}{\r{\textObardotlessj}}{\oe}{f}} \dicPos{f}[7] \dicFlx{(‑gjafar)}[19] \dicFieldCat{práv.} \dicDirectTranslationCS{legislativa, zákonodárství}
\dicEntry[löggæsla] \dicTerm{lög··gæsl|a} \dicIPA{{l}{\oe}{\textbabygamma}{\r{\textObardotlessj}}{a}{i}{s}{\textsubring{d}}{l}{a}} \dicPos{f}[1] \dicFlx{(‑u)}[5] \dicDirectTranslationCS{prosazování práva}
\dicEntry[löghald] \dicTerm{lög··hald} \dicIPA{{l}{\oe}{\textlengthmark}{x}{h}{a}{l}{\textsubring{d}}} \dicPos{n}[2] \dicFlx{(‑s)}[2] \dicFieldCat{práv.} \dicDirectTranslationCS{(úřední) zabavení, konfiskace};  \dicPhraseIS{leggja löghald á e‑u} \dicDirectTranslationCS{zkonfiskovat (co)}
\dicEntry[lögheimili] \dicTerm{lög··heimili} \dicIPA{{l}{\oe}{\textlengthmark}{x}{h}{ei}{m}{\textsci}{l}{\textsci}} \dicPos{n}[2] \dicFlx{(‑s, ‑)}[14] \dicFieldCat{práv.} \dicDirectTranslationCS{trvalé bydliště}
\dicEntry[löghlýðinn] \dicTerm{lög··hlýðinn} \dicIPA{{l}{\oe}{x}{\textsubring{l}}{i}{ð}{\textsci}{\textsubring{n}}} \dicPos{adj}[6]\dicFlx{}[-2] \dicDirectTranslationCS{dodržující zákony}
\dicEntry[löghyggja] \dicTerm{lög··hyggj|a} \dicIPA{{l}{\oe}{\textlengthmark}{x}{h}{\textsci}{\r{\textObardotlessj}}{a}} \dicPos{f}[1] \dicFlx{(‑u)}[5] \dicFieldCat{filos.} \dicDirectTranslationCS{determinismus}
\dicEntry[löglegur] \dicTerm{lög··legur} \dicIPA{{l}{\oe}{\textbabygamma}{l}{\textepsilon}{\textbabygamma}{\textscy}{\textsubring{r}}} \dicPos{adj}[1]\dicFlx{}[-8] \dicDirectTranslationCS{legální, právoplatný}
\dicEntry[lögmaður] \dicTerm{lög··|maður} \dicsymFrequent\  \dicIPA{{l}{\oe}{\textbabygamma}{m}{a}{ð}{\textscy}{\textsubring{r}}} \dicPos{m}[13] \dicFlx{(‑manns, ‑menn)}[2] \dicDirectTranslationCS{advokát(ka)} \dicExampleIS{Vantar þig lögmann?} \dicExampleCS{Potřebuješ advokáta?}
\dicEntry[lögmál] \dicTerm{lög··mál} \dicIPA{{l}{\oe}{\textbabygamma}{m}{au}{\textsubring{l}}} \dicPos{n}[2] \dicFlx{(‑s, ‑)}[5] \textbf{1.} \dicSynonym{lög} \dicDirectTranslationCS{zákon}  \textbf{2.} \dicDirectTranslationCS{zákon, řád (přírody ap.)} \dicExampleIS{lögmál lífsins} \dicExampleCS{zákon života}  \textbf{3.} \dicFieldCat{filos.} \dicSynonym{grundvallarregla} \dicDirectTranslationCS{princip, zákonitost}
\dicEntry[lögmæti] \dicTerm{lög··mæti} \dicIPA{{l}{\oe}{\textbabygamma}{m}{a}{i}{\textsubring{d}}{\textsci}} \dicPos{n}[2] \dicFlx{(‑s)}[20] \dicDirectTranslationCS{legálnost, zákonnost}
\dicEntry[lögmætur] \dicTerm{lög··mætur} \dicIPA{{l}{\oe}{\textbabygamma}{m}{a}{i}{\textsubring{d}}{\textscy}{\textsubring{r}}} \dicPos{adj}[1]\dicFlx{}[-1] \dicSynonym{gildur} \dicDirectTranslationCS{legální, zákonný}
\dicEntry[lögn] \dicTerm{lögn} \dicIPA{{l}{\oe}{\r{g}}{\textsubring{n}}} \dicPos{f}[7] \dicFlx{(lagnar, lagnir)}[16] \textbf{1.} \dicSynonym{leiðsla} \dicDirectTranslationCS{vedení, instalace} \dicExampleIS{raflögn} \dicExampleCS{elektrické vedení}  \textbf{2.} \dicSynonym*{það að leggja net} \dicDirectTranslationCS{pokládání sítě}
\dicEntry[lögregla] \dicTerm{lög··regl|a} \dicsymFrequent\  \dicIPA{{l}{\oe}{\textbabygamma}{r}{\textepsilon}{\r{g}}{l}{a}} \dicPos{f}[1] \dicFlx{(‑u, ‑ur)}[19] \textbf{1.} \dicSynonym*{stofnun til löggæslu} \dicDirectTranslationCS{policie} \dicExampleIS{Lögreglan leitar að vitnum að árekstri.} \dicExampleCS{Policie hledá svědky srážky vozidel.}  \textbf{2.} \dicSynonym{lögreglumaður} \dicDirectTranslationCS{policista, policistka}
\dicEntry[lögreglubíll] \dicTerm{lög·reglu··bíl|l} \dicsymFrequent\  \dicIPA{{l}{\oe}{\textbabygamma}{r}{\textepsilon}{\r{g}}{l}{\textscy}{\textsubring{b}}{i}{\textsubring{d}}{\textsubring{l}}} \dicPos{m}[6] \dicFlx{(‑s, ‑ar)}[48] \dicDirectTranslationCS{policejní auto\,/\addthin vůz} \dicExampleIS{Lögreglubílar lentu í árekstri.} \dicExampleCS{Policejní auta se srazila.}
\dicEntry[lögreglumaður] \dicTerm{lög·reglu··|maður} \dicsymFrequent\  \dicIPA{{l}{\oe}{\textbabygamma}{r}{\textepsilon}{\r{g}}{l}{\textscy}{m}{a}{ð}{\textscy}{\textsubring{r}}} \dicPos{m}[13] \dicFlx{(‑manns, ‑menn)}[2] \dicDirectTranslationCS{policista, policistka} \dicExampleIS{Lögreglumenn réðust inn í húsið.} \dicExampleCS{Policisté vtrhli do domu.}
\dicEntry[lögreglusamþykkt] \dicTerm{lög·reglu··sam·þykkt} \dicIPA{{l}{\oe}{\textbabygamma}{r}{\textepsilon}{\r{g}}{l}{\textscy}{s}{a}{m}{\texttheta}{\textsci}{x}{\textsubring{d}}} \dicPos{f}[7] \dicFlx{(‑ar, ‑ir)}[1] \dicDirectTranslationCS{policejní nařízení}
\dicEntry[lögreglustjóri] \dicTerm{lög·reglu··stjór|i} \dicIPA{{l}{\oe}{\textbabygamma}{r}{\textepsilon}{\r{g}}{l}{\textscy}{s}{\textsubring{d}}{j}{ou}{r}{\textsci}} \dicPos{m}[1] \dicFlx{(‑a, ‑ar)}[1] \dicDirectTranslationCS{policejní prezident(ka)}
\dicEntry[lögreglustöð] \dicTerm{lög·reglu··stöð} \dicIPA{{l}{\oe}{\textbabygamma}{r}{\textepsilon}{\r{g}}{l}{\textscy}{s}{\textsubring{d}}{\oe}{\texttheta}} \dicPos{f}[6] \dicFlx{(‑var, ‑var)}[1] \dicDirectTranslationCS{policejní stanice}
\dicEntry[lögregluvarðstjóri] \dicTerm{lög·reglu··varð·stjór|i} \dicIPA{{l}\-{\oe}\-{\textbabygamma}\-{r}\-{\textepsilon}\-{\r{g}}\-{l}\-{\textscy}\-{v}\-{a}\-{r}\-{ð}\-{s}\-{\textsubring{d}}\-{j}\-{ou}\-{r}\-{\textsci}\-} \dicPos{m}[1] \dicFlx{(‑a, ‑ar)}[1] \dicDirectTranslationCS{policejní komisař(ka)}
\dicEntry[lögregluvarðstofa] \dicTerm{lög·reglu··varð·stof|a} \dicIPA{{l}\-{\oe}\-{\textbabygamma}\-{r}\-{\textepsilon}\-{\r{g}}\-{l}\-{\textscy}\-{v}\-{a}\-{r}\-{ð}\-{s}\-{\textsubring{d}}\-{\textopeno}\-{v}\-{a}\-} \dicPos{f}[1] \dicFlx{(‑u, ‑ur)}[7] \dicDirectTranslationCS{okrsková služebna, policejní stanice}
\dicEntry[lögregluþjónn] \dicTerm{lög·reglu··þjón|n} \dicsymFrequent\  \dicIPA{{l}{\oe}{\textbabygamma}{r}{\textepsilon}{\r{g}}{l}{\textscy}{\texttheta}{j}{ou}{\textsubring{d}}{\textsubring{n}}} \dicPos{m}[6] \dicFlx{(‑s, ‑ar)}[42] \dicDirectTranslationCS{policista, policistka} \dicExampleIS{lögregluþjónar á næturvakt} \dicExampleCS{policisté na noční službě}
\dicEntry[Lögrétta] \dicTerm{Lög··rétt|a} \dicIPA{{l}{\oe}{\textbabygamma}{r}{j}{\textepsilon}{h}{\textsubring{d}}{a}} \dicPos{f}[1] \dicFlx{(‑u)}[6] \dicFieldCat{geog., hist.} \dicDirectTranslationCS{Lögrétta} \dicIndirectTranslationCS{(zákonodárný a~soudní sněm nacházející se v~Altingu)}
\dicEntry[lögræði] \dicTerm{lög··ræði} \dicIPA{{l}{\oe}{\textbabygamma}{r}{a}{i}{ð}{\textsci}} \dicPos{n}[2] \dicFlx{(‑s)}[20] \dicFieldCat{práv.} \dicDirectTranslationCS{svéprávnost};  \dicPhraseIS{svipta e‑n lögræði} \dicDirectTranslationCS{zbavit (koho) svéprávnosti}
\dicEntry[lögsótt] \dicTerm{lög··sótt} \dicIPA{{l}{\oe}{x}{s}{ou}{h}{\textsubring{d}}} \dicPos{v} \dicFlx{supin} \dicLink{lögsækja}
\dicEntry[lögsótti] \dicTerm{lög··sótti} \dicIPA{{l}{\oe}{x}{s}{ou}{h}{\textsubring{d}}{\textsci}} \dicPos{v} \dicFlx{ind pf sg 1 pers} \dicLink{lögsækja}
\dicEntry[lögsóttum] \dicTerm{lög··sóttum} \dicIPA{{l}{\oe}{x}{s}{ou}{h}{\textsubring{d}}{\textscy}{\textsubring{m}}} \dicPos{v} \dicFlx{ind pf pl 1 pers} \dicLink{lögsækja}
\dicEntry[lögsæki] \dicTerm{lög··sæki} \dicIPA{{l}{\oe}{x}{s}{a}{i}{\r{\textObardotlessj}}{\textsci}} \dicPos{v} \dicFlx{ind praes sg 1 pers} \dicLink{lögsækja}
\dicEntry[lögsækja] \dicTerm{lög··|sækja} \dicIPA{{l}{\oe}{x}{s}{a}{i}{\r{\textObardotlessj}}{a}} \dicPos{v}[5] \dicFlx{(‑sæki, ‑sótti, ‑sóttum, ‑sækti, ‑sótt)}[7] \dicFlx{acc} \dicFieldCat{práv.} \dicDirectTranslationCS{žalovat, podat\,/\addthin podávat žalobu} \dicExampleIS{lögsækja e‑n fyrir e‑ð} \dicExampleCS{žalovat (koho) za (co)}
\dicEntry[lögsækti] \dicTerm{lög··sækti} \dicIPA{{l}{\oe}{x}{s}{a}{i}{x}{\textsubring{d}}{\textsci}} \dicPos{v} \dicFlx{con pf sg 1 pers} \dicLink{lögsækja}
\dicEntry[lögtak] \dicTerm{lög··|tak} \dicIPA{{l}{\oe}{x}{t\smash{\textsuperscript{h}}}{a}{\r{g}}} \dicPos{n}[2] \dicFlx{(‑taks, ‑tök)}[8] \dicFieldCat{práv.} \dicDirectTranslationCS{exekuce}
\dicEntry[lögun] \dicTerm{lögun} \dicsymFrequent\  \dicIPA{{l}{\oe}{\textlengthmark}{\textbabygamma}{\textscy}{\textsubring{n}}} \dicPos{f}[7] \dicFlx{(‑ar, laganir)}[16] \textbf{1.} \dicSynonym{form} \dicDirectTranslationCS{tvar, podoba} \dicExampleIS{breyta löguninni} \dicExampleCS{změnit podobu}  \textbf{2.} \dicSynonym*{kaffilögun o.þ.h.} \dicDirectTranslationCS{příprava (kávy ap.)} \dicExampleIS{lögun á kaffi} \dicExampleCS{příprava kávy}
\dicEntry[lögur] \dicTerm{lögur} \dicIPA{{l}{\oe}{\textlengthmark}{\textbabygamma}{\textscy}{\textsubring{r}}} \dicPos{m}[11] \dicFlx{(lagar, legir)}[5] \textbf{1.} \dicDirectTranslationCS{kapalina, tekutina}  \textbf{2.} \dicSynonym{vatn} \dicDirectTranslationCS{voda, vodní plocha} \dicExampleIS{á láði og legi} \dicExampleCS{na souši a~na vodě}
\dicEntry[lök] \dicTerm{lök\smash{\textsuperscript{1}}} \dicIPA{{l}{\oe}{\textlengthmark}{\r{g}}} \dicPos{n} \dicFlx{pl nom} \dicLink{lak\smash{\textsuperscript{1}}}
\dicEntry[lök] \dicTerm{lök\smash{\textsuperscript{2}}} \dicIPA{{l}{\oe}{\textlengthmark}{\r{g}}} \dicPos{adj} \dicFlx{f sg nom pos} \dicLink{lakur}
\dicEntry[lökk] \dicTerm{lökk} \dicIPA{{l}{\oe}{h}{\r{g}}} \dicPos{n} \dicFlx{pl nom} \dicLink{lakk}
\dicEntry[lömb] \dicTerm{lömb} \dicIPA{{l}{\oe}{m}{\textsubring{b}}} \dicPos{n} \dicFlx{pl nom} \dicLink{lamb}
\dicEntry[lömdum] \dicTerm{lömdum} \dicIPA{{l}{\oe}{m}{\textsubring{d}}{\textscy}{\textsubring{m}}} \dicPos{v} \dicFlx{ind pf pl 1 pers} \dicLink{lemja}
\dicEntry[lömuð] \dicTerm{lömuð} \dicIPA{{l}{\oe}{\textlengthmark}{m}{\textscy}{\texttheta}} \dicPos{adj} \dicFlx{f sg nom pos} \dicLink{lamaður}
\dicEntry[lömun] \dicTerm{löm|un} \dicIPA{{l}{\oe}{\textlengthmark}{m}{\textscy}{\textsubring{n}}} \dicPos{f}[7] \dicFlx{(‑unar)}[12] \dicFieldCat{med.} \dicDirectTranslationCS{ochrnutí, paralýza} \dicExampleIS{fá lömun í fótinn} \dicExampleCS{ochrnout na nohu}
\dicEntry[lömunarveiki] \dicTerm{lömunar··veik|i} \dicIPA{{l}{\oe}{\textlengthmark}{m}{\textscy}{n}{a}{r}{v}{ei}{\r{\textObardotlessj}}{\textsci}} \dicPos{f}[3] \dicFlx{(‑i)}[3] \dicFieldCat{med.} \dicDirectTranslationCS{obrna}
\dicEntry[lönd] \dicTerm{lönd} \dicIPA{{l}{\oe}{n}{\textsubring{d}}} \dicPos{n} \dicFlx{pl nom} \dicLink{land}
\dicEntry[löndun] \dicTerm{löndun} \dicIPA{{l}{\oe}{n}{\textsubring{d}}{\textscy}{\textsubring{n}}} \dicPos{f}[7] \dicFlx{(‑ar, landanir)}[16] \dicSynonym*{uppskipun} \dicDirectTranslationCS{vykládka, vylodění}
\dicEntry[löng] \dicTerm{löng} \dicIPA{{l}{\oe i}{\ng}{\r{g}}} \dicPos{adj} \dicFlx{f sg nom pos} \dicLink{langur\smash{\textsuperscript{2}}}
\dicEntry[löngu] \dicTerm{löngu\smash{\textsuperscript{1}}} \dicIPA{{l}{\oe i}{\ng}{\r{g}}{\textscy}} \dicPos{f} \dicFlx{sg gen} \dicLink{langa\smash{\textsuperscript{1}}}
\dicEntry[löngu] \dicTerm{löngu\smash{\textsuperscript{2}}} \dicIPA{{l}{\oe i}{\ng}{\r{g}}{\textscy}} \dicPos{adv} \dicDirectTranslationCS{dlouho, dávno};  \dicPhraseIS{endur fyrir löngu} \dicFlx{adv} \dicDirectTranslationCS{dávno, před dávnými dobami};  \dicPhraseIS{fyrir (langa\,/\addthin margt) löngu} \dicFlx{adv} \dicDirectTranslationCS{dávno, před dávnou dobou};  \dicPhraseIS{fyrir lifandi löngu} \dicFlx{adv} \dicDirectTranslationCS{už\,/\addthin již (velmi) dávno};  \dicPhraseIS{löngu áður} \dicFlx{adv} \dicDirectTranslationCS{dlouho předtím};  \dicPhraseIS{löngu fyrr} \dicFlx{adv} \dicDirectTranslationCS{mnohem dříve, dlouho předtím};  \dicPhraseIS{löngu síðar\,/\addthin seinna} \dicFlx{adv} \dicDirectTranslationCS{dlouho poté, o~dost později}
\dicEntry[löngum] \dicTerm{löngum} \dicsymFrequent\  \dicIPA{{l}{\oe i}{\ng}{\r{g}}{\textscy}{\textsubring{m}}} \dicPos{adv} \dicDirectTranslationCS{dlouho} \dicExampleIS{Það hefur löngum verið vitað að sólböð eru hættuleg.} \dicExampleCS{Už se dlouho ví, že opalování je nebezpečné.}
\dicEntry[löngun] \dicTerm{löngun} \dicsymFrequent\  \dicIPA{{l}{\oe i}{\ng}{\r{g}}{\textscy}{\textsubring{n}}} \dicPos{f}[7] \dicFlx{(‑ar, langanir)}[16] \dicDirectTranslationCS{touha, toužení, tužba} \dicExampleIS{löngun í nýjan bíl} \dicExampleCS{touha po novém autě}
\dicEntry[löngutangar] \dicTerm{löngu··tangar} \dicIPA{{l}{\oe i}{\ng}{\r{g}}{\textscy}{t\smash{\textsuperscript{h}}}{au}{\ng}{\r{g}}{a}{\textsubring{r}}} \dicPos{f} \dicFlx{sg gen} \dicLink{langatöng}
\dicEntry[löngutengur] \dicTerm{löngu··tengur} \dicIPA{{l}{\oe i}{\ng}{\r{g}}{\textscy}{t\smash{\textsuperscript{h}}}{\textepsilon}{i}{\ng}{\r{g}}{\textscy}{\textsubring{r}}} \dicPos{f} \dicFlx{pl nom} \dicLink{langatöng}
\dicEntry[löngutöng] \dicTerm{löngu··|töng} \dicIPA{{l}{\oe i}{\ng}{\r{g}}{\textscy}{t\smash{\textsuperscript{h}}}{\oe i}{\ng}{\r{g}}} \dicPos{f}[8] \dicFlx{(‑tangar, ‑tengur\,/\addthin ‑tangir)}[9] \dicLink{langatöng}
\dicEntry[löpp] \dicTerm{löpp} \dicsymFrequent\  \dicIPA{{l}{\oe}{h}{\textsubring{b}}} \dicPos{f}[7] \dicFlx{(lappar, lappir)}[16] \dicDirectTranslationCS{noha, chodidlo} \dicExampleIS{vöðvar í löppum} \dicExampleCS{svaly na nohou};  \dicPhraseIS{koma sér á lappir} \dicDirectTranslationCS{vstát z~postele}
\dicEntry[löptum] \dicTerm{löptum} \dicIPA{{l}{\oe}{f}{\textsubring{d}}{\textscy}{\textsubring{m}}} \dicPos{v} \dicFlx{ind pf pl 1 pers} \dicLink{lepja}
\dicEntry[löstur] \dicTerm{löstur} \dicIPA{{l}{\oe}{s}{\textsubring{d}}{\textscy}{\textsubring{r}}} \dicPos{m}[11] \dicFlx{(lastar, lestir)}[5] \dicDirectTranslationCS{zlozvyk, neřest} \dicExampleIS{drykkjuskapur og aðrir lestir} \dicExampleCS{opilství a~jiné neřesti}
\dicEntry[löt] \dicTerm{löt} \dicIPA{{l}{\oe}{\textlengthmark}{\textsubring{d}}} \dicPos{adj} \dicFlx{f sg nom pos} \dicLink{latur}
\dicEntry[lötra] \dicTerm{lötr|a} \dicIPA{{l}{\oe}{\textlengthmark}{\textsubring{d}}{r}{a}} \dicPos{v}[1] \dicFlx{(‑aði)}[44] \dicDirectTranslationCS{loudat se, šinout se, courat se} \dicExampleIS{Hann lötraði á eftir hópnum.} \dicExampleCS{Coural se za skupinou.}
\dicEntry[löttum] \dicTerm{löttum} \dicIPA{{l}{\oe}{h}{\textsubring{d}}{\textscy}{\textsubring{m}}} \dicPos{v} \dicFlx{ind pf pl 1 pers} \dicLink{letja}
}

\ifPDFtest
\input{letters}
%\input{letters_problematic_3}
\else
\input{letters}
\fi
\restoregeometry
\pagestyle{plain}

\fi

% ================================================================ PHONETICS =

\ifinputphon

\cleardoublepage

\chapter{Seznam islandských fonémů}              \label{sec:phon_phonems}

\begin{table}[h]
\begin{tabular}{lllll} \toprule

\textbf{foném} & \textbf{vysvětlivka} & \textbf{příklad} & \textbf{fonetický přepis} & \textbf{překlad} \\
\midrule
{\textipa{[{a}]}} & otevřená přední nezaokrouhlená samohl. & rata & {\textipa{[{r}{a}{\textlengthmark}{\textsubring{d}}{a}]}} & trefit \\ 
{\textipa{[{\textepsilon}]}} & polootevřená přední nezaokrouhlená samohl. & meta & {\textipa{[{m}{\textepsilon}{\textlengthmark}{\textsubring{d}}{a}]}} & ocenit \\ 
{\textipa{[{e}]}} & polozavřená přední nezaokrouhlená samohl. & heiður & {\textipa{[{h}{ei}{\textlengthmark}{ð}{\textscy}{\textsubring{r}}]}} & čest \\
{\textipa{[{\textsci}]}} & téměř zavřená téměř přední nezaokrouhlená samohl. & vinna & {\textipa{[{v}{\textsci}{n}{\textlengthmark}{a}]}} & pracovat \\ 
{\textipa{[{i}]}} & zavřená přední nezaokrouhlená samohl. & fínn & {\textipa{[{f}{i}{\textsubring{d}}{\textsubring{n}}]}} & jemný \\ 
{\textipa{[{\textopeno}]}} & polootevřená zadní zaokrouhlená samohl. & lofa & {\textipa{[{l}{\textopeno}{\textlengthmark}{v}{a}]}} & slibovat \\ 
{\textipa{[{o}]}} & polozavřená zadní zaokrouhlená samohl. & bóndi & {\textipa{[{\textsubring{b}}{ou}{n}{\textsubring{d}}{\textsci}]}} & farmář \\ 
{\textipa{[{\textscy}]}} & téměř zavřená téměř přední zaokrouhlená samohl. & hundur & {\textipa{[{h}{\textscy}{n}{\textsubring{d}}{\textscy}{\textsubring{r}}]}} & pes \\ 
{\textipa{[{u}]}} & zavřená zadní zaokrouhlená samohl. & súkkulaði & {\textipa{[{s}{u}{h}{\r{g}}{\textscy}{l}{a}{ð}{\textsci}]}} & čokoláda \\ 
{\textipa{[{\oe}]}} & téměř otevřená přední nezaokrouhlená samohl. & lönd & {\textipa{[{l}{\oe}{n}{\textsubring{d}}]}} & země \\ 
{\textipa{[{m}]}} & bilabiální nazála & mýri & {\textipa{[{m}{i}{\textlengthmark}{r}{\textsci}]}} & mokřina \\ 
{\textipa{[{\textsubring{m}}]}} & neznělá bilabiální nazála & þyrmt & {\textipa{[{\texttheta}{\textsci}{\textsubring{r}}{\textsubring{m}}{\textsubring{d}}]}} & {\textit{$\shortrightarrow$ þyrma}} \\  
{\textipa{[{n}]}} & alveolární nazála & nemandi & {\textipa{[{n}{\textepsilon}{\textlengthmark}{m}{a}{n}{\textsubring{d}}{\textsci}]}} & žák \\ 
{\textipa{[{\textsubring{n}}]}} & neznělá alveolární nazála & mennt & {\textipa{[{m}{\textepsilon}{\textsubring{n}}{\textsubring{d}}]}} & vzdělání \\ 
{\textipa{[{\textltailn}]}} & palatální nazála & angi & {\textipa{[{au}{\textltailn}{\r{\textObardotlessj}}{\textsci}]}} & drobek \\ 
{\textipa{[{\r{\textltailn}}]}} & neznělá palatální nazála & banki & {\textipa{[{\textsubring{b}}{au}{\r{\textltailn}}{\r{\textObardotlessj}}{\textsci}]}} & banka \\
{\textipa{[{\ng}]}} & velární nazála & skyggndur & {\textipa{[{s}{\r{\textObardotlessj}}{\textsci}{\ng}{\textsubring{d}}{\textscy}{\textsubring{r}}]}} & {\textit{$\shortrightarrow$ skyggna}} \\ 
{\textipa{[{\r{\ng}}]}} & neznělá velární nazála & skyggnt & {\textipa{[{s}{\r{\textObardotlessj}}{\textsci}{\r{\ng}}{\textsubring{d}}]}} & {\textit{$\shortrightarrow$ skyggna}} \\ 
{\textipa{[{\textsubring{b}}]}} & neznělá bilabiální ploziva & bara & {\textipa{[{\textsubring{b}}{a}{\textlengthmark}{r}{a}]}} & jenom \\ 
{\textipa{[{p\textsuperscript{h}}]}} & neznělá bilabiální ploziva s přídechem & prestur & {\textipa{[{p\textsuperscript{h}}{r}{\textepsilon}{s}{\textsubring{d}}{\textscy}{\textsubring{r}}]}} & kněz \\ 
{\textipa{[{\textsubring{d}}]}} & neznělá alveolární ploziva & dæmi & {\textipa{[{\textsubring{d}}{a}{i}{\textlengthmark}{m}{\textsci}]}} & příklad \\ 
{\textipa{[{t\textsuperscript{h}}]}} & neznělá alveolární ploziva s přídechem & tómur & {\textipa{[{t\textsuperscript{h}}{ou}{\textlengthmark}{m}{\textscy}{\textsubring{r}}]}} & prázdný \\ 
{\textipa{[{\r{\textObardotlessj}}]}} & neznělá palatální ploziva & gefa & {\textipa{[{\r{\textObardotlessj}}{\textepsilon}{\textlengthmark}{v}{a}]}} & dát \\ 
{\textipa{[{c\textsuperscript{h}}]}} & neznělá palatální ploziva s přídechem & kenna & {\textipa{[{c\textsuperscript{h}}{\textepsilon}{n}{\textlengthmark}{a}]}} & učit \\
{\textipa{[{\r{g}}]}} & neznělá velární ploziva & gata & {\textipa{[{\r{g}}{a}{\textlengthmark}{\textsubring{d}}{a}]}} & ulice \\ 
{\textipa{[{k\textsuperscript{h}}]}} & neznělá velární ploziva s přídechem & koma & {\textipa{[{k\textsuperscript{h}}{\textopeno}{\textlengthmark}{m}{a}]}} & přijít \\
{\textipa{[{v}]}} & znělá labiodentální frikativa & hafa & {\textipa{[{h}{a}{\textlengthmark}{v}{a}]}} & mít \\ 
{\textipa{[{f}]}} & neznělá labiodentální frikativa & finna & {\textipa{[{f}{\textsci}{n}{\textlengthmark}{a}]}} & najít \\ 
{\textipa{[{ð}]}} & znělá dentální frikativa & borða & {\textipa{[{\textsubring{b}}{\textopeno}{r}{ð}{a}]}} & jíst \\ 
{\textipa{[{\texttheta}]}} & neznělá dentální frikativa & þola & {\textipa{[{\texttheta}{\textopeno}{\textlengthmark}{l}{a}]}} & snést \\
{\textipa{[{s}]}} & neznělá alveolární frikativa & systir & {\textipa{[{s}{\textsci}{s}{\textsubring{d}}{\textsci}{\textsubring{r}}]}} & sestra \\ 
{\textipa{[{\c{c}}]}} & neznělá palatální frikativa & hjóla & {\textipa{[{\c{c}}{ou}{\textlengthmark}{l}{a}]}} & jet na kole \\ 
{\textipa{[{\textbabygamma}]}} & znělá velární frikativa & sagði & {\textipa{[{s}{a}{\textbabygamma}{ð}{\textsci}]}} & řekl \\ 
{\textipa{[{x}]}} & neznělá velární frikativa & lag & {\textipa{[{l}{a}{\textlengthmark}{x}]}} & píseň \\
{\textipa{[{h}]}} & neznělá glotální frikativa & heima & {\textipa{[{h}{ei}{\textlengthmark}{m}{a}]}} & doma \\ 
{\textipa{[{l}]}} & alveolární laterální aproximanta & loforð & {\textipa{[{l}{\textopeno}{\textlengthmark}{v}{\textopeno}{r}{\texttheta}]}} & slib \\ 
{\textipa{[{\textsubring{l}}]}} & neznělá alveolární laterální aproximanta & hlutur & {\textipa{[{\textsubring{l}}{\textscy}{\textlengthmark}{\textsubring{d}}{\textscy}{\textsubring{r}}]}} & věc \\  
{\textipa{[{j}]}} & palatální aproximanta & jól & {\textipa{[{j}{ou}{\textlengthmark}{\textsubring{l}}]}} & Vánoce \\ 
{\textipa{[{r}]}} & alveolární vibranta & reynsla & {\textipa{[{r}{ei}{n}{s}{\textsubring{d}}{l}{a}]}} & zkušenost \\ 
{\textipa{[{\textsubring{r}}]}} & neznělá alveolární vibranta & hraun & {\textipa{[{\textsubring{r}}{\oe i}{\textlengthmark}{\textsubring{n}}]}} & láva \\ 
\bottomrule
\end{tabular}
\end{table}

\clearpage

%\section{Výslovnost islandských samohlásek}      \label{sec:phon_vowels}
%\LTXtable{\columnwidth}{phon/vowels}
%\clearpage

%\section{Výslovnost islandských souhlásek}       \label{sec:phon_consonants}
%\LTXtable{\columnwidth}{phon/consonants}
%\clearpage
\fi

% =============================================================== MORPHOLOGY =

\ifinputmorpho

\cleardoublepage

\chapter{Morfologický klíč}                      \label{sec:morpho}

\section{Podstatné jméno, rod mužský}            \label{sec:morpho_m}
{\small\LTXfw{morpho/m}}
\clearpage

\section{Podstatné jméno, rod ženský}            \label{sec:morpho_f}
{\small\LTXfw{morpho/f}}
\clearpage

\section{Podstatné jméno, rod střední}           \label{sec:morpho_n}
{\small\LTXfw{morpho/n}}
\clearpage

\section{Přídavné jméno}                         \label{sec:morpho_adj}
{\small\LTXfw{morpho/adj}}
\clearpage

\section{Zájmeno}                                \label{sec:morpho_pron}
{\small\LTXfw{morpho/pron}}
 \begin{center}
\begin{minipage}[t]{.45\textwidth}
\begin{tabular}{l>{\footnotesize\itshape}lll}
 {\textbf{\textit{třída}}} & {\textit{pád}}   & \textit{sg} & \textit{pl}  \\ 
 \hline
\multirow{3}{*}{{{\textbf{pron} \Large{\textbf{17}}}}}  &  nom & ég & við   \\*
 & acc &  mig  & okkur  \\*
 & dat & mér & okkur   \\*
 & gen & mín  & okkar  \\
\hline

\multirow{3}{*}{{{\textbf{pron} \Large{\textbf{17}}}}}  &  nom & hann & þeir   \\*
 & acc &  hann  & þá  \\*
 & dat & honum & þeim   \\*
 & gen & hans  & þeirra  \\
\hline

\multirow{3}{*}{{{\textbf{pron} \Large{\textbf{17}}}}}  &  nom & hún & þær   \\*
 & acc &  hana  & þær  \\*
 & dat & henni & þeim   \\*
 & gen & hennar  & þeirra  \\
\hline

\multirow{3}{*}{{{\textbf{pron} \Large{\textbf{17}}}}}  &  nom &  & vér   \\*
 & acc &    & oss  \\*
 & dat &  & oss   \\*
 & gen &   & vor  \\
\hline

\multirow{3}{*}{{{\textbf{pron} \Large{\textbf{17}}}}}  &  nom &  & þér   \\*
 & acc &    & yður  \\*
 & dat &  & yður   \\*
 & gen &   & yðar  \\
\hline

\multirow{3}{*}{{{\textbf{pron} \Large{\textbf{17}}}}}  &  nom & þú & þið   \\*
 & acc &  þig  & ykkur  \\*
 & dat & þér & ykkur   \\*
 & gen & þín  & ykkar  \\
\hline

\multirow{3}{*}{{{\textbf{pron} \Large{\textbf{17}}}}}  &  nom & það & þau   \\*
 & acc &  það  & þau  \\*
 & dat & því & þeim   \\*
 & gen & þess  & þeirra  \\
\hline
\end{tabular}
\end{minipage}\hfil
\begin{minipage}[t]{.45\textwidth}
\begin{tabular}{l>{\footnotesize\itshape}lll}
\hline
\multirow{3}{*}{{{\textbf{pron} \Large{\textbf{18}}}}}  &  nom &  &    \\*
 & acc &  sig  &   \\*
 & dat & sér &    \\*
 & gen & sín  &   \\
 \hline
\end{tabular}
\end{minipage}
\end{center}

\clearpage

\section{Číslovka}                               \label{sec:morpho_num}
 {\small\LTXfw{morpho/num}}
\clearpage

\begin{landscape}
\section{Sloveso}                                \label{sec:morpho_v}
{\small\LTXfw{morpho/v}}
\end{landscape}

\clearpage

\fi

% ============================================================== BACK MATTER =

\ifmakebackmatter

% ------------------------------------------------ Photographs authors index -

\cleardoublepage
%\pagestyle{empty}
%\pagestyle{indexstyle}
%\phantomsection
%\addcontentsline{toc}{chapter}
%  {Seznam autorů fotografií a ilustrací}        
\label{sec:photo}
\printindex[figures]

% ------------------------------------------------------------- Bibliography -
\cleardoublepage
\chapter{Bibliografie}
%\twocolumn
%\phantomsection
%\addcontentsline{toc}{chapter}{Bibliografie}
\begin{multicols}{2}
\nocite{*}\printbibliography[heading=none]
\end{multicols}

% ------------------------------------------------------------------ Apendix -
\cleardoublepage
\onecolumn
\ifPDF
\else
\chapter{Apendix: Islandská gramatika}
\clearpage
\covergeometry  
\makeatletter
\includepdf[
  pages={1,2,3,4,5,6,7,8,9,10,11,12,13,14,15,16,17,18,19,20,21,22,23,24,25,26,27,28,29,30,31,32,33,34,35,36,37,38,39,40,41,42,43,44,45,46,47,48,49,50,51,52,53,54,55,56,57,58,59,60,61,62,63,64,65,66,67,68,69,70,71,72,73},
  turn=false,
  scale=0.79,
  keepaspectratio,
  offset={\dimexpr(\paperwidth-\csname Gm@layoutwidth\endcsname)/2-\csname Gm@layouthoffset\endcsname\relax}
  {\dimexpr(\paperheight-\csname  Gm@layoutheight\endcsname)/2-\csname Gm@layoutvoffset\endcsname\relax}
   ]{gram.pdf}
\makeatother
\restoregeometry  
\fi

% ------------------------------------------------------------------ Barcode -
\pagenumbering{gobble}
\ifPDF
\else
\cleardoublepage\null\clearpage % clear to verso
\vspace*\fill
\EANisbn[SC3]
\fi
\fi

% --------------------------------------------------------------- Back cover -

\ifmakecovers
  \cleardoublepage\null\clearpage % inner side of the cover
\covergeometry  
  \makecoverwith\backcoverimages
\restoregeometry  
\fi

\end{document}

